\documentclass[oneside,10pt]{amsart}
\PassOptionsToPackage{hyperfootnotes=false}{hyperref}
\ifdefined\pdfvariable\pdfvariable suppressoptionalinfo 767\fi
\usepackage[no-math]{fontspec}
\setmainfont{Noto Serif}
\usepackage{amsmath,amssymb,amsthm}
\usepackage[all]{xy}
\usepackage{xstring,xparse,xcolor}
\definecolor{brightmaroon}{rgb}{0.76,0.13,0.28}
\usepackage[linktocpage=true,colorlinks=true,hyperindex,hypertexnames=false,citecolor=blue,linkcolor=brightmaroon]{hyperref}
\usepackage[left=30mm,right=30mm,top=22mm,bottom=25mm]{geometry}
\usepackage{marginnote}
\usepackage[shortlabels]{enumitem}
\usepackage{longtable,multicol,microtype}
\usepackage[brazilian]{babel}
\setlength{\emergencystretch}{5em}
\raggedbottom

\hypersetup{
  pdftitle={Elementos de geometria algébrica, EGA II — edição brasileira},
  pdfauthor={Alexander Grothendieck and Jean Dieudonné; AI translation and edition by OpenAI Codex},
  pdfsubject={EGA II em português brasileiro}
}

\newtheoremstyle{ega-env-style}{}{}{\rmfamily}{}{\bfseries}{.}{ }{\thmnote{(#3)}}
\newtheoremstyle{ega-thm-env-style}{}{}{\itshape}{}{\bfseries}{. --- }{ }{\thmname{#1}\thmnote{ (#3)}}
\newtheoremstyle{ega-defn-env-style}{}{}{\rmfamily}{}{\bfseries}{. --- }{ }{\thmname{#1}\thmnote{ (#3)}}
\theoremstyle{ega-env-style}
\newtheorem*{env}{---}
\theoremstyle{ega-thm-env-style}
\newtheorem*{theorem}{Teorema}
\newtheorem*{proposition}{Proposição}
\newtheorem*{lemma}{Lema}
\newtheorem*{corollary}{Corolário}
\theoremstyle{ega-defn-env-style}
\newtheorem*{definition}{Definição}
\newtheorem*{example}{Exemplo}
\newtheorem*{remark}{Observação}
\newtheorem*{remarks}{Observações}
\newtheorem*{notation}{Notação}
\renewcommand{\proofname}{Demonstração}

\def\sh{\mathcal}
\def\bb{\mathbf}
\def\cat{\mathtt}
\def\leq{\leqslant}
\def\geq{\geqslant}
\def\setmin{-}
\def\rad{\mathfrak{r}}
\def\nilrad{\mathfrak{N}}
\def\emp{\varnothing}
\def\vphi{\phi}
\def\HH{\mathrm{H}}
\def\CHH{\check{\HH}}
\def\RR{\mathrm{R}}
\def\LL{\mathrm{L}}
\def\dual#1{{#1}^\vee}
\def\kres{k}
\def\C{\cat{C}}
\def\op{^\cat{op}}
\def\Set{\cat{Set}}
\def\CHom{\cat{Hom}}
\def\supertilde{{\,\widetilde{\,}\,}}
\def\GL{\bb{GL}}
\def\bbGamma{\boldsymbol\Gamma}
\def\red{\mathrm{red}}
\def\rg{\operatorname{rg}}
\def\gr{\operatorname{gr}}
\def\Gr{\operatorname{Gr}}
\def\Hom{\operatorname{Hom}}
\def\Proj{\operatorname{Proj}}
\def\Tor{\operatorname{Tor}}
\def\Ext{\operatorname{Ext}}
\def\Supp{\operatorname{Supp}}
\def\Ker{\operatorname{Ker}}
\def\Im{\operatorname{Im}}
\def\Coker{\operatorname{Coker}}
\def\Spec{\operatorname{Spec}}
\def\Frac{\operatorname{Frac}}
\def\Spm{\operatorname{Spm}}
\def\Spf{\operatorname{Spf}}
\def\grad{\operatorname{grad}}
\def\dimc{\operatorname{dimc}}
\def\codim{\operatorname{codim}}
\def\id{\operatorname{id}}
\def\Der{\operatorname{Der}}
\def\Diff{\operatorname{Diff}}
\def\Hyp{\operatorname{Hyp}}
\def\Tr{\operatorname{Tr}}
\def\shHom{\mathcal{H}\!\mathit{om}}
\def\shProj{\mathcal{P}\!\mathit{roj}}
\def\shAut{\mathcal{A}\!\mathit{ut}}
\renewcommand{\to}{\mathchoice{\longrightarrow}{\rightarrow}{\rightarrow}{\rightarrow}}
\newcommand{\from}{\mathchoice{\longleftarrow}{\leftarrow}{\leftarrow}{\leftarrow}}
\let\mapstoo\mapsto
\renewcommand{\mapsto}{\mathchoice{\longmapsto}{\mapstoo}{\mapstoo}{\mapstoo}}
\def\isoto{\simeq}
\def\simto{\xrightarrow{\sim}}
\renewcommand{\hat}[1]{\widehat{#1}}
\renewcommand{\tilde}[1]{\widetilde{#1}}
\def\sectionbreak{\begin{center}***\end{center}}
\newcommand{\oldpage}[2][]{\marginnote{\normalfont\footnotesize\textbf{#1}~|~#2}\ignorespaces}

\let\EGAOriginalLabel\label
\renewcommand{\label}[1]{\phantomsection\EGAOriginalLabel{#1}}
\makeatletter
\def\EGARefNumberAux#1.#2\@nil{#2}
\newcommand{\EGARefNumber}[1]{\EGARefNumberAux#1\@nil}
\newcommand{\EGAMaybeRefLink}[2]{%
  \ifcsname r@#1\endcsname\hyperref[#1]{#2}\else#2\fi}
\NewDocumentCommand{\sref}{O{} m O{}}{%
  {\normalfont(\IfBlankF{#1}{\textbf{#1},~}%
  \EGAMaybeRefLink{#2}{\EGARefNumber{#2}}%
  \IfBlankF{#3}{,~#3})}}
\makeatother
\newcommand{\eref}[1]{\hyperref[#1]{\normalfont{(\StrBehind{#1}{.})}}}


\newcommand{\EGAResetVolumeCounters}{%
  \setcounter{section}{0}\setcounter{subsection}{0}\setcounter{equation}{0}%
  \setcounter{figure}{0}\setcounter{table}{0}\setcounter{footnote}{0}}
\newcommand{\EGAVolumeTitlePage}[3]{%
  \clearpage\EGAResetVolumeCounters\phantomsection\label{section:#1}%
  \pdfbookmark[0]{#2: #3}{volume.#1}\thispagestyle{empty}\vspace*{\fill}%
  \begin{center}
    {\Huge\bfseries #2\par}\vspace{1.5em}{\LARGE #3\par}\vspace{2em}
    {\large Alexander Grothendieck e Jean Dieudonné\par}\vspace{1.3em}
    {\large\bfseries Edição brasileira atual\par}\vspace{1.3em}
  \end{center}
  \begin{minipage}{0.88\textwidth}\small
    Tradução, edição e composição produzidas por OpenAI Codex — GPT-5.6 Sol, esforço Ultra. Não se atribui revisão humana.\par\medskip
    Fonte: Alexander Grothendieck e Jean Dieudonné, Publications Mathématiques de l'IHÉS 8 (1961), pp. 5–222; testemunho digital NUMDAM PMIHES\_1961\_\_8\_\_5\_0.\par\medskip
    Edição corrigida.\par\medskip
    Esta edição independente preserva a autoria e a relação com a fonte histórica; não implica endosso dos autores, de seus sucessores nem do NUMDAM.
  \end{minipage}
  \vspace*{\fill}\clearpage}

\begin{document}
\EGAVolumeTitlePage{ega2}{EGA II}{Estudo global elementar de algumas classes de morfismos}

\begingroup\small
\noindent\textbf{Fontes, direitos e autoria}

A obra original é de Alexandre Grothendieck, com a colaboração de Jean Dieudonné: EGA II, Publications Mathématiques de l’IHÉS 8 (1961), pp. 5–222, \href{https://numdam.org/item/PMIHES\_1961\_\_8\_\_5\_0/}{NUMDAM}. A referência inglesa e a estrutura TeX utilizadas procedem da tradução humana de \href{https://github.com/ryankeleti/ega/tree/3d7afe28c99d1887ccdb85acf2ca1afd93fb1971}{ryankeleti/ega}, versão 3d7afe28c99d1887ccdb85acf2ca1afd93fb1971. Os colaboradores nomeados são Mateo Carmona, Tim Holzschuh, Tim Hosgood, Ryan Keleti, Thiago Solovera \& Nery, Wayne Yang. Os \href{https://github.com/ryankeleti/ega/blob/3d7afe28c99d1887ccdb85acf2ca1afd93fb1971/README.md}{créditos por seção} especificam suas contribuições. A tradução inglesa herdada mantém essa autoria; a tradução brasileira e suas adições produzidas por IA não são atribuídas a esses colaboradores. A atribuição não constitui autorização ou endosso.

A tradução para o português é produzida por IA. As correções, a edição e a composição também são produzidas por IA. O responsável pela publicação não traduziu o texto. A autoria e as atribuições das obras originais e das contribuições herdadas são preservadas; não são atribuídas à IA. As correções documentadas desta edição incluem OpenAI Codex, GPT-6 Astra, esforço Ultra.
\par\endgroup\clearpage
\begingroup\small
Os direitos sobre as obras originais e sobre as contribuições protegidas a estas traduções e transcrições derivadas permanecem com seus respectivos titulares. Na medida em que o editor detenha direitos sobre a tradução, a transcrição, as correções e a composição produzidas por IA, ele os dedica ao domínio público sob CC0 1.0 Universal (https://creativecommons.org/publicdomain/zero/1.0/deed.pt-br); isso não afeta os direitos sobre as obras originais nem concede permissão para usá-las. Os avisos de direitos de terceiros e as autorizações já concedidas continuam aplicáveis.

Estas edições são disponibilizadas gratuitamente, para fins não comerciais e exclusivamente educativos, sob a presunção de que os respectivos titulares dos direitos não se oporiam à disponibilidade gratuita destas traduções e transcrições, ao lado dos originais franceses no \href{https://www.numdam.org/}{NUMDAM}. Essa presunção não significa que uma autorização tenha sido obtida; o acesso gratuito aos originais não autoriza, por si só, sua tradução, adaptação ou republicação.

Se você detém direitos sobre algum material apresentado aqui e se opõe à sua disponibilidade, poderá \href{https://github.com/KokunoYumeto/ega-fga-sga-pt-br/issues}{abrir uma questão no GitHub}. O responsável pela publicação retirará o material afetado a pedido do titular dos direitos.

\par\endgroup\clearpage

\section*{Sumário}

\begin{longtable}{@{}p{0.08\textwidth}p{0.86\textwidth}@{}}
  \hyperref[section:II.1]{\textsection1}. & Morfismos afins.\\
  \hyperref[section:II.2]{\textsection2}. & Espectros primos homogêneos.\\
  \hyperref[section:II.3]{\textsection3}. & Espectro homogêneo de um feixe de álgebras graduadas.\\
  \hyperref[section:II.4]{\textsection4}. & Fibrados projetivos; feixes amplos.\\
  \hyperref[section:II.5]{\textsection5}. & Morfismos quase-afins; morfismos quase-projetivos; morfismos próprios; morfismos projetivos.\\
  \hyperref[section:II.6]{\textsection6}. & Morfismos inteiros e morfismos finitos.\\
  \hyperref[section:II.7]{\textsection7}. & Critérios valorativos.\\
  \hyperref[section:II.8]{\textsection8}. & Esquemas de explosão; cones projetantes; fecho projetivo.
\end{longtable}
\bigskip

\oldpage[II]{5}
As diversas classes de morfismos estudadas neste capítulo são tratadas sem grande recurso aos métodos coomológicos; um estudo mais aprofundado que empregue esses métodos será feito no Capítulo~III, no qual utilizaremos sobretudo as seções~\textsection\textsection\hyperref[section:II.2]{2}, \hyperref[section:II.4]{4} e \hyperref[section:II.5]{5} do Capítulo~II.
A seção~\textsection\hyperref[section:II.8]{8} pode ser omitida numa primeira leitura: ela fornece alguns complementos ao formalismo desenvolvido nas seções~\textsection\textsection\hyperref[section:II.1]{1} a \hyperref[section:II.3]{3}, que se reduzem a aplicações fáceis desse formalismo, e dela faremos uso menos constante que dos demais resultados deste capítulo.
\bigskip

\section{Morfismos afins}
\label{section:II.1}

A maior parte dos resultados desta seção consiste nas contrapartidas ``globais'' dos resultados do Capítulo~I, \textsection1; eles não são, portanto, essencialmente novos e fornecem simplesmente uma linguagem cômoda para o que segue.
\bigskip

\subsection{$S$-pré-esquemas e $\mathcal{O}_S$-álgebras}
\label{subsection:II.1.1}

\begin{env}[1.1.1]
\label{II.1.1.1}
Sejam $S$ um pré-esquema, $X$ um $S$-pré-esquema e $f:X\to S$ seu morfismo estrutural.
Sabe-se por \sref[0]{0.4.2.4} que a imagem direta $f_*(\sh{O}_X)$ é uma $\sh{O}_S$-álgebra, que
\oldpage[II]{6}
denotaremos por $\sh{A}(X)$ quando disso não resultar confusão; se $U$ é um aberto de $S$, tem-se
\[
  \sh{A}(f^{-1}(U))=\sh{A}(X)|U.
\]
Do mesmo modo, para todo $\sh{O}_X$-módulo $\sh{F}$ (resp. toda $\sh{O}_X$-álgebra $\sh{B}$), denotaremos por $\sh{A}(\sh{F})$ (resp. $\sh{A}(\sh{B})$) a imagem direta $f_*(\sh{F})$ (resp. $f_*(\sh{B})$), que é um $\sh{A}(X)$-módulo (resp. uma $\sh{A}(X)$-álgebra), e não apenas um $\sh{O}_S$-módulo (resp. uma $\sh{O}_S$-álgebra).
\end{env}

\begin{env}[1.1.2]
\label{II.1.1.2}
Sejam $Y$ um segundo $S$-pré-esquema, $g:Y\to S$ seu morfismo estrutural e $h:X\to Y$ um $S$-morfismo; tem-se, portanto, o diagrama comutativo
\[
\label{II.1.1.2.1}
  \xymatrix{
    X\ar[rr]^h\ar[rd]_f & &
    Y\ar[ld]^g\\
    & S.
  }
  \tag{1.1.2.1}
\]

Por definição, $h=(\psi,\theta)$, onde $\theta:\sh{O}_Y\to h_*(\sh{O}_X)=\psi_*(\sh{O}_X)$ é um homomorfismo de feixes de anéis; deduz-se então de \sref[0]{0.4.2.2} um homomorfismo de $\sh{O}_S$-álgebras $g_*(\theta):g_*(\sh{O}_Y)\to g_*(h_*(\sh{O}_X))=f_*(\sh{O}_X)$, isto é, um homomorfismo de $\sh{O}_S$-álgebras $\sh{A}(Y)\to\sh{A}(X)$, que denotaremos também por $\sh{A}(h)$.
Se $h':Y\to Z$ é um segundo $S$-morfismo, é imediato que $\sh{A}(h'\circ h)=\sh{A}(h)\circ\sh{A}(h')$.
Definimos, assim, um \emph{funtor contravariante $X\mapsto\sh{A}(X)$} na categoria dos $S$-pré-esquemas, com valores na categoria das $\sh{O}_S$-álgebras.

Sejam agora $\sh{F}$ um $\sh{O}_X$-módulo, $\sh{G}$ um $\sh{O}_Y$-módulo e $u:\sh{G}\to\sh{F}$ um $h$-morfismo, isto é, por \sref[0]{0.4.4.1}, um homomorfismo de $\sh{O}_Y$-módulos $\sh{G}\to h_*(\sh{F})$.
Então $g_*(u):g_*(\sh{G})\to g_*(h_*(\sh{F}))=f_*(\sh{F})$ é um homomorfismo $\sh{A}(\sh{G})\to\sh{A}(\sh{F})$ de $\sh{O}_S$-módulos, que denotaremos por $\sh{A}(u)$; além disso, o par $(\sh{A}(h),\sh{A}(u))$ constitui um \emph{di-homomorfismo} do $\sh{A}(Y)$-módulo $\sh{A}(\sh{G})$ no $\sh{A}(X)$-módulo $\sh{A}(\sh{F})$.
\end{env}

\begin{env}[1.1.3]
\label{II.1.1.3}
Fixado o pré-esquema $S$, podem-se considerar os pares $(X,\sh{F})$, onde $X$ é um $S$-pré-esquema e $\sh{F}$ é um $\sh{O}_X$-módulo, como formando uma \emph{categoria}, definindo-se um \emph{morfismo} $(X,\sh{F})\to(Y,\sh{G})$ como um par $(h,u)$, onde $h:X\to Y$ é um $S$-morfismo e $u:\sh{G}\to\sh{F}$ é um $h$-morfismo.
Pode-se então dizer que $(X,\sh{F})\mapsto(\sh{A}(X),\sh{A}(\sh{F}))$ é um \emph{funtor contravariante} com valores na categoria cujos objetos são os pares formados por uma $\sh{O}_S$-álgebra e um módulo sobre essa álgebra, e cujos morfismos são os di-homomorfismos.
\end{env}

\subsection{Pré-esquemas afins sobre um pré-esquema}
\label{subsection:II.1.2}

\begin{definition}[1.2.1]
\label{II.1.2.1}
Sejam $X$ um $S$-pré-esquema e $f:X\to S$ seu morfismo estrutural.
Diz-se que $X$ é \emph{afim sobre $S$} se existe uma cobertura $(S_\alpha)$ de $S$ por abertos afins tal que, para todo $\alpha$, o pré-esquema induzido por $X$ no aberto $f^{-1}(S_\alpha)$ seja afim.
\end{definition}

\begin{example}[1.2.2]
\label{II.1.2.2}
Todo subpré-esquema fechado de $S$ é um $S$-pré-esquema afim sobre $S$ (\sref[I]{I.4.2.3} e \sref[I]{I.4.2.4}).
\end{example}

\begin{remark}[1.2.3]
\label{II.1.2.3}
Um pré-esquema $X$ afim sobre $S$ não é necessariamente um esquema afim, como mostra o exemplo $X=S$ \sref{II.1.2.2}.
Por outro lado, se um esquema afim $X$ é um $S$-pré-esquema, então $X$ não é necessariamente afim sobre
\oldpage[II]{7}
$S$ (ver o Exemplo~\sref{II.1.3.3}).
Contudo, recordemos que, se $S$ é um \emph{esquema}, todo $S$-pré-esquema que seja um esquema afim é afim sobre $S$ \sref[I]{I.5.5.10}.
\end{remark}

\begin{proposition}[1.2.4]
\label{II.1.2.4}
Todo $S$-pré-esquema que seja afim sobre $S$ é separado sobre $S$ (em outras palavras, é um $S$-esquema).
\end{proposition}

\begin{proof}
Isto decorre imediatamente de \sref[I]{I.5.5.5} e \sref[I]{I.5.5.8}.
\end{proof}

\begin{proposition}[1.2.5]
\label{II.1.2.5}
Sejam $X$ um $S$-esquema afim sobre $S$ e $f:X\to S$ seu morfismo estrutural.
Para todo aberto $U\subset S$, $f^{-1}(U)$ é afim sobre $U$.
\end{proposition}

\begin{proof}
Em virtude da Definição~\sref{II.1.2.1}, podemos reduzir-nos ao caso em que $S=\Spec(A)$ e $X=\Spec(B)$ são afins; então $f=({}^a\vphi,\widetilde{\vphi})$, onde $\vphi:A\to B$ é um homomorfismo.
Como os $D(g)$, para $g\in A$, formam uma base de $S$, podemos reduzir-nos ao caso em que $U=D(g)$; mas sabe-se então que $f^{-1}(U)=D(\vphi(g))$ (\textbf{I},~1.2.2.2), de onde resulta a proposição.
\end{proof}

\begin{proposition}[1.2.6]
\label{II.1.2.6}
Sejam $X$ um $S$-esquema afim sobre $S$ e $f:X\to S$ seu morfismo estrutural.
Para todo $\sh{O}_X$-módulo quase-coerente $\sh{F}$, $f_*(\sh{F})$ é um $\sh{O}_S$-módulo quase-coerente.
\end{proposition}

\begin{proof}
Levando em conta a Proposição~\sref{II.1.2.4}, isto decorre de \sref[I]{I.9.2.2}[a].
\end{proof}

Em particular, a $\sh{O}_S$-álgebra $\sh{A}(X)=f_*(\sh{O}_X)$ é \emph{quase-coerente}.

\begin{proposition}[1.2.7]
\label{II.1.2.7}
Seja $X$ um $S$-esquema afim sobre $S$.
Para todo $S$-pré-esquema $Y$, a aplicação $h\mapsto\sh{A}(h)$ do conjunto $\Hom_S(Y,X)$ no conjunto $\Hom(\sh{A}(X),\sh{A}(Y))$ \sref{II.1.1.2} é bijetiva.
\end{proposition}

\begin{proof}
Sejam $f:X\to S$ e $g:Y\to S$ os morfismos estruturais.
Suponhamos primeiro que $S=\Spec(A)$ e $X=\Spec(B)$ sejam afins; é preciso provar que, para todo homomorfismo $\omega:f_*(\sh{O}_X)\to g_*(\sh{O}_Y)$ de $\sh{O}_S$-álgebras, existe um único $S$-morfismo $h:Y\to X$ tal que $\sh{A}(h)=\omega$.
Por definição, para todo aberto $U\subset S$, $\omega$ define um homomorfismo $\omega_U=\Gamma(U,\omega):\Gamma(f^{-1}(U),\sh{O}_X)\to\Gamma(g^{-1}(U),\sh{O}_Y)$ de $\Gamma(U,\sh{O}_S)$-álgebras.
Em particular, para $U=S$, isto fornece um homomorfismo $\vphi:\Gamma(X,\sh{O}_X)\to\Gamma(Y,\sh{O}_Y)$ de $\Gamma(S,\sh{O}_S)$-álgebras, ao qual corresponde um $S$-morfismo bem determinado $h:Y\to X$, pois $X$ é afim \sref[I]{I.2.2.4}.
Resta provar que $\sh{A}(h)=\omega$, ou, em outras palavras, que, para todo aberto $U$ de uma base de $S$, $\omega_U$ coincide com o homomorfismo de álgebras $\vphi_U$ correspondente ao $S$-morfismo $g^{-1}(U)\to f^{-1}(U)$ que é restrição de $h$.
Podemos reduzir-nos ao caso em que $U=D(\lambda)$, com $\lambda\in A$; então, se $f=({}^a\rho,\widetilde{\rho})$, onde $\rho:A\to B$ é um homomorfismo de anéis, tem-se $f^{-1}(U)=D(\mu)$, onde $\mu=\rho(\lambda)$, e $\Gamma(f^{-1}(U),\sh{O}_X)$ é o anel de frações $B_\mu$; o diagrama
\[
  \xymatrix{
    B\ar[r]^\vphi\ar[d] &
    \Gamma(Y,\sh{O}_Y)\ar[d]\\
    B_\mu\ar[r]_{\vphi_U} &
    \Gamma(g^{-1}(U),\sh{O}_Y)
  }
\]
é comutativo, assim como o diagrama análogo no qual $\vphi_U$ é substituído por $\omega_U$; a igualdade $\vphi_U=\omega_U$ decorre então da propriedade universal dos anéis de frações \sref[0]{0.1.2.4}.

Passemos agora ao caso geral; seja $(S_\alpha)$ uma cobertura de $S$ por abertos afins
\oldpage[II]{8}
tal que os $f^{-1}(S_\alpha)$ sejam afins.
Então todo homomorfismo $\omega:\sh{A}(X)\to\sh{A}(Y)$ de $\sh{O}_S$-álgebras fornece, por restrição, uma família de homomorfismos
\[
  \omega_\alpha:\sh{A}(f^{-1}(S_\alpha))\to\sh{A}(g^{-1}(S_\alpha))
\]
de $\sh{O}_{S_\alpha}$-álgebras, portanto uma família de $S_\alpha$-morfismos $h_\alpha:g^{-1}(S_\alpha)\to f^{-1}(S_\alpha)$ pelo que precede.
Resta ver que, para todo aberto afim $U$ de uma base de $S_\alpha\cap S_\beta$, as restrições de $h_\alpha$ e $h_\beta$ a $g^{-1}(U)$ coincidem, o que é evidente, pois, pelo que precede, ambas as restrições correspondem ao homomorfismo $\sh{A}(X)|U\to\sh{A}(Y)|U$ que é restrição de $\omega$.
\end{proof}

\begin{corollary}[1.2.8]
\label{II.1.2.8}
Sejam $X$ e $Y$ dois $S$-esquemas afins sobre $S$.
Para que um $S$-morfismo $h:Y\to X$ seja um isomorfismo, é necessário e suficiente que $\sh{A}(h):\sh{A}(X)\to\sh{A}(Y)$ seja um isomorfismo.
\end{corollary}

\begin{proof}
Isto decorre imediatamente da Proposição~\sref{II.1.2.7} e do caráter funtorial de $\sh{A}(X)$.
\end{proof}

\subsection{Pré-esquema afim sobre $S$ associado a uma $\mathcal{O}_S$-álgebra}
\label{subsection:II.1.3}
\label{II.1.3}

\begin{proposition}[1.3.1]
\label{II.1.3.1}
Seja $S$ um pré-esquema.
Para toda $\sh{O}_S$-álgebra quase-coerente $\sh{B}$, existe um pré-esquema $X$ afim sobre $S$, determinado a menos de um $S$-isomorfismo único, tal que $\sh{A}(X)=\sh{B}$.
\end{proposition}

\begin{proof}
A unicidade decorre do Corolário~\sref{II.1.2.8}; demonstremos a existência de $X$.
Para todo aberto afim $U\subset S$, seja $X_U$ o pré-esquema $\Spec(\Gamma(U,\sh{B}))$; como $\Gamma(U,\sh{B})$ é uma $\Gamma(U,\sh{O}_S)$-álgebra, $X_U$ é um $S$-pré-esquema \sref[I]{I.1.6.1}.
Além disso, como $\sh{B}$ é quase-coerente, a $\sh{O}_S$-álgebra $\sh{A}(X_U)$ identifica-se canonicamente com $\sh{B}|U$ (\sref[I]{I.1.3.7}, \sref[I]{I.1.3.13} e \sref[I]{I.1.6.3}).
Seja $V$ um segundo aberto afim de $S$, e seja $X_{U,V}$ o pré-esquema induzido por $X_U$ em $f_U^{-1}(U\cap V)$, onde $f_U$ denota o morfismo estrutural $X_U\to S$; $X_{U,V}$ e $X_{V,U}$ são afins sobre $U\cap V$ \sref{II.1.2.5}, e, por definição, $\sh{A}(X_{U,V})$ e $\sh{A}(X_{V,U})$ identificam-se canonicamente com $\sh{B}|(U\cap V)$.
Existe, portanto, por \sref{II.1.2.8}, um $S$-isomorfismo canônico $\theta_{U,V}:X_{V,U}\to X_{U,V}$; além disso, se $W$ é um terceiro aberto afim de $S$, e se $\theta_{U,V}'$, $\theta_{V,W}'$ e $\theta_{U,W}'$ são as restrições de $\theta_{U,V}$, $\theta_{V,W}$ e $\theta_{U,W}$ às imagens inversas de $U\cap V\cap W$ em $X_V$, $X_W$ e $X_W$, respectivamente, pelos morfismos estruturais, então $\theta_{U,V}'\circ\theta_{V,W}'=\theta_{U,W}'$.
Por conseguinte, existe um pré-esquema $X$, uma cobertura $(T_U)$ de $X$ por abertos afins e, para cada $U$, um isomorfismo $\vphi_U:X_U\to T_U$, tais que $\vphi_U$ envia $f_U^{-1}(U\cap V)$ sobre $T_U\cap T_V$ e $\theta_{U,V}=\vphi_U^{-1}\circ\vphi_V$ \sref[I]{I.2.3.1}.
O morfismo $g_U=f_U\circ\vphi_U^{-1}$ torna $T_U$ um $S$-pré-esquema, e os morfismos $g_U$ e $g_V$ coincidem em $T_U\cap T_V$; logo, $X$ é um $S$-pré-esquema.
Além disso, é claro por definição que $X$ é afim sobre $S$ e que $\sh{A}(T_U)=\sh{B}|U$; portanto, $\sh{A}(X)=\sh{B}$.
\end{proof}

Diz-se que o $S$-esquema $X$ assim definido é \emph{associado à $\sh{O}_S$-álgebra $\sh{B}$}, ou que é o \emph{espectro de $\sh{B}$}, e denota-se por $\Spec(\sh{B})$.

\begin{corollary}[1.3.2]
\label{II.1.3.2}
Sejam $X$ um pré-esquema afim sobre $S$ e $f:X\to S$ seu morfismo estrutural.
Para todo aberto afim $U\subset S$, o pré-esquema induzido em $f^{-1}(U)$ é o esquema afim de anel $\Gamma(U,\sh{A}(X))$.
\end{corollary}

\begin{proof}
\oldpage[II]{9}
Como se pode supor que $X$ esteja associado a uma $\sh{O}_S$-álgebra, em virtude das Proposições~\sref{II.1.2.6} e \sref{II.1.3.1}, o corolário decorre da construção de $X$ descrita na Proposição~\sref{II.1.3.1}.
\end{proof}

\begin{example}[1.3.3]
\label{II.1.3.3}
Seja $S$ o plano afim sobre um corpo $K$, no qual o ponto $0$ foi duplicado \sref[I]{I.5.5.11}; com as notações de \sref[I]{I.5.5.11}, $S$ é a união de dois abertos afins $Y_1$ e $Y_2$; se $f$ é a imersão aberta $Y_1\to S$, então $f^{-1}(Y_2)=Y_1\cap Y_2$ não é um aberto afim de $Y_1$ (\emph{loc. cit.}), o que fornece um exemplo de esquema afim que não é afim sobre $S$.
\end{example}

\begin{corollary}[1.3.4]
\label{II.1.3.4}
Seja $S$ um esquema afim; para que um $S$-pré-esquema $X$ seja afim sobre $S$, é necessário e suficiente que $X$ seja um esquema afim.
\end{corollary}

\begin{corollary}[1.3.5]
\label{II.1.3.5}
Sejam $X$ um pré-esquema afim sobre um pré-esquema $S$ e $Y$ um $X$-pré-esquema.
Para que $Y$ seja afim sobre $S$, é necessário e suficiente que $Y$ seja afim sobre $X$.
\end{corollary}

\begin{proof}
Reduzimo-nos imediatamente ao caso em que $S$ é um esquema afim, e então $X$ também é um esquema afim \sref{II.1.3.4}; as duas condições do enunciado exprimem que $Y$ é um esquema afim \sref{II.1.3.4}.
\end{proof}

\begin{env}[1.3.6]
\label{II.1.3.6}
Seja $X$ um pré-esquema afim sobre $S$.
Definir um pré-esquema $Y$ afim \emph{sobre $X$} equivale, em virtude do Corolário~\sref{II.1.3.5}, a dar um pré-esquema $Y$ afim \emph{sobre $S$} e um $S$-morfismo $g:Y\to X$; em outras palavras (Proposição~\sref{II.1.3.1} e \sref{II.1.2.7}), equivale a dar uma $\sh{O}_S$-álgebra quase-coerente $\sh{B}$ e um homomorfismo $\sh{A}(X)\to\sh{B}$ de $\sh{O}_S$-álgebras (que se pode considerar como definindo em $\sh{B}$ uma estrutura de $\sh{A}(X)$-álgebra).
Se $f:X\to S$ é o morfismo estrutural, tem-se então $\sh{B}=f_*(g_*(\sh{O}_Y))$.
\end{env}

\begin{corollary}[1.3.7]
\label{II.1.3.7}
Seja $X$ um pré-esquema afim sobre $S$; para que $X$ seja de tipo finito sobre $S$, é necessário e suficiente que a $\sh{O}_S$-álgebra quase-coerente $\sh{A}(X)$ seja de tipo finito \sref[I]{I.9.6.2}.
\end{corollary}

\begin{proof}
Por definição \sref[I]{I.9.6.2}, podemos reduzir-nos ao caso em que $S$ é afim; então $X$ é um esquema afim \sref{II.1.3.4}, e, se $S=\Spec(A)$ e $X=\Spec(B)$, $\sh{A}(X)$ é a $\sh{O}_S$-álgebra $\widetilde{B}$; como $\Gamma(S,\widetilde{B})=B$, o corolário decorre de \sref[I]{I.9.6.2} e \sref[I]{I.6.3.3}.
\end{proof}

\begin{corollary}[1.3.8]
\label{II.1.3.8}
Seja $X$ um pré-esquema afim sobre $S$; para que $X$ seja reduzido, é necessário e suficiente que a $\sh{O}_S$-álgebra quase-coerente $\sh{A}(X)$ seja reduzida \sref[0]{0.4.1.4}.
\end{corollary}

\begin{proof}
A questão é local sobre $S$; em virtude do Corolário~\sref{II.1.3.2}, o resultado decorre de \sref[I]{I.5.1.1} e \sref[I]{I.5.1.4}.
\end{proof}

\subsection{Feixes quase-coerentes sobre um pré-esquema afim sobre $S$}
\label{subsection:II.1.4}

\begin{proposition}[1.4.1]
\label{II.1.4.1}
Sejam $X$ um pré-esquema afim sobre $S$, $Y$ um $S$-pré-esquema e $\sh{F}$ (resp. $\sh{G}$) um $\sh{O}_X$-módulo (resp. um $\sh{O}_Y$-módulo) quase-coerente.
Então a aplicação $(h,u)\mapsto(\sh{A}(h),\sh{A}(u))$ do conjunto dos morfismos $(Y,\sh{G})\to(X,\sh{F})$ no conjunto dos di-homomorfismos $(\sh{A}(X),\sh{A}(\sh{F}))\to(\sh{A}(Y),\sh{A}(\sh{G}))$ (\sref{II.1.1.2} e \sref{II.1.1.3}) é bijetiva.
\end{proposition}

\begin{proof}
A demonstração segue exatamente o mesmo procedimento que a da Proposição~\sref{II.1.2.7}, usando \sref[I]{I.2.2.5} e \sref[I]{I.2.2.4}; os detalhes são deixados ao leitor.
\end{proof}

\begin{corollary}[1.4.2]
\label{II.1.4.2}
Se, além das hipóteses da Proposição~\sref{II.1.4.1}, supusermos que $Y$ é afim sobre $S$, então, para que $(h,u)$ seja um isomorfismo, é necessário e suficiente que $(\sh{A}(h),\sh{A}(u))$ seja um di-isomorfismo.
\end{corollary}
\oldpage[II]{10}
\begin{proposition}[1.4.3]
\label{II.1.4.3}
Para todo par $(\sh{B},\sh{M})$ constituído por uma $\sh{O}_S$-álgebra quase-coerente $\sh{B}$ e um $\sh{B}$-módulo quase-coerente $\sh{M}$ \emph{(considerado como $\sh{O}_S$-módulo ou como $\sh{B}$-módulo, o que dá no mesmo~\sref[I]{I.9.6.1})}, existe um par $(X,\sh{F})$ constituído por um pré-esquema $X$ afim sobre $S$ e um $\sh{O}_X$-módulo quase-coerente $\sh{F}$, tal que $\sh{A}(X)=\sh{B}$ e $\sh{A}(\sh{F})=\sh{M}$; além disso, esse par é determinado a menos de isomorfismo único.
\end{proposition}

\begin{proof}
A unicidade resulta de \sref{II.1.4.1} e \sref{II.1.4.2}; a existência se demonstra como em \sref{II.1.3.1}, e deixamos novamente os detalhes ao leitor.
\end{proof}

Denotamos por $\widetilde{\sh{M}}$ o $\sh{O}_X$-módulo $\sh{F}$ e dizemos que ele está \emph{associado} ao $\sh{B}$-módulo quase-coerente $\sh{M}$; para todo aberto afim $U\subset S$, $\widetilde{\sh{M}}|p^{-1}(U)$ (onde $p$ é o morfismo estrutural $X\to S$) é canonicamente isomorfo a $(\Gamma(U,\sh{M}))\supertilde$.

\begin{corollary}[1.4.4]
\label{II.1.4.4}
Na categoria dos $\sh{B}$-módulos quase-coerentes, $\sh{M}\mapsto\widetilde{\sh{M}}$ é um funtor covariante aditivo exato que comuta com os limites indutivos e as somas diretas.
\end{corollary}

\begin{proof}
Reduzimo-nos imediatamente ao caso em que $S$ é afim, e o corolário resulta então de \sref[I]{I.1.3.5}, \sref[I]{I.1.3.9} e \sref[I]{I.1.3.11}.
\end{proof}

\begin{corollary}[1.4.5]
\label{II.1.4.5}
Sob as hipóteses de \sref{II.1.4.3}, para que $\widetilde{\sh{M}}$ seja um $\sh{O}_X$-módulo de tipo finito, é necessário e suficiente que $\sh{M}$ seja um $\sh{B}$-módulo de tipo finito.
\end{corollary}

\begin{proof}
Reduzimo-nos imediatamente ao caso em que $S=\Spec(A)$ é um esquema afim.
Então $\sh{B}=\widetilde{B}$, onde $B$ é uma $A$-álgebra, e $\sh{M}=\widetilde{M}$, onde $M$ é um $B$-módulo \sref[I]{I.1.3.13}; \emph{no pré-esquema $X$}, $\sh{O}_X$ está associado ao anel $B$ e $\widetilde{\sh{M}}$ ao $B$-módulo $M$; para que $\widetilde{\sh{M}}$ seja de tipo finito, é, portanto, necessário e suficiente que $M$ seja de tipo finito \sref[I]{I.1.3.13}, de onde resulta nossa afirmação.
\end{proof}

\begin{proposition}[1.4.6]
\label{II.1.4.6}
Sejam $Y$ um pré-esquema afim sobre $S$, e $X$ e $X'$ dois pré-esquemas afins sobre $Y$ \emph{(portanto, também sobre $S$~\sref{II.1.3.5})}.
Sejam $\sh{B}=\sh{A}(Y)$, $\sh{A}=\sh{A}(X)$ e $\sh{A}'=\sh{A}(X')$.
Então $X\times_Y X'$ é afim sobre $Y$ \emph{(portanto, também sobre $S$)}, e $\sh{A}(X\times_Y X')$ identifica-se canonicamente com $\sh{A}\otimes_\sh{B}\sh{A}'$.
\end{proposition}

\begin{proof}
Com efeito, por \sref[I]{I.9.6.1}, $\sh{A}\otimes_\sh{B}\sh{A}'$ é uma $\sh{B}$-álgebra quase-coerente e, portanto, também uma $\sh{O}_S$-álgebra quase-coerente \sref[I]{I.9.6.1}; seja $Z$ o espectro de $\sh{A}\otimes_\sh{B}\sh{A}'$ \sref{II.1.3.1}.
Os $\sh{B}$-homomorfismos canônicos $\sh{A}\to\sh{A}\otimes_\sh{B}\sh{A}'$ e $\sh{A}'\to\sh{A}\otimes_\sh{B}\sh{A}'$ correspondem, segundo \sref{II.1.2.7}, a $Y$-morfismos $p:Z\to X$ e $p':Z\to X'$.
Para ver que a terna $(Z,p,p')$ é um produto $X\times_Y X'$, podemos nos limitar ao caso em que $S$ é um esquema afim de anel $C$ \sref[I]{I.3.2.6.4}.
Mas então $Y$, $X$ e $X'$ são esquemas afins \sref{II.1.3.4}, cujos anéis $B$, $A$ e $A'$ são $C$-álgebras tais que $\sh{B}=\widetilde{B}$, $\sh{A}=\widetilde{A}$ e $\sh{A}'=\widetilde{A'}$.
Sabe-se então, por \sref[I]{I.1.3.13}, que $\sh{A}\otimes_\sh{B}\sh{A}'$ se identifica canonicamente com a $\sh{O}_S$-álgebra $(A\otimes_B A')\supertilde$; portanto, o anel $A(Z)$ se identifica com $A\otimes_B A'$, e os morfismos $p$ e $p'$ correspondem aos homomorfismos canônicos $A\to A\otimes_B A'$ e $A'\to A\otimes_B A'$.
A proposição resulta então de \sref[I]{I.3.2.2}.
\end{proof}

\begin{corollary}[1.4.7]
\label{II.1.4.7}
Sejam $\sh{F}$ (resp. $\sh{F}'$) um $\sh{O}_X$-módulo quase-coerente (resp. um $\sh{O}_{X'}$-módulo quase-coerente); então $\sh{A}(\sh{F}\otimes_Y\sh{F}')$ identifica-se canonicamente com $\sh{A}(\sh{F})\otimes_{\sh{A}(Y)}\sh{A}(\sh{F}')$.
\end{corollary}

\begin{proof}
Sabe-se que $\sh{F}\otimes_Y\sh{F}'$ é quase-coerente sobre $X\times_Y X'$ \sref[I]{I.9.1.2}.
Sejam $g:Y\to S$, $f:X\to Y$ e $f':X'\to Y$ os morfismos estruturais, de modo que o morfismo estrutural
\oldpage[II]{11}
$h:Z\to S$ é igual a $g\circ f\circ p$ e a $g\circ f'\circ p'$.
Definimos um homomorfismo canônico
\[
  \sh{A}(\sh{F})\otimes_{\sh{A}(Y)}\sh{A}(\sh{F}')
  \to
  \sh{A}(\sh{F}\otimes_Y\sh{F}')
\]
da maneira seguinte: para todo aberto $U\subset S$, temos homomorfismos canônicos
\[
  \begin{aligned}
  \Gamma(f^{-1}(g^{-1}(U)),\sh{F})
  &\longrightarrow\Gamma(h^{-1}(U),p^*(\sh{F})),
  \\[-2pt]
  \Gamma(f^{\prime-1}(g^{-1}(U)),\sh{F}')
  &\longrightarrow\Gamma(h^{-1}(U),p^{\prime*}(\sh{F}')).
  \end{aligned}
\]
Por \sref[0]{0.4.4.3}, deduzimos daí um homomorfismo canônico
\[
  \begin{aligned}
  &\Gamma(f^{-1}(g^{-1}(U)),\sh{F})
    \otimes_{\Gamma(g^{-1}(U),\sh{O}_Y)}
    \Gamma(f^{\prime-1}(g^{-1}(U)),\sh{F}')\longrightarrow\\
  &\hspace{17em}
    \Gamma(h^{-1}(U),p^*(\sh{F}))
    \otimes_{\Gamma(h^{-1}(U),\sh{O}_Z)}
    \Gamma(h^{-1}(U),p^{\prime*}(\sh{F}')).
  \end{aligned}
\]

Para ver que se trata de fato de um isomorfismo de $\sh{A}(Z)$-módulos, podemos nos limitar ao caso em que $S$ (e, por conseguinte, $X$, $X'$, $Y$ e $X\times_Y X'$) são esquemas afins e, com as notações de \sref{II.1.4.6}, $\sh{F}=\widetilde{M}$ e $\sh{F}'=\widetilde{M'}$, sendo $M$ (resp. $M'$) um $A$-módulo (resp. um $A'$-módulo).
Então $\sh{F}\otimes_Y\sh{F}'$ identifica-se com o feixe sobre $X\times_Y X'$ associado ao $(A\otimes_B A')$-módulo $M\otimes_B M'$ \sref[I]{I.9.1.3}, e o corolário resulta da identificação canônica dos $\sh{O}_S$-módulos $(M\otimes_B M')\supertilde$ e $\widetilde{M}\otimes_{\widetilde{B}}\widetilde{M'}$ (onde $M$, $M'$ e $B$ são considerados como $C$-módulos) (\sref[I]{I.1.3.13} e \sref[I]{I.1.6.3}).
\end{proof}

Aplicando, em particular, \sref{II.1.4.7} ao caso em que $X=Y$ e $\sh{F}'=\sh{O}_{X'}$, vemos que o $\sh{A}'$-módulo $\sh{A}(f^{\prime*}(\sh{F}))$ se identifica com $\sh{A}(\sh{F})\otimes_\sh{B}\sh{A}'$.

\begin{env}[1.4.8]
\label{II.1.4.8}
Em particular, quando $X=X'=Y$ (sendo $X$ afim sobre $S$), vemos que, se $\sh{F}$ e $\sh{G}$ são dois $\sh{O}_X$-módulos quase-coerentes, temos
\[
\label{II.1.4.8.1}
  \sh{A}(\sh{F}\otimes_{\sh{O}_X}\sh{G})
  =\sh{A}(\sh{F})\otimes_{\sh{A}(X)}\sh{A}(\sh{G})
  \tag{1.4.8.1}
\]
a menos de isomorfismo canônico funtorial.
Se, além disso, $\sh{F}$ admite uma apresentação finita, resulta de \sref[I]{I.1.6.3} e \sref[I]{I.1.3.12} que
\[
\label{II.1.4.8.2}
  \sh{A}(\shHom_X(\sh{F},\sh{G}))
  =\shHom_{\sh{A}(X)}(\sh{A}(\sh{F}),\sh{A}(\sh{G}))
  \tag{1.4.8.2}
\]
a menos de isomorfismo canônico.
\end{env}

\begin{remark}[1.4.9]
\label{II.1.4.9}
Se $X$ e $X'$ são dois pré-esquemas afins sobre $S$, a soma disjunta $X\sqcup X'$ também é afim sobre $S$, pois a soma disjunta de dois esquemas afins é um esquema afim.
\end{remark}

\begin{proposition}[1.4.10]
\label{II.1.4.10}
Sejam $S$ um pré-esquema, $\sh{B}$ uma $\sh{O}_S$-álgebra quase-coerente e $X=\Spec(\sh{B})$.
Para todo feixe quase-coerente de ideais $\sh{J}$ de $\sh{B}$, $\widetilde{\sh{J}}$ é um feixe quase-coerente de ideais de $\sh{O}_X$, e o subpré-esquema fechado $Y$ de $X$ definido por $\widetilde{\sh{J}}$ é canonicamente isomorfo a $\Spec(\sh{B}/\sh{J})$.
\end{proposition}

\begin{proof}
Resulta imediatamente de \sref[I]{I.4.2.3} que $Y$ é afim sobre $S$; em virtude de \sref{II.1.3.1}, reduzimo-nos, portanto, ao caso em que $S$ é afim, e a proposição resulta então imediatamente de \sref[I]{I.4.1.2}.
\end{proof}

Podemos também exprimir o resultado de \sref{II.1.4.10} dizendo que, se $h:\sh{B}\to\sh{B}'$ é um homomorfismo \emph{sobrejetivo} de $\sh{O}_S$-álgebras quase-coerentes, ${}^a h:\Spec(\sh{B}')\to\Spec(\sh{B})$ é uma \emph{imersão fechada}.

\begin{proposition}[1.4.11]
\label{II.1.4.11}
\oldpage[II]{12}
Sejam $S$ um pré-esquema, $\sh{B}$ uma $\sh{O}_S$-álgebra quase-coerente e $X=\Spec(\sh{B})$.
Para todo feixe quase-coerente de ideais $\sh{K}$ de $\sh{O}_S$, temos (denotando por $f$ o morfismo estrutural $X\to S$)
\[
  f^*(\sh{K})\sh{O}_X=(\sh{K}\sh{B})\supertilde
\]
a menos de isomorfismo canônico.
\end{proposition}

\begin{proof}
Sendo a questão local sobre $S$, reduzimo-nos ao caso em que $S=\Spec(A)$ é afim, e, nesse caso, a proposição não é outra senão \sref[I]{I.1.6.9}.
\end{proof}

\subsection{Mudança do pré-esquema de base}
\label{subsection:II.1.5}

\begin{proposition}[1.5.1]
\label{II.1.5.1}
Seja $X$ um pré-esquema afim sobre $S$.
Para toda extensão $g:S'\to S$ do pré-esquema de base, $X'=X_{(S')}=X\times_S S'$ é afim sobre $S'$.
\end{proposition}

\begin{proof}
Se $f'$ é a projeção $X'\to S'$, basta provar que $f^{\prime-1}(U')$ é um aberto afim para todo aberto afim $U'$ de $S'$ tal que $g(U')$ esteja contido em um aberto afim $U$ de $S$ \sref{II.1.2.1}; podemos, portanto, limitar-nos ao caso em que $S$ e $S'$ são afins, e basta provar que $X'$ é então um esquema afim \sref{II.1.3.4}.
Mas então $X$ é um esquema afim \sref{II.1.3.4} e, se $A$, $A'$ e $B$ são, respectivamente, os anéis de $S$, $S'$ e $X$, sabe-se que $X'$ é o esquema afim de anel $A'\otimes_A B$ \sref[I]{I.3.2.2}.
\end{proof}

\begin{corollary}[1.5.2]
\label{II.1.5.2}
Sob as hipóteses de \sref{II.1.5.1}, sejam $f:X\to S$ o morfismo estrutural, e $f':X'\to S'$ e $g':X'\to X$ as projeções, de modo que o diagrama
\[
  \xymatrix{
    X\ar[d] &
    X'\ar[l]_{g'}\ar[d]^{f'}\\
    S &
    S'\ar[l]_g
  }
\]
é comutativo.
Para todo $\sh{O}_X$-módulo quase-coerente $\sh{F}$, existe um isomorfismo canônico de $\sh{O}_{S'}$-módulos
\[
\label{II.1.5.2.1}
  u:g^*(f_*(\sh{F}))\isoto f'_*(g^{\prime*}(\sh{F})).
  \tag{1.5.2.1}
\]
Em particular, existe um isomorfismo canônico de $\sh{A}(X')$ sobre $g^*(\sh{A}(X))$.
\end{corollary}

\begin{proof}
Para definir $u$, basta definir um homomorfismo
\[
  v:f_*(\sh{F})\to
  g_*(f'_*(g^{\prime*}(\sh{F})))
  =f_*(g'_*(g^{\prime*}(\sh{F})))
\]
e tomar $u=v^\sharp$ \sref[0]{0.4.4.3}.
Tomaremos $v=f_*(\rho)$, onde $\rho$ é o homomorfismo canônico $\sh{F}\to g'_*(g^{\prime*}(\sh{F}))$ \sref[0]{0.4.4.3}.
Para provar que $u$ é um isomorfismo, podemos nos limitar ao caso em que $S$ e $S'$ e, portanto, $X$ e $X'$, são afins; com as notações de \sref{II.1.5.1}, temos então $\sh{F}=\widetilde{M}$, onde $M$ é um $B$-módulo.
Constata-se imediatamente que $g^*(f_*(\sh{F}))$ e $f'_*(g^{\prime*}(\sh{F}))$ são ambos iguais ao $\sh{O}_{S'}$-módulo associado ao $A'$-módulo $A'\otimes_A M$ (onde $M$ é considerado como $A$-módulo), e que $u$ é o homomorfismo associado à identidade (\sref[I]{I.1.6.3}, \sref[I]{I.1.6.5} e \sref[I]{I.1.6.7}).
\end{proof}

\begin{remark}[1.5.3]
\label{II.1.5.3}
Não se deve supor que \sref{II.1.5.2} continue válido quando já não se supõe que $X$ seja afim sobre $S$, mesmo quando $S'=\Spec(\kres(s))$ ($s\in S$) e $S'\to S$ é o morfismo canônico \sref[I]{I.2.4.5}---caso em que $X'$ não é senão a \emph{fibra $f^{-1}(s)$} \sref[I]{I.3.6.2}.
Em outras palavras, quando $X$ não é afim sobre $S$, a operação
\oldpage[II]{13}
«imagem direta de feixes quase-coerentes» não comuta com a operação de «passagem às fibras».
Veremos, contudo, no Capítulo~III \sref[III]{III.4.2.4}, um resultado nesse sentido, de natureza «asintótica», válido para os feixes \emph{coerentes} sobre $X$ quando $f$ é próprio (\hyperref[subsection:II.5.4]{5.4}) e $S$ é noetheriano.
\end{remark}

\begin{corollary}[1.5.4]
\label{II.1.5.4}
Para todo pré-esquema $X$ afim sobre $S$ e todo $s\in S$, a fibra $f^{-1}(s)$ (onde $f$ denota o morfismo estrutural $X\to S$) é um esquema afim.
\end{corollary}

\begin{proof}
Basta aplicar \sref{II.1.5.1} a $S'=\Spec(\kres(s))$ e utilizar \sref{II.1.3.4}.
\end{proof}

\begin{corollary}[1.5.5]
\label{II.1.5.5}
Sejam $X$ um $S$-pré-esquema e $S'$ um pré-esquema afim sobre $S$; então $X'=X_{(S')}$ é um pré-esquema afim sobre $X$.
Além disso, se $f:X\to S$ é o morfismo estrutural, existe um isomorfismo canônico de $\sh{O}_X$-álgebras
\[
  \sh{A}(X')\isoto f^*(\sh{A}(S')),
\]
e, para todo $\sh{A}(S')$-módulo quase-coerente $\sh{M}$, um di-isomorfismo canônico
\[
  f^*(\sh{M})\isoto\sh{A}(f^{\prime*}(\widetilde{\sh{M}})),
\]
denotando por $f'=f_{(S')}$ o morfismo estrutural $X'\to S'$.
\end{corollary}

\begin{proof}
Basta intercambiar os papéis de $X$ e $S'$ em \sref{II.1.5.1} e \sref{II.1.5.2}.
\end{proof}

\begin{env}[1.5.6]
\label{II.1.5.6}
Sejam agora $S$ e $S'$ dois pré-esquemas, $q:S'\to S$ um morfismo, $\sh{B}$ (resp. $\sh{B}'$) uma $\sh{O}_S$-álgebra quase-coerente (resp. uma $\sh{O}_{S'}$-álgebra quase-coerente), e $u:\sh{B}\to\sh{B}'$ um $q$-morfismo (isto é, um homomorfismo $\sh{B}\to q_*(\sh{B}')$ de $\sh{O}_S$-álgebras).
Se $X=\Spec(\sh{B})$ e $X'=\Spec(\sh{B}')$, deduz-se canonicamente um morfismo
\[
  v=\Spec(u):X'\to X
\]
tal que o diagrama
\[
\label{II.1.5.6.1}
  \xymatrix{
    X'\ar[r]^-{v}\ar[d] & X\ar[d]\\
    S'\ar[r]_-{q} & S
  }
  \tag{1.5.6.1}
\]
é comutativo (sendo as setas verticais os morfismos estruturais).

Com efeito, dar $u$ equivale a dar um homomorfismo de $\sh{O}_{S'}$-álgebras quase-coerentes $u^\sharp:q^*(\sh{B})\to\sh{B}'$ \sref[0]{0.4.4.3}; este define, portanto, canonicamente um $S'$-morfismo
\[
  w:\Spec(\sh{B}')\to\Spec(q^*(\sh{B}))
\]
tal que $\sh{A}(w)=u^\sharp$ \sref{II.1.2.7}.
Por outro lado, resulta de \sref{II.1.5.2} que $\Spec(q^*(\sh{B}))$ se identifica canonicamente com $X\times_S S'$; o morfismo $v$ é a composição
\[
  X'\xrightarrow{w}X\times_S S'\xrightarrow{p_1}X
\]
de $w$ com a primeira projeção, e a comutatividade de \sref{II.1.5.6.1} resulta das definições.

Sejam $U$ (resp. $U'$) um aberto afim de $S$ (resp. de $S'$) tal que $q(U')\subset U$, $A=\Gamma(U,\sh{O}_S)$ e $A'=\Gamma(U',\sh{O}_{S'})$ seus anéis, $B=\Gamma(U,\sh{B})$ e $B'=\Gamma(U',\sh{B}')$.
A restrição de $u$, como um $(q|U')$-morfismo $\sh{B}|U\to\sh{B}'|U'$, corresponde a um di-homomorfismo de álgebras $B\to B'$; se $V$ e $V'$ são as imagens inversas de $U$ e $U'$ em $X$ e $X'$, respectivamente, pelos morfismos estruturais, o morfismo $V'\to V$, restrição de $v$, corresponde, segundo \sref[I]{I.1.7.3}, ao di-homomorfismo precedente.
\end{env}

\begin{env}[1.5.7]
\label{II.1.5.7}
Sob as mesmas hipóteses que em \sref{II.1.5.6}, seja $\sh{M}$ um $\sh{B}$-módulo quase-coerente; existe então um isomorfismo canônico de $\sh{O}_{X'}$-módulos
\[
\label{II.1.5.7.1}
  v^*(\widetilde{\sh{M}})
  \isoto
  (q^*(\sh{M})\otimes_{q^*(\sh{B})}\sh{B}')\supertilde.
  \tag{1.5.7.1}
\]
\oldpage[II]{14}
Com efeito, o isomorfismo canônico \sref{II.1.5.2.1} fornece um isomorfismo canônico de $p_1^*(\widetilde{\sh{M}})$ sobre o feixe em $\Spec(q^*(\sh{B}))$ associado ao $q^*(\sh{B})$-módulo $q^*(\sh{M})$, e basta então aplicar \sref{II.1.4.7}.
\end{env}

\subsection{Morfismos afins}
\label{subsection:II.1.6}

\begin{env}[1.6.1]
\label{II.1.6.1}
Diremos que um morfismo $f:X\to Y$ de pré-esquemas é \emph{afim} se define $X$ como um pré-esquema afim sobre $Y$.
As propriedades dos pré-esquemas afins sobre outro se traduzem, nessa linguagem, da maneira seguinte:
\end{env}

\begin{proposition}[1.6.2]
\label{II.1.6.2}
\medskip\noindent
\begin{enumerate}
  \item[{\rm(i)}] Uma imersão fechada é afim.
  \item[{\rm(ii)}] A composição de dois morfismos afins é afim.
  \item[{\rm(iii)}] Se $f:X\to Y$ é um $S$-morfismo afim, então $f_{(S')}:X_{(S')}\to Y_{(S')}$ é afim para toda extensão $S'\to S$ da base.
  \item[{\rm(iv)}] Se $f:X\to Y$ e $f':X'\to Y'$ são dois $S$-morfismos afins, então
    \[
      f\times_S f':X\times_S X'\to Y\times_S Y'
    \]
    é afim.
  \item[{\rm(v)}] Se $f:X\to Y$ e $g:Y\to Z$ são dois morfismos tais que $g\circ f$ é afim e $g$ é separado, então $f$ é afim.
  \item[{\rm(vi)}] Se $f$ é afim, o mesmo vale para $f_\red$.
\end{enumerate}
\end{proposition}

\begin{proof}
Em virtude de \sref[I]{I.5.5.12}, basta demonstrar (i), (ii) e (iii).
Ora, (i) não é outra coisa senão \sref{II.1.2.2}, e (ii) não é outra coisa senão \sref{II.1.3.5}; por fim, (iii) decorre de \sref{II.1.5.1}, pois $X_{(S')}$ se identifica com o produto $X\times_Y Y_{(S')}$ \sref[I]{I.3.3.11}.
\end{proof}

\begin{corollary}[1.6.3]
\label{II.1.6.3}
Se $X$ é um esquema afim e $Y$ um esquema, todo morfismo $X\to Y$ é afim.
\end{corollary}

\begin{proposition}[1.6.4]
\label{II.1.6.4}
Sejam $Y$ um pré-esquema localmente noetheriano e $f:X\to Y$ um morfismo de tipo finito.
Para que $f$ seja afim, é necessário e suficiente que $f_\red$ o seja.
\end{proposition}

\begin{proof}
Em vista de \sref{II.1.6.2}[vi], temos apenas de demonstrar a suficiência da condição.
Basta provar que, se $Y$ é afim e noetheriano, então $X$ é afim; ora, $Y_\red$ é então afim, logo o mesmo vale para $X_\red$ por hipótese.
Ora, $X$ é noetheriano, portanto a conclusão resulta de \sref[I]{I.6.1.7}.
\end{proof}

\subsection{Fibrado vetorial associado a um feixe de módulos}
\label{subsection:II.1.7}

\begin{env}[1.7.1]
\label{II.1.7.1}
Sejam $A$ um anel e $E$ um $A$-módulo.
Recordemos que se denomina \emph{álgebra simétrica} sobre $E$, e se denota por $\bb{S}(E)$ (ou $\bb{S}_A(E)$), a álgebra quociente da álgebra tensorial $\bb{T}(E)$ pelo ideal bilateral gerado pelos elementos $x\otimes y-y\otimes x$, onde $x$ e $y$ percorrem $E$.
A álgebra $\bb{S}(E)$ é caracterizada pela seguinte propriedade universal: se $\sigma$ é a aplicação canônica $E\to\bb{S}(E)$ (obtida compondo $E\to\bb{T}(E)$ com a aplicação canônica $\bb{T}(E)\to\bb{S}(E)$), toda aplicação $A$-linear $E\to B$, onde $B$ é uma $A$-álgebra \emph{comutativa}, fatora-se de maneira única como
\[
  E\xrightarrow{\sigma}\bb{S}(E)\xrightarrow{g}B,
\]
onde $g$ é um $A$-homomorfismo \emph{de álgebras}.
Deduz-se imediatamente dessa caracterização que, para dois $A$-módulos $E$ e $F$, temos
\[
  \bb{S}(E\oplus F)=\bb{S}(E)\otimes\bb{S}(F)
\]
\oldpage[II]{15}
a menos de isomorfismo canônico; além disso, $\bb{S}(E)$ é um funtor covariante em $E$, da categoria dos $A$-módulos para a das $A$-álgebras comutativas; enfim, a caracterização precedente mostra também que, se $E=\varinjlim E_\lambda$, então $\bb{S}(E)=\varinjlim\bb{S}(E_\lambda)$ a menos de isomorfismo canônico.
Por abuso de linguagem, um produto $\sigma(x_1)\sigma(x_2)\cdots\sigma(x_n)$, em que $x_i\in E$, é frequentemente denotado por $x_1x_2\cdots x_n$ quando disso não resulta confusão.
A álgebra $\bb{S}(E)$ é \emph{graduada}, sendo $\bb{S}_n(E)$ o conjunto das combinações lineares dos produtos de $n$ elementos de $E$ ($n\geq0$); a álgebra $\bb{S}(A)$ é canonicamente isomorfa à álgebra de polinômios $A[T]$ em uma indeterminada, e a álgebra $\bb{S}(A^n)$, à álgebra de polinômios em $n$ indeterminadas $A[T_1,\dots,T_n]$.
\end{env}

\begin{env}[1.7.2]
\label{II.1.7.2}
Seja $\vphi:A\to B$ um homomorfismo de anéis.
Se $F$ é um $B$-módulo, a aplicação canônica $F\to\bb{S}(F)$ fornece uma aplicação canônica $F_{[\vphi]}\to\bb{S}(F)_{[\vphi]}$, que, portanto, se fatora em $F_{[\vphi]}\to\bb{S}(F_{[\vphi]})\to\bb{S}(F)_{[\vphi]}$; o homomorfismo canônico $\bb{S}(F_{[\vphi]})\to\bb{S}(F)_{[\vphi]}$ é sobrejetivo, mas não necessariamente bijetivo.
Se $E$ é um $A$-módulo, todo di-homomorfismo $E\to F$ (isto é, todo $A$-homomorfismo $E\to F_{[\vphi]}$) fornece, portanto, canonicamente, um $A$-homomorfismo de álgebras $\bb{S}(E)\to\bb{S}(F_{[\vphi]})\to\bb{S}(F)_{[\vphi]}$, isto é, um di-homomorfismo de álgebras $\bb{S}(E)\to\bb{S}(F)$.

Com as mesmas notações, para todo $A$-módulo $E$, $\bb{S}(E\otimes_A B)$ identifica-se canonicamente com a álgebra $\bb{S}(E)\otimes_A B$; isso resulta imediatamente da propriedade universal de $\bb{S}(E)$ \sref{II.1.7.1}.
\end{env}

\begin{env}[1.7.3]
\label{II.1.7.3}
Seja $R$ um subconjunto multiplicativo do anel $A$; aplicando \sref{II.1.7.2} ao anel $B=R^{-1}A$ e recordando que $R^{-1}E=E\otimes_A R^{-1}A$, vemos que $\bb{S}(R^{-1}E)=R^{-1}\bb{S}(E)$ a menos de isomorfismo canônico.
Além disso, se $R'\supset R$ é um segundo subconjunto multiplicativo de $A$, o diagrama
\[
  \xymatrix{
    {R^{-1}E}\ar[r]\ar[d] & {{R'}^{-1}E}\ar[d]\\
    {\bb{S}(R^{-1}E)}\ar[r] & {\bb{S}({R'}^{-1}E)}
  }
\]
é comutativo.
\end{env}

\begin{env}[1.7.4]
\label{II.1.7.4}
Seja agora $(S,\sh{A})$ um espaço anelado, e seja $\sh{E}$ um $\sh{A}$-módulo sobre $S$.
Se, a todo aberto $U\subset S$, associamos a $\Gamma(U,\sh{A})$-álgebra $\bb{S}(\Gamma(U,\sh{E}))$, obtemos, pela funtorialidade de $\bb{S}$ \sref{II.1.7.2}, um pré-feixe de álgebras; dizemos que o feixe associado, denotado por $\bb{S}(\sh{E})$ ou $\bb{S}_\sh{A}(\sh{E})$, é a \emph{álgebra simétrica sobre $\sh{A}$ do $\sh{A}$-módulo $\sh{E}$}.
Resulta imediatamente de \sref{II.1.7.1} que $\bb{S}(\sh{E})$ é solução de um problema universal: todo homomorfismo de $\sh{A}$-módulos $\sh{E}\to\sh{B}$, em que $\sh{B}$ é uma $\sh{A}$-álgebra comutativa, fatora-se de maneira única em $\sh{E}\to\bb{S}(\sh{E})\to\sh{B}$, sendo a segunda seta um homomorfismo de $\sh{A}$-álgebras.
Há, portanto, uma correspondência biunívoca entre os homomorfismos $\sh{E}\to\sh{B}$ de $\sh{A}$-módulos e os homomorfismos $\bb{S}(\sh{E})\to\sh{B}$ de $\sh{A}$-álgebras.
Em particular, todo homomorfismo $u:\sh{E}\to\sh{F}$ de $\sh{A}$-módulos define um homomorfismo $\bb{S}(u):\bb{S}(\sh{E})\to\bb{S}(\sh{F})$ de $\sh{A}$-álgebras, e $\bb{S}(\sh{E})$ é, portanto, um funtor covariante em $\sh{E}$.
\oldpage[II]{16}
Em virtude de \sref{II.1.7.2} e do fato de que $\bb{S}$ comuta com os limites indutivos, temos $(\bb{S}(\sh{E}))_s=\bb{S}(\sh{E}_s)$ para todo ponto $s\in S$.
Se $\sh{E}$ e $\sh{F}$ são dois $\sh{A}$-módulos, $\bb{S}(\sh{E}\oplus\sh{F})$ identifica-se canonicamente com $\bb{S}(\sh{E})\otimes_\sh{A}\bb{S}(\sh{F})$, como se vê nos pré-feixes correspondentes.

Observemos também que $\bb{S}(\sh{E})$ é uma $\sh{A}$-álgebra graduada, soma direta infinita dos $\bb{S}_n(\sh{E})$, em que o $\sh{A}$-módulo $\bb{S}_n(\sh{E})$ é o feixe associado ao pré-feixe $U\mapsto\bb{S}_n(\Gamma(U,\sh{E}))$.
Tomando, em particular, $\sh{E}=\sh{A}$, vemos que $\bb{S}_\sh{A}(\sh{A})$ se identifica com $\sh{A}[T]=\sh{A}\otimes_\bb{Z}\bb{Z}[T]$ ($T$ indeterminada e $\bb{Z}$ considerado como feixe constante).
\end{env}

\begin{env}[1.7.5]
\label{II.1.7.5}
Sejam $(T,\sh{B})$ um segundo espaço anelado e $f$ um morfismo $(S,\sh{A})\to(T,\sh{B})$.
Se $\sh{F}$ é um $\sh{B}$-módulo, $\bb{S}(f^*(\sh{F}))$ identifica-se canonicamente com $f^*(\bb{S}(\sh{F}))$; com efeito, se $f=(\psi,\theta)$, temos, por definição \sref[0]{0.4.3.1},
\[
  \bb{S}(f^*(\sh{F}))=\bb{S}(\psi^*(\sh{F})\otimes_{\psi^*(\sh{B})}\sh{A})=\bb{S}(\psi^*(\sh{F}))\otimes_{\psi^*(\sh{B})}\sh{A}
\]
\sref{II.1.7.2}; para todo aberto $U$ de $S$ e toda seção $h$ de $\bb{S}(\psi^*(\sh{F}))$ sobre $U$, $h$ coincide, em uma vizinhança $V$ de cada ponto $s\in U$, com um elemento de $\bb{S}(\Gamma(V,\psi^*(\sh{F})))$; se nos reportamos à definição de $\psi^*(\sh{F})$ \sref[0]{0.3.7.1} e levamos em conta que todo elemento de $\bb{S}(E)$, para um módulo $E$, é combinação linear de um número finito de produtos de elementos de $E$, vemos que existem uma vizinhança $W$ de $\psi(s)$ em $T$, uma seção $h'$ de $\bb{S}(\sh{F})$ sobre $W$ e uma vizinhança $V'\subset V\cap\psi^{-1}(W)$ de $s$, tais que $h$ coincide com $t\mapsto h'(\psi(t))$ sobre $V'$; daí nossa afirmação.
\end{env}

\begin{proposition}[1.7.6]
\label{II.1.7.6}
Sejam $A$ um anel, $S=\Spec(A)$ seu espectro primo, e $\sh{E}=\widetilde{M}$ o $\sh{O}_S$-módulo associado a um $A$-módulo $M$; então a $\sh{O}_S$-álgebra $\bb{S}(\sh{E})$ está associada à $A$-álgebra $\bb{S}(M)$.
\end{proposition}

\begin{proof}
Com efeito, para todo $f\in A$, $\bb{S}(M_f)=(\bb{S}(M))_f$ \sref{II.1.7.3}, e a proposição resulta, portanto, das definições \sref[I]{I.1.3.4}.
\end{proof}

\begin{corollary}[1.7.7]
\label{II.1.7.7}
Se $S$ é um pré-esquema e $\sh{E}$ um $\sh{O}_S$-módulo quase-coerente, a $\sh{O}_S$-álgebra $\bb{S}(\sh{E})$ é quase-coerente.
Se, além disso, $\sh{E}$ é de tipo finito, cada um dos $\sh{O}_S$-módulos $\bb{S}_n(\sh{E})$ é de tipo finito.
\end{corollary}

\begin{proof}
A primeira afirmação é consequência imediata de \sref{II.1.7.6} e de \sref[I]{I.1.4.1}; a segunda resulta do fato de que, se $E$ é um $A$-módulo de tipo finito, $\bb{S}_n(E)$ é um $A$-módulo de tipo finito; aplica-se então \sref[I]{I.1.3.13}.
\end{proof}

\begin{definition}[1.7.8]
\label{II.1.7.8}
Seja $\sh{E}$ um $\sh{O}_S$-módulo quase-coerente.
Chama-se \emph{fibrado vetorial sobre $S$ definido por $\sh{E}$}, e denota-se por $\bb{V}(\sh{E})$, o espectro \sref{II.1.3.1} da $\sh{O}_S$-álgebra quase-coerente $\bb{S}(\sh{E})$.
\end{definition}

Em virtude de \sref{II.1.2.7}, para todo $S$-pré-esquema $X$, há uma correspondência biunívoca canônica entre os $S$-morfismos $X\to\bb{V}(\sh{E})$ e os homomorfismos de \emph{$\sh{O}_S$-álgebras $\bb{S}(\sh{E})\to\sh{A}(X)$}, portanto também entre esses $S$-morfismos e os homomorfismos de \emph{$\sh{O}_S$-módulos $\sh{E}\to\sh{A}(X)=f_*(\sh{O}_X)$} (em que $f$ é o morfismo estrutural $X\to S$).
Em particular:
\begin{env}[1.7.9]
\label{II.1.7.9}
Tomemos para $X$ o subpré-esquema de $S$ induzido sobre um \emph{aberto $U\subset S$}.
Então, os $S$-morfismos $U\to\bb{V}(\sh{E})$ nada mais são do que as $U$-seções \sref[I]{I.2.5.5} do $U$-pré-esquema induzido por $\bb{V}(\sh{E})$ no aberto $p^{-1}(U)$ (em que $p$ é o morfismo estrutural $\bb{V}(\sh{E})\to S$).
Segundo o que acabamos de ver, essas $U$-seções correspondem biunivocamente aos homomorfismos de $\sh{O}_S$-módulos $\sh{E}\to j_*(\sh{O}_S|U)$ (em que $j$ é a injeção canônica $U\to S$), ou,
\oldpage[II]{17}
o que é equivalente, em virtude de \sref[0]{0.4.4.3}, aos homomorfismos de $(\sh{O}_S|U)$-módulos $j^*(\sh{E})=\sh{E}|U\to\sh{O}_S|U$.
Além disso, é imediato que à restrição a um aberto $U'\subset U$ de um $S$-morfismo $U\to\bb{V}(\sh{E})$ corresponde a restrição a $U'$ do homomorfismo correspondente $\sh{E}|U\to\sh{O}_S|U$.
Concluímos que \emph{o feixe dos germes de $S$-seções de $\bb{V}(\sh{E})$ se identifica canonicamente com o dual $\dual{\sh{E}}$ de $\sh{E}$}.

Em particular, tomando $X=U=S$, ao homomorfismo \emph{nulo} $\sh{E}\to\sh{O}_S$ corresponde uma $S$-seção canônica de $\bb{V}(\sh{E})$, chamada \emph{$S$-seção nula} (cf. \sref{II.8.3.3}).
\end{env}

\begin{env}[1.7.10]
\label{II.1.7.10}
Tomemos agora para $X$ o espectro $\{\xi\}$ de um corpo $K$; o morfismo estrutural $f:X\to S$ corresponde, portanto, a um monomorfismo $\kres(s)\to K$, em que $s=f(\xi)$ \sref[I]{I.2.4.6}; os $S$-morfismos $\{\xi\}\to\bb{V}(\sh{E})$ nada mais são, portanto, do que os \emph{pontos geométricos de $\bb{V}(\sh{E})$ a valores na extensão $K$ de $\kres(s)$} \sref[I]{I.3.4.5}, pontos que se localizam nos pontos de $p^{-1}(s)$.
O conjunto desses pontos, que ainda se pode chamar de \emph{fibra geométrica racional sobre $K$ de $\bb{V}(\sh{E})$ acima do ponto $s$}, identifica-se, segundo \sref{II.1.7.8}, com o conjunto dos homomorfismos de $\sh{O}_S$-módulos $\sh{E}\to f_*(\sh{O}_X)$, ou, o que é equivalente \sref[0]{0.4.4.3}, com o conjunto dos homomorfismos de $\sh{O}_X$-módulos $f^*(\sh{E})\to\sh{O}_X=K$.
Mas temos, por definição \sref[0]{0.4.3.1}, $f^*(\sh{E})=\sh{E}_s\otimes_{\sh{O}_s}K=\sh{E}^s\otimes_{\kres(s)}K$, pondo $\sh{E}^s=\sh{E}_s/\mathfrak{m}_s\sh{E}_s$; a fibra geométrica de $\bb{V}(\sh{E})$ racional sobre $K$ acima de $s$ identifica-se, portanto, com o \emph{dual do $K$-espaço vetorial $\sh{E}^s\otimes_{\kres(s)}K$}; se $\sh{E}^s$ ou $K$ tem dimensão finita sobre $\kres(s)$, esse dual identifica-se também com $\dual{(\sh{E}^s)}\otimes_{\kres(s)}K$, denotando por $\dual{(\sh{E}^s)}$ o dual do $\kres(s)$-espaço vetorial $\sh{E}^s$.
\end{env}

\begin{proposition}[1.7.11]
\label{II.1.7.11}
\medskip\noindent
\begin{enumerate}
  \item[{\rm(i)}] $\bb{V}(\sh{E})$ é um funtor contravariante em $\sh{E}$, da categoria dos $\sh{O}_S$-módulos quase-coerentes para a categoria dos $S$-esquemas afins.
  \item[{\rm(ii)}] Se $\sh{E}$ é um $\sh{O}_S$-módulo de tipo finito, $\bb{V}(\sh{E})$ é de tipo finito sobre $S$.
  \item[{\rm(iii)}] Se $\sh{E}$ e $\sh{F}$ são dois $\sh{O}_S$-módulos quase-coerentes, $\bb{V}(\sh{E}\oplus\sh{F})$ identifica-se canonicamente com $\bb{V}(\sh{E})\times_S\bb{V}(\sh{F})$.
  \item[{\rm(iv)}] Seja $g:S'\to S$ um morfismo; para todo $\sh{O}_S$-módulo quase-coerente $\sh{E}$, $\bb{V}(g^*(\sh{E}))$ identifica-se canonicamente com $\bb{V}(\sh{E})_{(S')}=\bb{V}(\sh{E})\times_S S'$.
  \item[{\rm(v)}] A um homomorfismo sobrejetivo $\sh{E}\to\sh{F}$ de $\sh{O}_S$-módulos quase-coerentes corresponde uma imersão fechada $\bb{V}(\sh{F})\to\bb{V}(\sh{E})$.
\end{enumerate}
\end{proposition}

\begin{proof}
(i) é consequência imediata de \sref{II.1.2.7}, levando em conta que todo homomorfismo de $\sh{O}_S$-módulos $\sh{E}\to\sh{F}$ define canonicamente um homomorfismo de $\sh{O}_S$-álgebras $\bb{S}(\sh{E})\to\bb{S}(\sh{F})$.
(ii) resulta imediatamente das definições \sref[I]{I.6.3.1} e do fato de que, se $E$ é um $A$-módulo de tipo finito, $\bb{S}(E)$ é uma $A$-álgebra de tipo finito.
Para demonstrar (iii), basta partir do isomorfismo canônico $\bb{S}(\sh{E}\oplus\sh{F})\isoto\bb{S}(\sh{E})\otimes_{\sh{O}_S}\bb{S}(\sh{F})$ \sref{II.1.7.4} e aplicar \sref{II.1.4.6}.
Do mesmo modo, para demonstrar (iv), basta partir do isomorfismo canônico $\bb{S}(g^*(\sh{E}))\isoto g^*(\bb{S}(\sh{E}))$ \sref{II.1.7.5} e aplicar \sref{II.1.5.2}.
Enfim, para estabelecer (v), basta observar que, se o homomorfismo $\sh{E}\to\sh{F}$ é sobrejetivo, o mesmo vale para o homomorfismo correspondente $\bb{S}(\sh{E})\to\bb{S}(\sh{F})$ de $\sh{O}_S$-álgebras, e a conclusão resulta de \sref{II.1.4.10}.
\end{proof}

\begin{env}[1.7.12]
\label{II.1.7.12}
Tomemos, em particular, $\sh{E}=\sh{O}_S$; o pré-esquema $\bb{V}(\sh{O}_S)$ é o $S$-esquema afim, espectro da $\sh{O}_S$-álgebra $\bb{S}(\sh{O}_S)$, que se identifica com a $\sh{O}_S$-álgebra $\sh{O}_S[T]=\sh{O}_S\otimes_\bb{Z}\bb{Z}[T]$
\oldpage[II]{18}
($T$ indeterminada); isso é evidente quando $S=\Spec(\bb{Z})$, em virtude de \sref{II.1.7.6}, e daí se passa ao caso geral considerando o morfismo estrutural $S\to\Spec(\bb{Z})$ e utilizando \sref{II.1.7.11}[iv].
Em razão desse resultado, ainda se põe $\bb{V}(\sh{O}_S)=S[T]$, e temos, portanto, a fórmula
\[
\label{II.1.7.12.1}
  S[T]=S\otimes_\bb{Z}\bb{Z}[T].
  \tag{1.7.12.1}
\]

A identificação do feixe de germes de $S$-seções de $S[T]$ com $\sh{O}_S$, já vista em \sref[I]{I.3.3.15}, reaparece aqui, em um contexto mais geral, como caso particular de \sref{II.1.7.9}.
\end{env}

\begin{env}[1.7.13]
\label{II.1.7.13}
Para todo $S$-pré-esquema $X$, vimos em \sref{II.1.7.8} que $\Hom_S(X,S[T])$ se identifica canonicamente com $\Hom_{\sh{O}_S}(\sh{O}_S,\sh{A}(X))$, o qual é canonicamente isomorfo a $\Gamma(S,\sh{A}(X))$ e, por conseguinte, munido de uma estrutura de anel; além disso, todo $S$-morfismo $h:X\to Y$ induz um homomorfismo de anéis $\Gamma(\sh{A}(h)):\Gamma(S,\sh{A}(Y))\to\Gamma(S,\sh{A}(X))$ \sref{II.1.1.2}.
Quando se mune $\Hom_S(X,S[T])$ da estrutura de anel assim definida, vemos, portanto, que se pode considerar $\Hom_S(X,S[T])$ como um \emph{funtor contravariante em $X$}, da categoria dos $S$-pré-esquemas para a dos anéis.
Por outro lado, $\Hom_S(X,\bb{V}(\sh{E}))$ identifica-se igualmente com $\Hom_{\sh{O}_S}(\sh{E},\sh{A}(X))$ (em que $\sh{A}(X)$ é considerado como um \emph{$\sh{O}_S$-módulo}); por conseguinte, pode-se munir canonicamente esse conjunto de uma estrutura de \emph{módulo} sobre o anel $\Hom_S(X,S[T])$, e vemos, como acima, que o par
\[
  (\Hom_S(X,S[T]),\Hom_S(X,\bb{V}(\sh{E})))
\]
é um funtor contravariante em $X$, com valores na categoria cujos objetos são os pares $(A,M)$ formados por um anel $A$ e um $A$-módulo $M$, sendo os morfismos os di-homomorfismos.

Interpretaremos esses fatos dizendo que $S[T]$ é um \emph{$S$-esquema de anéis} e que $\bb{V}(\sh{E})$ é um \emph{$S$-esquema de módulos} sobre o $S$-esquema de anéis $S[T]$ (cf. Capítulo~0, \textsection8).
\end{env}

\begin{env}[1.7.14]
\label{II.1.7.14}
Veremos que a estrutura de $S$-esquema de módulos definida sobre o $S$-esquema $\bb{V}(\sh{E})$ permite reconstruir o $\sh{O}_S$-módulo $\sh{E}$ a menos de isomorfismo único: para isso, mostraremos que $\sh{E}$ é canonicamente isomorfo a um $\sh{O}_S$-submódulo de $\bb{S}(\sh{E})=\sh{A}(\bb{V}(\sh{E}))$, definido por meio dessa estrutura.
Com efeito, por \sref{II.1.7.4}, o conjunto $\Hom_{\sh{O}_S}(\bb{S}(\sh{E}),\sh{A}(X))$ dos homomorfismos de \emph{$\sh{O}_S$-álgebras} identifica-se canonicamente com $\Hom_{\sh{O}_S}(\sh{E},\sh{A}(X))$, conjunto dos homomorfismos de \emph{$\sh{O}_S$-módulos}: se $h$ e $h'$ são dois elementos deste último conjunto, $s_i$ ($1\leq i\leq n$) são seções de $\sh{E}$ acima de um aberto $U\subset S$, e $t$ é uma seção de $\sh{A}(X)$ acima de $U$, temos, por definição,
\[
  (h+h')(s_1s_2\cdots s_n)=\prod_{i=1}^n(h(s_i)+h'(s_i))
\]
e
\[
  (t\cdot h)(s_1s_2\cdots s_n)=t^n\prod_{i=1}^n h(s_i).
\]

Sendo assim, se $z$ é uma seção de $\bb{S}(\sh{E})$ acima de $U$, $h\mapsto h(z)$ é uma aplicação de $\Hom_S(X,\bb{V}(\sh{E}))=\Hom_{\sh{O}_S}(\bb{S}(\sh{E}),\sh{A}(X))$ em $\Gamma(U,\sh{A}(X))$.
Vamos
\oldpage[II]{19}
mostrar que $\sh{E}$ se identifica com o submódulo de $\bb{S}(\sh{E})$ tal que, \emph{para todo aberto $U\subset S$, toda seção $z$ desse $\sh{O}_S$-submódulo acima de $U$ e todo $S$-pré-esquema $X$, a aplicação $h\mapsto h(z)$ de $\Hom_{\sh{O}_S|U}(\bb{S}(\sh{E})|U,\sh{A}(X)|U)$ em $\Gamma(U,\sh{A}(X))$ seja um homomorfismo de $\Gamma(U,\sh{A}(X))$-módulos}.

É, com efeito, imediato que $\sh{E}$ possui essa propriedade; para demonstrar a recíproca, podemos limitar-nos a provar que, quando $S=\Spec(A)$ e $\sh{E}=\widetilde{M}$, uma seção $z$ de $\bb{S}(\sh{E})$ acima de $S$ que, para $U=S$, possui a propriedade enunciada acima é necessariamente uma seção de $\sh{E}$; temos então $z=\sum_{n=0}^{\infty}z_n$, em que $z_n\in\bb{S}_n(M)$, e trata-se de provar que $z_n=0$ para $n\neq1$.
Ponhamos $B=\bb{S}(M)$ e tomemos para $X$ o pré-esquema $\Spec(B[T])$, em que $T$ é uma indeterminada.
O conjunto $\Hom_{\sh{O}_S}(\bb{S}(\sh{E}),\sh{A}(X))$ identifica-se com o conjunto dos homomorfismos de anéis $h:B\to B[T]$ \sref[I]{I.1.3.13}, e, segundo o que vimos acima, temos $(T\cdot h)(z)=\sum_{n=0}^{\infty}T^n h(z_n)$; a hipótese sobre $z$ implica que $\sum_{n=0}^{\infty}T^n h(z_n)=T\cdot\sum_{n=0}^{\infty}h(z_n)$ para todo homomorfismo $h$.
Tomando, em particular, para $h$ a injeção canônica, obtemos $\sum_{n=0}^{\infty}T^n z_n=T\cdot\sum_{n=0}^{\infty}z_n$, o que acarreta de fato a conclusão $z_n=0$ para $n\neq1$.
\end{env}

\begin{proposition}[1.7.15]
\label{II.1.7.15}
Seja $Y$ um pré-esquema cujo espaço subjacente é noetheriano, ou um esquema quase-compacto.
Todo $Y$-esquema afim $X$ de tipo finito sobre $Y$ é $Y$-isomorfo a um $Y$-subesquema fechado de um $Y$-esquema da forma $\bb{V}(\sh{E})$, em que $\sh{E}$ é um $\sh{O}_Y$-módulo quase-coerente de tipo finito.
\end{proposition}

\begin{proof}
Com efeito, a $\sh{O}_Y$-álgebra quase-coerente $\sh{A}(X)$ é de tipo finito \sref{II.1.3.7}.
As hipóteses feitas acarretam que $\sh{A}(X)$ é gerada por um $\sh{O}_Y$-submódulo quase-coerente de tipo finito $\sh{E}$ \sref[I]{I.9.6.5}; por definição, isso acarreta que o homomorfismo canônico $\bb{S}(\sh{E})\to\sh{A}(X)$ que prolonga canonicamente a injeção $\sh{E}\to\sh{A}(X)$ é \emph{sobrejetivo}; a conclusão resulta, então, de \sref{II.1.4.10}.
\end{proof}

\section{Espectros primos homogêneos}
\label{section:II.2}

\subsection{Generalidades sobre anéis e módulos graduados}
\label{subsection:II.2.1}

\begin{notation}[2.1.1]
\label{II.2.1.1}
Dado um anel $S$ \emph{graduado em graus positivos}, denotaremos por $S_n$ a parte de $S$ formada pelos elementos homogêneos de grau $n$ ($n\geq0$), e por $S_+$ a soma (direta) dos $S_n$ tais que $n>0$; temos $1\in S_0$, $S_0$ é um subanel de $S$, $S_+$ é um ideal graduado de $S$, e $S$ é soma direta de $S_0$ e de $S_+$.
Se $M$ é um módulo \emph{graduado} sobre $S$ (com graus positivos ou negativos), denotaremos, do mesmo modo, por $M_n$ o $S_0$-módulo formado pelos elementos homogêneos de grau $n$ de $M$ (com $n\in\bb{Z}$).

Para todo inteiro $d>0$, denotaremos por $S^{(d)}$ a soma direta dos $S_{nd}$; considerando os elementos de $S_{nd}$ como homogêneos de grau $n$, os $S_{nd}$ definem sobre $S^{(d)}$ uma estrutura de anel graduado.

Para todo inteiro $k$ tal que $0\leq k\leq d-1$, denotaremos por $M^{(d,k)}$ a soma direta
\oldpage[II]{20}
dos $M_{nd+k}$ ($n\in\bb{Z}$); trata-se de um $S^{(d)}$-módulo graduado quando se consideram os elementos de $M_{nd+k}$ como homogêneos de grau $n$.
Escreveremos $M^{(d)}$ em lugar de $M^{(d,0)}$.

Com as notações precedentes, para todo inteiro $n$ (positivo ou negativo), denotaremos por $M(n)$ o $S$-módulo graduado definido por $(M(n))_k=M_{n+k}$ para todo $k\in\bb{Z}$.
Em particular, $S(n)$ será um $S$-módulo graduado tal que $(S(n))_k=S_{n+k}$, convencionando-se que $S_n=0$ para $n<0$.
Diz-se que um $S$-módulo graduado $M$ é \emph{livre} se é isomorfo, como \emph{módulo graduado}, a uma soma direta de módulos da forma $S(n)$; como $S(n)$ é um $S$-módulo cíclico, gerado pelo elemento $1$ de $S$ considerado como elemento de grau $-n$, isso equivale a dizer que $M$ admite uma \emph{base} sobre $S$ formada por elementos \emph{homogêneos}.

Diz-se que um $S$-módulo graduado $M$ \emph{admite uma apresentação finita} se existe uma sequência exata $P\to Q\to M\to0$, na qual $P$ e $Q$ são somas diretas finitas de módulos da forma $S(n)$ e os homomorfismos são de grau $0$ (cf. \sref{II.2.1.2}).
\end{notation}

\begin{env}[2.1.2]
\label{II.2.1.2}
Sejam $M$ e $N$ dois $S$-módulos graduados; define-se sobre $M\otimes_S N$ uma estrutura de $S$-módulo \emph{graduado} da seguinte maneira.
Sobre o produto tensorial $M\otimes_\bb{Z}N$, pode-se definir uma estrutura de $\bb{Z}$-módulo graduado (onde $\bb{Z}$ é graduado por $\bb{Z}_0=\bb{Z}$, $\bb{Z}_n=0$ para $n\neq0$), pondo
\[
  (M\otimes_\bb{Z}N)_q
  =\bigoplus_{m+n=q}M_m\otimes_\bb{Z}N_n.
\]
Como $M$ e $N$ são, respectivamente, somas diretas dos $M_m$ e dos $N_n$, sabe-se que $M\otimes_\bb{Z}N$ se identifica canonicamente com a soma direta de todos os $M_m\otimes_\bb{Z}N_n$.
Posto isso, tem-se $M\otimes_S N=(M\otimes_\bb{Z}N)/P$, onde $P$ é o $\bb{Z}$-submódulo de $M\otimes_\bb{Z}N$ gerado pelos elementos $(xs)\otimes y-x\otimes(sy)$ para $x\in M$, $y\in N$ e $s\in S$; é claro que $P$ é um $\bb{Z}$-submódulo \emph{graduado} de $M\otimes_\bb{Z}N$, e vê-se imediatamente que, passando ao quociente, obtém-se de fato uma estrutura de $S$-módulo graduado sobre $M\otimes_S N$.

Para dois $S$-módulos graduados $M$ e $N$, lembremos que um homomorfismo $u:M\to N$ de $S$-módulos se diz \emph{de grau $k$} se $u(M_j)\subset N_{j+k}$ para todo $j\in\bb{Z}$.
Se $H_n$ denota o conjunto de todos os homomorfismos de grau $n$ de $M$ em $N$, denota-se por $\Hom_S(M,N)$ a \emph{soma} (direta) dos $H_n$ ($n\in\bb{Z}$) no $S$-módulo $H$ de todos os homomorfismos (de $S$-módulos) de $M$ em $N$; em geral, $\Hom_S(M,N)$ não é igual a este último.
Todavia, tem-se $H=\Hom_S(M,N)$ quando $M$ é \emph{de tipo finito}; com efeito, pode-se então supor que $M$ seja gerado por um número finito de elementos homogêneos $x_i$ ($1\leq i\leq n$), e todo homomorfismo $u\in H$ se escreve de maneira única como $\sum_{k\in\bb{Z}}u_k$, onde, para cada $k$, $u_k(x_i)$ é igual à componente homogênea de grau $k+\deg(x_i)$ de $u(x_i)$ ($1\leq i\leq n$), o que implica $u_k=0$ salvo para um número finito de índices; por definição, $u_k\in H_k$, de onde resulta a conclusão.

Os elementos de grau $0$ de $\Hom_S(M,N)$ são chamados \emph{homomorfismos de $S$-módulos graduados}.
É claro que $S_mH_n\subset H_{m+n}$, logo os $H_n$ definem sobre $\Hom_S(M,N)$ uma estrutura de $S$-módulo graduado.

Resulta imediatamente dessas definições que
\[
\label{II.2.1.2.1}
  M(m)\otimes_S N(n)=(M\otimes_S N)(m+n),
\tag{2.1.2.1}
\]
\[
\label{II.2.1.2.2}
  \Hom_S(M(m),N(n))=(\Hom_S(M,N))(n-m),
\tag{2.1.2.2}
\]
para dois $S$-módulos graduados $M$ e $N$.
\oldpage[II]{21}
Sejam $S$ e $S'$ dois anéis graduados; um homomorfismo de \emph{anéis graduados $\vphi:S\to S'$} é um homomorfismo de anéis tal que $\vphi(S_n)\subset S_n'$ para todo $n\in\bb{Z}$ (em outros termos, $\vphi$ deve ser um homomorfismo \emph{de grau $0$} de $\bb{Z}$-módulos graduados).
O dado de tal homomorfismo define sobre $S'$ uma estrutura de $S$-módulo \emph{graduado}; munido dessa estrutura e de sua estrutura de anel graduado, diz-se que $S'$ é uma \emph{$S$-álgebra graduada}.

Se $M$ é um $S$-módulo graduado, o produto tensorial $M\otimes_S S'$ de $S$-módulos \emph{graduados} fica munido naturalmente de uma estrutura de $S'$-módulo \emph{graduado}, sendo a graduação a definida acima.
\end{env}

\begin{lemma}[2.1.3]
\label{II.2.1.3}
Seja $S$ um anel graduado em graus positivos.
Para que um subconjunto $E$ de $S_+$ formado por elementos homogêneos gere $S_+$ como $S$-módulo, é necessário e suficiente que $E$ gere $S$ como $S_0$-álgebra.
\end{lemma}

\begin{proof}
A condição é evidentemente suficiente; mostremos que é necessária.
Seja $E_n$ (resp. $E^n$) o conjunto dos elementos de $E$ de grau igual a $n$ (resp. de grau $\leq n$); bastará mostrar, por indução sobre $n>0$, que $S_n$ é o $S_0$-módulo gerado pelos elementos de grau $n$ que são produtos de elementos de $E^n$.
Ora, isso é evidente para $n=1$ em virtude da hipótese; esta mostra, além disso, que
\[
  S_n=\sum_{p=0}^{n-1}S_pE_{n-p},
\]
e o raciocínio por indução é então imediato.
\end{proof}

\begin{corollary}[2.1.4]
\label{II.2.1.4}
Para que $S_+$ seja um ideal de tipo finito, é necessário e suficiente que $S$ seja uma $S_0$-álgebra de tipo finito.
\end{corollary}

\begin{proof}
Pode-se sempre supor que um sistema finito de geradores da $S_0$-álgebra $S$ (resp. do $S$-ideal $S_+$) seja formado por elementos homogêneos, substituindo cada um dos geradores considerados por suas componentes homogêneas.
\end{proof}

\begin{corollary}[2.1.5]
\label{II.2.1.5}
Para que $S$ seja noetheriano, é necessário e suficiente que $S_0$ seja noetheriano e que $S$ seja uma $S_0$-álgebra de tipo finito.
\end{corollary}

\begin{proof}
A condição é evidentemente suficiente; ela é necessária, pois $S_0$ é isomorfo a $S/S_+$ e $S_+$ deve ser um ideal de tipo finito \sref{II.2.1.4}.
\end{proof}

\begin{lemma}[2.1.6]
\label{II.2.1.6}
Seja $S$ um anel graduado em graus positivos que é uma $S_0$-álgebra de tipo finito.
Seja $M$ um $S$-módulo graduado de tipo finito.
Então:
\begin{enumerate}
  \item[{\rm(i)}] Os $M_n$ são $S_0$-módulos de tipo finito, e existe um inteiro $n_0$ tal que $M_n=0$ para $n\leq n_0$.
  \item[{\rm(ii)}] Existem um inteiro $n_1$ e um inteiro $h>0$ tais que, para todo inteiro $n\geq n_1$, tem-se $M_{n+h}=S_hM_n$.
  \item[{\rm(iii)}] Para todo par de inteiros $(d,k)$ tais que $d>0$ e $0\leq k\leq d-1$, $M^{(d,k)}$ é um $S^{(d)}$-módulo de tipo finito.
  \item[{\rm(iv)}] Para todo inteiro $d>0$, $S^{(d)}$ é uma $S_0$-álgebra de tipo finito.
  \item[{\rm(v)}] Existe um inteiro $h>0$ tal que $S_{mh}=(S_h)^m$ para todo $m>0$.
  \item[{\rm(vi)}] Para todo inteiro $n>0$, existe um inteiro $m_0$ tal que $S_m\subset S_+^n$ para todo $m\geq m_0$.
\end{enumerate}
\end{lemma}

\begin{proof}
Pode-se supor que $S$ seja gerado (como $S_0$-álgebra) por elementos homogêneos $f_i$, de graus $h_i$ ($1\leq i\leq r$), e que $M$ seja gerado (como $S$-módulo) por elementos homogêneos $x_j$ de graus $k_j$ ($1\leq j\leq s$).
É claro que $M_n$ é formado por combinações
\oldpage[II]{22}
lineares, com coeficientes em $S_0$, de elementos $f_1^{\alpha_1}\cdots f_r^{\alpha_r}x_j$ tais que os $\alpha_i$ sejam inteiros $\geq0$ que satisfaçam
\[
  k_j+\sum_i\alpha_i h_i=n.
\]
Para cada $j$, há apenas um número finito de sistemas $(\alpha_i)$ que satisfazem essa equação, pois os $h_i$ são $>0$, de onde resulta a primeira afirmação de~(i); a segunda é evidente.

Por outro lado, seja $h$ o mínimo múltiplo comum dos $h_i$ e ponhamos $g_i=f_i^{h/h_i}$ ($1\leq i\leq r$), de modo que todos os $g_i$ sejam de grau $h$; sejam $z_\mu$ os elementos de $M$ da forma $f_1^{\alpha_1}\cdots f_r^{\alpha_r}x_j$, com $0\leq\alpha_i<h/h_i$ para $1\leq i\leq r$.
Esses elementos são em número finito; seja $n_1$ o maior de seus graus.
É claro que, para $n\geq n_1$, todo elemento de $M_{n+h}$ é combinação linear dos $z_\mu$ cujos coeficientes são monômios de grau $>0$ nos $g_i$; portanto, tem-se $M_{n+h}=S_hM_n$, o que estabelece~(ii).
Da mesma maneira, vê-se (para todo $d>0$) que um elemento de $M^{(d,k)}$ é combinação linear, com coeficientes em $S_0$, de elementos da forma $g^d f_1^{\alpha_1}\cdots f_r^{\alpha_r}x_j$, com $0\leq\alpha_i<d$, sendo $g$ um elemento homogêneo de $S$; daí resulta~(iii).
(iv) resulta então de~(iii) e do Lema~\sref{II.2.1.3}, tomando $M=S_+$, pois $(S_+)^{(d)}=(S^{(d)})_+$.
A afirmação~(v) deduz-se de~(ii), tomando $M=S$.
Por fim, para um $n$ dado, há apenas um número finito de sistemas $(\alpha_i)$ tais que $\alpha_i\geq0$ e $\sum_i\alpha_i<n$; se $m_0$ é escolhido estritamente maior que o maior valor de $\sum_i\alpha_i h_i$ para esses sistemas, tem-se $S_m\subset S_+^n$ para $m\geq m_0$, o que prova~(vi).
\end{proof}

\begin{corollary}[2.1.7]
\label{II.2.1.7}
Se $S$ é noetheriano, o mesmo vale para $S^{(d)}$ para todo inteiro $d>0$.
\end{corollary}

\begin{proof}
Isso resulta de \sref{II.2.1.5} e de \sref{II.2.1.6}[iv].
\end{proof}

\begin{env}[2.1.8]
\label{II.2.1.8}
Seja $\mathfrak{p}$ um ideal primo \emph{graduado} do anel graduado $S$; $\mathfrak{p}$ é, portanto, soma direta dos subgrupos $\mathfrak{p}_n=\mathfrak{p}\cap S_n$.
Suponhamos que \emph{$\mathfrak{p}$ não contenha $S_+$}.
Então, se $f\in S_+$ não pertence a $\mathfrak{p}$, a relação $f^nx\in\mathfrak{p}$ equivale a $x\in\mathfrak{p}$; em particular, se $f\in S_d$ ($d>0$), para todo $x\in S_{m-nd}$, a relação $f^nx\in\mathfrak{p}_m$ equivale a $x\in\mathfrak{p}_{m-nd}$.
\end{env}

\begin{proposition}[2.1.9]
\label{II.2.1.9}
Seja $n_0$ um inteiro $>0$; para todo $n\geq n_0$, seja $\mathfrak{p}_n$ um subgrupo de $S_n$.
Para que exista um ideal primo graduado $\mathfrak{p}$ de $S$ que não contenha $S_+$ e tal que $\mathfrak{p}\cap S_n=\mathfrak{p}_n$ para todo $n\geq n_0$, é necessário e suficiente que sejam verificadas as seguintes condições:
\begin{enumerate}
  \item[{\rm($1^\circ$)}] $S_m\mathfrak{p}_n\subset\mathfrak{p}_{m+n}$ para todo $m\geq0$ e todo $n\geq n_0$.
  \item[{\rm($2^\circ$)}] Para $m\geq n_0$, $n\geq n_0$, $f\in S_m$ e $g\in S_n$, a relação $fg\in\mathfrak{p}_{m+n}$ implica $f\in\mathfrak{p}_m$ ou $g\in\mathfrak{p}_n$.
  \item[{\rm($3^\circ$)}] $\mathfrak{p}_n\neq S_n$ para pelo menos um $n\geq n_0$.
\end{enumerate}
Além disso, o ideal primo graduado $\mathfrak{p}$ é então único.
\end{proposition}

\begin{proof}
É evidente que as condições $1^\circ$ e $2^\circ$ são necessárias.
Além disso, se $\mathfrak{p}\not\supset S_+$, existe pelo menos um $k>0$ tal que $\mathfrak{p}\cap S_k\neq S_k$; se $f\in S_k$ não pertence a $\mathfrak{p}$, a relação $\mathfrak{p}\cap S_n=S_n$ implica $\mathfrak{p}\cap S_{n-mk}=S_{n-mk}$ segundo \sref{II.2.1.8}.
Portanto, se $\mathfrak{p}\cap S_n=S_n$ a partir de certo valor de $n$, ter-se-ia $\mathfrak{p}\supset S_+$, contrariamente à hipótese, o que prova que $3^\circ$ é necessária.
Reciprocamente, suponhamos satisfeitas as condições $1^\circ$, $2^\circ$ e $3^\circ$.
Observemos que, se, para um inteiro $d\geq n_0$, $f\in S_d$ não pertence a $\mathfrak{p}_d$, então, caso $\mathfrak{p}$ exista, $\mathfrak{p}_m$, para $m<n_0$, é necessariamente igual ao conjunto dos $x\in S_m$ tais que $f^rx\in\mathfrak{p}_{m+rd}$ salvo para um número finito de valores de $r$.
Isso já prova que, se $\mathfrak{p}$ existe, é único.
Resta mostrar que, se se definem os $\mathfrak{p}_m$ para $m<n_0$ pela condição precedente, $\mathfrak{p}=\sum_{n=0}^{\infty}\mathfrak{p}_n$ é um ideal primo.
Observemos primeiro que, em virtude de $2^\circ$, para $m\geq n_0$, $\mathfrak{p}_m$ também é definido como o conjunto dos $x\in S_m$ tais que $f^rx\in\mathfrak{p}_{m+rd}$ salvo para um número finito de valores de $r$.
Dito
\oldpage[II]{23}
isso, se $g\in S_m$ e $x\in\mathfrak{p}_n$, tem-se $f^rgx\in\mathfrak{p}_{m+n+rd}$ salvo para um número finito de valores de $r$, logo $gx\in\mathfrak{p}_{m+n}$, o que prova que $\mathfrak{p}$ é um ideal de $S$.
Para estabelecer que esse ideal é primo, em outros termos, que o anel $S/\mathfrak{p}$, graduado pelos subgrupos $S_n/\mathfrak{p}_n$, é um anel íntegro, basta (considerando as componentes de maior grau de dois elementos de $S/\mathfrak{p}$) provar que, se $x\in S_m$ e $y\in S_n$ são tais que $x\notin\mathfrak{p}_m$ e $y\notin\mathfrak{p}_n$, então $xy\notin\mathfrak{p}_{m+n}$.
Caso contrário, para $r$ suficientemente grande, ter-se-ia $f^{2r}xy\in\mathfrak{p}_{m+n+2rd}$; mas $f^ry\notin\mathfrak{p}_{n+rd}$ para todo $r>0$.
Resulta então de $2^\circ$ que, salvo para um número finito de valores de $r$, tem-se $f^rx\in\mathfrak{p}_{m+rd}$, e concluir-se-ia que $x\in\mathfrak{p}_m$, contrariamente à hipótese.
\end{proof}

\begin{env}[2.1.10]
\label{II.2.1.10}
Diz-se que uma parte $\mathfrak{J}$ de $S_+$ é um \emph{ideal de $S_+$} se é um ideal de $S$, e que $\mathfrak{J}$ é um \emph{ideal primo graduado de $S_+$} se é a interseção de $S_+$ com um ideal primo graduado de $S$ \emph{que não contém $S_+$} (esse ideal primo é, aliás, único segundo \sref{II.2.1.9}).
Se $\mathfrak{J}$ é um ideal de $S_+$, o \emph{radical de $\mathfrak{J}$ em $S_+$} é o conjunto dos elementos de $S_+$ que têm alguma potência pertencente a $\mathfrak{J}$, isto é, o conjunto
\[
  \rad_+(\mathfrak{J})=\rad(\mathfrak{J})\cap S_+.
\]
Em particular, o radical de $0$ em $S_+$ é também chamado \emph{nilradical} de $S_+$ e denotado por $\nilrad_+$: é, portanto, o conjunto dos elementos nilpotentes de $S_+$.
Se $\mathfrak{J}$ é um ideal \emph{graduado} de $S_+$, seu radical $\rad_+(\mathfrak{J})$ é um ideal \emph{graduado}: passando ao anel quociente $S/\mathfrak{J}$, pode-se reduzir ao caso $\mathfrak{J}=0$, e tudo se reduz a ver que, se $x=x_h+x_{h+1}+\cdots+x_k$ é nilpotente, o mesmo vale para os $x_i\in S_i$ ($1\leq h\leq i\leq k$).
Pode-se supor $x_k\neq0$, e a componente de maior grau de $x^n$ é então $x_k^n$; logo $x_k$ é nilpotente, e raciocina-se em seguida por indução sobre $k$.
Diz-se que o anel graduado $S$ é \emph{essencialmente reduzido} se $\nilrad_+=0$, em outros termos, se $S_+$ não contém elementos nilpotentes $\neq0$.
\end{env}

\begin{env}[2.1.11]
\label{II.2.1.11}
Observemos que, se, no anel graduado $S$, um elemento $x$ é divisor de zero, o mesmo vale para sua componente de maior grau.
Diz-se que um anel $S$ é \emph{essencialmente íntegro} se o anel $S_+$ (\emph{sem elemento unidade}) não contém divisores de zero e não é nulo; basta, para isso, que nenhum elemento homogêneo não nulo de $S_+$ seja divisor de zero nesse anel.
É claro que, se $\mathfrak{p}$ é um ideal primo graduado de $S_+$, $S/\mathfrak{p}$ é essencialmente íntegro.

Seja $S$ um anel graduado essencialmente íntegro, e seja $x_0\in S_0$; se existe então \emph{um} elemento homogêneo $f\neq0$ de $S_+$ tal que $x_0f=0$, tem-se $x_0S_+=0$, pois $(x_0g)f=(x_0f)g=0$ para todo $g\in S_+$, e a hipótese implica, portanto, $x_0g=0$.
Para que $S$ seja íntegro, é necessário e suficiente que $S_0$ seja íntegro e que o anulador de $S_+$ em $S_0$ se reduza a $0$.
\end{env}

\subsection{Anéis de frações de um anel graduado}
\label{subsection:II.2.2}

\begin{env}[2.2.1]
\label{II.2.2.1}
Seja $S$ um anel graduado em graus positivos, e seja $f$ um elemento \emph{homogêneo} de $S$, de grau $d>0$; então o anel de frações $S'=S_f$ é graduado, tomando como $S_n'$ o conjunto dos $x/f^k$, onde $x\in S_{n+kd}$ e $k\geq0$ (observaremos aqui que $n$ pode assumir valores negativos arbitrários); denotaremos por $S_{(f)}$ o subanel $S_0'=(S_f)_0$ de $S'$ formado pelos elementos \emph{de grau $0$}.

Se $f\in S_d$, os monômios $(f/1)^h$ em $S_f$ ($h$ inteiro positivo ou negativo) formam uma \emph{família linearmente independente} sobre o anel $S_{(f)}$, e o conjunto de suas combinações lineares não é outro senão
\oldpage[II]{24}
o anel $(S^{(d)})_f$, que é, portanto, isomorfo a $S_{(f)}[T,T^{-1}]=S_{(f)}\otimes_\bb{Z}\bb{Z}[T,T^{-1}]$ (onde $T$ é uma indeterminada).
Com efeito, se se tem uma relação
\[
  \sum_{h=-a}^{b}z_h(f/1)^h=0,
\]
com $z_h=x_h/f^m$, onde todos os $x_h$ pertencem a $S_{md}$, essa relação equivale, por definição, à existência de um $k\geq a$ tal que
\[
  \sum_{h=-a}^{b}f^{h+k}x_h=0;
\]
como os graus dos termos dessa soma são distintos, tem-se $f^{h+k}x_h=0$ para todo $h$, de onde $z_h=0$ para todo $h$.

Se $M$ é um $S$-módulo graduado, $M'=M_f$ é um $S_f$-módulo graduado, sendo $M_n'$ o conjunto dos $z/f^k$ com $z\in M_{n+kd}$ ($k\geq0$); denota-se por $M_{(f)}$ o conjunto dos elementos homogêneos de grau $0$ de $M'$.
É imediato que $M_{(f)}$ é um $S_{(f)}$-módulo e que se tem
\[
  (M^{(d)})_f=M_{(f)}\otimes_{S_{(f)}}(S^{(d)})_f.
\]
\end{env}

\begin{lemma}[2.2.2]
\label{II.2.2.2}
Sejam $d$ e $e$ dois inteiros $>0$, $f\in S_d$ e $g\in S_e$.
Existe um isomorfismo canônico de anéis
\[
  S_{(fg)}\isoto(S_{(f)})_{g^d/f^e};
\]
se se identificam canonicamente esses dois anéis, existe um isomorfismo canônico de módulos
\[
  M_{(fg)}\isoto(M_{(f)})_{g^d/f^e}.
\]
\end{lemma}

\begin{proof}
Com efeito, $fg$ divide $f^eg^d$, e este último elemento divide $(fg)^{de}$; logo os anéis graduados $S_{fg}$ e $S_{f^eg^d}$ se identificam canonicamente.
Por outro lado, $S_{f^eg^d}$ também se identifica com $(S_{f^e})_{g^d/1}$ \sref[0]{0.1.4.6}; e, como $f^e/1$ é invertível em $S_{f^e}$, $S_{f^eg^d}$ também se identifica com $(S_{f^e})_{g^d/f^e}$.
Ora, o elemento $g^d/f^e$ é de grau $0$ em $S_{f^e}$; conclui-se imediatamente que o subanel de $(S_{f^e})_{g^d/f^e}$ formado pelos elementos de grau $0$ é $(S_{(f^e)})_{g^d/f^e}$ e, como se tem evidentemente $S_{(f^e)}=S_{(f)}$, isso demonstra a primeira parte do lema; a segunda estabelece-se da mesma maneira.
\end{proof}

\begin{env}[2.2.3]
\label{II.2.2.3}
Sob as hipóteses de \sref{II.2.2.2}, é claro que o homomorfismo canônico $S_f\to S_{fg}$ \sref[0]{0.1.4.1}, que transforma $x/f^k$ em $g^kx/(fg)^k$, é de grau $0$ e, portanto, fornece por restrição um \emph{homomorfismo canônico $S_{(f)}\to S_{(fg)}$}, tal que o diagrama
\[
  \xymatrix{
    & S_{(f)}\ar[dl]\ar[dr]\\
    S_{(fg)}\ar[rr]^-{\sim} & &
    (S_{(f)})_{g^d/f^e}
  }
\]
é comutativo.
Define-se da mesma maneira um homomorfismo canônico $M_{(f)}\to M_{(fg)}$.
\end{env}

\begin{lemma}[2.2.4]
\label{II.2.2.4}
Se $f$ e $g$ são dois elementos homogêneos de $S_+$, o anel $S_{(fg)}$ é gerado pela reunião das imagens canônicas de $S_{(f)}$ e $S_{(g)}$.
\end{lemma}

\begin{proof}
Em virtude de \sref{II.2.2.2}, basta ver que
\[
  1/(g^d/f^e)=f^{d+e}/(fg)^d
\]
pertence à imagem canônica de $S_{(g)}$ em $S_{(fg)}$, o que é evidente por definição.
\end{proof}

\begin{proposition}[2.2.5]
\label{II.2.2.5}
Seja $d$ um inteiro $>0$ e seja $f\in S_d$.
Então existe um isomorfismo canônico de anéis
\[
  S_{(f)}\isoto S^{(d)}/(f-1)S^{(d)};
\]
se se identificam esses dois anéis por esse isomorfismo, existe um isomorfismo canônico de módulos
\[
  M_{(f)}\isoto M^{(d)}/(f-1)M^{(d)}.
\]
\end{proposition}

\begin{proof}
O primeiro desses isomorfismos define-se fazendo corresponder a $x/f^n$, onde $x\in S_{nd}$, o elemento $\overline{x}$, classe de $x$ módulo $(f-1)S^{(d)}$; essa aplicação está bem definida, pois se tem a congruência
\[
  f^hx\equiv x\pmod{(f-1)S^{(d)}}
\]
para todo $x\in S^{(d)}$; portanto, se $f^hx=0$ para algum $h>0$,
\oldpage[II]{25}
tem-se $\overline{x}=0$.
Por outro lado, se $x\in S_{nd}$ é tal que $x=(f-1)y$, com
$y=y_{hd}+y_{(h+1)d}+\cdots+y_{kd}$, onde $y_{jd}\in S_{jd}$ e $y_{hd}\neq0$, tem-se necessariamente $h=n$ e $x=-y_{hd}$, assim como as relações $y_{(j+1)d}=fy_{jd}$ para $h\leq j\leq k-1$ e $fy_{kd}=0$, o que dá finalmente $f^{k-n+1}x=0$.
A toda classe $\overline{x}\text{ mod. }(f-1)S^{(d)}$ de um elemento $x\in S_{nd}$ pode-se, portanto, fazer corresponder o elemento $x/f^n$ de $S_{(f)}$, pois a observação precedente mostra que essa aplicação está bem definida.
É imediato que as duas aplicações assim definidas são homomorfismos de anéis, inversos um do outro.
Procede-se exatamente da mesma maneira para $M$.
\end{proof}

\begin{corollary}[2.2.6]
\label{II.2.2.6}
Se $S$ é noetheriano, o mesmo vale para $S_{(f)}$ para todo $f$ homogêneo de grau $>0$.
\end{corollary}

\begin{proof}
Isso decorre imediatamente de \sref{II.2.1.7} e \sref{II.2.2.5}.
\end{proof}

\begin{env}[2.2.7]
\label{II.2.2.7}
Seja $T$ um subconjunto multiplicativo de $S_+$ formado por elementos \emph{homogêneos}; então $T_0=T\cup\{1\}$ é um subconjunto multiplicativo de $S$.
Como os elementos de $T_0$ são homogêneos, o anel $T_0^{-1}S$ continua graduado de maneira evidente; denotaremos por $S_{(T)}$ o subanel de $T_0^{-1}S$ formado pelos elementos de grau $0$, isto é, pelos elementos da forma $x/h$, onde $h\in T$ e $x$ é homogêneo de grau igual ao de $h$.
Sabe-se por \sref[0]{0.1.4.5} que $T_0^{-1}S$ se identifica canonicamente com o limite indutivo dos anéis $S_f$, onde $f$ percorre $T$ (relativamente aos homomorfismos canônicos $S_f\to S_{fg}$); como essa identificação respeita os graus, ela identifica $S_{(T)}$ com o \emph{limite indutivo} dos $S_{(f)}$ para $f\in T$.
Para todo $S$-módulo graduado $M$, define-se do mesmo modo o módulo $M_{(T)}$ (sobre o anel $S_{(T)}$), formado pelos elementos de grau $0$ de $T_0^{-1}M$, e vê-se que esse módulo é o limite indutivo dos $M_{(f)}$ para $f\in T$.

Se $\mathfrak{p}$ é um ideal primo graduado de $S_+$, denotaremos respectivamente por $S_{(\mathfrak{p})}$ e $M_{(\mathfrak{p})}$ o anel $S_{(T)}$ e o módulo $M_{(T)}$, onde $T$ é o conjunto dos elementos \emph{homogêneos} de $S_+$ que não pertencem a $\mathfrak{p}$.
\end{env}

\subsection{Espectro primo homogêneo de um anel graduado}
\label{subsection:II.2.3}

\begin{env}[2.3.1]
\label{II.2.3.1}
Dado um anel $S$ graduado em graus positivos, chama-se \emph{espectro primo homogêneo} de $S$, e denota-se por $\Proj(S)$, o conjunto dos ideais primos graduados de $S_+$ \sref{II.2.1.10}, ou, o que é equivalente, o conjunto dos ideais primos graduados de $S$ \emph{que não contêm $S_+$}; definiremos uma estrutura de \emph{esquema} cujo conjunto subjacente é $\Proj(S)$.
\end{env}

\begin{env}[2.3.2]
\label{II.2.3.2}
Para todo subconjunto $E$ de $S$, seja $V_+(E)$ o conjunto dos ideais primos graduados de $S$ que contêm $E$ e não contêm $S_+$; trata-se, portanto, do subconjunto $V(E)\cap\Proj(S)$ de $\Spec(S)$.
De \sref[I]{I.1.1.2} deduz-se:
\[
\label{II.2.3.2.1}
  V_+(0)=\Proj(S),\qquad V_+(S)=V_+(S_+)=\emp,
\tag{2.3.2.1}
\]
\[
\label{II.2.3.2.2}
  V_+\big(\textstyle\bigcup_\lambda E_\lambda\big)
  =\textstyle\bigcap_\lambda V_+(E_\lambda),
\tag{2.3.2.2}
\]
\[
\label{II.2.3.2.3}
  V_+(EE')=V_+(E)\cup V_+(E').
\tag{2.3.2.3}
\]

Não se altera $V_+(E)$ substituindo $E$ pelo ideal graduado gerado por $E$; além disso, se $\mathfrak{J}$ é um ideal graduado de $S$, tem-se
\[
\label{II.2.3.2.4}
  V_+(\mathfrak{J})
  =V_+\big(\textstyle\bigcup_{q\geq n}(\mathfrak{J}\cap S_q)\big)
\tag{2.3.2.4}
\]
\oldpage[II]{26}
para todo $n>0$.
Com efeito, se $\mathfrak{p}\in\Proj(S)$ contém os elementos homogêneos de $\mathfrak{J}$ de grau $\geq n$, como, por hipótese, existe um elemento homogêneo $f\in S_d$ não contido em $\mathfrak{p}$, então, para todo $m\geq0$ e todo $x\in S_m\cap\mathfrak{J}$, tem-se $f^rx\in\mathfrak{J}\cap S_{m+rd}$ para todo $r\geq0$; salvo para um número finito de valores de $r$, tem-se $m+rd\geq n$ e, portanto, $f^rx\in\mathfrak{p}\cap S_{m+rd}$, o que implica $x\in\mathfrak{p}\cap S_m$ \sref{II.2.1.8}.

Por fim, para todo ideal graduado $\mathfrak{J}$ de $S$, tem-se
\[
  V_+(\mathfrak{J})=V_+(\rad_+(\mathfrak{J})).
\tag{2.3.2.5}
\]
\end{env}

\begin{env}[2.3.3]
\label{II.2.3.3}
Por definição, os $V_+(E)$ são os subconjuntos fechados de $X=\Proj(S)$ para a topologia induzida pela topologia espectral de $\Spec(S)$, que também chamaremos de \emph{topologia espectral} em $X$.
Para todo $f\in S$, põe-se
\[
\label{II.2.3.3.1}
  D_+(f)=D(f)\cap\Proj(S)=\Proj(S)\setmin V_+(f)
\tag{2.3.3.1}
\]
e, por conseguinte, para quaisquer dois elementos $f$, $g$ de $S$ \sref[I]{I.1.1.9.1},
\[
\label{II.2.3.3.2}
  D_+(fg)=D_+(f)\cap D_+(g).
\tag{2.3.3.2}
\]
\end{env}

\begin{proposition}[2.3.4]
\label{II.2.3.4}
Quando $f$ percorre o conjunto dos elementos homogêneos de $S_+$, os $D_+(f)$ formam uma base da topologia de $X=\Proj(S)$.
\end{proposition}

\begin{proof}
Com efeito, decorre de \sref{II.2.3.2.2} e \sref{II.2.3.2.4} que todo subconjunto fechado de $X$ é interseção de conjuntos da forma $V_+(f)$, onde $f$ é homogêneo de grau $>0$.
\end{proof}

\begin{env}[2.3.5]
\label{II.2.3.5}
Seja $f$ um elemento \emph{homogêneo} de $S_+$, de grau $d>0$.
Para todo ideal primo graduado $\mathfrak{p}$ de $S$ que não contém $f$, sabe-se que o conjunto dos $x/f^n$, onde $x\in\mathfrak{p}$ e $n\geq0$, é um ideal primo do anel de frações $S_f$ \sref[0]{0.1.2.6}; sua interseção com $S_{(f)}$ é, portanto, um ideal primo desse anel, que denotaremos por $\psi_f(\mathfrak{p})$: é o conjunto dos $x/f^n$ para $n\geq0$ e $x\in\mathfrak{p}\cap S_{nd}$.
Define-se assim uma aplicação
\[
  \psi_f:D_+(f)\to\Spec(S_{(f)});
\]
Além disso, se $g\in S_e$ é um segundo elemento homogêneo de $S_+$, tem-se um diagrama comutativo
\[
\label{II.2.3.5.1}
  \xymatrix{
    D_+(f) \ar[r]^{\psi_f}
    & \Spec(S_{(f)})
  \\D_+(fg) \ar[u] \ar[r]_{\psi_{fg}}
    & \Spec(S_{(fg)}) \ar[u]
  }
\tag{2.3.5.1}
\]
onde a seta vertical da esquerda é a inclusão, e a da direita é a aplicação ${}^a\!\omega_{fg,f}$ deduzida do homomorfismo canônico $\omega=\omega_{fg,f}:S_{(f)}\to S_{(fg)}$ \sref[I]{I.1.2.1}.
Com efeito, se $x/f^n\in\omega^{-1}(\psi_{fg}(\mathfrak{p}))$, com $fg\notin\mathfrak{p}$, então, por definição, $g^nx/(fg)^n\in\psi_{fg}(\mathfrak{p})$, logo $g^nx\in\mathfrak{p}$ e, portanto, $x\in\mathfrak{p}$; a recíproca é evidente.
\end{env}

\begin{proposition}[2.3.6]
\label{II.2.3.6}
A aplicação $\psi_f$ é um homeomorfismo de $D_+(f)$ sobre $\Spec(S_{(f)})$.
\end{proposition}

\begin{proof}
Em primeiro lugar, $\psi_f$ é contínua; com efeito, se $h\in S_{nd}$, então $h/f^n\in\psi_f(\mathfrak{p})$ se e somente se $h\in\mathfrak{p}$, e, portanto, $\psi_f^{-1}(D(h/f^n))=D_+(hf)$; nossa afirmação decorre então de \sref{II.2.3.3.2}.
Além disso, os $D_+(hf)$, onde $h$ percorre os conjuntos $S_{nd}$, formam uma base da topologia de $D_+(f)$, em virtude de \sref{II.2.3.4} e \sref{II.2.3.3.2}.
O
\oldpage[II]{27}
precedente demonstra, portanto, levando em conta o axioma $(T_0)$, válido em $D_+(f)$ e em $\Spec(S_{(f)})$, que $\psi_f$ é injetiva e que a aplicação inversa $\psi_f(D_+(f))\to D_+(f)$ é contínua.
Por fim, para ver que $\psi_f$ é sobrejetiva, observa-se que, se $\mathfrak{q}_0$ é um ideal primo de $S_{(f)}$ e, para todo $n>0$, denota-se por $\mathfrak{p}_n$ o conjunto dos $x\in S_n$ tais que $x^d/f^n\in\mathfrak{q}_0$, então os $\mathfrak{p}_n$ verificam as condições de \sref{II.2.1.9}.
Com efeito, se $x,y\in S_n$ são tais que $x^d/f^n,y^d/f^n\in\mathfrak{q}_0$, tem-se $(x+y)^{2d}/f^{2n}\in\mathfrak{q}_0$, de onde $(x+y)^d/f^n\in\mathfrak{q}_0$, pois $\mathfrak{q}_0$ é primo.
Isso demonstra que os $\mathfrak{p}_n$ são subgrupos dos $S_n$, e a verificação das demais condições de \sref{II.2.1.9} é imediata, levando em conta que $\mathfrak{q}_0$ é primo.
Se $\mathfrak{p}$ é o ideal primo graduado de $S$ assim definido, tem-se de fato $\psi_f(\mathfrak{p})=\mathfrak{q}_0$, pois, para $x\in S_{nd}$, as relações $x/f^n\in\mathfrak{q}_0$ e $x^d/f^{nd}\in\mathfrak{q}_0$ são equivalentes, visto que $\mathfrak{q}_0$ é primo.
\end{proof}

\begin{corollary}[2.3.7]
\label{II.2.3.7}
Para que $D_+(f)=\emp$, é necessário e suficiente que $f$ seja nilpotente.
\end{corollary}

\begin{proof}
Com efeito, para que $\Spec(S_{(f)})=\emp$, é necessário e suficiente que $S_{(f)}=0$, ou ainda que $1=0$ em $S_f$, o que significa, por definição, que $f$ é nilpotente.
\end{proof}

\begin{corollary}[2.3.8]
\label{II.2.3.8}
Seja $E$ um subconjunto de $S_+$.
As condições seguintes são equivalentes:
\begin{enumerate}
  \item[\rm(a)] $V_+(E)=X=\Proj(S)$.
  \item[\rm(b)] Todo elemento de $E$ é nilpotente.
  \item[\rm(c)] As componentes homogêneas de todo elemento de $E$ são nilpotentes.
\end{enumerate}
\end{corollary}

\begin{proof}
É claro que (c) implica (b), e que (b) implica (a).
Se $\mathfrak{J}$ é o ideal graduado de $S$ gerado por $E$, a condição (a) equivale a $V_+(\mathfrak{J})=X$; \emph{a fortiori}, (a) implica que todo elemento homogêneo $f\in\mathfrak{J}$ é tal que $V_+(f)=X$, logo $f$ é nilpotente em virtude de \sref{II.2.3.7}.
\end{proof}

\begin{corollary}[2.3.9]
\label{II.2.3.9}
Se $\mathfrak{J}$ é um ideal graduado de $S_+$, então $\rad_+(\mathfrak{J})$ é a interseção dos ideais primos graduados de $S_+$ que contêm $\mathfrak{J}$.
\end{corollary}

\begin{proof}
Considerando o anel graduado $S/\mathfrak{J}$, pode-se reduzir ao caso em que $\mathfrak{J}=0$.
É preciso provar que, se $f\in S_+$ não é nilpotente, existe um ideal primo graduado de $S$ que não contém $f$; ora, pelo menos uma das componentes homogêneas de $f$ não é nilpotente e, portanto, pode-se supor $f$ homogêneo.
A afirmação decorre então de \sref{II.2.3.7}.
\end{proof}

\begin{env}[2.3.10]
\label{II.2.3.10}
Para todo subconjunto $Y$ de $X=\Proj(S)$, seja $\mathfrak{j}_+(Y)$ o conjunto dos $f\in S_+$ tais que $Y\subset V_+(f)$; isso equivale a dizer que $\mathfrak{j}_+(Y)=\mathfrak{j}(Y)\cap S_+$.
Assim, $\mathfrak{j}_+(Y)$ é um ideal de $S_+$ igual ao seu radical em $S_+$.
\end{env}

\begin{proposition}[2.3.11]
\label{II.2.3.11}
\begin{enumerate}
  \item[\rm(i)] Para todo subconjunto $E$ de $S_+$, $\mathfrak{j}_+(V_+(E))$ é o radical em $S_+$ do ideal graduado de $S_+$ gerado por $E$.
  \item[\rm(ii)] Para todo subconjunto $Y$ de $X$, $V_+(\mathfrak{j}_+(Y))=\overline{Y}$, onde $\overline{Y}$ é o fecho de $Y$ em $X$.
\end{enumerate}
\end{proposition}

\begin{proof}
\medskip\noindent
\begin{enumerate}
  \item[\rm(i)] Se $\mathfrak{J}$ é o ideal graduado de $S_+$ gerado por $E$, tem-se $V_+(E)=V_+(\mathfrak{J})$, e a afirmação decorre então de \sref{II.2.3.9}.
  \item[\rm(ii)] Como $V_+(\mathfrak{J})=\bigcap_{f\in\mathfrak{J}}V_+(f)$, a relação $Y\subset V_+(\mathfrak{J})$ implica $Y\subset V_+(f)$ para todo $f\in\mathfrak{J}$ e, por conseguinte, $\mathfrak{j}_+(Y)\supset\mathfrak{J}$, de onde $V_+(\mathfrak{j}_+(Y))\subset V_+(\mathfrak{J})$; isso demonstra (ii) pela definição dos subconjuntos fechados.
\end{enumerate}
\end{proof}

\begin{corollary}[2.3.12]
\label{II.2.3.12}
Os subconjuntos fechados $Y$ de $X=\Proj(S)$ e os ideais graduados de $S_+$ iguais ao seu radical em $S_+$ correspondem-se bijetivamente pelas aplicações que invertem inclusões $Y\mapsto\mathfrak{j}_+(Y)$ e $\mathfrak{J}\mapsto V_+(\mathfrak{J})$; à união $Y_1\cup Y_2$ de dois subconjuntos fechados de $X$ corresponde
\oldpage[II]{28}
$\mathfrak{j}_+(Y_1)\cap\mathfrak{j}_+(Y_2)$, e à interseção de uma família arbitrária $(Y_\lambda)$ de subconjuntos fechados corresponde o radical em $S_+$ da soma dos $\mathfrak{j}_+(Y_\lambda)$.
\end{corollary}

\begin{corollary}[2.3.13]
\label{II.2.3.13}
Seja $\mathfrak{J}$ um ideal graduado de $S_+$.
Para que $V_+(\mathfrak{J})=\emp$, é necessário e suficiente que todo elemento de $S_+$ tenha uma potência em $\mathfrak{J}$.
\end{corollary}

Este último corolário também se exprime por uma das formas equivalentes seguintes:

\begin{corollary}[2.3.14]
\label{II.2.3.14}
Seja $(f_\alpha)$ uma família de elementos homogêneos de $S_+$.
Para que os $D_+(f_\alpha)$ formem um recobrimento de $X=\Proj(S)$, é necessário e suficiente que todo elemento de $S_+$ tenha uma potência no ideal gerado pelos $f_\alpha$.
\end{corollary}

\begin{corollary}[2.3.15]
\label{II.2.3.15}
Sejam $(f_\alpha)$ uma família de elementos homogêneos de $S_+$ e $f$ um elemento de $S_+$.
As relações seguintes são equivalentes:
\begin{enumerate}
  \item[\rm(a)] $D_+(f)\subset\bigcup_\alpha D_+(f_\alpha)$;
  \item[\rm(b)] $V_+(f)\supset\bigcap_\alpha V_+(f_\alpha)$;
  \item[\rm(c)] uma potência de $f$ pertence ao ideal gerado pelos $f_\alpha$.
\end{enumerate}
\end{corollary}

\begin{corollary}[2.3.16]
\label{II.2.3.16}
Para que $X=\Proj(S)$ seja vazio, é necessário e suficiente que todo elemento de $S_+$ seja nilpotente.
\end{corollary}

\begin{corollary}[2.3.17]
\label{II.2.3.17}
Na correspondência bijetiva descrita em \sref{II.2.3.12}, aos subconjuntos fechados \emph{irredutíveis} de $X$ correspondem os ideais graduados \emph{primos} de $S_+$.
\end{corollary}

\begin{proof}
Com efeito, se $Y=Y_1\cup Y_2$, onde $Y_1$ e $Y_2$ são fechados e distintos de $Y$, tem-se
\[
  \mathfrak{j}_+(Y)=\mathfrak{j}_+(Y_1)\cap\mathfrak{j}_+(Y_2)
\]
os ideais $\mathfrak{j}_+(Y_1)$ e $\mathfrak{j}_+(Y_2)$ sendo distintos de $\mathfrak{j}_+(Y)$; portanto, $\mathfrak{j}_+(Y)$ não é primo.
Reciprocamente, se $\mathfrak{J}$ é um ideal graduado não primo de $S_+$, existem dois elementos $f$, $g$ de $S_+$ tais que $f\notin\mathfrak{J}$, $g\notin\mathfrak{J}$ e $fg\in\mathfrak{J}$.
Então $V_+(f)\not\supset V_+(\mathfrak{J})$ e $V_+(g)\not\supset V_+(\mathfrak{J})$, mas $V_+(\mathfrak{J})\subset V_+(f)\cup V_+(g)$ por \sref{II.2.3.2.3}; conclui-se que $V_+(\mathfrak{J})$ é a união dos subconjuntos fechados $V_+(f)\cap V_+(\mathfrak{J})$ e $V_+(g)\cap V_+(\mathfrak{J})$, que são distintos de $V_+(\mathfrak{J})$.
\end{proof}

\subsection{A estrutura de esquema em $\Proj(S)$}
\label{subsection:II.2.4}

\begin{env}[2.4.1]
\label{II.2.4.1}
Sejam $f$, $g$ dois elementos homogêneos de $S_+$; consideremos os esquemas afins $Y_f=\Spec(S_{(f)})$, $Y_g=\Spec(S_{(g)})$ e $Y_{fg}=\Spec(S_{(fg)})$.
Em virtude de \sref{II.2.2.2}, o morfismo $w_{fg,f}=({}^a\!\omega_{fg,f},\widetilde{\omega}_{fg,f})$ de $Y_{fg}$ em $Y_f$, correspondente ao homomorfismo canônico $\omega_{fg,f}:S_{(f)}\to S_{(fg)}$, é uma \emph{imersão aberta} \sref[I]{I.1.3.6}.
Por meio do homeomorfismo inverso de $\psi_f:D_+(f)\to Y_f$ \sref{II.2.3.6}, pode-se transportar para $D_+(f)$ a estrutura de esquema afim de $Y_f$; em virtude da comutatividade do diagrama \sref{II.2.3.5.1}, o esquema afim $D_+(fg)$ identifica-se assim com o esquema induzido no subconjunto aberto $D_+(fg)$ do espaço subjacente ao esquema afim $D_+(f)$.
Fica claro então, levando em conta \sref{II.2.3.4}, que $X=\Proj(S)$ é munido de uma única estrutura de \emph{pré-esquema}, cuja restrição a cada $D_+(f)$ é o esquema afim que acaba de ser definido.
Além disso:
\end{env}

\begin{proposition}[2.4.2]
\label{II.2.4.2}
O pré-esquema $\Proj(S)$ é um esquema.
\end{proposition}

\begin{proof}
Basta \sref[I]{I.5.5.6} demonstrar que, quaisquer que sejam $f$, $g$ homogêneos em $S_+$, $D_+(f)\cap D_+(g)=D_+(fg)$ é afim e seu anel é gerado pelas imagens canônicas dos anéis de $D_+(f)$ e $D_+(g)$; o primeiro ponto é evidente por definição, e o segundo decorre de \sref{II.2.2.4}.
\end{proof}
\oldpage[II]{29}
Quando se falar do espectro primo homogêneo $\Proj(S)$ como um \emph{esquema}, tratar-se-á sempre da estrutura que acaba de ser definida.

\begin{example}[2.4.3]
\label{II.2.4.3}
Tomemos $S=K[T_1,T_2]$, sendo $K$ um corpo, $T_1$, $T_2$ duas indeterminadas, e sendo $S$ graduado pelo grau total.
Decorre de \sref{II.2.3.14} que $\Proj(S)$ é a união de $D_+(T_1)$ e $D_+(T_2)$; vê-se imediatamente que esses esquemas afins são canonicamente isomorfos a $\Spec(K[T])$, e que $\Proj(S)$ se obtém pela colagem desses dois esquemas afins descrita em \sref[I]{I.2.3.2} (cf. \sref{II.7.4.14}).
\end{example}

\begin{proposition}[2.4.4]
\label{II.2.4.4}
Seja $S$ um anel graduado em graus positivos, e seja $X$ o esquema $\Proj(S)$.
\begin{enumerate}
  \item[\rm(i)] Se $\mathfrak{N}_+$ é o nilradical de $S_+$ \sref{II.2.1.10}, então o esquema $X_\red$ é canonicamente isomorfo a $\Proj(S/\mathfrak{N}_+)$; em particular, se $S$ é essencialmente reduzido, $\Proj(S)$ é reduzido.
  \item[\rm(ii)] Suponhamos $S$ essencialmente reduzido. Então, para que $X$ seja integral, é necessário e suficiente que $S$ seja essencialmente íntegro.
\end{enumerate}
\end{proposition}

\begin{proof}
\medskip\noindent
\begin{enumerate}
  \item[\rm(i)] Seja $\overline{S}$ o anel graduado $S/\mathfrak{N}_+$, e denotemos por $x\mapsto\overline{x}$ o homomorfismo canônico $S\to\overline{S}$, de grau $0$.
  Para todo $f\in S_d$ ($d>0$), o homomorfismo canônico $S_f\to\overline{S}_{\overline{f}}$ \sref[0]{0.1.5.1} é sobrejetivo e de grau $0$ e, portanto, por restrição, dá um homomorfismo sobrejetivo $S_{(f)}\to\overline{S}_{(\overline{f})}$.
  Se se supõe $f\notin\mathfrak{N}_+$, verifica-se imediatamente que $\overline{S}_{(\overline{f})}$ é reduzido e que o núcleo do homomorfismo precedente é o nilradical de $S_{(f)}$, isto é,
  \[
    \overline{S}_{(\overline{f})}=(S_{(f)})_\red.
  \]
  A esse homomorfismo corresponde, portanto, uma imersão fechada $D_+(\overline{f})\to D_+(f)$ que identifica $D_+(\overline{f})$ com $(D_+(f))_\red$ \sref[I]{I.5.1.2} e que é, em particular, um homeomorfismo dos espaços subjacentes a esses dois esquemas afins.
  Além disso, se $g\notin\mathfrak{N}_+$ é um segundo elemento homogêneo de $S_+$, o diagrama
  \[
    \xymatrix{
      S_{(f)} \ar[r] \ar[d]
      & \overline{S}_{(\overline{f})} \ar[d]
    \\S_{(fg)} \ar[r]
      & \overline{S}_{(\overline{fg})}
    }
  \]
  é comutativo.
  Como, além disso, os $D_+(f)$, para $f$ homogêneo de grau $>0$ e $f\notin\mathfrak{N}_+$, formam um recobrimento de $X=\Proj(S)$ \sref{II.2.3.7}, vê-se que os morfismos $D_+(\overline{f})\to D_+(f)$ são as restrições de uma imersão fechada $\Proj(\overline{S})\to\Proj(S)$ que é um homeomorfismo dos espaços subjacentes; daí a conclusão \sref[I]{I.5.1.2}.
  \item[\rm(ii)] Suponhamos que $S$ seja essencialmente íntegro, isto é, que $(0)$ seja um ideal primo graduado de $S_+$ distinto de $S_+$; então $X$ é reduzido em virtude de (i) e irredutível em virtude de \sref{II.2.3.17}.
  Reciprocamente, suponhamos que $S$ seja essencialmente reduzido e $X$ integral; então, para $f\neq0$ homogêneo em $S_+$, tem-se $D_+(f)\neq\emp$ \sref{II.2.3.7}.
  A hipótese de que $X$ é irredutível implica que $D_+(f)\cap D_+(g)\neq\emp$ para $f$, $g$ homogêneos e $\neq0$ em $S_+$; portanto, $fg\neq0$ em virtude de \sref{II.2.3.3.2}, e conclui-se que $S_+$ não tem divisores de zero, de onde resulta a primeira afirmação.
\end{enumerate}
\end{proof}

\begin{env}[2.4.5]
\label{II.2.4.5}
Dado um anel comutativo $A$, recordemos que se diz que um anel graduado $S$ é uma \emph{$A$-álgebra graduada} se é munido de uma estrutura de $A$-álgebra tal que cada um de seus subgrupos $S_n$ seja um $A$-módulo; para isso, basta, aliás, que $S_0$ seja
\oldpage[II]{30}
uma $A$-álgebra; dito de outro modo, define-se em $S$ uma estrutura de $A$-álgebra graduada definindo uma estrutura de $A$-álgebra em $S_0$ e pondo $\alpha\cdot x=(\alpha\cdot1)x$ para $\alpha\in A$ e $x\in S_n$.
\end{env}

\begin{proposition}[2.4.6]
\label{II.2.4.6}
Suponhamos que $S$ seja uma $A$-álgebra graduada.
Então, em $X=\Proj(S)$, o feixe estrutural $\sh{O}_X$ é um feixe de $A$-álgebras (considerando-se $A$ como feixe constante em $X$); dito de outro modo, $X$ é um esquema sobre $\Spec(A)$.
\end{proposition}

\begin{proof}
Basta observar que, para todo $f$ homogêneo em $S_+$, $S_{(f)}$ é uma álgebra sobre $A$, e que o diagrama
\[
  \xymatrix{
    S_{(f)} \ar[rr] && S_{(fg)}
  \\ & A \ar[ul] \ar[ur] &
  }
\]
é comutativo para $f,g$ homogêneos em $S_+$.
\end{proof}

\begin{proposition}[2.4.7]
\label{II.2.4.7}
Seja $S$ um anel graduado em graus positivos.
\begin{enumerate}
  \item[\rm(i)] Para todo inteiro $d>0$, existe um isomorfismo canônico do esquema $\Proj(S)$ sobre o esquema $\Proj(S^{(d)})$.
  \item[\rm(ii)] Seja $S'$ o anel graduado tal que $S'_0=\bb{Z}$ e $S'_n=S_n$ (considerado como $\bb{Z}$-módulo) para $n>0$. Existe um isomorfismo canônico do esquema $\Proj(S)$ sobre o esquema $\Proj(S')$.
\end{enumerate}
\end{proposition}

\begin{proof}
\medskip\noindent
\begin{enumerate}
  \item[\rm(i)] Mostremos primeiro que a aplicação $\mathfrak{p}\mapsto\mathfrak{p}\cap S^{(d)}$ é uma bijeção do conjunto $\Proj(S)$ sobre o conjunto $\Proj(S^{(d)})$.
  Com efeito, suponhamos dado um ideal primo graduado $\mathfrak{p}'\in\Proj(S^{(d)})$ e ponhamos $\mathfrak{p}_{nd}=\mathfrak{p}'\cap S_{nd}$ ($n\geq0$).
  Para todo $n>0$ que não seja múltiplo de $d$, definamos $\mathfrak{p}_n$ como o conjunto dos $x\in S_n$ tais que $x^d\in\mathfrak{p}_{nd}$.
  Se $x,y\in\mathfrak{p}_n$, tem-se $(x+y)^{2d}\in\mathfrak{p}_{2nd}$ e, portanto, $(x+y)^d\in\mathfrak{p}_{nd}$, pois $\mathfrak{p}'$ é primo.
  É imediato que os $\mathfrak{p}_n$ assim definidos para todo $n\geq0$ satisfazem as condições de \sref{II.2.1.9}; existe, portanto, um único ideal primo $\mathfrak{p}\in\Proj(S)$ tal que $\mathfrak{p}\cap S^{(d)}=\mathfrak{p}'$.
  Como, para todo $f$ homogêneo em $S_+$, se tem $V_+(f)=V_+(f^d)$ \sref{II.2.3.2.3}, vê-se que a bijeção precedente é um \emph{homeomorfismo} de espaços topológicos.
  Finalmente, com as mesmas notações, $S_{(f)}$ e $S^{(d)}_{(f^d)}$ identificam-se canonicamente \sref{II.2.2.2}; logo, $\Proj(S)$ e $\Proj(S^{(d)})$ identificam-se canonicamente como \emph{esquemas}.
  \item[\rm(ii)] Se, a cada $\mathfrak{p}\in\Proj(S)$, associarmos o único ideal primo $\mathfrak{p}'\in\Proj(S')$ tal que $\mathfrak{p}'\cap S_n=\mathfrak{p}\cap S_n$ para todo $n>0$ \sref{II.2.1.9}, é claro que se define um homeomorfismo canônico $\Proj(S)\simto\Proj(S')$ dos espaços subjacentes, pois $V_+(f)$ é o mesmo conjunto para $S$ e $S'$ quando $f$ é um elemento homogêneo de $S_+$.
  Como, além disso, $S_{(f)}=S'_{(f)}$, $\Proj(S)$ e $\Proj(S')$ identificam-se como \emph{esquemas}.
\end{enumerate}
\end{proof}

\begin{corollary}[2.4.8]
\label{II.2.4.8}
Se $S$ é uma $A$-álgebra graduada e $S'_A$ é a $A$-álgebra graduada tal que $(S'_A)_0=A$ e $(S'_A)_n=S_n$ para $n>0$, existe um isomorfismo canônico de $\Proj(S)$ sobre $\Proj(S'_A)$.
\end{corollary}

\begin{proof}
Com efeito, esses dois esquemas são canonicamente isomorfos a $\Proj(S')$, com as notações de \sref{II.2.4.7}[(ii)].
\end{proof}

\subsection{O feixe associado a um módulo graduado}
\label{subsection:II.2.5}

\begin{env}[2.5.1]
\label{II.2.5.1}
Seja $M$ um $S$-módulo \emph{graduado}.
Para todo $f$ homogêneo em $S_+$, $M_{(f)}$ é um $S_{(f)}$-módulo e, portanto, corresponde-lhe um feixe quase-coerente associado $(M_{(f)})\supertilde$ no esquema afim $\Spec(S_{(f)})$, identificado com $D_+(f)$ \sref[I]{I.1.3.4}.
\end{env}
\oldpage[II]{31}
\begin{proposition}[2.5.2]
\label{II.2.5.2}
Existe em $X=\Proj(S)$ um único $\sh{O}_X$-módulo quase-coerente $\widetilde{M}$ tal que, para todo $f$ homogêneo em $S_+$, se tenha $\Gamma(D_+(f),\widetilde{M})=M_{(f)}$, sendo o homomorfismo de restrição $\Gamma(D_+(f),\widetilde{M})\to\Gamma(D_+(fg),\widetilde{M})$, para $f,g$ homogêneos em $S_+$, o homomorfismo canônico $M_{(f)}\to M_{(fg)}$ \sref{II.2.2.3}.
\end{proposition}

\begin{proof}
Suponhamos $f\in S_d$ e $g\in S_e$.
Como $D_+(fg)$ se identifica com o espectro primo de $(S_{(f)})_{g^d/f^e}$ em virtude de \sref{II.2.2.2}, a restrição a $D_+(fg)$ do feixe $(M_{(f)})\supertilde$ em $D_+(f)$ identifica-se canonicamente com o feixe associado ao módulo $(M_{(f)})_{g^d/f^e}$ \sref[I]{I.1.3.6} e, portanto, também com $(M_{(fg)})\supertilde$ \sref{II.2.2.2}.
Conclui-se daí que existe um isomorfismo canônico
\[
  \theta_{g,f}:(M_{(f)})\supertilde|D_+(fg)
  \simto(M_{(g)})\supertilde|D_+(fg)
\]
tal que, se $h$ é um terceiro elemento homogêneo de $S_+$, tem-se $\theta_{f,h}=\theta_{f,g}\circ\theta_{g,h}$ em $D_+(fgh)$.
Por conseguinte \sref[0]{0.3.3.1}, existe em $X$ um $\sh{O}_X$-módulo quase-coerente $\sh{F}$ e, para cada $f$ homogêneo em $S_+$, um isomorfismo $\eta_f$ de $\sh{F}|D_+(f)$ sobre $(M_{(f)})\supertilde$ tal que $\theta_{g,f}=\eta_g\circ\eta_f^{-1}$.
Se considerarmos então o feixe $\sh{G}$ associado ao pré-feixe (na base da topologia de $X$ formada pelos $D_+(f)$) definido por $D_+(f)\mapsto M_{(f)}$, tomando como homomorfismos de restrição os homomorfismos canônicos $M_{(f)}\to M_{(fg)}$, o que precede prova que $\sh{F}$ e $\sh{G}$ são isomorfos (levando em conta \sref[I]{I.1.3.7}); o feixe $\sh{G}$ será denotado por $\widetilde{M}$ e satisfaz de fato as condições do enunciado.
Tem-se, em particular, $\widetilde{S}=\sh{O}_X$.
\end{proof}

\begin{definition}[2.5.3]
\label{II.2.5.3}
Diz-se que o $\sh{O}_X$-módulo quase-coerente $\widetilde{M}$ definido em \sref{II.2.5.2} é \emph{associado} ao $S$-módulo graduado $M$.
\end{definition}

Recordemos que os $S$-módulos graduados formam uma categoria quando se restringem os homomorfismos de módulos graduados aos homomorfismos \emph{de grau $0$}.
Com essa convenção:

\begin{proposition}[2.5.4]
\label{II.2.5.4}
O funtor $M\mapsto\widetilde{M}$ é um funtor covariante aditivo exato da categoria dos $S$-módulos graduados para a dos $\sh{O}_X$-módulos quase-coerentes, e comuta com os limites indutivos e as somas diretas.
\end{proposition}

\begin{proof}
Com efeito, como essas propriedades são locais, basta verificá-las nos feixes da forma $\widetilde{M}|D_+(f)=(M_{(f)})\supertilde$.
Ora, o funtor de localização $M\mapsto M_f$, o funtor de grau zero $N\mapsto N_0$ (na categoria dos $S_f$-módulos graduados) e o funtor de feixe associado $P\mapsto\widetilde{P}$ (na categoria dos $S_{(f)}$-módulos) têm, todos os três, as propriedades de exatidão e de comutação com os limites indutivos e as somas diretas \sref[I]{I.1.3.5} e \sref[I]{I.1.3.9}; daí a proposição.
\end{proof}

Denotaremos por $\widetilde{u}$ o homomorfismo $\widetilde{M}\to\widetilde{N}$ correspondente a um homomorfismo de grau $0$, $u:M\to N$.
Deduz-se imediatamente de \sref{II.2.5.4} que os resultados de \sref[I]{I.1.3.9} e \sref[I]{I.1.3.10} continuam válidos para os $S$-módulos graduados e os homomorfismos de grau $0$ (com o sentido dado aqui a $\widetilde{M}$), pois as demonstrações são puramente formais.

\begin{proposition}[2.5.5]
\label{II.2.5.5}
Para todo $\mathfrak{p}\in X=\Proj(S)$, tem-se $\widetilde{M}_\mathfrak{p}=M_{(\mathfrak{p})}$.
\end{proposition}

\begin{proof}
Por definição, $\widetilde{M}_\mathfrak{p}=\varinjlim\Gamma(D_+(f),\widetilde{M})$, onde $f$ percorre o conjunto dos elementos homogêneos de $S_+$ tais que $f\notin\mathfrak{p}$.
Como $\Gamma(D_+(f),\widetilde{M})=M_{(f)}$, a proposição decorre da definição de $M_{(\mathfrak{p})}$ \sref{II.2.2.7}.
\end{proof}
\oldpage[II]{32}
Em particular, o \emph{anel local $(\sh{O}_X)_\mathfrak{p}$} não é outro senão o anel $S_{(\mathfrak{p})}$, conjunto das frações $x/f$, em que $f$ é homogêneo em $S_+$ e não pertence a $\mathfrak{p}$, e $x$ é homogêneo do mesmo grau que $f$.

Mais particularmente, se $S$ é \emph{essencialmente íntegro}, de modo que $\Proj(S)=X$ é \emph{integral} \sref{II.2.4.4}, e se $\xi=(0)$ é o ponto genérico de $X$, o \emph{corpo das funções racionais $R(X)=\sh{O}_\xi$} é o corpo formado pelos elementos $fg^{-1}$, em que $f$ e $g$ são homogêneos do mesmo grau em $S_+$ e $g\neq0$.

\begin{proposition}[2.5.6]
\label{II.2.5.6}
Se, para todo $z\in M$ e todo $f$ homogêneo em $S_+$, existe uma potência de $f$ que anula $z$, então $\widetilde{M}=0$.
Essa condição suficiente é também necessária quando a $S_0$-álgebra $S$ é gerada pelo conjunto $S_1$ dos elementos homogêneos de grau $1$.
\end{proposition}

\begin{proof}
A condição $\widetilde{M}=0$ equivale a $M_{(f)}=0$ para todo $f$ homogêneo em $S_+$.
Por outro lado, se $f\in S_d$, dizer que $M_{(f)}=0$ significa que, para todo $z\in M$ homogêneo e de grau múltiplo de $d$, há uma potência $f^n$ tal que $f^nz=0$.
Se $M_{(f)}=0$ para todo $f\in S_1$, a condição do enunciado é, portanto, satisfeita para todo $f\in S_1$; ela é \emph{a fortiori} satisfeita para todo $f$ homogêneo em $S_+$ quando $S_1$ gera $S$, pois todo elemento homogêneo de $S_+$ é então combinação linear de produtos de elementos de $S_1$.
\end{proof}

\begin{proposition}[2.5.7]
\label{II.2.5.7}
Sejam $d>0$ um inteiro e $f\in S_d$.
Então, para todo $n\in\bb{Z}$, o $(\sh{O}_X|D_+(f))$-módulo $(S(nd))\supertilde|D_+(f)$ é canonicamente isomorfo a $\sh{O}_X|D_+(f)$.
\end{proposition}

\begin{proof}
Com efeito, a multiplicação pelo elemento invertível $(f/1)^n$ de $S_f$ é uma bijeção de $S_{(f)}=(S_f)_0$ sobre
\[
  (S_f)_{nd}=(S_f(nd))_0=(S(nd)_f)_0=S(nd)_{(f)}.
\]
Dito de outro modo, os $S_{(f)}$-módulos $S_{(f)}$ e $S(nd)_{(f)}$ são canonicamente isomorfos; daí a proposição.
\end{proof}

\begin{corollary}[2.5.8]
\label{II.2.5.8}
No subconjunto aberto $U=\bigcup_{f\in S_d}D_+(f)$, a restrição do $\sh{O}_X$-módulo $(S(nd))\supertilde$ é um $(\sh{O}_X|U)$-módulo invertível \sref[0]{0.5.4.1}.
\end{corollary}

\begin{corollary}[2.5.9]
\label{II.2.5.9}
Se o ideal $S_+$ de $S$ é gerado pelo conjunto $S_1$ dos elementos homogêneos de grau $1$, os $\sh{O}_X$-módulos $(S(n))\supertilde$ são invertíveis para todo $n\in\bb{Z}$.
\end{corollary}

\begin{proof}
Basta observar que $X=\bigcup_{f\in S_1}D_+(f)$ em virtude da hipótese \sref{II.2.3.14}, e aplicar \sref{II.2.5.8} com $U=X$.
\end{proof}

\begin{env}[2.5.10]
\label{II.2.5.10}
Poremos, em todo o restante deste parágrafo,
\[
\label{II.2.5.10.1}
  \sh{O}_X(n)=(S(n))\supertilde
\tag{2.5.10.1}
\]
para todo $n\in\bb{Z}$ e, para todo subconjunto aberto $U$ de $X$ e todo $(\sh{O}_X|U)$-módulo $\sh{F}$,
\[
\label{II.2.5.10.2}
  \sh{F}(n)=\sh{F}\otimes_{\sh{O}_X|U}(\sh{O}_X(n)|U)
\tag{2.5.10.2}
\]
para todo $n\in\bb{Z}$.
Se o ideal $S_+$ é gerado por $S_1$, o funtor $\sh{F}\mapsto\sh{F}(n)$ é \emph{exato} em $\sh{F}$ para todo $n\in\bb{Z}$, pois $\sh{O}_X(n)$ é então um $\sh{O}_X$-módulo \emph{invertível}.
\end{env}

\begin{env}[2.5.11]
\label{II.2.5.11}
Sejam $M$ e $N$ dois $S$-módulos graduados.
Para todo $f\in S_d$ ($d>0$), define-se um homomorfismo canônico funtorial de $S_{(f)}$-módulos
\[
\label{II.2.5.11.1}
  \lambda_f:M_{(f)}\otimes_{S_{(f)}}N_{(f)}
  \to(M\otimes_S N)_{(f)}
\tag{2.5.11.1}
\]
\oldpage[II]{33}
compondo o homomorfismo $M_{(f)}\otimes_{S_{(f)}}N_{(f)}\to M_f\otimes_{S_f}N_f$ (proveniente das injeções $M_{(f)}\to M_f$, $N_{(f)}\to N_f$ e $S_{(f)}\to S_f$) com o isomorfismo canônico $M_f\otimes_{S_f}N_f\simto(M\otimes_S N)_f$ \sref[0]{0.1.3.4}, e observando que, pela definição do produto tensorial de dois módulos graduados, este último isomorfismo conserva os graus.
Para $x\in M_{md}$ e $y\in N_{nd}$ ($m,n\geq0$), tem-se, portanto,
\[
  \lambda_f((x/f^m)\otimes(y/f^n))=(x\otimes y)/f^{m+n}.
\]

Decorre imediatamente dessa definição que, se $g\in S_e$ ($e>0$), o diagrama
\[
  \xymatrix{
    M_{(f)}\otimes_{S_{(f)}}N_{(f)} \ar[r]^{\lambda_f} \ar[d]
    & (M\otimes_S N)_{(f)} \ar[d]
  \\M_{(fg)}\otimes_{S_{(fg)}}N_{(fg)} \ar[r]_{\lambda_{fg}}
    & (M\otimes_S N)_{(fg)}
  }
\]
(em que a seta vertical da direita é o homomorfismo canônico e a da esquerda provém dos homomorfismos canônicos) é comutativo.
Os homomorfismos $\lambda_f$ induzem, portanto, um homomorfismo canônico funtorial de $\sh{O}_X$-módulos
\[
\label{II.2.5.11.2}
  \lambda:\widetilde{M}\otimes_{\sh{O}_X}\widetilde{N}
  \to(M\otimes_S N)\supertilde.
\tag{2.5.11.2}
\]

Consideremos, em particular, dois ideais graduados $\mathfrak{J},\mathfrak{K}$ de $S$.
Como $\widetilde{\mathfrak{J}}$ e $\widetilde{\mathfrak{K}}$ são feixes de ideais de $\sh{O}_X$, tem-se um homomorfismo canônico $\widetilde{\mathfrak{J}}\otimes_{\sh{O}_X}\widetilde{\mathfrak{K}}\to\sh{O}_X$; o diagrama
\[
\label{II.2.5.11.3}
  \xymatrix{
    \widetilde{\mathfrak{J}}\otimes_{\sh{O}_X}\widetilde{\mathfrak{K}} \ar[rr]^\lambda \ar[dr]
    &&(\mathfrak{J}\otimes_S\mathfrak{K})\supertilde \ar[dl]
  \\&\sh{O}_X&
  }
\tag{2.5.11.3}
\]
é então comutativo.
Com efeito, reduzimo-nos a verificá-lo em cada subconjunto aberto $D_+(f)$ (para $f$ homogêneo em $S_+$), e isso decorre imediatamente da definição \sref{II.2.5.11.1} de $\lambda_f$ e de \sref[I]{I.1.3.13}.

Observemos, por fim, que, se $M,N,P$ são três $S$-módulos graduados, o diagrama
\[
\label{II.2.5.11.4}
  \xymatrix{
    \widetilde{M}\otimes_{\sh{O}_X}\widetilde{N}\otimes_{\sh{O}_X}\widetilde{P} \ar[r]^{\lambda\otimes1} \ar[d]_{1\otimes\lambda}
    &(M\otimes_S N)\supertilde\otimes_{\sh{O}_X}\widetilde{P} \ar[d]^\lambda
  \\widetilde{M}\otimes_{\sh{O}_X}(N\otimes_S P)\supertilde \ar[r]_\lambda
    &(M\otimes_S N\otimes_S P)\supertilde
  }
\tag{2.5.11.4}
\]
\oldpage[II]{34}
é comutativo.
Basta novamente verificá-lo em cada $D_+(f)$, e isso decorre imediatamente das definições e de \sref[I]{I.1.3.13}.
\end{env}

\begin{env}[2.5.12]
\label{II.2.5.12}
Sob as hipóteses de \sref{II.2.5.11}, define-se um homomorfismo canônico funtorial de $S_{(f)}$-módulos
\[
\label{II.2.5.12.1}
  \mu_f:(\Hom_S(M,N))_{(f)}
  \to\Hom_{S_{(f)}}(M_{(f)},N_{(f)})
\tag{2.5.12.1}
\]
fazendo corresponder a $u/f^n$, em que $u$ é um homomorfismo de grau $nd$, o homomorfismo $M_{(f)}\to N_{(f)}$ que transforma $x/f^m$ ($x\in M_{md}$) em $u(x)/f^{m+n}$.
Para $g\in S_e$ ($e>0$), tem-se ainda um diagrama comutativo:
\[
  \xymatrix{
    (\Hom_S(M,N))_{(f)} \ar[r]^{\mu_f} \ar[d]
    & \Hom_{S_{(f)}}(M_{(f)},N_{(f)}) \ar[d]
  \\(\Hom_S(M,N))_{(fg)} \ar[r]_{\mu_{fg}}
    & \Hom_{S_{(fg)}}(M_{(fg)},N_{(fg)})
  }
\]
(em que a seta da esquerda é o homomorfismo canônico e a da direita provém dos homomorfismos canônicos).
Conclui-se ainda (levando em conta \sref[I]{I.1.3.8}) que os $\mu_f$ definem um homomorfismo canônico funtorial de $\sh{O}_X$-módulos
\[
\label{II.2.5.12.2}
  \mu:(\Hom_S(M,N))\supertilde
  \to\shHom_{\sh{O}_X}(\widetilde{M},\widetilde{N}).
\tag{2.5.12.2}
\]
\end{env}

\begin{proposition}[2.5.13]
\label{II.2.5.13}
Suponhamos que o ideal $S_+$ seja gerado por $S_1$.
Então, o homomorfismo $\lambda$ \sref{II.2.5.11.2} é um isomorfismo; o mesmo vale para o homomorfismo $\mu$ \sref{II.2.5.12.2} quando o $S$-módulo graduado $M$ admite uma apresentação finita \sref{II.2.1.1}.
\end{proposition}

\begin{proof}
Como $X$ é reunião dos $D_+(f)$ para $f\in S_1$ \sref{II.2.3.14}, reduzimo-nos a provar que $\lambda_f$ e $\mu_f$ são isomorfismos, sob as hipóteses consideradas, quando $f$ é homogêneo e \emph{de grau $1$}.
Ora, define-se então uma aplicação $\bb{Z}$-bilinear $M_m\times N_n\to M_{(f)}\otimes_{S_{(f)}}N_{(f)}$ fazendo corresponder a $(x,y)$ o elemento $(x/f^m)\otimes(y/f^n)$ (quando $m<0$, escrevemos $x/f^m$ para $f^{-m}x/1$).
Essas aplicações definem uma aplicação $\bb{Z}$-linear $M\otimes_{\bb{Z}}N\to M_{(f)}\otimes_{S_{(f)}}N_{(f)}$ e, se $s\in S_q$, essa aplicação transforma $(sx)\otimes y$ em
\[
  (s/f^q)((x/f^m)\otimes(y/f^n))
\]
para $x\in M_m$ e $y\in N_n$.
Obtém-se, portanto, um di-homomorfismo de módulos
\[
  \gamma_f:M\otimes_S N\to M_{(f)}\otimes_{S_{(f)}}N_{(f)}
\]
relativo ao homomorfismo canônico $S\to S_{(f)}$ que transforma $s\in S_q$ em $s/f^q$.
Suponhamos, além disso, que, para um elemento $\sum_i(x_i\otimes y_i)$ de $M\otimes_S N$, em que $x_i,y_i$ são homogêneos de graus respectivos $m_i,n_i$, se tenha
\[
  f^r\sum_i(x_i\otimes y_i)=0,
\]
dito de outro modo, $\sum_i(f^rx_i\otimes y_i)=0$.
Deduz-se então de \sref[0]{0.1.3.4} que
\[
  \sum_i(f^rx_i/f^{m_i+r})\otimes(y_i/f^{n_i})=0,
\]
isto é, $\gamma_f(\sum_i(x_i\otimes y_i))=0$.
Por conseguinte, $\gamma_f$ fatora-se em
\[
  M\otimes_S N\to(M\otimes_S N)_f
  \xrightarrow{\gamma'_f}M_{(f)}\otimes_{S_{(f)}}N_{(f)}.
\]
Se $\lambda'_f$ é a restrição de $\gamma'_f$
\oldpage[II]{35}
a $(M\otimes_S N)_{(f)}$, verifica-se imediatamente que $\lambda_f$ e $\lambda'_f$ são $S_{(f)}$-homomorfismos recíprocos, de onde resulta a primeira parte da proposição.

Para demonstrar a segunda parte, suponhamos que $M$ seja o conúcleo de um homomorfismo $P\to Q$ de $S$-módulos graduados, sendo $P$ e $Q$ somas diretas de um número finito de módulos da forma $S(n)$.
Usando a exatidão à esquerda do funtor $L\mapsto\Hom_S(L,N)$ e a exatidão do funtor $M\mapsto M_{(f)}$, reduzimo-nos imediatamente a provar que $\mu_f$ é um isomorfismo quando $M=S(n)$.
Ora, para todo $z$ homogêneo em $N$, seja $u_z$ o homomorfismo de $S(n)$ em $N$ tal que $u_z(1)=z$; vê-se imediatamente que $\eta:z\mapsto u_z$ é um isomorfismo de grau $0$ de $N(-n)$ sobre $\Hom_S(S(n),N)$.
Corresponde-lhe um isomorfismo
\[
  \eta_f:(N(-n))_{(f)}\to(\Hom_S(S(n),N))_{(f)}.
\]

Por outro lado, seja $\eta'_f$ o isomorfismo $N_{(f)}\to\Hom_{S_{(f)}}(S(n)_{(f)},N_{(f)})$ que, a todo $z'\in N_{(f)}$, faz corresponder o homomorfismo $v_{z'}$ tal que
\[
  v_{z'}(s/f^k)=sz'/f^{n+k}
\]
para $s\in S_{n+k}=(S(n))_k$.
Constata-se facilmente que a aplicação composta
\[
  (N(-n))_{(f)}
  \xrightarrow{\eta_f}(\Hom_S(S(n),N))_{(f)}
  \xrightarrow{\mu_f}\Hom_{S_{(f)}}(S(n)_{(f)},N_{(f)})
  \xrightarrow{{\eta'_f}^{-1}}N_{(f)}
\]
é o isomorfismo $z/f^h\mapsto z/f^{h-n}$ de $(N(-n))_{(f)}$ sobre $N_{(f)}$; portanto, $\mu_f$ é um isomorfismo.
\end{proof}

Se o ideal $S_+$ é gerado por $S_1$, deduz-se de \sref{II.2.5.13} que, para todo ideal graduado $\mathfrak{J}$ de $S$ e todo $S$-módulo graduado $M$, tem-se
\[
\label{II.2.5.13.1}
  \widetilde{\mathfrak{J}}\cdot\widetilde{M}
  =(\mathfrak{J}\cdot M)\supertilde
\tag{2.5.13.1}
\]
a menos de isomorfismo canônico; com efeito, isso resulta da comutatividade do diagrama
\[
  \xymatrix{
    \widetilde{\mathfrak{J}}\otimes_{\sh{O}_X}\widetilde{M} \ar[rr]^\lambda \ar[dr]
    &&(\mathfrak{J}\otimes_S M)\supertilde \ar[dl]
  \\&\widetilde{M}&
  }
\]
que se verifica como o de \sref{II.2.5.11.3}.

\begin{corollary}[2.5.14]
\label{II.2.5.14}
Suponhamos que $S$ seja gerado por $S_1$.
Quaisquer que sejam $m,n\in\bb{Z}$, tem-se então:
\[
\label{II.2.5.14.1}
  \sh{O}_X(m)\otimes_{\sh{O}_X}\sh{O}_X(n)=\sh{O}_X(m+n)
\tag{2.5.14.1}
\]
\[
\label{II.2.5.14.2}
  \sh{O}_X(n)=(\sh{O}_X(1))^{\otimes n}
\tag{2.5.14.2}
\]
a menos de isomorfismos canônicos.
\end{corollary}

\begin{proof}
A primeira fórmula resulta de \sref{II.2.5.13} e da existência do isomorfismo canônico de grau $0$
\[
  S(m)\otimes_S S(n)\simto S(m+n),
\]
que faz corresponder ao elemento $1\otimes1$ (onde o primeiro fator $1$ está em $(S(m))_{-m}$ e o segundo em $(S(n))_{-n}$) o elemento $1\in(S(m+n))_{-(m+n)}$.
Basta então demonstrar a segunda fórmula para $n=-1$ e, em virtude de \sref{II.2.5.13}, isso se reduz a ver que $\Hom_S(S(1),S)$ é canonicamente isomorfo a $S(-1)$, o que se verifica imediatamente remontando às definições \sref{II.2.1.2} e lembrando que $S(1)$ é um $S$-módulo cíclico.
\end{proof}
\oldpage[II]{36}
\begin{corollary}[2.5.15]
\label{II.2.5.15}
Suponhamos que $S$ seja gerado por $S_1$.
Então, para todo $S$-módulo graduado $M$ e todo $n\in\bb{Z}$, tem-se
\[
\label{II.2.5.15.1}
  (M(n))\supertilde=\widetilde{M}(n)
\tag{2.5.15.1}
\]
a menos de isomorfismo canônico.
\end{corollary}

\begin{proof}
Isso resulta das definições \sref{II.2.5.10.2} e \sref{II.2.5.10.1}, da Proposição~\sref{II.2.5.13} e da existência de um isomorfismo canônico $M(n)\simto M\otimes_S S(n)$ de grau $0$ que, a todo $z\in(M(n))_h=M_{n+h}$, faz corresponder
\[
  z\otimes1\in M_{n+h}\otimes(S(n))_{-n}
  \subset(M\otimes_S S(n))_h.
\]
\end{proof}

\begin{env}[2.5.16]
\label{II.2.5.16}
Designemos por $S'$ o anel graduado tal que $S'_0=\bb{Z}$ e $S'_n=S_n$ para $n>0$.
Então, se $f\in S_d$ ($d>0$), tem-se $(S(n))_{(f)}=(S'(n))_{(f)}$ para todo $n\in\bb{Z}$, pois um elemento de $(S'(n))_{(f)}$ é da forma $x/f^k$, com $x\in S'_{n+kd}$ ($k>0$), e sempre se pode tomar $k$ tal que $n+kd\neq0$.
Como $X=\Proj(S)$ e $X'=\Proj(S')$ se identificam canonicamente \sref{II.2.4.7}[(ii)], vê-se que, para todo $n\in\bb{Z}$, $\sh{O}_X(n)$ e $\sh{O}_{X'}(n)$ são imagens um do outro por essa identificação.

Notemos, por outro lado, que, para todo $d>0$ e todo $n\in\bb{Z}$, tem-se
\[
  (S^{(d)}(n))_h=S_{(n+h)d}=(S(nd))_{hd}
\]
para $f\in S_d$; tem-se, portanto, $(S^{(d)}(n))_{(f)}=(S(nd))_{(f)}$.
Sabe-se que os esquemas $X=\Proj(S)$ e $X^{(d)}=\Proj(S^{(d)})$ se identificam canonicamente \sref{II.2.4.7}[(i)].
O que precede mostra que, se a $S_0$-álgebra $S^{(d)}$ é gerada por $S_d$, então $\sh{O}_X(nd)$ e $\sh{O}_{X^{(d)}}(n)$ são imagens um do outro por essa identificação, para todo $n\in\bb{Z}$.
\end{env}

\begin{proposition}[2.5.17]
\label{II.2.5.17}
Seja $d>0$ um inteiro, e seja $U=\bigcup_{f\in S_d}D_+(f)$.
Então, para todo inteiro $n$, a restrição a $U$ do homomorfismo canônico
\[
  \sh{O}_X(nd)\otimes_{\sh{O}_X}\sh{O}_X(-nd)\to\sh{O}_X
\]
é um isomorfismo.
\end{proposition}

\begin{proof}
Em virtude de \sref{II.2.5.16}, podemos limitar-nos ao caso em que $d=1$, e a conclusão resulta então da demonstração de \sref{II.2.5.13}.
\end{proof}

\subsection{O $S$-módulo graduado associado a um feixe sobre $\operatorname{Proj}(S)$}
\label{subsection:II.2.6}

\emph{Em toda esta seção, supõe-se que o ideal $S_+$ de $S$ é gerado pelo conjunto $S_1$ dos elementos homogêneos de grau $1$.}

\begin{env}[2.6.1]
\label{II.2.6.1}
O $\sh{O}_X$-módulo $\sh{O}_X(1)$ é então \emph{invertível} \sref{II.2.5.9}; põe-se, para todo $\sh{O}_X$-módulo $\sh{F}$ \sref[0]{0.5.4.6},
\[
\label{II.2.6.1.1}
  \Gamma_\bullet(\sh{F})
  =\Gamma_\bullet(\sh{O}_X(1),\sh{F})
  =\bigoplus_{n\in\bb{Z}}\Gamma(X,\sh{F}(n))
\tag{2.6.1.1}
\]
tendo em conta \sref{II.2.5.14.2}.
Recordemos \sref[0]{0.5.4.6} que $\Gamma_\bullet(\sh{O}_X)$ é munido de uma estrutura de \emph{anel graduado}, e $\Gamma_\bullet(\sh{F})$ de uma estrutura de \emph{$\Gamma_\bullet(\sh{O}_X)$-módulo graduado}.

Como $\sh{O}_X(n)$ é localmente livre, $\Gamma_\bullet(\sh{F})$ é um funtor covariante aditivo \emph{exato à esquerda} em $\sh{F}$; em particular, se $\sh{J}$ é um feixe de ideais de $\sh{O}_X$, $\Gamma_\bullet(\sh{J})$ identifica-se canonicamente com um \emph{ideal graduado} de $\Gamma_\bullet(\sh{O}_X)$.
\end{env}

\begin{env}[2.6.2]
\label{II.2.6.2}
Seja $M$ um $S$-módulo graduado.
Para todo $f\in S_d$ ($d>0$), $x\mapsto x/1$ é um homomorfismo de grupos abelianos $M_0\to M_{(f)}$ e, como $M_{(f)}$ se identifica canonicamente
\oldpage[II]{37}
com $\Gamma(D_+(f),\widetilde{M})$, obtém-se assim um homomorfismo de grupos abelianos
\[
  \alpha_0^f:M_0\to\Gamma(D_+(f),\widetilde{M}).
\]
É claro que, para todo $g\in S_e$ ($e>0$), o diagrama
\[
  \xymatrix{
    & \Gamma(D_+(f),\widetilde{M}) \ar[dd]
  \\M_0 \ar[ur]^{\alpha_0^f} \ar[dr]_{\alpha_0^{fg}}
  \\& \Gamma(D_+(fg),\widetilde{M})
  }
\]
é comutativo; isso significa que, para todo $x\in M_0$, as seções $\alpha_0^f(x)$ e $\alpha_0^g(x)$ de $\widetilde{M}$ coincidem em $D_+(f)\cap D_+(g)$ e, por conseguinte, existe uma única seção $\alpha_0(x)\in\Gamma(X,\widetilde{M})$ cuja restrição a cada $D_+(f)$ é $\alpha_0^f(x)$.
Definiu-se assim, sem utilizar a hipótese de que $S$ é gerado por $S_1$, um homomorfismo de grupos abelianos
\[
\label{II.2.6.2.1}
  \alpha_0:M_0\to\Gamma(X,\widetilde{M}).
\tag{2.6.2.1}
\]

Aplicando esse resultado ao $S$-módulo graduado $M(n)$, para todo $n\in\bb{Z}$, obtém-se, para cada $n\in\bb{Z}$, um homomorfismo de grupos abelianos
\[
\label{II.2.6.2.2}
  \alpha_n:M_n=(M(n))_0\to\Gamma(X,\widetilde{M}(n))
\tag{2.6.2.2}
\]
tendo em conta \sref{II.2.5.15}; daí se obtém um homomorfismo funtorial (de grau $0$) de grupos abelianos graduados
\[
\label{II.2.6.2.3}
  \alpha:M\to\Gamma_\bullet(\widetilde{M})
\tag{2.6.2.3}
\]
(também denotado por $\alpha_M$) que, em cada $M_n$, coincide com $\alpha_n$.

Se tomarmos, em particular, $M=S$, verifica-se imediatamente, tendo em conta a definição da multiplicação em $\Gamma_\bullet(\sh{O}_X)$ \sref[0]{0.5.4.6}, que $\alpha:S\to\Gamma_\bullet(\sh{O}_X)$ é um homomorfismo de anéis graduados e que, para todo $S$-módulo graduado $M$, \sref{II.2.6.2.3} é um di-homomorfismo de módulos graduados.
\end{env}

\begin{proposition}[2.6.3]
\label{II.2.6.3}
Para todo $f\in S_d$ ($d>0$), $D_+(f)$ é idêntico ao conjunto dos $\mathfrak{p}\in X$ nos quais a seção $\alpha_d(f)$ de $\sh{O}_X(d)$ não se anula \sref[0]{0.5.5.2}.
\end{proposition}

\begin{proof}
Como $X=\bigcup_{g\in S_1}D_+(g)$ por hipótese, basta verificar que, para todo $g\in S_1$, o conjunto dos $\mathfrak{p}\in D_+(g)$ nos quais $\alpha_d(f)$ não se anula é idêntico a $D_+(fg)$.
Ora, a restrição de $\alpha_d(f)$ a $D_+(g)$ é, por definição, a seção correspondente ao elemento $f/1$ de $(S(d))_{(g)}$.
Pelo isomorfismo canônico $(S(d))_{(g)}\simto S_{(g)}$ \sref{II.2.5.7}, essa seção de $\sh{O}_X(d)$ sobre $D_+(g)$ identifica-se com a seção de $\sh{O}_X$ sobre $D_+(g)$ que corresponde ao elemento $f/g^d$ de $S_{(g)}$.
Dizer que essa seção se anula em $\mathfrak{p}\in D_+(g)$ significa que $f/g^d\in\mathfrak{q}$, onde $\mathfrak{q}$ é o ideal primo de $S_{(g)}$ correspondente a $\mathfrak{p}$ \sref{II.2.3.6}; por definição, isso quer dizer que $f\in\mathfrak{p}$, de onde resulta a proposição.
\end{proof}

\begin{env}[2.6.4]
\label{II.2.6.4}
Seja agora $\sh{F}$ um $\sh{O}_X$-módulo, e ponhamos $M=\Gamma_\bullet(\sh{F})$.
Em virtude da existência do homomorfismo de anéis graduados $\alpha:S\to\Gamma_\bullet(\sh{O}_X)$, $M$ pode ser considerado um $S$-módulo graduado.
Para todo $f\in S_d$ ($d>0$), resulta de \sref{II.2.6.3} que a restrição a $D_+(f)$ da seção $\alpha_d(f)$ de $\sh{O}_X(d)$ é invertível; o mesmo vale, portanto, para a restrição a $D_+(f)$ da seção $\alpha_{nd}(f^n)$ de $\sh{O}_X(nd)$, para todo $n>0$.
Seja então $z\in M_{nd}=\Gamma(X,\sh{F}(nd))$ ($n>0$); se existe um inteiro $k\geq0$ tal que a restrição a $D_+(f)$ de $f^kz$, isto é, a
\oldpage[II]{38}
seção $(z|D_+(f))(\alpha_{kd}(f^k)|D_+(f))$ de $\sh{F}((n+k)d)$ é nula, então, em razão da observação precedente, tem-se também $z|D_+(f)=0$.
Isso mostra que se define um $S_{(f)}$-homomorfismo $\beta_f:M_{(f)}\to\Gamma(D_+(f),\sh{F})$ fazendo corresponder ao elemento $z/f^n$ a seção
\[
  (z|D_+(f))(\alpha_{nd}(f^n)|D_+(f))^{-1}
\]
de $\sh{F}$ sobre $D_+(f)$.
Verifica-se ainda imediatamente que o diagrama
\[
\label{II.2.6.4.1}
  \xymatrix{
    M_{(f)} \ar[r]^{\beta_f} \ar[d]
    & \Gamma(D_+(f),\sh{F}) \ar[d]
  \\M_{(fg)} \ar[r]_{\beta_{fg}}
    & \Gamma(D_+(fg),\sh{F})
  }
\tag{2.6.4.1}
\]
é comutativo para $g\in S_e$ ($e>0$).
Se lembrarmos que $M_{(f)}$ se identifica canonicamente com $\Gamma(D_+(f),\widetilde{M})$ e que os $D_+(f)$ formam uma base da topologia de $X$ \sref{II.2.3.4}, vê-se que os $\beta_f$ provêm de um único homomorfismo canônico de $\sh{O}_X$-módulos
\[
\label{II.2.6.4.2}
  \beta:(\Gamma_\bullet(\sh{F}))\supertilde\to\sh{F}
\tag{2.6.4.2}
\]
(também denotado por $\beta_{\sh{F}}$), que é evidentemente funtorial.
\end{env}

\begin{proposition}[2.6.5]
\label{II.2.6.5}
Sejam $M$ um $S$-módulo graduado e $\sh{F}$ um $\sh{O}_X$-módulo; então os homomorfismos compostos
\[
\label{II.2.6.5.1}
  \widetilde{M}
  \xrightarrow{\widetilde{\alpha}}(\Gamma_\bullet(\widetilde{M}))\supertilde
  \xrightarrow{\beta}\widetilde{M}
\tag{2.6.5.1}
\]
\[
\label{II.2.6.5.2}
  \Gamma_\bullet(\sh{F})
  \xrightarrow{\alpha}\Gamma_\bullet((\Gamma_\bullet(\sh{F}))\supertilde)
  \xrightarrow{\Gamma_\bullet(\beta)}\Gamma_\bullet(\sh{F})
\tag{2.6.5.2}
\]
são os isomorfismos identidade.
\end{proposition}

\begin{proof}
A verificação de \sref{II.2.6.5.1} é local: em um aberto $D_+(f)$, ela resulta imediatamente das definições e do fato de que $\beta$, aplicado a feixes quase-coerentes, é determinado por sua ação nas seções sobre $D_+(f)$ \sref[I]{I.1.3.8}.
A verificação de \sref{II.2.6.5.2} faz-se separadamente em cada grau: se pusermos $M=\Gamma_\bullet(\sh{F})$, teremos $M_n=\Gamma(X,\sh{F}(n))$ e
\[
  (\Gamma_\bullet(\widetilde{M}))_n
  =\Gamma(X,\widetilde{M}(n))
  =\Gamma(X,(M(n))\supertilde).
\]
Ora, se $f\in S_1$ e $z\in M_n$, $\alpha_n^f(z)$ é o elemento $z/1$ de $(M(n))_{(f)}$, igual a $(f/1)^n(z/f^n)$; corresponde-lhe, por $\beta_f$, a seção
\[
  \Big((\alpha_1(f))^n|D_+(f)\Big)
  \Big((z|D_+(f))((\alpha_1(f))^n|D_+(f))^{-1}\Big)
\]
sobre $D_+(f)$, isto é, a restrição de $z$ a $D_+(f)$, o que verifica \sref{II.2.6.5.2}.
\end{proof}

\subsection{Condições de finitude}
\label{subsection:II.2.7}

\begin{proposition}[2.7.1]
\label{II.2.7.1}
\begin{enumerate}
  \item[{\rm(i)}] Se $S$ é um anel graduado noetheriano, então $X=\Proj(S)$ é um esquema noetheriano.
  \item[{\rm(ii)}] Se $S$ é uma $A$-álgebra graduada de tipo finito, então $X=\Proj(S)$ é um esquema de tipo finito sobre $Y=\Spec(A)$.
\end{enumerate}
\end{proposition}
\oldpage[II]{39}
\begin{proof}
\medskip\noindent
\begin{enumerate}
  \item[{\rm(i)}] Se $S$ é noetheriano, o ideal $S_+$ admite um sistema finito de geradores homogêneos $f_i$ ($1\leq i\leq p$); portanto, por \sref{II.2.3.14}, o espaço subjacente $X$ é a união dos $D_+(f_i)=\Spec(S_{(f_i)})$, e tudo se reduz a ver que cada um dos $S_{(f_i)}$ é noetheriano, o que resulta de \sref{II.2.2.6}.
  \item[{\rm(ii)}] A hipótese implica que $S_0$ é uma $A$-álgebra de tipo finito e que $S$ é uma $S_0$-álgebra de tipo finito; portanto, $S_+$ é um ideal de tipo finito \sref{II.2.1.4}.
    Reduzimo-nos então, como em (i), a provar que, para todo $f\in S_d$, $S_{(f)}$ é uma $A$-álgebra de tipo finito.
    Em virtude de \sref{II.2.2.5}, basta mostrar que $S^{(d)}$ é uma $A$-álgebra de tipo finito, o que resulta de \sref{II.2.1.6}.
\end{enumerate}
\end{proof}

\begin{env}[2.7.2]
\label{II.2.7.2}
Consideraremos, no que segue, as seguintes condições de finitude para um $S$-módulo graduado $M$:
\begin{enumerate}
  \item[{\rm(TF)}] Existe um inteiro $n$ tal que o submódulo $\bigoplus_{k\geq n}M_k$ seja um $S$-módulo de tipo finito.
  \item[{\rm(TN)}] Existe um inteiro $n$ tal que $M_k=0$ para $k\geq n$.
\end{enumerate}

Se $M$ satisfaz (TN), tem-se $M_{(f)}=0$ para todo $f$ homogêneo em $S_+$; portanto, $\widetilde{M}=0$.

Sejam $M$ e $N$ dois $S$-módulos graduados; diz-se que um homomorfismo $u:M\to N$ de grau $0$ é (TN)-\emph{injetivo} (resp. (TN)-\emph{sobrejetivo}, (TN)-\emph{bijetivo}) se existe um inteiro $n$ tal que $u_k:M_k\to N_k$ seja injetivo (resp. sobrejetivo, bijetivo) para $k\geq n$.
Dizer que $u$ é (TN)-injetivo (resp. (TN)-sobrejetivo) equivale, portanto, a dizer que $\Ker u$ (resp. $\Coker u$) satisfaz (TN).
Em virtude de \sref{II.2.5.4}, se $u$ é (TN)-injetivo (resp. (TN)-sobrejetivo, (TN)-bijetivo), então $\widetilde{u}$ é injetivo (resp. sobrejetivo, bijetivo); quando $u$ é (TN)-bijetivo, diz-se ainda que $u$ é um (TN)-\emph{isomorfismo}.
\end{env}

\begin{proposition}[2.7.3]
\label{II.2.7.3}
Sejam $S$ um anel graduado tal que o ideal $S_+$ seja de tipo finito, e $M$ um $S$-módulo graduado.
\begin{enumerate}
  \item[{\rm(i)}] Se $M$ satisfaz a condição (TF), o $\sh{O}_X$-módulo $\widetilde{M}$ é de tipo finito.
  \item[{\rm(ii)}] Suponhamos que $M$ satisfaça (TF); para que $\widetilde{M}=0$, é necessário e suficiente que $M$ satisfaça (TN).
\end{enumerate}
\end{proposition}

\begin{proof}
Acabamos de ver que a condição (TN) implica $\widetilde{M}=0$.
Se $M$ satisfaz (TF), o submódulo graduado $M'=\bigoplus_{k\geq n}M_k$, que é, por hipótese, de tipo finito, é tal que $M/M'$ satisfaz (TN); tem-se, por conseguinte, $(M/M')\supertilde=0$, e a exatidão do funtor $M\mapsto\widetilde{M}$ \sref{II.2.5.4} implica $\widetilde{M}=\widetilde{M'}$; para provar que $\widetilde{M}$ é de tipo finito, reduzimo-nos, portanto, ao caso em que $M$ é \emph{de tipo finito}.
Sendo a questão local, basta provar que $M_{(f)}$ é um $S_{(f)}$-módulo de tipo finito \sref[I]{I.1.3.9}; mas $M^{(d)}$ é um $S^{(d)}$-módulo de tipo finito \sref{II.2.1.6}[iii], e nossa afirmação resulta então de \sref{II.2.2.5}.

Suponhamos agora que $M$ satisfaça (TF) e que $\widetilde{M}=0$; então, com as mesmas notações acima, tem-se $\widetilde{M'}=0$, e a condição (TN) para $M'$ equivale à condição (TN) para $M$; portanto, para provar que $\widetilde{M}=0$ implica que $M$ satisfaz (TN), podemos ainda limitar-nos ao caso em que $M$ é gerado por um número finito de elementos homogêneos $x_i$ ($1\leq i\leq p$); seja, por outro lado, $(f_j)_{1\leq j\leq q}$ um sistema de geradores homogêneos do ideal $S_+$.
Por hipótese, tem-se $M_{(f_j)}=0$ para todo $j$; logo, existe um inteiro $n$ tal que $f_j^nx_i=0$, quaisquer que sejam $i$ e $j$.
Seja $n_j$ o grau de $f_j$, e seja $m$ o maior valor de $\sum_j r_jn_j$ sobre o conjunto finito dos sistemas de inteiros não negativos $(r_j)$ tais que $\sum_j r_j\leq nq$; é claro então
\oldpage[II]{40}
que, se $k>m$, então $S_kx_i=0$ para todo $i$;
se $h$ é o maior dos graus dos $x_i$, conclui-se que $M_k=0$ para $k>h+m$, o que termina a demonstração.
\end{proof}

\begin{corollary}[2.7.4]
\label{II.2.7.4}
Seja $S$ um anel graduado tal que o ideal $S_+$ seja de tipo finito;
para que $X=\Proj(S)=\emp$, é necessário e suficiente que exista $n$ tal que $S_k=0$ para $k\geq n$.
\end{corollary}

\begin{proof}
A condição $X=\emp$ é, com efeito, equivalente a $\sh{O}_X=\widetilde{S}=0$, e $S$ é um $S$-módulo cíclico.
\end{proof}

\begin{theorem}[2.7.5]
\label{II.2.7.5}
Suponha-se que o ideal $S_+$ seja gerado por um número finito de elementos homogêneos de grau~$1$;
seja $X=\Proj(S)$.
Então, para todo $\sh{O}_X$-módulo quase-coerente $\sh{F}$, o homomorfismo canônico $\beta:(\Gamma_\bullet(\sh{F}))\supertilde\to\sh{F}$ \sref{II.2.6.4} é um isomorfismo.
\end{theorem}

\begin{proof}
Se $S_+$ é gerado por um número finito de elementos $f_i\in S_1$, então $X$ é a união dos subespaços $\Spec(S_{(f_i)})$ \sref{II.2.3.6}, que são quase-compactos, e portanto $X$ é quase-compacto;
além disso, $X$ é um esquema \sref{II.2.4.2};
por \sref[I]{I.9.3.2}, \sref{II.2.5.14.2} e \sref{II.2.6.3}, tem-se, para todo $f\in S_d$ ($d>0$), um isomorfismo canônico $(\Gamma_\bullet(\sh{F}))_{(\alpha_d(f))}\simto\Gamma(D_+(f),\sh{F})$;
por outro lado, por definição, $(\Gamma_\bullet(\sh{F}))_{(\alpha_d(f))}$ (onde $\Gamma_\bullet(\sh{F})$ é considerado como $\Gamma_\bullet(\sh{O}_X)$-módulo) não é senão $(\Gamma_\bullet(\sh{F}))_{(f)}$ (onde $\Gamma_\bullet(\sh{F})$ é considerado como $S$-módulo);
reportando-se à definição \sref[I]{I.9.3.1} do isomorfismo canônico precedente, vê-se que ele coincide com o homomorfismo $\beta_f$, donde o teorema.
\end{proof}

\begin{remark}[2.7.6]
\label{II.2.7.6}
Se se supõe que o anel graduado $S$ é \emph{noetheriano}, a condição de \sref{II.2.7.5} verifica-se \emph{ipso facto} desde que se suponha que o ideal $S_+$ seja gerado pelo conjunto $S_1$ dos elementos homogêneos de grau~$1$.
\end{remark}

\begin{corollary}[2.7.7]
\label{II.2.7.7}
Sob as hipóteses de \sref{II.2.7.5}, todo $\sh{O}_X$-módulo quase-coerente $\sh{F}$ é isomorfo a um $\sh{O}_X$-módulo da forma $\widetilde{M}$, onde $M$ é um $S$-módulo graduado.
\end{corollary}

\begin{corollary}[2.7.8]
\label{II.2.7.8}
Sob as hipóteses de \sref{II.2.7.5}, todo $\sh{O}_X$-módulo quase-coerente de tipo finito $\sh{F}$ é isomorfo a um $\sh{O}_X$-módulo da forma $\widetilde{N}$, onde $N$ é um $S$-módulo graduado de tipo finito.
\end{corollary}

\begin{proof}
Pode-se supor que $\sh{F}=\widetilde{M}$, onde $M$ é um $S$-módulo graduado \sref{II.2.7.7}.
Seja $(f_\lambda)_{\lambda\in L}$ um sistema de geradores homogêneos de $M$;
para todo subconjunto finito $H$ de $L$, seja $M_H$ o submódulo graduado de $M$ gerado pelos $f_\lambda$ tais que $\lambda\in H$;
é claro que $M$ é o limite indutivo de seus submódulos $M_H$, e portanto $\sh{F}$ é o limite indutivo de seus sub-$\sh{O}_X$-módulos $\widetilde{M_H}$ \sref{II.2.5.4}.
Mas, como $\sh{F}$ é de tipo finito e o espaço subjacente $X$ é quase-compacto, resulta de \sref[0]{0.5.2.3} que $\sh{F}=\widetilde{M_H}$ para algum subconjunto finito $H$ de $L$.
\end{proof}

\begin{corollary}[2.7.9]
\label{II.2.7.9}
Sob as hipóteses de \sref{II.2.7.5}, seja $\sh{F}$ um $\sh{O}_X$-módulo quase-coerente de tipo finito.
Existe então um inteiro $n_0$ tal que, para todo $n\geq n_0$, $\sh{F}(n)$ seja isomorfo a um quociente de um $\sh{O}_X$-módulo da forma $\sh{O}_X^k$ ($k>0$ dependendo de $n$) e, por conseguinte, seja gerado por um número finito de suas seções sobre $X$ \sref[0]{0.5.1.1}.
\end{corollary}

\begin{proof}
Em virtude de \sref{II.2.7.8}, pode-se supor que $\sh{F}=\widetilde{M}$, onde $M$ é quociente de uma soma direta finita de $S$-módulos da forma $S(m_i)$;
em virtude de \sref{II.2.5.4}, pode-se portanto reduzir ao caso em que $M=S(m)$, e assim $\sh{F}(n)=(S(m+n))\supertilde=\sh{O}_X(m+n)$ \sref{II.2.5.15}.
Basta, portanto, demonstrar o
\oldpage[II]{41}
  \begin{lemma}[2.7.9.1]
  \label{II.2.7.9.1}
  Sob as hipóteses de \sref{II.2.7.5}, para todo $n\geq0$, existe um inteiro $k$ (dependendo de $n$) e um homomorfismo sobrejetivo $\sh{O}_X^k\to\sh{O}_X(n)$.
  \end{lemma}

Com efeito, basta \sref{II.2.7.2} mostrar que, para um $k$ conveniente, existe um homomorfismo (TN)-sobrejetivo $u$, \emph{de grau~$0$}, do $S$-módulo graduado produto $S^k$ para $S(n)$.
Ora, $(S(n))_0=S_n$ e, por hipótese, $S_h=S_1^h$ para todo $h>0$, logo $SS_n=S_n+S_{n+1}+\ldots$.
Como $S_n$ é um $S_0$-módulo de tipo finito (\sref{II.2.1.5} e \sref{II.2.1.6}[i]), considere-se um sistema de geradores $(a_i)_{1\leq i\leq k}$ desse módulo;
o homomorfismo $u$ enviará o $i$-ésimo elemento da base canônica de $S^k$ em $a_i$ ($1\leq i\leq k$);
como $\Coker u$ se identifica então com $(S(n))_{-n}+\ldots+(S(n))_{-1}$, $u$ responde efetivamente à questão.
\end{proof}

\begin{corollary}[2.7.10]
\label{II.2.7.10}
Sob as hipóteses de \sref{II.2.7.5}, seja $\sh{F}$ um $\sh{O}_X$-módulo quase-coerente de tipo finito.
Existe um inteiro $n_0$ tal que, para todo $n\geq n_0$, $\sh{F}$ seja isomorfo a um quociente de um $\sh{O}_X$-módulo da forma $(\sh{O}_X(-n))^k$ ($k$ dependendo de $n$).
\end{corollary}

\begin{proposition}[2.7.11]
\label{II.2.7.11}
Suponham-se verificadas as hipóteses de \sref{II.2.7.5}, e seja $M$ um $S$-módulo graduado.
Então:
\begin{enumerate}
  \item[\rm{(i)}] O homomorfismo canônico $\widetilde{\alpha}:\widetilde{M}\to(\Gamma_\bullet(\widetilde{M}))\supertilde$ é um isomorfismo.
  \item[\rm{(ii)}] Seja $\sh{G}$ um sub-$\sh{O}_X$-módulo quase-coerente de $\widetilde{M}$, e seja $N$ o sub-$S$-módulo graduado de $M$, imagem inversa de $\Gamma_\bullet(\sh{G})$ por $\alpha$.
    Então $\widetilde{N}=\sh{G}$ (sendo $\widetilde{N}$ identificado, em virtude de \sref{II.2.5.4}, com um sub-$\sh{O}_X$-módulo de $\widetilde{M}$).
\end{enumerate}
\end{proposition}

\begin{proof}
Como $\beta:(\Gamma_\bullet(\widetilde{M}))\supertilde\to\widetilde{M}$ é um isomorfismo em virtude de \sref{II.2.7.5}, $\widetilde{\alpha}$ é o isomorfismo inverso em virtude de \sref{II.2.6.5.1}, donde (i).
Seja $P$ o submódulo graduado $\alpha(M)$ de $\Gamma_\bullet(\widetilde{M})$;
como o funtor $M\mapsto\widetilde{M}$ é exato \sref{II.2.5.4}, a imagem de $\widetilde{M}$ por $\widetilde{\alpha}$ é igual a $\widetilde{P}$ e, portanto, em virtude de (i), $\widetilde{P}=(\Gamma_\bullet(\widetilde{M}))\supertilde$.
Ponha-se $Q=\Gamma_\bullet(\sh{G})\cap P$;
em virtude do que precede e de \sref{II.2.5.4}, tem-se $\widetilde{Q}=(\Gamma_\bullet(\sh{G}))\supertilde$; portanto, a restrição de $\beta$ a $\widetilde{Q}$ é um \emph{isomorfismo} desse $\sh{O}_X$-módulo sobre $\sh{G}$, por \sref{II.2.7.5}.
Mas, por definição de $N$ e por \sref{II.2.5.4}, a restrição do isomorfismo $\widetilde{\alpha}$ a $\widetilde{N}$ é um isomorfismo de $\widetilde{N}$ sobre $\widetilde{Q}$, donde a conclusão por \sref{II.2.6.5.1}.
\end{proof}

\subsection{Comportamentos funtoriais}
\label{subsection:II.2.8}

\begin{env}[2.8.1]
\label{II.2.8.1}
Sejam $S$ e $S'$ dois anéis graduados em graus positivos, e seja $\vphi:S'\to S$ um homomorfismo de anéis graduados.
Designaremos por $G(\vphi)$ a parte aberta de $X=\Proj(S)$ complementar de $V_+(\vphi(S'_+))$ ou, o que dá no mesmo, a reunião dos $D_+(\vphi(f'))$, onde $f'$ percorre o conjunto dos elementos homogêneos de $S'_+$.
A restrição a $G(\vphi)$ da aplicação contínua ${}^a\vphi$ de $\Spec(S)$ em $\Spec(S')$ \sref[I]{I.1.2.1} é, portanto, uma aplicação contínua de $G(\vphi)$ em $\Proj(S')$, que ainda denotaremos por ${}^a\vphi$, por abuso de linguagem.
Se $f'\in S'_+$ é homogêneo, tem-se
\[
\label{II.2.8.1.1}
  {}^a\vphi^{-1}(D_+(f'))=D_+(\vphi(f'))
\tag{2.8.1.1}
\]
levando-se em conta que ${}^a\vphi$ aplica $G(\vphi)$ em $\Proj(S')$ e \sref[I]{I.1.2.2.2}.
Por outro lado, o homomorfismo $\vphi$ define canonicamente (com as mesmas notações) um homomorfismo de anéis graduados $S'_{f'}\to S_f$, donde, por restrição aos elementos de grau~$0$,
\oldpage[II]{42}
um homomorfismo $S'_{(f')}\to S_{(f)}$, que designaremos por $\vphi_{(f)}$;
a ele corresponde \sref[I]{I.1.6.1} um morfismo de esquemas afins $({}^a\vphi_{(f)},\widetilde{\vphi}_{(f)}):\Spec(S_{(f)})\to\Spec(S'_{(f')})$.
Se se identifica canonicamente $\Spec(S_{(f)})$ com o esquema induzido por $\Proj(S)$ sobre $D_+(f)$ \sref{II.2.3.6}, fica definido um morfismo $\Phi_f:D_+(f)\to D_+(f')$, e ${}^a\vphi_{(f)}$ identifica-se com a restrição de ${}^a\vphi$ a $D_+(f)$.
Por outro lado, é imediato que, se $g'$ é um segundo elemento homogêneo de $S'_+$ e $g=\vphi(g')$, o diagrama
\[
  \xymatrix{
    D_+(f) \ar[r]^{\Phi_f}
    & D_+(f')
  \\D_+(fg) \ar[u] \ar[r]_{\Phi_{fg}}
    & D_+(f'g') \ar[u]
  }
\]
é comutativo, em razão da comutatividade do diagrama
\[
  \xymatrix{
    S'_{(f')} \ar[r]^{\vphi_{(f)}} \ar[d]_{\omega_{f'g',f'}}
    & S_{(f)} \ar[d]^{\omega_{fg,f}}
  \\S'_{(f'g')} \ar[r]_{\vphi_{(fg)}}
    & S_{(fg)}
  }
\]
Levando-se em conta a definição de $G(\vphi)$ e \sref{II.2.3.3.2}, vê-se, portanto, que:
\end{env}

\begin{proposition}[2.8.2]
\label{II.2.8.2}
Dado um homomorfismo de anéis graduados $\vphi:S'\to S$, existe um único morfismo $({}^a\vphi,\widetilde{\vphi})$ do pré-esquema induzido $G(\vphi)$ em $\Proj(S')$ (dito \emph{associado a $\vphi$} e denotado por $\Proj(\vphi)$), tal que, para todo elemento homogêneo $f'\in S'_+$, a restrição desse morfismo a $D_+(\vphi(f'))$ coincida com o morfismo associado ao homomorfismo $S'_{(f')}\to S_{(\vphi(f'))}$ correspondente a $\vphi$.
\end{proposition}

\begin{proof}
Observe-se que, com as notações precedentes, se $f'\in S'_d$, o diagrama
\[
\label{II.2.8.2.1}
  \xymatrix{
    S'_{(f')} \ar[r]^{\vphi_{(f)}} \ar[d]_{\sim}
    & S_{(f)} \ar[d]^{\sim}
  \\{S'}^{(d)}/(f'-1){S'}^{(d)} \ar[r]
    & S^{(d)}/(f-1)S^{(d)}
  }
\tag{2.8.2.1}
\]
é comutativo (sendo as setas verticais os isomorfismos de \sref{II.2.2.5}).
\end{proof}

\begin{corollary}[2.8.3]
\label{II.2.8.3}
\begin{enumerate}
  \item[\rm{(i)}] O morfismo $\Proj(\vphi)$ é afim.
  \item[\rm{(ii)}] Se $\Ker(\vphi)$ é nilpotente (e, em particular, se $\vphi$ é injetivo), o morfismo $\Proj(\vphi)$ é dominante.
\end{enumerate}
\end{corollary}

\begin{proof}
A afirmação~(i) é consequência imediata de \sref{II.2.8.2} e de \sref{II.2.8.1.1}.
Por outro lado, se $\Ker(\vphi)$ é nilpotente, para todo $f'$ homogêneo em $S'_+$ verifica-se imediatamente que $\Ker(\vphi_f)$ (com $f=\vphi(f')$) é nilpotente, e portanto também $\Ker(\vphi_{(f)})$;
a conclusão resulta, pois, de \sref{II.2.8.2} e de \sref[I]{I.1.2.7}.
\end{proof}
\oldpage[II]{43}
Observe-se que há, em geral, morfismos de $\Proj(S)$ em $\Proj(S')$ que não são afins e que, portanto, não provêm de homomorfismos de anéis graduados $S'\to S$;
um exemplo é o morfismo estrutural $\Proj(S)\to\Spec(A)$ quando $A$ é um corpo (sendo $\Spec(A)$ identificado com $\Proj(A[T])$, cf.~\sref{II.3.1.7});
isso resulta, com efeito, de \sref[I]{I.2.3.2}.

\begin{env}[2.8.4]
\label{II.2.8.4}
Seja $S''$ um terceiro anel graduado em graus positivos, seja $\vphi':S''\to S'$ um homomorfismo de anéis graduados, e ponha-se $\vphi''=\vphi\circ\vphi'$.
Em razão de \sref{II.2.8.1.1} e da fórmula ${}^a\vphi''=({}^a\vphi')\circ({}^a\vphi)$, verifica-se imediatamente que $G(\vphi'')\subset G(\vphi)$ e que, se $\Phi$, $\Phi'$ e $\Phi''$ são os morfismos associados respectivamente a $\vphi$, $\vphi'$ e $\vphi''$, tem-se $\Phi''=\Phi'\circ(\Phi|G(\vphi''))$.
\end{env}

\begin{env}[2.8.5]
\label{II.2.8.5}
Suponha-se que $S$ (resp. $S'$) seja uma $A$-álgebra graduada (resp. uma $A'$-álgebra graduada), e seja $\psi:A'\to A$ um homomorfismo de anéis tal que o diagrama
\[
\label{II.2.8.5.1}
  \xymatrix{
    A' \ar[r]^{\psi} \ar[d]
    & A \ar[d]
  \\S' \ar[r]_{\vphi}
    & S
  }
\tag{2.8.5.1}
\]
seja comutativo.
Pode-se então considerar $G(\vphi)$ e $\Proj(S')$ como esquemas sobre $\Spec(A)$ e $\Spec(A')$, respectivamente;
se $\Phi$ e $\Psi$ são os morfismos associados respectivamente a $\vphi$ e $\psi$, o diagrama
\[
  \xymatrix{
    G(\vphi) \ar[r]^{\Phi} \ar[d]
    & \Proj(S') \ar[d]
  \\\Spec(A) \ar[r]_{\Psi}
    & \Spec(A')
  }
\]
é então comutativo:
com efeito, basta demonstrá-lo para a restrição de $\Phi$ a $D_+(f)$, onde $f=\vphi(f')$, sendo $f'$ homogêneo em $S'_+$;
e isso resulta então da comutatividade do diagrama
\[
  \xymatrix{
    A' \ar[r]^{\psi} \ar[d]
    & A \ar[d]
  \\S'_{(f')} \ar[r]_{\vphi_{(f)}}
    & S_{(f)}
  }
\]
\end{env}

\begin{env}[2.8.6]
\label{II.2.8.6}
Seja agora $M$ um $S$-módulo graduado, e considere-se o $S'$-módulo $M_{[\vphi]}$, que é evidentemente graduado.
Seja $f'$ homogêneo em $S'_+$, e seja $f=\vphi(f')$;
sabe-se \sref[0]{0.1.5.2} que há um isomorfismo canônico $(M_{[\vphi]})_{f'}\simto(M_f)_{[\vphi_f]}$, e é imediato que esse isomorfismo conserva os graus, logo há um isomorfismo canônico $(M_{[\vphi]})_{(f')}\simto(M_{(f)})_{[\vphi_{(f)}]}$.
A esse isomorfismo corresponde canonicamente um isomorfismo de feixes $(M_{[\vphi]})\supertilde|D_+(f')\simto(\Phi_f)_*(\widetilde{M}|D_+(f))$ (\sref{II.2.5.2} e \sref[I]{I.1.6.3}).
Além disso,
\oldpage[II]{44}
se $g'$ é um segundo elemento homogêneo de $S'_+$ e $g=\vphi(g')$, o diagrama
\[
  \xymatrix{
    (M_{[\vphi]})_{(f')} \ar[r]^{\sim} \ar[d]
    & (M_{(f)})_{[\vphi_{(f)}]} \ar[d]
  \\(M_{[\vphi]})_{(f'g')} \ar[r]^{\sim}
    & (M_{(fg)})_{[\vphi_{(fg)}]}
  }
\]
é comutativo, donde se conclui imediatamente que o isomorfismo
\[
  (M_{[\vphi]})\supertilde|D_+(f'g')\simto(\Phi_{fg})_*(\widetilde{M}|D_+(fg))
\]
é a restrição a $D_+(f'g')$ do isomorfismo $(M_{[\vphi]})\supertilde|D_+(f')\simto(\Phi_f)_*(\widetilde{M}|D_+(f))$.
Como $\Phi_f$ é a restrição a $D_+(f)$ do morfismo $\Phi$, vê-se, portanto, levando-se em conta \sref{II.2.8.1.1} e pondo $X'=\Proj(S')$, que:
\end{env}

\begin{proposition}[2.8.7]
\label{II.2.8.7}
Existe um isomorfismo canônico funtorial do $\sh{O}_{X'}$-módulo $(M_{[\vphi]})\supertilde$ sobre o $\sh{O}_{X'}$-módulo $\Phi_*(\widetilde{M}|G(\vphi))$.
\end{proposition}

Deduz-se imediatamente uma aplicação canônica funtorial do conjunto dos $\vphi$-morfismos $M'\to M$ de um $S'$-módulo graduado no $S$-módulo graduado $M$ no conjunto dos $\Phi$-morfismos $\widetilde{M'}\to\widetilde{M}|G(\vphi)$.
Com as notações de \sref{II.2.8.4}, se $M''$ é um $S''$-módulo graduado, ao composto de um $\vphi$-morfismo $M'\to M$ e de um $\vphi'$-morfismo $M''\to M'$ corresponde canonicamente o composto de $\widetilde{M'}|G(\vphi')\to\widetilde{M}|G(\vphi'')$ e de $\widetilde{M''}\to\widetilde{M'}|G(\vphi')$.

\begin{proposition}[2.8.8]
\label{II.2.8.8}
Sob as hipóteses de \sref{II.2.8.1}, seja $M'$ um $S'$-módulo graduado.
Existe um homomorfismo canônico funtorial $\nu$ do $(\sh{O}_X|G(\vphi))$-módulo $\Phi^*(\widetilde{M'})$ no $(\sh{O}_X|G(\vphi))$-módulo $(M'\otimes_{S'}S)\supertilde|G(\vphi)$.
Se o ideal $S'_+$ é gerado por $S'_1$, $\nu$ é um isomorfismo.
\end{proposition}

\begin{proof}
Com efeito, para $f'\in S'_d$ ($d>0$), define-se um homomorfismo canônico funtorial de $S_{(f)}$-módulos (onde $f=\vphi(f')$)
\[
\label{II.2.8.8.1}
  \nu_f:M'_{(f')}\otimes_{S'_{(f')}}S_{(f)}\to(M'\otimes_{S'}S)_{(f)}
\tag{2.8.8.1}
\]
compondo o homomorfismo $M'_{(f')}\otimes_{S'_{(f')}}S_{(f)}\to M'_{f'}\otimes_{S'_{f'}}S_f$ com o isomorfismo canônico $M'_{f'}\otimes_{S'_{f'}}S_f\simto(M'\otimes_{S'}S)_f$ \sref[0]{0.1.5.4}, e observando que este último conserva os graus.
Verifica-se imediatamente a compatibilidade dos $\nu_f$ com os operadores de restrição de $D_+(f)$ a $D_+(fg)$ (para um segundo $g'\in S'_+$ e $g=\vphi(g')$), donde a definição do homomorfismo
\[
  \nu:\Phi^*(\widetilde{M'})\to(M'\otimes_{S'}S)\supertilde|G(\vphi)
\]
levando-se em conta \sref[I]{I.1.6.5}.
Para provar a segunda afirmação, basta mostrar que $\nu_f$ é um isomorfismo para todo $f'\in S'_1$, pois $G(\vphi)$ é então reunião dos $D_+(\vphi(f'))$.
Define-se primeiro uma aplicação $\bb{Z}$-bilinear $M'_m\times S_n\to M'_{(f')}\otimes_{S'_{(f')}}S_{(f)}$, fazendo corresponder a $(x',s)$ o elemento $(x'/{f'}^m)\otimes(s/f^n)$ (com a convenção de que $x'/{f'}^m$ é ${f'}^{-m}x'/1$
\oldpage[II]{45}
quando $m<0$).
Constata-se, como na demonstração de \sref{II.2.5.13}, que essa aplicação dá origem a um di-homomorfismo de módulos
\[
  \eta_f: M'\otimes_{S'}S \to M'_{(f')}\otimes_{S'_{(f')}}S_{(f)}.
\]
Além disso, se, para $r>0$, tem-se $f^r\sum_i(x'_i\otimes s_i)=0$, isso também se escreve $\sum_i({f'}^rx'_i\otimes s_i)=0$, donde, por \sref[0]{0.1.5.4}, $\sum_i({f'}^rx'_i/{f'}^{m_i+r})\otimes(s_i/f^{n_i})=0$, isto é, $\eta_f(\sum_i x'_i\otimes s_i)=0$, o que prova que $\eta_f$ se fatora em $M'\otimes_{S'}S\to(M'\otimes_{S'}S)_f\xrightarrow{\eta'_f}M'_{(f')}\otimes_{S'_{(f')}}S_{(f)}$;
verifica-se, enfim, que $\eta'_f$ e $\nu_f$ são isomorfismos inversos um do outro.

Em particular, resulta de \sref{II.2.1.2.1} que se tem um homomorfismo canônico
\[
\label{II.2.8.8.2}
  \Phi^*(\sh{O}_{X'}(n)) \simto \sh{O}_X(n)|G(\vphi)
\tag{2.8.8.2}
\]
para todo $n\in\bb{Z}$.
\end{proof}

\begin{env}[2.8.9]
\label{II.2.8.9}
Sejam $A$ e $A'$ dois anéis, e seja $\psi:A'\to A$ um homomorfismo de anéis que define um morfismo $\Psi:\Spec(A)\to\Spec(A')$.
Seja $S'$ uma $A'$-álgebra graduada em graus positivos, e ponhamos $S=S'\otimes_{A'}A$, que é evidentemente uma $A$-álgebra graduada pelos $S'_n\otimes_{A'}A$;
a aplicação $\vphi:s'\mapsto s'\otimes1$ é então um homomorfismo de anéis graduados que torna comutativo o diagrama \sref{II.2.8.5.1}.
Como, aqui, $S_+$ é o $A$-módulo gerado por $\vphi(S'_+)$, tem-se $G(\vphi)=\Proj(S)=X$;
daí, pondo $X'=\Proj(S')$, obtém-se o diagrama comutativo
\[
\label{II.2.8.9.1}
  \xymatrix{
    X \ar[r]^{\Phi} \ar[d]_p
    & X' \ar[d]
  \\Y \ar[r]_{\Psi}
    & Y'
  }
\tag{2.8.9.1}
\]

Seja agora $M'$ um $S'$-módulo graduado, e ponhamos $M=M'\otimes_{A'}A=M'\otimes_{S'}S$.
Nessas condições:
\end{env}

\begin{proposition}[2.8.10]
\label{II.2.8.10}
O diagrama \sref{II.2.8.9.1} identifica o esquema $X$ com o produto $X'\times_{Y'}Y$;
além disso, o homomorfismo canônico $\nu:\Phi^*(\widetilde{M'})\to\widetilde{M}$ \sref{II.2.8.8} é um isomorfismo.
\end{proposition}

\begin{proof}
A primeira afirmação estará demonstrada se provarmos que, para todo $f'$ homogêneo em $S'_+$, pondo $f=\vphi(f')$, as restrições de $\Phi$ e de $p$ a $D_+(f)$ identificam esse esquema com o produto $D_+(f')\times_{Y'}Y$ \sref[I]{I.3.2.6.2};
em outras palavras, basta \sref[I]{I.3.2.2} provar que $S_{(f)}$ se identifica canonicamente com $S'_{(f')}\otimes_{A'}A$, o que é imediato em virtude da existência do isomorfismo canônico $S_f\simto S'_{f'}\otimes_{A'}A$ que preserva os graus \sref[0]{0.1.5.4}.
A segunda afirmação resulta de que $M'_{(f')}\otimes_{S'_{(f')}}S_{(f)}$ se identifica, pelo que precede, com $M'_{(f')}\otimes_{A'}A$, e este último com $M_{(f)}$, pois $M_f$ se identifica canonicamente com $M'_{f'}\otimes_{A'}A$ por um isomorfismo que preserva os graus.
\end{proof}

\begin{corollary}[2.8.11]
\label{II.2.8.11}
Para todo inteiro $n\in\bb{Z}$, $\widetilde{M}(n)$ se identifica com $\Phi^*(\widetilde{M'}(n))=\widetilde{M'}(n)\otimes_{Y'}\sh{O}_Y$;
em particular, $\sh{O}_X(n)=\Phi^*(\sh{O}_{X'}(n))=\sh{O}_{X'}(n)\otimes_{Y'}\sh{O}_Y$.
\end{corollary}

\begin{proof}
Isso resulta de \sref{II.2.8.10} e de \sref{II.2.5.15}.
\end{proof}
\oldpage[II]{46}

\begin{env}[2.8.12]
\label{II.2.8.12}
Sob as hipóteses de \sref{II.2.8.9}, para $f'\in S'_d$ ($d>0$) e $f=\vphi(f')$, o diagrama
\[
\label{II.2.8.12.1}
  \xymatrix{
    M'_{(f')} \ar[r]^-{\sim} \ar[d]
    & {M'}^{(d)}/(f'-1){M'}^{(d)} \ar[d]
  \\M_{(f)} \ar[r]^-{\sim}
    & M^{(d)}/(f-1)M^{(d)}
  }
  \tag{2.8.12.1}
\]
(cf. \sref{II.2.2.5}) é comutativo.
\end{env}

\begin{env}[2.8.13]
\label{II.2.8.13}
Conservemos as notações e as hipóteses de \sref{II.2.8.9}, e seja $\sh{F}'$ um $\sh{O}_{X'}$-módulo;
se pusermos $\sh{F}=\Phi^*(\sh{F}')$, teremos, para todo $n\in\bb{Z}$, $\sh{F}(n)\simeq\Phi^*(\sh{F}'(n))$, em virtude de \sref{II.2.8.11} e de \sref[0]{0.4.3.3}.
Por conseguinte, segundo \sref[0]{0.3.7.1}, tem-se um homomorfismo canônico
\[
  \Gamma(\rho): \Gamma(X',\sh{F}'(n)) \to \Gamma(X,\sh{F}(n))
\]
que fornece um di-homomorfismo canônico de módulos graduados
\[
  \Gamma_\bullet(\sh{F}') \to \Gamma_\bullet(\sh{F}).
\]

Suponhamos que o ideal $S_+$ seja gerado por $S_1$ e que $\sh{F}'=\widetilde{M'}$, de modo que $\sh{F}=\widetilde{M}$, com $M=M'\otimes_{A'}A$.
Se $f'$ é homogêneo em $S'_+$ e $f=\vphi(f')$, vimos que $M_{(f)}=M'_{(f')}\otimes_{A'}A$, e o diagrama
\[
  \xymatrix{
    M'_0 \ar[r] \ar[d]
    & M'_{(f')} \ar[d]
    & =\Gamma(D_+(f'),\widetilde{M'})
  \\M_0 \ar[r]
    & M_{(f)}
    & =\Gamma(D_+(f),\widetilde{M})
  }
\]
é, portanto, comutativo;
conclui-se imediatamente dessa observação e da definição do homomorfismo $\alpha$ \sref{II.2.6.2} que o diagrama
\[
\label{II.2.8.13.1}
  \xymatrix{
    M' \ar[r]^-{\alpha_{M'}} \ar[d]
    & \Gamma_\bullet(\widetilde{M'}) \ar[d]
  \\M \ar[r]_-{\alpha_M}
    & \Gamma_\bullet(\widetilde{M})
  }
\tag{2.8.13.1}
\]
é comutativo.
Do mesmo modo, o diagrama
\[
\label{II.2.8.13.2}
  \xymatrix{
    (\Gamma_\bullet(\sh{F}'))\supertilde \ar[r]^-{\beta_{\sh{F}'}} \ar[d]
    & \sh{F}' \ar[d]
  \\(\Gamma_\bullet(\sh{F}))\supertilde \ar[r]_-{\beta_{\sh{F}}}
    & \sh{F}
  }
\tag{2.8.13.2}
\]
é comutativo (sendo a seta vertical da direita o $\Phi$-morfismo canônico $\sh{F}'\to\Phi^*(\sh{F}')=\sh{F}$).
\end{env}
\oldpage[II]{47}

\begin{env}[2.8.14]
\label{II.2.8.14}
Conservando ainda as notações e as hipóteses de \sref{II.2.8.9}, seja $N'$ outro $S'$-módulo graduado, e seja $N=N'\otimes_{A'}A$.
É imediato que os di-homomorfismos canônicos $M'\to M$ e $N'\to N$ fornecem um di-homomorfismo $M'\otimes_{S'}N'\to M\otimes_S N$ (relativo ao homomorfismo canônico de anéis $S'\to S$) e, por conseguinte, também um $S$-homomorfismo de grau~$0$, $(M'\otimes_{S'}N')\otimes_{A'}A\to M\otimes_S N$, ao qual corresponde (levando em conta \sref{II.2.8.10}) um homomorfismo de $\sh{O}_X$-módulos
\[
\label{II.2.8.14.1}
  \Phi^*((M'\otimes_{S'}N')\supertilde) \to (M\otimes_S N)\supertilde.
\tag{2.8.14.1}
\]

Além disso, verifica-se imediatamente que o diagrama
\[
\label{II.2.8.14.2}
  \xymatrix{
    \Phi^*(\widetilde{M'}\otimes_{\sh{O}_{X'}}\widetilde{N'}) \ar[r]^\sim \ar[d]_{\Phi^*(\lambda)}
    & \widetilde{M}\otimes_{\sh{O}_X}\widetilde{N} \ar[d]^\lambda
    & =\Phi^*(\widetilde{M'})\otimes_{\sh{O}_X}\Phi^*(\widetilde{N'})
  \\\Phi^*((M'\otimes_{S'}N')\supertilde) \ar[r]
    & (M\otimes_S N)\supertilde
  }
\tag{2.8.14.2}
\]
é comutativo, sendo a primeira linha o isomorfismo canônico \sref[0]{0.4.3.3}.
Se o ideal $S'_+$ é gerado por $S'_1$, é claro que $S_+$ é gerado por $S_1$, e as duas setas verticais de \sref{II.2.8.14.2} são então isomorfismos \sref{II.2.5.13};
o mesmo ocorre, portanto, com \sref{II.2.8.14.1}.

Tem-se igualmente um di-homomorfismo canônico $\Hom_{S'}(M',N')\to\Hom_S(M,N)$ fazendo corresponder a um homomorfismo $u'$ de grau~$k$ o homomorfismo $u'\otimes1$, que também é de grau~$k$;
deduz-se daí, ainda, um $S$-homomorfismo de grau~$0$
\[
  (\Hom_{S'}(M',N'))\otimes_{A'}A \to \Hom_S(M,N)
\]
e, portanto, um homomorfismo de $\sh{O}_X$-módulos
\[
\label{II.2.8.14.3}
  \Phi^*((\Hom_{S'}(M',N'))\supertilde) \to (\Hom_S(M,N))\supertilde.
\tag{2.8.14.3}
\]

Além disso, o diagrama
\[
  \xymatrix{
    \Phi^*((\Hom_{S'}(M',N'))\supertilde) \ar[r] \ar[d]_{\Phi^*(\mu)}
    & (\Hom_S(M,N))\supertilde \ar[d]^\mu
  \\\Phi^*(\shHom_{\sh{O}_{X'}}(\widetilde{M'},\widetilde{N'})) \ar[r]
    & \shHom_{\sh{O}_X}(\widetilde{M},\widetilde{N})
  }
\]
é comutativo (sendo a segunda linha horizontal o homomorfismo canônico \sref[0]{0.4.4.6}).
\end{env}
\oldpage[II]{48}

\begin{env}[2.8.15]
\label{II.2.8.15}
Com as notações e as hipóteses de \sref{II.2.8.1}, resulta de \sref{II.2.4.7} que não se altera o morfismo $\Phi$, a menos de isomorfismos, substituindo $S_0$ e $S'_0$ por $\bb{Z}$ e $\vphi_0$ pela aplicação identidade e, em seguida, substituindo $S$ e $S'$ por $S^{(d)}$ e ${S'}^{(d)}$, respectivamente ($d>0$), e $\vphi$ por sua restrição $\vphi^{(d)}$ a $S^{(d)}$.
\end{env}

\subsection{Subpré-esquemas fechados de um esquema $\operatorname{Proj}(S)$}
\label{subsection:II.2.9}

\begin{env}[2.9.1]
\label{II.2.9.1}
Se $\vphi:S\to S'$ é um homomorfismo de anéis graduados, diz-se que $\vphi$ é (TN)-\emph{sobrejetivo} (resp. (TN)-\emph{injetivo}, (TN)-\emph{bijetivo}) se existe um inteiro $n$ tal que, para $k\geq n$, $\vphi_k:S_k\to S'_k$ seja \emph{sobrejetivo} (resp. \emph{injetivo}, \emph{bijetivo}).
Resulta então de \sref{II.2.8.15} que o estudo de $\Phi$ pode ser reduzido ao caso em que $\vphi$ é sobrejetivo (resp. injetivo, bijetivo).
Em vez de dizer que $\vphi$ é (TN)-bijetivo, diz-se também que se trata de um (TN)-\emph{isomorfismo}.
\end{env}

\begin{proposition}[2.9.2]
\label{II.2.9.2}
Seja $S$ um anel graduado em graus positivos, e seja $X=\Proj(S)$.
\begin{enumerate}
  \item[{\rm(i)}] Se $\vphi:S\to S'$ é um homomorfismo (TN)-sobrejetivo de anéis graduados, o morfismo correspondente $\Phi$ \sref{II.2.8.1} está definido em todo $\Proj(S')$ e é uma imersão fechada de $\Proj(S')$ em $X$.
    Se $\mathfrak{J}$ é o núcleo de $\vphi$, o subpré-esquema fechado de $X$ associado a $\Phi$ é definido pelo feixe quase-coerente de ideais $\widetilde{\mathfrak{J}}$ de $\sh{O}_X$.
  \item[{\rm(ii)}] Suponhamos, além disso, que o ideal $S_+$ seja gerado por um número finito de elementos homogêneos de grau~$1$.
    Seja $X'$ um subpré-esquema fechado de $X$, definido por um feixe quase-coerente de ideais $\sh{J}$ de $\sh{O}_X$.
    Seja $\mathfrak{J}$ o ideal graduado de $S$ que é a imagem inversa de $\Gamma_\bullet(\sh{J})$ pelo homomorfismo canônico $\alpha:S\to\Gamma_\bullet(\sh{O}_X)$ \sref{II.2.6.2}, e ponhamos $S'=S/\mathfrak{J}$.
    Então $X'$ é o subpré-esquema associado à imersão fechada $\Proj(S')\to X$ correspondente ao homomorfismo canônico de anéis graduados $S\to S'$.
\end{enumerate}
\end{proposition}

\begin{proof}
\medskip\noindent
\begin{enumerate}
  \item[{\rm(i)}] Pode-se supor que $\vphi$ seja sobrejetivo \sref{II.2.9.1}.
    Como, por hipótese, $\vphi(S_+)$ gera $S'_+$, tem-se $G(\vphi)=\Proj(S')$.
    Por outro lado, a segunda afirmação verifica-se localmente em $X$;
    seja, portanto, $f$ um elemento homogêneo de $S_+$, e ponhamos $f'=\vphi(f)$.
    Como $\vphi$ é um homomorfismo sobrejetivo de anéis graduados, verifica-se imediatamente que $\vphi_{(f')}:S_{(f)}\to S'_{(f')}$ é sobrejetivo e que seu núcleo é $\mathfrak{J}_{(f)}$, o que conclui a prova de (i) \sref[I]{I.4.2.3}.
  \item[{\rm(ii)}] Em virtude de (i), reduzimo-nos a verificar que o homomorfismo $\widetilde{j}:\widetilde{\mathfrak{J}}\to\sh{O}_X$ deduzido da injeção canônica $j:\mathfrak{J}\to S$ é um isomorfismo de $\widetilde{\mathfrak{J}}$ sobre $\sh{J}$, o que resulta de \sref{II.2.7.11}.
\end{enumerate}
\end{proof}

Note-se que $\mathfrak{J}$ é o \emph{maior} dos ideais graduados $\mathfrak{J}'$ de $S$ tais que $\widetilde{j}(\widetilde{\mathfrak{J'}})=\sh{J}$, pois se verifica imediatamente, remontando às definições \sref{II.2.6.2}, que essa relação implica $\alpha(\mathfrak{J}')\subset\Gamma_\bullet(\sh{J})$.

\begin{corollary}[2.9.3]
\label{II.2.9.3}
Suponhamos satisfeitas as hipóteses de \sref{II.2.9.2}[(i)] e, além disso, que o ideal $S_+$ seja gerado por $S_1$;
então $\Phi^*((S(n))\supertilde)$ é canonicamente isomorfo a $(S'(n))\supertilde$ para todo $n\in\bb{Z}$ e, por conseguinte, $\Phi^*(\sh{F}(n))$ é canonicamente isomorfo a $\Phi^*(\sh{F})(n)$ para todo $\sh{O}_X$-módulo $\sh{F}$.
\end{corollary}

\begin{proof}
Esse é um caso particular de \sref{II.2.8.8}, levando em conta \sref{II.2.5.10.2}.
\end{proof}

\begin{corollary}[2.9.4]
\label{II.2.9.4}
Suponhamos satisfeitas as hipóteses de \sref{II.2.9.2}[(ii)].
Para que o subpré-esquema fechado $X'$ de $X$ seja integral, é necessário e suficiente que o ideal graduado $\mathfrak{J}$ seja primo em $S$.
\end{corollary}
\oldpage[II]{49}

\begin{proof}
Como $X'$ é isomorfo a $\Proj(S/\mathfrak{J})$, a condição é suficiente em virtude de \sref{II.2.4.4}.
Para ver que ela é necessária, consideremos a sequência exata $0\to\sh{J}\to\sh{O}_X\to\sh{O}_X/\sh{J}$, que fornece a sequência exata
\[
  0 \to \Gamma_\bullet(\sh{J}) \to \Gamma_\bullet(\sh{O}_X) \to \Gamma_\bullet(\sh{O}_X/\sh{J}).
\]

Basta provar que, se $f\in S_m$ e $g\in S_n$ são tais que a imagem de $\alpha_{n+m}(fg)$ em $\Gamma_\bullet(\sh{O}_X/\sh{J})$ é nula, então uma das imagens de $\alpha_m(f)$ e $\alpha_n(g)$ é nula.
Ora, por definição, essas imagens são seções dos $(\sh{O}_X/\sh{J})$-módulos invertíveis $\sh{L}=(\sh{O}_X/\sh{J})(m)$ e $\sh{L}'=(\sh{O}_X/\sh{J})(n)$ sobre o esquema integral $X'$;
a hipótese implica que o produto dessas duas seções é nulo em $\sh{L}\otimes\sh{L}'$ (\sref{II.2.9.3} e \sref{II.2.5.14.1}) e, portanto, uma delas é nula em virtude de \sref[I]{I.7.4.4}.
\end{proof}

\begin{corollary}[2.9.5]
\label{II.2.9.5}
Sejam $A$ um anel, $M$ um $A$-módulo, $S$ uma $A$-álgebra graduada gerada pelo conjunto $S_1$ dos elementos homogêneos de grau~$1$, $u:M\to S_1$ um homomorfismo sobrejetivo de $A$-módulos, e $\overline{u}:\bb{S}(M)\to S$ o homomorfismo (de $A$-álgebras) da álgebra simétrica $\bb{S}(M)$ de $M$ em $S$ que prolonga $u$.
O morfismo correspondente a $\overline{u}$ é então uma imersão fechada de $\Proj(S)$ em $\Proj(\bb{S}(M))$.
\end{corollary}

\begin{proof}
Com efeito, $\overline{u}$ é sobrejetivo por hipótese, e basta aplicar \sref{II.2.9.2}.
\end{proof}

\section{Espectro homogêneo de um feixe de álgebras graduadas}
\label{section:II.3}

\subsection{Espectro homogêneo de uma $\sh{O}_Y$-álgebra graduada quase-coerente}
\label{subsection:II.3.1}
\label{II.3.1}

\begin{env}[3.1.1]
\label{II.3.1.1}
Sejam $Y$ um pré-esquema, $\sh{S}$ uma $\sh{O}_Y$-álgebra graduada e $\sh{M}$ um $\sh{S}$-módulo graduado.
Se $\sh{S}$ é \emph{quase-coerente}, cada uma de suas componentes homogêneas $\sh{S}_n$ é um $\sh{O}_Y$-módulo \emph{quase-coerente}, por ser a imagem de $\sh{S}$ por um homomorfismo de $\sh{S}$ em si mesma (\sref[I]{I.1.3.8} e \sref[I]{I.1.3.9});
do mesmo modo, se $\sh{M}$ é quase-coerente como $\sh{O}_Y$-módulo, suas componentes homogêneas $\sh{M}_n$ também o são, e reciprocamente.
Se $d$ é um inteiro $>0$, designaremos por $\sh{S}^{(d)}$ a soma direta dos $\sh{O}_Y$-módulos $\sh{S}_{nd}$ ($n\in\bb{Z}$), que é quase-coerente se $\sh{S}$ o for \sref[I]{I.1.3.9};
para todo inteiro $k$ tal que $0\leq k\leq d-1$, designaremos por $\sh{M}^{(d,k)}$ (ou $\sh{M}^{(d)}$, para $k=0$) a soma direta dos $\sh{M}_{nd+k}$ ($n\in\bb{Z}$), que é um $\sh{S}^{(d)}$-módulo graduado, quase-coerente quando $\sh{S}$ e $\sh{M}$ são quase-coerentes \sref[I]{I.9.6.1}.
Denotaremos por $\sh{M}(n)$ o $\sh{S}$-módulo graduado tal que $(\sh{M}(n))_k=\sh{M}_{n+k}$ para todo $k\in\bb{Z}$;
se $\sh{S}$ e $\sh{M}$ são quase-coerentes, $\sh{M}(n)$ é um $\sh{S}$-módulo graduado quase-coerente \sref[I]{I.9.6.1}.

Diremos que $\sh{M}$ é um $\sh{S}$-módulo graduado \emph{de tipo finito} (resp. que admite uma \emph{apresentação finita}) se, para todo $y\in Y$, existem uma vizinhança aberta $U$ de $y$ e inteiros $n_i$ (resp. inteiros $m_i$ e $n_j$) tais que haja um homomorfismo sobrejetivo de grau~$0$, $\bigoplus_{i=1}^r(\sh{S}(n_i)|U)\to\sh{M}|U$ (resp. tais que $\sh{M}|U$ seja isomorfo ao conúcleo de um homomorfismo de grau~$0$, $\bigoplus_{i=1}^r(\sh{S}(m_i)|U)\to\bigoplus_{j=1}^s(\sh{S}(n_j)|U)$).

Seja $U$ um aberto afim de $Y$, e seja $A=\Gamma(U,\sh{O}_Y)$ seu anel;
por hipótese, a $(\sh{O}_Y|U)$-álgebra graduada $\sh{S}|U$ é isomorfa a $\widetilde{S}$, onde $S=\Gamma(U,\sh{S})$ é uma $A$-álgebra
\oldpage[II]{50}
graduada \sref[I]{I.1.4.3};
ponhamos $X_U=\Proj(\Gamma(U,\sh{S}))$.
Seja $U'\subset U$ um segundo aberto afim de $Y$, de anel $A'$, e seja $j:U'\to U$ a injeção canônica, que corresponde ao homomorfismo de restrição $A\to A'$;
temos $\sh{S}|U'=j^*(\sh{S}|U)$ e, por conseguinte, $S'=\Gamma(U',\sh{S})$ identifica-se canonicamente com $S\otimes_AA'$ \sref[I]{I.1.6.4}.
Conclui-se de \sref{II.2.8.10} que $X_{U'}$ se identifica canonicamente com $X_U\times_U U'$ e, portanto, também com $f_U^{-1}(U')$, onde $f_U$ denota o morfismo estrutural $X_U\to U$ \sref[I]{I.4.4.1}.
Denotaremos por $\sigma_{U',U}$ o isomorfismo canônico $f_U^{-1}(U')\simto X_{U'}$ assim definido, e por $\rho_{U',U}$ a imersão aberta $X_{U'}\to X_U$ obtida compondo $\sigma_{U',U}^{-1}$ com a injeção canônica $f_U^{-1}(U')\to X_U$.
É imediato que, se $U''\subset U'$ é um terceiro aberto afim de $Y$, então $\rho_{U'',U}=\rho_{U',U}\circ\rho_{U'',U'}$.
\end{env}

\begin{proposition}[3.1.2]
\label{II.3.1.2}
Seja $Y$ um pré-esquema.
Para toda $\sh{O}_Y$-álgebra graduada quase-coerente em graus positivos $\sh{S}$, existe um pré-esquema $X$ sobre $Y$, único a menos de $Y$-isomorfismo, que possui a seguinte propriedade:
se $f:X\to Y$ é o morfismo estrutural, então, para todo aberto afim $U$ de $Y$, existe um \emph{isomorfismo} $\eta_U$ do pré-esquema induzido $f^{-1}(U)$ sobre $X_U=\Proj(\Gamma(U,\sh{S}))$ tal que, se $V$ é um segundo aberto afim de $Y$ contido em $U$, o diagrama
\[
\label{II.3.1.2.1}
  \xymatrix{
    f^{-1}(V) \ar[r]^{\eta_V} \ar[d]
    & X_V \ar[d]^{\rho_{V,U}}
  \\f^{-1}(U) \ar[r]_{\eta_U}
    & X_U
  }
\tag{3.1.2.1}
\]
é comutativo.
\end{proposition}

\begin{proof}
Para dois abertos afins $U$ e $V$ de $Y$, seja $X_{U,V}$ o pré-esquema induzido por $X_U$ sobre $f_U^{-1}(U\cap V)$;
vamos definir um $Y$-isomorfismo $\theta_{U,V}:X_{V,U}\simto X_{U,V}$.
Para isso, consideremos um aberto afim $W\subset U\cap V$:
compondo os isomorfismos
\[
  f_U^{-1}(W)
  \xrightarrow{\sigma_{W,U}} X_W
  \xrightarrow{\sigma_{W,V}^{-1}} f_V^{-1}(W),
\]
obtemos um isomorfismo $\tau_W$, e verifica-se imediatamente que, se $W'\subset W$ é um aberto afim, então $\tau_{W'}$ é a restrição de $\tau_W$ a $f_U^{-1}(W')$;
os $\tau_W$ são, portanto, as restrições de um $Y$-isomorfismo $\theta_{V,U}$.
Além disso, se $U$, $V$ e $W$ são três abertos afins de $Y$, e se $\theta'_{U,V}$, $\theta'_{V,W}$ e $\theta'_{U,W}$ são, respectivamente, as restrições de $\theta_{U,V}$, $\theta_{V,W}$ e $\theta_{U,W}$ às imagens inversas de $U\cap V\cap W$ em $X_V$, $X_W$ e $X_W$, resulta das definições anteriores que $\theta'_{U,V}\circ\theta'_{V,W}=\theta'_{U,W}$.
A existência de um $X$ que satisfaça as propriedades enunciadas resulta, portanto, de \sref[I]{I.2.3.1};
sua unicidade a menos de $Y$-isomorfismo é trivial, levando-se em conta \sref{II.3.1.2.1}.
\end{proof}

\begin{env}[3.1.3]
\label{II.3.1.3}
Diremos que o pré-esquema $X$ definido em \sref{II.3.1.2} é o \emph{espectro homogêneo} da $\sh{O}_Y$-álgebra graduada quase-coerente $\sh{S}$, e o denotaremos por $\Proj(\sh{S})$.
É imediato que $\Proj(\sh{S})$ é \emph{separado sobre $Y$} (\sref{II.2.4.2} e \sref[I]{I.5.5.5});
se $\sh{S}$ é uma $\sh{O}_Y$-álgebra \emph{de tipo finito} \sref[I]{I.9.6.2}, então $\Proj(\sh{S})$ é \emph{de tipo finito} sobre $Y$ (\sref{II.2.7.1}[(ii)] e \sref[I]{I.6.3.1}).

Se $f$ é o morfismo estrutural $X\to Y$, é imediato que, para todo pré-esquema induzido por $Y$ sobre um aberto $U$ de $Y$, $f^{-1}(U)$ se identifica com o espectro homogêneo $\Proj(\sh{S}|U)$.
\end{env}

\begin{proposition}[3.1.4]
\label{II.3.1.4}
Seja $f\in\Gamma(Y,\sh{S}_d)$, com $d>0$.
Existe uma parte aberta $X_f$ do espaço subjacente a $X=\Proj(\sh{S})$ que possui a seguinte propriedade:
para todo aberto afim $U$ de $Y$,
\oldpage[II]{51}
tem-se $X_f\cap\vphi^{-1}(U)=D_+(f|U)$ em $\vphi^{-1}(U)$, identificado com $X_U=\Proj(\Gamma(U,\sh{S}))$, onde $\vphi$ denota o morfismo estrutural $X\to Y$.
Além disso, o pré-esquema induzido sobre $X_f$ é afim sobre $Y$ e é canonicamente isomorfo a $\Spec(\sh{S}^{(d)}/(f-1)\sh{S}^{(d)})$ \sref{II.1.3.1}.
\end{proposition}

\begin{proof}
Temos $f|U\in\Gamma(U,\sh{S}_d)=(\Gamma(U,\sh{S}))_d$.
Se $U$ e $U'$ são dois abertos afins de $Y$ tais que $U'\subset U$, então $f|U'$ é a imagem de $f|U$ pelo homomorfismo de restrição
\[
  \Gamma(U,\sh{S}) \to \Gamma(U',\sh{S});
\]
portanto, $D_+(f|U')$ é igual (com as notações de \sref{II.3.1.1}) ao pré-esquema induzido sobre a imagem inversa $\rho_{U',U}^{-1}(D_+(f|U))$ em $X_{U'}$ \sref{II.2.8.1};
daí a primeira afirmação.
Além disso, o pré-esquema induzido sobre $D_+(f|U)$ por $X_U$ identifica-se canonicamente com $\Spec((\Gamma(U,\sh{S}))_{(f|U)})$, e essas identificações são compatíveis com os homomorfismos de restrição \sref{II.2.8.1};
a segunda afirmação resulta então de \sref{II.2.2.5} e da comutatividade do diagrama \sref{II.2.8.2.1}.
\end{proof}

Diremos também que $X_f$ (como aberto no espaço subjacente $X$) é o conjunto dos $x\in X$ nos quais $f$ \emph{não se anula}.

\begin{corollary}[3.1.5]
\label{II.3.1.5}
Se $f\in\Gamma(Y,\sh{S}_d)$ e $g\in\Gamma(Y,\sh{S}_e)$, então
\[
\label{II.3.1.5.1}
  X_{fg}=X_f\cap X_g.
\tag{3.1.5.1}
\]
\end{corollary}

\begin{proof}
Basta considerar a interseção dos dois membros com um conjunto $\vphi^{-1}(U)$, onde $U$ é um aberto afim de $Y$, e aplicar a fórmula \sref{II.2.3.3.2}.
\end{proof}

\begin{corollary}[3.1.6]
\label{II.3.1.6}
Seja $(f_\alpha)$ uma família de seções de $\sh{S}$ sobre $Y$ tal que $f_\alpha\in\Gamma(Y,\sh{S}_{d_\alpha})$;
se o feixe de ideais de $\sh{S}$ gerado por essa família \sref[0]{0.5.1.1} contém todos os $\sh{S}_n$ a partir de certo índice, então o espaço subjacente $X$ é a união dos $X_{f_\alpha}$.
\end{corollary}

\begin{proof}
Para todo aberto afim $U$ de $Y$, $\vphi^{-1}(U)$ é a união dos $X_{f_\alpha}\cap\vphi^{-1}(U)$ \sref{II.2.3.14}.
\end{proof}

\begin{corollary}[3.1.7]
\label{II.3.1.7}
Seja $\sh{A}$ uma $\sh{O}_Y$-álgebra quase-coerente;
ponhamos
\[
  \sh{S}=\sh{A}[T]=\sh{A}\otimes_{\bb{Z}}\bb{Z}[T],
\]
onde $T$ é uma indeterminada (sendo $\bb{Z}$ e $\bb{Z}[T]$ considerados como feixes simples sobre $Y$).
Então $X=\Proj(\sh{S})$ identifica-se canonicamente com $\Spec(\sh{A})$.
Em particular, $\Proj(\sh{O}_Y[T])$ identifica-se com $Y$.
\end{corollary}

\begin{proof}
Aplicando \sref{II.3.1.6} à única seção $f\in\Gamma(Y,\sh{S})$ que é igual a $T$ em cada ponto de $Y$, vemos que $X_f=X$.
Além disso, aqui temos $d=1$, e $\sh{S}^{(1)}/(f-1)\sh{S}^{(1)}=\sh{S}/(f-1)\sh{S}$ é canonicamente isomorfa a $\sh{A}$, donde resulta o corolário \sref{II.3.1.4}.
\end{proof}

Seja $g\in\Gamma(Y,\sh{O}_Y)$;
se tomarmos $\sh{S}=\sh{O}_Y[T]$, então $g\in\Gamma(Y,\sh{S}_0)$;
seja
\[
  h=gT\in\Gamma(Y,\sh{S}_1).
\]
Se $X=\Proj(\sh{S})$, a identificação canônica definida em \sref{II.3.1.7} identifica $X_h$ com o aberto $Y_g$ de $Y$ (no sentido de \sref[0]{0.5.5.2}):
com efeito, podemos restringir-nos ao caso em que $Y=\Spec(A)$ é afim, e tudo se reduz então (levando-se em conta \sref{II.2.2.5}) ao fato de que o anel de frações $A_g$ se identifica canonicamente com $A[T]/(gT-1)A[T]$ \sref[0]{0.1.2.3}.

\begin{proposition}[3.1.8]
\label{II.3.1.8}
Seja $\sh{S}$ uma $\sh{O}_Y$-álgebra graduada quase-coerente em graus positivos.
Então:
\begin{enumerate}
  \item[(i)] Para todo $d>0$, existe um $Y$-isomorfismo canônico de $\Proj(\sh{S})$ sobre $\Proj(\sh{S}^{(d)})$.
\oldpage[II]{52}
  \item[(ii)] Seja $\sh{S}'$ a $\sh{O}_Y$-álgebra graduada que é a soma direta de $\sh{O}_Y$ e dos $\sh{S}_n$ ($n>0$);
    então $\Proj(\sh{S}')$ e $\Proj(\sh{S})$ são canonicamente $Y$-isomorfos.
  \item[(iii)] Seja $\sh{L}$ um $\sh{O}_Y$-módulo invertível \sref[0]{0.5.4.1}, e seja $\sh{S}_{(\sh{L})}$ a $\sh{O}_Y$-álgebra graduada que é a soma direta dos $\sh{S}_d\otimes\sh{L}^{\otimes d}$ ($d\geq0$);
    então $\Proj(\sh{S})$ e $\Proj(\sh{S}_{(\sh{L})})$ são canonicamente $Y$-isomorfos.
\end{enumerate}
\end{proposition}

\begin{proof}
Em cada um dos três casos, basta definir o isomorfismo localmente sobre $Y$, pois a verificação das compatibilidades com as operações de restrição de um aberto a um aberto menor é trivial.
Podemos, portanto, supor que $Y$ é afim, e então (i) resulta de \sref{II.2.4.7}[(i)] e (ii), de \sref{II.2.4.8}.
Quanto a (iii), se supusermos, além disso, que $\sh{L}$ é isomorfo a $\sh{O}_Y$ (o que é permitido, pois a questão é local sobre $Y$), o isomorfismo entre $\Proj(\sh{S})$ e $\Proj(\sh{S}_{(\sh{L})})$ é evidente;
para definir um isomorfismo \emph{canônico}, sejam $Y=\Spec(A)$ e $\sh{S}=\widetilde{S}$, onde $S$ é uma $A$-álgebra graduada, e seja $c$ um gerador do $A$-módulo livre $L$ tal que $\sh{L}=\widetilde{L}$;
então, para todo $n>0$, $x_n\mapsto x_n\otimes c^{\otimes n}$ é um $A$-isomorfismo de $S_n$ sobre $S_n\otimes L^{\otimes n}$, e esses $A$-isomorfismos definem um $A$-isomorfismo de álgebras graduadas $\vphi_c:S\to S_{(L)}=\bigoplus_{n\geq0}S_n\otimes L^{\otimes n}$.
Seja então $f\in S_+$ homogêneo de grau $d$;
para todo $x\in S_{nd}$, temos
\[
  \frac{x\otimes c^{\otimes nd}}{(f\otimes c^{\otimes d})^n}
  =
  \frac{x\otimes(\varepsilon c)^{\otimes nd}}{(f\otimes(\varepsilon c)^{\otimes d})^n}
\]
para todo elemento inversível $\varepsilon\in A$, o que mostra que o isomorfismo $S_{(f)}\to(S_{(L)})_{(f\otimes c^{\otimes d})}$ deduzido de $\vphi_c$ é \emph{independente} do gerador $c$ considerado de $L$, e assim conclui a demonstração.
\end{proof}

\begin{env}[3.1.9]
\label{II.3.1.9}
Recordemos (\sref[0]{0.4.1.3} e \sref[I]{I.1.3.14}) que, para que a $\sh{O}_Y$-álgebra graduada quase-coerente $\sh{S}$ seja \emph{gerada pelo $\sh{O}_Y$-módulo $\sh{S}_1$}, é necessário e suficiente que exista um recobrimento $(U_\alpha)$ de $Y$ por abertos afins tais que a álgebra graduada $\Gamma(U_\alpha,\sh{S})$ sobre $\Gamma(U_\alpha,\sh{S}_0)$ seja gerada pelo conjunto $\Gamma(U_\alpha,\sh{S}_1)$ de seus elementos homogêneos de grau $1$.
Para todo aberto $V$ de $Y$, $\sh{S}|V$ é então gerada pelo $(\sh{O}_Y|V)$-módulo $\sh{S}_1|V$.
\end{env}

\begin{proposition}[3.1.10]
\label{II.3.1.10}
Suponhamos que exista um recobrimento finito de $Y$ por abertos afins $(U_i)$ tal que cada álgebra graduada $\Gamma(U_i,\sh{S})$ seja de tipo finito sobre $\Gamma(U_i,\sh{O}_Y)$.
Então existe $d>0$ tal que $\sh{S}^{(d)}$ é gerada por $\sh{S}_d$, sendo $\sh{S}_d$ um $\sh{O}_Y$-módulo de tipo finito.
\end{proposition}

\begin{proof}
Com efeito, resulta de \sref{II.2.1.6}[(v)] que, para cada $i$, existe um inteiro $m_i$ tal que $\Gamma(U_i,\sh{S}_{nm_i})=(\Gamma(U_i,\sh{S}_{m_i}))^n$ para todo $n>0$;
basta tomar para $d$ um múltiplo comum de todos os $m_i$, levando-se em conta \sref{II.2.1.6}[(i)].
\end{proof}

\begin{corollary}[3.1.11]
\label{II.3.1.11}
Sob as hipóteses de \sref{II.3.1.10}, $\Proj(\sh{S})$ é $Y$-isomorfo a um espectro homogêneo $\Proj(\sh{S}')$, onde $\sh{S}'$ é uma $\sh{O}_Y$-álgebra graduada gerada por $\sh{S}'_1$, sendo $\sh{S}'_1$ um $\sh{O}_Y$-módulo de tipo finito.
\end{corollary}

\begin{proof}
Basta tomar $\sh{S}'=\sh{S}^{(d)}$, onde $d$ satisfaz a propriedade de \sref{II.3.1.10}, e aplicar \sref{II.3.1.8}[(i)].
\end{proof}

\begin{env}[3.1.12]
\label{II.3.1.12}
Se $\sh{S}$ é uma $\sh{O}_Y$-álgebra graduada quase-coerente em graus positivos, sabemos, por \sref[I]{I.5.1.1}, que seu \emph{nilradical} $\sh{N}$ é um $\sh{O}_Y$-módulo quase-coerente;
diremos que $\sh{N}_+=\sh{N}\cap\sh{S}_+$ é o \emph{nilradical de $\sh{S}_+$};
trata-se de um $\sh{S}_0$-módulo graduado quase-coerente, pois podemos reduzir-nos imediatamente ao caso em que $Y$ é afim, e a afirmação resulta então de \sref{II.2.1.10}.
Para todo $y\in Y$, $(\sh{N}_+)_y$ é, portanto, o nilradical de $(\sh{S}_+)_y=(\sh{S}_y)_+$ \sref[I]{I.5.1.1}.
Diremos que a $\sh{O}_Y$-álgebra graduada $\sh{S}$ é \emph{essencialmente reduzida} se $\sh{N}_+=0$, o que equivale
\oldpage[II]{53}
a dizer que $\sh{S}_y$ é uma $\sh{O}_y$-álgebra graduada essencialmente reduzida para todo $y\in Y$.
Para toda $\sh{O}_Y$-álgebra graduada $\sh{S}$, $\sh{S}/\sh{N}_+$ é essencialmente reduzida.

Diremos que $\sh{S}$ é \emph{íntegra} se, para todo $y\in Y$, $\sh{S}_y$ é um anel íntegro e, além disso, $(\sh{S}_y)_+=(\sh{S}_+)_y\neq0$.
\end{env}

\begin{proposition}[3.1.13]
\label{II.3.1.13}
Seja $\sh{S}$ uma $\sh{O}_Y$-álgebra graduada em graus positivos.
Se $X=\Proj(\sh{S})$, então o $Y$-esquema $X_\red$ é canonicamente isomorfo a $\Proj(\sh{S}/\sh{N}_+)$;
em particular, se $\sh{S}$ é essencialmente reduzida, então $X$ é reduzido.
\end{proposition}

\begin{proof}
O fato de $X'=\Proj(\sh{S}/\sh{N}_+)$ ser reduzido resulta imediatamente de \sref{II.2.4.4}[(i)], pois essa propriedade é local;
além disso, para todo aberto afim $U\subset Y$, ${\vphi'}^{-1}(U)$ é igual a $(\vphi^{-1}(U))_\red$ (onde $\vphi$ e $\vphi'$ denotam, respectivamente, os morfismos estruturais $X\to Y$ e $X'\to Y$);
verifica-se imediatamente que os $U$-morfismos canônicos ${\vphi'}^{-1}(U)\to\vphi^{-1}(U)$ são compatíveis com as operações de restrição e definem, portanto, uma imersão fechada $X'\to X$ que é um homeomorfismo dos espaços subjacentes;
daí a conclusão \sref[I]{I.5.1.2}.
\end{proof}

\begin{proposition}[3.1.14]
\label{II.3.1.14}
Sejam $Y$ um pré-esquema integral e $\sh{S}$ uma $\sh{O}_Y$-álgebra graduada quase-coerente tal que $\sh{S}_0=\sh{O}_Y$.
\begin{enumerate}
  \item[(i)] Se $\sh{S}$ é íntegra \sref{II.3.1.12}, então $X=\Proj(\sh{S})$ é integral, e o morfismo estrutural $\vphi:X\to Y$ é dominante.
  \item[(ii)] Suponhamos, além disso, que $\sh{S}$ seja essencialmente reduzida.
    Então, reciprocamente, se $X$ é integral e $\vphi$ é dominante, $\sh{S}$ é íntegra.
\end{enumerate}
\end{proposition}

\begin{proof}
\medskip\noindent
\begin{enumerate}
  \item[(i)] Se $(U_\alpha)$ é uma base de $Y$ formada por abertos afins não vazios, basta demonstrar a proposição quando se substitui $Y$ por um dos $U_\alpha$ e $\sh{S}$ por $\sh{S}|U_\alpha$:
    com efeito, por um lado, resultará que os espaços subjacentes $\vphi^{-1}(U_\alpha)$ são abertos irredutíveis (e, portanto, não vazios) de $X$ tais que $\vphi^{-1}(U_\alpha)\cap\vphi^{-1}(U_\beta)\neq\emp$ para todo par de índices (pois $U_\alpha\cap U_\beta$ contém algum $U_\gamma$), de modo que $X$ é irredutível \sref[0]{0.2.1.4};
    por outro lado, $X$ será reduzido, pois essa propriedade é local, e portanto $X$ será integral, com $\vphi(X)$ denso em $Y$.

    Suponhamos, então, que $Y=\Spec(A)$, onde $A$ é íntegro \sref[I]{I.5.1.4}, e que $\sh{S}=\widetilde{S}$, onde $S$ é uma $A$-álgebra graduada;
    a hipótese é que, para todo $y\in Y$, $\widetilde{S}_y=S_y$ é um anel graduado íntegro tal que $(S_y)_+\neq0$.
    Basta provar que $S$ é um anel \emph{íntegro}, pois então teremos $S_+\neq0$ e poderemos aplicar \sref{II.2.4.4}[(ii)].
    Sejam $f,g\neq0$ dois elementos de $S$, e suponhamos que $fg=0$;
    para todo $y\in Y$, teremos então $(f/1)(g/1)=0$ em $S_y$ e, portanto, por hipótese, $f/1=0$ ou $g/1=0$.
    Suponhamos, por exemplo, que $f/1=0$ em $S_y$;
    isso significa que existe $a\in A$ tal que $a\notin\mathfrak{j}_y$ e $af=0$;
    temos então, para \emph{todo} $z\in Y$, $(a/1)(f/1)=0$ no anel \emph{íntegro} $S_z$ e, como $a/1\neq0$ (pois $A$ é íntegro), $f/1=0$, o que implica $f=0$.
  \item[(ii)] Como a questão é local sobre $Y$, podemos supor novamente que $Y=\Spec(A)$, com $A$ íntegro, e que $\sh{S}=\widetilde{S}$.
    Por hipótese, para todo $y\in Y$, $(S_y)_+$ não contém nenhum elemento nilpotente não nulo, e o mesmo vale para $(S_0)_y=A_y$ por hipótese;
    logo, $S_y$ é um anel reduzido para todo $y\in Y$, e concluímos, antes de tudo, que o próprio $S$ é reduzido \sref[I]{I.5.1.1}.
    A hipótese de que $X$ é integral implica, por outro lado, que $S$ é essencialmente íntegro \sref{II.2.4.4}[(ii)], e tudo se reduz então a demonstrar que o anulador $\mathfrak{J}$ de $S_+$ em $A=S_0$ se reduz
\oldpage[II]{54}
a $0$ \sref{II.2.1.11}.
Se não fosse assim, teríamos $(S_h)_+=0$ para algum $h\neq0$ em $\mathfrak{J}$ e, portanto, por \sref{II.3.1.1}, $\vphi^{-1}(D(h))=\emp$, e $\vphi(X)$ não seria denso em $Y$, contrariamente à hipótese (pois $D(h)\neq\emp$, já que $h$ não é nilpotente).
\end{enumerate}
\end{proof}

\subsection{Feixe sobre $\operatorname{Proj}(\mathcal{S})$ associado a um $\mathcal{S}$-módulo graduado}
\label{subsection:II.3.2}

\begin{env}[3.2.1]
\label{II.3.2.1}
Sejam $Y$ um pré-esquema, $\sh{S}$ uma $\sh{O}_Y$-álgebra graduada quase-coerente em graus positivos e $\sh{M}$ um $\sh{S}$-módulo graduado quase-coerente (sobre $(Y,\sh{O}_Y)$ ou, o que é equivalente \sref[I]{I.9.6.1}, sobre o espaço anelado $(Y,\sh{S})$).
Com as notações de \sref{II.3.1.1}, denotaremos por $\widetilde{\sh{M}}_U$ o $\sh{O}_{X_U}$-módulo quase-coerente $(\Gamma(U,\sh{M}))\supertilde$;
para $U'\subset U$, $\Gamma(U',\sh{M})$ identifica-se canonicamente com $\Gamma(U,\sh{M})\otimes_AA'$ \sref[I]{I.1.6.4};
portanto, temos $\widetilde{\sh{M}}_{U'}=\rho_{U',U}^*(\widetilde{\sh{M}}_U)$ \sref{II.2.8.11}.
\end{env}

\begin{proposition}[3.2.2]
\label{II.3.2.2}
Existe sobre $\Proj(\sh{S})=X$ um único $\sh{O}_X$-módulo quase-coerente $\widetilde{\sh{M}}$ tal que, para todo aberto afim $U$ de $Y$, temos $\eta_U^*((\Gamma(U,\sh{M}))\supertilde)=\widetilde{\sh{M}}|f^{-1}(U)$ (onde $\eta_U$ denota o isomorfismo $f^{-1}(U)\simto\Proj(\Gamma(U,\sh{S}))$), sendo $f$ o morfismo estrutural $X\to Y$.
\end{proposition}

\begin{proof}
Como $\rho_{U',U}$ se identifica com o morfismo de injeção $f^{-1}(U')\to f^{-1}(U)$ \sref{II.3.1.2.1}, a proposição resulta imediatamente da relação $\widetilde{\sh{M}}_{U'}=\rho_{U',U}^*(\widetilde{\sh{M}}_U)$ e do princípio de colagem dos feixes \sref[0]{0.3.3.1}.
\end{proof}

Diremos que $\widetilde{\sh{M}}$ é o $\sh{O}_X$-módulo \emph{associado ao} $\sh{S}$-módulo graduado quase-coerente $\sh{M}$.

\begin{proposition}[3.2.3]
\label{II.3.2.3}
Seja $\sh{M}$ um $\sh{S}$-módulo graduado quase-coerente, e seja $f\in\Gamma(Y,\sh{S}_d)$ ($d>0$).
Se $\xi_f$ é o isomorfismo canônico de $X_f$ sobre o $Y$-pré-esquema $Z_f=\Spec(\sh{S}^{(d)}/(f-1)\sh{S}^{(d)})$ \sref{II.3.1.4}, então $(\xi_f)_*(\widetilde{\sh{M}}|X_f)$ é o $\sh{O}_{Z_f}$-módulo $(\sh{M}^{(d)}/(f-1)\sh{M}^{(d)})\supertilde$ \sref{II.1.4.3}.
\end{proposition}

\begin{proof}
Como a questão é local sobre $Y$, somos imediatamente reduzidos a \sref{II.2.2.5}, levando-se em conta a comutatividade do diagrama \sref{II.2.8.12.1}.
\end{proof}

\begin{proposition}[3.2.4]
\label{II.3.2.4}
A correspondência $\sh{M}\mapsto\widetilde{\sh{M}}$ é um funtor covariante aditivo exato da categoria dos $\sh{S}$-módulos graduados quase-coerentes para a categoria dos $\sh{O}_X$-módulos quase-coerentes, e comuta com os limites indutivos e as somas diretas.
\end{proposition}

\begin{proof}
Como a questão é local sobre $Y$, ela se reduz a \sref[I]{I.1.3.11}, \sref[I]{I.1.3.9} e \sref{II.2.5.4}.
\end{proof}

Em particular, se $\sh{N}$ é um sub-$\sh{S}$-módulo graduado quase-coerente de $\sh{M}$, então $\widetilde{\sh{N}}$ identifica-se canonicamente com um sub-$\sh{O}_X$-módulo quase-coerente de $\widetilde{\sh{M}}$;
mais particularmente, para todo feixe graduado quase-coerente $\sh{J}$ de ideais de $\sh{S}$, $\widetilde{\sh{J}}$ é um feixe quase-coerente de ideais de $\sh{O}_X$.

Se $\sh{M}$ é um $\sh{S}$-módulo graduado quase-coerente e $\sh{I}$ é um feixe quase-coerente de ideais de $\sh{O}_Y$, então $\sh{I}\sh{M}$ é um sub-$\sh{S}$-módulo graduado quase-coerente de $\sh{M}$, e temos
\[
\label{II.3.2.4.1}
  (\sh{I}\sh{M})\supertilde=\sh{I}\cdot\widetilde{\sh{M}}
\tag{3.2.4.1}
\]
(onde o segundo membro tem o sentido definido em \sref[0]{0.4.3.5}).
Basta, com efeito, verificar essa fórmula quando $Y=\Spec(A)$ é afim, $\sh{S}=\widetilde{S}$, onde $S$ é uma $A$-álgebra graduada, e $\sh{M}=\widetilde{M}$,
\oldpage[II]{55}
onde $M$ é um $S$-módulo graduado, e $\sh{I}=\widetilde{\mathfrak{I}}$, onde $\mathfrak{I}$ é um ideal de $A$.
Para todo elemento homogêneo $f$ de $S_+$, a restrição a $D_+(f)=\Spec(S_{(f)})$ do primeiro membro de \sref{II.3.2.4.1} está associada a $(\mathfrak{I}M)_{(f)}=\mathfrak{I}\cdot M_{(f)}$, e o mesmo se dá com a restrição do segundo membro, em virtude de \sref[I]{I.1.3.13} e \sref[I]{I.1.6.9}.

\begin{proposition}[3.2.5]
\label{II.3.2.5}
Seja $f\in\Gamma(Y,\sh{S}_d)$ ($d>0$).
Sobre o aberto $X_f$, o $(\sh{O}_X|X_f)$-módulo $(\sh{S}(nd))\supertilde|X_f$ é canonicamente isomorfo a $\sh{O}_X|X_f$ para todo $n\in\bb{Z}$.
Em particular, se a $\sh{O}_Y$-álgebra $\sh{S}$ é gerada por $\sh{S}_1$ \sref{II.3.1.9}, os $\sh{O}_X$-módulos $(\sh{S}(n))\supertilde$ são invertíveis para todo $n\in\bb{Z}$.
\end{proposition}

\begin{proof}
Com efeito, para todo aberto afim $U$ de $Y$, definiu-se em \sref{II.2.5.7} um isomorfismo canônico de $(\sh{S}(nd))\supertilde|(X_f\cap\vphi^{-1}(U))$ sobre $\sh{O}_X|(X_f\cap\vphi^{-1}(U))$, levando em conta \sref{II.3.1.4} (onde $\vphi$ é o morfismo estrutural $X\to Y$);
é imediato verificar que esses isomorfismos são compatíveis com a restrição de $U$ a um aberto afim $U'\subset U$, donde a primeira afirmação.
Para estabelecer a segunda, basta observar que, se $\sh{S}$ é gerada por $\sh{S}_1$, existe um recobrimento $(U_\alpha)$ de $Y$ por abertos afins tais que $\Gamma(U_\alpha,\sh{S})$ seja gerada por $\Gamma(U_\alpha,\sh{S})_1=\Gamma(U_\alpha,\sh{S}_1)$;
somos então reduzidos a aplicar o resultado de \sref{II.2.5.9}, pois a propriedade de ser invertível é local.
\end{proof}

Poremos ainda, para todo $n\in\bb{Z}$,
\[
\label{II.3.2.5.1}
  \sh{O}_X(n)=(\sh{S}(n))\supertilde
\tag{3.2.5.1}
\]
e, para todo $\sh{O}_X$-módulo $\sh{F}$,
\[
\label{II.3.2.5.2}
  \sh{F}(n)=\sh{F}\otimes_{\sh{O}_X}\sh{O}_X(n).
\tag{3.2.5.2}
\]

Resulta imediatamente dessas definições que, para todo aberto $U$ de $Y$, se tem
\[
  ((\sh{S}|U)(n))\supertilde=\sh{O}_X(n)|f^{-1}(U),
\]
onde $f$ é o morfismo estrutural $X\to Y$.

\begin{proposition}[3.2.6]
\label{II.3.2.6}
Sejam $\sh{M}$ e $\sh{N}$ dois $\sh{S}$-módulos graduados quase-coerentes.
Existe um homomorfismo canônico funtorial (em $\sh{M}$ e $\sh{N}$)
\[
\label{II.3.2.6.1}
  \lambda:\widetilde{\sh{M}}\otimes_{\sh{O}_X}\widetilde{\sh{N}}\to(\sh{M}\otimes_{\sh{S}}\sh{N})\supertilde
\tag{3.2.6.1}
\]
e um homomorfismo canônico funtorial (em $\sh{M}$ e $\sh{N}$)
\[
\label{II.3.2.6.2}
  \mu:(\shHom_{\sh{S}}(\sh{M},\sh{N}))\supertilde\to\shHom_{\sh{O}_X}(\widetilde{\sh{M}},\widetilde{\sh{N}}).
\tag{3.2.6.2}
\]

Além disso, se $\sh{S}$ é gerada por $\sh{S}_1$ \sref{II.3.1.9}, $\lambda$ é um isomorfismo;
se, além disso, $\sh{M}$ admite uma apresentação finita \sref{II.3.1.1}, $\mu$ é um isomorfismo.
\end{proposition}

\begin{proof}
Os homomorfismos $\lambda$ e $\mu$ foram definidos em \sref{II.2.5.11.2} e \sref{II.2.5.12.2} quando $Y$ é afim;
como essas definições são locais, elas se transportam imediatamente ao caso geral considerado aqui, levando em conta \sref{II.2.8.14}.
\end{proof}

\begin{corollary}[3.2.7]
\label{II.3.2.7}
Se $\sh{S}$ é gerada por $\sh{S}_1$, tem-se, quaisquer que sejam $m,n\in\bb{Z}$,
\[
\label{II.3.2.7.1}
  \sh{O}_X(m)\otimes_{\sh{O}_X}\sh{O}_X(n)=\sh{O}_X(m+n)
\tag{3.2.7.1}
\]
\[
\label{II.3.2.7.2}
  \sh{O}_X(n)=(\sh{O}_X(1))^{\otimes n}
\tag{3.2.7.2}
\]
a menos de isomorfismos canônicos.
\end{corollary}
\oldpage[II]{56}

\begin{corollary}[3.2.8]
\label{II.3.2.8}
Suponhamos que $\sh{S}$ seja gerada por $\sh{S}_1$.
Para todo $\sh{S}$-módulo graduado $\sh{M}$ e todo $n\in\bb{Z}$, tem-se
\[
\label{II.3.2.8.1}
  (\sh{M}(n))\supertilde=\widetilde{\sh{M}}(n)
\tag{3.2.8.1}
\]
a menos de isomorfismo canônico.
\end{corollary}

\begin{proof}
Isso resulta das propriedades correspondentes para $Y$ afim (\sref{II.2.5.14} e \sref{II.2.5.15}), bem como de \sref{II.2.8.11}.
\end{proof}

\begin{remarks}[3.2.9]
\label{II.3.2.9}
\medskip\noindent
\begin{enumerate}
  \item[\rm{(i)}] Se $\sh{S}=\sh{A}[T]$, sendo $\sh{A}$ uma $\sh{O}_Y$-álgebra quase-coerente \sref{II.3.1.7}, verifica-se imediatamente que todos os $\sh{O}_X$-módulos invertíveis $\sh{O}_X(n)$ são canonicamente isomorfos a $\sh{O}_X$.

  Além disso, seja $\sh{N}$ um $\sh{A}$-módulo quase-coerente, e ponhamos $\sh{M}=\sh{N}\otimes_{\sh{A}}\sh{A}[T]$.
  Resulta então de \sref{II.3.2.3} e \sref{II.3.1.7} que, na identificação canônica de $X=\Proj(\sh{A}[T])$ com $X'=\Spec(\sh{A})$, o $\sh{O}_X$-módulo $\widetilde{\sh{M}}$ se identifica com o $\sh{O}_{X'}$-módulo $\widetilde{\sh{N}}$ associado a $\sh{N}$ (no sentido de \sref{II.1.4.3}).
  \item[\rm{(ii)}] Seja $\sh{S}$ uma $\sh{O}_Y$-álgebra graduada qualquer, e seja $\sh{S}'$ a $\sh{O}_Y$-álgebra graduada tal que $\sh{S}'_0=\sh{O}_Y$ e $\sh{S}'_n=\sh{S}_n$ para todo $n>0$;
  então o isomorfismo canônico de $X=\Proj(\sh{S})$ sobre $X'=\Proj(\sh{S}')$ \sref{II.3.1.8}[(ii)] identifica $\sh{O}_X(n)$ com $\sh{O}_{X'}(n)$ para todo $n\in\bb{Z}$.
  Isso resulta da mesma proposição para o caso afim \sref{II.2.5.16} e do fato de que essas identificações, para os abertos afins de $Y$, comutam com as operações de restrição.
  Do mesmo modo, seja $X^{(d)}=\Proj(\sh{S}^{(d)})$;
  então o isomorfismo canônico de $X$ sobre $X^{(d)}$ \sref{II.3.1.8}[(i)] identifica $\sh{O}_X(nd)$ com $\sh{O}_{X^{(d)}}(n)$ para todo $n\in\bb{Z}$.
\end{enumerate}
\end{remarks}

\begin{proposition}[3.2.10]
\label{II.3.2.10}
Seja $\sh{L}$ um $\sh{O}_Y$-módulo invertível, e seja $g$ o isomorfismo canônico de $X_{(\sh{L})}=\Proj(\sh{S}_{(\sh{L})})$ sobre $X=\Proj(\sh{S})$ \sref{II.3.1.8}[(iii)].
Então, para todo $n\in\bb{Z}$, $g_*(\sh{O}_{X_{(\sh{L})}}(n))$ é canonicamente isomorfo a $\sh{O}_X(n)\otimes_Y\sh{L}^{\otimes n}$.
\end{proposition}

\begin{proof}
Suponhamos primeiro que $Y$ seja afim, de anel $A$, e que $\sh{L}=\widetilde{L}$, onde $L$ é um $A$-módulo livre de posto um.
Com as notações da demonstração de \sref{II.3.1.8}[(iii)], define-se, para $f\in S_d$, um isomorfismo de $(S(n))_{(f)}\otimes_A L^{\otimes n}$ sobre $(S_{(L)}(n))_{(f\otimes c^{\otimes d})}$ fazendo corresponder a $(x/f^k)\otimes c^{\otimes n}$, onde $x\in S_{kd+n}$, o elemento $(x\otimes c^{\otimes(n+kd)})/(f\otimes c^{\otimes d})^k$;
é imediato que esse isomorfismo independe do gerador $c$ escolhido para $L$;
além disso, os isomorfismos assim definidos para cada $f\in S_+$ são compatíveis com as operações de restrição $D_+(f)\to D_+(fg)$.
Por fim, no caso geral, vê-se sem dificuldade, retomando as definições de \sref{II.3.1.1}, que os isomorfismos assim definidos para cada aberto afim $U$ de $Y$ são compatíveis com a passagem de $U$ a um aberto afim $U'\subset U$.
\end{proof}

\subsection{$\mathcal{S}$-módulo graduado associado a um feixe sobre $\operatorname{Proj}(\mathcal{S})$}
\label{subsection:II.3.3}

\emph{Em todo este número, supomos que a $\sh{O}_Y$-álgebra graduada $\sh{S}$ seja gerada por $\sh{S}_1$ \sref{II.3.1.9}.}
Recordemos que, em virtude de \sref{II.3.1.8}[(i)], essa restrição não é essencial, mediante as condições de finitude de \sref{II.3.1.10}.

\begin{env}[3.3.1]
\label{II.3.3.1}
Seja $p$ o morfismo estrutural $X=\Proj(\sh{S})\to Y$.
Para todo $\sh{O}_X$-módulo $\sh{F}$, poremos
\[
\label{II.3.3.1.1}
  \bbGamma_*(\sh{F})=\bigoplus_{n\in\bb{Z}}p_*(\sh{F}(n))
\tag{3.3.1.1}
\]
\oldpage[II]{57}
e, em particular,
\[
\label{II.3.3.1.2}
  \bbGamma_*(\sh{O}_X)=\bigoplus_{n\in\bb{Z}}p_*(\sh{O}_X(n)).
\tag{3.3.1.2}
\]

Sabe-se \sref[0]{0.4.2.2} que existe um homomorfismo canônico
\[
  p_*(\sh{F})\otimes_{\sh{O}_Y}p_*(\sh{G})\to p_*(\sh{F}\otimes_{\sh{O}_X}\sh{G})
\]
para dois $\sh{O}_X$-módulos $\sh{F}$ e $\sh{G}$;
deduz-se, pois, de \sref{II.3.2.7.1} que $\bbGamma_*(\sh{O}_X)$ é munido de uma estrutura de \emph{$\sh{O}_Y$-álgebra graduada}, e \sref{II.3.2.5.2} define igualmente sobre $\bbGamma_*(\sh{F})$ uma estrutura de \emph{módulo graduado sobre $\bbGamma_*(\sh{O}_X)$}.

Em virtude de \sref{II.3.2.5} e da exatidão à esquerda do funtor $p_*$ \sref[0]{0.4.2.1}, a correspondência $\sh{F}\mapsto\bbGamma_*(\sh{F})$ é um funtor covariante aditivo e \emph{exato à esquerda} da categoria dos $\sh{O}_X$-módulos para a categoria dos $\sh{O}_Y$-módulos graduados (onde os morfismos são os homomorfismos de grau~$0$).
Em particular, se $\sh{J}$ é um feixe de ideais de $\sh{O}_X$, então $\bbGamma_*(\sh{J})$ se identifica com um \emph{feixe graduado de ideais} de $\bbGamma_*(\sh{O}_X)$.
\end{env}

\begin{env}[3.3.2]
\label{II.3.3.2}
Seja $\sh{M}$ um $\sh{S}$-módulo graduado quase-coerente.
Para todo aberto afim $U$ de $Y$, definiu-se em \sref{II.2.6.2} um homomorfismo de grupos abelianos
\[
  \alpha_{0,U}:\Gamma(U,\sh{M}_0)\to\Gamma(p^{-1}(U),\widetilde{\sh{M}}).
\]
É imediato que esses homomorfismos comutam com as operações de restrição \sref{II.2.8.13.1} e definem, portanto (sem utilizar a hipótese de que $\sh{S}$ seja gerada por $\sh{S}_1$), um homomorfismo de feixes de grupos abelianos
\[
\label{II.3.3.2.1}
  \alpha_0:\sh{M}_0\to p_*(\widetilde{\sh{M}}).
\tag{3.3.2.1}
\]

Aplicando esse resultado a cada um dos $\sh{M}_n=(\sh{M}(n))_0$, e levando em conta \sref{II.3.2.8.1}, define-se um homomorfismo de feixes de grupos abelianos
\[
\label{II.3.3.2.2}
  \alpha_n:\sh{M}_n\to p_*(\widetilde{\sh{M}}(n))
  \qquad\text{para todo }n\in\bb{Z},
\tag{3.3.2.2}
\]
donde um homomorfismo funtorial (de grau~$0$) de feixes graduados de grupos abelianos
\[
\label{II.3.3.2.3}
  \alpha:\sh{M}\to\bbGamma_*(\widetilde{\sh{M}})
\tag{3.3.2.3}
\]
(também denotado por $\alpha_{\sh{M}}$).

Tomando, em particular, $\sh{M}=\sh{S}$, verifica-se que $\alpha:\sh{S}\to\bbGamma_*(\sh{O}_X)$ é um homomorfismo de $\sh{O}_Y$-álgebras graduadas e que \sref{II.3.3.2.3} é um di-homomorfismo de módulos graduados relativo a esse homomorfismo de álgebras graduadas.

Observemos ainda que a cada $\alpha_n$ corresponde \sref[0]{0.4.4.3} um homomorfismo canônico de $\sh{O}_X$-módulos
\[
\label{II.3.3.2.4}
  \alpha_n^\sharp:p^*(\sh{M}_n)\to\widetilde{\sh{M}}(n).
\tag{3.3.2.4}
\]

Verifica-se sem dificuldade que esse homomorfismo não é outro senão aquele que corresponde funtorialmente \sref{II.3.2.4} ao homomorfismo canônico (de grau~$0$) de $\sh{O}_Y$-módulos \emph{graduados}
\[
\label{II.3.3.2.5}
  \sh{M}_n\otimes_{\sh{O}_Y}\sh{S}\to\sh{M}(n),
\tag{3.3.2.5}
\]
\oldpage[II]{58}
onde a graduação do segundo membro provém naturalmente da de $\sh{M}$.
Podemos, com efeito, limitar-nos ao caso em que $Y=\Spec(A)$ é afim, $\sh{M}=\widetilde{M}$ e $\sh{S}=\widetilde{S}$, sendo a $A$-álgebra graduada $S$ gerada por $S_1$, de modo que, quando $f$ percorre $S_1$, os $D_+(f)$ formam um recobrimento de $X$.
Retomando as definições de \sref{II.2.6.2}, vê-se então, levando em conta \sref[I]{I.1.6.7}, que a restrição a $D_+(f)$ do homomorfismo \sref{II.3.3.2.4} corresponde \sref[I]{I.1.3.8} ao homomorfismo de $S_{(f)}$-módulos $M_n\otimes_A S_{(f)}\to(M(n))_{(f)}$ que faz corresponder a $x\otimes1$, onde $x\in M_n$, o elemento $x/1$;
isso demonstra a afirmação.
\end{env}

\begin{proposition}[3.3.3]
\label{II.3.3.3}
Para toda seção $f\in\Gamma(Y,\sh{S}_d)$ ($d>0$), $X_f$ é idêntico ao conjunto dos pontos de $X$ onde $\alpha_d(f)$ (considerada como seção de $\sh{O}_X(d)$) não se anula \sref[0]{0.5.5.2}.
\end{proposition}

\begin{proof}
($\alpha_d(f)$ é uma seção de $p_*(\sh{O}_X(d))$ sobre $Y$, mas, por definição, tal seção é também uma seção de $\sh{O}_X(d)$ sobre $X$ \sref[0]{0.4.2.1}.)
A definição de $X_f$ \sref{II.3.1.4} reduz ao caso em que $Y$ é afim, já tratado em \sref{II.2.6.3}.
\end{proof}

\begin{env}[3.3.4]
\label{II.3.3.4}
Daqui por diante, suporemos, além da hipótese do início deste número, que, para todo $\sh{O}_X$-módulo quase-coerente $\sh{F}$, os $p_*(\sh{F}(n))$ sejam quase-coerentes sobre $Y$, e que, por conseguinte, $\bbGamma_*(\sh{F})=\bigoplus_{n\in\bb{Z}}p_*(\sh{F}(n))$ seja também um $\sh{O}_Y$-módulo quase-coerente (\sref[I]{I.1.4.1} e \sref[I]{I.1.3.9});
essa circunstância ocorrerá sempre, em particular, se $X$ for de tipo finito sobre $Y$ \sref[I]{I.9.2.2}.
Conclui-se daí que $(\bbGamma_*(\sh{F}))\supertilde$ está definido e é um $\sh{O}_X$-módulo quase-coerente.
Para todo aberto afim $U$ de $Y$, tem-se (\sref[I]{I.1.3.9} e \sref[I]{I.2.5.4})
\begin{align*}
  \Big(\Gamma(U,\bigoplus_{n\in\bb{Z}}p_*(\sh{F}(n)))\Big)\supertilde
  &=\bigoplus_{n\in\bb{Z}}\Big(\Gamma(U,p_*(\sh{F}(n)))\Big)\supertilde
\\&=\bigoplus_{n\in\bb{Z}}\Big(\Gamma(p^{-1}(U),\sh{F}(n))\Big)\supertilde
\\&=\Big(\bigoplus_{n\in\bb{Z}}\Gamma(p^{-1}(U),\sh{F}(n))\Big)\supertilde
\\&=(\bbGamma_*(\sh{F}|p^{-1}(U)))\supertilde,
\end{align*}
e, por conseguinte, em virtude de \sref{II.2.6.4}, tem-se um homomorfismo canônico
\[
  \beta_U:\Big(\Gamma(U,\bigoplus_{n\in\bb{Z}}p_*(\sh{F}(n)))\Big)\supertilde\to\sh{F}|p^{-1}(U).
\]

Além disso, a comutatividade de \sref{II.2.8.13.2} mostra que esses homomorfismos comutam com as operações de restrição sobre $Y$;
deduz-se, pois, um homomorfismo canônico funtorial
\[
\label{II.3.3.4.1}
  \beta:(\bbGamma_*(\sh{F}))\supertilde\to\sh{F}
\tag{3.3.4.1}
\]
(também denotado por $\beta_{\sh{F}}$) para os $\sh{O}_X$-módulos quase-coerentes.
\end{env}

\begin{proposition}[3.3.5]
\label{II.3.3.5}
Sejam $\sh{M}$ um $\sh{S}$-módulo graduado quase-coerente e $\sh{F}$ um $\sh{O}_X$-módulo quase-coerente;
os homomorfismos compostos
\[
\label{II.3.3.5.1}
  \widetilde{\sh{M}}\xrightarrow{\widetilde{\alpha}}(\bbGamma_*(\widetilde{\sh{M}}))\supertilde\xrightarrow{\beta}\widetilde{\sh{M}}
\tag{3.3.5.1}
\]
\[
\label{II.3.3.5.2}
  \bbGamma_*(\sh{F})\xrightarrow{\alpha}\bbGamma_*((\bbGamma_*(\sh{F}))\supertilde)\xrightarrow{\bbGamma_*(\beta)}\bbGamma_*(\sh{F})
\tag{3.3.5.2}
\]
são os isomorfismos identidade.
\end{proposition}

\begin{proof}
A questão é, com efeito, local sobre $Y$, e somos reduzidos a \sref{II.2.6.5}.
\end{proof}
\subsection{Condições de finitude}
\label{subsection:II.3.4}

\oldpage[II]{59}
\begin{proposition}[3.4.1]
\label{II.3.4.1}
Sejam $Y$ um pré-esquema e $\sh{S}$ uma $\sh{O}_Y$-álgebra quase-coerente gerada por $\sh{S}_1$ \sref{II.3.1.9};
suponhamos, além disso, que $\sh{S}_1$ seja de tipo finito.
Então $X=\Proj(\sh{S})$ é de tipo finito sobre $Y$.
\end{proposition}

\begin{proof}
Com efeito, como a questão é local sobre $Y$, podemos supor que $Y$ seja afim, de anel $A$;
temos então $\sh{S}=\widetilde{S}$, onde $S=\Gamma(Y,\sh{S})$, e, por hipótese, $S$ é uma $A$-álgebra gerada por $S_1=\Gamma(Y,\sh{S}_1)$, podendo-se supor, além disso, que $S_1$ seja um $A$-módulo de tipo finito (\sref[I]{I.1.3.9} e \sref[I]{I.1.3.12}).
Por conseguinte, $S$ é uma $A$-álgebra graduada de tipo finito, e somos reduzidos a \sref{II.2.7.1}[(ii)].
\end{proof}

\begin{env}[3.4.2]
\label{II.3.4.2}
Seja $\sh{S}$ uma $\sh{O}_Y$-álgebra graduada quase-coerente;
para um $\sh{S}$-módulo graduado quase-coerente $\sh{M}$, consideraremos as condições de finitude seguintes:
\begin{enumerate}
  \item[(\textbf{TF})] Existe um inteiro $n$ tal que o $\sh{S}$-módulo $\bigoplus_{k\geq n}\sh{M}_k$ seja de tipo finito.
  \item[(\textbf{TN})] Existe um inteiro $n$ tal que $\sh{M}_k=0$ para $k\geq n$.
\end{enumerate}

Se $\sh{M}$ satisfaz (\textbf{TN}), então $\widetilde{\sh{M}}=0$, pois a questão é local sobre $Y$ \sref{II.2.7.2}.

Sejam $\sh{M}$ e $\sh{N}$ dois $\sh{S}$-módulos graduados quase-coerentes;
diz-se que um homomorfismo $u:\sh{M}\to\sh{N}$ de grau~$0$ é \emph{(\textbf{TN})-injetivo} (resp. \emph{(\textbf{TN})-sobrejetivo}, \emph{(\textbf{TN})-bijetivo}) se existe um inteiro $n$ tal que $u_k:\sh{M}_k\to\sh{N}_k$ seja injetivo (resp. sobrejetivo, bijetivo) para $k\geq n$;
então $\widetilde{u}:\widetilde{\sh{M}}\to\widetilde{\sh{N}}$ é injetivo (resp. sobrejetivo, bijetivo) em virtude de \sref{II.2.7.2}, pois a questão é local sobre $Y$, e levando em conta \sref[I]{I.1.3.9};
quando $u$ é (\textbf{TN})-bijetivo, diz-se também que $u$ é um \emph{(\textbf{TN})-isomorfismo}.
\end{env}

\begin{proposition}[3.4.3]
\label{II.3.4.3}
Sejam $Y$ um pré-esquema e $\sh{S}$ uma $\sh{O}_Y$-álgebra graduada quase-coerente gerada por $\sh{S}_1$, supondo-se que $\sh{S}_1$ seja de tipo finito.
Seja $\sh{M}$ um $\sh{S}$-módulo graduado quase-coerente.
\begin{enumerate}
  \item[(i)] Se $\sh{M}$ satisfaz a condição (\textbf{TF}), então $\widetilde{\sh{M}}$ é de tipo finito.
  \item[(ii)] Suponhamos que $\sh{M}$ satisfaça (\textbf{TF}); para que $\widetilde{\sh{M}}=0$, é necessário e suficiente que $\sh{M}$ satisfaça (\textbf{TN}).
\end{enumerate}
\end{proposition}

\begin{proof}
Como as questões são locais sobre $Y$, somos reduzidos ao caso em que $Y$ é afim, de anel $A$, $\sh{S}=\widetilde{S}$, onde $S$ é uma $A$-álgebra graduada tal que o ideal $S_+$ seja de tipo finito, e $\sh{M}=\widetilde{M}$, onde $M$ é um $S$-módulo graduado;
a proposição resulta então de \sref{II.2.7.3}.
\end{proof}

\begin{theorem}[3.4.4]
\label{II.3.4.4}
Sejam $Y$ um pré-esquema e $\sh{S}$ uma $\sh{O}_Y$-álgebra graduada quase-coerente gerada por $\sh{S}_1$, supondo-se que $\sh{S}_1$ seja de tipo finito;
seja $X=\Proj(\sh{S})$.
Para todo $\sh{O}_X$-módulo quase-coerente $\sh{F}$, o homomorfismo canônico $\beta$ \sref{II.3.3.4} é um isomorfismo.
\end{theorem}

\begin{proof}
Observemos primeiro que $\beta$ está definido em virtude de \sref{II.3.4.1}.
Para ver que $\beta$ é um isomorfismo, somos reduzidos ao caso em que $Y$ é afim, de anel $A$, $\sh{S}=\widetilde{S}$, onde $S$ é uma $A$-álgebra graduada gerada por $S_1$, e $S_1$ é um $A$-módulo de tipo finito.
Basta então aplicar \sref{II.2.7.5}.
\end{proof}

\begin{corollary}[3.4.5]
\label{II.3.4.5}
Sob as hipóteses de \sref{II.3.4.4}, todo $\sh{O}_X$-módulo quase-coerente $\sh{F}$ é isomorfo a um $\sh{O}_X$-módulo da forma $\widetilde{\sh{M}}$, onde $\sh{M}$ é um $\sh{S}$-módulo graduado quase-coerente.
Se, além disso, $\sh{F}$ é de tipo finito, e se supomos que $Y$ seja um esquema quase-compacto ou que o espaço subjacente a $Y$ seja noetheriano, podemos supor que $\sh{M}$ seja de tipo finito.
\end{corollary}
\oldpage[II]{60}
\begin{proof}
A primeira afirmação resulta imediatamente de \sref{II.3.4.4}, tomando $\sh{M}=\bbGamma_*(\sh{F})$.
Para estabelecer a segunda, basta provar que $\sh{M}$ é o limite indutivo de seus sub-$\sh{S}$-módulos \emph{graduados} de tipo finito $\sh{N}_\lambda$: com efeito, daí resultará que $\widetilde{\sh{M}}$ é o limite indutivo dos $\widetilde{\sh{N}}_\lambda$ \sref{II.3.2.4} e, portanto, que $\sh{F}$ é o limite indutivo dos $\beta(\widetilde{\sh{N}}_\lambda)$;
como $X$ é quase-compacto (\sref{II.3.4.1} e \sref[I]{I.6.3.1}) e $\sh{F}$ é de tipo finito, $\sh{F}$ será necessariamente igual a um dos $\beta(\widetilde{\sh{N}}_\lambda)$ \sref[0]{0.5.2.3}.

Para definir os $\sh{N}_\lambda$ que têm $\sh{M}$ por limite indutivo, basta considerar, para cada $n\in\bb{Z}$, o $\sh{O}_Y$-módulo quase-coerente $\sh{M}_n$, que é o limite indutivo de seus sub-$\sh{O}_Y$-módulos $\sh{M}_n^{(\mu_n)}$ de tipo finito, em virtude das hipóteses sobre $Y$ \sref[I]{I.9.4.9};
é imediato que $\sh{P}_{\mu_n}=\sh{S}\cdot\sh{M}_n^{(\mu_n)}$ é um $\sh{S}$-módulo graduado de tipo finito, e verifica-se imediatamente que, tomando como $\sh{N}_\lambda$ as somas finitas de $\sh{S}$-módulos da forma $\sh{P}_{\mu_n}$, obtêm-se os objetos procurados.
\end{proof}

\begin{corollary}[3.4.6]
\label{II.3.4.6}
Suponhamos verificadas as hipóteses de \sref{II.3.4.4} e, além disso, que o espaço subjacente a $Y$ seja quase-compacto;
seja $\sh{F}$ um $\sh{O}_X$-módulo quase-coerente de tipo finito.
Existe um inteiro $n_0$ tal que, para $n\geq n_0$, o homomorfismo canônico
\[
  \sigma:p^*(p_*(\sh{F}(n)))\to\sh{F}(n)
\]
\sref[0]{0.4.4.3} seja sobrejetivo.
\end{corollary}

\begin{proof}
Para todo $y\in Y$, seja $U$ uma vizinhança aberta afim de $y$ em $Y$.
Existe um inteiro $n_0(U)$ tal que, para $n\geq n_0(U)$, $\sh{F}(n)|p^{-1}(U)$ seja gerado por um número finito de suas seções sobre $p^{-1}(U)$ \sref{II.2.7.9};
mas estas últimas são imagens canônicas de seções de $p^*(p_*(\sh{F}(n)))$ sobre $p^{-1}(U)$ (\sref[0]{0.3.7.1} e \sref[0]{0.4.4.3}), logo $\sh{F}(n)|p^{-1}(U)$ é igual à imagem canônica de $p^*(p_*(\sh{F}(n)))|p^{-1}(U)$.
Finalmente, como $Y$ é quase-compacto, existe uma cobertura finita de $Y$ por abertos afins $U_i$ e, tomando como $n_0$ o maior dos $n_0(U_i)$, conclui-se a demonstração.
\end{proof}

\begin{remarks}[3.4.7]
\label{II.3.4.7}
Se $p=(\psi,\theta):X\to Y$ é um morfismo de espaços anelados e $\sh{F}$ um $\sh{O}_X$-módulo, o fato de o homomorfismo canônico $\sigma:p^*(p_*(\sh{F}))\to\sh{F}$ ser sobrejetivo explicita-se da seguinte maneira \sref[0]{0.4.4.1}: para todo $x\in X$ e toda seção $s$ de $\sh{F}$ sobre uma vizinhança aberta $V$ de $x$, existem uma vizinhança aberta $U$ de $p(x)$ em $Y$, um número finito de seções $t_i$ ($1\leq i\leq m$) de $\sh{F}$ sobre $p^{-1}(U)$, uma vizinhança $W\subset V\cap p^{-1}(U)$ de $x$ e seções $a_i$ ($1\leq i\leq m$) de $\sh{O}_X$ sobre $W$, tais que
\[
  s|W=\sum_{i=1}^m a_i\cdot(t_i|W).
\]
Quando $Y$ é um \emph{esquema afim} e $p_*(\sh{F})$ é \emph{quase-coerente}, essa condição equivale ao fato de $\sh{F}$ ser \emph{gerado por suas seções sobre $X$} \sref[0]{0.5.5.1}: com efeito, se $Y=\Spec(A)$, pode-se supor que $U=D(f)$, com $f\in A$;
existem um inteiro $n>0$ e seções $s_i$ de $\sh{F}$ sobre $X$ tais que $t_i$ seja a restrição a $p^{-1}(U)$ de $s_i/g^n$, com $g=\theta(f)$ (aplicando \sref[I]{I.1.4.1} a $p_*(\sh{F})$);
como $g$ é invertível sobre $p^{-1}(U)$, tem-se, portanto,
\[
  s|W=\sum_i b_i\cdot(s_i|W),
\]
com $b_i=a_i(g|W)^{-n}$, de onde resulta a afirmação.
Quando $Y$ é afim, o Corolário~\sref{II.3.4.6} recupera, portanto, \sref{II.2.7.9}.

Conclui-se daí que, quando $Y$ é um pré-esquema qualquer, as três condições
\oldpage[II]{61}
seguintes, para um $\sh{O}_X$-módulo quase-coerente $\sh{F}$ tal que $p_*(\sh{F})$ seja \emph{quase-coerente}, são equivalentes:
\begin{enumerate}
  \item[a)] \emph{O homomorfismo canônico $\sigma:p^*(p_*(\sh{F}))\to\sh{F}$ é sobrejetivo.}
  \item[b)] \emph{Existe um $\sh{O}_Y$-módulo quase-coerente $\sh{G}$ e um homomorfismo sobrejetivo $p^*(\sh{G})\to\sh{F}$.}
  \item[c)] \emph{Para todo aberto afim $U$ de $Y$, $\sh{F}|p^{-1}(U)$ é gerado por suas seções sobre $p^{-1}(U)$.}
\end{enumerate}

Com efeito, acaba-se de provar a equivalência entre \emph{a)} e \emph{c)}.
Por outro lado, é claro que \emph{a)} implica \emph{b)}, pois $p_*(\sh{F})$ é quase-coerente por hipótese.
Reciprocamente, todo homomorfismo $u:p^*(\sh{G})\to\sh{F}$ se fatora como
\[
  p^*(\sh{G})\to p^*(p_*(\sh{F}))\xrightarrow{\sigma}\sh{F}
\]
\sref[0]{0.3.5.4.4};
se $u$ é sobrejetivo, o mesmo ocorre com $\sigma$, o que prova que \emph{b)} implica \emph{a)}.
\end{remarks}

\begin{corollary}[3.4.8]
\label{II.3.4.8}
Suponhamos verificadas as hipóteses de \sref{II.3.4.4} e suponhamos, além disso, que $Y$ seja um esquema quase-compacto ou que o espaço subjacente a $Y$ seja noetheriano.
Seja $\sh{F}$ um $\sh{O}_X$-módulo quase-coerente de tipo finito;
existe então um inteiro $n_0$ tal que, para $n\geq n_0$, $\sh{F}$ seja isomorfo a um quociente de um $\sh{O}_X$-módulo da forma $(p^*(\sh{G}))(-n)$, onde $\sh{G}$ é um $\sh{O}_Y$-módulo quase-coerente de tipo finito (dependente de $n$).
\end{corollary}

\begin{proof}
Como o morfismo estrutural $X\to Y$ é separado e de tipo finito, $p_*(\sh{F}(n))$ é quase-coerente \sref[I]{I.9.2.2}[b)] e, portanto, limite indutivo de seus sub-$\sh{O}_Y$-módulos quase-coerentes de tipo finito, em virtude da hipótese sobre $Y$ \sref[I]{I.9.4.9}.
Deduz-se então de \sref{II.3.4.6}, \sref[0]{0.4.3.2} e \sref[0]{0.5.2.3} que $\sh{F}(n)$ é a imagem canônica de um $\sh{O}_X$-módulo da forma $p^*(\sh{G})$, onde $\sh{G}$ é um sub-$\sh{O}_Y$-módulo quase-coerente de tipo finito de $p_*(\sh{F}(n))$;
o corolário resulta então de \sref{II.3.2.5.2} e \sref{II.3.2.7.1}.
\end{proof}

\subsection{Comportamento funtorial}
\label{subsection:II.3.5}

\begin{env}[3.5.1]
\label{II.3.5.1}
Sejam $Y$ um pré-esquema, e $\sh{S}$ e $\sh{S}'$ duas $\sh{O}_Y$-álgebras graduadas quase-coerentes em graus positivos;
ponhamos $X=\Proj(\sh{S})$ e $X'=\Proj(\sh{S}')$, e sejam $p$ e $p'$ os morfismos estruturais de $X$ e $X'$ em $Y$.
Seja $\vphi:\sh{S}'\to\sh{S}$ um $\sh{O}_Y$-homomorfismo de álgebras graduadas.
Para todo aberto afim $U$ de $Y$, ponhamos $S_U=\Gamma(U,\sh{S})$ e $S'_U=\Gamma(U,\sh{S}')$;
o homomorfismo $\vphi$ define um homomorfismo $\vphi_U:S'_U\to S_U$ de $A_U$-álgebras graduadas, pondo $A_U=\Gamma(U,\sh{O}_Y)$.
Correspondem-lhe, em $p^{-1}(U)$, um subconjunto aberto $G(\vphi_U)$ e um morfismo $\Phi_U:G(\vphi_U)\to p'^{-1}(U)$ \sref{II.2.8.1}.
Além disso, se $V\subset U$ é um aberto afim, o diagrama
\[
\label{II.3.5.1.1}
  \xymatrix{
    S'_U \ar[r]^{\vphi_U} \ar[d]
    & S_U \ar[d]
  \\S'_V \ar[r]_{\vphi_V}
    & S_V
  }
\tag{3.5.1.1}
\]
é comutativo, e verifica-se imediatamente, a partir das definições \sref{II.2.8.1}, que
\[
  G(\vphi_V)=G(\vphi_U)\cap p^{-1}(V)
\]
e que $\Phi_V$ é a restrição de $\Phi_U$ a $G(\vphi_V)$.
Assim, definiu-se um subconjunto aberto $G(\vphi)$ de $X$ tal que $G(\vphi)\cap p^{-1}(U)=G(\vphi_U)$ para todo aberto afim $U\subset Y$, e um $Y$-morfismo
\oldpage[II]{62}
\emph{afim} $\Phi:G(\vphi)\to X'$, que se diz \emph{associado} a $\vphi$ e que denotaremos por $\Proj(\vphi)$.
Quando, para todo $y\in Y$, existe uma vizinhança afim $U$ de $y$ tal que o $\Gamma(U,\sh{O}_Y)$-módulo $\Gamma(U,\sh{S}_+)$ seja gerado por $\vphi(\Gamma(U,\sh{S}'_+))$, tem-se $G(\vphi_U)=p^{-1}(U)$ e, por conseguinte, $G(\vphi)=X$.
\end{env}

\begin{proposition}[3.5.2]
\label{II.3.5.2}
\medskip\noindent
\begin{enumerate}
  \item[(i)] Se $\sh{M}$ é um $\sh{S}$-módulo graduado quase-coerente, existe um isomorfismo canônico funtorial do $\sh{O}_{X'}$-módulo $(\sh{M}_{[\vphi]})\supertilde$ sobre o $\sh{O}_{X'}$-módulo $\Phi_*(\widetilde{\sh{M}}|G(\vphi))$.
  \item[(ii)] Se $\sh{M}'$ é um $\sh{S}'$-módulo graduado quase-coerente, existe um homomorfismo canônico funtorial $\nu$ do $(\sh{O}_X|G(\vphi))$-módulo $\Phi^*(\widetilde{\sh{M}'})$ para o $(\sh{O}_X|G(\vphi))$-módulo $(\sh{M}'\otimes_{\sh{S}'}\sh{S})\supertilde|G(\vphi)$.
    Se $\sh{S}'$ é gerada por $\sh{S}'_1$, $\nu$ é um isomorfismo.
\end{enumerate}
\end{proposition}

\begin{proof}
Os homomorfismos considerados já foram, com efeito, definidos quando $Y$ é afim (\sref{II.2.8.7} e \sref{II.2.8.8});
no caso geral, basta verificar que são compatíveis com as operações de restrição de um aberto afim de $Y$ a um aberto menor, o que resulta imediatamente da comutatividade de \sref{II.3.5.1.1}.
\end{proof}

Em particular, para todo $n\in\bb{Z}$, tem-se um homomorfismo canônico
\[
\label{II.3.5.2.1}
  \Phi^*(\sh{O}_{X'}(n))\to\sh{O}_X(n)|G(\vphi).
\tag{3.5.2.1}
\]

\begin{proposition}[3.5.3]
\label{II.3.5.3}
Sejam $Y$ e $Y'$ dois pré-esquemas, $\psi:Y'\to Y$ um morfismo e $\sh{S}$ uma $\sh{O}_Y$-álgebra graduada quase-coerente, e ponhamos $\sh{S}'=\psi^*(\sh{S})$.
Então o $Y'$-esquema $X'=\Proj(\sh{S}')$ identifica-se canonicamente com $\Proj(\sh{S})\times_Y Y'$.
Além disso, se $\sh{M}$ é um $\sh{S}$-módulo graduado quase-coerente, o $\sh{O}_{X'}$-módulo $(\psi^*(\sh{M}))\supertilde$ identifica-se com $\widetilde{\sh{M}}\otimes_Y\sh{O}_{Y'}$.
\end{proposition}

\begin{proof}
Observemos, em primeiro lugar, que $\psi^*(\sh{S})$ e $\psi^*(\sh{M})$ são $\sh{O}_{Y'}$-módulos quase-coerentes, assim como seus componentes homogêneos \sref[0]{0.5.1.4}.
Sejam $U$ um aberto afim de $Y$, $U'\subset\psi^{-1}(U)$ um aberto afim de $Y'$, e $A$, $A'$ os anéis de $U$ e $U'$, respectivamente;
tem-se então $\sh{S}|U=\widetilde{S}$, onde $S$ é uma $A$-álgebra graduada, e $\sh{S}'|U'$ identifica-se com $(S\otimes_A A')\supertilde$ \sref[I]{I.1.6.5};
a primeira afirmação resulta então de \sref{II.2.8.10} e \sref[I]{I.3.2.6.2}, pois se verifica imediatamente que a projeção $\Proj(\sh{S}'|U')\to\Proj(\sh{S}|U)$ definida pela identificação precedente é compatível com as operações de restrição em $U$ e $U'$ e define, portanto, um morfismo $\Proj(\sh{S}')\to\Proj(\sh{S})$.
Sejam agora
\begin{align*}
  q &: \Proj(\sh{S})\to Y,
\\q'&: \Proj(\sh{S}')\to Y'
\end{align*}
os morfismos estruturais;
$q'^{-1}(U')$ identifica-se então com $q^{-1}(U)\times_U U'$, e os dois feixes $(\psi^*(\sh{M}))\supertilde|q'^{-1}(U')$ e $(\widetilde{\sh{M}}\otimes_Y\sh{O}_{Y'})|q'^{-1}(U')$ identificam-se canonicamente, ambos, com $(M\otimes_A A')\supertilde$, onde se pôs $M=\Gamma(U,\sh{M})$, em virtude de \sref{II.2.8.10} e \sref[I]{I.1.6.5};
daí resulta a segunda afirmação, pois se verifica ainda, de modo imediato, a compatibilidade das identificações precedentes com as operações de restrição.
\end{proof}

\begin{corollary}[3.5.4]
\label{II.3.5.4}
Com as notações de \sref{II.3.5.3}, $\sh{O}_{X'}(n)$ identifica-se canonicamente com $\sh{O}_X(n)\otimes_Y\sh{O}_{Y'}$ para todo $n\in\bb{Z}$ (com $X=\Proj(\sh{S})$).
\end{corollary}

\begin{proof}
Com efeito, com as notações de \sref{II.3.5.3}, é claro que $\psi^*(\sh{S}(n))=\sh{S}'(n)$ para todo $n\in\bb{Z}$.
\end{proof}

\begin{env}[3.5.5]
\label{II.3.5.5}
Mantendo as notações precedentes, designemos por $\Psi$ a projeção canônica $X'\to X$ e ponhamos $\sh{M}'=\psi^*(\sh{M})$;
suporemos, além disso, que $\sh{S}$ é gerada por $\sh{S}_1$
\oldpage[II]{63}
e que $X$ é de tipo finito sobre $Y$;
daí resulta que $\sh{S}'$ é gerada por $\sh{S}'_1$ (como se vê reduzindo-se ao caso em que $Y$ e $Y'$ são afins) e que $X'$ é de tipo finito sobre $Y'$ \sref[I]{I.6.3.4}.
Seja $\sh{F}$ um $\sh{O}_X$-módulo e ponhamos $\sh{F}'=\Psi^*(\sh{F})$;
resulta então de \sref{II.3.5.4} e \sref[0]{0.4.3.3} que $\sh{F}'(n)=\Psi^*(\sh{F}(n))$ para todo $n\in\bb{Z}$.
Define-se, além disso, um $\psi$-homomorfismo canônico $\theta_n:q_*(\sh{F}(n))\to q'_*(\sh{F}'(n))$ da seguinte maneira: em vista da comutatividade do diagrama
\[
  \xymatrix{
    X \ar[d]_{q}
    & X' \ar[l]_{\Psi} \ar[d]^{q'}
  \\Y
    & Y' \ar[l]^{\psi}
  }
\]
trata-se de definir um homomorfismo $q_*(\sh{F}(n))\to\psi_*(q'_*(\Psi^*(\sh{F}(n))))=q_*(\Psi_*(\Psi^*(\sh{F}(n))))$, e basta tomar o homomorfismo $\theta_n=q_*(\rho_n)$, onde $\rho_n$ é o homomorfismo canônico $\sh{F}(n)\to\Psi_*(\Psi^*(\sh{F}(n)))$ \sref[0]{0.4.4.3}.
É imediato que, para todo aberto afim $U$ de $Y$ e todo aberto afim $U'$ de $Y'$ tal que $U'\subset\psi^{-1}(U)$, o homomorfismo $\theta_n$ induz nas seções o homomorfismo canônico \sref[0]{0.3.7.2} $\Gamma(q^{-1}(U),\sh{F}(n))\to\Gamma(q'^{-1}(U'),\sh{F}'(n))$.
A comutatividade de \sref{II.2.8.13.2} mostra então que, se $\sh{F}$ é quase-coerente, o diagrama
\[
  \xymatrix{
    \sh{F} \ar[r]^{\rho}
    & \sh{F}'
  \\(\bbGamma_*(\sh{F}))\supertilde \ar[u]^{\beta_{\sh{F}}} \ar[r]_{\widetilde{\theta}}
    & (\bbGamma_*(\sh{F}'))\supertilde \ar[u]_{\beta_{\sh{F}'}}
  }
\]
é comutativo (sendo a seta horizontal superior o $\Psi$-morfismo canônico $\sh{F}\to\Psi^*(\sh{F})$).

Do mesmo modo, a comutatividade de \sref{II.2.8.13.1} mostra que o diagrama
\[
  \xymatrix{
    \bbGamma_*(\widetilde{\sh{M}}) \ar[r]^{\theta}
    & \bbGamma_*(\widetilde{\sh{M}'})
  \\ \sh{M} \ar[r]_{\rho} \ar[u]^{\alpha_{\sh{M}}}
    & \sh{M}' \ar[u]_{\alpha_{\sh{M}'}}
  }
\]
é comutativo (sendo a seta horizontal inferior o $\psi$-morfismo canônico $\sh{M}\to\psi^*(\sh{M})$).
\end{env}

\begin{env}[3.5.6]
\label{II.3.5.6}
Consideremos agora dois pré-esquemas $Y$, $Y'$, um morfismo $g:Y'\to Y$, uma $\sh{O}_Y$-álgebra graduada quase-coerente $\sh{S}$, uma $\sh{O}_{Y'}$-álgebra graduada quase-coerente $\sh{S}'$ e um $g$-morfismo de álgebras graduadas $u:\sh{S}\to\sh{S}'$, isto é, um $\sh{O}_Y$-homomorfismo de álgebras graduadas $\sh{S}\to g_*(\sh{S}')$;
sabe-se, aliás, que isso equivale a dar um $\sh{O}_{Y'}$-homomorfismo de álgebras graduadas $u^\sharp:g^*(\sh{S})\to\sh{S}'$.
Deduz-se então canonicamente de $u^\sharp$ um $Y'$-morfismo $W=\Proj(u^\sharp):G(u^\sharp)\to\Proj(g^*(\sh{S}))$, onde $G(u^\sharp)$ é um aberto de $X'=\Proj(\sh{S}')$ \sref{II.3.5.1}.
Por outro lado, $X''=\Proj(g^*(\sh{S}))$ identifica-se
\oldpage[II]{64}
canonicamente com $X\times_Y Y'$, pondo $X=\Proj(\sh{S})$ \sref{II.3.5.3};
compondo com $\Proj(u^\sharp)$ a primeira projeção $p:X\times_Y Y'\to X$, obtém-se, portanto, um morfismo $v:G(u^\sharp)\to X$, que denotaremos por $\Proj(u)$ e tal que o diagrama
\[
  \xymatrix{
    G(u^\sharp) \ar[r]^{v} \ar[d]
    & X \ar[d]
  \\Y' \ar[r]_{g}
    & Y
  }
\]
é comutativo.

Além disso, para todo $\sh{S}$-módulo graduado quase-coerente $\sh{M}$, tem-se um $v$-morfismo canônico
\[
\label{II.3.5.6.1}
  \nu:\widetilde{\sh{M}}\to(g^*(\sh{M})\otimes_{g^*(\sh{S})}\sh{S}')\supertilde|G(u^\sharp).
\tag{3.5.6.1}
\]

Com efeito, $\nu^\sharp$ obtém-se compondo os homomorfismos
\[
  v^*(\widetilde{\sh{M}})
  =W^*(p^*(\widetilde{\sh{M}}))
  \to W^*((g^*(\sh{M}))\supertilde)
  \to (g^*(\sh{M})\otimes_{g^*(\sh{S})}\sh{S}')\supertilde|G(u^\sharp)
\]
onde a primeira seta provém do isomorfismo \sref{II.3.5.3}, e a segunda é o homomorfismo \sref{II.3.5.2}[(ii)];
quando $\sh{S}$ é gerada por $\sh{S}_1$, resulta de \sref{II.3.5.2} que $\nu^\sharp$ é um \emph{isomorfismo}.

Como caso particular de \sref{II.3.5.6.1}, tem-se, para todo $n\in\bb{Z}$, um $v$-morfismo canônico
\[
\label{II.3.5.6.2}
  \nu:\sh{O}_X(n)\to\sh{O}_{X'}(n)|G(u^\sharp).
\tag{3.5.6.2}
\]
\end{env}

\subsection{Subpré-esquemas fechados de um pré-esquema $\Proj(\sh{S})$}
\label{subsection:II.3.6}

\begin{env}[3.6.1]
\label{II.3.6.1}
Sejam $Y$ um pré-esquema e $\vphi:\sh{S}\to\sh{S}'$ um homomorfismo de grau~$0$ de $\sh{O}_Y$-álgebras graduadas quase-coerentes.
Diz-se que $\vphi$ é \emph{(\textbf{TN})-sobrejetivo} (resp. \emph{(\textbf{TN})-injetivo}, \emph{(\textbf{TN})-bijetivo}) se existe um inteiro $n$ tal que, para $k\geq n$, $\vphi_k:\sh{S}_k\to\sh{S}'_k$ seja sobrejetivo (resp. injetivo, bijetivo).
Quando isso ocorre, reduz-se o estudo do morfismo correspondente $\Phi:\Proj(\sh{S}')\to\Proj(\sh{S})$ ao caso em que $\vphi$ é sobrejetivo (resp. injetivo, bijetivo).
Isso se demonstra como em \sref{II.2.9.1} (que é o caso particular correspondente a $Y$ afim), utilizando \sref{II.3.1.8}.
Em vez de dizer que $\vphi$ é (\textbf{TN})-bijetivo, diz-se também que é um \emph{(\textbf{TN})-isomorfismo}.
\end{env}

\begin{proposition}[3.6.2]
\label{II.3.6.2}
Sejam $Y$ um pré-esquema e $\sh{S}$ uma $\sh{O}_Y$-álgebra graduada quase-coerente, e ponhamos $X=\Proj(\sh{S})$.
\begin{enumerate}
  \item[(i)] Se $\vphi:\sh{S}\to\sh{S}'$ é um homomorfismo (\textbf{TN})-sobrejetivo de $\sh{O}_Y$-álgebras graduadas, o morfismo correspondente $\Phi=\Proj(\vphi)$ \sref{II.3.5.1} está definido em todo $\Proj(\sh{S}')$ e é uma imersão fechada de $\Proj(\sh{S}')$ em $X$.
    Se $\sh{J}$ é o núcleo de $\vphi$, o subpré-esquema fechado de $X$ associado a $\Phi$ é definido pelo feixe quase-coerente de ideais $\widetilde{\sh{J}}$ de $\sh{O}_X$.
  \item[(ii)] Suponhamos, além disso, que $\sh{S}_0=\sh{O}_Y$, que $\sh{S}$ seja gerada por $\sh{S}_1$ e que $\sh{S}_1$ seja de tipo finito.
    Seja $X'$ um subpré-esquema fechado de $X=\Proj(\sh{S})$, definido por um feixe quase-coerente de ideais $\sh{I}$ de $\sh{O}_X$.
    Seja $\sh{J}$
\oldpage[II]{65}

o feixe graduado quase-coerente de ideais de $\sh{S}$, imagem inversa de $\bbGamma_*(\sh{I})$ pelo homomorfismo canônico $\alpha:\sh{S}\to\bbGamma_*(\sh{O}_X)$ \sref{II.3.3.2}, e ponhamos $\sh{S}'=\sh{S}/\sh{J}$.
Então $X'$ é o subpré-esquema associado \sref[I]{I.4.2.1} à imersão fechada $\Proj(\sh{S}')\to X$ correspondente ao homomorfismo canônico $\sh{S}\to\sh{S}'$ de $\sh{O}_Y$-álgebras graduadas.
\end{enumerate}
\end{proposition}

\begin{proof}
\medskip\noindent
\begin{enumerate}
  \item[(i)] Podemos supor que $\vphi$ seja sobrejetivo \sref{II.3.6.1}.
    Então, para todo aberto afim $U$ de $Y$, $\Gamma(U,\sh{S})\to\Gamma(U,\sh{S}')$ é sobrejetivo \sref[I]{I.1.3.9}; logo, $G(\vphi)=\Proj(\sh{S}')$ por \sref{II.3.5.1}.
    Reduzimo-nos imediatamente a demonstrar a proposição quando $Y$ é afim, e então ela decorre de \sref{II.2.9.2}[(i)].
  \item[(ii)] Reduzimo-nos a provar que o homomorfismo $\widetilde{\sh{J}}\to\sh{O}_X$ deduzido da injeção canônica $\sh{J}\to\sh{S}$ é um isomorfismo de $\widetilde{\sh{J}}$ sobre $\sh{I}$;
    como a questão é local sobre $Y$, podemos supor que $Y$ seja afim de anel $A$, o que implica $\sh{S}=\widetilde{S}$, onde $S$ é uma $A$-álgebra graduada gerada por $S_1$, sendo $S_1$ de tipo finito sobre $A$.
    Basta então aplicar \sref{II.2.9.2}[(ii)].
\end{enumerate}
\end{proof}

\begin{corollary}[3.6.3]
\label{II.3.6.3}
Sob as condições de \sref{II.3.6.2}[(i)], e supondo, além disso, que $\sh{S}$ seja gerada por $\sh{S}_1$, $\Phi^*(\sh{O}_X(n))$ identifica-se canonicamente com $\sh{O}_{X'}(n)$ para todo $n\in\bb{Z}$.
\end{corollary}

\begin{proof}
Definiu-se um tal isomorfismo canônico quando $Y$ é afim \sref{II.2.9.3};
no caso geral, basta verificar que os isomorfismos assim definidos para cada aberto afim $U$ de $Y$ são compatíveis com a passagem de $U$ a um aberto afim $U'\subset U$, o que é imediato.
\end{proof}

\begin{corollary}[3.6.4]
\label{II.3.6.4}
Sejam $Y$ um pré-esquema, $\sh{S}$ uma $\sh{O}_Y$-álgebra graduada quase-coerente gerada por $\sh{S}_1$, $\sh{M}$ um $\sh{O}_Y$-módulo quase-coerente, $u$ um $\sh{O}_Y$-homomorfismo sobrejetivo $\sh{M}\to\sh{S}_1$, e $\overline{u}:\bb{S}_{\sh{O}_Y}(\sh{M})\to\sh{S}$ o homomorfismo de $\sh{O}_Y$-álgebras graduadas que prolonga $u$ \sref{II.1.7.4}.
Então o morfismo correspondente a $\overline{u}$ é uma imersão fechada de $\Proj(\sh{S})$ em $\Proj(\bb{S}_{\sh{O}_Y}(\sh{M}))$.
\end{corollary}

\begin{proof}
Com efeito, $\overline{u}$ é sobrejetivo por hipótese, e aplica-se \sref{II.3.6.2}[(i)].
\end{proof}

\subsection{Morfismos de um pré-esquema em um espectro homogêneo}
\label{subsection:II.3.7}

\begin{env}[3.7.1]
\label{II.3.7.1}
Sejam $q:X\to Y$ um morfismo de pré-esquemas, $\sh{L}$ um $\sh{O}_X$-módulo invertível e $\sh{S}$ uma $\sh{O}_Y$-álgebra graduada quase-coerente em graus positivos;
então $q^*(\sh{S})$ é uma $\sh{O}_X$-álgebra graduada quase-coerente em graus positivos.
Consideremos a $\sh{O}_X$-álgebra graduada quase-coerente $\sh{S}'=\bigoplus_{n\geq0}\sh{L}^{\otimes n}$ e suponhamos dado um $\sh{O}_X$-homomorfismo de álgebras graduadas
\[
  \psi:q^*(\sh{S})\longrightarrow
  \sh{S}'=\bigoplus_{n\geq0}\sh{L}^{\otimes n},
\]
o que equivale, aliás, a dar um $q$-morfismo de álgebras graduadas
\[
  \psi^\flat:\sh{S}\longrightarrow q_*(\sh{S}').
\]
Sabemos que $\Proj(\sh{S}')$ se identifica canonicamente com $X$ (\sref{II.3.1.7} e \sref{II.3.1.8}[(iii)]);
deduzimos canonicamente de $\psi$ um subconjunto aberto $G(\psi)$ de $X$ e um $Y$-morfismo
\[
\label{II.3.7.1.1}
  r_{\sh{L},\psi}:G(\psi)\longrightarrow \Proj(\sh{S})=P
\tag{3.7.1.1}
\]
\oldpage[II]{66}

que chamaremos de morfismo \emph{associado a $\sh{L}$ e a $\psi$};
recordemos \sref{II.3.5.6} que esse morfismo é obtido compondo o $Y$-morfismo
\[
  \tau=\Proj(\psi):G(\psi)\longrightarrow\Proj(q^*(\sh{S}))
\]
com a primeira projeção $\pi:\Proj(q^*(\sh{S}))=P\times_YX\to P$.
\end{env}

\begin{env}[3.7.2]
\label{II.3.7.2}
Explicitemos $r=r_{\sh{L},\psi}$ quando $Y=\Spec(A)$ é afim e, por conseguinte, $\sh{S}=\widetilde{S}$, onde $S$ é uma $A$-álgebra graduada em graus positivos.
Suponhamos primeiro que $X=\Spec(B)$ seja também afim e que $\sh{L}=\widetilde{L}$, onde $L$ é um $B$-módulo livre de posto~$1$.
Temos então $q^*(\sh{S})=(S\otimes_A B)\supertilde$ \sref[I]{I.1.6.5};
se $c$ é um gerador de $L$, $\psi_n:q^*(\sh{S}_n)\to\sh{L}^{\otimes n}$ corresponde a um homomorfismo $w_n:s\otimes b\mapsto bv_n(s)c^{\otimes n}$ de $S_n\otimes_A B$ em $L^{\otimes n}$, onde $v_n:S_n\to B$ é um homomorfismo de $A$-módulos, e os $v_n$ constituem um homomorfismo de álgebras $S\to B$.
Seja $f\in S_d$ ($d>0$) e ponhamos $g=v_d(f)$;
temos $\pi^{-1}(D_+(f))=D_+(f\otimes1)$ em virtude de \sref{II.2.8.10} e da identificação de $D_+(f)$ com $\Spec(S_{(f)})$ \sref{II.2.3.6};
por outro lado, a fórmula \sref{II.2.8.1.1} mostra (levando em conta a identificação canônica de $X$ com $\Proj(\sh{S}')$) que
\[
  \tau^{-1}(D_+(f\otimes1))=D(g),
\]
de onde
\[
\label{II.3.7.2.1}
  r^{-1}(D_+(f))=D(g).
\tag{3.7.2.1}
\]

Além disso, o morfismo $\tau=\Proj(\psi)$, restrito a $D(g)$, corresponde ao homomorfismo que envia $(s\otimes1)/(f\otimes1)^n$ (para $s\in S_{nd}$) em $v_{nd}(s)/g^n$ \sref{II.2.8.1}, e a projeção $\pi$, restrita a $D_+(f\otimes1)$, corresponde ao homomorfismo $s/f^n\mapsto(s\otimes1)/(f\otimes1)^n$;
concluímos que $r$, restrito a $D(g)$, corresponde ao homomorfismo de $A$-álgebras $\omega:S_{(f)}\to B_g$ tal que $\omega(s/f^n)=v_{nd}(s)/g^n$ para $s\in S_{nd}$ ($n>0$).
Passando ao caso em que $X$ é arbitrário (permanecendo $Y$ afim), obtemos, portanto, levando em conta \sref{II.2.8.1}:
\end{env}

\begin{proposition}[3.7.3]
\label{II.3.7.3}
Se $Y=\Spec(A)$ é afim e $\sh{S}=\widetilde{S}$, onde $S$ é uma $A$-álgebra graduada, então, para todo $f\in S_d=\Gamma(Y,\sh{S}_d)$, tem-se
\[
\label{II.3.7.3.1}
  r_{\sh{L},\psi}^{-1}(D_+(f))=X_{\psi^\flat(f)}
  \qquad\bigl(\text{onde }\psi^\flat(f)\in\Gamma(X,\sh{L}^{\otimes d})\bigr)
\tag{3.7.3.1}
\]
e a restrição $X_{\psi^\flat(f)}\to D_+(f)=\Spec(S_{(f)})$ de $r_{\sh{L},\psi}$ corresponde \sref[I]{I.2.2.4} ao homomorfismo de álgebras
\[
\label{II.3.7.3.2}
  \psi_{(f)}^\flat:S_{(f)}\longrightarrow
  \Gamma(X_{\psi^\flat(f)},\sh{O}_X)
\tag{3.7.3.2}
\]
tal que, para $s\in S_{nd}=\Gamma(Y,\sh{S}_{nd})$,
\[
\label{II.3.7.3.3}
  \psi_{(f)}^\flat(s/f^n)
  =(\psi^\flat(s)|X_{\psi^\flat(f)})
   (\psi^\flat(f)|X_{\psi^\flat(f)})^{-n}.
\tag{3.7.3.3}
\]
\end{proposition}

Diremos que $r_{\sh{L},\psi}$ está \emph{definido em toda parte} se $G(\psi)=X$.
Para que assim seja, é evidentemente necessário e suficiente que $G(\psi)\cap q^{-1}(U)=q^{-1}(U)$ para todo aberto afim $U\subset Y$;
em outras palavras, a questão é \emph{local sobre $Y$}.
Se $Y$ é afim, $G(\psi)$ é a reunião dos $r^{-1}(D_+(f))$, para $f$ homogêneo em $S_+$ \sref{II.2.8.1};
em virtude de \sref{II.3.7.3.1}, os $X_{\psi^\flat(f)}$ devem, portanto, formar um \emph{recobrimento} de $X$;
em outras palavras:

\begin{corollary}[3.7.4]
\label{II.3.7.4}
Sob as hipóteses de \sref{II.3.7.3}, para que $r_{\sh{L},\psi}$ esteja definido em toda parte, é necessário e suficiente que, para todo $x\in X$, existam um inteiro $n>0$ e uma seção $s$ de $\sh{S}_n$ sobre $Y$ tais que, pondo $t=\psi^\flat(s)\in\Gamma(X,\sh{L}^{\otimes n})$, se tenha $t(x)\neq0$.
\end{corollary}
\oldpage[II]{67}

Observe-se que esta condição é sempre satisfeita se $\psi$ é (\textbf{TN})-sobrejetivo.

Do mesmo modo, a questão de saber se $r_{\sh{L},\psi}$ é \emph{dominante} é local sobre $Y$, e tem-se:

\begin{corollary}[3.7.5]
\label{II.3.7.5}
Sob as hipóteses de \sref{II.3.7.3}, para que $r_{\sh{L},\psi}$ seja dominante, é necessário e suficiente que, para todo inteiro $n>0$, toda seção $s\in S_n$ tal que $\psi^\flat(s)\in\Gamma(X,\sh{L}^{\otimes n})$ seja localmente nilpotente seja ela própria nilpotente.
\end{corollary}

\begin{proof}
Devemos, com efeito, exprimir que $r_{\sh{L},\psi}^{-1}(D_+(s))$ só é vazio se $D_+(s)$ for vazio, e o corolário resulta, portanto, de \sref{II.3.7.3.1} e \sref{II.2.3.7}.
\end{proof}

\begin{proposition}[3.7.6]
\label{II.3.7.6}
Sejam $q:X\to Y$ um morfismo, $\sh{L}$ um $\sh{O}_X$-módulo invertível, $\sh{S}$ e $\sh{S}'$ duas $\sh{O}_Y$-álgebras graduadas quase-coerentes, $u:\sh{S}'\to\sh{S}$ um homomorfismo de álgebras graduadas, $\psi:q^*(\sh{S})\to\bigoplus_{n\geq0}\sh{L}^{\otimes n}$ um homomorfismo de álgebras graduadas, e $\psi'=\psi\circ q^*(u)$ o homomorfismo composto.
Se $r_{\sh{L},\psi'}$ está definido em toda parte, o mesmo vale para $r_{\sh{L},\psi}$;
se $u$ é (\textbf{TN})-sobrejetivo e $r_{\sh{L},\psi'}$ é dominante, o mesmo vale para $r_{\sh{L},\psi}$;
inversamente, se $u$ é (\textbf{TN})-injetivo e $r_{\sh{L},\psi}$ é dominante, $r_{\sh{L},\psi'}$ é dominante.
\end{proposition}

\begin{proof}
Temos, com efeito, $G(\psi')\subset G(\psi)$ \sref{II.2.8.4}, de onde resulta a primeira afirmação;
se $u$ é (\textbf{TN})-sobrejetivo, $\Proj(u):\Proj(\sh{S})\to\Proj(\sh{S}')$ está definido em toda parte e é uma imersão fechada;
como $r_{\sh{L},\psi'}$ é a composição de $\Proj(u)$ com a restrição de $r_{\sh{L},\psi}$ a $G(\psi')$, concluímos que, se $r_{\sh{L},\psi'}$ é dominante, o mesmo vale para $r_{\sh{L},\psi}$.
Por fim, se $u$ é (\textbf{TN})-injetivo, sabemos que $\Proj(u)$ é um morfismo dominante de $G(u)$ em $\Proj(\sh{S}')$ \sref{II.2.8.3};
como $G(\psi')$ é a imagem inversa de $G(u)$ por $r_{\sh{L},\psi}$, vemos que, se $r_{\sh{L},\psi}$ é dominante, o mesmo vale para $r_{\sh{L},\psi'}$.
\end{proof}

\begin{proposition}[3.7.7]
\label{II.3.7.7}
Sejam $Y$ um pré-esquema quase-compacto, $q:X\to Y$ um morfismo quase-compacto, $\sh{L}$ um $\sh{O}_X$-módulo invertível e $\sh{S}$ uma $\sh{O}_Y$-álgebra graduada quase-coerente, limite indutivo filtrante de um sistema indutivo $(\sh{S}^\lambda)$ de $\sh{O}_Y$-álgebras graduadas quase-coerentes.
Sejam $\vphi_\lambda:\sh{S}^\lambda\to\sh{S}$ o homomorfismo canônico, $\psi:q^*(\sh{S})\to\bigoplus_{n\geq0}\sh{L}^{\otimes n}$ um homomorfismo de álgebras graduadas, e ponhamos $\psi_\lambda=\psi\circ q^*(\vphi_\lambda)$.
Para que $r_{\sh{L},\psi}$ esteja definido em toda parte, é necessário e suficiente que exista um $\lambda$ tal que $r_{\sh{L},\psi_\lambda}$ esteja definido em toda parte;
então $r_{\sh{L},\psi_\mu}$ está definido em toda parte para $\mu\geq\lambda$.
\end{proposition}

\begin{proof}
A condição é suficiente em virtude de \sref{II.3.7.6}.
Inversamente, suponhamos que $r_{\sh{L},\psi}$ esteja definido em toda parte;
podemos reduzir-nos ao caso em que $Y$ é afim, pois, se para todo aberto afim $U\subset Y$ existe $\lambda(U)$ tal que a restrição de $r_{\sh{L},\psi_{\lambda(U)}}$ a $q^{-1}(U)$ esteja definida em toda parte, bastará (sendo $Y$ quase-compacto) recobrir $Y$ por um número finito de abertos afins $U_i$ e tomar $\lambda\geq\lambda(U_i)$ para todos os índices $i$, em virtude de \sref{II.3.7.6}.
Se $Y$ é afim, a hipótese implica que, para todo $x\in X$, há uma seção $s^{(x)}$ de um $S_n$ tal que, pondo $t^{(x)}=\psi^\flat(s^{(x)})$, se tenha $t^{(x)}(x)\neq0$ (sendo $t^{(x)}$ considerada como seção de $\sh{L}^{\otimes n}$ sobre $X$), o que implica $t^{(x)}(z)\neq0$ para todo $z$ em uma vizinhança $V(x)$ de $x$.
Recubramos $X$ por um número finito de $V(x_i)$ e sejam $s^{(i)}$ as seções correspondentes de $S$;
existe então um $\lambda$ tal que todas as $s^{(i)}$ sejam da forma $\vphi_\lambda({s'_\lambda}^{(i)})$, com ${s'_\lambda}^{(i)}\in S^\lambda$ para todo $i$;
resulta, portanto, de \sref{II.3.7.4} que $r_{\sh{L},\psi_\lambda}$ está definido em toda parte.
A última afirmação é consequência trivial de \sref{II.3.7.6}.
\end{proof}

\begin{corollary}[3.7.8]
\label{II.3.7.8}
Sob as hipóteses de \sref{II.3.7.7}, se os $r_{\sh{L},\psi_\lambda}$ são dominantes, o mesmo vale para $r_{\sh{L},\psi}$;
a recíproca é verdadeira quando os $\vphi_\lambda$ são injetivos.
\end{corollary}
\oldpage[II]{68}

\begin{proof}
A segunda afirmação é um caso particular de \sref{II.3.7.6};
por outro lado, para mostrar que $r_{\sh{L},\psi}$ é dominante, podemos restringir-nos ao caso em que $Y$ é afim;
se $s\in S$ é tal que $\psi^\flat(s)$ seja localmente nilpotente, como podemos escrever $s=\vphi_\lambda(s_\lambda)$ para ao menos um $\lambda$, concluímos da hipótese e de \sref{II.3.7.5} que $s_\lambda$ é nilpotente, logo o mesmo vale para $s$, e aplica-se o critério de \sref{II.3.7.5}.
\end{proof}

\begin{remarks}[3.7.9]
\label{II.3.7.9}
\medskip\noindent
\begin{enumerate}
  \item[(i)] Com as notações de \sref{II.3.7.1}, e levando em conta \sref{II.3.2.10}, temos, para todo $n\in\bb{Z}$, um homomorfismo canônico
    \[
    \label{II.3.7.9.1}
      \theta:r_{\sh{L},\psi}^*(\sh{O}_P(n))\longrightarrow
      \sh{L}^{\otimes n}
    \tag{3.7.9.1}
    \]
    definido de maneira geral em \sref{II.3.5.6.2}.
    Vê-se imediatamente que, sob as condições de \sref{II.3.7.3}, a restrição desse homomorfismo a $X_{\psi^\flat(f)}$ faz corresponder a $s/f^k$ ($s\in S_{n+kd}$) o elemento $(\psi^\flat(s)|X_{\psi^\flat(f)})(\psi^\flat(f)|X_{\psi^\flat(f)})^{-k}$.
  \item[(ii)] Seja $\sh{F}$ um $\sh{O}_X$-módulo quase-coerente, e suponhamos que $q$ seja quase-compacto e separado, de modo que, para todo $n\geq0$, $q_*(\sh{F}\otimes\sh{L}^{\otimes n})$ seja um $\sh{O}_Y$-módulo quase-coerente \sref[I]{I.9.2.2}.
    Seja $\sh{M}'=\bigoplus_{n\geq0}\sh{F}\otimes\sh{L}^{\otimes n}$, que é um $\sh{S}'$-módulo graduado quase-coerente, e consideremos sua imagem direta $\sh{M}=q_*(\sh{M}')=\bigoplus_{n\geq0}q_*(\sh{F}\otimes\sh{L}^{\otimes n})$ (que é um $\sh{S}$-módulo quase-coerente por meio do homomorfismo $\psi^\flat$).
    Veremos que existe um homomorfismo canônico de $\sh{O}_X$-módulos
    \[
    \label{II.3.7.9.2}
      \xi:r_{\sh{L},\psi}^*(\widetilde{\sh{M}})
      \longrightarrow\sh{F}|G(\psi).
    \tag{3.7.9.2}
    \]

    Com efeito, definiu-se em \sref{II.3.5.6.1} um homomorfismo canônico
    \[
    \label{II.3.7.9.3}
      r_{\sh{L},\psi}^*(\widetilde{\sh{M}})
      \longrightarrow
      (q^*(\sh{M})\otimes_{q^*(\sh{S})}\sh{S}')\supertilde|G(\psi),
    \tag{3.7.9.3}
    \]
    onde o segundo membro é considerado como um feixe quase-coerente sobre $\Proj(\sh{S}')$.
    Por outro lado, há um homomorfismo canônico
    \[
    \label{II.3.7.9.4}
      q^*(q_*(\sh{M}'))\otimes_{q^*(\sh{S})}\sh{S}'
      \longrightarrow\sh{M}'
    \tag{3.7.9.4}
    \]
    para todo $\sh{S}'$-módulo graduado quase-coerente $\sh{M}'$:
    para todo aberto $U$ de $X$, toda seção $t'$ de $q^*(q_*(\sh{M}'_h))$ sobre $U$ e toda seção $b'$ de $\sh{S}'_k$ sobre $U$, faz-se corresponder a $t'\otimes b'$ a seção $b'\sigma(t')$ de $\sh{M}'_{h+k}$, onde $\sigma(t')$ é a seção de $\sh{M}'_h$ sobre $U$ que corresponde canonicamente a $t'$ \sref[0]{0.4.4.3}.
    Deduzimos daí um homomorfismo canônico
    \[
    \label{II.3.7.9.5}
      (q^*(q_*(\sh{M}'))\otimes_{q^*(\sh{S})}\sh{S}')\supertilde|G(\psi)
      \longrightarrow\widetilde{\sh{M}'}|G(\psi)
    \tag{3.7.9.5}
    \]
    e, como finalmente $\widetilde{\sh{M}'}$ se identifica canonicamente com $\sh{F}$ \sref{II.3.2.9}[(i)], obtemos precisamente o homomorfismo canônico procurado.

    Sob as condições de \sref{II.3.7.3}, a restrição desse homomorfismo a $X_{\psi^\flat(f)}$ explicita-se da seguinte maneira:
    dada uma seção $t_{nd}$ de $\sh{F}\otimes\sh{L}^{\otimes nd}$ sobre $X$, se $t'_{nd}$ é a seção $t_{nd}$ considerada como seção de $q_*(\sh{F}\otimes\sh{L}^{\otimes nd})$ sobre $Y$, ao elemento $t'_{nd}/f^n$ corresponde a seção $(t_{nd}|X_{\psi^\flat(f)})(\psi^\flat(f)|X_{\psi^\flat(f)})^{-n}$ de $\sh{F}$ sobre $X_{\psi^\flat(f)}$.
\end{enumerate}
\end{remarks}
\oldpage[II]{69}

\subsection{Critérios de imersão em um espectro homogêneo}
\label{subsection:II.3.8}

\begin{env}[3.8.1]
\label{II.3.8.1}
Com as notações de \sref{II.3.7.1}, a questão de saber se $r_{\sh{L},\psi}$ é uma imersão (resp. uma imersão aberta, uma imersão fechada) é evidentemente \emph{local sobre $Y$}.
\end{env}

\begin{proposition}[3.8.2]
\label{II.3.8.2}
Sob as hipóteses de \sref{II.3.7.3}, para que $r_{\sh{L},\psi}$ esteja definido em toda parte e seja uma imersão, é necessário e suficiente que exista uma família de seções $s_\alpha\in S_{n_\alpha}$ ($n_\alpha>0$) tal que, pondo $f_\alpha=\psi^\flat(s_\alpha)$, sejam satisfeitas as condições seguintes:
\begin{enumerate}
  \item[(i)] Os $X_{f_\alpha}$ formam um recobrimento de $X$.
  \item[(ii)] Os $X_{f_\alpha}$ são abertos afins.
  \item[(iii)] Para todo $\alpha$ e todo $t\in\Gamma(X_{f_\alpha},\sh{O}_X)$, existem um inteiro $m>0$ e um $s\in S_{mn_\alpha}$ tais que $t=(\psi^\flat(s)|X_{f_\alpha})(f_\alpha|X_{f_\alpha})^{-m}$.
\end{enumerate}

Para que $r_{\sh{L},\psi}$ esteja definido em toda parte e seja uma imersão aberta, é necessário e suficiente que exista uma família $(s_\alpha)$ que satisfaça as condições (i), (ii), (iii) e:
\begin{enumerate}
  \item[(iv)] Para todo $n>0$ e todo $s\in S_{nn_\alpha}$ tal que $\psi^\flat(s)|X_{f_\alpha}=0$, existe um inteiro $k>0$ tal que $s_\alpha^ks=0$.
\end{enumerate}

Para que $r_{\sh{L},\psi}$ esteja definido em toda parte e seja uma imersão fechada, é necessário e suficiente que exista uma família $(s_\alpha)$ que satisfaça (i), (ii), (iii) e:
\begin{enumerate}
  \item[(v)] Os $D_+(s_\alpha)$ formam um recobrimento de $P=\Proj(S)$.
\end{enumerate}
\end{proposition}

\begin{proof}
Com efeito, para que $r$ seja uma imersão (resp. uma imersão fechada), é necessário e suficiente que exista um recobrimento de $r(G(\psi))$ (resp. de $P$) pelos $D_+(s_\alpha)$ tal que, pondo $V_\alpha=r^{-1}(D_+(s_\alpha))$, a restrição de $r$ a $V_\alpha$ seja uma imersão fechada de $V_\alpha$ em $D_+(s_\alpha)$ \sref[I]{I.4.2.4}.
A condição (i) exprime ao mesmo tempo que $r$ está definido em toda parte e que os $D_+(s_\alpha)$ recobrem $r(X)$, por \sref{II.3.7.3.1};
como $D_+(s_\alpha)$ é afim, as condições (ii) e (iii) exprimem que a restrição de $r$ a $X_{f_\alpha}$ é uma imersão fechada em $D_+(s_\alpha)$, por \sref[I]{I.4.2.3};
por fim, como (iii) e (iv) exprimem que $\psi_{(s_\alpha)}^\flat$ é bijetivo (notação de \sref{II.3.7.3.2}), (ii), (iii) e (iv) exprimem que a restrição de $r$ a $X_{f_\alpha}$ é um isomorfismo sobre $D_+(s_\alpha)$ para todo $\alpha$;
portanto, (i), (ii), (iii) e (iv) exprimem que $r$ é uma imersão aberta.
\end{proof}

\begin{corollary}[3.8.3]
\label{II.3.8.3}
Sob as hipóteses de \sref{II.3.7.6}, se $r_{\sh{L},\psi'}$ está definido em toda parte e é uma imersão, o mesmo vale para $r_{\sh{L},\psi}$.
Se, além disso, $u$ é (\textbf{TN})-sobrejetivo e $r_{\sh{L},\psi'}$ é uma imersão aberta (resp. fechada), o mesmo vale para $r_{\sh{L},\psi}$.
\end{corollary}

\begin{proof}
Com efeito, por \sref{II.3.8.2}, há uma família $s'_\alpha\in S'_{n_\alpha}$ tal que, pondo $f_\alpha=\psi'^\flat(s'_\alpha)$, sejam satisfeitas as condições (i), (ii), (iii).
Ora, se pomos $s_\alpha=u(s'_\alpha)$, temos também $f_\alpha=\psi^\flat(s_\alpha)$, e, se $t=(\psi'^\flat(s')|X_{f_\alpha})(f_\alpha|X_{f_\alpha})^{-m}$, temos também $t=(\psi^\flat(s)|X_{f_\alpha})(f_\alpha|X_{f_\alpha})^{-m}$, pondo $s=u(s')$, de onde resulta a primeira afirmação.
A segunda resulta imediatamente do fato de que $\Proj(u)$ é uma imersão fechada.
\end{proof}

\begin{proposition}[3.8.4]
\label{II.3.8.4}
Suponhamos satisfeitas as hipóteses de \sref{II.3.7.7} e, além disso, que $q:X\to Y$ seja um morfismo de tipo finito.
Então, para que $r_{\sh{L},\psi}$ esteja definido em toda parte e seja uma imersão, é necessário e suficiente que exista $\lambda$ tal que $r_{\sh{L},\psi_\lambda}$ esteja definido em toda parte e seja uma imersão;
nesse caso, $r_{\sh{L},\psi_\mu}$ está definido em toda parte e é uma imersão para $\mu\geq\lambda$.
\end{proposition}
\oldpage[II]{70}
\begin{proof}
Levando em conta \sref{II.3.8.3}, basta demonstrar que, se $r_{\sh{L},\psi}$ está definido em toda parte e é uma imersão, então o mesmo vale para $r_{\sh{L},\psi_\lambda}$ para ao menos um $\lambda$.
Pelo mesmo argumento que em \sref{II.3.7.7}, usando a quase-compacidade de $Y$, reduzimo-nos primeiro ao caso em que $Y=\Spec(A)$ é afim.
Como $X$ é quase-compacto, \sref{II.3.8.2} mostra então a existência de uma família \emph{finita} $(s_i\in S_{n_i})$ de elementos de $S$ que satisfazem as condições~(i), (ii) e (iii).
O morfismo $X_{f_i}\to Y$ (onde $f_i=\psi^\flat(s_i)$) é de tipo finito: com efeito, é um morfismo de esquemas afins, portanto é quase-compacto \sref[I]{I.6.6.1}, e é localmente de tipo finito, pois $q$ é de tipo finito \sref[I]{I.6.3.2}; a conclusão resulta então de \sref[I]{I.6.6.3}.
O anel $B_i$ de $X_{f_i}$ é, por conseguinte, uma $A$-álgebra de tipo finito \sref[I]{I.6.3.3};
seja $(t_{ij})$ uma família de geradores dessa álgebra.
Por hipótese, existem elementos $s'_{ij}\in S_{m_{ij}n_i}$ tais que
\[
  t_{ij} = (\psi^\flat(s'_{ij})|X_{f_i})(\psi^\flat(s_i)|X_{f_i})^{-m_{ij}}.
\]
Ainda por hipótese, existem um $\lambda$ e elementos $s_{i\lambda}\in S_{n_i}^\lambda$ e $s'_{ij\lambda}\in S_{m_{ij}n_i}^\lambda$ cujas imagens por $\vphi_\lambda$ são, respectivamente, $s_i$ e $s'_{ij}$;
é claro que $r_{\sh{L},\psi_\lambda}$ satisfaz as condições~(i), (ii) e (iii) de \sref{II.3.8.2}.
\end{proof}

\begin{proposition}[3.8.5]
\label{II.3.8.5}
Seja $Y$ um esquema quase-compacto, ou um pré-esquema cujo espaço subjacente seja noetheriano, seja $q:X\to Y$ um morfismo \emph{de tipo finito}, seja $\sh{L}$ um $\sh{O}_X$-módulo invertível, seja $\sh{S}$ uma $\sh{O}_Y$-álgebra graduada quase-coerente, e seja $\psi:q^*(\sh{S})\to\bigoplus_{n\geq0}\sh{L}^{\otimes n}$ um homomorfismo de álgebras graduadas.
Para que $r_{\sh{L},\psi}$ esteja definido em toda parte e seja uma imersão, é necessário e suficiente que existam um inteiro $n>0$ e um sub-$\sh{O}_Y$-módulo quase-coerente \emph{de tipo finito} $\sh{E}$ de $\sh{S}_n$ tais que
\begin{enumerate}
  \item[a)] o homomorfismo $\psi_n\circ q^*(j_n):q^*(\sh{E})\to\sh{L}^{\otimes n}$ (onde $j_n:\sh{E}\to\sh{S}_n$ é a injeção canônica) seja sobrejetivo; e
  \item[b)] se denotarmos por $\sh{S}'$ a sub-$\sh{O}_Y$-álgebra (graduada) de $\sh{S}$ gerada por $\sh{E}$, e por $\psi'$ o homomorfismo $\psi\circ q^*(j')$ (onde $j'$ é a injeção de $\sh{S}'$ em $\sh{S}$), então $r_{\sh{L},\psi'}$ esteja definido em toda parte e seja uma imersão.
\end{enumerate}

Quando isso ocorre, todo sub-$\sh{O}_Y$-módulo quase-coerente $\sh{E}'$ de $\sh{S}_n$ que contenha $\sh{E}$ possui a mesma propriedade, e o mesmo vale para a imagem de $\otimes^k\sh{E}$ em $\sh{S}_{kn}$ para todo $k>0$.
\end{proposition}

\begin{proof}
A suficiência da condição e as duas últimas afirmações são casos particulares de \sref{II.3.8.3}, levando em conta o isomorfismo canônico entre $\Proj(\sh{S})$ e $\Proj(\sh{S}^{(k)})$ \sref{II.3.1.8}.

Seja $(U_i)$ uma cobertura aberta finita de $Y$ por abertos afins, e ponhamos $A_i=A(U_i)$.
Como $q^{-1}(U_i)$ é quase-compacto, a hipótese de que $r_{\sh{L},\psi}$ esteja definido em toda parte e seja uma imersão implica, juntamente com \sref{II.3.8.2}, a existência de uma família finita $(s_{ij}\in S_{n_{ij}}^{(i)})$ de elementos de $S^{(i)}=\Gamma(U_i,\sh{S})$ que satisfazem as condições~(i), (ii) e (iii).
Vê-se, como na demonstração de \sref{II.3.8.4}, que o morfismo $X_{f_{ij}}\to U_i$ (onde $f_{ij}=\psi^\flat(s_{ij})$) é de tipo finito, e que o anel $B_{ij}$ de $X_{f_{ij}}$ é, por conseguinte, uma $A_i$-álgebra de tipo finito, com um sistema de geradores da forma $(\psi^\flat(t_{ijk})|X_{f_{ij}})(f_{ij}|X_{f_{ij}})^{-m_{ijk}}$, onde $t_{ijk}\in S_{m_{ijk}n_{ij}}^{(i)}$.
Seja $n$ um múltiplo comum dos $m_{ijk}n_{ij}$;
substituindo (para cada $(i,j,k)$) $s_{ij}$ por uma potência $s_{ij}^\rho$ tal que $\rho n_{ij}=n$, e multiplicando $t_{ijk}$ por $s_{ij}^{\rho-m_{ijk}}$, podemos supor que, para cada $i$, todos os $s_{ij}$ e $t_{ijk}$ pertencem a $S_n^{(i)}$ e que $m_{ijk}=1$.
Seja $E_i$ o sub-$A_i$-módulo de $S_n^{(i)}$ gerado por esses elementos (para $i$ fixado).
Existe um sub-$\sh{O}_Y$-módulo coerente de tipo finito $\sh{E}_i$ de $\sh{S}_n$ tal que $\sh{E}_i|U_i=(E_i)\supertilde$ \sref[I]{I.9.4.7}.
\oldpage[II]{71}
É claro que o sub-$\sh{O}_Y$-módulo $\sh{E}$ de $\sh{S}_n$ dado pela soma dos $\sh{E}_i$ é o objeto procurado (pois cada seção $f_{ij}$ é tal que, para todo $x\in X_{f_{ij}}$, existe uma vizinhança afim $V\subset X_{f_{ij}}$ de $x$ tal que $f_{ij}|V$ seja uma base de $\Gamma(V,\sh{L}^{\otimes n})$).
\end{proof}

\section{Fibrados projetivos; feixes amplos}
\label{section:II.4}


\subsection{Definição dos fibrados projetivos}
\label{subsection:II.4.1}

\begin{definition}[4.1.1]
\label{II.4.1.1}
Sejam $Y$ um pré-esquema, $\sh{E}$ um $\sh{O}_Y$-módulo quase-coerente e $\bb{S}_{\sh{O}_Y}(\sh{E})$ a $\sh{O}_Y$-álgebra simétrica de $\sh{E}$ \sref{II.1.7.4}, que é quase-coerente \sref{II.1.7.7}.
Chama-se \emph{fibrado projetivo sobre $Y$ definido por $\sh{E}$}, e denota-se por $\bb{P}(\sh{E})$, o $Y$-esquema $P=\Proj(\bb{S}_{\sh{O}_Y}(\sh{E}))$.
O $\sh{O}_P$-módulo $\sh{O}_P(1)$ é chamado \emph{feixe fundamental sobre $P$}.
\end{definition}

Quando $Y$ é afim de anel $A$, tem-se $\sh{E}=\widetilde{E}$, onde $E$ é um $A$-módulo, e escreve-se então $\bb{P}(E)$ em vez de $\bb{P}(\widetilde{E})$.

Quando se toma $\sh{E}=\sh{O}_Y^n$, escreve-se $\bb{P}_Y^{n-1}$ em vez de $\bb{P}(\sh{E})$;
se, além disso, $Y$ é afim de anel $A$, escreve-se também $\bb{P}_A^{n-1}$ em vez de $\bb{P}_Y^{n-1}$.
Como $\bb{S}_{\sh{O}_Y}(\sh{O}_Y)$ se identifica canonicamente com $\sh{O}_Y[T]$ \sref{II.1.7.4}, $\bb{P}_Y^0$ se identifica canonicamente com $Y$ \sref{II.3.1.7};
o exemplo~\sref{II.2.4.3} não é senão $\bb{P}_K^1$.

\begin{env}[4.1.2]
\label{II.4.1.2}
Sejam $\sh{E}$ e $\sh{F}$ dois $\sh{O}_Y$-módulos quase-coerentes, e seja $u:\sh{E}\to\sh{F}$ um $\sh{O}_Y$-homomorfismo;
a ele corresponde canonicamente um homomorfismo $\bb{S}(u):\bb{S}_{\sh{O}_Y}(\sh{E})\to\bb{S}_{\sh{O}_Y}(\sh{F})$ de $\sh{O}_Y$-álgebras graduadas \sref{II.1.7.4}.
Se $u$ é \emph{sobrejetivo}, o mesmo vale para $\bb{S}(u)$ e, por conseguinte, \sref{II.3.6.2}[(i)], $\Proj(\bb{S}(u))$ é uma \emph{imersão fechada} $\bb{P}(\sh{F})\to\bb{P}(\sh{E})$, que denotaremos por $\bb{P}(u)$.
Pode-se, portanto, dizer que $\bb{P}(\sh{E})$ é um \emph{funtor contravariante} em $\sh{E}$, com a condição de só se tomarem como morfismos dos $\sh{O}_Y$-módulos quase-coerentes os homomorfismos \emph{sobrejetivos}.

Supondo ainda que $u$ seja sobrejetivo, e pondo $P=\bb{P}(\sh{E})$, $Q=\bb{P}(\sh{F})$ e $j=\bb{P}(u)$, tem-se, a menos de isomorfismo,
\[
\label{II.4.1.2.1}
  j^*(\sh{O}_P(n)) = \sh{O}_Q(n)
  \qquad\mbox{para todo $n\in\bb{Z}$}
  \tag{4.1.2.1}
\]
como resulta de \sref{II.3.6.3}.
\end{env}

\begin{env}[4.1.3]
\label{II.4.1.3}
Seja agora $\psi:Y'\to Y$ um morfismo, e seja $\sh{E}'=\psi^*(\sh{E})$;
tem-se então $\bb{S}_{\sh{O}_{Y'}}(\sh{E}')=\psi^*(\bb{S}_{\sh{O}_Y}(\sh{E}))$ \sref{II.1.7.5};
por conseguinte, \sref{II.3.5.3}, tem-se
\[
\label{II.4.1.3.1}
  \bb{P}(\psi^*(\sh{E})) = \bb{P}(\sh{E})\times_Y Y'
  \tag{4.1.3.1}
\]
a menos de isomorfismo canônico;
além disso, tem-se evidentemente
\[
  \psi^*((\bb{S}_{\sh{O}_Y}(\sh{E}))(n)) = (\bb{S}_{\sh{O}_{Y'}}(\sh{E}'))(n)
\]
para todo $n\in\bb{Z}$; logo, pondo $P=\bb{P}(\sh{E})$ e $P'=\bb{P}(\sh{E}')$, tem-se, por \sref{II.3.5.4}, a menos de isomorfismo,
\[
\label{II.4.1.3.2}
  \sh{O}_{P'}(n) = \sh{O}_P(n)\otimes_Y\sh{O}_{Y'}
  \qquad\mbox{para todo $n\in\bb{Z}$.}
  \tag{4.1.3.2}
\]
\end{env}
\oldpage[II]{72}
\begin{proposition}[4.1.4]
\label{II.4.1.4}
Seja $\sh{L}$ um $\sh{O}_Y$-módulo invertível.
Para todo $\sh{O}_Y$-módulo quase-coerente $\sh{E}$, existe um $Y$-isomorfismo canônico $i_\sh{L}:\bb{P}(\sh{E})\simto\bb{P}(\sh{E}\otimes\sh{L})$;
além disso, se pusermos $P=\bb{P}(\sh{E})$ e $Q=\bb{P}(\sh{E}\otimes\sh{L})$, então $i_\sh{L}^*(\sh{O}_Q(n))$ é canonicamente isomorfo a $\sh{O}_P(n)\otimes_Y\sh{L}^{\otimes n}$ para todo $n\in\bb{Z}$.
\end{proposition}

\begin{proof}
Observemos primeiro que, se $A$ é um anel, $E$ um $A$-módulo e $L$ um $A$-módulo \emph{livre de posto um}, define-se canonicamente um homomorfismo de $A$-módulos
\[
  \bb{S}_n(E\otimes L) \to \bb{S}_n(E)\otimes L^{\otimes n}
\]
fazendo corresponder a $(x_1\otimes y_1)\ldots(x_n\otimes y_n)$ o elemento
\[
  (x_1x_2\ldots x_n)\otimes(y_1\otimes y_2\otimes\ldots\otimes y_n)
  \qquad\mbox{($x_i\in E$, $y_i\in L$, para $1\leq i\leq n$);}
\]
verifica-se imediatamente (reduzindo-se ao caso em que $L=A$) que esse homomorfismo é, de fato, um isomorfismo.
Conclui-se daí um isomorfismo canônico de $A$-álgebras graduadas $\bb{S}_A(E\otimes L)\simto\bigoplus_{n\geq0}\bb{S}_n(E)\otimes L^{\otimes n}$.
Voltando às condições de \sref{II.4.1.4}, as observações precedentes permitem definir um isomorfismo canônico de $\sh{O}_Y$-álgebras graduadas
\[
\label{II.4.1.4.1}
  \bb{S}_{\sh{O}_Y}(\sh{E}\otimes_{\sh{O}_Y}\sh{L}) \simto \bigoplus_{n\geq0}\bb{S}_n(\sh{E})\otimes_{\sh{O}_Y}\sh{L}^{\otimes n}
  \tag{4.1.4.1}
\]
definindo esse isomorfismo como um isomorfismo de pré-feixes e levando em conta \sref{II.1.7.4}, \sref[I]{I.1.3.9} e \sref[I]{I.1.3.12}.
A proposição resulta então de \sref{II.3.1.8}[(iii)] e \sref{II.3.2.10}.
\end{proof}

\begin{env}[4.1.5]
\label{II.4.1.5}
Com as hipóteses de \sref{II.4.1.1}, ponhamos $P=\bb{P}(\sh{E})$ e denotemos por $p$ o morfismo estrutural $P\to Y$.
Como, por definição, $\sh{E}=(\bb{S}_{\sh{O}_Y}(\sh{E}))_1$, tem-se um homomorfismo canônico $\alpha_1:\sh{E}\to p_*(\sh{O}_P(1))$ \sref{II.3.3.2.2} e, por conseguinte, \sref[0]{0.4.4.3}, também um homomorfismo canônico
\[
\label{II.4.1.5.1}
  \alpha_1^\sharp: p^*(\sh{E}) \to \sh{O}_P(1).
  \tag{4.1.5.1}
\]
\end{env}

\begin{proposition}[4.1.6]
\label{II.4.1.6}
O homomorfismo canônico \sref{II.4.1.5.1} é sobrejetivo.
\end{proposition}

\begin{proof}
Vimos em \sref{II.3.3.2} que $\alpha_1^\sharp$ corresponde funtorialmente ao homomorfismo canônico $\sh{E}\otimes_{\sh{O}_Y}\bb{S}_{\sh{O}_Y}(\sh{E})\to(\bb{S}_{\sh{O}_Y}(\sh{E}))(1)$;
como, por definição, $\sh{E}$ gera $\bb{S}_{\sh{O}_Y}(\sh{E})$, esse homomorfismo é sobrejetivo, donde a conclusão em virtude de \sref{II.3.2.4}.
\end{proof}


\subsection{Morfismos de um pré-esquema em um fibrado projetivo}
\label{subsection:II.4.2}

\begin{env}[4.2.1]
\label{II.4.2.1}
Conservando as notações de \sref{II.4.1.5}, sejam $X$ um $Y$-pré-esquema, $q:X\to Y$ o morfismo estrutural e $r:X\to P$ um $Y$-\emph{morfismo}, de modo que o diagrama seguinte seja comutativo:
\[
  \xymatrix{
    P \ar[d]_p & X \ar[l]_r \ar[dl]^q
  \\Y
  }
\]
\oldpage[II]{73}

Como o funtor $r^*$ é exato à direita \sref[0]{0.4.3.1}, deduz-se do homomorfismo de \sref{II.4.1.5.1}, que é sobrejetivo por \sref{II.4.1.6}, um homomorfismo sobrejetivo
\[
  r^*(\alpha_1^\sharp): r^*(p^*(\sh{E})) \to r^*(\sh{O}_P(1)).
\]

Mas $r^*(p^*(\sh{E}))=q^*(\sh{E})$, e $r^*(\sh{O}_P(1))$ é localmente isomorfo a $r^*(\sh{O}_P)=\sh{O}_X$; em outras palavras, é um $\sh{O}_X$-módulo \emph{invertível} $\sh{L}_r$.
Assim, a partir de $r$, definiu-se um $\sh{O}_X$-homomorfismo canônico sobrejetivo
\[
\label{II.4.2.1.1}
  \vphi_r:q^*(\sh{E}) \to \sh{L}_r.
  \tag{4.2.1.1}
\]

Quando $Y=\Spec(A)$ é afim e $\sh{E}=\widetilde{E}$, esse homomorfismo pode ser descrito ainda da maneira seguinte.
Dado $f\in E$, resulta primeiro de \sref{II.2.6.3} que
\[
\label{II.4.2.1.2}
  r^{-1}(D_+(f)) = X_{\vphi_r^\flat(f)}.
  \tag{4.2.1.2}
\]

Por outro lado, seja $V$ um aberto afim de $X$ contido em $r^{-1}(D_+(f))$, e seja $B$ seu anel, que é uma $A$-álgebra;
ponhamos $S=\bb{S}_A(E)$;
a restrição de $r$ a $V$ corresponde a um $A$-homomorfismo $\omega:S_{(f)}\to B$, e tem-se $q^*(\sh{E})|V=(E\otimes_A B)\supertilde$ e $\sh{L}_r|V=\widetilde{L_r}$, onde $L_r=(S(1))_{(f)}\otimes_{S_{(f)}}B_{[\omega]}$ \sref[I]{I.1.6.5}.
A restrição de $\vphi_r$ a $q^*(\sh{E})|V$ corresponde ao $B$-homomorfismo $u:E\otimes_A B\to L_r$ que transforma $x\otimes1$ em $(x/1)\otimes1=(f/1)\otimes\omega(x/f)$.
A extensão canônica de $\vphi_r$ a um homomorfismo de $\sh{O}_X$-álgebras
\[
  \psi_r: q^*(\bb{S}(\sh{E})) = \bb{S}(q^*(\sh{E})) \to \bb{S}(\sh{L}_r) = \bigoplus_{n\geq0}\sh{L}_r^{\otimes n}
\]
é, portanto, tal que sua restrição a $q^*(\bb{S}_n(\sh{E}))|V$ corresponde ao homomorfismo $\bb{S}_n(E\otimes_A B)=\bb{S}_n(E)\otimes_A B\to L_r^{\otimes n}$ que transforma $s\otimes1$ em $(f/1)^{\otimes n}\otimes\omega(s/f^n)$.
\end{env}

\begin{env}[4.2.2]
\label{II.4.2.2}
Reciprocamente, suponhamos dados um morfismo $q:X\to Y$, um $\sh{O}_X$-módulo invertível $\sh{L}$ e um $\sh{O}_Y$-módulo quase-coerente $\sh{E}$;
a um homomorfismo $\vphi:q^*(\sh{E})\to\sh{L}$ corresponde canonicamente um homomorfismo de $\sh{O}_X$-álgebras quase-coerentes
\[
  \psi: \bb{S}(q^*(\sh{E})) = q^*(\bb{S}(\sh{E})) \to \bigoplus_{n\geq0}\sh{L}^{\otimes n},
\]
e, por conseguinte, \sref{II.3.7.1}, um $Y$-morfismo $r_{\sh{L},\psi}:G(\psi)\to\Proj(\bb{S}(\sh{E}))=\bb{P}(\sh{E})$, que denotaremos por $r_{\sh{L},\vphi}$.
Se $\vphi$ é \emph{sobrejetivo}, o mesmo vale para $\psi$ e, portanto, \sref{II.3.7.4}, $r_{\sh{L},\vphi}$ está \emph{definido em toda parte}.
Além disso, com as notações de \sref{II.4.2.1} e \sref{II.4.2.2}:
\end{env}

\begin{proposition}[4.2.3]
\label{II.4.2.3}
Dados um morfismo $q:X\to Y$ e um $\sh{O}_Y$-módulo quase-coerente $\sh{E}$, as aplicações $r\mapsto(\sh{L}_r,\vphi_r)$ e $(\sh{L},\vphi)\mapsto r_{\sh{L},\vphi}$ põem em correspondência biunívoca o conjunto dos $Y$-morfismos $r:X\to\bb{P}(\sh{E})$ e o conjunto das classes de equivalência dos pares $(\sh{L},\vphi)$ formados por um $\sh{O}_X$-módulo invertível $\sh{L}$ e um homomorfismo sobrejetivo $\vphi:q^*(\sh{E})\to\sh{L}$, sendo dois pares $(\sh{L},\vphi)$ e $(\sh{L}',\vphi')$ equivalentes se existir um $\sh{O}_X$-isomorfismo $\tau:\sh{L}\simto\sh{L}'$ tal que $\vphi'=\tau\circ\vphi$.
\end{proposition}

\begin{proof}
Partamos primeiro de um $Y$-morfismo $r:X\to\bb{P}(\sh{E})$, formemos $\sh{L}_r$ e $\vphi_r$ \sref{II.4.2.1}, e depois $r'=r_{\sh{L}_r,\vphi_r}$;
resulta imediatamente de \sref{II.4.2.1} e \sref{II.3.7.2} que os morfismos $r$ e $r'$ são idênticos (tomando como gerador de $\sh{L}_r$ o elemento $(f/1)\otimes1$ para definir os homomorfismos $v_n$ de \sref{II.3.7.2}).
Reciprocamente, partamos de um par $(\sh{L},\vphi)$ e formemos
\oldpage[II]{74}
$r''=r_{\sh{L},\vphi}$, e depois $\sh{L}_{r''}$ e $\vphi_{r''}$;
mostremos que existe um isomorfismo canônico $\tau:\sh{L}_{r''}\simto\sh{L}$ tal que $\vphi=\tau\circ\vphi_{r''}$;
para defini-lo, podemos restringir-nos ao caso em que $Y=\Spec(A)$ e $X=\Spec(B)$ sejam afins e, com as notações de \sref{II.4.2.1} e \sref{II.3.7.2}, fazer corresponder a cada elemento $(x/1)\otimes1$ de $L_{r''}$ (com $x\in E$) o elemento $v_1(x)c$ de $L$.
Verifica-se imediatamente que $\tau$ não depende do gerador escolhido $c$ de $L$;
como $v_1$ é sobrejetivo por hipótese, para provar que $\tau$ é um isomorfismo basta ver que, se $x/1=0$ em $(S(1))_{(f)}$, então $v_1(x)/1=0$ em $B_g$;
mas a primeira relação significa que $f^nx=0$ em $\bb{S}_{n+1}(E)$ para algum $n$, e daí se deduz que $v_{n+1}(f^nx)=g^nv_1(x)=0$ em $B$, donde a conclusão.
Enfim, é imediato que, para dois pares equivalentes $(\sh{L},\vphi)$ e $(\sh{L}',\vphi')$, tem-se $r_{\sh{L},\vphi}=r_{\sh{L}',\vphi'}$.
\end{proof}

Em particular, para $X=Y$:
\begin{theorem}[4.2.4]
\label{II.4.2.4}
O conjunto das seções sobre $Y$ de $\bb{P}(\sh{E})$ está em correspondência biunívoca canônica com o conjunto dos sub-$\sh{O}_Y$-módulos quase-coerentes $\sh{F}$ de $\sh{E}$ tais que $\sh{E}/\sh{F}$ seja invertível.
\end{theorem}

Observe-se que essa propriedade corresponde à definição clássica do ``espaço projetivo'' como conjunto dos hiperplanos de um espaço vetorial (o caso clássico corresponde a $Y=\Spec(K)$, onde $K$ é um corpo, e $\sh{E}=\widetilde{E}$, sendo $E$ um $K$-espaço vetorial de dimensão finita; os $\sh{F}$ que possuem a propriedade enunciada em \sref{II.4.2.4} correspondem então aos hiperplanos de $E$, e sabe-se, por outro lado, que as seções sobre $Y$ de $\bb{P}(\sh{E})$ são então os \emph{pontos $K$-racionais de $\bb{P}(\sh{E})$} \sref[I]{I.3.4.5}).

\begin{remark}[4.2.5]
\label{II.4.2.5}
Como há uma correspondência biunívoca canônica entre os $Y$-morfismos de $X$ em $P$ e seus morfismos gráficos, que são seções sobre $X$ de $P\times_Y X$ \sref[I]{I.3.3.14}, vê-se que, reciprocamente, \sref{II.4.2.3} pode ser deduzida de \sref{II.4.2.4}.
Denotemos por $\Hyp_Y(X,\sh{E})$ o conjunto dos sub-$\sh{O}_X$-módulos quase-coerentes $\sh{F}$ de $\sh{E}\otimes_Y\sh{O}_X=q^*(\sh{E})$ tais que $q^*(\sh{E})/\sh{F}$ seja um $\sh{O}_X$-módulo invertível.
Se $g:X'\to X$ é um $Y$-morfismo, resulta da exatidão à direita de $g^*$ que
\[
  g^*(q^*(\sh{E})/\sh{F})=g^*q^*(\sh{E})/g^*(\sh{F}),
\]
de modo que este último feixe é invertível. Por conseguinte, $X\mapsto\Hyp_Y(X,\sh{E})$ é um \emph{funtor contravariante} na categoria dos $Y$-pré-esquemas.
Pode-se então interpretar o teorema \sref{II.4.2.4} como definindo um \emph{isomorfismo canônico} entre os funtores $X\mapsto\Hom_Y(X,\bb{P}(\sh{E}))$ e $X\mapsto\Hyp_Y(X,\sh{E})$, ambos contravariantes em $X$ variável na categoria dos $Y$-pré-esquemas.
Isso fornece também uma caracterização do fibrado projetivo $P=\bb{P}(\sh{E})$ pela seguinte \emph{propriedade universal}, mais próxima da intuição geométrica que as construções dos §§\hyperref[section:II.2]{2}--\hyperref[section:II.3]{3}:
para todo morfismo $q:X\to Y$ e todo $\sh{O}_X$-módulo invertível $\sh{L}$ que seja um quociente de $\sh{E}\otimes_{\sh{O}_Y}\sh{O}_X$, existe um único $Y$-morfismo $r:X\to P$ tal que $\sh{L}=r^*(\sh{O}_P(1))$.

Veremos mais adiante como se podem definir, da mesma maneira, entre outros, os esquemas ``grassmannianos''.
\end{remark}

\begin{corollary}[4.2.6]
\label{II.4.2.6}
Suponhamos que todo $\sh{O}_Y$-módulo invertível seja trivial \sref[I]{I.2.4.8}.
Seja $V$ o grupo $\Hom_{\sh{O}_Y}(\sh{E},\sh{O}_Y)$, considerado como módulo sobre o anel $A=\Gamma(Y,\sh{O}_Y)$, e seja $V^*$ o subconjunto de $V$ formado pelos homomorfismos sobrejetivos.
Então o conjunto das seções sobre $Y$ de $\bb{P}(\sh{E})$ identifica-se canonicamente com $V^*/A^*$, onde $A^*$ é o grupo das unidades de $A$.
\end{corollary}
\oldpage[II]{75}
Em particular:
\begin{enumerate}
  \item O corolário \sref{II.4.2.6} aplica-se sempre que $Y$ é um \emph{esquema local} \sref[I]{I.2.4.8}.
    Sejam $Y$ um pré-esquema qualquer, $y$ um ponto de $Y$ e $Y'=\Spec(\kres(y))$;
    a fibra $p^{-1}(y)$ de $\bb{P}(\sh{E})$ identifica-se, em virtude de \sref{II.4.1.3.1}, com $\bb{P}(\sh{E}^y)$, onde $\sh{E}^y=\sh{E}_y\otimes_{\sh{O}_y}\kres(y)=\sh{E}_y/\mathfrak{m}_y\sh{E}_y$ é considerado como espaço vetorial sobre $\kres(y)$.
    Mais geralmente, se $K$ é uma extensão de $\kres(y)$, $p^{-1}(y)\otimes_{\kres(y)}K$ identifica-se com $\bb{P}(\sh{E}^y\otimes_{\kres(y)}K)$.
    O corolário \sref{II.4.2.6} mostra, portanto, que o conjunto dos \emph{pontos geométricos de $\bb{P}(\sh{E})$ a valores na extensão $K$ de $\kres(y)$} \sref[I]{I.3.4.5}, que se pode ainda chamar de \emph{fibra geométrica racional sobre $K$ de $\bb{P}(\sh{E})$ acima do ponto $y$}, identifica-se com o \emph{espaço projetivo} associado ao \emph{dual} do $K$-espaço vetorial $\sh{E}^y\otimes_{\kres(y)}K$.
  \item Suponhamos que $Y$ seja afim de anel $A$ e, além disso, que todo $\sh{O}_Y$-módulo invertível seja trivial;
    tomemos ainda $\sh{E}=\sh{O}_Y^n$;
    então, em \sref{II.4.2.6}, $V$ identifica-se com $A^n$ \sref[I]{I.1.3.8}, e $V^*$, com o conjunto dos sistemas $(f_i)_{1\leq i\leq n}$ de elementos de $A$ que geram o ideal $A$;
    dois desses sistemas definem a mesma $Y$-seção de $\bb{P}_Y^{n-1}=\bb{P}_A^{n-1}$, em outras palavras, \emph{o mesmo ponto de $\bb{P}_A^{n-1}$ a valores em $A$}, se, e somente se, um se deduz do outro por multiplicação por um elemento invertível de $A$.
\end{enumerate}

Essas propriedades justificam a terminologia de «fibrado projetivo» para $\bb{P}(\sh{E})$.
Observaremos que a definição assim obtida de «espaço projetivo» é, de fato, \emph{dual} da definição clássica;
isso nos é imposto pela necessidade de poder definir $\bb{P}(\sh{E})$ para um $\sh{O}_Y$-módulo quase-coerente \emph{qualquer} $\sh{E}$, não necessariamente localmente livre.

\begin{remark}[4.2.7]
\label{II.4.2.7}
Veremos, no Capítulo~V, que, se $Y$ é não vazio, localmente noetheriano e conexo, e se $\sh{E}$ é localmente livre de posto pelo menos dois, então, pondo $P=\bb{P}(\sh{E})$, todo $\sh{O}_P$-módulo invertível $\sh{L}$ é isomorfo a um $\sh{O}_P$-módulo da forma $\sh{L}'\otimes_{\sh{O}_Y}\sh{O}_P(m)$, onde $\sh{L}'$ é um $\sh{O}_Y$-módulo invertível, bem determinado a menos de isomorfismo, e $m$ é um inteiro bem determinado.
Em outras palavras, $\HH^1(P,\sh{O}_P^*)$ é isomorfo a $\bb{Z}\times\HH^1(Y,\sh{O}_Y^*)$ \sref[0]{0.5.4.7}.
Veremos também (\sref[III]{III.2.1.15}, levando em conta \sref[0]{0.5.4.10}) que $p_*(\sh{L})=0$ se $m<0$, e que $p_*(\sh{L})$ é isomorfo a $\sh{L}'\otimes_{\sh{O}_Y}(\bb{S}_{\sh{O}_Y}(\sh{E}))_m$ se $m\geq0$.
Se $\sh{F}$ é um $\sh{O}_Y$-módulo quase-coerente, todo $Y$-morfismo $\bb{P}(\sh{E})\to\bb{P}(\sh{F})$ é, portanto, determinado pela escolha de um $\sh{O}_Y$-módulo invertível $\sh{L}'$, de um inteiro $m\geq0$ e de um $\sh{O}_Y$-homomorfismo $\psi:\sh{F}\to\sh{L}'\otimes_{\sh{O}_Y}(\bb{S}_{\sh{O}_Y}(\sh{E}))_m$ tal que o homomorfismo correspondente $\psi^\sharp$ de $\sh{O}_{\bb{P}(\sh{E})}$-módulos seja sobrejetivo.
Veremos também que, se o $Y$-morfismo considerado é um isomorfismo, então $m=1$ e $\sh{F}$ é isomorfo a $\sh{E}\otimes_{\sh{O}_Y}\sh{L}'$ (recíproca de \sref{II.4.1.4}).
Isso permitirá determinar o feixe dos germes de automorfismos de $\bb{P}(\sh{E})$ como quociente do feixe de grupos $\shAut(\sh{E})$ (localmente isomorfo a $\GL(n,\sh{O}_Y)$ se $\sh{E}$ é de posto $n$) por $\sh{O}_Y^*$.
\end{remark}

\begin{env}[4.2.8]
\label{II.4.2.8}
Conservando a notação de \sref{II.4.2.1}, seja $u:X'\to X$ um morfismo;
se o $Y$-morfismo $r:X\to P$ corresponde ao homomorfismo $\vphi:q^*(\sh{E})\to\sh{L}$, o $Y$-morfismo $r\circ u$ corresponde a $u^*(\vphi):u^*(q^*(\sh{E}))\to u^*(\sh{L})$, como resulta imediatamente das definições.
\end{env}

\begin{env}[4.2.9]
\label{II.4.2.9}
Sejam $\sh{E}$ e $\sh{F}$ dois $\sh{O}_Y$-módulos quase-coerentes, $v:\sh{E}\to\sh{F}$ um homomorfismo sobrejetivo e $j=\bb{P}(v)$ a imersão fechada correspondente $\bb{P}(\sh{F})\to\bb{P}(\sh{E})$ \sref{II.4.1.2}.
Se o $Y$-morfismo $r:X\to\bb{P}(\sh{F})$ corresponde ao homomorfismo $\vphi:q^*(\sh{F})\to\sh{L}$, então o
\oldpage[II]{76}
$Y$-morfismo $j\circ r$ corresponde a $q^*(\sh{E})\xrightarrow{q^*(v)}q^*(\sh{F})\xrightarrow{\vphi}\sh{L}$;
isso resulta ainda da definição dada em \sref{II.4.2.1}.
\end{env}

\begin{env}[4.2.10]
\label{II.4.2.10}
Seja $\psi:Y'\to Y$ um morfismo, e ponhamos $\sh{E}'=\psi^*(\sh{E})$.
Se o $Y$-morfismo $r:X\to P$ corresponde ao homomorfismo $\vphi:q^*(\sh{E})\to\sh{L}$, então o $Y'$-morfismo
\[
  r_{(Y')}:X_{(Y')}\to P'=\bb{P}(\sh{E}')
\]
corresponde a $\vphi_{(Y')}:q_{(Y')}^*(\sh{E}')=q^*(\sh{E})\otimes_Y\sh{O}_{Y'}\to\sh{L}\otimes_Y\sh{O}_{Y'}$.
Com efeito, em virtude de \sref{II.4.1.3.1}, temos o diagrama comutativo
\[
  \xymatrix{
    Y' \ar[d]
    & P'=P_{(Y')} \ar[l]_{p_{(Y')}} \ar[d]^{u}
    & X_{(Y')} \ar[l]_{r_{(Y')}} \ar[d]^{v}
  \\Y
    & P \ar[l]_{p}
    & X \ar[l]_{r}
  }
\]

Por \sref{II.4.1.3.2}, temos
\[
  (r_{(Y')})^*(\sh{O}_{P'}(1))=(r_{(Y')})^*(u^*(\sh{O}_P(1)))=v^*(r^*(\sh{O}_P(1)))=v^*(\sh{L})=\sh{L}\otimes_Y\sh{O}_{Y'};
\]
por outro lado, $u^*(\alpha_1^\sharp)$ não é senão o homomorfismo canônico $\alpha_1^\sharp:(p_{(Y')})^*(\sh{E}')\to\sh{O}_{P'}(1)$;
isso se vê explicitando os homomorfismos canônicos $\alpha_1^\sharp$ em $P$ e $P'$ como em \sref{II.4.1.6}.
Daí nossa afirmação.
\end{env}

\subsection{O morfismo de Segre}
\label{subsection:II.4.3}

\begin{env}[4.3.1]
\label{II.4.3.1}
Sejam $Y$ um pré-esquema, $\sh{E}$ e $\sh{F}$ dois $\sh{O}_Y$-módulos quase-coerentes;
ponhamos $P_1=\bb{P}(\sh{E})$ e $P_2=\bb{P}(\sh{F})$, e designemos por $p_1:P_1\to Y$ e $p_2:P_2\to Y$ os morfismos estruturais.
Seja $Q=P_1\times_Y P_2$, e sejam $q_1:Q\to P_1$ e $q_2:Q\to P_2$ as projeções canônicas;
o $\sh{O}_Q$-módulo $\sh{L}=\sh{O}_{P_1}(1)\otimes_Y\sh{O}_{P_2}(1)=q_1^*(\sh{O}_{P_1}(1))\otimes_{\sh{O}_Q}q_2^*(\sh{O}_{P_2}(1))$ é invertível como produto tensorial de dois $\sh{O}_Q$-módulos invertíveis \sref[0]{0.5.4.4}.
Por outro lado, se $r=p_1\circ q_1=p_2\circ q_2$ é o morfismo estrutural $Q\to Y$, temos $r^*(\sh{E}\otimes_{\sh{O}_Y}\sh{F})=q_1^*(p_1^*(\sh{E}))\otimes_{\sh{O}_Q}q_2^*(p_2^*(\sh{F}))$ \sref[0]{0.4.3.3};
os homomorfismos canônicos sobrejetivos \sref{II.4.1.5.1}, $p_1^*(\sh{E})\to\sh{O}_{P_1}(1)$ e $p_2^*(\sh{F})\to\sh{O}_{P_2}(1)$, dão, portanto, por produto tensorial, um homomorfismo canônico
\[
\label{II.4.3.1.1}
  s:r^*(\sh{E}\otimes_{\sh{O}_Y}\sh{F})\to\sh{L}
  \tag{4.3.1.1}
\]
evidentemente sobrejetivo;
dele deduzimos, portanto, por \sref{II.4.2.2}, um morfismo canônico, chamado \emph{morfismo de Segre}:
\[
\label{II.4.3.1.2}
  \varsigma:\bb{P}(\sh{E})\times_Y\bb{P}(\sh{F})\to\bb{P}(\sh{E}\otimes_{\sh{O}_Y}\sh{F}).
  \tag{4.3.1.2}
\]

Explicitemos o morfismo $\varsigma$ quando $Y=\Spec(A)$ é afim, $\sh{E}=\widetilde{E}$ e $\sh{F}=\widetilde{F}$, sendo $E$ e $F$ dois $A$-módulos, donde $\sh{E}\otimes_{\sh{O}_Y}\sh{F}=(E\otimes_A F)^\sim$ \sref[I]{I.1.3.12};
ponhamos $R=\bb{S}_A(E)$, $S=\bb{S}_A(F)$ e $T=\bb{S}_A(E\otimes_A F)$;
sejam $f\in E$, $g\in F$, e consideremos o aberto afim
\[
  D_+(f)\times_YD_+(g)=\Spec(B)
\]
\oldpage[II]{77}
de $Q$, onde $B=R_{(f)}\otimes_A S_{(g)}$;
a restrição de $\sh{L}$ a esse aberto afim é $\widetilde{L}$, onde
\[
  L=(R(1))_{(f)}\otimes_A(S(1))_{(g)},
\]
e o elemento $c=(f/1)\otimes(g/1)$ é um gerador de $L$ considerado como $B$-módulo livre \sref{II.2.5.7}.
O homomorfismo \sref{II.4.3.1.1} corresponde ao homomorfismo
\[
  (x\otimes y)\otimes b\longmapsto b((x/1)\otimes(y/1))
\]
de $(E\otimes_A F)\otimes_A B$ em $L$.
Com a notação de \sref{II.3.7.2}, temos, portanto, $v_1(x\otimes y)=(x/f)\otimes(y/g)$;
a restrição de $\varsigma$ a $D_+(f)\times_YD_+(g)$ é um morfismo desse esquema afim em $D_+(f\otimes g)$, correspondente ao homomorfismo de anéis $\omega:T_{(f\otimes g)}\to R_{(f)}\otimes_A S_{(g)}$ tal que
\[
\label{II.4.3.1.3}
  \omega((x\otimes y)/(f\otimes g))=(x/f)\otimes(y/g)
  \tag{4.3.1.3}
\]
para $x\in E$ e $y\in F$.
\end{env}

\begin{env}[4.3.2]
\label{II.4.3.2}
Resulta de \sref{II.4.2.3} que temos um isomorfismo canônico
\[
\label{II.4.3.2.1}
  \tau:\varsigma^*(\sh{O}_P(1))\simto\sh{O}_{P_1}(1)\otimes_Y\sh{O}_{P_2}(1)
  \tag{4.3.2.1}
\]
onde pusemos $P=\bb{P}(\sh{E}\otimes_{\sh{O}_Y}\sh{F})$.
Mostremos que, para $x\in\Gamma(Y,\sh{E})$ e $y\in\Gamma(Y,\sh{F})$, temos
\[
\label{II.4.3.2.2}
  \tau(\alpha_1(x\otimes y))=\alpha_1(x)\otimes\alpha_1(y).
  \tag{4.3.2.2}
\]

Com efeito, somos reduzidos ao caso em que $Y$ é afim, e temos então, com a notação de \sref{II.4.3.1} e de \sref{II.2.6.2}, $\alpha_1^{f\otimes g}(x\otimes y)=(x\otimes y)/1$ em $(T(1))_{(f\otimes g)}$, $\alpha_1^f(x)=x/1$ em $(R(1))_{(f)}$ e $\alpha_1^g(y)=y/1$ em $(S(1))_{(g)}$.
A definição de $\tau$ dada em \sref{II.4.2.3} e o cálculo de $v_1$ feito em \sref{II.4.3.1} provam então imediatamente \sref{II.4.3.2.2}.
Daí deduzimos a fórmula
\[
\label{II.4.3.2.3}
  \varsigma^{-1}(P_{x\otimes y})=(P_1)_x\times_Y(P_2)_y
  \tag{4.3.2.3}
\]
com a notação de \sref{II.3.1.4}.
Com efeito, levando em conta \sref{II.3.3.3}, a fórmula \sref{II.4.3.2.2} reduz-se (voltando ao caso afim, com o auxílio de \sref[I]{I.3.2.7} e \sref[I]{I.3.2.3}) a provar o lema seguinte:

\begin{lemma}[4.3.2.4]
\label{II.4.3.2.4}
Sejam $B$ e $B'$ duas $A$-álgebras, e sejam $Y=\Spec(A)$, $Z=\Spec(B)$ e $Z'=\Spec(B')$;
quaisquer que sejam $t\in B$ e $t'\in B'$, temos $D(t\otimes t')=D(t)\times_YD(t')$.
\end{lemma}

\begin{proof}
Com efeito, se $p:Z\times_YZ'\to Z$ e $p':Z\times_YZ'\to Z'$ são as projeções canônicas, resulta de \sref[I]{I.1.2.2.2} que $p^{-1}(D(t))=D(t\otimes1)$ e ${p'}^{-1}(D(t'))=D(1\otimes t')$;
a conclusão resulta de \sref[I]{I.3.2.7} e \sref[I]{I.1.1.9.1}, pois $(t\otimes1)(1\otimes t')=t\otimes t'$.
\end{proof}
\end{env}

\begin{proposition}[4.3.3]
\label{II.4.3.3}
O morfismo de Segre é uma imersão fechada.
\end{proposition}

\begin{proof}
Sendo a questão local sobre $Y$, somos reduzidos ao caso em que $Y$ é afim.
Com a notação de \sref{II.4.3.1} e \sref{II.4.3.2}, os $D_+(f\otimes g)$ formam então uma base da topologia de $P$, pois os $f\otimes g$ geram $T$ quando $f$ percorre $E$ e $g$ percorre $F$.
Por outro lado, temos, em virtude de \sref{II.4.3.2.3}, $\varsigma^{-1}(D_+(f\otimes g))=D_+(f)\times_YD_+(g)$.
Basta, portanto, por \sref[I]{I.4.2.4}, provar que a restrição de $\varsigma$ a $D_+(f)\times_YD_+(g)$ é uma imersão fechada em $D_+(f\otimes g)$.
Ora, com a mesma notação, a fórmula \sref{II.4.3.1.3} mostra que $\omega$ é \emph{sobrejetivo}, o que termina a demonstração.
\end{proof}

\begin{env}[4.3.4]
\label{II.4.3.4}
O morfismo de Segre é \emph{funtorial} em $\sh{E}$ e $\sh{F}$, quando se restringem os
\oldpage[II]{78}
homomorfismos de $\sh{O}_Y$-módulos quase-coerentes aos homomorfismos \emph{sobrejetivos}.
Com efeito, é preciso mostrar que, se $\sh{E}\to\sh{E}'$ é um $\sh{O}_Y$-homomorfismo sobrejetivo, o diagrama
\[
  \xymatrix{
    \bb{P}(\sh{E}')\times_Y\bb{P}(\sh{F}) \ar[r]^{j\times1} \ar[d]_{\varsigma}
    & \bb{P}(\sh{E})\times_Y\bb{P}(\sh{F}) \ar[d]^{\varsigma}
  \\ \bb{P}(\sh{E}'\otimes_{\sh{O}_Y}\sh{F}) \ar[r]
    & \bb{P}(\sh{E}\otimes_{\sh{O}_Y}\sh{F})
  }
\]
é comutativo, denotando $j$ a imersão fechada canônica $\bb{P}(\sh{E}')\to\bb{P}(\sh{E})$.
Ponhamos $P'_1=\bb{P}(\sh{E}')$ e conservemos, de resto, a notação de \sref{II.4.3.1};
$j\times1$ é uma imersão fechada \sref[I]{I.4.3.1} e temos, a menos de isomorfismos,
\[
  (j\times1)^*(\sh{O}_{P_1}(1)\otimes\sh{O}_{P_2}(1))
  =j^*(\sh{O}_{P_1}(1))\otimes\sh{O}_{P_2}(1)
  =\sh{O}_{P'_1}(1)\otimes\sh{O}_{P_2}(1)
\]
em virtude de \sref{II.4.1.2.1} e de \sref[I]{I.9.1.5};
nossa afirmação resulta, portanto, de \sref{II.4.2.8} e \sref{II.4.2.9}.
\end{env}

\begin{env}[4.3.5]
\label{II.4.3.5}
Com a notação de \sref{II.4.3.1}, seja $\psi:Y'\to Y$ um morfismo, e ponhamos $\sh{E}'=\psi^*(\sh{E})$ e $\sh{F}'=\psi^*(\sh{F})$;
então o morfismo de Segre $\bb{P}(\sh{E}')\times\bb{P}(\sh{F}')\to\bb{P}(\sh{E}'\otimes\sh{F}')$ identifica-se com $\varsigma_{(Y')}$.
Com efeito, conservando a notação de \sref{II.4.3.1}, ponhamos, além disso, $P'_1=\bb{P}(\sh{E}')$ e $P'_2=\bb{P}(\sh{F}')$;
sabemos por \sref{II.4.1.3.1} que $P'_i$ se identifica com $(P_i)_{(Y')}$ ($i=1,2$), de modo que o morfismo estrutural $P'_1\times P'_2\to Y'$ se identifica com $r_{(Y')}$.
Por outro lado, $\sh{E}'\otimes\sh{F}'$ identifica-se com $\psi^*(\sh{E}\otimes\sh{F})$, portanto $\bb{P}(\sh{E}'\otimes\sh{F}')$ identifica-se com $(\bb{P}(\sh{E}\otimes\sh{F}))_{(Y')}$ \sref{II.4.1.3.1}.
Finalmente, $\sh{O}_{P'_1}(1)\otimes_{Y'}\sh{O}_{P'_2}(1)=\sh{L}'$ identifica-se com $\sh{L}\otimes_Y\sh{O}_{Y'}$, em virtude de \sref{II.4.1.3.2} e de \sref[I]{I.9.1.11}.
O homomorfismo canônico $(r_{(Y')})^*(\sh{E}'\otimes\sh{F}')\to\sh{L}'$ identifica-se então com $s_{(Y')}$, e nossa afirmação resulta de \sref{II.4.2.10}.
\end{env}

\begin{remark}[4.3.6]
\label{II.4.3.6}
O pré-esquema soma de $\bb{P}(\sh{E})$ e $\bb{P}(\sh{F})$ é igualmente canonicamente isomorfo a um \emph{subpré-esquema fechado de $\bb{P}(\sh{E}\oplus\sh{F})$}.
Com efeito, os homomorfismos sobrejetivos $\sh{E}\oplus\sh{F}\to\sh{E}$ e $\sh{E}\oplus\sh{F}\to\sh{F}$ correspondem a imersões fechadas $\bb{P}(\sh{E})\to\bb{P}(\sh{E}\oplus\sh{F})$ e $\bb{P}(\sh{F})\to\bb{P}(\sh{E}\oplus\sh{F})$;
tudo se reduz a ver que os espaços subjacentes dos subpré-esquemas fechados de $\bb{P}(\sh{E}\oplus\sh{F})$ assim obtidos não têm ponto em comum.
Sendo a questão local sobre $Y$, podemos adotar a notação de \sref{II.4.3.1};
ora, $\bb{S}_n(E)$ e $\bb{S}_n(F)$ identificam-se com submódulos de $\bb{S}_n(E\oplus F)$ cuja interseção se reduz a $0$;
e, se $\mathfrak{p}$ é um ideal primo graduado de $\bb{S}(E)$ tal que $\mathfrak{p}\cap\bb{S}_n(E)\neq\bb{S}_n(E)$ para todo $n\geq0$, corresponde-lhe em $\bb{S}(E\oplus F)$ um ideal primo graduado cuja interseção com $\bb{S}_n(E)$ é $\mathfrak{p}\cap\bb{S}_n(E)$, mas que \emph{contém} $\bb{S}_+(F)$, como se vê imediatamente;
dois pontos de $\Proj(\bb{S}(E))$ e $\Proj(\bb{S}(F))$, respectivamente, não podem, portanto, ter a mesma imagem em $\Proj(\bb{S}(E\oplus F))$.
\end{remark}

\subsection{Imersões nos fibrados projetivos; feixes muito amplos}
\label{subsection:II.4.4}

\begin{proposition}[4.4.1]
\label{II.4.4.1}
Sejam $Y$ um esquema quase-compacto, ou um pré-esquema cujo espaço subjacente é noetheriano, $q:X\to Y$ um morfismo \emph{de tipo finito} e $\sh{L}$ um $\sh{O}_X$-módulo invertível.
\begin{enumerate}
  \item[\rm{(i)}] Sejam $\sh{S}$ uma $\sh{O}_Y$-álgebra graduada quase-coerente em graus positivos e $\psi:q^*(\sh{S})\to\bigoplus_{n\geq0}\sh{L}^{\otimes n}$ um homomorfismo de álgebras graduadas.
    Para que $r_{\sh{L},\psi}$ esteja definido em toda parte e seja uma imersão, é necessário e
\oldpage[II]{79}
suficiente que existam um inteiro $n>0$ e um $\sh{O}_Y$-submódulo quase-coerente \emph{de tipo finito} $\sh{E}$ de $\sh{S}_n$ tais que o homomorfismo $\psi'=\psi_n\circ q^*(j):q^*(\sh{E})\to\sh{L}^{\otimes n}=\sh{L}'$ (sendo $j$ a injeção $\sh{E}\to\sh{S}_n$) seja sobrejetivo e que o morfismo $r_{\sh{L}',\psi'}:X\to\bb{P}(\sh{E})$ seja uma imersão.
  \item[\rm{(ii)}] Sejam $\sh{F}$ um $\sh{O}_Y$-módulo quase-coerente e $\vphi:q^*(\sh{F})\to\sh{L}$ um homomorfismo sobrejetivo.
    Para que o morfismo $r_{\sh{L},\vphi}$ seja uma imersão $X\to\bb{P}(\sh{F})$, é necessário e suficiente que exista um $\sh{O}_Y$-submódulo quase-coerente \emph{de tipo finito} $\sh{E}$ de $\sh{F}$ tal que o homomorfismo $\vphi'=\vphi\circ q^*(j):q^*(\sh{E})\to\sh{L}$ (onde $j:\sh{E}\to\sh{F}$ é a injeção canônica) seja sobrejetivo e que $r_{\sh{L},\vphi'}:X\to\bb{P}(\sh{E})$ seja uma imersão.
\end{enumerate}
\end{proposition}

\begin{proof}
\medskip\noindent
\begin{enumerate}
  \item[\rm{(i)}] O fato de $r_{\sh{L},\psi}$ estar definido em toda parte e ser uma imersão equivale, por \sref{II.3.8.5}, à existência de $n>0$ e de $\sh{E}$ tais que, se $\sh{S}'$ é a subálgebra de $\sh{S}$ gerada por $\sh{E}$, o homomorfismo $q^*(\sh{E})\to\sh{L}^{\otimes n}$ seja sobrejetivo e que $r_{\sh{L}',\psi'}:X\to\Proj(\sh{S}')$ esteja definido em toda parte e seja uma imersão.
    Temos, por outro lado, um homomorfismo canônico sobrejetivo $\bb{S}(\sh{E})\to\sh{S}'$, ao qual corresponde uma imersão fechada $\Proj(\sh{S}')\to\bb{P}(\sh{E})$ \sref{II.3.6.2};
    daí a conclusão.
  \item[\rm{(ii)}] Como $\sh{F}$ é limite indutivo de seus submódulos quase-coerentes de tipo finito $\sh{E}_\lambda$ \sref[I]{I.9.4.9}, $\bb{S}(\sh{F})$ é limite indutivo dos $\bb{S}(\sh{E}_\lambda)$;
    a conclusão resulta então de \sref{II.3.8.4}, observando que, na demonstração de \sref{II.3.8.4}, podemos tomar todos os $n_i$ iguais a $1$:
    com efeito (supondo $Y$ afim), se $r=r_{\sh{L},\vphi}$ é uma imersão, $r(X)$ é um subespaço quase-compacto de $\bb{P}(\sh{F})$ que se pode recobrir por um número finito de abertos de $\bb{P}(\sh{F})$ da forma $D_+(f)$, com $f\in F$, tais que $D_+(f)\cap r(X)$ seja fechado.
\end{enumerate}
\end{proof}

\begin{definition}[4.4.2]
\label{II.4.4.2}
Sejam $Y$ um pré-esquema e $q:X\to Y$ um morfismo.
Diz-se que um $\sh{O}_X$-módulo invertível $\sh{L}$ é \emph{muito amplo para $q$} (ou \emph{muito amplo para $Y$}, ou simplesmente \emph{muito amplo}, se daí não resultar confusão) se existem um $\sh{O}_Y$-módulo quase-coerente $\sh{E}$ e uma $Y$-imersão $i$ de $X$ em $P=\bb{P}(\sh{E})$ tais que $\sh{L}$ seja isomorfo a $i^*(\sh{O}_P(1))$.
\end{definition}

É equivalente \sref{II.4.2.3} dizer que existem um $\sh{O}_Y$-módulo quase-coerente $\sh{E}$ e um homomorfismo \emph{sobrejetivo} $\vphi:q^*(\sh{E})\to\sh{L}$ tais que $r_{\sh{L},\vphi}:X\to\bb{P}(\sh{E})$ seja uma \emph{imersão}.

Observaremos que a existência de um $\sh{O}_X$-módulo muito amplo relativamente a $Y$ implica que $q$ é \emph{separado} (\sref{II.3.1.3} e \sref[I]{I.5.5.1}[(i) e (ii)]).

\begin{corollary}[4.4.3]
\label{II.4.4.3}
Suponhamos que existam uma $\sh{O}_Y$-álgebra graduada quase-coerente $\sh{S}$ gerada por $\sh{S}_1$ e uma $Y$-imersão $i:X\to P=\Proj(\sh{S})$ tal que $\sh{L}$ seja isomorfo a $i^*(\sh{O}_P(1))$;
então $\sh{L}$ é muito amplo relativamente a $q$.
\end{corollary}

\begin{proof}
Com efeito, se $\sh{F}=\sh{S}_1$, o homomorfismo canônico $\bb{S}(\sh{F})\to\sh{S}$ é sobrejetivo; portanto, compondo a imersão fechada correspondente $\Proj(\sh{S})\to\bb{P}(\sh{F})$ \sref{II.3.6.2} e a imersão $i$, obtemos uma imersão $j:X\to\bb{P}(\sh{F})=P'$ tal que $\sh{L}$ é isomorfo a $j^*(\sh{O}_{P'}(1))$ \sref{II.3.6.3}.
\end{proof}

\begin{proposition}[4.4.4]
\label{II.4.4.4}
Sejam $q:X\to Y$ um morfismo quase-compacto e $\sh{L}$ um $\sh{O}_X$-módulo invertível.
As propriedades seguintes são equivalentes:
\begin{enumerate}
  \item[\rm{(a)}] $\sh{L}$ é muito amplo relativamente a $q$.
  \item[\rm{(b)}] $q_*(\sh{L})$ é quase-coerente, o homomorfismo canônico $\sigma:q^*(q_*(\sh{L}))\to\sh{L}$ é sobrejetivo, e o morfismo $r_{\sh{L},\sigma}:X\to\bb{P}(q_*(\sh{L}))$ é uma imersão.
\end{enumerate}
\end{proposition}

\begin{proof}
Como $q$ é quase-compacto, sabemos que $q_*(\sh{L})$ é quase-coerente se $q$ é separado \sref[I]{I.9.2.2}.
\oldpage[II]{80}
Sabemos, por \sref{II.3.4.7}, que a existência de um homomorfismo sobrejetivo $\vphi:q^*(\sh{E})\to\sh{L}$ (sendo $\sh{E}$ um $\sh{O}_Y$-módulo quase-coerente) implica que $\sigma$ é sobrejetivo;
além disso, à fatoração $q^*(\sh{E})\to q^*(q_*(\sh{L}))\xrightarrow{\sigma}\sh{L}$ de $\vphi$ corresponde canonicamente uma fatoração
\[
  q^*(\bb{S}(\sh{E}))\to q^*(\bb{S}(q_*(\sh{L})))\to\bigoplus_{n\geq0}\sh{L}^{\otimes n}
\]
portanto, por \sref{II.3.8.3}, a hipótese de que $r_{\sh{L},\vphi}$ seja uma imersão implica que o mesmo vale para $j=r_{\sh{L},\sigma}$;
além disso, por \sref{II.4.2.4}, $\sh{L}$ é isomorfo a $j^*(\sh{O}_{P'}(1))$, pondo $P'=\bb{P}(q_*(\sh{L}))$.
Vemos, pois, que (a) e (b) são equivalentes.
\end{proof}

\begin{corollary}[4.4.5]
\label{II.4.4.5}
Suponhamos que $q$ seja quase-compacto.
Para que $\sh{L}$ seja muito amplo relativamente a $Y$, é necessário e suficiente que exista um recobrimento aberto $(U_\alpha)$ de $Y$ tal que $\sh{L}|_{q^{-1}(U_\alpha)}$ seja muito amplo relativamente a $U_\alpha$, para todo $\alpha$.
\end{corollary}

\begin{proof}
Com efeito, a condição~(b) de \sref{II.4.4.4} é local sobre $Y$.
\end{proof}

\begin{proposition}[4.4.6]
\label{II.4.4.6}
Sejam $Y$ um esquema quase-compacto, ou um pré-esquema cujo espaço subjacente seja noetheriano, $q:X\to Y$ um morfismo \emph{de tipo finito}, e $\sh{L}$ um $\sh{O}_X$-módulo invertível.
Então as condições~(a) e~(b) de \sref{II.4.4.4} são equivalentes às seguintes:
\begin{enumerate}
  \item[\rm{(a')}] Existe um $\sh{O}_Y$-módulo quase-coerente $\sh{E}$ \emph{de tipo finito} e um homomorfismo sobrejetivo $\vphi:q^*(\sh{E})\to\sh{L}$ tal que $r_{\sh{L},\vphi}$ seja uma imersão.
  \item[\rm{(b')}] Existe um $\sh{O}_Y$-submódulo quase-coerente $\sh{E}$ de $q_*(\sh{L})$, \emph{de tipo finito}, que possui as propriedades enunciadas em~(a').
\end{enumerate}
\end{proposition}

\begin{proof}
É claro que (a') ou (b') implica (a);
por outro lado, (a) implica (a') em virtude de \sref{II.4.4.1}, e, do mesmo modo, (b) implica (b').
\end{proof}

\begin{corollary}[4.4.7]
\label{II.4.4.7}
Suponhamos que $Y$ seja um esquema quase-compacto ou um pré-esquema noetheriano.
Se $\sh{L}$ é muito amplo para $q$, existe uma $\sh{O}_Y$-álgebra graduada quase-coerente $\sh{S}$ tal que $\sh{S}_1$ seja de tipo finito e gere $\sh{S}$, e uma $Y$-imersão \emph{aberta dominante} $i:X\to P=\Proj(\sh{S})$ tal que $\sh{L}$ seja isomorfo a $i^*(\sh{O}_P(1))$.
\end{corollary}

\begin{proof}
Com efeito, verifica-se a condição~(b') de \sref{II.4.4.6};
o morfismo estrutural $p:\bb{P}(\sh{E})=P'\to Y$ é então separado e de tipo finito \sref{II.3.1.3}, de modo que $P'$ é um esquema quase-compacto (resp. um pré-esquema noetheriano) se $Y$ é um esquema quase-compacto (resp. um pré-esquema noetheriano).
Seja $Z$ o fecho \sref[I]{I.9.5.10} do subpré-esquema $X'$ de $P'$ associado à imersão $j=r_{\sh{L},\vphi}$ de $X$ em $P'$;
é claro que $j$ se fatora em uma imersão aberta dominante $i:X\to Z$, seguida da injeção canônica $Z\to P'$.
Mas $Z$ identifica-se com um pré-esquema $\Proj(\sh{S})$, onde $\sh{S}$ é uma $\sh{O}_Y$-álgebra graduada, quociente de $\bb{S}(\sh{E})$ por um feixe graduado quase-coerente de ideais \sref{II.3.6.2}, e é claro que $\sh{S}_1$ é de tipo finito e gera $\sh{S}$;
além disso, $\sh{O}_Z(1)$ é a imagem inversa de $\sh{O}_{P'}(1)$ pela injeção canônica \sref{II.3.6.3}, logo $\sh{L}=i^*(\sh{O}_Z(1))$.
\end{proof}

\begin{proposition}[4.4.8]
\label{II.4.4.8}
Sejam $q:X\to Y$ um morfismo, $\sh{L}$ um $\sh{O}_X$-módulo muito amplo relativamente a $q$, e $\sh{L}'$ um $\sh{O}_X$-módulo invertível, tais que exista um $\sh{O}_Y$-módulo quase-coerente $\sh{E}'$ e um homomorfismo sobrejetivo $q^*(\sh{E}')\to\sh{L}'$.
Então $\sh{L}\otimes_{\sh{O}_X}\sh{L}'$ é muito amplo relativamente a $q$.
\end{proposition}

\begin{proof}
A hipótese implica a existência de um $Y$-morfismo $r':X\to P'=\bb{P}(\sh{E}')$ tal que $\sh{L}'=r'^*(\sh{O}_{P'}(1))$ \sref{II.4.2.1}.
Existe, por hipótese, um $\sh{O}_Y$-módulo quase-coerente $\sh{E}$ e uma
\oldpage[II]{81}
$Y$-imersão $r:X\to P=\bb{P}(\sh{E})$ tal que $\sh{L}=r^*(\sh{O}_P(1))$.
Ponhamos $Q=\bb{P}(\sh{E}\otimes\sh{E}')$ e consideremos o morfismo de Segre $\varsigma:P\times_Y P'\to Q$ \sref{II.4.3.1}.
Como $r$ é uma imersão, o mesmo vale para $(r,r')_Y:X\to P\times_Y P'$ \sref[I]{I.5.3.14};
mas, como $\varsigma$ é uma imersão \sref{II.4.3.3}, o mesmo vale para $r'':X\xrightarrow{(r,r')}P\times_Y P'\xrightarrow{\varsigma}Q$.
Por outro lado, por \sref{II.4.3.2.1}, $\varsigma^*(\sh{O}_Q(1))$ é isomorfo a $\sh{O}_P(1)\otimes_Y\sh{O}_{P'}(1)$; logo, por \sref[I]{I.9.1.4}, $r''^*(\sh{O}_Q(1))$ é isomorfo a $\sh{L}\otimes\sh{L}'$, o que demonstra a proposição.
\end{proof}

\begin{corollary}[4.4.9]
\label{II.4.4.9}
Seja $q:X\to Y$ um morfismo.
\begin{enumerate}
  \item[\rm{(i)}] Seja $\sh{L}$ um $\sh{O}_X$-módulo invertível, e $\sh{K}$ um $\sh{O}_Y$-módulo invertível.
  Para que $\sh{L}$ seja muito amplo relativamente a $q$, é necessário e suficiente que $\sh{L}\otimes q^*(\sh{K})$ o seja.
  \item[\rm{(ii)}] Se $\sh{L}$ e $\sh{L}'$ são dois $\sh{O}_X$-módulos muito amplos relativamente a $q$, o mesmo vale para $\sh{L}\otimes\sh{L}'$;
  em particular, $\sh{L}^{\otimes n}$ é muito amplo relativamente a $q$ para todo $n>0$.
\end{enumerate}
\end{corollary}

\begin{proof}
A afirmação~(ii) é consequência imediata de \sref{II.4.4.8}, assim como a necessidade da condição~(i);
por outro lado, se $\sh{L}\otimes q^*(\sh{K})$ é muito amplo, o mesmo vale para $(\sh{L}\otimes q^*(\sh{K}))\otimes q^*(\sh{K}^{-1})$ em virtude do que precede, e este último $\sh{O}_X$-módulo é isomorfo a $\sh{L}$ (\sref[0]{0.5.4.3} e \sref[0]{0.5.4.5}).
\end{proof}

\begin{proposition}[4.4.10]
\label{II.4.4.10}
\medskip\noindent
\begin{enumerate}
  \item[\rm{(i)}] Para todo pré-esquema $Y$, todo $\sh{O}_Y$-módulo invertível $\sh{L}$ é muito amplo relativamente ao morfismo identidade $1_Y$.
  \item[\rm{(i \emph{bis})}] Sejam $f:X\to Y$ um morfismo e $j:X'\to X$ uma imersão.
  Se $\sh{L}$ é um $\sh{O}_X$-módulo muito amplo relativamente a $f$, então $j^*(\sh{L})$ é muito amplo relativamente a $f\circ j$.
  \item[\rm{(ii)}] Sejam $Z$ um pré-esquema quase-compacto, $f:X\to Y$ um morfismo de tipo finito, $g:Y\to Z$ um morfismo quase-compacto, $\sh{L}$ um $\sh{O}_X$-módulo muito amplo relativamente a $f$, e $\sh{K}$ um $\sh{O}_Y$-módulo muito amplo relativamente a $g$.
  Então existe um inteiro $n_0>0$ tal que $\sh{L}\otimes f^*(\sh{K}^{\otimes n})$ seja muito amplo relativamente a $g\circ f$ para todo $n\geq n_0$.
  \item[\rm{(iii)}] Sejam $f:X\to Y$ e $g:Y'\to Y$ dois morfismos, e ponhamos $X'=X_{(Y')}$.
  Se $\sh{L}$ é um $\sh{O}_X$-módulo muito amplo relativamente a $f$, então $\sh{L}'=\sh{L}\otimes_Y\sh{O}_{Y'}$ é um $\sh{O}_{X'}$-módulo muito amplo relativamente a $f_{(Y')}$.
  \item[\rm{(iv)}] Sejam $f_i:X_i\to Y_i$ ($i=1,2$) dois $S$-morfismos.
  Se $\sh{L}_i$ é um $\sh{O}_{X_i}$-módulo muito amplo relativamente a $f_i$ ($i=1,2$), então $\sh{L}_1\otimes_S\sh{L}_2$ é muito amplo relativamente a $f_1\times_S f_2$.
  \item[\rm{(v)}] Sejam $f:X\to Y$ e $g:Y\to Z$ dois morfismos.
  Se um $\sh{O}_X$-módulo $\sh{L}$ é muito amplo relativamente a $g\circ f$, então $\sh{L}$ é muito amplo relativamente a $f$.
  \item[\rm{(vi)}] Sejam $f:X\to Y$ um morfismo e $j$ a injeção canônica $X_\red\to X$.
  Se um $\sh{O}_X$-módulo $\sh{L}$ é muito amplo relativamente a $f$, então $j^*(\sh{L})$ é muito amplo relativamente a $f_\red$.
\end{enumerate}
\end{proposition}

\begin{proof}
A propriedade~(i~\emph{bis}) decorre imediatamente da definição \sref{II.4.4.2}, e é imediato que (vi) se deduz formalmente de (i~\emph{bis}) e de (v) por um raciocínio calcado no da demonstração de \sref[I]{I.5.5.12}, que deixamos ao leitor.
Para demonstrar (v), consideremos, como em \sref[I]{I.5.5.12}, a fatoração $X\xrightarrow{\Gamma_f}X\times_Z Y\xrightarrow{p_2}Y$, observando que $p_2=(g\circ f)\times1_Y$.
Resulta da hipótese e de (i) e (iv) que $\sh{L}\otimes_{\sh{O}_Z}\sh{O}_Y$ é muito amplo para $p_2$;
por outro lado, $\sh{L}=\Gamma_f^*(\sh{L}\otimes_{\sh{O}_Z}\sh{O}_Y)$ \sref[I]{I.9.1.4}, e $\Gamma_f$ é uma imersão \sref[I]{I.5.3.11};
podemos, pois, aplicar (i~\emph{bis}).
\oldpage[II]{82}
Para demonstrar (i), aplicamos a definição \sref{II.4.4.2} com $\sh{E}=\sh{L}$, observando que, então, $\bb{P}(\sh{E})$ se identifica com $Y$ \sref{II.4.1.4}.

Demonstremos (iii).
Existe um $\sh{O}_Y$-módulo quase-coerente $\sh{E}$ e uma $Y$-imersão $i:X\to\bb{P}(\sh{E})=P$ tal que $\sh{L}=i^*(\sh{O}_P(1))$;
então, se pusermos $\sh{E}'=g^*(\sh{E})$, $\sh{E}'$ é um $\sh{O}_{Y'}$-módulo quase-coerente, e temos $P'=\bb{P}(\sh{E}')=P_{(Y')}$ \sref{II.4.1.3.1}; $i_{(Y')}$ é uma imersão de $X_{(Y')}$ em $P'$ \sref[I]{I.4.3.2}, e $\sh{L}'$ é isomorfo a $(i_{(Y')})^*(\sh{O}_{P'}(1))$ \sref{II.4.2.10}.

Para demonstrar (iv), observemos que existe, por hipótese, uma $Y_i$-imersão $r_i:X_i\to P_i=\bb{P}(\sh{E}_i)$, onde $\sh{E}_i$ é um $\sh{O}_{Y_i}$-módulo quase-coerente, e $\sh{L}_i=r_i^*(\sh{O}_{P_i}(1))$ ($i=1,2$);
$r_1\times_S r_2$ é uma $S$-imersão de $X_1\times_S X_2$ em $P_1\times_S P_2$ \sref[I]{I.4.3.1}, e a imagem inversa de $\sh{O}_{P_1}(1)\otimes_S\sh{O}_{P_2}(1)$ por essa imersão é $\sh{L}_1\otimes_S\sh{L}_2$.
Por outro lado, ponhamos $T=Y_1\times_S Y_2$, e sejam $p_1,p_2$ as projeções respectivas de $T$ em $Y_1,Y_2$.
Se pusermos $P'_i=\bb{P}(p_i^*(\sh{E}_i))$ ($i=1,2$), temos, por \sref{II.4.1.3.1}, $P'_i=P_i\times_{Y_i}T$, de onde
\[
  P'_1\times_T P'_2
  =(P_1\times_{Y_1}T)\times_T(P_2\times_{Y_2}T)
  =P_1\times_{Y_1}(T\times_{Y_2}P_2)
  =P_1\times_{Y_1}(Y_1\times_S P_2)
  =P_1\times_S P_2
\]
salvo isomorfismo canônico.
Do mesmo modo, temos $\sh{O}_{P'_i}(1)=\sh{O}_{P_i}(1)\otimes_{Y_i}\sh{O}_T$ \sref{II.4.1.3.2}, e um cálculo análogo (baseado, em particular, em \sref[I]{I.9.1.9.1} e \sref[I]{I.9.1.2}) mostra que, na identificação precedente, $\sh{O}_{P'_1}(1)\otimes_T\sh{O}_{P'_2}(1)$ se identifica com $\sh{O}_{P_1}(1)\otimes_S\sh{O}_{P_2}(1)$.
Podemos, pois, considerar $r_1\times_S r_2$ como uma $T$-imersão de $X_1\times_S X_2$ em $P'_1\times_T P'_2$, sendo a imagem inversa de $\sh{O}_{P'_1}(1)\otimes_T\sh{O}_{P'_2}(1)$ por essa imersão igual a $\sh{L}_1\otimes_S\sh{L}_2$.
Terminamos então o raciocínio como em \sref{II.4.4.8}, utilizando o morfismo de Segre.

Resta demonstrar (ii).
Podemos primeiro nos restringir ao caso em que $Z$ é um esquema afim, pois, em geral, existe uma cobertura finita $(U_i)$ de $Z$ por abertos afins;
se a proposição estiver demonstrada para $\sh{K}|_{g^{-1}(U_i)}$, $\sh{L}|_{f^{-1}(g^{-1}(U_i))}$ e um inteiro $n_i$, bastará tomar para $n_0$ o maior dos $n_i$ a fim de demonstrar a proposição para $\sh{K}$ e $\sh{L}$ \sref{II.4.4.5}.
A hipótese implica que $f$ e $g$ são morfismos separados, logo $X$ e $Y$ são \emph{esquemas} quase-compactos.

Existe uma imersão $r:X\to P=\bb{P}(\sh{E})$, onde $\sh{E}$ é um $\sh{O}_Y$-módulo quase-coerente \emph{de tipo finito}, e $\sh{L}=r^*(\sh{O}_P(1))$, em virtude de \sref{II.4.4.6}.
Veremos que existe um $\sh{O}_P$-módulo $\sh{M}$ muito amplo relativamente ao morfismo composto $P\xrightarrow{h}Y\xrightarrow{g}Z$, tal que $\sh{O}_P(1)$ seja isomorfo a $\sh{M}\otimes_Y\sh{K}^{\otimes(-m)}$ para algum inteiro $m$.
Para $n\geq m+1$, $\sh{O}_P(1)\otimes_Y\sh{K}^{\otimes n}$ será, portanto, muito amplo para $Z$ em virtude da hipótese e de (iv) aplicado aos morfismos $h:P\to Y$ e $1_Y$;
como $r$ é uma imersão e $\sh{L}\otimes f^*(\sh{K}^{\otimes n})=r^*(\sh{O}_P(1)\otimes_Y\sh{K}^{\otimes n})$, a conclusão resultará de (i~\emph{bis}).
Para estabelecer nossa afirmação sobre $\sh{O}_P(1)$, utilizaremos o lema seguinte:

\begin{lemma}[4.4.10.1]
\label{II.4.4.10.1}
Sejam $Z$ um esquema quase-compacto, ou um pré-esquema cujo espaço subjacente seja noetheriano, $g:Y\to Z$ um morfismo quase-compacto, $\sh{K}$ um $\sh{O}_Y$-módulo invertível muito amplo para $g$, e $\sh{E}$ um $\sh{O}_Y$-módulo quase-coerente de tipo finito.
Existe então um inteiro $m_0$ tal que, para todo $m\geq m_0$, $\sh{E}$ seja isomorfo a um quociente de um $\sh{O}_Y$-módulo da forma $g^*(\sh{F})\otimes\sh{K}^{\otimes(-m)}$, onde $\sh{F}$ é um $\sh{O}_Z$-módulo quase-coerente de tipo finito (dependente de $m$).
\end{lemma}

Esse lema será demonstrado em \sref{II.4.5.10.1};
o leitor poderá verificar que \sref{II.4.4.10} não é utilizado em parte alguma de \sref{subsection:II.4.5}.
\oldpage[II]{83}
Admitido esse lema, existe uma imersão fechada $j_1$ de $P$ em
\[
  P_1=\bb{P}(g^*(\sh{F})\otimes\sh{K}^{\otimes(-m)})
\]
tal que $\sh{O}_P(1)$ seja isomorfo a $j_1^*(\sh{O}_{P_1}(1))$ \sref{II.4.1.2}.
Por outro lado, existe um isomorfismo de $P_1$ sobre $P_2=\bb{P}(g^*(\sh{F}))$ que identifica $\sh{O}_{P_1}(1)$ com $\sh{O}_{P_2}(1)\otimes_Y\sh{K}^{\otimes(-m)}$ \sref{II.4.1.4};
temos, portanto, uma imersão fechada $j_2:P\to P_2$ tal que $\sh{O}_P(1)$ seja isomorfo a $j_2^*(\sh{O}_{P_2}(1))\otimes_Y\sh{K}^{\otimes(-m)}$.
Finalmente, $P_2$ identifica-se com $P_3\times_Z Y$, onde $P_3=\bb{P}(\sh{F})$, e $\sh{O}_{P_2}(1)$ com $\sh{O}_{P_3}(1)\otimes_Z\sh{O}_Y$ \sref{II.4.1.3}.
Por definição, $\sh{O}_{P_3}(1)$ é muito amplo para $Z$;
como o mesmo vale para $\sh{K}$, concluímos de (iv) que $\sh{O}_{P_2}(1)\otimes_Y\sh{K}$ é muito amplo para $Z$;
o mesmo vale para $\sh{M}=j_2^*(\sh{O}_{P_2}(1)\otimes_Y\sh{K})$ em virtude de (i~\emph{bis}), e $\sh{O}_P(1)$ é isomorfo a $\sh{M}\otimes_Y\sh{K}^{\otimes(-m-1)}$, o que termina a demonstração.
\end{proof}

\begin{proposition}[4.4.11]
\label{II.4.4.11}
Sejam $f:X\to Y$ e $f':X'\to Y$ dois morfismos, $X''$ o pré-esquema soma $X\sqcup X'$, e $f''$ o morfismo $X''\to Y$ que coincide com $f$ (resp. $f'$) em $X$ (resp. $X'$).
Seja $\sh{L}$ (resp. $\sh{L}'$) um $\sh{O}_X$-módulo (resp. um $\sh{O}_{X'}$-módulo) invertível, e seja $\sh{L}''$ o $\sh{O}_{X''}$-módulo invertível que coincide com $\sh{L}$ em $X$ e com $\sh{L}'$ em $X'$.
Para que $\sh{L}''$ seja muito amplo relativamente a $f''$, é necessário e suficiente que $\sh{L}$ seja muito amplo relativamente a $f$ e que $\sh{L}'$ seja muito amplo relativamente a $f'$.
\end{proposition}

\begin{proof}
Reduzimo-nos imediatamente ao caso em que $Y$ é afim.
Se $\sh{L}''$ é muito amplo, o mesmo vale para $\sh{L}$ e $\sh{L}'$ em virtude de \sref{II.4.4.10}[(i~\emph{bis})].
Reciprocamente, se $\sh{L}$ e $\sh{L}'$ são muito amplos, resulta imediatamente da definição \sref{II.4.4.2} e de \sref{II.4.3.6} que $\sh{L}''$ é muito amplo.
\end{proof}

\subsection{Feixes amplos}
\label{subsection:II.4.5}

\begin{env}[4.5.1]
\label{II.4.5.1}
Dados um pré-esquema $X$ e um $\sh{O}_X$-módulo invertível $\sh{L}$, poremos, para todo $\sh{O}_X$-módulo $\sh{F}$ (quando não for possível nenhuma confusão acerca de $\sh{L}$), $\sh{F}(n)=\sh{F}\otimes_{\sh{O}_X}\sh{L}^{\otimes n}$ ($n\in\bb{Z}$);
poremos, por outro lado, $S=\bigoplus_{n\geq0}\Gamma(X,\sh{L}^{\otimes n})$ (subanel graduado do anel $\Gamma_\bullet(\sh{L})$ definido em \sref[0]{0.5.4.6}).
Se considerarmos $X$ como um $\bb{Z}$-pré-esquema e designarmos por $p$ o morfismo estrutural $X\to\Spec(\bb{Z})$, há uma correspondência biunívoca entre os homomorfismos $p^*(\widetilde{S})\to\bigoplus_{n\geq0}\sh{L}^{\otimes n}$ de $\sh{O}_X$-álgebras graduadas e os endomorfismos do anel graduado $S$ \sref[I]{I.2.2.5};
o homomorfismo $\varepsilon:p^*(\widetilde{S})\to\bigoplus_{n\geq0}\sh{L}^{\otimes n}$ que corresponde assim ao automorfismo \emph{identidade} de $S$ será dito \emph{canônico}.
A ele corresponde \sref{II.3.7.1} um morfismo $G(\varepsilon)\to\Proj(S)$, que também será dito \emph{canônico}.
\end{env}

\begin{theorem}[4.5.2]
\label{II.4.5.2}
Sejam $X$ um esquema quase-compacto ou um pré-esquema cujo espaço subjacente seja noetheriano, $\sh{L}$ um $\sh{O}_X$-módulo invertível, e $S$ o anel graduado $\bigoplus_{n\geq0}\Gamma(X,\sh{L}^{\otimes n})$.
As condições seguintes são equivalentes:
\begin{enumerate}
  \item[\rm{(a)}] Quando $f$ percorre o conjunto dos elementos homogêneos de $S_+$, os $X_f$ formam uma base da topologia de $X$.
  \item[\rm{(a')}] Quando $f$ percorre o conjunto dos elementos homogêneos de $S_+$, aqueles $X_f$ que são afins formam um recobrimento de $X$.
  \item[\rm{(b)}] O morfismo canônico $G(\varepsilon)\to\Proj(S)$ \sref{II.4.5.1} está definido em toda parte e é uma imersão aberta dominante.
\oldpage[II]{84}
  \item[\rm{(b')}] O morfismo canônico $G(\varepsilon)\to\Proj(S)$ está definido em toda parte e é um homeomorfismo do espaço subjacente a $X$ sobre um subespaço de $\Proj(S)$.
  \item[\rm{(c)}] Para todo $\sh{O}_X$-módulo quase-coerente $\sh{F}$, se designarmos por $\sh{F}_n$ o $\sh{O}_X$-submódulo de $\sh{F}(n)$ gerado pelas seções de $\sh{F}(n)$ sobre $X$, então $\sh{F}$ é a soma dos $\sh{O}_X$-submódulos $\sh{F}_n(-n)$, para os inteiros $n>0$.
  \item[\rm{(c')}] A propriedade~\rm{(c)} verifica-se para todo feixe quase-coerente de ideais de $\sh{O}_X$.
\end{enumerate}

Além disso, se $(f_\alpha)$ é uma família de elementos homogêneos de $S_+$ tal que os $X_{f_\alpha}$ sejam afins, então a restrição a $\bigcup_\alpha X_{f_\alpha}$ do morfismo canônico $X\to\Proj(S)$ é um isomorfismo de $\bigcup_\alpha X_{f_\alpha}$ sobre $\bigcup_\alpha(\Proj(S))_{f_\alpha}$.
\end{theorem}

\begin{proof}
É claro que (b) implica (b'), e (b') implica (a) em virtude da fórmula \sref{II.3.7.3.1} (levando em conta que $\varepsilon^\flat$ é a identidade).
A condição~(a) implica (a'), pois todo $x\in X$ possui uma vizinhança afim $U$ tal que $\sh{L}|U$ seja isomorfo a $\sh{O}_X|U$;
se $f\in\Gamma(X,\sh{L}^{\otimes n})$ é tal que $x\in X_f\subset U$, então $X_f$ é também o conjunto dos $x'\in U$ tais que $(f|U)(x')\neq0$, sendo, portanto, um aberto afim \sref[I]{I.1.3.6}.
Para demonstrar que (a') implica (b), basta demonstrar a última afirmação do enunciado e verificar, além disso, que, se $X=\bigcup_\alpha X_{f_\alpha}$, a condição~(iv) de \sref{II.3.8.2} é satisfeita.
Este último ponto resulta imediatamente de \sref[I]{I.9.3.1}[(i)].
Quanto à última afirmação de \sref{II.4.5.2}, como $X_{f_\alpha}$ é a imagem inversa de $(\Proj(S))_{f_\alpha}$ por $G(\varepsilon)\to\Proj(S)$, basta aplicar \sref[I]{I.9.3.2}.
Logo (a), (a'), (b) e (b') são equivalentes.

Para mostrar que (a') implica (c), observemos que, se $X_f$ é afim (com $f\in S_k$), então $\sh{F}|X_f$ é gerado por suas seções sobre $X_f$ \sref[I]{I.1.3.9};
por outro lado, por \sref[I]{I.9.3.1}[(ii)], tal seção $s$ é da forma $(t|X_f)\otimes(f|X_f)^{-m}$, onde $t\in\Gamma(X,\sh{F}(km))$;
por definição, $t$ é também uma seção de $\sh{F}_{km}$, logo $s$ é de fato uma seção de $\sh{F}_{km}(-km)$ sobre $X_f$, o que demonstra (c).
É claro que (c) implica (c'), e resta demonstrar que (c') implica (a).
Ora, seja $U$ uma vizinhança aberta de $x\in X$, e seja $\sh{J}$ um feixe quase-coerente de ideais de $\sh{O}_X$ que define um subpré-esquema fechado de $X$ cujo espaço subjacente seja $X\setmin U$ \sref[I]{I.5.2.1}.
A hipótese~(c') implica que existe um inteiro $n>0$ e uma seção $f$ de $\sh{J}(n)$ sobre $X$ tal que $f(x)\neq0$.
Mas temos evidentemente $f\in S_n$ e $x\in X_f\subset U$, o que demonstra (a).
\end{proof}

Quando $X$ é um pré-esquema cujo espaço subjacente é noetheriano, as condições equivalentes de \sref{II.4.5.2} implicam que $X$ é um \emph{esquema}, pois é isomorfo a um subpré-esquema do esquema $\Proj(S)$, em virtude de \sref{II.4.5.2}[(b)].

\begin{definition}[4.5.3]
\label{II.4.5.3}
Diz-se que um $\sh{O}_X$-módulo invertível $\sh{L}$ é \emph{amplo} se $X$ é um esquema quase-compacto e se as condições equivalentes de \sref{II.4.5.2} são verificadas.
\end{definition}

Resulta evidentemente do critério~(a) de \sref{II.4.5.2} que, se $\sh{L}$ é um $\sh{O}_X$-módulo amplo, então, para todo aberto quase-compacto $U$ de $X$, $\sh{L}|U$ é um $(\sh{O}_X|U)$-módulo amplo.

Resulta da demonstração de \sref{II.4.5.2} que os $X_f$ \emph{afins} já formam uma base da topologia de $X$.
Além disso:

\begin{corollary}[4.5.4]
\label{II.4.5.4}
Seja $\sh{L}$ um $\sh{O}_X$-módulo amplo.
Para todo subespaço finito $Z$ de $X$ e toda vizinhança $U$ de $Z$, existem um inteiro $n$ e um $f\in\Gamma(X,\sh{L}^{\otimes n})$ tais que $X_f$ seja uma vizinhança afim de $Z$ contida em $U$.
\end{corollary}

\begin{proof}
\oldpage[II]{85}
Em virtude de \sref{II.4.5.2}[(b)], podemos limitar-nos a demonstrar que, para todo subconjunto finito $Z'$ de $\Proj(S)$ e toda vizinhança aberta $U$ de $Z'$, existe um elemento homogêneo $f\in S_+$ tal que $Z'\subset(\Proj(S))_f\subset U$ \sref{II.2.4.1}.
Ora, por definição, o conjunto fechado $Y$, complementar de $U$ em $\Proj(S)$, é da forma $V_+(\mathfrak{I})$, onde $\mathfrak{I}$ é um ideal graduado de $S$ que não contém $S_+$ \sref{II.2.3.2};
por outro lado, os pontos de $Z'$ são, por definição, ideais primos graduados $\mathfrak{p}_i$ de $S$ que não contêm $\mathfrak{I}$ \sref{II.2.3.1}.
Existe, portanto, um elemento $f\in\mathfrak{I}$ que não pertence a nenhum dos $\mathfrak{p}_i$ (Bourbaki, \emph{Alg. comm.}, cap.~II, \S1, n.º~1, prop.~2), e, como os $\mathfrak{p}_i$ são graduados, o raciocínio feito \emph{loc. cit.} mostra que podemos até mesmo tomar $f$ homogêneo;
esse elemento satisfaz então o que se queria.
\end{proof}

\begin{proposition}[4.5.5]
\label{II.4.5.5}
Suponhamos que $X$ seja um esquema quase-compacto, ou um pré-esquema cujo espaço subjacente seja noetheriano.
As condições~(a) a (c') de \sref{II.4.5.2} são também equivalentes às seguintes:
\begin{enumerate}
  \item[\rm{(d)}] Para todo $\sh{O}_X$-módulo quase-coerente $\sh{F}$ de tipo finito, existe um inteiro $n_0$ tal que, para todo $n\geq n_0$, $\sh{F}(n)$ seja gerado por suas seções sobre $X$.
  \item[\rm{(d')}] Para todo $\sh{O}_X$-módulo quase-coerente $\sh{F}$ de tipo finito, existem inteiros $n>0$ e $k>0$ tais que $\sh{F}$ seja isomorfo a um quociente do $\sh{O}_X$-módulo $\sh{L}^{\otimes(-n)}\otimes\sh{O}_X^k$.
  \item[\rm{(d'')}] A propriedade~(d') verifica-se para todo feixe quase-coerente de ideais de $\sh{O}_X$ de tipo finito.
\end{enumerate}
\end{proposition}

\begin{proof}
Como $X$ é quase-compacto, se um $\sh{O}_X$-módulo quase-coerente de tipo finito $\sh{F}$ é tal que $\sh{F}(n)$ (que é de tipo finito) é gerado por suas seções sobre $X$, então $\sh{F}(n)$ é gerado por um número \emph{finito} dessas seções \sref[0]{0.5.2.3};
logo (d) implica (d'), e é claro que (d') implica (d'').
Como todo $\sh{O}_X$-módulo quase-coerente $\sh{G}$ é limite indutivo de seus sub-$\sh{O}_X$-módulos de tipo finito \sref[I]{I.9.4.9}, para verificar a condição~(c') de \sref{II.4.5.2}, basta verificá-la para um feixe quase-coerente de ideais de $\sh{O}_X$ que seja de tipo finito, e (d'') implica, portanto, (c').
Resta ver que, se $\sh{L}$ é amplo, a propriedade~(d) é verificada.
Consideremos um recobrimento finito de $X$ pelos $X_{f_i}$ ($f_i\in S_{n_i}$), que podemos supor afins;
substituindo os $f_i$ por potências convenientes (o que não altera os $X_{f_i}$), podemos supor que todos os $n_i$ sejam iguais a um mesmo inteiro $m$.
O feixe $\sh{F}|X_{f_i}$, sendo de tipo finito por hipótese, é gerado por um número finito de suas seções $h_{ij}$ sobre $X_{f_i}$ \sref[I]{I.1.3.13};
existe, portanto, um inteiro $k_0$ tal que a seção $h_{ij}\otimes f_i^{\otimes k_0}$ se estenda a uma seção de $\sh{F}(k_0m)$ sobre $X$ para todo par $(i,j)$ \sref[I]{I.9.3.1}.
\emph{A fortiori}, os $h_{ij}\otimes f_i^{\otimes k}$ estendem-se a seções de $\sh{F}(km)$ sobre $X$ para todo $k\geq k_0$ e, para esses valores de $k$, $\sh{F}(km)$ é, portanto, gerado por suas seções sobre $X$.
Para todo $p$ tal que $0<p<m$, $\sh{F}(p)$ também é de tipo finito; logo existe um inteiro $k_p$ tal que $\sh{F}(p)(km)=\sh{F}(p+km)$ seja gerado por suas seções sobre $X$ para todo $k\geq k_p$.
Tomando $n_0$ maior do que todos os $k_pm$, concluímos que $\sh{F}(n)$ é gerado por suas seções sobre $X$ para todo $n\geq n_0$, pois tal $n$ se escreve $n=km+p$, com $k\geq k_p$ e $0\leq p<m$.
\end{proof}

\begin{proposition}[4.5.6]
\label{II.4.5.6}
Sejam $X$ um esquema quase-compacto e $\sh{L}$ um $\sh{O}_X$-módulo invertível.
\begin{enumerate}
  \item[\rm{(i)}] Seja $n>0$ um inteiro. Para que $\sh{L}$ seja amplo, é necessário e suficiente que $\sh{L}^{\otimes n}$ seja amplo.
  \item[\rm{(ii)}] Seja $\sh{L}'$ um $\sh{O}_X$-módulo invertível tal que, para todo $x\in X$, exista um inteiro $n>0$
\oldpage[II]{86}
    e uma seção $s'$ de $\sh{L}'^{\otimes n}$ sobre $X$ tais que $s'(x)\neq0$.
    Então, se $\sh{L}$ é amplo, $\sh{L}\otimes\sh{L}'$ também é amplo.
\end{enumerate}
\end{proposition}

\begin{proof}
A propriedade~(i) é consequência evidente do critério~(a) de \sref{II.4.5.2}, pois $X_{f^{\otimes n}}=X_f$.
Por outro lado, se $\sh{L}$ é amplo, então, para todo $x\in X$ e toda vizinhança $U$ de $x$, existem $m>0$ e $f\in\Gamma(X,\sh{L}^{\otimes m})$ tais que $x\in X_f\subset U$ \sref{II.4.5.2}[(a)];
se $f'\in\Gamma(X,\sh{L}'^{\otimes n})$ é tal que $f'(x)\neq0$, então $s(x)\neq0$ para $s=f^{\otimes n}\otimes f'^{\otimes m}\in\Gamma(X,(\sh{L}\otimes\sh{L}')^{\otimes mn})$; logo $x\in X_s\subset X_f\subset U$, o que demonstra que $\sh{L}\otimes\sh{L}'$ é amplo \sref{II.4.5.2}[(a)].
\end{proof}

\begin{corollary}[4.5.7]
\label{II.4.5.7}
O produto tensorial de dois $\sh{O}_X$-módulos amplos é amplo.
\end{corollary}

\begin{corollary}[4.5.8]
\label{II.4.5.8}
Sejam $\sh{L}$ um $\sh{O}_X$-módulo amplo e $\sh{L}'$ um $\sh{O}_X$-módulo invertível;
existe então um inteiro $n_0>0$ tal que $\sh{L}^{\otimes n}\otimes\sh{L}'$ seja amplo e gerado por suas seções sobre $X$ para $n\geq n_0$.
\end{corollary}

\begin{proof}
Com efeito, resulta de \sref{II.4.5.5} que existe um inteiro $m_0$ tal que $\sh{L}^{\otimes m}\otimes\sh{L}'$ seja gerado por suas seções sobre $X$ para todo $m\geq m_0$;
em virtude de \sref{II.4.5.6}, podemos, portanto, tomar $n_0=m_0+1$.
\end{proof}

\begin{remark}[4.5.9]
\label{II.4.5.9}
Seja $P=\HH^1(X,\sh{O}_X^\times)$ o grupo das classes de $\sh{O}_X$-módulos invertíveis \sref[0]{0.5.4.7}, e seja $P^+$ o subconjunto de $P$ formado pelas classes de feixes amplos.
Suponhamos que $P^+$ seja \emph{não vazio}.
Resulta então de \sref{II.4.5.7} e \sref{II.4.5.8} que
\[
  P^+ + P^+ \subset P^+
  \quad\text{e}\quad
  P^+ - P^+ = P
\]
Em outras palavras, $P^+\cup\{0\}$ é o conjunto dos elementos \emph{positivos} de $P$ para uma estrutura de \emph{pré-ordem} sobre $P$ compatível com sua estrutura de grupo, e que é até mesmo \emph{arquimediana}, em virtude de \sref{II.4.5.8}.
É por isso que às vezes se diz ``feixe positivo'' em vez de feixe amplo, e ``feixe negativo'' para o inverso de um feixe amplo (mas não adotaremos essa terminologia).
\end{remark}

\begin{proposition}[4.5.10]
\label{II.4.5.10}
Sejam $Y$ um esquema afim, $q:X\to Y$ um morfismo separado quase-compacto e $\sh{L}$ um $\sh{O}_X$-módulo invertível.
\begin{enumerate}
  \item[\rm{(i)}] Se $\sh{L}$ é muito amplo relativamente a $q$, então $\sh{L}$ é amplo.
  \item[\rm{(ii)}] Suponhamos, além disso, que o morfismo $q$ seja \emph{de tipo finito}.
    Para que $\sh{L}$ seja amplo, é necessário e suficiente que possua uma das seguintes propriedades equivalentes:
    \begin{enumerate}
      \item[\rm{(e)}] Existe $n_0>0$ tal que, para todo inteiro $n\geq n_0$, $\sh{L}^{\otimes n}$ seja muito amplo relativamente a $q$.
      \item[\rm{(e')}] Existe $n>0$ tal que $\sh{L}^{\otimes n}$ seja muito amplo relativamente a $q$.
    \end{enumerate}
\end{enumerate}
\end{proposition}

\begin{proof}
A primeira afirmação resulta da definição \sref{II.4.4.2} de um $\sh{O}_X$-módulo muito amplo:
se $A$ é o anel de $Y$, existem um $A$-módulo $E$ e um homomorfismo sobrejetivo
\[
  \psi:q^*((\bb{S}(E))^{\supertilde})\to\bigoplus_{n\geq0}\sh{L}^{\otimes n}
\]
tais que $i=r_{\sh{L},\psi}$ seja uma imersão definida em toda parte $X\to P=\bb{P}(\widetilde{E})$ e $\sh{L}=i^*(\sh{O}_P(1))$;
como os $D_+(f)$, para $f$ homogêneo em $(\bb{S}(E))_+$, formam uma base da topologia de $P$, e como $i^{-1}(D_+(f))=X_{\psi^\flat(f)}$ por \sref{II.3.7.3.1}, vemos que a condição~(a) de \sref{II.4.5.2} é verificada, e, portanto, $\sh{L}$ é amplo.

Demonstremos agora que, se $q$ é de tipo finito e $\sh{L}$ é amplo, então a condição~(e) é verificada.
Em primeiro lugar, resulta do critério~(b) de \sref{II.4.5.2} e de \sref{II.4.4.1}[(i)] que existe
\oldpage[II]{87}
um inteiro $k_0>0$ tal que $\sh{L}^{\otimes k_0}$ seja muito amplo relativamente a $q$.
Por outro lado, em virtude de \sref{II.4.5.5}, existe um inteiro $m_0$ tal que, para todo $m\geq m_0$, $\sh{L}^{\otimes m}$ seja gerado por suas seções sobre $X$.
Ponhamos $n_0=k_0+m_0$;
se $n\geq n_0$, podemos escrever $n=k_0+m$, com $m\geq m_0$, donde $\sh{L}^{\otimes n}=\sh{L}^{\otimes k_0}\otimes\sh{L}^{\otimes m}$.
Como $\sh{L}^{\otimes m}$ é gerado por suas seções sobre $X$, resulta de \sref{II.4.4.8} e \sref{II.3.4.7} que $\sh{L}^{\otimes n}$ é muito amplo relativamente a $q$.
Por fim, é claro que (e) implica (e'), e (e') implica que $\sh{L}$ é amplo em virtude de (i) e de \sref{II.4.5.6}[(i)].

  \begin{env}[4.5.10.1]
  \label{II.4.5.10.1}
  \emph{[Demonstração do Lema~\sref{II.4.4.10.1}].}
  Ponhamos $\sh{E}(n)=\sh{E}\otimes\sh{K}^{\otimes n}$;
  como $g$ é separado \sref{II.4.4.2}, o raciocínio de \sref{II.3.4.8} aplica-se e reduz a afirmação a demonstrar que o homomorfismo canônico $g^*(g_*(\sh{E}(n)))\to\sh{E}(n)$ seja sobrejetivo para todo $n$ suficientemente grande.
  Além disso, como $Z$ é quase-compacto, o raciocínio de \sref{II.3.4.6} reduz-nos a demonstrar a afirmação quando $Z$ é afim.
  Ora, $\sh{K}$ é então amplo em virtude de \sref{II.4.5.10}[(i)], e a conclusão resulta de \sref{II.4.5.5}[(d')].
  \end{env}
\end{proof}

\begin{corollary}[4.5.11]
\label{II.4.5.11}
Sejam $Y$ um esquema afim, $q:X\to Y$ um morfismo separado de tipo finito, $\sh{L}$ um $\sh{O}_X$-módulo amplo e $\sh{L}'$ um $\sh{O}_X$-módulo invertível.
Existe um inteiro $n_0$ tal que, para todo $n\geq n_0$, $\sh{L}^{\otimes n}\otimes\sh{L}'$ seja muito amplo relativamente a $q$.
\end{corollary}

\begin{proof}
Existe $m_0$ tal que, para $m\geq m_0$, $\sh{L}^{\otimes m}\otimes\sh{L}'$ seja gerado por suas seções sobre $X$ \sref{II.4.5.8};
existe também $k_0$ tal que $\sh{L}^{\otimes k}$ seja muito amplo relativamente a $q$ para $k\geq k_0$.
Então $\sh{L}^{\otimes(k+m_0)}\otimes\sh{L}'$ é muito amplo se $k\geq k_0$ (\sref{II.4.4.8} e \sref{II.3.4.7}).
\end{proof}

\begin{remark}[4.5.12]
\label{II.4.5.12}
Ignora-se se a hipótese de que um $\sh{O}_X$-módulo $\sh{L}$ seja tal que $\sh{L}^{\otimes n}$ seja muito amplo (relativamente a $q$) implica a mesma conclusão para $\sh{L}^{\otimes(n+1)}$.
\end{remark}

\begin{proposition}[4.5.13]
\label{II.4.5.13}
Sejam $X$ um pré-esquema quase-compacto, $Z$ um subpré-esquema fechado de $X$ definido por um feixe quase-coerente \emph{nilpotente} $\sh{J}$ de ideais de $\sh{O}_X$, e $j$ a injeção canônica $Z\to X$.
Para que um $\sh{O}_X$-módulo invertível $\sh{L}$ seja amplo, é necessário e suficiente que $\sh{L}'=j^*(\sh{L})$ seja um $\sh{O}_Z$-módulo amplo.
\end{proposition}

\begin{proof}
A condição é \emph{necessária}.
Com efeito, para toda seção $f$ de $\sh{L}^{\otimes n}$ sobre $X$, seja $f'$ sua imagem canônica $f\otimes1$, que é uma seção de $\sh{L}'^{\otimes n}=\sh{L}^{\otimes n}\otimes_{\sh{O}_X}(\sh{O}_X/\sh{J})$ sobre o espaço $Z$ (que é idêntico a $X$);
é claro que $X_f=Z_{f'}$, e, portanto, o critério~(a) de \sref{II.4.5.2} mostra que $\sh{L}'$ é amplo.

Para ver que a condição é \emph{suficiente}, observemos primeiro que podemos reduzir-nos ao caso em que $\sh{J}^2=0$, considerando a sequência (finita) de pré-esquemas $X_k=(X,\sh{O}_X/\sh{J}^{k+1})$, cada um dos quais é um subpré-esquema fechado do seguinte, definido por um feixe de ideais de quadrado nulo.
Além disso, $X$ é um esquema, pois $X_\red$ o é por hipótese (\sref{II.4.5.3} e \sref[I]{I.5.5.1}).
O critério~(a) de \sref{II.4.5.2} mostra que bastará demonstrar o
  \begin{lemma}[4.5.13.1]
  \label{II.4.5.13.1}
  Sob as hipóteses de \sref{II.4.5.13}, suponhamos, além disso, que $\sh{J}$ seja de quadrado nulo;
  sendo $\sh{L}$ um $\sh{O}_X$-módulo invertível, seja $g$ uma seção de $\sh{L}'^{\otimes n}$ sobre $Z$ tal que $Z_g$ seja afim.
  Então existe um inteiro $m>0$ tal que $g^{\otimes m}$ seja a imagem canônica de uma seção $f$ de $\sh{L}^{\otimes nm}$ sobre $X$.
  \end{lemma}

  Temos a sequência exata de $\sh{O}_X$-módulos
  \[
    0\to\sh{J}(n)\to\sh{O}_X(n)=\sh{L}^{\otimes n}\to\sh{O}_Z(n)=\sh{L}'^{\otimes n}\to0
  \]
\oldpage[II]{88}
  pois $\sh{F}(n)$ é um funtor exato em $\sh{F}$;
  daí a sequência exata de cohomologia
  \[
    0\to\Gamma(X,\sh{J}(n))\to\Gamma(X,\sh{L}^{\otimes n})\to\Gamma(X,\sh{L}'^{\otimes n})\xrightarrow{\partial}\HH^1(X,\sh{J}(n))
  \]
  que associa, em particular, a $g$ um elemento $\partial g\in\HH^1(X,\sh{J}(n))$.
  Observemos que, como $\sh{J}^2=0$, $\sh{J}$ pode ser considerado um $\sh{O}_Z$-módulo quase-coerente e, para todo $k$, $\sh{L}'^{\otimes k}\otimes_{\sh{O}_Z}\sh{J}(n)=\sh{J}(n+k)$;
  para toda seção $s\in\Gamma(X,\sh{L}'^{\otimes k})$, a multiplicação tensorial por $s$ é, portanto, um homomorfismo $\sh{J}(n)\to\sh{J}(n+k)$ de $\sh{O}_Z$-módulos, que induz um homomorfismo $\HH^i(X,\sh{J}(n))\xrightarrow{s}\HH^i(X,\sh{J}(n+k))$ de grupos de cohomologia.

  Isso posto, veremos que
  \[
  \label{II.4.5.13.2}
    g^{\otimes m}\otimes\partial g=0
  \tag{4.5.13.2}
  \]
  para $m>0$ suficientemente grande.
  Com efeito, $Z_g$ é um \emph{aberto afim} de $Z$, e, portanto, $\HH^1(Z_g,\sh{J}(n))=0$ quando $\sh{J}(n)$ é considerado um \emph{$\sh{O}_Z$-módulo} \sref[I]{I.5.1.9.2}.
  Em particular, se pusermos $g'=g|Z_g$ e considerarmos sua imagem pela aplicação $\partial:\Gamma(Z_g,\sh{L}'^{\otimes n})\to\HH^1(Z_g,\sh{J}(n))$, teremos $\partial g'=0$.
  Para explicitar essa relação, observemos que, em dimensão~$1$, a cohomologia de um feixe de grupos abelianos coincide com sua cohomologia de \v{C}ech (G, II, 5.9);
  para calcular $\partial g$, devemos, portanto, considerar um recobrimento aberto $(U_\alpha)$ suficientemente fino de $X$, que podemos supor \emph{finito} e formado por abertos afins, tomar, para cada $\alpha$, uma seção $g_\alpha\in\Gamma(U_\alpha,\sh{L}^{\otimes n})$ cuja imagem canônica em $\Gamma(U_\alpha,\sh{L}'^{\otimes n})$ seja $g|U_\alpha$, e considerar a classe do cociclo $(g_{\alpha|\beta}-g_{\beta|\alpha})$, onde $g_{\alpha|\beta}$ denota a restrição de $g_\alpha$ a $U_\alpha\cap U_\beta$ (esse cociclo toma valores em $\sh{J}(n)$).
  Podemos, além disso, supor que $\partial g'$ seja calculado da mesma maneira, por meio do recobrimento formado pelos $U_\alpha\cap Z_g$ e das restrições $g_\alpha|(U_\alpha\cap Z_g)$ (substituindo, se necessário, $(U_\alpha)$ por um recobrimento mais fino);
  a relação $\partial g'=0$ implica então que existe, para cada $\alpha$, uma seção $h_\alpha\in\Gamma(U_\alpha\cap Z_g,\sh{J}(n))$ tal que $(g_{\alpha|\beta}-g_{\beta|\alpha})|(U_\alpha\cap U_\beta\cap Z_g)=h_{\alpha|\beta}-h_{\beta|\alpha}$, onde $h_{\alpha|\beta}$ denota a restrição de $h_\alpha$ a $U_\alpha\cap U_\beta\cap Z_g$ (G, II, 5.11).
  Existe então um inteiro $m>0$ tal que $g^{\otimes m}\otimes h_\alpha$ seja a restrição a $U_\alpha\cap Z_g$ de uma seção $t_\alpha\in\Gamma(U_\alpha,\sh{J}(n+nm))$ para todo $\alpha$ \sref[I]{I.9.3.1};
  temos, portanto, $g^{\otimes m}\otimes(g_{\alpha|\beta}-g_{\beta|\alpha})=t_{\alpha|\beta}-t_{\beta|\alpha}$ para todo par de índices, o que demonstra \sref{II.4.5.13.2}.

  Observemos, por outro lado, que, se $s\in\Gamma(X,\sh{O}_Z(p))$ e $t\in\Gamma(X,\sh{O}_Z(q))$, temos, no grupo $\HH^1(X,\sh{J}(p+q))$,
  \[
  \label{II.4.5.13.3}
    \partial(s\otimes t)=(\partial s)\otimes t+s\otimes(\partial t).
  \tag{4.5.13.3}
  \]

  Com efeito, podemos novamente calcular os dois membros considerando um recobrimento aberto $(U_\alpha)$ de $X$ e, para cada $\alpha$, uma seção $s_\alpha\in\Gamma(U_\alpha,\sh{O}_X(p))$ (resp. $t_\alpha\in\Gamma(U_\alpha,\sh{O}_X(q))$) cuja imagem canônica em $\Gamma(U_\alpha,\sh{O}_Z(p))$ (resp. $\Gamma(U_\alpha,\sh{O}_Z(q))$) seja $s|U_\alpha$ (resp. $t|U_\alpha$);
  a relação~\sref{II.4.5.13.3} resulta então das relações
  \[
    (s_{\alpha|\beta}\otimes t_{\alpha|\beta})-(s_{\beta|\alpha}\otimes t_{\beta|\alpha})
    =(s_{\alpha|\beta}-s_{\beta|\alpha})\otimes t_{\alpha|\beta}+s_{\beta|\alpha}\otimes(t_{\alpha|\beta}-t_{\beta|\alpha})
  \]
  com as mesmas notações acima.
  Por indução sobre $k$, temos, portanto,
  \[
  \label{II.4.5.13.4}
    \partial(g^{\otimes k})=(kg^{\otimes(k-1)})\otimes(\partial g)
  \tag{4.5.13.4}
  \]

\oldpage[II]{89}
  e concluímos de \sref{II.4.5.13.2} e \sref{II.4.5.13.4} que $\partial(g^{\otimes(m+1)})=0$;
  portanto, $g^{\otimes(m+1)}$ é a imagem canônica de uma seção $f$ de $\sh{L}^{\otimes n(m+1)}$ sobre $X$, o que conclui a demonstração de \sref{II.4.5.13}.
\end{proof}

\begin{corollary}[4.5.14]
\label{II.4.5.14}
Sejam $X$ um esquema noetheriano e $j$ a injeção canônica $X_\red\to X$.
Para que um $\sh{O}_X$-módulo invertível $\sh{L}$ seja amplo, é necessário e suficiente que $j^*(\sh{L})$ seja um $\sh{O}_{X_\red}$-módulo amplo.
\end{corollary}

\begin{proof}
Isso resulta de \sref[I]{I.6.1.6}.
\end{proof}


\subsection{Feixes relativamente amplos}
\label{subsection:II.4.6}

\begin{definition}[4.6.1]
\label{II.4.6.1}
Sejam $f:X\to Y$ um morfismo quase-compacto e $\sh{L}$ um $\sh{O}_X$-módulo invertível.
Diz-se que $\sh{L}$ é \emph{amplo relativamente a $f$}, ou \emph{relativamente a $Y$}, ou \emph{$f$-amplo}, ou \emph{$Y$-amplo} (ou até simplesmente \emph{amplo}, se não puder resultar nenhuma confusão com a noção definida em \sref{II.4.5.3}) se existir um recobrimento $(U_\alpha)$ de $Y$ formado por abertos afins tal que, pondo $X_\alpha=f^{-1}(U_\alpha)$, $\sh{L}|X_\alpha$ seja um $\sh{O}_{X_\alpha}$-módulo amplo para todo $\alpha$.
\end{definition}

A existência de um $\sh{O}_X$-módulo $f$-amplo implica necessariamente que $f$ é \emph{separado} (\sref{II.4.5.3} e \sref[I]{I.5.5.5}).

\begin{proposition}[4.6.2]
\label{II.4.6.2}
Sejam $f:X\to Y$ um morfismo quase-compacto e $\sh{L}$ um $\sh{O}_X$-módulo invertível.
Se $\sh{L}$ é muito amplo relativamente a $f$, então é amplo relativamente a $f$.
\end{proposition}

\begin{proof}
Isso resulta do caráter local (sobre $Y$) da noção de feixe muito amplo \sref{II.4.4.5}, da definição \sref{II.4.6.1} e do critério~\sref{II.4.5.10}[(i)].
\end{proof}

\begin{proposition}[4.6.3]
\label{II.4.6.3}
Sejam $f:X\to Y$ um morfismo quase-compacto, $\sh{L}$ um $\sh{O}_X$-módulo invertível e $\sh{S}$ a $\sh{O}_Y$-álgebra graduada $\bigoplus_{n\geq0}f_*(\sh{L}^{\otimes n})$.
As condições seguintes são equivalentes:
\begin{enumerate}
  \item[\rm{(a)}] $\sh{L}$ é $f$-amplo.
  \item[\rm{(b)}] $\sh{S}$ é quase-coerente, e o homomorfismo canônico $\sigma:f^*(\sh{S})\to\bigoplus_{n\geq0}\sh{L}^{\otimes n}$ \sref[0]{0.4.4.3} é tal que o $Y$-morfismo $r_{\sh{L},\sigma}:G(\sigma)\to\Proj(\sh{S})=P$ esteja definido em toda parte e seja uma imersão aberta dominante.
  \item[\rm{(b')}] O morfismo $f$ é separado, e o $Y$-morfismo $r_{\sh{L},\sigma}$ está definido em toda parte e é um homeomorfismo do espaço subjacente a $X$ sobre um subespaço de $\Proj(\sh{S})$.
\end{enumerate}

Além disso, quando essas condições são satisfeitas, para todo $n\in\bb{Z}$, o homomorfismo canônico
\[
\label{II.4.6.3.1}
  r_{\sh{L},\sigma}^*(\sh{O}_P(n))\to\sh{L}^{\otimes n}
\tag{4.6.3.1}
\]
definido em \sref{II.3.7.9.1} é um isomorfismo.

Por fim, para todo $\sh{O}_X$-módulo quase-coerente $\sh{F}$, se pusermos $\sh{M}=\bigoplus_{n\geq0}f_*(\sh{F}\otimes\sh{L}^{\otimes n})$, o homomorfismo canônico
\[
\label{II.4.6.3.2}
  r_{\sh{L},\sigma}^*(\widetilde{\sh{M}})\to\sh{F}
\tag{4.6.3.2}
\]
definido em \sref{II.3.7.9.2} é um isomorfismo.
\end{proposition}

\begin{proof}
Observamos que (a) implica que $f$ é separado e, portanto, que $\sh{S}$ é quase-coerente \sref[I]{I.9.2.2}[(a)].
Como o fato de $r_{\sh{L},\sigma}$ ser uma imersão definida em toda parte é de caráter local sobre $Y$, para demonstrar que (a) implica (b), podemos supor que $Y$ seja afim e $\sh{L}$ amplo;
\oldpage[II]{90}
a afirmação resulta então de \sref{II.4.5.2}[(b)].
É claro que (b) implica (b');
finalmente, para demonstrar que (b') implica (a), basta considerar um recobrimento de $Y$ por abertos afins $(U_\alpha)$ e aplicar o critério~\sref{II.4.5.2}[(b')] a cada feixe $\sh{L}|f^{-1}(U_\alpha)$.

Para as duas últimas afirmações, utiliza-se o fato de que $\sigma^\flat$ é aqui a identidade e a explicitação dos homomorfismos \sref{II.3.7.9.1} e \sref{II.3.7.9.2};
resulta imediatamente que \sref{II.4.6.3.1} é um isomorfismo.
Quanto a \sref{II.4.6.3.2}, podemos reduzir-nos ao caso em que $Y$ é afim, portanto $\sh{L}$ é amplo;
é claro que o homomorfismo \sref{II.4.6.3.2} é injetivo, e o critério~\sref{II.4.5.2}[(c)] mostra que ele é sobrejetivo, de onde a conclusão.
\end{proof}

\begin{corollary}[4.6.4]
\label{II.4.6.4}
Seja $(U_\alpha)$ um recobrimento aberto de $Y$.
Para que $\sh{L}$ seja amplo relativamente a $Y$, é necessário e suficiente que $\sh{L}|f^{-1}(U_\alpha)$ seja amplo relativamente a $U_\alpha$ para todo $\alpha$.
\end{corollary}

\begin{proof}
Com efeito, a condição~(b) é local sobre $Y$.
\end{proof}

\begin{corollary}[4.6.5]
\label{II.4.6.5}
Seja $\sh{K}$ um $\sh{O}_Y$-módulo invertível.
Para que $\sh{L}$ seja $Y$-amplo, é necessário e suficiente que $\sh{L}\otimes f^*(\sh{K})$ o seja.
\end{corollary}

\begin{proof}
É uma consequência evidente de \sref{II.4.6.4}, tomando os $U_\alpha$ tais que $\sh{K}|U_\alpha$ seja isomorfo a $\sh{O}_Y|U_\alpha$ para todo $\alpha$.
\end{proof}

\begin{corollary}[4.6.6]
\label{II.4.6.6}
Suponhamos que $Y$ seja afim;
para que $\sh{L}$ seja $Y$-amplo, é necessário e suficiente que $\sh{L}$ seja amplo.
\end{corollary}

\begin{proof}
É uma consequência imediata da definição \sref{II.4.6.1} e dos critérios \sref{II.4.6.3}[(b)] e \sref{II.4.5.2}[(b)], pois aqui $\Proj(\sh{S})=\Proj(\Gamma(Y,\sh{S}))$ por definição.
\end{proof}

\begin{corollary}[4.6.7]
\label{II.4.6.7}
Seja $f:X\to Y$ um morfismo quase-compacto.
Suponhamos que existam um $\sh{O}_Y$-módulo quase-coerente $\sh{E}$ e um $Y$-morfismo $g:X\to P=\bb{P}(\sh{E})$ que seja um homeomorfismo do espaço subjacente a $X$ sobre um subespaço de $P$;
então $\sh{L}=g^*(\sh{O}_P(1))$ é $Y$-amplo.
\end{corollary}

\begin{proof}
Podemos supor que $Y$ seja afim;
o corolário resulta então do critério~\sref{II.4.5.2}[(a)], da fórmula~\sref{II.3.7.3.1} e de \sref{II.4.2.3}.
\end{proof}

\begin{proposition}[4.6.8]
\label{II.4.6.8}
Sejam $X$ um esquema quase-compacto ou um pré-esquema cujo espaço subjacente seja noetheriano, e $f:X\to Y$ um morfismo separado quase-compacto.
Para que um $\sh{O}_X$-módulo invertível $\sh{L}$ seja $f$-amplo, é necessário e suficiente que uma das condições equivalentes seguintes seja verificada:
\begin{enumerate}
  \item[\rm{(c)}] Para todo $\sh{O}_X$-módulo quase-coerente $\sh{F}$ de tipo finito, existe um inteiro $n_0>0$ tal que, para todo $n\geq n_0$, o homomorfismo canônico $\sigma:f^*(f_*(\sh{F}\otimes\sh{L}^{\otimes n}))\to\sh{F}\otimes\sh{L}^{\otimes n}$ seja sobrejetivo.
  \item[\rm{(c')}] Para todo feixe quase-coerente $\sh{J}$ de ideais de $\sh{O}_X$ de tipo finito, existe um inteiro $n>0$ tal que o homomorfismo canônico $\sigma:f^*(f_*(\sh{J}\otimes\sh{L}^{\otimes n}))\to\sh{J}\otimes\sh{L}^{\otimes n}$ seja sobrejetivo.
\end{enumerate}
\end{proposition}

\begin{proof}
Como $X$ é quase-compacto, $f(X)$ também o é, e existe, portanto, um recobrimento finito $(U_i)$ de $f(X)$ por abertos afins de $Y$.
Para demonstrar a condição~(c) quando $\sh{L}$ é $f$-amplo, podemos substituir $Y$ pelos $U_i$ e $X$ pelos $f^{-1}(U_i)$, pois, se obtivermos para cada $i$ um inteiro $n_i$ tal que (c) seja verificada (para $U_i$, $f^{-1}(U_i)$ e $\sh{L}|f^{-1}(U_i)$) desde que $n\geq n_i$, bastará tomar como $n_0$ o maior dos $n_i$ para obter (c) para $Y$, $X$ e $\sh{L}$.
Ora, quando $Y$ é afim, a condição~(c) resulta de \sref{II.4.5.5}[(d)], levando em conta \sref{II.4.6.6}.
É trivial que (c) implica (c').
Finalmente, para demonstrar que (c') implica que $\sh{L}$ é $f$-amplo, podemos novamente nos limitar ao caso em que $Y$ é afim:
com efeito, todo feixe quase-coerente $\sh{J}_i$ de ideais de $\sh{O}_X|f^{-1}(U_i)$ de tipo finito é a restrição de um feixe quase-coerente de ideais de $\sh{O}_X$ de tipo finito \sref[I]{I.9.4.7}, e a hipótese~(c') implica que
\oldpage[II]{91}
$\sh{J}_i\otimes(\sh{L}^{\otimes n}|f^{-1}(U_i))$ seja gerado por suas seções (levando em conta \sref[I]{I.9.2.2} e \sref{II.3.4.7});
basta, portanto, aplicar o critério~\sref{II.4.5.5}[(d'')].
\end{proof}

\begin{proposition}[4.6.9]
\label{II.4.6.9}
Sejam $f:X\to Y$ um morfismo quase-compacto e $\sh{L}$ um $\sh{O}_X$-módulo invertível.
\begin{enumerate}
  \item[\rm{(i)}] Seja $n>0$ um inteiro.
    Para que $\sh{L}$ seja $f$-amplo, é necessário e suficiente que $\sh{L}^{\otimes n}$ seja $f$-amplo.
  \item[\rm{(ii)}] Seja $\sh{L}'$ um $\sh{O}_X$-módulo invertível, e suponhamos que exista um inteiro $n>0$ tal que o homomorfismo canônico $\sigma:f^*(f_*(\sh{L}'^{\otimes n}))\to\sh{L}'^{\otimes n}$ seja sobrejetivo.
    Então, se $\sh{L}$ é $f$-amplo, o mesmo vale para $\sh{L}\otimes\sh{L}'$.
\end{enumerate}
\end{proposition}

\begin{proof}
Reduzimo-nos imediatamente ao caso em que $Y$ é afim, e a proposição é então uma consequência imediata de \sref{II.4.5.6}.
\end{proof}

\begin{corollary}[4.6.10]
\label{II.4.6.10}
O produto tensorial de dois $\sh{O}_X$-módulos $f$-amplos é $f$-amplo.
\end{corollary}

\begin{proposition}[4.6.11]
\label{II.4.6.11}
Sejam $Y$ um pré-esquema quase-compacto, $f:X\to Y$ um morfismo de tipo finito e $\sh{L}$ um $\sh{O}_X$-módulo invertível.
Para que $\sh{L}$ seja $f$-amplo, é necessário e suficiente que possua uma das propriedades equivalentes seguintes:
\begin{enumerate}
  \item[\rm{(d)}] Existe $n_0>0$ tal que, para todo inteiro $n\geq n_0$, $\sh{L}^{\otimes n}$ seja muito amplo relativamente a $f$.
  \item[\rm{(d')}] Existe $n>0$ tal que $\sh{L}^{\otimes n}$ seja muito amplo relativamente a $f$.
\end{enumerate}
\end{proposition}

\begin{proof}
Se $\sh{L}$ é amplo relativamente a $f$, existe um recobrimento \emph{finito} $(U_i)$ de $Y$ por abertos afins tal que os $\sh{L}|f^{-1}(U_i)$ sejam amplos.
Concluímos então de \sref{II.4.5.10} que existe um inteiro $n_0$ tal que $\sh{L}^{\otimes n}|f^{-1}(U_i)$ seja muito amplo relativamente a $f^{-1}(U_i)\to U_i$ para todo $n\geq n_0$ e todo $i$, e, portanto, $\sh{L}^{\otimes n}$ seja muito amplo relativamente a $f$ \sref{II.4.4.5}.
Reciprocamente, (d') implica já que $\sh{L}^{\otimes n}$ é $f$-amplo \sref{II.4.6.2}, e, portanto, o mesmo vale para $\sh{L}$ \sref{II.4.6.9}[(i)].
\end{proof}

\begin{corollary}[4.6.12]
\label{II.4.6.12}
Sejam $Y$ um pré-esquema quase-compacto, $f:X\to Y$ um morfismo de tipo finito, e $\sh{L}$ e $\sh{L}'$ dois $\sh{O}_X$-módulos invertíveis.
Se $\sh{L}$ é $f$-amplo, existe $n_0$ tal que $\sh{L}^{\otimes n}\otimes\sh{L}'$ seja muito amplo relativamente a $f$ para todo $n\geq n_0$.
\end{corollary}

\begin{proof}
Raciocina-se como em \sref{II.4.6.11}, utilizando um recobrimento aberto afim finito de $Y$ e \sref{II.4.5.11}.
\end{proof}

\begin{proposition}[4.6.13]
\label{II.4.6.13}
\begin{enumerate}
  \item[\rm{(i)}] Para todo pré-esquema $Y$, todo $\sh{O}_Y$-módulo invertível $\sh{L}$ é amplo relativamente ao morfismo identidade $1_Y$.
  \item[\rm{(i \emph{bis})}] Sejam $f:X\to Y$ um morfismo quase-compacto e $j:X'\to X$ um morfismo quase-compacto que seja um homeomorfismo do espaço subjacente a $X'$ sobre um subespaço de $X$.
    Se $\sh{L}$ é um $\sh{O}_X$-módulo amplo relativamente a $f$, então $j^*(\sh{L})$ é amplo relativamente a $f\circ j$.
  \item[\rm{(ii)}] Sejam $Z$ um pré-esquema quase-compacto, $f:X\to Y$ e $g:Y\to Z$ dois morfismos quase-compactos, $\sh{L}$ um $\sh{O}_X$-módulo amplo relativamente a $f$, e $\sh{K}$ um $\sh{O}_Y$-módulo amplo relativamente a $g$.
    Então existe um inteiro $n_0>0$ tal que $\sh{L}\otimes f^*(\sh{K}^{\otimes n})$ seja amplo relativamente a $g\circ f$ para todo $n\geq n_0$.
  \item[\rm{(iii)}] Sejam $f:X\to Y$ um morfismo quase-compacto e $g:Y'\to Y$ um morfismo, e ponhamos $X'=X_{(Y')}$.
    Se $\sh{L}$ é um $\sh{O}_X$-módulo amplo relativamente a $f$, então $\sh{L}'=\sh{L}\otimes_Y\sh{O}_{Y'}$ é um $\sh{O}_{X'}$-módulo amplo relativamente a $f_{(Y')}$.
  \item[\rm{(iv)}] Sejam $f_i:X_i\to Y_i$ ($i=1,2$) dois $S$-morfismos quase-compactos.
    Se $\sh{L}_i$ é um $\sh{O}_{X_i}$-módulo amplo relativamente a $f_i$ ($i=1,2$), então $\sh{L}_1\otimes_S\sh{L}_2$ é amplo relativamente a $f_1\times_S f_2$.
  \item[\rm{(v)}] Sejam $f:X\to Y$ e $g:Y\to Z$ dois morfismos tais que $g\circ f$ seja quase-compacto.
    Se um
\oldpage[II]{92}
    $\sh{O}_X$-módulo $\sh{L}$ é amplo relativamente a $g\circ f$, e se $g$ é separado ou o espaço subjacente a $X$ é localmente noetheriano, então $\sh{L}$ é amplo relativamente a $f$.
  \item[\rm{(vi)}] Sejam $f:X\to Y$ um morfismo quase-compacto e $j$ a injeção canônica $X_\red\to X$.
    Se $\sh{L}$ é um $\sh{O}_X$-módulo amplo relativamente a $f$, então $j^*(\sh{L})$ é amplo relativamente a $f_\red$.
\end{enumerate}
\end{proposition}

\begin{proof}
Observemos, antes de tudo, que (v) e (vi) se deduzem de (i), (i~\emph{bis}) e (iv) pelo mesmo raciocínio que em \sref{II.4.4.10}, utilizando \sref{II.4.6.4} em vez de \sref{II.4.4.5};
deixamos ao leitor o detalhe do raciocínio.
A afirmação~(i) é consequência trivial de \sref{II.4.4.10}[(i)] e de \sref{II.4.6.2}.
Para demonstrar (i~\emph{bis}), (iii) e (iv), utilizaremos o lema seguinte:
  \begin{lemma}[4.6.13.1]
  \label{II.4.6.13.1}
  \begin{enumerate}
    \item[\rm{(i)}] Sejam $u:Z\to S$ um morfismo, $\sh{L}$ um $\sh{O}_S$-módulo invertível, $s$ uma seção de $\sh{L}$ sobre $S$, e $s'$ a seção de $u^*(\sh{L})=\sh{L}'$ sobre $Z$ que lhe corresponde canonicamente.
      Então $Z_{s'}=u^{-1}(S_s)$.
    \item[\rm{(ii)}] Sejam $Z$ e $Z'$ dois $S$-pré-esquemas, $p$ e $p'$ as projeções de $T=Z\times_S Z'$, $\sh{L}$ (resp. $\sh{L}'$) um $\sh{O}_Z$-módulo invertível (resp. um $\sh{O}_{Z'}$-módulo invertível), $t$ (resp. $t'$) uma seção de $\sh{L}$ (resp. $\sh{L}'$) sobre $Z$ (resp. $Z'$), e $s$ (resp. $s'$) a seção de $p^*(\sh{L})$ (resp. $p'^*(\sh{L}')$) sobre $Z\times_S Z'$ que lhe corresponde.
      Então $T_{s\otimes s'}=Z_t\times_S Z'_{t'}$.
  \end{enumerate}
  \end{lemma}

  \begin{proof}
  Resulta das definições que podemos reduzir-nos ao caso em que todos os pré-esquemas considerados são afins.
  Além disso, em (i), podemos supor que $\sh{L}=\sh{O}_S$;
  a afirmação~(i) resulta então imediatamente de \sref[I]{I.1.2.2.2}.
  Do mesmo modo, em (ii), podemos nos limitar ao caso em que $\sh{L}=\sh{O}_Z$ e $\sh{L}'=\sh{O}_{Z'}$, e então a afirmação se reduz ao lema~\sref{II.4.3.2.4}.
  \end{proof}

Demonstremos agora (i~\emph{bis}).
Podemos supor que $Y$ seja afim \sref{II.4.6.4}, portanto $\sh{L}$ é amplo \sref{II.4.6.6};
quando $s$ percorre a união dos conjuntos $\Gamma(X,\sh{L}^{\otimes n})$ ($n>0$), os $X_s$ formam uma base da topologia de $X$ \sref{II.4.5.2}[(a)], e, portanto, por hipótese, os $j^{-1}(X_s)$ formam uma base da topologia de $X'$;
concluímos então, pelo lema~\sref{II.4.6.13.1}[(i)] e por \sref{II.4.5.2}[(a)], que $j^*(\sh{L})$ é amplo.

Demonstremos em seguida (iii).
Podemos novamente supor que $Y$ e $Y'$ sejam afins \sref{II.4.6.4}, de onde resulta que a projeção $h:X'\to X$ é afim \sref{II.1.5.5}.
Como $\sh{L}$ é amplo \sref{II.4.6.6}, quando $s$ percorre o conjunto das seções dos $\sh{L}^{\otimes n}$ ($n>0$) sobre $X$ tais que $X_s$ seja afim, os $X_s$ recobrem $X$ \sref{II.4.5.2}[(a')], e, portanto, os $h^{-1}(X_s)$ são afins \sref{II.1.2.5} e recobrem $X'$;
resulta novamente do lema~\sref{II.4.6.13.1}[(i)] e de \sref{II.4.5.2}[(a')] que $\sh{L}'$ é amplo, sendo o morfismo $f_{(Y')}$ quase-compacto \sref[I]{I.6.6.4}[(iii)].

Para demonstrar (iv), observemos primeiro que $f_1\times_S f_2$ é quase-compacto \sref[I]{I.6.6.4}[(iv)].
Podemos, além disso, supor que $S$, $Y_1$ e $Y_2$ sejam afins (\sref{II.4.6.4} e \sref[I]{I.3.2.7}), portanto $\sh{L}_i$ é amplo ($i=1,2$) \sref{II.4.6.6}.
Os abertos $(X_1)_{s_1}\times_S(X_2)_{s_2}$ formam um recobrimento de $X_1\times_S X_2$, quando $s_i$ percorre as seções de $\sh{L}_i^{\otimes n_i}$ tais que $(X_i)_{s_i}$ seja afim \sref{II.4.5.2}[(a')].
Além disso, substituindo $s_1$ e $s_2$ por potências convenientes, o que não altera os $(X_i)_{s_i}$, podemos sempre supor que $n_1=n_2$.
Deduzimos então de \sref{II.4.6.13.1}[(ii)] e de \sref{II.4.5.2}[(a')] que $\sh{L}_1\otimes_S\sh{L}_2$ é amplo, de onde a afirmação, pois $Y_1\times_S Y_2$ é afim \sref{II.4.6.6}.

Resta demonstrar (ii).
Pelo mesmo raciocínio que em \sref{II.4.4.10}, mas utilizando aqui \sref{II.4.6.4}, podemos nos limitar ao caso em que $Z$ é afim.
Como $\sh{K}$ é então amplo e $Y$ é quase-compacto, existe um número finito de seções $s_i\in\Gamma(Y,\sh{K}^{\otimes k_i})$ tais que os $Y_{s_i}$
\oldpage[II]{93}
sejam \emph{afins} e recubram $Y$ \sref{II.4.5.2}[(a')];
substituindo os $s_i$ por potências convenientes, podemos, além disso, supor que todos os $k_i$ sejam iguais a um mesmo inteiro $k$.
Sejam $s'_i$ as seções de $f^*(\sh{K}^{\otimes k})$ sobre $X$ que correspondem canonicamente aos $s_i$, de modo que os $X_{s'_i}=f^{-1}(Y_{s_i})$ \sref{II.4.6.13.1}[(i)] recubram $X$.
Como $\sh{L}|X_{s'_i}$ é amplo (\sref{II.4.6.4} e \sref{II.4.6.6}), existe, para cada $i$, um número finito de seções $t_{ij}\in\Gamma(X,\sh{L}^{\otimes n_{ij}})$ tais que os $X_{t_{ij}}$ sejam afins, estejam contidos em $X_{s'_i}$ e recubram $X_{s'_i}$ \sref{II.4.5.2}[(a')];
podemos, além disso, supor que todos os $n_{ij}$ sejam iguais a um mesmo número $n$.
Nessas condições, $X$ é separado e quase-compacto, de modo que existem um inteiro $m>0$ e, para todo $(i,j)$, uma seção
\[
  u_{ij}\in\Gamma(X,\sh{L}^{\otimes n}\otimes_X f^*(\sh{K}^{\otimes mk}))
\]
tal que $t_{ij}\otimes s'^{\otimes m}_i$ seja a restrição de $u_{ij}$ a $X_{s'_i}$ \sref[I]{I.9.3.1};
além disso, $X_{u_{ij}}=X_{t_{ij}}$, logo os $X_{u_{ij}}$ são afins e recobrem $X$.
Podemos, além disso, supor que $m$ seja da forma $nr$;
se pusermos $n_0=rk$, vemos então por \sref{II.4.5.2}[(a')] que $\sh{L}\otimes_{\sh{O}_X}f^*(\sh{K}^{\otimes n_0})$ é amplo.
Além disso, existe $h_0>0$ tal que $\sh{K}^{\otimes h}$ seja gerado por suas seções sobre $Y$ desde que $h\geq h_0$ \sref{II.4.5.5};
\emph{a fortiori}, $f^*(\sh{K}^{\otimes h})$ é gerado por suas seções sobre $X$ para $h\geq h_0$, pela definição das imagens inversas (\sref[0]{0.3.7.1} e \sref{II.4.4.1}).
Concluímos que $\sh{L}\otimes f^*(\sh{K}^{\otimes n_0+h})$ é amplo desde que $h\geq h_0$ \sref{II.4.5.6}, o que termina a demonstração.
\end{proof}

\begin{remark}[4.6.14]
\label{II.4.6.14}
Nas condições de (ii), não se deve crer que $\sh{L}\otimes f^*(\sh{K})$ seja amplo para $g\circ f$;
com efeito, como $\sh{L}\otimes f^*(\sh{K}^{-1})$ também é amplo para $f$ \sref{II.4.6.5}, concluiríamos que $\sh{L}$ seria amplo para $g\circ f$;
tomando, em particular, como $g$ o morfismo identidade, \emph{todo} $\sh{O}_X$-módulo invertível seria amplo para $f$, o que não ocorre em geral (ver \sref{II.5.1.6}, \sref{II.5.3.4}[(i)] e \sref{II.5.3.1}).
\end{remark}

\begin{proposition}[4.6.15]
\label{II.4.6.15}
Sejam $f:X\to Y$ um morfismo quase-compacto, $\sh{J}$ um feixe quase-coerente de ideais de $\sh{O}_X$, localmente nilpotente, $Z$ o subpré-esquema fechado de $X$ definido por $\sh{J}$, e $j:Z\to X$ a injeção canônica.
Para que um $\sh{O}_X$-módulo invertível $\sh{L}$ seja amplo para $f$, é necessário e suficiente que $j^*(\sh{L})$ seja amplo para $f\circ j$.
\end{proposition}

\begin{proof}
Com efeito, como a questão é local sobre $Y$ \sref{II.4.6.4}, podemos supor que $Y$ seja afim;
como $X$ é então quase-compacto, podemos supor que $\sh{J}$ seja nilpotente.
Levando em conta \sref{II.4.6.6}, a proposição não é então outra coisa senão \sref{II.4.5.13}.
\end{proof}

\begin{corollary}[4.6.16]
\label{II.4.6.16}
Sejam $X$ um pré-esquema localmente noetheriano e $f:X\to Y$ um morfismo quase-compacto.
Para que um $\sh{O}_X$-módulo invertível $\sh{L}$ seja amplo para $f$, é necessário e suficiente que sua imagem inversa $\sh{L}'$ pela injeção canônica $X_\red\to X$ seja ampla para $f_\red$.
\end{corollary}

\begin{proof}
Já vimos que a condição é necessária \sref{II.4.6.13}[(vi)];
reciprocamente, se ela é verificada, podemos nos limitar, para provar que $\sh{L}$ é amplo para $f$, ao caso em que $Y$ é afim \sref{II.4.6.4};
então $Y_\red$ também é afim, portanto $\sh{L}'$ é amplo \sref{II.4.6.6}, e o mesmo vale para $\sh{L}$ em virtude de \sref{II.4.5.13}, pois $X$ é então noetheriano e $X_\red$ é um subpré-esquema fechado de $X$ definido por um feixe quase-coerente nilpotente de ideais \sref[I]{I.6.1.6}.
\end{proof}

\begin{proposition}[4.6.17]
\label{II.4.6.17}
Sendo as notações e hipóteses as de \sref{II.4.4.11}, para que $\sh{L}''$ seja amplo relativamente a $f''$, é necessário e suficiente que $\sh{L}$ seja amplo relativamente a $f$ e $\sh{L}'$ seja amplo relativamente a $f'$.
\end{proposition}
\oldpage[II]{94}
\begin{proof}
A necessidade da condição resulta de \sref{II.4.6.13}[(i~\emph{bis})].
Para ver que a condição é suficiente, podemos nos limitar ao caso em que $Y$ é afim, e então o fato de que $\sh{L}''$ é amplo resulta do critério~\sref{II.4.5.2}[(a)] aplicado a $\sh{L}$, $\sh{L}'$ e $\sh{L}''$, observando que uma seção de $\sh{L}$ sobre $X$ se prolonga (por $0$) a uma seção de $\sh{L}''$ sobre $X''$.
\end{proof}

\begin{proposition}[4.6.18]
\label{II.4.6.18}
Sejam $Y$ um pré-esquema quase-compacto, $\sh{S}$ uma $\sh{O}_Y$-álgebra graduada quase-coerente de tipo finito, $X=\Proj(\sh{S})$, e $f:X\to Y$ o morfismo estrutural.
Então $f$ é de tipo finito, e existe um inteiro $d>0$ tal que $\sh{O}_X(d)$ seja invertível e $f$-amplo.
\end{proposition}

\begin{proof}
Em virtude de \sref{II.3.1.10}, existe um inteiro $d>0$ tal que $\sh{S}^{(d)}$ seja gerada por $\sh{S}_d$.
Sabemos que, no isomorfismo canônico entre $X$ e $X^{(d)}=\Proj(\sh{S}^{(d)})$, $\sh{O}_X(d)$ identifica-se com $\sh{O}_{X^{(d)}}(1)$ \sref{II.3.2.9}[(ii)].
Vemos, portanto, que nos reduzimos ao caso em que $\sh{S}$ é gerada por $\sh{S}_1$;
a proposição resulta então de \sref{II.4.4.3} e \sref{II.4.6.2} (levando em conta que $f$ é um morfismo de tipo finito \sref{II.3.4.1}).
\end{proof}

\section{Morfismos quase-afins; morfismos quase-projetivos; morfismos próprios; morfismos projetivos}
\label{section:II.5}

\subsection{Morfismos quase-afins}
\label{subsection:II.5.1}

\begin{definition}[5.1.1]
\label{II.5.1.1}
Chama-se esquema quase-afim todo esquema isomorfo a um subesquema induzido em um aberto quase-compacto de um esquema afim.
Diz-se que um morfismo $f:X\to Y$ é quase-afim, ou ainda que $X$ (considerado como $Y$-pré-esquema por meio de $f$) é um $Y$-esquema quase-afim, se existe um recobrimento $(U_\alpha)$ de $Y$ por abertos afins tal que os $f^{-1}(U_\alpha)$ sejam esquemas quase-afins.
\end{definition}

É claro que um morfismo quase-afim é \emph{separado} (\sref[I]{I.5.5.5} e \sref[I]{I.5.5.8}) e \emph{quase-compacto} \sref[I]{I.6.6.1};
todo morfismo afim é evidentemente quase-afim.

Recordemos que, para todo pré-esquema $X$, se pusermos $A=\Gamma(X,\sh{O}_X)$, o homomorfismo identidade $A\to A=\Gamma(X,\sh{O}_X)$ define um morfismo $X\to\Spec(A)$, dito \emph{canônico} \sref[I]{I.2.2.4};
este não é senão o morfismo canônico definido em \sref{II.4.5.1} para o caso particular em que $\sh{L}=\sh{O}_X$, se nos lembrarmos de que $\Proj(A[T])$ se identifica canonicamente com $\Spec(A)$ \sref{II.3.1.7}.

\begin{proposition}[5.1.2]
\label{II.5.1.2}
Sejam $X$ um esquema quase-compacto ou um pré-esquema cujo espaço subjacente seja noetheriano, e $A$ o anel $\Gamma(X,\sh{O}_X)$.
As condições seguintes são equivalentes:
\begin{enumerate}
  \item[{\rm(a)}] $X$ é um esquema quase-afim.
  \item[{\rm(b)}] O morfismo canônico $u:X\to\Spec(A)$ é uma imersão aberta.
  \item[{\rm(b')}] O morfismo canônico $u:X\to\Spec(A)$ é um homeomorfismo de $X$ sobre um subespaço do espaço subjacente a $\Spec(A)$.
  \item[{\rm(c)}] O $\sh{O}_X$-módulo $\sh{O}_X$ é muito amplo relativamente a $u$ \sref{II.4.4.2}.
  \item[{\rm(c')}] O $\sh{O}_X$-módulo $\sh{O}_X$ é amplo \sref{II.4.5.1}.
  \item[{\rm(d)}] Quando $f$ percorre $A$, os $X_f$ formam uma base da topologia de $X$.
  \item[{\rm(d')}] Quando $f$ percorre $A$, aqueles $X_f$ que são afins formam um recobrimento de $X$.
\oldpage[II]{95}
  \item[{\rm(e)}] Todo $\sh{O}_X$-módulo quase-coerente é gerado por suas seções sobre $X$.
  \item[{\rm(e')}] Todo feixe quase-coerente de ideais de $\sh{O}_X$ de tipo finito é gerado por suas seções sobre $X$.
\end{enumerate}
\end{proposition}

\begin{proof}
É claro que (b) implica (a), e (a) implica (c) em virtude do critério~\sref{II.4.4.4}[(b)] aplicado ao morfismo identidade (levando em conta a observação que precede o enunciado);
por outro lado, (c) implica (c$'$) \sref{II.4.5.10}[(i)], e (c$'$), (b) e (b$'$) são equivalentes em virtude de \sref{II.4.5.2}[(b)] e \sref{II.4.5.2}[(b$'$)].
Por fim, (c$'$) equivale a cada uma das condições (d), (d$'$), (e) e (e$'$), em virtude de \sref{II.4.5.2}[(a)], \sref{II.4.5.2}[(a$'$)], \sref{II.4.5.2}[(c)] e \sref{II.4.5.5}[(d$''$)].
\end{proof}

Observaremos, além disso, que, com as notações anteriores, os $X_f$ que são afins formam uma \emph{base} da topologia de $X$, e que o morfismo canônico $u$ é \emph{dominante} \sref{II.4.5.2}.

\begin{corollary}[5.1.3]
\label{II.5.1.3}
Seja $X$ um pré-esquema quase-compacto.
Se existir um morfismo $v:X\to Y$ de $X$ em um esquema afim $Y$ que seja um homeomorfismo de $X$ sobre um subespaço aberto de $Y$, então $X$ é quase-afim.
\end{corollary}

\begin{proof}
Existe uma família $(g_\alpha)$ de seções de $\sh{O}_Y$ sobre $Y$ tal que os $D(g_\alpha)$ formam uma base da topologia de $v(X)$;
se $v=(\psi,\theta)$ e se pusermos $f_\alpha=\theta(g_\alpha)$, teremos $X_{f_\alpha}=\psi^{-1}(D(g_\alpha))$ \sref[I]{I.2.2.4.1}; portanto, os $X_{f_\alpha}$ formam uma base da topologia de $X$, e o critério~\sref{II.5.1.2}[(d)] é satisfeito.
\end{proof}

\begin{corollary}[5.1.4]
\label{II.5.1.4}
Se $X$ é um esquema quase-afim, então \emph{todo} $\sh{O}_X$-módulo invertível é muito amplo (relativamente ao morfismo canônico) e, \emph{a fortiori}, amplo.
\end{corollary}

\begin{proof}
Com efeito, tal módulo $\sh{L}$ é gerado por suas seções sobre $X$ \sref{II.5.1.2}[(e)]; portanto, $\sh{L}\otimes\sh{O}_X=\sh{L}$ é muito amplo \sref{II.4.4.8}.
\end{proof}

\begin{corollary}[5.1.5]
\label{II.5.1.5}
Seja $X$ um pré-esquema quase-compacto.
Se existir um $\sh{O}_X$-módulo invertível $\sh{L}$ tal que $\sh{L}$ e $\sh{L}^{-1}$ sejam amplos, então $X$ é um esquema quase-afim.
\end{corollary}

\begin{proof}
Com efeito, $\sh{O}_X=\sh{L}\otimes\sh{L}^{-1}$ é então amplo \sref{II.4.5.7}.
\end{proof}

\begin{proposition}[5.1.6]
\label{II.5.1.6}
Seja $f:X\to Y$ um morfismo quase-compacto.
As condições seguintes são equivalentes:
\begin{enumerate}
  \item[{\rm(a)}] O morfismo $f$ é quase-afim.
  \item[{\rm(b)}] A $\sh{O}_Y$-álgebra $f_*(\sh{O}_X)=\sh{A}(X)$ é quase-coerente, e o morfismo canônico $X\to\Spec(\sh{A}(X))$ correspondente ao homomorfismo identidade $\sh{A}(X)\to\sh{A}(X)$ \sref{II.1.2.7} é uma imersão aberta.
  \item[{\rm(b')}] A $\sh{O}_Y$-álgebra $\sh{A}(X)$ é quase-coerente, e o morfismo canônico $X\to\Spec(\sh{A}(X))$ é um homeomorfismo de $X$ sobre um subespaço de $\Spec(\sh{A}(X))$.
  \item[{\rm(c)}] O $\sh{O}_X$-módulo $\sh{O}_X$ é muito amplo para $f$.
  \item[{\rm(c')}] O $\sh{O}_X$-módulo $\sh{O}_X$ é amplo para $f$.
  \item[{\rm(d)}] O morfismo $f$ é separado e, para todo $\sh{O}_X$-módulo quase-coerente $\sh{F}$, o homomorfismo canônico $\sigma:f^*(f_*(\sh{F}))\to\sh{F}$ \sref[0]{0.4.4.3} é sobrejetivo.
\end{enumerate}

Além disso, quando $f$ é quase-afim, todo $\sh{O}_X$-módulo invertível $\sh{L}$ é muito amplo relativamente a $f$.
\end{proposition}

\begin{proof}
A equivalência entre (a) e (c$'$) decorre do caráter local sobre $Y$ da $f$-amplitude \sref{II.4.6.4}, da definição~\sref{II.5.1.1} e do critério~\sref{II.5.1.2}[(c$'$)].
As outras propriedades são locais sobre $Y$

\oldpage[II]{96}
e decorrem, portanto, imediatamente de \sref{II.5.1.2} e \sref{II.5.1.4}, levando em conta que $f_*(\sh{F})$ é quase-coerente quando $f$ é separado \sref[I]{I.9.2.2}[(a)].
\end{proof}

\begin{corollary}[5.1.7]
\label{II.5.1.7}
Seja $f:X\to Y$ um morfismo quase-afim.
Para todo aberto $U$ de $Y$, a restrição $f^{-1}(U)\to U$ de $f$ é quase-afim.
\end{corollary}

\begin{corollary}[5.1.8]
\label{II.5.1.8}
Sejam $Y$ um esquema afim e $f:X\to Y$ um morfismo quase-compacto.
Para que $f$ seja quase-afim, é necessário e suficiente que $X$ seja um esquema quase-afim.
\end{corollary}

\begin{proof}
É uma consequência imediata de \sref{II.5.1.6} e \sref{II.4.6.6}.
\end{proof}

\begin{corollary}[5.1.9]
\label{II.5.1.9}
Sejam $Y$ um esquema quase-compacto ou um pré-esquema cujo espaço subjacente seja noetheriano, e $f:X\to Y$ um morfismo de \emph{tipo finito}.
Se $f$ é quase-afim, existe uma sub-$\sh{O}_Y$-álgebra quase-coerente $\sh{B}$ de $\sh{A}(X)=f_*(\sh{O}_X)$, de \emph{tipo finito} \sref[I]{I.9.6.2}, tal que o morfismo $X\to\Spec(\sh{B})$ correspondente à injeção canônica $\sh{B}\to\sh{A}(X)$ \sref{II.1.2.7} seja uma imersão.
Além disso, toda sub-$\sh{O}_Y$-álgebra quase-coerente de tipo finito $\sh{B}'$ de $\sh{A}(X)$ que contenha $\sh{B}$ possui a mesma propriedade.
\end{corollary}

\begin{proof}
Com efeito, $\sh{A}(X)$ é o limite indutivo de suas sub-$\sh{O}_Y$-álgebras quase-coerentes de tipo finito \sref[I]{I.9.6.5};
o resultado é então um caso particular de \sref{II.3.8.4}, levando em conta a identificação de $\Spec(\sh{A}(X))$ com $\Proj(\sh{A}(X)[T])$ \sref{II.3.1.7}.
\end{proof}

\begin{proposition}[5.1.10]
\label{II.5.1.10}
\medskip\noindent
\begin{enumerate}
  \item[{\rm(i)}] Todo morfismo quase-compacto $X\to Y$ que seja um homeomorfismo do espaço subjacente a $X$ sobre um subespaço do espaço subjacente a $Y$ (em particular, toda imersão fechada) é quase-afim.
  \item[{\rm(ii)}] O composto de dois morfismos quase-afins é quase-afim.
  \item[{\rm(iii)}] Se $f:X\to Y$ é um $S$-morfismo quase-afim, então $f_{(S')}:X_{(S')}\to Y_{(S')}$ é um morfismo quase-afim para toda extensão $S'\to S$ do pré-esquema de base.
  \item[{\rm(iv)}] Se $f:X\to Y$ e $g:X'\to Y'$ são dois $S$-morfismos quase-afins, então $f\times_S g$ é quase-afim.
  \item[{\rm(v)}] Se $f:X\to Y$ e $g:Y\to Z$ são dois morfismos tais que $g\circ f$ seja quase-afim, e se $g$ é separado ou o espaço subjacente a $X$ é localmente noetheriano, então $f$ é quase-afim.
  \item[{\rm(vi)}] Se $f$ é um morfismo quase-afim, então $f_\red$ também o é.
\end{enumerate}
\end{proposition}

\begin{proof}
Levando em conta o critério~\sref{II.5.1.6}[(c$'$)], as afirmações (i), (iii), (iv), (v) e (vi) decorrem imediatamente, respectivamente, de \sref{II.4.6.13}[(i~\emph{bis})], \sref{II.4.6.13}[(iii)], \sref{II.4.6.13}[(iv)], \sref{II.4.6.13}[(v)] e \sref{II.4.6.13}[(vi)].
Para demonstrar (ii), podemos nos reduzir ao caso em que $Z$ é afim, e então a afirmação decorre diretamente de \sref{II.4.6.13}[(ii)], aplicado a $\sh{L}=\sh{O}_X$ e $\sh{K}=\sh{O}_Y$.
\end{proof}

\begin{remark}[5.1.11]
\label{II.5.1.11}
Sejam $f:X\to Y$ e $g:Y\to Z$ dois morfismos tais que $X\times_Z Y$ seja localmente noetheriano.
Então a imersão do gráfico $\Gamma_f:X\to X\times_Z Y$ é quase-afim, pois é quase-compacta \sref[I]{I.6.3.5}, e \sref[I]{I.5.5.12} mostra que, em (v), a conclusão continua válida se suprimirmos a hipótese de que $g$ é separado.
\end{remark}

\begin{proposition}[5.1.12]
\label{II.5.1.12}
Sejam $f:X\to Y$ um morfismo quase-compacto e $g:X'\to X$ um morfismo quase-afim.
Se $\sh{L}$ é um $\sh{O}_X$-módulo amplo para $f$, então $g^*(\sh{L})$ é um $\sh{O}_{X'}$-módulo amplo para $f\circ g$.
\end{proposition}

\begin{proof}
Como $\sh{O}_{X'}$ é muito amplo para $g$ e a questão é local sobre $Y$ \sref{II.4.6.4}, decorre de \sref{II.4.6.13}[(ii)] que existe (para $Y$ afim) um inteiro $n$ tal que
\[
  g^*(\sh{L}^{\otimes n})=(g^*(\sh{L}))^{\otimes n}
\]
seja amplo para $f\circ g$; logo, $g^*(\sh{L})$ é amplo para $f\circ g$ \sref{II.4.6.9}.
\end{proof}

\oldpage[II]{97}
\subsection{O critério de Serre}
\label{subsection:II.5.2}

\begin{theorem}[5.2.1]
\label{II.5.2.1}
\emph{(Critério de Serre).}
Seja $X$ um esquema quase-compacto ou um pré-esquema cujo espaço subjacente seja noetheriano.
As condições seguintes são equivalentes:
\begin{enumerate}
  \item[{\rm(a)}] $X$ é um esquema afim.
  \item[{\rm(b)}] Existe uma família de elementos $f_\alpha\in A=\Gamma(X,\sh{O}_X)$ tal que os $X_{f_\alpha}$ sejam afins e o ideal gerado pelos $f_\alpha$ em $A$ seja igual a $A$.
  \item[{\rm(c)}] O funtor $\Gamma(X,\sh{F})$ é exato em $\sh{F}$ na categoria dos $\sh{O}_X$-módulos quase-coerentes; em outras palavras, se
    \[
      0\to\sh{F}'\to\sh{F}\to\sh{F}''\to0
      \tag{*}
    \]
    é uma sequência exata de $\sh{O}_X$-módulos quase-coerentes, então a sequência
    \[
      0\to\Gamma(X,\sh{F}')\to\Gamma(X,\sh{F})\to\Gamma(X,\sh{F}'')\to0
    \]
    é exata.
  \item[{\rm(c')}] A propriedade~{\rm(c)} é válida para toda sequência exata {\rm($*$)} de $\sh{O}_X$-módulos quase-coerentes tal que $\sh{F}$ seja isomorfo a um $\sh{O}_X$-submódulo de $\sh{O}_X^n$ para algum $n$ finito.
  \item[{\rm(d)}] $\HH^1(X,\sh{F})=0$ para todo $\sh{O}_X$-módulo quase-coerente $\sh{F}$.
  \item[{\rm(d$'$)}] $\HH^1(X,\sh{J})=0$ para todo feixe quase-coerente de ideais $\sh{J}$ de $\sh{O}_X$.
\end{enumerate}
\end{theorem}

\begin{proof}
É evidente que (a) implica (b);
por outro lado, (b) implica que os $X_{f_\alpha}$ recobrem $X$, pois, por hipótese, a seção $1$ é combinação linear dos $f_\alpha$, e os $D(f_\alpha)$ recobrem então $\Spec(A)$.
A última afirmação de \sref{II.4.5.2} implica, portanto, que $X\to\Spec(A)$ é um isomorfismo.

Sabemos que (a) implica (c) \sref[I]{I.1.3.11}, e (c) implica trivialmente (c$'$).
Demonstremos que (c$'$) implica (b).
Primeiramente, (c$'$) implica que, para todo ponto \emph{fechado} $x\in X$ e toda vizinhança aberta $U$ de $x$, existe $f\in A$ tal que $x\in X_f\subset U$.
Com efeito, seja $\sh{J}$ (resp. $\sh{J}'$) o feixe quase-coerente de ideais de $\sh{O}_X$ que define o subpré-esquema fechado reduzido de $X$ cujo espaço subjacente é $X\setmin U$ (resp. $(X\setmin U)\cup\{x\}$) \sref[I]{I.5.2.1};
é claro que $\sh{J}'\subset\sh{J}$ e que $\sh{J}''=\sh{J}/\sh{J}'$ é um $\sh{O}_X$-módulo quase-coerente cujo suporte é $\{x\}$ e tal que $\sh{J}_x''=\kres(x)$.
A hipótese~(c$'$), aplicada à sequência exata $0\to\sh{J}'\to\sh{J}\to\sh{J}''\to0$, mostra que $\Gamma(X,\sh{J})\to\Gamma(X,\sh{J}'')$ é sobrejetivo.
A seção de $\sh{J}''$ cujo germe em $x$ é $1_x$ é, portanto, a imagem de uma seção $f\in\Gamma(X,\sh{J})\subset\Gamma(X,\sh{O}_X)$, e temos, por definição, $f(x)=1_x$ e $f(y)=0$ em $X\setmin U$, o que estabelece nossa afirmação.
Além disso, se $U$ é afim, então $X_f$ também o é \sref[I]{I.1.3.6}; logo, a união dos $X_f$ que são afins ($f\in A$) é um aberto $Z$ que contém \emph{todos os pontos fechados} de $X$;
como $X$ é um espaço de Kolmogoroff quase-compacto, temos necessariamente $Z=X$ \sref[0]{0.2.1.3}.
Como $X$ é quase-compacto, existe um número \emph{finito} de elementos $f_i\in A$ ($1\leq i\leq n$) tais que os $X_{f_i}$ sejam afins e recubram $X$.
Consideremos então o homomorfismo $\sh{O}_X^n\to\sh{O}_X$ definido pelas seções $f_i$ \sref[0]{0.5.1.1};
como, para todo $x\in X$, pelo menos um dos $(f_i)_x$ é invertível, esse homomorfismo é \emph{sobrejetivo}, e temos, portanto, uma sequência exata $0\to\sh{R}\to\sh{O}_X^n\to\sh{O}_X\to0$, em que $\sh{R}$ é um $\sh{O}_X$-submódulo quase-coerente de $\sh{O}_X^n$.
Resulta então
\oldpage[II]{98}
de (c$'$) que o homomorfismo correspondente $\Gamma(X,\sh{O}_X^n)\to\Gamma(X,\sh{O}_X)$ é sobrejetivo, o que demonstra (b).

Por fim, (a) implica (d) \sref[I]{I.5.1.9.2}, e (d) implica trivialmente (d$'$).
Resta mostrar que (d$'$) implica (c$'$).
Ora, se $\sh{F}'$ é um $\sh{O}_X$-submódulo quase-coerente de $\sh{O}_X^n$, a filtração $0\subset\sh{O}_X\subset\sh{O}_X^2\subset\ldots\subset\sh{O}_X^n$ define em $\sh{F}'$ uma filtração formada pelos $\sh{F}'_k=\sh{F}'\cap\sh{O}_X^k$ ($0\leq k\leq n$), que são $\sh{O}_X$-módulos quase-coerentes \sref[I]{I.4.1.1}, e $\sh{F}'_{k+1}/\sh{F}'_k$ é isomorfo a um $\sh{O}_X$-submódulo quase-coerente de $\sh{O}_X^{k+1}/\sh{O}_X^k=\sh{O}_X$, isto é, a um feixe quase-coerente de ideais de $\sh{O}_X$.
A hipótese~(d$'$) implica, portanto, que $\HH^1(X,\sh{F}'_{k+1}/\sh{F}'_k)=0$;
a sequência exata de cohomologia $\HH^1(X,\sh{F}'_k)\to\HH^1(X,\sh{F}'_{k+1})\to\HH^1(X,\sh{F}'_{k+1}/\sh{F}'_k)=0$ permite então demonstrar, por indução sobre $k$, que $\HH^1(X,\sh{F}'_k)=0$ para todo $k$.
\end{proof}

\begin{remark}[5.2.1.1]
\label{II.5.2.1.1}
Quando $X$ é um pré-esquema \emph{noetheriano}, pode-se substituir ``quase-coerente'' por ``coerente'' nos enunciados de (c$'$) e (d$'$).
Com efeito, na demonstração de que (c$'$) implica (b), $\sh{J}$ e $\sh{J}'$ são então feixes \emph{coerentes} de ideais e, por outro lado, todo submódulo quase-coerente de um módulo coerente é coerente \sref[I]{I.6.1.1};
daí a conclusão.
\end{remark}

\begin{corollary}[5.2.2]
\label{II.5.2.2}
Seja $f:X\to Y$ um morfismo separado quase-compacto.
As condições seguintes são equivalentes:
\begin{enumerate}
  \item[{\rm(a)}] $f$ é um morfismo afim.
  \item[{\rm(b)}] O funtor $f_*$ é exato na categoria dos $\sh{O}_X$-módulos quase-coerentes.
  \item[{\rm(c)}] Para todo $\sh{O}_X$-módulo quase-coerente $\sh{F}$, temos $\RR^1f_*(\sh{F})=0$.
  \item[{\rm(c')}] Para todo feixe quase-coerente de ideais $\sh{J}$ de $\sh{O}_X$, temos $\RR^1f_*(\sh{J})=0$.
\end{enumerate}
\end{corollary}

\begin{proof}
Todas essas condições são locais sobre $Y$, pela definição do funtor $\RR^1f_*$ (T,~3.7.3); portanto, podemos supor que $Y$ seja afim.
Se $f$ é afim, então $X$ é afim, e a propriedade~(b) não é outra senão \sref[I]{I.1.6.4}.
Reciprocamente, mostremos que (b) implica (a):
para todo $\sh{O}_X$-módulo quase-coerente $\sh{F}$, $f_*(\sh{F})$ é um $\sh{O}_Y$-módulo quase-coerente \sref[I]{I.9.2.2}[(a)].
Por hipótese, o funtor $f_*$ é exato em $\sh{F}$, e o funtor $\Gamma(Y,\sh{G})$ é exato em $\sh{G}$ (na categoria dos $\sh{O}_Y$-módulos quase-coerentes), pois $Y$ é afim \sref[I]{I.1.3.11};
logo, $\Gamma(Y,f_*(\sh{F}))=\Gamma(X,\sh{F})$ é exato em $\sh{F}$, o que demonstra nossa afirmação em virtude de \sref{II.5.2.1}[(c)].

Se $f$ é afim, $f^{-1}(U)$ é afim para todo aberto afim $U$ de $Y$ \sref{II.1.3.2}; portanto, $\HH^1(f^{-1}(U),\sh{F})=0$ \sref{II.5.2.1}[(d)], o que, por definição, implica $\RR^1f_*(\sh{F})=0$.
Suponhamos, por fim, que a condição~(c$'$) seja satisfeita;
a sequência exata dos termos de baixo grau da sequência espectral de Leray (G,~II,~4.17.1 e I,~4.5.1) fornece, em particular, a sequência exata
\[
  0\to\HH^1(Y,f_*(\sh{J}))\to\HH^1(X,\sh{J})\to\HH^0(Y,\RR^1f_*(\sh{J})).
\]
Como $Y$ é afim e $f_*(\sh{J})$ é quase-coerente \sref[I]{I.9.2.2}[(a)], temos $\HH^1(Y,f_*(\sh{J}))=0$ \sref{II.5.2.1};
a hipótese~(c$'$) implica, portanto, $\HH^1(X,\sh{J})=0$, e concluímos por \sref{II.5.2.1} que $X$ é um esquema afim.
\end{proof}

\begin{corollary}[5.2.3]
\label{II.5.2.3}
Se $f:X\to Y$ é um morfismo afim, então, para todo $\sh{O}_X$-módulo quase-coerente $\sh{F}$, o homomorfismo canônico $\HH^1(Y,f_*(\sh{F}))\to\HH^1(X,\sh{F})$ é bijetivo.
\end{corollary}

\oldpage[II]{99}
\begin{proof}
Temos a sequência exata
\[
  0\to\HH^1(Y,f_*(\sh{F}))\to\HH^1(X,\sh{F})\to\HH^0(Y,\RR^1f_*(\sh{F}))
\]
dos termos de baixo grau da sequência espectral de Leray, e a conclusão decorre de \sref{II.5.2.2}.
\end{proof}

\begin{remark}[5.2.4]
\label{II.5.2.4}
No Capítulo~III,~\textsection1, provaremos que, se $X$ é afim, então $\HH^i(X,\sh{F})=0$ para todo $i>0$ e todo $\sh{O}_X$-módulo quase-coerente $\sh{F}$.
\end{remark}

\subsection{Morfismos quase-projetivos}
\label{subsection:II.5.3}

\begin{definition}[5.3.1]
\label{II.5.3.1}
Diz-se que um morfismo $f:X\to Y$ é \emph{quase-projetivo}, ou que $X$ (considerado como $Y$-pré-esquema por meio de $f$) é \emph{quase-projetivo sobre $Y$}, ou que $X$ é um \emph{$Y$-esquema quase-projetivo}, se $f$ é de tipo finito e existe um $\sh{O}_X$-módulo invertível $f$-amplo.
\end{definition}

Observaremos que essa noção \emph{não é local sobre $Y$}:
os contraexemplos de Nagata~\cite{II-26} e de Hironaka mostram que, mesmo que $X$ e $Y$ sejam esquemas algébricos não singulares sobre um corpo algebricamente fechado, todo ponto de $Y$ pode ter uma vizinhança afim $U$ tal que $f^{-1}(U)$ seja quase-projetivo sobre $U$, sem que $f$ seja quase-projetivo.

Observaremos que um morfismo quase-projetivo é necessariamente \emph{separado} \sref{II.4.6.1}.
Quando $Y$ é quase-compacto, é equivalente dizer que $f$ é quase-projetivo ou que $f$ é de tipo finito e existe um $\sh{O}_X$-módulo \emph{muito amplo} relativamente a $f$ (\sref{II.4.6.2} e \sref{II.4.6.11}).
Além disso:

\begin{proposition}[5.3.2]
\label{II.5.3.2}
Sejam $Y$ um esquema quase-compacto ou um pré-esquema cujo espaço subjacente seja noetheriano, e $X$ um $Y$-pré-esquema.
As condições seguintes são equivalentes:
\begin{enumerate}
  \item[{\rm(a)}] $X$ é um $Y$-esquema quase-projetivo.
  \item[{\rm(b)}] $X$ é de tipo finito sobre $Y$, e existe um $\sh{O}_Y$-módulo quase-coerente de tipo finito $\sh{E}$ tal que $X$ seja $Y$-isomorfo a um subpré-esquema de $\bb{P}(\sh{E})$.
  \item[{\rm(c)}] $X$ é de tipo finito sobre $Y$, e existe uma $\sh{O}_Y$-álgebra graduada quase-coerente $\sh{S}$ tal que $\sh{S}_1$ seja de tipo finito e gere $\sh{S}$, e tal que $X$ seja $Y$-isomorfo a um subpré-esquema induzido sobre um aberto denso em toda parte de $\Proj(\sh{S})$.
\end{enumerate}
\end{proposition}

\begin{proof}
Isso decorre imediatamente da observação anterior e de \sref{II.4.4.3}, \sref{II.4.4.6} e \sref{II.4.4.7}.
\end{proof}

Observaremos que, quando $Y$ é um pré-esquema \emph{noetheriano}, pode-se, nas condições~(b) e (c) de \sref{II.5.3.2}, suprimir a hipótese de que $X$ seja de tipo finito sobre $Y$, pois ela é então automaticamente satisfeita \sref[I]{I.6.3.5}.

\begin{corollary}[5.3.3]
\label{II.5.3.3}
Seja $Y$ um esquema quase-compacto tal que exista um $\sh{O}_Y$-módulo amplo $\sh{L}$ \sref{II.4.5.3}.
Para que um $Y$-esquema $X$ seja quase-projetivo, é necessário e suficiente que ele seja de tipo finito sobre $Y$ e isomorfo a um $Y$-subesquema de um fibrado projetivo da forma $\bb{P}_Y^r$.
\end{corollary}

\begin{proof}
Com efeito, se $\sh{E}$ é um $\sh{O}_Y$-módulo quase-coerente de tipo finito, $\sh{E}$ é isomorfo a um quociente de um $\sh{O}_Y$-módulo $\sh{L}^{\otimes(-n)}\otimes_{\sh{O}_Y}\sh{O}_Y^k$ \sref{II.4.5.5}; portanto, $\bb{P}(\sh{E})$ é isomorfo a um subesquema fechado de $\bb{P}_Y^{k-1}$ (\sref{II.4.1.2} e \sref{II.4.1.4}).
\end{proof}

\begin{proposition}[5.3.4]
\label{II.5.3.4}
\medskip\noindent
\begin{enumerate}
  \item[{\rm(i)}] Todo morfismo quase-afim de tipo finito (e, em particular, toda imersão quase-compacta ou todo morfismo afim de tipo finito) é quase-projetivo.
  \item[{\rm(ii)}] Se $f:X\to Y$ e $g:Y\to Z$ são quase-projetivos, e se $Z$ é quase-compacto, então $g\circ f$ é quase-projetivo.

\oldpage[II]{100}
  \item[{\rm(iii)}] Se $f:X\to Y$ é um $S$-morfismo quase-projetivo, então $f_{(S')}:X_{(S')}\to Y_{(S')}$ é quase-projetivo para toda extensão $S'\to S$ do pré-esquema de base.
  \item[{\rm(iv)}] Se $f:X\to Y$ e $g:X'\to Y'$ são dois $S$-morfismos quase-projetivos, então $f\times_S g$ é quase-projetivo.
  \item[{\rm(v)}] Se $f:X\to Y$ e $g:Y\to Z$ são dois morfismos tais que $g\circ f$ seja quase-projetivo, e se $g$ é separado ou $X$ localmente noetheriano, então $f$ é quase-projetivo.
  \item[{\rm(vi)}] Se $f$ é um morfismo quase-projetivo, o mesmo vale para $f_\red$.
\end{enumerate}
\end{proposition}

\begin{proof}
(i) decorre de \sref{II.5.1.6} e \sref{II.5.1.10}[(i)].
As outras afirmações são consequências imediatas da definição~\sref{II.5.3.1}, das propriedades dos morfismos de tipo finito \sref[I]{I.6.3.4} e de \sref{II.4.6.13}.
\end{proof}

\begin{remark}[5.3.5]
\label{II.5.3.5}
Observemos que pode ocorrer que $f_\red$ seja quase-projetivo sem que $f$ o seja, mesmo supondo que $Y$ seja o espectro de uma álgebra de posto finito sobre $\bb{C}$ e que $f$ seja próprio.
\end{remark}

\begin{corollary}[5.3.6]
\label{II.5.3.6}
Se $X$ e $X'$ são dois $Y$-esquemas quase-projetivos, então $X\sqcup X'$ é um $Y$-esquema quase-projetivo.
\end{corollary}

\begin{proof}
Isso decorre de \sref{II.4.6.18}.
\end{proof}

\subsection{Morfismos próprios e morfismos universalmente fechados}
\label{subsection:II.5.4}

\begin{definition}[5.4.1]
\label{II.5.4.1}
Diz-se que um morfismo de pré-esquemas $f:X\to Y$ é \emph{próprio} se satisfaz as duas condições seguintes:
\begin{enumerate}
  \item[(a)] $f$ é separado e de tipo finito;
  \item[(b)] para todo pré-esquema $Y'$ e todo morfismo $Y'\to Y$, a projeção $f_{(Y')}:X\times_Y Y'\to Y'$ é um morfismo fechado \sref[I]{I.2.2.6}.
\end{enumerate}

Quando assim ocorre, diz-se também que $X$ (considerado como um $Y$-pré-esquema cujo morfismo estrutural é $f$) é próprio sobre $Y$.
\end{definition}

É imediato que as condições~(a) e (b) são \emph{locais} sobre $Y$.
Para verificar que a imagem de um subconjunto fechado $Z$ de $X\times_Y Y'$ pela projeção $q:X\times_Y Y'\to Y'$ é fechada em $Y'$, basta ver que $q(Z)\cap U'$ é fechado em $U'$ para todo aberto afim $U'$ de $Y'$;
como $q(Z)\cap U'=q(Z\cap q^{-1}(U'))$, e como $q^{-1}(U')$ se identifica com $X\times_Y U'$ \sref[I]{I.4.4.1}, vê-se que, para verificar a condição~(b) da Definição~\sref{II.5.4.1}, podemos nos limitar ao caso em que $Y'$ é um esquema \emph{afim}.
Veremos mais adiante \sref{II.5.6.3} que, se $Y$ é localmente noetheriano, podemos até nos restringir a verificar (b) quando $Y'$ é de tipo finito sobre $Y$.

É claro que todo morfismo próprio é \emph{fechado}.

\begin{proposition}[5.4.2]
\label{II.5.4.2}
\medskip\noindent
\begin{enumerate}
  \item[{\rm(i)}] Uma imersão fechada é um morfismo próprio.
  \item[{\rm(ii)}] O composto de dois morfismos próprios é próprio.
  \item[{\rm(iii)}] Se $X$ e $Y$ são $S$-pré-esquemas e $f:X\to Y$ é um $S$-morfismo próprio, então $f_{(S')}:X_{(S')}\to Y_{(S')}$ é próprio para toda extensão $S'\to S$ do pré-esquema de base.
  \item[{\rm(iv)}] Se $f:X\to Y$ e $g:X'\to Y'$ são dois $S$-morfismos próprios, então $f\times_S g:X\times_S X'\to Y\times_S Y'$ é um $S$-morfismo próprio.
\end{enumerate}
\end{proposition}
\oldpage[II]{101}
\begin{proof}
Basta demonstrar (i), (ii) e (iii) \sref[I]{I.3.5.1}.
Em cada um desses três casos, a verificação da condição~(a) de \sref{II.5.4.1} decorre de resultados anteriores (\sref[I]{I.5.5.1} e \sref[I]{I.6.3.4}); resta verificar a condição~(b).
Isso é imediato no caso (i), pois, se $X\to Y$ é uma imersão fechada, o mesmo vale para $X\times_Y Y'\to Y\times_Y Y'=Y'$ (\sref[I]{I.4.3.2} e \sref[I]{I.3.3.3}).
Para demonstrar (ii), consideremos dois morfismos próprios $X\to Y$ e $Y\to Z$, e um morfismo $Z'\to Z$.
Podemos escrever $X\times_Z Z'=X\times_Y(Y\times_Z Z')$ \sref[I]{I.3.3.9.1}; por conseguinte, a projeção $X\times_Z Z'\to Z'$ se fatora como $X\times_Y(Y\times_Z Z')\to Y\times_Z Z'\to Z'$.
Levando em conta a observação inicial, (ii) decorre, portanto, do fato de que o composto de dois morfismos fechados é fechado.
Finalmente, para todo morfismo $S'\to S$, $X_{(S')}$ se identifica com $X\times_Y Y_{(S')}$ \sref[I]{I.3.3.11}; para todo morfismo $Z\to Y_{(S')}$, podemos escrever
\[
  X_{(S')}\times_{Y_{(S')}}Z=(X\times_Y Y_{(S')})\times_{Y_{(S')}}Z=X\times_Y Z;
\]
como, por hipótese, $X\times_Y Z\to Z$ é fechado, isso demonstra (iii).
\end{proof}

\begin{corollary}[5.4.3]
\label{II.5.4.3}
Sejam $f:X\to Y$ e $g:Y\to Z$ dois morfismos tais que $g\circ f$ seja próprio.
\begin{enumerate}
  \item[{\rm(i)}] Se $g$ é separado, então $f$ é próprio.
  \item[{\rm(ii)}] Se $g$ é separado e de tipo finito, e se $f$ é sobrejetivo, então $g$ é próprio.
\end{enumerate}
\end{corollary}

\begin{proof}
(i) decorre de \sref{II.5.4.2} pelo procedimento geral \sref[I]{I.5.5.12}.
Para demonstrar (ii), basta verificar a condição~(b) da Definição~\sref{II.5.4.1}.
Para todo morfismo $Z'\to Z$, o diagrama
\[
  \xymatrix{
    X\times_Z Z'\ar[r]^{f\times1_{Z'}}\ar[dr]_p &
    Y\times_Z Z'\ar[d]^{p'}\\
    & Z'
  }
\]
(onde $p$ e $p'$ são as projeções) é comutativo \sref[I]{I.3.2.1};
além disso, $f\times1_{Z'}$ é sobrejetivo porque $f$ o é \sref[I]{I.3.5.2}, e $p$ é um morfismo fechado por hipótese.
Todo subconjunto fechado $F$ de $Y\times_Z Z'$ é, então, a imagem por $f\times1_{Z'}$ de um subconjunto fechado $E$ de $X\times_Z Z'$; logo, $p'(F)=p(E)$ é fechado em $Z'$ por hipótese, donde o corolário.
\end{proof}

\begin{corollary}[5.4.4]
\label{II.5.4.4}
Se $X$ é um pré-esquema próprio sobre $Y$, e $\sh{S}$ é uma $\sh{O}_Y$-álgebra quase-coerente, então todo $Y$-morfismo $f:X\to\Proj(\sh{S})$ é próprio (e, \emph{a fortiori}, fechado).
\end{corollary}

\begin{proof}
Com efeito, o morfismo estrutural $p:\Proj(\sh{S})\to Y$ é separado, e $p\circ f$ é próprio por hipótese.
\end{proof}

\begin{corollary}[5.4.5]
\label{II.5.4.5}
Seja $f:X\to Y$ um morfismo separado de tipo finito.
Sejam $(X_i)_{1\leq i\leq n}$ (resp. $(Y_i)_{1\leq i\leq n}$) uma família finita de subpré-esquemas fechados de $X$ (resp. de $Y$), e $j_i$ (resp. $h_i$) a injeção canônica $X_i\to X$ (resp. $Y_i\to Y$).
Suponhamos que o espaço subjacente a $X$ seja a união dos $X_i$ e que, para todo $i$, exista um morfismo $f_i:X_i\to Y_i$ tal que o diagrama
\[
  \xymatrix{
    X_i\ar[r]^{f_i}\ar[d]_{j_i} &
    Y_i\ar[d]^{h_i}\\
    X\ar[r]_f &
    Y
  }
\]
seja comutativo.
Então, para que $f$ seja próprio, é necessário e suficiente que cada um dos $f_i$ o seja.
\end{corollary}
\oldpage[II]{102}
\begin{proof}
Se $f$ é próprio, o mesmo vale para $f\circ j_i$, pois $j_i$ é uma imersão fechada \sref{II.5.4.2};
como $h_i$ é uma imersão fechada, e portanto um morfismo separado, $f_i$ é próprio por \sref{II.5.4.3}.
Reciprocamente, suponhamos que cada um dos $f_i$ seja próprio e consideremos o pré-esquema $Z$, \emph{soma} dos $X_i$; seja $u$ o morfismo $Z\to X$ que se reduz a $j_i$ em cada $X_i$.
A restrição de $f\circ u$ a cada $X_i$ é igual a $f\circ j_i=h_i\circ f_i$ e, portanto, é própria, pois tanto $h_i$ quanto $f_i$ o são \sref{II.5.4.2};
decorre então imediatamente da Definição~\sref{II.5.4.1} que $f\circ u$ é próprio.
Mas, como $u$ é sobrejetivo por hipótese, concluímos que $f$ é próprio por \sref{II.5.4.3}.
\end{proof}

\begin{corollary}[5.4.6]
\label{II.5.4.6}
Seja $f:X\to Y$ um morfismo separado e de tipo finito; para que $f$ seja próprio, é necessário e suficiente que $f_\red:X_\red\to Y_\red$ o seja.
\end{corollary}

\begin{proof}
É um caso particular de \sref{II.5.4.5}, com $n=1$, $X_1=X_\red$ e $Y_1=Y_\red$ \sref[I]{I.5.1.5}.
\end{proof}

\begin{env}[5.4.7]
\label{II.5.4.7}
Se $X$ e $Y$ são pré-esquemas noetherianos e $f:X\to Y$ é um morfismo separado de tipo finito, podemos, para verificar que $f$ seja próprio, reduzir-nos ao caso de morfismos \emph{dominantes} de pré-esquemas \emph{integrais}.
Com efeito, sejam $X_i$ ($1\leq i\leq n$) as componentes irredutíveis, em número finito, de $X$ e consideremos, para cada $i$, o único subpré-esquema fechado reduzido de $X$ que tem $X_i$ como espaço subjacente, o qual continuaremos a denotar por $X_i$ \sref[I]{I.5.2.1}.
Seja, por outro lado, $Y_i$ o único subpré-esquema fechado reduzido de $Y$ que tem $\overline{f(X_i)}$ como espaço subjacente.
Se $g_i$ (resp. $h_i$) é o morfismo de injeção $X_i\to X$ (resp. $Y_i\to Y$), concluímos que $f\circ g_i=h_i\circ f_i$, onde $f_i$ é um morfismo dominante $X_i\to Y_i$ \sref[I]{I.5.2.2};
estamos, então, nas condições de aplicação de \sref{II.5.4.5}, e, para que $f$ seja próprio, é necessário e suficiente que todos os $f_i$ o sejam.
\end{env}

\begin{corollary}[5.4.8]
\label{II.5.4.8}
Sejam $X$ e $Y$ dois $S$-pré-esquemas separados de tipo finito sobre $S$, e seja $f:X\to Y$ um $S$-morfismo.
Para que $f$ seja próprio, é necessário e suficiente que, para todo $S$-pré-esquema $S'$, o morfismo $f\times_S 1_{S'}:X\times_S S'\to Y\times_S S'$ seja fechado.
\end{corollary}

\begin{proof}
Observemos primeiro que, se $g:X\to S$ e $h:Y\to S$ são os morfismos estruturais, temos, por definição, $g=h\circ f$; portanto, $f$ é separado e de tipo finito (\sref[I]{I.5.5.1} e \sref[I]{I.6.3.4}).
Se $f$ é próprio, o mesmo vale para $f\times_S 1_{S'}$ \sref{II.5.4.2}; \emph{a fortiori}, $f\times_S 1_{S'}$ é fechado.
Reciprocamente, suponhamos satisfeita a condição do enunciado e seja $Y'$ um $Y$-pré-esquema;
$Y'$ pode também ser considerado como um $S$-pré-esquema e, como $Y\to S$ é separado, $X\times_Y Y'$ se identifica com um subpré-esquema fechado de $X\times_S Y'$ \sref[I]{I.5.4.2}.
No diagrama comutativo
\[
  \xymatrix{
    X\times_Y Y'\ar[r]^{f\times1_{Y'}}\ar[d] &
    Y\times_Y Y'=Y'\ar[d]\\
    X\times_S Y'\ar[r]_{f\times_S1_{Y'}} &
    Y\times_S Y',
  }
\]
as setas verticais são imersões fechadas; decorre imediatamente que, se $f\times_S1_{Y'}$ é um morfismo fechado, o mesmo vale para $f\times_Y1_{Y'}$.
\end{proof}

\begin{remark}[5.4.9]
\label{II.5.4.9}
Diremos que um morfismo $f:X\to Y$ é \emph{universalmente fechado} se satisfaz a condição~(b) da Definição~\sref{II.5.4.1}.
O leitor observará que,
\oldpage[II]{103}
nas proposições \sref{II.5.4.2} a \sref{II.5.4.8}, podemos substituir em toda parte ``próprio'' por ``universalmente fechado'' sem modificar a validade dos resultados (e, nas hipóteses de \sref{II.5.4.3}, \sref{II.5.4.5}, \sref{II.5.4.6} e \sref{II.5.4.8}, podemos omitir as condições de finitude).
\end{remark}

\begin{env}[5.4.10]
\label{II.5.4.10}
Seja $f:X\to Y$ um morfismo de tipo finito.
Diremos que um subconjunto fechado $Z$ de $X$ é \emph{próprio sobre $Y$} (ou \emph{$Y$-próprio}, ou \emph{próprio para $f$}) se a restrição de $f$ a um subpré-esquema fechado de $X$ cujo espaço subjacente seja $Z$ \sref[I]{I.5.2.1} é própria.
Como essa restrição é então separada, decorre de \sref{II.5.4.6} e \sref[I]{I.5.5.1}[(vi)] que a propriedade precedente \emph{não depende} do subpré-esquema fechado de $X$ que tem $Z$ como espaço subjacente.
Se $g:X'\to X$ é um morfismo \emph{próprio}, então $g^{-1}(Z)$ é um subconjunto \emph{próprio} de $X'$:
se $T$ é um subpré-esquema de $X$ cujo espaço subjacente é $Z$, basta observar que a restrição de $g$ ao subpré-esquema fechado $g^{-1}(T)$ de $X'$ é um morfismo próprio $g^{-1}(T)\to T$ por \sref{II.5.4.2}[(iii)], e aplicar em seguida \sref{II.5.4.2}[(ii)].
Por outro lado, se $X''$ é um $Y$-\emph{esquema} de tipo finito e $u:X\to X''$ é um $Y$-morfismo, então $u(Z)$ é um subconjunto \emph{próprio} de $X''$;
com efeito, tomemos para $T$ o subpré-esquema fechado reduzido de $X$ cujo espaço subjacente é $Z$;
como a restrição de $f$ a $T$ é própria, o mesmo vale para a restrição de $u$ a $T$ \sref{II.5.4.3}[(i)], logo $u(Z)$ é fechado em $X''$;
seja $T''$ um subpré-esquema fechado de $X''$ cujo espaço subjacente é $u(Z)$ \sref[I]{I.5.2.1}, de modo que $u|T$ se fatore como $T\xrightarrow{v}T''\xrightarrow{j}X''$, onde $j$ é a injeção canônica \sref[I]{I.5.2.2}; então $v$ é próprio por \sref{II.5.4.3}[(i)] e sobrejetivo pela escolha de $T''$;
se $g$ é a restrição a $T''$ do morfismo estrutural $X''\to Y$, então $g$ é separado e de tipo finito, e temos $f|T=g\circ v$;
decorre, portanto, de \sref{II.5.4.3}[(ii)] que $g$ é próprio, donde nossa afirmação.
\end{env}

Decorre, em particular, dessas observações que, se $Z$ é um subconjunto $Y$-próprio de $X$, então:
\begin{enumerate}
  \item para todo subpré-esquema fechado $X'$ de $X$, $Z\cap X'$ é um subconjunto $Y$-próprio de $X'$; e
  \item se $X$ é um subpré-esquema de um $Y$-esquema de tipo finito $X''$, então $Z$ é também um subconjunto $Y$-próprio de $X''$ (e, em particular, é \emph{fechado em $X''$}).
\end{enumerate}

\subsection{Morfismos projetivos}
\label{subsection:II.5.5}

\begin{proposition}[5.5.1]
\label{II.5.5.1}
Seja $X$ um $Y$-pré-esquema.
As condições seguintes são equivalentes:
\begin{enumerate}
  \item[{\rm(a)}] $X$ é $Y$-isomorfo a um subpré-esquema \emph{fechado} de um fibrado projetivo $\bb{P}(\sh{E})$, onde $\sh{E}$ é um $\sh{O}_Y$-módulo quase-coerente de tipo finito.
  \item[{\rm(b)}] Existe uma $\sh{O}_Y$-álgebra graduada quase-coerente $\sh{S}$ tal que $\sh{S}_1$ seja de tipo finito e gere $\sh{S}$, e tal que $X$ seja $Y$-isomorfo a $\Proj(\sh{S})$.
\end{enumerate}
\end{proposition}

\begin{proof}
A condição~(a) implica (b) em virtude de \sref{II.3.6.2}[(ii)]: se $\sh{J}$ é um feixe graduado quase-coerente de ideais de $\bb{S}(\sh{E})$, a $\sh{O}_Y$-álgebra graduada quase-coerente $\sh{S}=\bb{S}(\sh{E})/\sh{J}$ é gerada por $\sh{S}_1$, e $\sh{S}_1$, imagem canônica de $\sh{E}$, é um $\sh{O}_Y$-módulo de tipo finito.
A condição~(b) implica (a) em virtude de \sref{II.3.6.2}, aplicado ao caso em que $\sh{M}\to\sh{S}_1$ é a aplicação identidade.
\end{proof}
\oldpage[II]{104}
\begin{definition}[5.5.2]
\label{II.5.5.2}
Diz-se que um $Y$-pré-esquema $X$ é \emph{projetivo} sobre $Y$, ou que é um $Y$-esquema projetivo, se satisfaz as condições equivalentes~(a) e (b) de \sref{II.5.5.1}.
Diz-se que um morfismo $f:X\to Y$ é projetivo se faz de $X$ um $Y$-esquema projetivo.
\end{definition}

É claro que, se $f:X\to Y$ é projetivo, existe um $\sh{O}_X$-módulo \emph{muito amplo} relativamente a $f$ \sref{II.4.4.2}.

\begin{theorem}[5.5.3]
\label{II.5.5.3}
\medskip\noindent
\begin{enumerate}
  \item[{\rm(i)}] Todo morfismo projetivo é quase-projetivo e próprio.
  \item[{\rm(ii)}] Reciprocamente, seja $Y$ um esquema quase-compacto ou um pré-esquema cujo espaço subjacente seja noetheriano; então todo morfismo $f:X\to Y$ que seja quase-projetivo e próprio é projetivo.
\end{enumerate}
\end{theorem}

\begin{proof}
\medskip\noindent
\begin{enumerate}
  \item[(i)] É claro que, se $f:X\to Y$ é projetivo, então é de tipo finito e quase-projetivo (portanto, em particular, separado); além disso, decorre imediatamente de \sref{II.5.5.1}[(b)] e de \sref{II.3.5.3} que, se $f$ é projetivo, o mesmo vale para $f\times_Y 1_{Y'}:X\times_Y Y'\to Y'$ para todo morfismo $Y'\to Y$.
    Para mostrar que $f$ é universalmente fechado, basta, portanto, provar que todo morfismo projetivo $f$ é \emph{fechado}.
    Como a questão é local sobre $Y$, podemos supor que $Y=\Spec(A)$ e, por \sref{II.5.5.1}, que $X=\Proj(S)$, onde $S$ é uma $A$-álgebra graduada gerada por um número finito de elementos de $S_1$.
    Para todo $y\in Y$, a fibra $f^{-1}(y)$ se identifica com $\Proj(S)\times_Y\Spec(\kres(y))$ \sref[I]{I.3.6.1} e, portanto, também com $\Proj(S\otimes_A\kres(y))$ \sref{II.2.8.10};
    assim, $f^{-1}(y)$ é vazia se, e somente se, $S\otimes_A\kres(y)$ satisfaz a condição~(TN) \sref{II.2.7.4}, isto é, se $S_n\otimes_A\kres(y)=0$ para $n$ suficientemente grande.
    Mas, como $(S_n)_y$ é um $\sh{O}_y$-módulo de tipo finito, a condição precedente implica que $(S_n)_y=0$ para $n$ suficientemente grande, em virtude do lema de Nakayama.
    Se $\mathfrak{a}_n$ é o anulador, em $A$, do $A$-módulo $S_n$, a condição precedente implica ainda que $\mathfrak{a}_n\not\subset\mathfrak{j}_y$ para $n$ suficientemente grande \sref[0]{0.1.7.4}.
    Ora, como $S_nS_1=S_{n+1}$ por hipótese, temos $\mathfrak{a}_n\subset\mathfrak{a}_{n+1}$ e, se $\mathfrak{a}$ é a união dos $\mathfrak{a}_n$, vemos, portanto, que $f(X)=V(\mathfrak{a})$, o que demonstra que $f(X)$ é fechado em $Y$.
    Se, agora, $X'$ é um subconjunto fechado qualquer de $X$, existe um subpré-esquema fechado de $X$ cujo espaço subjacente é $X'$ \sref[I]{I.5.2.1}, e é claro, por \sref{II.5.5.1}[(a)], que o morfismo $X'\to X\xrightarrow{f}Y$ é projetivo; portanto, $f(X')$ é fechado em $Y$.
  \item[(ii)] A hipótese sobre $Y$ e o fato de que $f$ é quase-projetivo implicam a existência de um $\sh{O}_Y$-módulo quase-coerente de tipo finito $\sh{E}$ e de uma $Y$-imersão $j:X\to\bb{P}(\sh{E})$ \sref{II.5.3.2}.
    Mas, como $f$ é próprio, $j$ é \emph{fechada} por \sref{II.5.4.4}; portanto, $f$ é projetivo.
\end{enumerate}
\end{proof}

\begin{remark}[5.5.4]
\label{II.5.5.4}
\medskip\noindent
\begin{enumerate}
  \item[(i)] Seja $f:X\to Y$ um morfismo tal que: $1^\circ$ $f$ seja próprio; $2^\circ$ exista um $\sh{O}_X$-módulo $\sh{L}$ \emph{muito amplo} relativamente a $f$; $3^\circ$ o $\sh{O}_Y$-módulo quase-coerente $\sh{E}=f_*(\sh{L})$ seja \emph{de tipo finito}.
    Então $f$ é um morfismo \emph{projetivo}: com efeito, por \sref{II.4.4.4}, existe uma $Y$-imersão $r:X\to\bb{P}(\sh{E})$ e, como $f$ é próprio, $r$ é uma imersão \emph{fechada} \sref{II.5.4.4}.
    Veremos no Capítulo~III, \textsection3, que, quando $Y$ é \emph{localmente noetheriano}, a terceira condição acima ($\sh{E}$ ser de tipo finito) é consequência das outras duas; portanto, nesse caso, as duas primeiras condições \emph{caracterizam} os morfismos projetivos e, se $Y$ é quase-compacto, podemos substituir a segunda condição (a existência de um $\sh{O}_X$-módulo $\sh{L}$ muito amplo relativamente a $f$) pela hipótese de que existe um $\sh{O}_X$-módulo \emph{amplo} relativamente a $f$ \sref{II.4.6.11}.
  \item[(ii)] Seja $Y$ um esquema quase-compacto tal que exista um $\sh{O}_Y$-módulo amplo.
    Para que um $Y$-esquema $X$ seja \emph{projetivo}, é necessário e suficiente que seja $Y$-isomorfo a um $Y$-subesquema \emph{fechado} de um fibrado projetivo da forma $\bb{P}_Y^r$.
    A condição é evidentemente suficiente.
\oldpage[II]{105}
Reciprocamente, se $X$ é projetivo sobre $Y$, então $X$ é quase-projetivo; portanto, existe uma $Y$-imersão $j$ de $X$ em algum $\bb{P}_Y^r$ \sref{II.5.3.3}, que é \emph{fechada} por \sref{II.5.4.4} e \sref{II.5.5.3}.
  \item[\rm{(iii)}] O raciocínio de \sref{II.5.5.3} mostra que, para todo pré-esquema $Y$ e todo inteiro $r\geq0$, o morfismo estrutural $\bb{P}_Y^r\to Y$ é \emph{sobrejetivo}, pois, se pusermos $\sh{S}=\bb{S}_{\sh{O}_Y}(\sh{O}_Y^{r+1})$, teremos evidentemente $\sh{S}_y=\bb{S}_{\kres(y)}(\kres(y)^{r+1})$ \sref{II.1.7.3}, e portanto $(\sh{S}_n)_y\neq0$ para todo $y\in Y$ e todo $n\geq0$.
  \item[\rm{(iv)}] Resulta dos exemplos de Nagata~\cite{II-26} que existem morfismos próprios que não são quase-projetivos.
\end{enumerate}
\end{remark}

\begin{proposition}[5.5.5]
\label{II.5.5.5}
\medskip\noindent
\begin{enumerate}
  \item[\rm{(i)}] Uma imersão fechada é um morfismo projetivo.
  \item[\rm{(ii)}] Se $f:X\to Y$ e $g:Y\to Z$ são morfismos projetivos, e se $Z$ é um esquema quase-compacto ou um pré-esquema cujo espaço subjacente é noetheriano, então $g\circ f$ é projetivo.
  \item[\rm{(iii)}] Se $f:X\to Y$ é um $S$-morfismo projetivo, então $f_{(S')}:X_{(S')}\to Y_{(S')}$ é projetivo para toda extensão $S'\to S$ do pré-esquema de base.
  \item[\rm{(iv)}] Se $f:X\to Y$ e $g:X'\to Y'$ são $S$-morfismos projetivos, o mesmo vale para $f\times_S g$.
  \item[\rm{(v)}] Se $g\circ f$ é um morfismo projetivo e $g$ é separado, então $f$ é projetivo.
  \item[\rm{(vi)}] Se $f$ é projetivo, o mesmo vale para $f_\red$.
\end{enumerate}
\end{proposition}

\begin{proof}
(i) resulta imediatamente de \sref{II.3.1.7}.
É preciso demonstrar aqui (iii) e (iv) separadamente, por causa da restrição introduzida sobre $Z$ em (ii) (cf.~\sref[I]{I.3.5.1}).
Para provar (iii), reduzimo-nos ao caso em que $S=Y$ \sref[I]{I.3.3.11}, e a afirmação resulta então imediatamente de \sref{II.5.5.1}[b] e \sref{II.3.5.3}.
Para demonstrar (iv), reduzimo-nos imediatamente ao caso em que $X=\bb{P}(\sh{E})$ e $X'=\bb{P}(\sh{E}')$, sendo $\sh{E}$ (resp. $\sh{E}'$) um $\sh{O}_Y$-módulo (resp. um $\sh{O}_{Y'}$-módulo) quase-coerente de tipo finito.
Sejam $p$ e $p'$ as projeções canônicas de $T=Y\times_S Y'$ em $Y$ e $Y'$, respectivamente; por \sref{II.4.1.3.1}, temos $\bb{P}(p^*(\sh{E}))=\bb{P}(\sh{E})\times_Y T$ e $\bb{P}(p^{\prime *}(\sh{E}'))=\bb{P}(\sh{E}')\times_{Y'}T$; daí
\begin{align*}
  \bb{P}(p^*(\sh{E}))\times_T\bb{P}(p^{\prime *}(\sh{E}'))&=(\bb{P}(\sh{E})\times_Y T)\times_T(T\times_{Y'}\bb{P}(\sh{E}'))\\
                                               &=\bb{P}(\sh{E})\times_Y(T\times_{Y'}\bb{P}(\sh{E}'))=\bb{P}(\sh{E})\times_S\bb{P}(\sh{E}')
\end{align*}
substituindo $T$ por $Y\times_S Y'$ e utilizando \sref[I]{I.3.3.9.1}.
Ora, $p^*(\sh{E})$ e $p^{\prime *}(\sh{E}')$ são de tipo finito sobre $T$ \sref[0]{0.5.2.4}, e portanto o mesmo vale para $p^*(\sh{E})\otimes_{\sh{O}_T}p^{\prime *}(\sh{E}')$;
como $\bb{P}(p^*(\sh{E}))\times_T\bb{P}(p^{\prime *}(\sh{E}'))$ se identifica com um subpré-esquema fechado de $\bb{P}(p^*(\sh{E})\otimes_{\sh{O}_T}p^{\prime *}(\sh{E}'))$ \sref{II.4.3.3}, isso termina de provar (iv).
Para obter (v) e (vi), podemos aplicar \sref[I]{I.5.5.13}, pois todo subpré-esquema fechado de um $Y$-esquema projetivo é um $Y$-esquema projetivo por \sref{II.5.5.1}[a].

Resta provar (ii); em virtude da hipótese sobre $Z$, isso resulta de \sref{II.5.5.3}, \sref{II.5.3.4}[ii] e \sref{II.5.4.2}[ii].
\end{proof}

\begin{proposition}[5.5.6]
\label{II.5.5.6}
Se $X$ e $X'$ são dois $Y$-esquemas projetivos, então $X\sqcup X'$ é um $Y$-esquema projetivo.
\end{proposition}

\begin{proof}
É uma consequência evidente de \sref{II.5.5.2} e \sref{II.4.3.6}.
\end{proof}

\begin{proposition}[5.5.7]
\label{II.5.5.7}
Sejam $X$ um $Y$-esquema projetivo e $\sh{L}$ um $\sh{O}_X$-módulo amplo relativamente a $Y$;
para toda seção $f$ de $\sh{L}$ sobre $X$, $X_f$ é afim sobre $Y$.
\end{proposition}
\oldpage[II]{106}
\begin{proof}
Como a questão é local sobre $Y$, podemos supor que $Y=\Spec(A)$; além disso, $X_{f^{\otimes n}}=X_f$, logo, substituindo $\sh{L}$ por um $\sh{L}^{\otimes n}$ conveniente, podemos supor que $\sh{L}$ é muito amplo relativamente ao morfismo estrutural $q:X\to Y$ \sref{II.4.6.11}.
O homomorfismo canônico $\sigma:q^*(q_*(\sh{L}))\to\sh{L}$ é, portanto, sobrejetivo, e o morfismo correspondente
\[
  r=r_{\sh{L},\sigma}:X\to P=\bb{P}(q_*(\sh{L}))
\]
é uma imersão tal que $\sh{L}=r^*(\sh{O}_P(1))$ \sref{II.4.4.4}; além disso, como $X$ é próprio sobre $Y$, a imersão $r$ é fechada \sref{II.5.4.4}.
Por outro lado, temos, por definição, $f\in\Gamma(Y,q_*(\sh{L}))$, e $\sigma^\flat$ é a identidade de $q_*(\sh{L})$; resulta então da fórmula \sref{II.3.7.3.1} que $X_f=r^{-1}(D_+(f))$;
$X_f$ é, por conseguinte, um subpré-esquema fechado do esquema afim $D_+(f)$ e é, portanto, um esquema afim.
\end{proof}

Tomando, em particular, $Y=X$, obtemos (levando em conta \sref{II.4.6.13}[i]) o corolário seguinte, cuja demonstração direta é, aliás, imediata:
\begin{corollary}[5.5.8]
\label{II.5.5.8}
Sejam $X$ um pré-esquema e $\sh{L}$ um $\sh{O}_X$-módulo invertível.
Para toda seção $f$ de $\sh{L}$ sobre $X$, $X_f$ é afim sobre $X$ (e, portanto, um esquema afim quando $X$ é um esquema afim).
\end{corollary}

\subsection{O lema de Chow}
\label{subsection:II.5.6}

\begin{theorem}[5.6.1]
\label{II.5.6.1}
\emph{(Lema de Chow)}.
Sejam $S$ um pré-esquema e $X$ um $S$-esquema de tipo finito.
Suponhamos verificada uma das hipóteses seguintes:
\begin{enumerate}
  \item[\rm{(a)}] $S$ é noetheriano;
  \item[\rm{(b)}] $S$ é um esquema quase-compacto, e $X$ tem um número finito de componentes irredutíveis.
\end{enumerate}
Nessas condições:
\begin{enumerate}
  \item[\rm{(i)}] existe um $S$-esquema $X'$ \emph{quase-projetivo} sobre $S$ e um $S$-morfismo $f:X'\to X$ \emph{projetivo} e \emph{sobrejetivo};
  \item[\rm{(ii)}] podem-se tomar $X'$ e $f$ tais que exista um aberto $U\subset X$ para o qual $U'=f^{-1}(U)$ seja denso em $X'$ e a restrição de $f$ a $U'$ seja um isomorfismo $U'\isoto U$;
  \item[\rm{(iii)}] se $X$ é reduzido (resp. irredutível, integral), pode-se supor que $X'$ é reduzido (resp. irredutível, integral).
\end{enumerate}
\end{theorem}

\begin{proof}
A demonstração faz-se em várias etapas.
\begin{enumerate}
  \item[(A)] Podemos primeiro reduzir-nos ao caso em que $X$ é \emph{irredutível}.
    Com efeito, sob a hipótese~(a), $X$ é noetheriano; portanto, sob ambas as hipóteses, as componentes irredutíveis $X_i$ de $X$ são em número finito.
    Se o teorema foi demonstrado para cada um dos subpré-esquemas fechados reduzidos de $X$ que têm os $X_i$ por espaços subjacentes, e se $X'_i$ e $f_i:X'_i\to X_i$ são, respectivamente, o pré-esquema e o morfismo correspondentes a $X_i$, então o pré-esquema $X'$, soma dos $X'_i$, e o morfismo $f:X'\to X$ cuja restrição a cada $X'_i$ é $j_i\circ f_i$ (sendo $j_i$ a injeção canônica $X_i\to X$) satisfazem a conclusão do teorema.
    É imediato, com efeito, que $X'$ é reduzido se todos os $X'_i$ o são; por outro lado, satisfaz-se (ii) tomando para $U$ a união dos conjuntos $U_i\cap\complement\left(\bigcup_{j\neq i}X_j\right)$.
    Finalmente, como os $X'_i$ são quase-projetivos sobre $S$, também o é $X'$
\oldpage[II]{107}
    \sref{II.5.3.6}; do mesmo modo, os morfismos $X'_i\to X$ são projetivos por \sref{II.5.5.5}[i] e \sref{II.5.5.5}[ii], logo $f$ é projetivo \sref{II.5.5.6}, e é evidentemente sobrejetivo por definição.
  \item[(B)] Suponhamos agora $X$ \emph{irredutível}.
    Como o morfismo estrutural $r:X\to S$ é de tipo finito, há um recobrimento finito $(S_i)$ de $S$ por abertos afins e, para cada $i$, um recobrimento finito $(T_{ij})$ de $r^{-1}(S_i)$ por abertos afins; como o morfismo $T_{ij}\to S_i$ é afim e de tipo finito, ele é quase-projetivo \sref{II.5.3.4}[i];
    como, em cada uma das hipóteses~(a) e (b), a imersão $S_i\to S$ é quase-compacta, ela é um morfismo quase-projetivo \sref{II.5.3.4}[i], e portanto a restrição de $r$ a $T_{ij}$ é um morfismo quase-projetivo \sref{II.5.3.4}[ii].
    Denotemos os $T_{ij}$ por $U_k$ ($1\leq k\leq n$).
    Existe, para cada índice $k$, uma imersão aberta $\vphi_k:U_k\to P_k$, onde $P_k$ é projetivo sobre $S$ (\sref{II.5.3.2} e \sref{II.5.5.2}).
    Seja $U=\bigcap_k U_k$; como $X$ é irredutível e os $U_k$ são não vazios, $U$ é não vazio e, portanto, denso em toda parte em $X$; as restrições dos $\vphi_k$ a $U$ definem um morfismo
    \[
      \vphi:U\to P=P_1\times_S P_2\times_S\cdots\times_S P_n
    \]
    de modo que os diagramas
    \[
    \label{II.5.6.1.1}
      \xymatrix{
        U\ar[r]^\vphi\ar[d]_{j_k} &
        P\ar[d]^{p_k}\\
        U_k\ar[r]_{\vphi_k} &
        P_k
      }
      \tag{5.6.1.1}
    \]
    sejam comutativos, sendo $j_k$ a injeção canônica $U\to U_k$ e $p_k$ a projeção canônica $P\to P_k$.
    Se $j$ é a injeção canônica $U\to X$, o morfismo $\psi=(j,\vphi)_S:U\to X\times_S P$ é uma \emph{imersão} \sref[I]{I.5.3.14}.
    Sob a hipótese~(a), $X\times_S P$ é localmente noetheriano (\sref{II.3.4.1}, \sref[I]{I.6.3.7} e \sref[I]{I.6.3.8});
    sob a hipótese~(b), $X\times_S P$ é um esquema quase-compacto (\sref[I]{I.5.5.1} e \sref[I]{I.6.6.4});
    em ambos os casos, existe o \emph{fecho} $X'$ em $X\times_S P$ do subpré-esquema $Z$ associado a $\psi$ (e cujo espaço subjacente é $\psi(U)$), e $\psi$ se fatoriza como
    \[
    \label{II.5.6.1.2}
      \psi:U\xrightarrow{\psi'}X'\xrightarrow{h}X\times_S P,
      \tag{5.6.1.2}
    \]
    onde $\psi'$ é uma \emph{imersão aberta} e $h$ uma \emph{imersão fechada} \sref[I]{I.9.5.10}.
    Sejam $q_1:X\times_S P\to X$ e $q_2:X\times_S P\to P$ as projeções canônicas; poremos
    \[
    \label{II.5.6.1.3}
      f:X'\xrightarrow{h}X\times_S P\xrightarrow{q_1}X,
      \tag{5.6.1.3}
    \]
    \[
    \label{II.5.6.1.4}
      g:X'\xrightarrow{h}X\times_S P\xrightarrow{q_2}P.
      \tag{5.6.1.4}
    \]
    Veremos que $X'$ e $f$ satisfazem a conclusão do teorema.
  \item[(C)] Mostremos primeiro que $f$ é \emph{projetivo} e \emph{sobrejetivo}, e que a restrição de $f$ a $U'=f^{-1}(U)$ é um \emph{isomorfismo} de $U'$ sobre $U$.
    Como os $P_k$ são projetivos sobre $S$, o mesmo vale para $P$ \sref{II.5.5.5}[iv], logo $X\times_S P$ é projetivo sobre $X$ \sref{II.5.5.5}[iii], e o mesmo vale para $X'$, que é um subpré-esquema fechado de $X\times_S P$.
    Por outro lado, temos $f\circ\psi'=q_1\circ(h\circ\psi')=q_1\circ\psi=j$, logo $f(X')$ contém o aberto denso em toda parte $U$ de $X$; mas $f$ é um morfismo \emph{fechado} \sref{II.5.5.3}, portanto $f(X')=X$.
    Observemos agora que $q_1^{-1}(U)=U\times_S P$ é induzido sobre um aberto de $X\times_S P$ e, por definição, o pré-esquema $U'=h^{-1}(U\times_S P)$ é induzido por $X'$ sobre o aberto $U'$; ele é, por conseguinte, o fecho \emph{relativo}
\oldpage[II]{108}
a $U\times_S P$ do subpré-esquema $Z$ \sref[I]{I.9.5.8}.
Ora, a imersão $\psi$ se fatoriza como $U\xrightarrow{\Gamma_\vphi}U\times_S P\xrightarrow{j\times1}X\times_S P$ e, como $P$ é separado sobre $S$, o morfismo gráfico $\Gamma_\vphi$ é uma imersão fechada \sref[I]{I.5.4.3}; logo $Z$ é um subpré-esquema \emph{fechado} de $U\times_S P$, de onde se conclui que $U'=Z$.
Como, além disso, $\psi$ é uma imersão, a restrição de $f$ a $U'$ é um isomorfismo sobre $U$, inverso de $\psi'$; finalmente, por definição de $X'$, $U'$ é denso em $X'$.
  \item[(D)] Provemos agora que $g$ é uma \emph{imersão}, o que estabelecerá que $X'$ é \emph{quase-projetivo} sobre $S$, pois $P$ é projetivo sobre $S$.
    Ponhamos
    \begin{align*}
      V_k  &=\vphi_k(U_k) \quad(\text{aberto de }P_k)\\
      W_k  &=p_k^{-1}(V_k)\quad(\text{aberto de }P)\\
      U_k' &=f^{-1}(U_k)  \quad(\text{aberto de }X')\\
      U_k''&=g^{-1}(W_k)  \quad(\text{aberto de }X').
    \end{align*}
    É claro que os $U'_k$ formam um recobrimento aberto de $X'$; veremos primeiro que o mesmo vale para os $U_k''$, provando que $U_k'\subset U_k''$.
    Para isso, bastará mostrar que o diagrama
    \[
    \label{II.5.6.1.5}
      \xymatrix{
        U_k'\ar[r]^{g|U_k'}\ar[d]_{f|U_k'} &
        P\ar[d]^{p_k}\\
        U_k\ar[r]_{\vphi_k} &
        P_k
      }
      \tag{5.6.1.5}
    \]
    é comutativo.
    Ora, o pré-esquema $U_k'=h^{-1}(U_k\times_S P)$ é induzido por $X'$ sobre o aberto $U_k'$ e é, portanto, o fecho de $Z=U'\subset U_k'$ relativo a $U_k'$ \sref[I]{I.9.5.8}.
    Para mostrar a comutatividade do diagrama \sref{II.5.6.1.5}, basta, portanto, como $P_k$ é um $S$-esquema, mostrar que, compondo esse diagrama com a injeção canônica $U'\to U_k'$ (ou, o que dá no mesmo em virtude do isomorfismo de $U'$ e $U$, com $\psi$), obtemos um diagrama comutativo \sref[I]{I.9.5.6}.
    Mas, por definição, o diagrama que então obtemos não é outro senão \sref{II.5.6.1.1}, donde nossa afirmação.

    Os $W_k$ formam, portanto, um recobrimento aberto de $g(X')$; para provar que $g$ é uma imersão, basta ver que cada uma das restrições $g|U_k''$ é uma imersão em $W_k$ \sref[I]{I.4.2.4}.
    Para isso, consideremos o morfismo $u_k:W_k\xrightarrow{p_k}V_k\xrightarrow{\vphi_k^{-1}}U_k\to X$; como $X$ é separado sobre $S$, o morfismo gráfico $\Gamma_{u_k}:W_k\to X\times_S W_k$ é uma imersão fechada \sref[I]{I.5.4.3}, logo o gráfico $T_k=\Gamma_{u_k}(W_k)$ é um subpré-esquema fechado de $X\times_S W_k$;
    se mostrarmos que o morfismo $U'\to X\times_S W_k$ se fatoriza por esse subpré-esquema, então o morfismo para $X\times_S W_k$ proveniente do subpré-esquema induzido por $X'$ sobre o aberto $U_k''$ de $X'$ também se fatorizará por esse gráfico, por \sref[I]{I.9.5.8}.
    Como a restrição de $q_2$ a $T_k$ é um isomorfismo sobre $W_k$, a restrição de $g$ a $U''_k$ será uma imersão em $W_k$, e nossa afirmação estará demonstrada.
    Seja $v_k$ a injeção canônica $U'\to X\times_S W_k$; devemos mostrar que existe um morfismo $w_k:U'\to W_k$ tal que $v_k=\Gamma_{u_k}\circ w_k$.
    Por definição do produto, basta provar que $q_1\circ v_k=u_k\circ q_2\circ v_k$ \sref[I]{I.3.2.1}, ou, compondo à direita
\oldpage[II]{109}
com o isomorfismo $\psi':U\to U'$, que $q_1\circ\psi=u_k\circ q_2\circ\psi$.
Mas, como $q_1\circ\psi=j$ e $q_2\circ\psi=\vphi$, nossa afirmação resulta da comutatividade do diagrama \sref{II.5.6.1.1}, levando em conta a definição de $u_k$.
  \item[(E)] É claro que, como $U$ e, portanto, $U'$ são irredutíveis, o mesmo vale para $X'$ na construção precedente, e o morfismo $f$ é, pois, \emph{birracional} \sref[I]{I.2.2.9}.
    Se, além disso, $X$ é reduzido, o mesmo vale para $U'$, logo $X'$ é reduzido \sref[I]{I.9.5.9}.
    Isso conclui a demonstração.
\end{enumerate}
\end{proof}

\begin{corollary}[5.6.2]
\label{II.5.6.2}
Suponhamos verificada uma das hipóteses \emph{(a), (b)} de \sref{II.5.6.1}.
Para que $X$ seja próprio sobre $S$, é necessário e suficiente que exista um esquema $X'$ projetivo sobre $S$ e um $S$-morfismo sobrejetivo $f:X'\to X$ (que é então projetivo por \sref{II.5.5.5}[v]).
Quando isso ocorre, pode-se ainda escolher $f$ de modo que exista um aberto denso $U$ de $X$ tal que a restrição de $f$ a $f^{-1}(U)$ seja um isomorfismo $f^{-1}(U)\isoto U$ e que $f^{-1}(U)$ seja denso em $X'$.
Se, além disso, $X$ é irredutível (resp. reduzido), pode-se supor que o mesmo vale para $X'$; quando $X$ e $X'$ são irredutíveis, $f$ é um morfismo birracional.
\end{corollary}

\begin{proof}
A condição é suficiente em virtude de \sref{II.5.5.3} e \sref{II.5.4.3}[ii].
Ela é necessária, pois, com a notação de \sref{II.5.6.1}, se $X$ é próprio sobre $S$, então $X'$ é próprio sobre $S$, já que é projetivo sobre $X$, logo próprio sobre $X$ \sref{II.5.5.3}, e nossa afirmação resulta de \sref{II.5.4.2}[ii]; além disso, como $X'$ é quase-projetivo sobre $S$, ele é projetivo sobre $S$ em virtude de \sref{II.5.5.3}.
\end{proof}

\begin{corollary}[5.6.3]
\label{II.5.6.3}
Sejam $S$ um pré-esquema localmente noetheriano e $X$ um $S$-esquema de tipo finito sobre $S$, e seja $f_0:X\to S$ seu morfismo estrutural.
Para que $X$ seja próprio sobre $S$, é necessário e suficiente que, para todo morfismo \emph{de tipo finito} $S'\to S$, $(f_0)_{(S')}:X_{(S')}\to S'$ seja um morfismo fechado.
Basta mesmo que essa condição seja verificada para todo $S$-pré-esquema da forma $S'=S\otimes_\bb{Z}\bb{Z}[T_1,\ldots,T_n]$ (sendo $T_1,\ldots,T_n$ indeterminadas).
\end{corollary}

\begin{proof}
Sendo a condição evidentemente necessária, mostremos que ela é suficiente.
Como a questão é local sobre $S$ e $S'$ \sref{II.5.4.1}, podemos supor $S$ e $S'$ afins e noetherianos.
Em virtude do lema de Chow, existe um $S$-esquema projetivo $P$, uma imersão $j:X'\to P$ e um morfismo projetivo e sobrejetivo $f:X'\to X$, tais que o diagrama
\[
  \xymatrix{
    X\ar[d]_{f_0} &
    X'\ar[l]_f\ar[d]^j\\
    S &
    P\ar[l]_r
  }
\]
seja comutativo.
Como $P$ é de tipo finito sobre $S$, a primeira hipótese implica que a projeção $q_2:X\times_S P\to P$ é um morfismo \emph{fechado}.
O morfismo $(f,j):X'\to X\times_S P$ se fatoriza como a imersão gráfica $\Gamma_j:X'\to X'\times_S P$, que é fechada porque $P$ é separado sobre $S$, seguida da mudança de base $f\times1_P:X'\times_S P\to X\times_S P$.
Ora, $f$, sendo projetivo, é próprio \sref{II.5.5.3}, logo $f\times1_P$ é fechado; portanto, $(f,j)$ é fechado.
Como $j=q_2\circ(f,j)$ e $j$ já é uma imersão, concluímos que $j$ é uma imersão fechada e, por conseguinte, própria \sref{II.5.4.2}[i].
Além disso, o morfismo estrutural $r:P\to S$ é projetivo, logo próprio \sref{II.5.5.3}; portanto, $f_0\circ f=r\circ j$ é próprio \sref{II.5.4.2}[ii]; finalmente, como $f$ é sobrejetivo, $f_0$ é próprio em virtude de \sref{II.5.4.3}.

Para demonstrar a proposição utilizando somente a segunda hipótese, basta mostrar que ela implica a primeira.
Ora, se $S'$ é afim e de tipo finito sobre $S=\Spec(A)$,
\oldpage[II]{110}
então temos $S'=\Spec(A[c_1,\ldots,c_n])$ \sref[I]{I.6.3.3}, e $S'$ é, portanto, isomorfo a um subpré-esquema fechado de $S''=\Spec(A[T_1,\ldots,T_n])$ (sendo os $T_i$ indeterminados).
No diagrama comutativo
\[
  \xymatrix{
    X\times_S S'\ar[r]^{1_X\times j}\ar[d]_{(f_0)_{(S')}} &
    X\times_S S''\ar[d]^{(f_0)_{(S'')}}\\
    S'\ar[r]_j &
    S''
  }
\]
tanto $j$ como $1_X\times j$ são imersões fechadas \sref[I]{I.4.3.1}, e $(f_0)_{(S'')}$ é fechado por hipótese;
portanto, $(f_0)_{(S')}$ também é fechado.
\end{proof}

\section{Morfismos inteiros e morfismos finitos}
\label{section:II.6}

\subsection{Pré-esquemas inteiros sobre outro}
\label{subsection:II.6.1}

\begin{definition}[6.1.1]
\label{II.6.1.1}
Sejam $X$ um $S$-pré-esquema e $f:X\to S$ seu morfismo estrutural.
Dizemos que $X$ é \emph{inteiro sobre $S$}, ou que $f$ é um \emph{morfismo inteiro}, se existe um recobrimento $(S_\alpha)$ de $S$ por abertos afins tal que, para todo $\alpha$, o pré-esquema induzido $f^{-1}(S_\alpha)$ é um esquema afim cujo anel $B_\alpha$ é uma álgebra inteira sobre o anel $A_\alpha$ de $S_\alpha$.
Dizemos que $X$ é \emph{finito sobre $S$}, ou que $f$ é um \emph{morfismo finito}, se $X$ é inteiro sobre $S$ e de tipo finito sobre $S$.
\end{definition}

Quando $S$ é afim, de anel $A$, diremos também ``inteiro (resp. finito) sobre $A$'' em vez de ``inteiro (resp. finito) sobre $S$''.

\begin{env}[6.1.2]
\label{II.6.1.2}
É claro que, se $X$ é inteiro sobre $S$, então é \emph{afim} sobre $S$.
Para que um pré-esquema $X$ afim sobre $S$ seja inteiro (resp. finito) sobre $S$, é necessário e suficiente que a $\sh{O}_S$-álgebra quase-coerente associada $\sh{A}(X)$ seja tal que exista um recobrimento $(S_\alpha)$ de $S$ por abertos afins com a propriedade de que, para todo $\alpha$, $\Gamma(S_\alpha,\sh{A}(X))$ seja uma álgebra inteira (resp. inteira e de tipo finito) sobre $\Gamma(S_\alpha,\sh{O}_S)$.
Uma $\sh{O}_S$-álgebra quase-coerente com essa propriedade é dita \emph{inteira} (resp. \emph{finita}) sobre $\sh{O}_S$.
Dar um pré-esquema inteiro (resp. finito) sobre $S$ equivale, portanto, \sref{II.1.3.1} a dar uma $\sh{O}_S$-álgebra quase-coerente inteira (resp. finita) sobre $\sh{O}_S$.
Notemos que uma $\sh{O}_S$-álgebra quase-coerente $\sh{B}$ é finita se e somente se é um \emph{$\sh{O}_S$-módulo de tipo finito} \sref[I]{I.1.3.9};
equivalentemente, $\sh{B}$ é uma $\sh{O}_S$-álgebra \emph{inteira} e \emph{de tipo finito}, pois uma álgebra inteira e de tipo finito sobre um anel $A$ é um $A$-módulo de tipo finito.
\end{env}

\begin{proposition}[6.1.3]
\label{II.6.1.3}
Seja $S$ um pré-esquema localmente noetheriano.
Para que um pré-esquema $X$ afim sobre $S$ seja finito sobre $S$, é necessário e suficiente que a $\sh{O}_S$-álgebra $\sh{A}(X)$ seja coerente.
\end{proposition}
\begin{proof}
Levando em conta a observação precedente, isso equivale a notar que, se $S$ é localmente noetheriano, os $\sh{O}_S$-módulos quase-coerentes de tipo finito são precisamente os $\sh{O}_S$-módulos coerentes \sref[I]{I.1.5.1}.
\end{proof}

\begin{proposition}[6.1.4]
\label{II.6.1.4}
Seja $X$ um pré-esquema inteiro (resp. finito) sobre $S$, e seja $f:X\to S$ o morfismo estrutural.
Então, para todo aberto afim $U\subset S$, de anel $A$, $f^{-1}(U)$ é um esquema afim cujo anel $B$ é uma álgebra inteira (resp. finita) sobre $A$.
\end{proposition}
\oldpage[II]{111}
\begin{proof}
Demonstraremos primeiro o lema seguinte:

  \begin{lemma}[6.1.4.1]
  \label{II.6.1.4.1}
  Sejam $A$ um anel, $M$ um $A$-módulo e $(g_i)_{1\leq i\leq m}$ um sistema finito de elementos de $A$ tal que os $D(g_i)$ (para $1\leq i\leq m$) recubram $\Spec(A)$.
  Se, para todo $i$, $M_{g_i}$ é um $A_{g_i}$-módulo de tipo finito, então $M$ é um $A$-módulo de tipo finito.
  \end{lemma}

  \begin{proof}
  Podemos, com efeito, supor que $M_{g_i}$ admite um sistema finito de geradores $(m_{ij}/g_i^n)$, com $m_{ij}\in M$, sendo $n$ o mesmo para todos os índices $i$.
  Mostremos que os $m_{ij}$ formam um sistema de geradores de $M$.
  Seja $M'$ o sub-$A$-módulo de $M$ gerado por esses elementos, e seja $m$ um elemento de $M$.
  Por hipótese, para cada $i$, existem $a_{ij}\in A$ e um inteiro $p$ (independente de $i$) tais que, em $M_{g_i}$, $m/1=(\sum_j a_{ij}m_{ij})/g_i^p$;
  isso implica que existe um inteiro $r\geq p$ tal que, para todo $i$, temos $g_i^rm\in M'$.
  Ora, como os $D(g_i^r)=D(g_i)$ recobrem $\Spec(A)$, o ideal de $A$ gerado pelos $g_i^r$ é igual a $A$; em outras palavras, existem elementos $a_i\in A$ tais que $\sum_i a_ig_i^r=1$.
  Por conseguinte, $m=(\sum_i a_i g_i^r)m\in M'$, de onde resulta o lema.
  \end{proof}

Sabemos já \sref{II.1.3.2} que $f^{-1}(U)$ é afim.
Se $\vphi$ é o homomorfismo $A\to B$ correspondente a $f$, existe um recobrimento finito de $U$ por abertos $D(g_i)$ (com $g_i\in A$) tal que, pondo $h_i=\vphi(g_i)$, $B_{h_i}$ seja uma álgebra inteira (resp. inteira finita) sobre $A_{g_i}$.
Com efeito, existe um recobrimento de $U$ por abertos afins $V_\alpha\subset U$ tal que, pondo $A_\alpha=A(V_\alpha)$ e $B_\alpha=A(f^{-1}(V_\alpha))$, $B_\alpha$ seja uma álgebra inteira (resp. finita) sobre $A_\alpha$.
Todo $x\in U$ pertence a algum $V_\alpha$; logo, existe $g\in A$ tal que $x\in D(g)\subset V_\alpha$.
Se $g_\alpha$ é a imagem de $g$ em $A_\alpha$, temos $A(D(g))=A_g=(A_\alpha)_{g_\alpha}$;
seja $h=\vphi(g)$, e seja $h_\alpha$ a imagem de $g_\alpha$ em $B_\alpha$;
temos
\[
  A(D(h))=B_h=(B_\alpha)_{h_\alpha},
\]
e, como $B_\alpha$ é inteira (resp. finita) sobre $A_\alpha$, $(B_\alpha)_{h_\alpha}$ é inteira (resp. finita) sobre $(A_\alpha)_{g_\alpha}$.
Basta agora usar o fato de que $U$ é quase-compacto para obter o recobrimento procurado.

Suponhamos primeiro que os $B_{h_i}$ sejam inteiras finitas sobre os $A_{g_i}$;
como, enquanto $A_{g_i}$-módulo, $B_{h_i}$ também se escreve $B_{g_i}$, o Lema~\sref{II.6.1.4.1} mostra que, nesse caso, $B$ é um $A$-módulo de tipo finito.

Suponhamos agora somente que cada $B_{h_i}$ seja inteira sobre $A_{g_i}$;
seja $b\in B$, e seja $C$ a sub-$A$-álgebra de $B$ gerada por $b$.
Para todo $i$, $C_{h_i}$ é a álgebra sobre $A_{g_i}$ gerada por $b/1$ em $B_{h_i}$;
resulta da hipótese que cada $C_{h_i}$ é um $A_{g_i}$-módulo de tipo finito e, portanto, \sref{II.6.1.4.1} $C$ é um $A$-módulo de tipo finito, o que prova que $B$ é inteira sobre $A$.
\end{proof}

\begin{proposition}[6.1.5]
\label{II.6.1.5}
\medskip\noindent
\begin{enumerate}
  \item[\rm{(i)}] Uma imersão fechada é finita (e, \emph{a fortiori}, inteira).
  \item[\rm{(ii)}] A composta de dois morfismos finitos (resp. inteiros) é finita (resp. inteira).
  \item[\rm{(iii)}] Se $f:X\to Y$ é um $S$-morfismo finito (resp. inteiro), então $f_{(S')}:X_{(S')}\to Y_{(S')}$ é finito (resp. inteiro) para toda extensão da base $S'\to S$.
  \item[\rm{(iv)}] Se $f:X\to Y$ e $g:X'\to Y'$ são dois $S$-morfismos finitos (resp. inteiros), o mesmo vale para $f\times_S g:X\times_S X'\to Y\times_S Y'$.
  \item[\rm{(v)}] Se $f:X\to Y$ e $g:Y\to Z$ são dois morfismos tais que $g\circ f$ seja finito (resp. inteiro), e se $g$ é separado, então $f$ é finito (resp. inteiro).
  \item[\rm{(vi)}] Se $f:X\to Y$ é um morfismo finito (resp. inteiro), o mesmo vale para $f_\red$.
\end{enumerate}
\end{proposition}
\oldpage[II]{112}
\begin{proof}
Em virtude de \sref[I]{I.5.5.12}, basta demonstrar (i), (ii) e (iii).
Para provar que uma imersão fechada $X\to S$ é finita, podemos nos limitar ao caso em que $S=\Spec(A)$, e tudo se reduz então a observar que um anel quociente $A/\mathfrak{J}$ é um $A$-módulo gerado por um elemento.
Para provar que a composta de dois morfismos finitos (resp. inteiros) $X\to Y$, $Y\to Z$ é finita (resp. inteira), podemos ainda supor que $Z$ (e, por conseguinte, $X$ e $Y$ \sref{II.1.3.4}) seja afim; a afirmação equivale então a dizer que, se $B$ é uma $A$-álgebra finita (resp. inteira) e $C$ uma $B$-álgebra finita (resp. inteira), então $C$ é uma $A$-álgebra finita (resp. inteira), o que é imediato.
Por fim, para demonstrar (iii), podemos nos limitar ao caso em que $S=Y$, pois $X_{(S')}$ se identifica com $X\times_Y Y_{(S')}$ \sref[I]{I.3.3.11};
podemos, além disso, supor que $S=\Spec(A)$ e $S'=\Spec(A')$.
Então $X$ é afim, de anel $B$ \sref{II.1.3.4}, e $X_{(S')}$ é afim, de anel $A'\otimes_A B$;
basta observar que, se $B$ é uma $A$-álgebra finita (resp. inteira), então $A'\otimes_A B$ é uma $A'$-álgebra finita (resp. inteira).
\end{proof}

Notemos também que, se $X$ e $Y$ são dois $S$-pré-esquemas finitos (resp. inteiros) sobre $S$, sua \emph{soma} $X\sqcup Y$ é um pré-esquema finito (resp. inteiro) sobre $S$, pois isso se reduz a ver que, se $B$ e $C$ são duas $A$-álgebras finitas (resp. inteiras) sobre $A$, o mesmo vale para $B\times C$.

\begin{corollary}[6.1.6]
\label{II.6.1.6}
Se $X$ é um pré-esquema inteiro (resp. finito) sobre $S$, então, para todo aberto $U\subset S$, $f^{-1}(U)$ é inteiro (resp. finito) sobre $U$.
\end{corollary}

\begin{proof}
É um caso particular de \sref{II.6.1.5}[(iii)].
\end{proof}

\begin{corollary}[6.1.7]
\label{II.6.1.7}
Seja $f:X\to Y$ um morfismo finito.
Então, para todo $y\in Y$, a fibra $f^{-1}(y)$ é um esquema algébrico finito sobre $\kres(y)$ e, \emph{a fortiori}, seu espaço subjacente é discreto e finito.
\end{corollary}

\begin{proof}
Com efeito, enquanto $\kres(y)$-pré-esquema, $f^{-1}(y)$ identifica-se com $X\times_Y\Spec(\kres(y))$ \sref[I]{I.3.6.1}, que é finito sobre $\Spec(\kres(y))$ \sref{II.6.1.5}[(iii)];
é, portanto, um esquema afim cujo anel é uma álgebra de posto finito sobre $\kres(y)$ \sref{II.6.1.4}.
O corolário resulta então de \sref[I]{I.6.4.4}.
\end{proof}

\begin{corollary}[6.1.8]
\label{II.6.1.8}
Sejam $X$ e $S$ dois pré-esquemas integrais, e seja $f:X\to S$ um morfismo \emph{dominante}.
Se $f$ é inteiro (resp. finito), então o corpo $R(X)$ das funções racionais sobre $X$ é algébrico (resp. algébrico de grau finito) sobre o corpo $R(S)$ das funções racionais sobre $S$.
\end{corollary}

\begin{proof}
Seja $s$ o ponto genérico de $S$;
o $\kres(s)$-pré-esquema $f^{-1}(s)$ é inteiro (resp. finito) sobre $\Spec(\kres(s))$ \sref{II.6.1.5}[(iii)] e contém, por hipótese, o ponto genérico $x$ de $X$;
o anel local de $x$ em $f^{-1}(s)$ é igual a $\kres(x)$ \sref[I]{I.3.6.5} e é, portanto, um anel local de uma álgebra inteira (resp. finita) sobre $\kres(s)$ \sref{II.6.1.4}, de onde resulta o corolário.
\end{proof}

\begin{remark}[6.1.9]
\label{II.6.1.9}
A hipótese de que $g$ seja \emph{separado} é essencial para a validade de \sref{II.6.1.5}[(v)]: com efeito, se $Y$ não é separado sobre $Z$, a identidade $1_Y$ é o morfismo composto $Y\xrightarrow{\Delta_Y}Y\times_Z Y\xrightarrow{p_1}Y$, mas $\Delta_Y$ não é um morfismo inteiro, como resulta de \sref{II.6.1.10}:
\end{remark}

\begin{proposition}[6.1.10]
\label{II.6.1.10}
Todo morfismo inteiro é universalmente fechado.
\end{proposition}

\begin{proof}
Seja $f:X\to Y$ um morfismo inteiro;
em virtude de \sref{II.6.1.5}[(iii)], basta mostrar que $f$ é \emph{fechado}.
Seja $Z$ uma parte fechada de $X$;
\oldpage[II]{113}
existe um subpré-esquema de $X$ cujo espaço subjacente é $Z$ \sref[I]{I.5.2.1}, e resulta, portanto, de \sref{II.6.1.5}[(i) e (ii)] que podemos nos limitar a provar que $f(X)$ é \emph{fechado} em $Y$.
Em virtude de \sref{II.6.1.5}[(vi)], podemos supor que $X$ e $Y$ sejam \emph{reduzidos};
além disso, se $T$ é o subpré-esquema fechado reduzido de $Y$ cujo espaço subjacente é $\overline{f(X)}$ \sref[I]{I.5.2.1}, sabemos que $f$ se fatora como $X\xrightarrow{g}T\xrightarrow{j}Y$, onde $j$ é o morfismo de injeção \sref[I]{I.5.2.2}, e, como $j$ é separado \sref[I]{I.5.5.1}[(i)], resulta de \sref{II.6.1.5}[(v)] que $g$ é um morfismo inteiro.
Podemos, portanto, supor que $f(X)$ seja \emph{denso} em $Y$.
Por fim, sendo a questão local sobre $Y$, podemos nos limitar ao caso em que $Y=\Spec(A)$.
Então $X=\Spec(B)$, onde $B$ é uma $A$-álgebra inteira sobre $A$ \sref{II.6.1.4};
além disso, $A$ é reduzido \sref[I]{I.5.1.4}, e a hipótese de que $f(X)$ seja denso em $Y$ implica que o homomorfismo $\vphi:A\to B$ correspondente a $f$ é \emph{injetivo} \sref[I]{I.1.2.7}.
Dizer que, nessas condições, $f(X)=Y$ significa que todo ideal primo de $A$ é a interseção com $A$ de um ideal primo de $B$, o que não é senão o primeiro teorema de Cohen--Seidenberg (\cite[t.~I, p.~257, th.~3]{I-13}).
\end{proof}

\begin{corollary}[6.1.11]
\label{II.6.1.11}
Todo morfismo finito $f:X\to Y$ é projetivo.
\end{corollary}

\begin{proof}
Com efeito, como $f$ é afim, $\sh{O}_X$ é um $\sh{O}_X$-módulo \emph{muito amplo} relativamente a $f$ \sref{II.5.1.2};
além disso, $f_*(\sh{O}_X)$ é um $\sh{O}_Y$-módulo quase-coerente \emph{de tipo finito} \sref{II.6.1.2};
por fim, $f$ é separado, de tipo finito e universalmente fechado \sref{II.6.1.10}, e, portanto, satisfaz as condições do critério~\sref{II.5.5.4}[(i)].
\end{proof}

\begin{proposition}[6.1.12]
\label{II.6.1.12}
Seja $f:X'\to X$ um morfismo finito, e seja $\sh{B}=f_*(\sh{O}_{X'})$ (que é uma $\sh{O}_X$-álgebra quase-coerente e um $\sh{O}_X$-módulo de tipo finito).
Seja $\sh{F}'$ um $\sh{O}_{X'}$-módulo quase-coerente;
para que $\sh{F}'$ seja localmente livre de posto~$r$, é necessário e suficiente que $f_*(\sh{F}')$ seja um $\sh{B}$-módulo localmente livre de posto~$r$.
\end{proposition}

\begin{proof}
É claro que, se $f_*(\sh{F}')|U$ é isomorfo a $\sh{B}^r|U$ (sendo $U\subset X$ aberto), então $\sh{F}'|f^{-1}(U)$ é isomorfo a $\sh{O}_{X'}^r|f^{-1}(U)$ \sref{II.1.4.2}.
Reciprocamente, suponhamos que $\sh{F}'$ seja localmente livre de posto~$r$ e mostremos que $f_*(\sh{F}')$ é localmente isomorfo a $\sh{B}^r$ enquanto $\sh{B}$-módulo.
Seja $x$ um ponto de $X$;
quando $U$ percorre um sistema fundamental de vizinhanças afins de $x$, $f^{-1}(U)$ percorre um sistema fundamental de vizinhanças afins \sref{II.1.2.5} do conjunto finito $f^{-1}(x)$, pois $f$ é fechado \sref{II.6.1.10}.
A proposição resulta então do lema seguinte:
\end{proof}

\begin{lemma}[6.1.12.1]
\label{II.6.1.12.1}
Sejam $Y$ um pré-esquema, $\sh{E}$ um $\sh{O}_Y$-módulo localmente livre de posto~$r$, e $Z$ uma parte finita de $Y$ contida em um aberto afim $V$.
Existe então uma vizinhança $U\subset V$ de $Z$ tal que $\sh{E}|U$ seja isomorfo a $\sh{O}_Y^r|U$.
\end{lemma}

\begin{proof}
Podemos evidentemente supor que $Y$ seja afim;
para todo $z_i\in Z$, existe no fecho $\overline{\{z_i\}}$ pelo menos um ponto fechado $z'_i$ \sref[0]{0.2.1.3};
se $Z'$ é o conjunto dos $z'_i$, toda vizinhança de $Z'$ é uma vizinhança de $Z$, e podemos, portanto, supor que $Z$ seja fechado em $Y$ e discreto.
Consideremos o subpré-esquema fechado reduzido de $Y$ cujo espaço subjacente é $Z$ \sref[I]{I.5.2.1}, e seja $j:Z\to Y$ a injeção canônica;
$j^*(\sh{E})=\sh{E}\otimes_Y\sh{O}_Z$ é localmente livre de posto~$r$ sobre o esquema discreto $Z$ e, por conseguinte, isomorfo a $\sh{O}_Z^r$;
em outras palavras, existem $r$ seções $s_i$ (para $1\leq i\leq r$) de $\sh{E}\otimes_Y\sh{O}_Z$ sobre $Z$ tais que o homomorfismo $\sh{O}_Z^r\to\sh{E}\otimes_Y\sh{O}_Z$ definido por essas seções seja bijetivo.
Mas $Y=\Spec(A)$ é afim, $Z$ é definido por um ideal $\mathfrak{J}$ de $A$, e temos $\sh{E}=\widetilde{M}$, onde $M$ é um $A$-módulo;
os $s_i$ são elementos de $M\otimes_A(A/\mathfrak{J})$ e são, portanto, imagens de $r$ elementos $t_i\in M=\Gamma(Y,\sh{E})$.

\oldpage[II]{114}
Para todo $z_j\in Z$, existe então uma vizinhança $V_j$ de $z_j$ tal que as restrições dos $t_i$ a $V_j$ definam um isomorfismo $\sh{O}_Y^r|V_j\to\sh{E}|V_j$ \sref[0]{0.5.5.4};
a vizinhança $U$ dada pela união dos $V_j$ é a vizinhança procurada.
\end{proof}

\begin{proposition}[6.1.13]
\label{II.6.1.13}
Sejam $g:X'\to X$ um morfismo inteiro de pré-esquemas, $Y$ um pré-esquema normal localmente integral, e $f$ uma aplicação racional de $Y$ em $X'$ tal que $g\circ f$ seja uma aplicação racional definida em toda parte \sref[I]{I.7.2.1}.
Então $f$ está definida em toda parte.
\end{proposition}

\begin{proof}
Se $f_1$ e $f_2$ são dois morfismos (de abertos densos de $Y$ em $X'$) da classe $f$, é claro que $g\circ f_1$ e $g\circ f_2$ são morfismos equivalentes, o que justifica a notação $g\circ f$ para sua classe de equivalência.
Recordemos também que, quando se supõe, além disso, que $Y$ é localmente noetheriano, a hipótese de que $Y$ seja normal já implica que $Y$ é localmente integral \sref[I]{I.6.1.13}.

Para demonstrar \sref{II.6.1.13}, notemos primeiro que a questão é local sobre $Y$ e podemos, por conseguinte, supor que existe na classe de $g\circ f$ um \emph{morfismo} $h:Y\to X$.
Consideremos a imagem inversa $Y'=X'_{(h)}=X'_{(Y)}$, e notemos que o morfismo $g'=g_{(Y)}:Y'\to Y$ é \emph{inteiro} \sref{II.6.1.5}[(iii)].
Em vista da correspondência entre aplicações racionais de $Y$ em $X'$ e $Y$-seções racionais de $Y'$ \sref[I]{I.7.1.2}, vemos que nos reduzimos a provar o caso particular de \sref{II.6.1.13} em que $X=Y$; em outras palavras, o seguinte:
\end{proof}

\begin{corollary}[6.1.14]
\label{II.6.1.14}
Sejam $X$ um pré-esquema normal localmente integral, $g:X'\to X$ um morfismo \emph{inteiro}, e $f$ uma $X$-seção racional de $X'$.
Então $f$ está definida em toda parte.
\end{corollary}

\begin{proof}
Sendo a questão local sobre $X$, podemos supor que $X$ seja integral, e $f$ identifica-se então com um morfismo de um aberto $U$ de $X$ em $X'$ \sref[I]{I.7.2.2} que é uma $U$-seção de $g^{-1}(U)$.
Como $g$ é separado, $f$ é uma imersão fechada de $U$ em $g^{-1}(U)$ \sref[I]{I.5.4.6};
seja $Z$ o subpré-esquema fechado de $g^{-1}(U)$ associado a $f$ \sref[I]{I.4.2.1}, que é isomorfo a $U$ e, portanto, integral;
seja $X_1$ o subpré-esquema reduzido de $X'$ cujo espaço subjacente é o fecho $\overline{Z}$ de $Z$ em $X'$ \sref[I]{I.5.2.1};
$Z$ é então um subpré-esquema induzido sobre um aberto de $X_1$ \sref[I]{I.5.2.3} e, como é irredutível, o mesmo vale para $X_1$, que é, por conseguinte, integral.
O morfismo $f$ pode então ser considerado como uma $X$-seção racional de $X_1$;
como a restrição de $g$ a $X_1$ é um morfismo inteiro \sref{II.6.1.5}[(i) e (ii)], reduzimo-nos finalmente a provar \sref{II.6.1.14} no caso particular em que $X'=X_1$; em outras palavras, o seguinte:
\end{proof}

\begin{corollary}[6.1.15]
\label{II.6.1.15}
Sejam $X$ um pré-esquema integral normal, $X'$ um pré-esquema integral, e $g:X'\to X$ um morfismo \emph{inteiro}.
Se existe uma $X$-seção racional $f$ de $X'$, então $g$ é um isomorfismo.
\end{corollary}

\begin{proof}
Sendo a questão local sobre $X$, podemos supor que $X$ seja afim, de anel íntegro $A$, e então $X'$ é afim, de anel íntegro $A'$ inteiro sobre $A$ \sref{II.6.1.4};
além disso, o raciocínio de \sref{II.6.1.14} mostra que existe um aberto denso de $X$ isomorfo a um aberto denso de $X'$, de modo que $A$ e $A'$ têm o mesmo corpo de frações.
Por outro lado, em virtude de \sref[I]{I.8.2.1.1} e da hipótese de que os $\sh{O}_x$ sejam integralmente fechados, o anel $A$ é integralmente fechado;
portanto, $A'=A$, o que conclui a demonstração de \sref{II.6.1.13}.
\end{proof}

\subsection{Morfismos quase-finitos}
\label{subsection:II.6.2}

\begin{proposition}[6.2.1]
\label{II.6.2.1}
Sejam $f:X\to Y$ um morfismo localmente de tipo finito e $x$ um ponto de $X$.
As condições seguintes são equivalentes:
\oldpage[II]{115}
\begin{enumerate}
  \item[\rm{(a)}] O ponto $x$ é isolado em sua fibra $f^{-1}(f(x))$.
  \item[\rm{(b)}] O anel $\sh{O}_x$ é um $\sh{O}_{f(x)}$-módulo quase-finito \sref[0]{0.7.4.1}.
\end{enumerate}
\end{proposition}

\begin{proof}
Sendo a questão evidentemente local em $X$ e em $Y$, podemos supor que $X=\Spec(A)$ e $Y=\Spec(B)$ sejam afins, sendo $A$ uma $B$-álgebra de tipo finito \sref[I]{I.6.3.3}.
Além disso, podemos substituir $X$ por $X\times_Y\Spec(\sh{O}_{f(x)})$ sem alterar nem a fibra $f^{-1}(f(x))$ nem o anel local $\sh{O}_x$ \sref[I]{I.3.6.5};
podemos, portanto, supor que $B$ seja um anel local (igual a $\sh{O}_{f(x)}$).
Se $\mathfrak{n}$ é o ideal maximal de $B$, então $f^{-1}(f(x))$ é um esquema afim de anel $A/\mathfrak{n}A$, de tipo finito sobre $\kres(f(x))=B/\mathfrak{n}$ \sref[I]{I.6.4.11}.
Posto isso, se (a) é verificada, podemos ainda supor que $f^{-1}(f(x))$ se reduza ao ponto $x$;
logo, $A/\mathfrak{n}A$ tem posto finito sobre $B/\mathfrak{n}$ \sref[I]{I.6.4.4}, ou seja, $A$ é um $B$-módulo quase-finito.
Reciprocamente, se assim for, $f^{-1}(f(x))$ é um esquema afim artiniano, portanto discreto \sref[I]{I.6.4.4};
por conseguinte, $x$ é isolado em sua fibra, e isto mostra que (b) implica (a).
\end{proof}

\begin{corollary}[6.2.2]
\label{II.6.2.2}
Seja $f:X\to Y$ um morfismo de tipo finito.
As condições seguintes são equivalentes:
\begin{enumerate}
  \item[\rm{(a)}] Todo ponto $x\in X$ é isolado em sua fibra $f^{-1}(f(x))$ (ou seja, o subespaço $f^{-1}(f(x))$ é \emph{discreto}).
  \item[\rm{(b)}] Para todo $x\in X$, o pré-esquema $f^{-1}(f(x))$ é um $\kres(f(x))$-pré-esquema finito.
  \item[\rm{(c)}] Para todo $x\in X$, o anel $\sh{O}_x$ é um $\sh{O}_{f(x)}$-módulo quase-finito.
\end{enumerate}
\end{corollary}

\begin{proof}
A equivalência entre (a) e (c) resulta de \sref{II.6.2.1}.
Por outro lado, como $f^{-1}(f(x))$ é um $\kres(f(x))$-pré-esquema algébrico \sref[I]{I.6.4.11}, a equivalência entre (a) e (b) resulta de \sref[I]{I.6.4.4}.
\end{proof}

\begin{definition}[6.2.3]
\label{II.6.2.3}
Quando $f:X\to Y$ é um morfismo de tipo finito que verifica as condições equivalentes de \sref{II.6.2.2}, diz-se que $f$ é \emph{quase-finito}, ou que $X$ é \emph{quase-finito sobre $Y$}.
\end{definition}

É claro que todo morfismo \emph{finito} é quase-finito \sref{II.6.1.7}.

\begin{proposition}[6.2.4]
\label{II.6.2.4}
\medskip\noindent
\begin{enumerate}
  \item[\rm{(i)}] Uma imersão $X\to Y$ que seja fechada, ou tal que $X$ seja noetheriano, é um morfismo quase-finito.
  \item[\rm{(ii)}] Se $f:X\to Y$ e $g:Y\to Z$ são morfismos quase-finitos, então $g\circ f$ é quase-finito.
  \item[\rm{(iii)}] Se $X$ e $Y$ são $S$-pré-esquemas e $f:X\to Y$ é um $S$-morfismo quase-finito, então $f_{(S')}:X_{(S')}\to Y_{(S')}$ é quase-finito para toda extensão da base $g:S'\to S$.
  \item[\rm{(iv)}] Se $f:X\to Y$ e $g:X'\to Y'$ são dois $S$-morfismos quase-finitos, então
    \[
      f\times_S g:X\times_S X'\to Y\times_S Y'
    \]
    é quase-finito.
  \item[\rm{(v)}] Sejam $f:X\to Y$ e $g:Y\to Z$ dois morfismos tais que $g\circ f$ seja quase-finito; então, se, além disso, $g$ é separado, ou $X$ é noetheriano, ou $X\times_ZY$ é localmente noetheriano, $f$ é quase-finito.
  \item[\rm{(vi)}] Se $f$ é quase-finito, o mesmo vale para $f_\red$.
\end{enumerate}
\end{proposition}

\begin{proof}
Se $f:X\to Y$ é uma imersão, toda fibra tem no máximo um ponto, e a afirmação~(i) resulta, portanto, de \sref[I]{I.6.3.4}[(i)] e \sref[I]{I.6.3.5}.
Para provar (ii), observemos primeiro que $h=g\circ f$ é de tipo finito \sref[I]{I.6.3.4}[(ii)];
além disso, se $z=h(x)$ e $y=f(x)$, então $y$ é isolado em $g^{-1}(z)$, logo existe uma vizinhança aberta $V$ de $y$ em $Y$ que não encontra nenhum ponto de $g^{-1}(z)$ distinto de $y$;
$f^{-1}(V)$ é, portanto, uma vizinhança aberta de $x$ que não encontra nenhum

\oldpage[II]{116}
dos $f^{-1}(y')$, em que $y'\neq y$ pertence a $g^{-1}(z)$;
como $x$ é isolado em $f^{-1}(y)$, ele é isolado em $h^{-1}(z)=f^{-1}(g^{-1}(z))$.
Para provar (iii), podemos limitar-nos ao caso em que $Y=S$ \sref[I]{I.3.3.11};
observemos ainda, antes de tudo, que $f'=f_{(S')}$ é de tipo finito \sref[I]{I.6.3.4}[(iii)];
por outro lado, se $x'\in X'=X_{(S')}$ e se pusermos $y'=f'(x')$ e $y=g(y')$, então ${f'}^{-1}(y')$ identifica-se com $f^{-1}(y)\otimes_{\kres(y)}\kres(y')$ \sref[I]{I.3.6.5};
como $f^{-1}(y)$ tem posto finito sobre $\kres(y)$ por hipótese, ${f'}^{-1}(y')$ tem posto finito sobre $\kres(y')$, portanto é discreto.
As afirmações (iv), (v) e (vi) deduzem-se de (i), (ii) e (iii) pelo método geral \sref[I]{I.5.5.12}, salvo quando as hipóteses de (v) são distintas da hipótese ``$g$ separado'';
para tratar esse caso, observemos primeiro que, se $x$ é isolado em $f^{-1}(g^{-1}(g(f(x))))$, ele o é \emph{a fortiori} em $f^{-1}(f(x))$;
o fato de $f$ ser de tipo finito resulta então de \sref[I]{I.6.3.6}.
\end{proof}

\begin{proposition}[6.2.5]
\label{II.6.2.5}
Sejam $A$ um anel local noetheriano \emph{completo}, $X$ um $A$-esquema localmente de tipo finito e $x$ um ponto de $X$ sobre o ponto fechado $y$ de $Y=\Spec(A)$.
Suponha-se que $x$ seja isolado em sua fibra $f^{-1}(y)$ (sendo $f$ o morfismo estrutural $X\to Y$).
Então $\sh{O}_x$ é um $A$-módulo de tipo finito, e $X$ é $Y$-isomorfo à soma \sref[I]{I.3.1} de $X'=\Spec(\sh{O}_x)$ (que é um $Y$-esquema finito) e de um $A$-esquema $X''$.
\end{proposition}

\begin{proof}
Resulta de \sref{II.6.2.1} que $\sh{O}_x$ é um $A$-módulo quase-finito.
Como $\sh{O}_x$ é noetheriano \sref[I]{I.6.3.7} e o homomorfismo $A\to\sh{O}_x$ é local, a hipótese de que $A$ é \emph{completo} implica que $\sh{O}_x$ é um $A$-módulo \emph{de tipo finito} \sref[0]{0.7.4.3}.
Seja $X'=\Spec(\sh{O}_x)$ o esquema local de $X$ no ponto $x$ \sref[I]{I.2.4.1}, e seja $g:X'\to X$ o morfismo canônico.
Como o composto $X'\xrightarrow{g}X\xrightarrow{f}Y$ é finito \sref{II.6.1.1} e $f$ é separado, $g$ é finito \sref{II.6.1.5}[(v)], logo $g(X')$ é \emph{fechado} em $X$ \sref{II.6.1.10};
por outro lado, como $g$ é de tipo finito, $g$ é um isomorfismo local no ponto fechado $x'$ de $X'$ em virtude da definição de $g$ e de \sref[I]{I.6.5.4};
mas, como $X'$ é a única vizinhança aberta de $x'$, isso significa que $g$ é uma imersão aberta, logo $g(X')$ também é \emph{aberto} em $X$, o que termina a demonstração.
\end{proof}

\begin{corollary}[6.2.6]
\label{II.6.2.6}
Sejam $A$ um anel local noetheriano \emph{completo}, $Y=\Spec(A)$ e $f:X\to Y$ um morfismo separado e quase-finito.
Então $X$ é $Y$-isomorfo a uma soma $X'\sqcup X''$, em que $X'$ é um $Y$-esquema finito e $X''$ um $Y$-esquema quase-finito tal que, se $y$ é o ponto fechado de $Y$, então $X''\cap f^{-1}(y)=\emp$.
\end{corollary}

\begin{proof}
Com efeito, a fibra $f^{-1}(y)$ é finita e discreta por hipótese, e o corolário resulta, portanto, por indução sobre o número de pontos dessa fibra, de \sref{II.6.2.5}.
\end{proof}

\begin{remark}[6.2.7]
\label{II.6.2.7}
No Capítulo~V, veremos que, se $Y$ é \emph{localmente noetheriano}, um morfismo quase-finito separado $X\to Y$ é necessariamente \emph{quase-afim}.
\end{remark}

\subsection{Fecho integral de um pré-esquema}
\label{subsection:II.6.3}

\begin{proposition}[6.3.1]
\label{II.6.3.1}
Sejam $(X,\sh{A})$ um espaço anelado, $\sh{B}$ uma $\sh{A}$-álgebra (comutativa) e $f$ uma seção de $\sh{B}$ sobre $X$.
As propriedades seguintes são equivalentes:
\begin{enumerate}
  \item[\rm{(a)}] O sub-$\sh{A}$-módulo de $\sh{B}$ gerado pelos $f^n$, para $n\geq0$ \sref[0]{0.5.1.1}, é de tipo finito.
  \item[\rm{(b)}] Existe uma sub-$\sh{A}$-álgebra $\sh{C}$ de $\sh{B}$ que é um $\sh{A}$-módulo de tipo finito e tal que $f\in\Gamma(X,\sh{C})$.
  \item[\rm{(c)}] Para todo $x\in X$, $f_x$ é inteiro sobre o talo $\sh{A}_x$.
\end{enumerate}
\end{proposition}

\oldpage[II]{117}
\begin{proof}
Como o sub-$\sh{A}$-módulo de $\sh{B}$ gerado pelos $f^n$ é uma $\sh{A}$-álgebra, é claro que (a) implica (b).
Por outro lado, (b) implica que, para todo $x\in X$, o $\sh{A}_x$-módulo $\sh{C}_x$ é de tipo finito, o que implica que todo elemento da álgebra $\sh{C}_x$, e em particular $f_x$, é inteiro sobre $\sh{A}_x$.
Finalmente, se, para algum $x\in X$, temos uma relação da forma
\[
  f_x^n+(a_1)_xf_x^{n-1}+\ldots+(a_n)_x=0
\]
em que os $a_i$ ($1\leq i\leq n$) são seções de $\sh{A}$ sobre uma vizinhança aberta $U$ de $x$, a seção $f^n|U+a_1\cdot f^{n-1}|U+\ldots+a_n$ é nula sobre uma vizinhança $V\subset U$ de $x$;
daí resulta imediatamente que todos os $f^k|V$ ($k\geq0$) são combinações lineares, com coeficientes em $\Gamma(V,\sh{A})$, dos $f^j|V$ para $0\leq j\leq n-1$;
concluímos que (c) implica (a).
\end{proof}

Quando as condições equivalentes de \sref{II.6.3.1} são satisfeitas, diz-se que a seção $f$ é \emph{inteira sobre $\sh{A}$}.

\begin{corollary}[6.3.2]
\label{II.6.3.2}
Sob as hipóteses de \sref{II.6.3.1}, existe um único sub-$\sh{A}$-módulo $\sh{A}'$ de $\sh{B}$ tal que, para todo $x\in X$, $\sh{A}'_x$ seja o conjunto dos germes $f_x\in\sh{B}_x$ que são inteiros sobre $\sh{A}_x$.
Para todo aberto $U\subset X$, as seções de $\sh{A}'$ sobre $U$ são as seções $f\in\Gamma(U,\sh{B})$ que são inteiras sobre $\sh{A}|U$.
\end{corollary}

\begin{proof}
A existência de $\sh{A}'$ é imediata, tomando para $\Gamma(U,\sh{A}')$ o conjunto dos $f\in\Gamma(U,\sh{B})$ tais que $f_x$ seja inteiro sobre $\sh{A}_x$ para todo $x\in U$.
A segunda afirmação resulta imediatamente de \sref{II.6.3.1}.
\end{proof}

É claro que $\sh{A}'$ é uma sub-$\sh{A}$-álgebra de $\sh{B}$;
diz-se que ela é o \emph{fecho integral de $\sh{A}$ em $\sh{B}$}.

\begin{env}[6.3.3]
\label{II.6.3.3}
Sejam $(X,\sh{A})$ e $(Y,\sh{B})$ dois espaços anelados, e seja
\[
  g=(\psi,\theta):X\to Y
\]
um morfismo.
Seja $\sh{C}$ (resp. $\sh{D}$) uma $\sh{A}$-álgebra (resp. uma $\sh{B}$-álgebra), e seja
\[
  u:\sh{D}\to\sh{C}
\]
um $g$-morfismo \sref[0]{0.4.4.1}.
Então, se $\sh{A}'$ (resp. $\sh{B}'$) é o fecho integral de $\sh{A}$ (resp. $\sh{B}$) em $\sh{C}$ (resp. $\sh{D}$), a \emph{restrição} de $u$ a $\sh{B}'$ é um $g$-morfismo
\[
  u':\sh{B}'\to\sh{A}'.
\]

Com efeito, se $j$ é a injeção canônica $\sh{B}'\to\sh{D}$, basta mostrar que
\[
  v=u^\sharp\circ g^*(j):g^*(\sh{B}')\to\sh{C}
\]
aplica $g^*(\sh{B}')$ em $\sh{A}'$.
Ora, um elemento de $g^*(\sh{B}')_x=\sh{B}'_{\psi(x)}\otimes_{\sh{B}_{\psi(x)}}\sh{A}_x$ é inteiro sobre $\sh{A}_x$ pela definição de $\sh{B}'$, logo o mesmo vale para sua imagem por $v_x$, o que estabelece nossa afirmação.
\end{env}

\begin{proposition}[6.3.4]
\label{II.6.3.4}
Sejam $X$ um pré-esquema e $\sh{A}$ uma $\sh{O}_X$-álgebra quase-coerente.
Então o fecho integral $\sh{O}'_X$ de $\sh{O}_X$ em $\sh{A}$ é uma $\sh{O}_X$-álgebra quase-coerente e, para todo aberto afim $U$ de $X$, $\Gamma(U,\sh{O}'_X)$ é o fecho integral de $\Gamma(U,\sh{O}_X)$ em $\Gamma(U,\sh{A})$.
\end{proposition}

\begin{proof}
Podemos limitar-nos ao caso em que $X=\Spec(B)$ é afim e $\sh{A}=\widetilde{A}$, sendo $A$ uma
\oldpage[II]{118}
$B$-álgebra; seja $B'$ o fecho integral de $B$ em $A$.
Tudo se reduz a ver que, para todo $x\in X$, um elemento de $A_x$ inteiro sobre $B_x$ pertence necessariamente a $B'_x$, o que resulta do fato de que, para um anel comutativo $C$, as operações de fecho integral em uma $C$-álgebra e de passagem a um anel de frações (com respeito a uma parte multiplicativa de $C$) comutam \cite[t.~I, pp.~261 e 257]{I-13}.
\end{proof}

O $X$-esquema $X'=\Spec(\sh{O}'_X)$ é então chamado \emph{fecho integral de $X$ relativamente a $\sh{A}$} (ou \emph{relativamente a $\Spec(\sh{A})$});
é claro que $X'$ é \emph{inteiro} sobre $X$ \sref{II.6.1.2}.

Deduz-se imediatamente de \sref{II.6.3.4} que, se $f:X'\to X$ é o morfismo estrutural, então, para todo aberto $U$ de $X$, $f^{-1}(U)$ é o \emph{fecho integral do pré-esquema induzido por $X$ sobre $U$, relativamente a $\sh{A}|U$}.

\begin{env}[6.3.5]
\label{II.6.3.5}
Sejam $X$ e $Y$ dois pré-esquemas, $f:X\to Y$ um morfismo, $\sh{A}$ (resp. $\sh{B}$) uma $\sh{O}_X$-álgebra (resp. uma $\sh{O}_Y$-álgebra) quase-coerente, e seja $u:\sh{B}\to\sh{A}$ um $f$-morfismo.
Vimos em \sref{II.6.3.3} que daí se deduz um $f$-morfismo $u':\sh{O}'_Y\to\sh{O}'_X$, em que $\sh{O}'_X$ (resp. $\sh{O}'_Y$) é o fecho integral de $\sh{O}_X$ (resp. $\sh{O}_Y$) em $\sh{A}$ (resp. $\sh{B}$).
Por conseguinte, se $X'$ (resp. $Y'$) é o fecho integral de $X$ (resp. $Y$) relativamente a $\sh{A}$ (resp. $\sh{B}$), deduz-se canonicamente de $u$ um morfismo $f'=\Spec(u'):X'\to Y'$ \sref{II.1.5.6} que torna comutativo o diagrama
\[
\label{II.6.3.5.1}
  \xymatrix{
    X' \ar[r]^{f'} \ar[d]
    & Y' \ar[d]
  \\X \ar[r]_{f}
    & Y
  }
\tag{6.3.5.1}
\]
\end{env}

\begin{env}[6.3.6]
\label{II.6.3.6}
Suponhamos que $X$ tenha apenas um número \emph{finito} de componentes irredutíveis $X_i$ ($1\leq i\leq r$), de pontos genéricos $\xi_i$, e consideremos, em particular, o fecho integral de $X$ relativamente a uma $\sh{R}(X)$-álgebra quase-coerente $\sh{A}$ (como $\sh{O}_X$-álgebra ou como $\sh{R}(X)$-álgebra, o que dá no mesmo).
Sabemos \sref[I]{I.7.3.5} que $\sh{A}$ é soma direta de $r$ $\sh{O}_X$-álgebras quase-coerentes $\sh{A}_i$, estando o suporte de $\sh{A}_i$ contido em $X_i$, e sendo o feixe induzido por $\sh{A}_i$ sobre $X_i$ um feixe simples cuja fibra $A_i$ é uma álgebra sobre $\sh{O}_{\xi_i}$.
É claro então, por \sref{II.6.3.4}, que o fecho integral $\sh{O}'_X$ de $\sh{O}_X$ em $\sh{A}$ é a \emph{soma direta} dos fechos integrais $\sh{O}_X^{(i)}$ de $\sh{O}_X$ em cada uma das $\sh{A}_i$ e, por conseguinte, o fecho integral $X'=\Spec(\sh{O}'_X)$ de $X$ relativamente a $\sh{A}$ é um $X$-esquema soma dos $\Spec(\sh{O}_X^{(i)})=X'_i$ ($1\leq i\leq r$).

Suponhamos, além disso, que a $\sh{O}_X$-álgebra $\sh{A}$ seja \emph{reduzida}, ou, o que dá no mesmo, que cada uma das álgebras $A_i$ seja reduzida e possa, por conseguinte, ser considerada uma álgebra sobre o \emph{corpo} $\kres(\xi_i)$ (igual ao corpo das funções racionais do subpré-esquema reduzido de $X$ cujo espaço subjacente é $X_i$);
então, por \sref{II.1.3.8}, cada um dos $X'_i$ é um $X$-esquema \emph{reduzido}, e $X'$ é também o fecho integral de $X_\red$.
Suponhamos ainda que cada uma das álgebras $A_i$ seja \emph{soma direta de um número finito de corpos $K_{ij}$} ($1\leq j\leq s_i$);
se $\sh{K}_{ij}$ é a subálgebra de $\sh{A}_i$ correspondente a $K_{ij}$, é claro que $\sh{O}_X^{(i)}$ é a \emph{soma direta} dos fechos integrais $\sh{O}_X^{(ij)}$ de $\sh{O}_X$ em cada uma das $\sh{K}_{ij}$.
Por conseguinte, $X'_i$ é então soma dos $X$-esquemas $X'_{ij}=\Spec(\sh{O}_X^{(ij)})$ ($1\leq j\leq s_i$).
Além disso, sob essas hipóteses e com essas notações:
\end{env}

\oldpage[II]{119}
\begin{proposition}[6.3.7]
\label{II.6.3.7}
Cada um dos $X'_{ij}$ é um $X$-pré-esquema integral e normal, e seu corpo de funções racionais $R(X'_{ij})$ identifica-se canonicamente com o fecho algébrico $K'_{ij}$ de $\kres(\xi_i)$ em $K_{ij}$.
\end{proposition}

\begin{proof}
Em virtude do que precede, podemos supor $X$ integral, logo $r=1$ e $s_1=1$, de modo que a única álgebra $A_1$ seja um corpo $K$;
seja $\xi$ o ponto genérico de $X$, e seja $f:X'\to X$ o morfismo estrutural.
Para todo aberto afim não vazio $U$ de $X$, $f^{-1}(U)$ identifica-se com o espectro do fecho integral $B'_U$ do anel integral $B_U=\Gamma(U,\sh{O}_X)$ no corpo $K$ \sref{II.6.3.4};
como o anel $B'_U$ é integral e integralmente fechado, o mesmo vale para os anéis locais dos pontos de seu espectro, logo $f^{-1}(U)$ é, por definição, um esquema \emph{integral} e \emph{normal} (\sref[0]{0.4.1.4} e \sref[I]{I.5.1.4}).
Além disso, como $(0)$ é o único ideal primo de $B'_U$ acima do ideal primo $(0)$ de $B_U$ \cite[t.~I, p.~259]{I-13}, $f^{-1}(\xi)$ reduz-se a um único ponto $\xi'$, e $\kres(\xi')$ é o corpo de frações $K'$ de $B'_U$, que não é senão o fecho algébrico de $\kres(\xi)$ em $K$.
Por fim, $X'$ é irredutível, pois, quando $U$ percorre o conjunto dos abertos afins não vazios de $X$, os $f^{-1}(U)$ constituem um recobrimento aberto de $X'$ formado por abertos irredutíveis;
além disso, a interseção $f^{-1}(U\cap V)$ de dois desses abertos contém $\xi'$, logo não é vazia, e concluímos por \sref[0]{0.2.1.4}.
\end{proof}

\begin{corollary}[6.3.8]
\label{II.6.3.8}
Seja $X$ um pré-esquema reduzido que tenha apenas um número finito de componentes irredutíveis $X_i$ ($1\leq i\leq r$), e seja $\xi_i$ o ponto genérico de $X_i$.
O fecho integral $X'$ de $X$ relativamente a $\sh{R}(X)$ é soma de $r$ $X$-esquemas $X'_i$ que são integrais e normais.
Se $f:X'\to X$ é o morfismo estrutural, $f^{-1}(\xi_i)$ reduz-se ao ponto genérico $\xi'_i$ de $X'_i$, e temos $\kres(\xi'_i)=\kres(\xi_i)$; em outras palavras, $f$ é birracional.
\end{corollary}

Diz-se, nesse caso, que $X'$ é a \emph{normalização} do pré-esquema reduzido $X$;
observemos que $f$, sendo birracional e inteiro, é \emph{sobrejetivo} \sref{II.6.1.10}.
Para que se tenha então $X'=X$, é necessário e suficiente que $X$ seja \emph{normal}.
Quando $X$ é um pré-esquema integral, resulta de \sref{II.6.3.8} que sua normalização $X'$ é integral.

\begin{env}[6.3.9]
\label{II.6.3.9}
Sejam $X$ e $Y$ dois pré-esquemas integrais, $f:X\to Y$ um morfismo dominante, e sejam $L=R(X)$ e $K=R(Y)$ os corpos das funções racionais de $X$ e $Y$;
a $f$ corresponde canonicamente uma injeção $K\to L$ e, quando identificamos $K$ (resp. $L$) com o feixe simples $\sh{R}(Y)$ (resp. $\sh{R}(X)$), essa injeção é um $f$-morfismo.
Seja $K_1$ (resp. $L_1$) uma extensão de $K$ (resp. $L$), e suponhamos dado um monomorfismo $K_1\to L_1$ tal que o diagrama
\[
  \xymatrix{
    K_1 \ar[r]
    & L_1
  \\K \ar[u] \ar[r]
    & L \ar[u]
  }
\]
seja comutativo;
se $K_1$ (resp. $L_1$) é considerado um feixe simples sobre $Y$ (resp. $X$), portanto uma $\sh{R}(Y)$-álgebra (resp. uma $\sh{R}(X)$-álgebra), isso significa que $K_1\to L_1$ é um $f$-morfismo.
Posto isso, se $X'$ (resp. $Y'$) é o fecho integral de $X$ (resp. $Y$) relativamente a $L_1$ (resp. $K_1$), $X'$ (resp. $Y'$) é um pré-esquema integral e normal \sref{II.6.3.7} cujo corpo de funções racionais se identifica canonicamente com o fecho algébrico

\oldpage[II]{120}
$L'$ (resp. $K'$) de $L$ (resp. $K$) em $L_1$ (resp. $K_1$), e existe um morfismo canônico (necessariamente dominante) $f':X'\to Y'$ que torna comutativo o diagrama \sref{II.6.3.5.1}.
O caso mais importante é aquele em que se toma $L_1=L$, sendo então $K_1$ uma extensão de $K$ contida em $L$, e em que se supõe $X$ integral e \emph{normal}, de modo que $X'=X$.
O que precede mostra então que, quando $X$ é normal e $Y'$ é o fecho integral de $Y$ relativamente a um corpo $K_1\subset L=R(X)$, todo morfismo dominante $f:X\to Y$ se fatora como
\[
  f:X\xrightarrow{f'}Y'\to Y
\]
onde $f'$ é dominante;
além disso, uma vez dado o monomorfismo $K_1\to L$, $f'$ é necessariamente único, como se vê reduzindo-se ao caso em que $X$ e $Y$ são afins.
Vê-se, portanto, que, para os dados de $Y$, $L$ e de um $K$-monomorfismo $K_1\to L$, o fecho integral $Y'$ de $Y$ relativamente a $K_1$ é solução de um \emph{problema universal}.
\end{env}

\begin{remark}[6.3.10]
\label{II.6.3.10}
Retomemos as hipóteses de \sref{II.6.3.6}, supondo ainda que cada uma das álgebras $A_i$ tenha \emph{posto finito} sobre $\kres(\xi_i)$ (o que implica que $A_i$ é soma direta de um número finito de corpos);
pode-se então afirmar, em certos casos, que o morfismo estrutural $X'\to X$ não é apenas inteiro, mas até \emph{finito}.
Restrinjamo-nos ao caso em que $X$ é \emph{reduzido};
como a questão é local sobre $X$, pode-se supor ainda que $X$ é afim, de anel $C$, e que $C$ tem apenas um número finito de ideais minimais $\mathfrak{p}_i$ ($1\leq i\leq r$), tais que os $C_i=C/\mathfrak{p}_i$ sejam íntegros;
então $X'$ será finito sobre $X$ se o fecho integral de cada $C_i$ em toda extensão de grau finito de seu corpo de frações for um \emph{$C$-módulo de tipo finito} \sref{II.6.3.4}.
Sabe-se que essa condição é sempre satisfeita quando $C$ é uma \emph{álgebra de tipo finito sobre um corpo} \cite[t.~I, p.~267, th.~9]{I-13}, ou \emph{sobre $\bb{Z}$} \cite[I,p.~93, th.~3]{I-9}, ou \emph{sobre um anel local noetheriano completo} \cite[p.~298]{II-25}.
Conclui-se que $X'\to X$ será um morfismo finito quando $X$ for um esquema \emph{de tipo finito} sobre um corpo, sobre $\bb{Z}$ ou sobre um anel local noetheriano completo.
\end{remark}

\subsection{Determinante de um endomorfismo de $\mathcal{O}_X$-módulos}
\label{subsection:II.6.4}

\begin{env}[6.4.1]
\label{II.6.4.1}
Sejam $A$ um anel (comutativo), $E$ um $A$-módulo livre de posto~$n$ e $u$ um endomorfismo de $E$;
recordemos que, para definir o \emph{polinômio característico} de $u$, considera-se o endomorfismo $u\otimes1$ do $A[T]$-módulo livre de posto~$n$, $E\otimes_A A[T]$ (sendo $T$ uma indeterminada), e põe-se
\[
\label{II.6.4.1.1}
  P(u,T) = \det(T\cdot I-(u\otimes 1))
\tag{6.4.1.1}
\]
(sendo $I$ o automorfismo identidade de $E\otimes_A A[T]$).
Tem-se
\[
\label{II.6.4.1.2}
  P(u,T) = T^n - \sigma_1(u)T^{n-1} + \ldots + (-1)^n\sigma_n(u)
\tag{6.4.1.2}
\]
onde $\sigma_i(u)$ é um elemento de $A$, igual a um polinômio homogêneo de grau~$i$ (com coeficientes inteiros) em relação aos elementos da matriz de $u$ relativa a uma base qualquer de $E$;
diz-se que os $\sigma_i(u)$ são as \emph{funções simétricas elementares} de $u$;
em particular, tem-se $\sigma_1(u)=\Tr u$ e $\sigma_n(u)=\det u$.
Recordemos que, em virtude do teorema de Hamilton--Cayley, tem-se
\[
\label{II.6.4.1.3}
  P(u,u) = u^n - \sigma_1(u)u^{n-1} + \ldots + (-1)^n\sigma_n(u) = 0
\tag{6.4.1.3}
\]
\oldpage[II]{121}
Isso também se escreve
\[
\label{II.6.4.1.4}
  (\det u)\cdot I_E = uQ(u)
\tag{6.4.1.4}
\]
(sendo $I_E$ o automorfismo identidade de $E$), com
\[
\label{II.6.4.1.5}
  Q(u) = (-1)^{n+1}(u^{n-1}-\sigma_1(u)u^{n-2}+\ldots+(-1)^{n-1}\sigma_{n-1}(u)).
\tag{6.4.1.5}
\]

Seja $\vphi:A\to B$ um homomorfismo de anéis que faz de $B$ uma $A$-álgebra;
consideremos o $B$-módulo $E_{(B)}=E\otimes_A B$, que é livre de posto~$n$, e a extensão $u\otimes1$ de $u$ a um endomorfismo de $E_{(B)}$;
é imediato que $\sigma_i(u\otimes1)=\vphi(\sigma_i(u))$ para todo índice $i$.
\end{env}

\begin{env}[6.4.2]
\label{II.6.4.2}
Suponhamos agora que $A$ seja um anel \emph{íntegro}, de corpo de frações $K$, e seja $E$ um $A$-módulo \emph{de tipo finito} (mas não necessariamente livre).
Seja $n$ o \emph{posto} de $E$, isto é, o posto do $K$-módulo livre $E\otimes_A K$;
a todo endomorfismo $u$ de $E$ corresponde canonicamente o endomorfismo $u\otimes1$ de $E\otimes_A K$;
por abuso de linguagem, chama-se ainda \emph{polinômio característico} de $u$, e denota-se por $P(u,T)$, o polinômio $P(u\otimes1,T)$, cujos coeficientes $\sigma_i(u\otimes1)$ (que pertencem a $K$) são ainda denotados por $\sigma_i(u)$ e chamados \emph{funções simétricas elementares de $u$};
em particular, $\det u=\det(u\otimes1)$ por definição.
Com essas notações, as equações~\sref{II.6.4.1.3} a \sref{II.6.4.1.5} têm sentido e continuam válidas, interpretando-se $u^j$ como o composto
\[
  E\xrightarrow{x\mapsto x\otimes1}E\otimes_A K
   \xrightarrow{(u\otimes1)^j}E\otimes_A K.
\]

Se $F$ é o submódulo de torção de $E$ e $E_0=E/F$, então $u(F)\subset F$ e, passando aos quocientes, $u$ induz um endomorfismo $u_0$ de $E_0$;
além disso, $E\otimes_A K$ identifica-se com $E_0\otimes_A K$, e $u\otimes1$ com $u_0\otimes1$, de modo que $\sigma_i(u)=\sigma_i(u_0)$ para $1\leq i\leq n$.

Se $E$ é \emph{sem torção}, então $E$ identifica-se canonicamente com um sub-$A$-módulo de $E\otimes_A K$, e a relação $u\otimes1=0$ equivale a $u=0$.
Se $E$ é um $A$-módulo livre, as duas definições de $\sigma_i(u)$ dadas em \sref{II.6.4.1} e \sref{II.6.4.2} coincidem pelo que precede, o que justifica as notações adotadas.

Observemos também que, quando $E$ é um \emph{módulo de torção}, ou, dito de outro modo, quando $E_0=\{0\}$, a álgebra exterior de $E_0$ reduz-se a $K$, e o determinante do único endomorfismo $u_0$ de $E_0$ é \emph{igual a $1$}.
\end{env}

\begin{proposition}[6.4.3]
\label{II.6.4.3}
Sejam $A$ um anel íntegro, $E$ um $A$-módulo de tipo finito e $u$ um endomorfismo de $E$;
então as funções simétricas elementares $\sigma_i(u)$ de $u$ (e, em particular, $\det u$) são elementos de $K$ inteiros sobre $A$.
\end{proposition}

\begin{proof}
Seja $A'$ o fecho integral de $A$;
como $A'[T]$ é um anel integralmente fechado \cite[p.~99]{II-24}, ele é o fecho integral de $A[T]$ em seu corpo de frações $K(T)$.
Substituindo $u$ por $T\cdot I-u\otimes1$ e $A$ por $A[T]$, vê-se que se pode reduzir à demonstração de que $\det u$ é inteiro sobre $A$.
Se $n$ é o posto de $E$, então $\det(u)=\det(\wedge^n u)$ e $(\wedge^n u)\otimes1=\wedge^n(u\otimes1)$, de modo que se pode supor $n=1$.
Mas então a aplicação $u\mapsto\det u$ é um homomorfismo do $A$-módulo $\Hom_A(E,E)$ em $K$;
como $E$ é de tipo finito, $\Hom_A(E,E)$ é isomorfo a um sub-$A$-módulo do $A$-módulo de tipo finito $E^m$ (se $E$ admite um sistema de $m$ geradores); portanto, os elementos $\det u$ pertencem a um sub-$A$-módulo de tipo finito de $K$ e são, por conseguinte, inteiros sobre $A$.
\end{proof}
\oldpage[II]{122}
\begin{corollary}[6.4.4]
\label{II.6.4.4}
Sob as hipóteses de \sref{II.6.4.3}, se se supõe ainda que $A$ seja normal, então os $\sigma_i(u)$ (e, em particular, $\Tr u$ e $\det u$) pertencem a $A$.
\end{corollary}

\begin{proposition}[6.4.5]
\label{II.6.4.5}
Sejam $A$ um anel íntegro, $E$ um $A$-módulo de tipo finito e de posto~$n$, e $u$ um endomorfismo de $E$ tal que os $\sigma_i(u)$ pertençam a $A$.
Para que $u$ seja um automorfismo de $E$, é necessário que $\det u$ seja invertível em $A$;
essa condição é suficiente se $E$ for sem torção.
\end{proposition}

\begin{proof}
A condição é \emph{suficiente} quando $E$ é sem torção, pois a hipótese e as equações~\sref{II.6.4.1.4} e \sref{II.6.4.1.5} (válidas em $E$, e não apenas em $E\otimes_A K$, visto que $E$ é sem torção) provam que $(\det u)^{-1}Q(u)$ é o inverso de $u$.

A condição é \emph{necessária}, pois, se $u$ é invertível, resulta de \sref{II.6.4.3} que $\det(u^{-1})$ pertence ao fecho integral $A'$ de $A$ em seu corpo de frações $K$ e é evidentemente o inverso de $\det u$ \emph{em $A'$}.
Nossa afirmação resulta então do seguinte lema:
\end{proof}

\begin{lemma}[6.4.5.1]
\label{II.6.4.5.1}
Seja $A$ um subanel de um anel $A'$ tal que $A'$ seja \emph{inteiro} sobre $A$.
Se um elemento $x\in A$ é invertível em $A'$, ele também é invertível em $A$.
\end{lemma}

\begin{proof}
Caso contrário, $x$ pertenceria a um ideal maximal $\mathfrak{m}$ de $A$, e resultaria do primeiro teorema de Cohen--Seidenberg \cite[t.~I, p.~257, th.~3]{I-13} que existiria um ideal maximal $\mathfrak{m}'$ de $A'$ tal que $\mathfrak{m}=\mathfrak{m}'\cap A$;
ter-se-ia, portanto, $x\in\mathfrak{m}'$, o que é absurdo.
\end{proof}

\begin{corollary}[6.4.6]
\label{II.6.4.6}
Sejam $A$ um anel íntegro e integralmente fechado, $E$ um $A$-módulo sem torção de tipo finito e $u$ um endomorfismo de $E$.
Para que $u$ seja um automorfismo de $E$, é necessário e suficiente que $\det u$ seja invertível em $A$.
\end{corollary}

\begin{proof}
Isso resulta de \sref{II.6.4.4} e \sref{II.6.4.5}.
\end{proof}

\begin{remark}[6.4.7]
\label{II.6.4.7}
Mais adiante, precisaremos de uma generalização dos resultados precedentes.
Consideremos um anel noetheriano \emph{reduzido} $A$;
sejam $\mathfrak{p}_\alpha$ ($1\leq\alpha\leq r$) seus ideais minimais, $K_\alpha$ o corpo de frações do anel íntegro $A/\mathfrak{p}_\alpha$, e $K$ o anel total de frações de $A$, que é a \emph{soma direta} dos corpos $K_\alpha$.
Seja $E$ um $A$-módulo de tipo finito, e \emph{suponhamos} que $E\otimes_A K$ seja um $K$-módulo \emph{livre} de posto~$n$ (o que aqui já não é consequência das demais hipóteses);
isso equivale a dizer que todos os $K_\alpha$-espaços vetoriais $E\otimes_A K_\alpha=E_\alpha$ têm \emph{a mesma dimensão~$n$};
se $u$ é então um endomorfismo de $E$, põe-se ainda $P(u,T)=P(u\otimes1,T)$ e $\sigma_j(u)=\sigma_j(u\otimes1)$ e, em particular, $\det u=\det(u\otimes1)$;
os $\sigma_j(u)$ são, portanto, elementos de $K$.
É imediato que $E\otimes_A K$ é a soma direta dos $E_\alpha$, e que estes são estáveis por $u\otimes1$;
a restrição de $u\otimes1$ a $E_\alpha$ é precisamente a extensão $u_\alpha$ de $u$ a esse $K_\alpha$-espaço vetorial;
conclui-se que $\sigma_j(u)$ é o elemento de $K$ cujas componentes nos $K_\alpha$ são os $\sigma_j(u_\alpha)$.
Como o fecho integral de $A$ em $K$ é a \emph{soma direta} dos fechos integrais de $A$ nos $K_\alpha$, os $\sigma_j(u)$ são \emph{inteiros} sobre $A$.

\begin{lemma}[6.4.7.1]
\label{II.6.4.7.1}
A sub-$A$-álgebra de $K$ gerada por todos os elementos $\sigma_j(u)$ ($1\leq j\leq n$), quando $u$ percorre $\Hom_A(E,E)$, é um $A$-módulo de tipo finito.
\end{lemma}

\begin{proof}
Basta demonstrar que a sub-$A[T]$-álgebra de $K[T]$ gerada pelos $P(u,T)$ é um $A[T]$-módulo de tipo finito, pois, se os $F_i(T)$ ($1\leq i\leq m$) formam um sistema de geradores desse $A[T]$-módulo, os coeficientes dos $F_i(T)$ são inteiros sobre $A$ e, portanto, geram uma $A$-álgebra que é um $A$-módulo de tipo finito \cite[t.~I, p.~255, th.~1]{I-13}.
Pode-se, portanto,
\oldpage[II]{123}
substituir $A$ por $A[T]$ (que é noetheriano) e $E$ por $E\otimes_A A[T]=E'$, de modo que $E'\otimes_{A[T]}K[T]=E\otimes_A K[T]$ seja um $K[T]$-módulo livre de posto~$n$.
Retomando as notações iniciais, vê-se, portanto, que basta provar que o $A$-módulo gerado pelos elementos $\det u$, quando $u$ percorre $\Hom_A(E,E)$, é de tipo finito;
\emph{a fortiori} (pois todo submódulo de um $A$-módulo de tipo finito é de tipo finito), basta provar que, quando $v$ percorre o conjunto dos endomorfismos de $\wedge^n E$, o $A$-módulo gerado pelos $\det v$ é de tipo finito;
dito de outro modo, pode-se novamente reduzir ao caso em que $n=1$.
Mas então a proposição resulta do fato de que $\Hom_A(E,E)$ é um $A$-módulo de tipo finito e de que $v\mapsto\det v$ é um homomorfismo desse $A$-módulo em $K$.
\end{proof}

Seja $F$ o núcleo do homomorfismo canônico $E\to E\otimes_A K$, e seja $E_0=E/F$;
vê-se, como acima, que $E\otimes_A K$ se identifica com $E_0\otimes_A K$, que $u(F)\subset F$ e que, se $u_0$ é o endomorfismo de $E_0$ deduzido de $u$ por passagem aos quocientes, então $u\otimes1$ se identifica com $u_0\otimes1$, e $\sigma_j(u)=\sigma_j(u_0)$ para todo $j$.
\emph{Se $F=0$}, as equações~\sref{II.6.4.1.3} e \sref{II.6.4.1.5} ainda têm sentido e continuam válidas quando $E$ é identificado com um submódulo de $E\otimes_A K$, e os $u^j$ com homomorfismos $E\to E\otimes_A K$;
a Proposição~\sref{II.6.4.5} estende-se, portanto, também a esse caso, com a mesma demonstração.
\end{remark}

\begin{env}[6.4.8]
\label{II.6.4.8}
Sejam $(X,\sh{A})$ um espaço anelado, $\sh{E}$ um $\sh{A}$-módulo \emph{localmente livre} (de posto finito) e $u$ um endomorfismo de $\sh{E}$.
Por hipótese, existe uma base $\mathfrak{B}$ da topologia de $X$ tal que, para todo $V\in\mathfrak{B}$, $\sh{E}|V$ seja isomorfo a algum $\sh{A}^n|V$ (podendo $n$ variar com $V$).
Seja $u$ um endomorfismo de $\sh{E}$;
para todo $V\in\mathfrak{B}$, $u_V$ é, portanto, um endomorfismo do $\Gamma(V,\sh{A})$-módulo $\Gamma(V,\sh{E})$, que é livre por hipótese;
o determinante $\det u_V$ está, portanto, definido e pertence a $\Gamma(V,\sh{A})$.
Além disso, se $e_1,\ldots,e_n$ formam uma base de $\Gamma(V,\sh{E})$, suas restrições a todo aberto $W\subset V$ formam uma base de $\Gamma(W,\sh{E})$ sobre $\Gamma(W,\sh{A})$, de modo que $\det u_W$ é a restrição de $\det u_V$ a $W$.
Existe, portanto, uma única seção de $\sh{A}$ sobre $X$, denotada por $\det u$ e chamada \emph{determinante} de $u$, tal que a restrição de $\det u$ a todo $V\in\mathfrak{B}$ seja $\det u_V$.
É claro que, para todo $x\in X$, tem-se $(\det u)_x=\det u_x$;
para dois endomorfismos $u$ e $v$ de $\sh{E}$, tem-se
\[
\label{II.6.4.8.1}
  \det(u\circ v) = (\det u)(\det v)
\tag{6.4.8.1}
\]
assim como
\[
\label{II.6.4.8.2}
  \det(1_{\sh{E}}) = 1_{\sh{A}}
\tag{6.4.8.2}
\]
e, se o posto de $\sh{E}$ é constante (o que será o caso \sref[0]{0.5.4.1} se $X$ for conexo) e igual a $n$,
\[
\label{II.6.4.8.3}
  \det(s\cdot u) = s^n\det u
\tag{6.4.8.3}
\]
para todo $s\in\Gamma(X,\sh{A})$ (observemos que $\det(0)=0_{\sh{A}}$ se $n\geq1$, mas $\det(0)=1_{\sh{A}}$ se $n=0$).
Além disso, para que $u$ seja um \emph{automorfismo} de $\sh{E}$, é necessário e suficiente que $\det u$ seja \emph{invertível} em $\Gamma(X,\sh{A})$.

Se o posto de $\sh{E}$ é constante, podem-se definir da mesma maneira as funções simétricas elementares $\sigma_i(u)$, que são elementos de $\Gamma(X,\sh{A})$;
têm-se novamente as relações de \sref{II.6.4.1.3} a \sref{II.6.4.1.5}.
\oldpage[II]{124}
Definiu-se, assim, um homomorfismo $\det:\Hom_{\sh{A}}(\sh{E},\sh{E})\to\Gamma(X,\sh{A})$ de monoides multiplicativos;
observando-se que $\Hom_{\sh{A}}(\sh{E},\sh{E})=\Gamma(X,\shHom_{\sh{A}}(\sh{E},\sh{E}))$ por definição, vê-se que se pode, nessa definição, substituir $X$ por um aberto qualquer $U$ de $X$, o que define imediatamente um homomorfismo $\det:\shHom_{\sh{A}}(\sh{E},\sh{E})\to\sh{A}$ de feixes de monoides multiplicativos.
Quando $\sh{E}$ tem posto constante, definem-se da mesma maneira homomorfismos $\sigma_i:\shHom_{\sh{A}}(\sh{E},\sh{E})\to\sh{A}$ de feixes de conjuntos;
para $i=1$, o homomorfismo $\sigma_1=\Tr$ é um homomorfismo de $\sh{A}$-módulos.

Seja $(Y,\sh{B})$ um segundo espaço anelado, e seja $f:(X,\sh{A})\to(Y,\sh{B})$ um morfismo de espaços anelados;
se $\sh{F}$ é um $\sh{B}$-módulo localmente livre, então $f^*(\sh{F})$ é um $\sh{A}$-módulo localmente livre (de mesmo posto que $\sh{F}$ se este último tiver posto constante) \sref[0]{0.5.4.5}.
Para todo endomorfismo $v$ de $\sh{F}$, $f^*(v)$ é um endomorfismo de $f^*(\sh{F})$, e resulta imediatamente das definições que $\det f^*(v)$ é a seção de $\sh{A}=f^*(\sh{B})$ sobre $X$ que corresponde canonicamente a $\det v\in\Gamma(Y,\sh{B})$.
Pode-se ainda dizer que o homomorfismo $f^*(\det):f^*(\shHom_{\sh{B}}(\sh{F},\sh{F}))\to f^*(\sh{B})=\sh{A}$ é o composto
\[
  f^*(\shHom_{\sh{B}}(\sh{F},\sh{F}))
  \xrightarrow{\gamma^\sharp} \shHom_{\sh{A}}(f^*(\sh{F}),f^*(\sh{F}))
  \xrightarrow{\det} \sh{A}
\]
\sref[0]{0.4.4.6}.
Têm-se resultados análogos para os $\sigma_i$.
\end{env}

\begin{env}[6.4.9]
\label{II.6.4.9}
Suponhamos agora que $X$ seja um \emph{pré-esquema localmente integral}, de modo que o feixe $\sh{R}(X)$ de funções racionais sobre $X$ seja um feixe de corpos localmente simples \sref[I]{I.7.3.4}, e quase-coerente como $\sh{O}_X$-módulo.
Se $\sh{E}$ é um $\sh{O}_X$-módulo quase-coerente \emph{de tipo finito}, então $\sh{E}'=\sh{E}\otimes_{\sh{O}_X}\sh{R}(X)$ é um $\sh{R}(X)$-módulo localmente livre \sref[I]{I.7.3.6};
para todo endomorfismo $u$ de $\sh{E}$, $u\otimes1_{\sh{R}(X)}$ é então um endomorfismo de $\sh{E}'$, e $\det(u\otimes1)$ é uma seção de $\sh{R}(X)$ sobre $X$, chamada ainda \emph{determinante} de $u$ e denotada ainda por $\det u$.
Resulta de \sref{II.6.4.3} que $\det u$ é uma seção do \emph{fecho integral} de $\sh{O}_X$ em $\sh{R}(X)$ \sref{II.6.3.2};
se, além disso, $X$ é \emph{normal}, então $\det u$ é uma \emph{seção de $\sh{O}_X$} sobre $X$ e, se se supõe ainda que $\sh{E}$ seja \emph{sem torção}, para que $u$ seja um automorfismo de $\sh{E}$ é necessário e suficiente que $\det u$ seja \emph{invertível}, em virtude de \sref{II.6.4.6}.
As equações~\sref{II.6.4.8.1} a \sref{II.6.4.8.3} continuam válidas;
do homomorfismo $u\mapsto\det u$, aplicado aos módulos de seções de $\shHom_{\sh{O}_X}(\sh{E},\sh{E})$, deduz-se um homomorfismo de feixes $\det:\shHom_{\sh{O}_X}(\sh{E},\sh{E})\to\sh{R}(X)$, que toma seus valores em $\sh{O}_X$ quando $X$ é normal.
Têm-se definições e resultados análogos para as demais funções simétricas elementares $\sigma_j(u)$, se $\sh{E}'$ tiver \emph{posto constante};
se, além disso, $X$ é \emph{normal}, então os $\sigma_j(u)$ são \emph{seções de $\sh{O}_X$} sobre $X$.

Por fim, sejam $X$ e $Y$ dois pré-esquemas integrais, e seja $f:X\to Y$ um morfismo \emph{dominante}.
Sabe-se que existe um homomorfismo canônico $f^*(\sh{R}(Y))\to\sh{R}(X)$ \sref[I]{I.7.3.8}\footnote{Ver \emph{Errata et Addenda}, lista 1.}, de onde resulta, para todo $\sh{O}_Y$-módulo quase-coerente de tipo finito $\sh{F}$, um homomorfismo canônico $\theta:f^*(\sh{F}\otimes_{\sh{O}_Y}\sh{R}(Y))\to f^*(\sh{F})\otimes_{\sh{O}_X}\sh{R}(X)$.

\oldpage[II]{125}
Se $v$ é um endomorfismo de $\sh{F}$, então $f^*(v\otimes1_{\sh{R}(Y)})$ é um endomorfismo de $f^*(\sh{F}\otimes_{\sh{O}_Y}\sh{R}(Y))$, e tem-se um diagrama comutativo
\[
  \xymatrix{
    f^*(\sh{F}\otimes_{\sh{O}_Y}\sh{R}(Y))
      \ar[r]^{f^*(v\otimes1)}
      \ar[d]_{\theta}
    & f^*(\sh{F}\otimes_{\sh{O}_Y}\sh{R}(Y))
      \ar[d]^{\theta}
  \\f^*(\sh{F})\otimes_{\sh{O}_X}\sh{R}(X)
      \ar[r]_{f^*(v)\otimes1}
    & f^*(\sh{F})\otimes_{\sh{O}_X}\sh{R}(X)
  }
\]

Conclui-se facilmente que $\det f^*(v)$ é a imagem canônica, pelo homomorfismo $f^*(\sh{R}(Y))\to\sh{R}(X)$, da seção $\det v$ de $\sh{R}(Y)$;
com efeito, reduzimo-nos imediatamente ao caso em que $X=\Spec(A)$ e $Y=\Spec(B)$ são afins, sendo $A$ e $B$ anéis integrais de corpos de frações respectivos $K$ e $L$, e sendo o homomorfismo $B\to A$ injetivo e, portanto, prolongável a um monomorfismo $L\to K$;
se $\sh{F}=\widetilde{M}$, onde $M$ é um $B$-módulo de tipo finito, então o posto de $M\otimes_B L$ sobre $L$ é \emph{igual} ao de $(M\otimes_B A)\otimes_A K$ sobre $K$, e $\det((u\otimes1)\otimes1)$ é a imagem em $K$ de $\det(u\otimes1)$ para todo endomorfismo $u$ de $M$, donde a conclusão.
\end{env}

\begin{env}[6.4.10]
\label{II.6.4.10}
Suponhamos, por fim, que $X$ seja um \emph{pré-esquema localmente noetheriano reduzido}, de modo que o feixe $\sh{R}(X)$ de funções racionais sobre $X$ seja ainda um $\sh{O}_X$-módulo quase-coerente \sref[I]{I.7.3.4};
seja, por outro lado, $\sh{E}$ um $\sh{O}_X$-módulo \emph{coerente} tal que $\sh{E}'=\sh{E}\otimes_{\sh{O}_X}\sh{R}(X)$ seja \emph{localmente livre} (de posto finito).
Em virtude de \sref{II.6.4.7}, se $\sh{E}'$ é de posto constante, podem-se definir, para todo endomorfismo $u$ de $\sh{E}$, os $\sigma_j(u)$, que são seções de $\sh{R}(X)$ sobre $X$.
Quando não se supõe mais que $\sh{E}'$ seja de posto constante, ainda se pode definir o homomorfismo $\det:\shHom_{\sh{O}_X}(\sh{E},\sh{E})\to\sh{R}(X)$.
\end{env}

\subsection{Norma de um feixe invertível}
\label{subsection:II.6.5}

\begin{env}[6.5.1]
\label{II.6.5.1}
Sejam $(X,\sh{A})$ um espaço anelado e $\sh{B}$ uma $\sh{A}$-álgebra (comutativa).
O $\sh{A}$-módulo $\sh{B}$ identifica-se canonicamente com um $\sh{A}$-submódulo de $\shHom_{\sh{A}}(\sh{B},\sh{B})$, identificando-se uma seção $f$ de $\sh{B}$ sobre um aberto $U$ de $X$ com a multiplicação por essa seção.
Se $(X,\sh{A})$ e $\sh{B}$ satisfazem uma das condições enumeradas em \sref{II.6.4.8}, \sref{II.6.4.9} ou \sref{II.6.4.10}, pode-se, portanto, definir $\det(f)$ (e, em certos casos, os $\sigma_j(f)$), que são seções de $\sh{A}$ ou de $\sh{R}(X)$ sobre $U$; chamaremos a primeira de \emph{norma} de $f$ (e as demais, de \emph{funções simétricas elementares} de $f$), e a denotaremos por $N_{\sh{B}/\sh{A}}(f)$.
Suporemos satisfeita \emph{uma} das condições seguintes:
\begin{enumerate}
  \item[{\rm(I)}] \emph{$\sh{B}$ é um $\sh{A}$-módulo localmente livre (de posto finito).}
  \item[{\rm(II)}] \emph{$(X,\sh{A})$ é um pré-esquema localmente noetheriano reduzido, $\sh{B}$ é um $\sh{A}$-módulo coerente tal que $\sh{B}\otimes_{\sh{A}}\sh{R}(X)$ é um $\sh{R}(X)$-módulo localmente livre e, para toda seção $f\in\Gamma(U,\sh{B})$ sobre um aberto $U$ de $X$, a norma $N_{\sh{B}/\sh{A}}(f)$ é uma seção de $\sh{A}$ sobre $U$.}
\end{enumerate}
\oldpage[II]{126}
Observemos que esta última condição é automaticamente satisfeita quando o pré-esquema localmente noetheriano $X$ é \emph{normal} \sref{II.6.4.9}.

Por outro lado, a hipótese de que $\sh{B}\otimes_{\sh{A}}\sh{R}(X)$ seja localmente livre pode ainda ser expressa da maneira seguinte: denotemos por $X_\alpha$ os subpré-esquemas fechados reduzidos de $X$ cujos espaços subjacentes são as componentes irredutíveis de $X$ \sref[I]{I.5.2.1}, que são, portanto, pré-esquemas integrais localmente noetherianos.
Todo $x\in X$ pertence a um número finito dos subespaços $X_\alpha$;
por outro lado, $\sh{B}\otimes_{\sh{A}}\sh{R}(X_\alpha)$ é um $\sh{R}(X_\alpha)$-módulo localmente livre de posto constante~$k_\alpha$ \sref[I]{I.7.3.6};
dizer que $\sh{B}\otimes_{\sh{A}}\sh{R}(X)$ é um $\sh{R}(X)$-módulo localmente livre significa que, para todo $x\in X$, \emph{são iguais os postos $k_\alpha$ correspondentes aos índices tais que $x\in X_\alpha$}.
A questão é, com efeito, local, e reduzimo-nos ao caso $X=\Spec(C)$, onde $C$ é um anel noetheriano reduzido, e $\sh{B}=\widetilde{D}$, onde $D$ é uma $C$-álgebra que é um $C$-módulo de tipo finito;
se $\mathfrak{p}_i$ ($1\leq i\leq m$) são os ideais primos minimais de $C$, o anel total de frações $L$ de $C$ é a soma direta dos corpos de frações $K_i$ dos anéis integrais $C/\mathfrak{p}_i$, e $D\otimes_C L$ é a soma direta dos $D\otimes_C K_i$, donde a conclusão.

É claro que, sob as hipóteses~(I) ou (II), define-se assim um \emph{homomorfismo de feixes de monoides multiplicativos} $N_{\sh{B}/\sh{A}}:\sh{B}\to\sh{A}$, também denotado por $N$ se disso não resultar confusão, e chamado homomorfismo \emph{norma}.
Para duas seções $f$ e $g$ de $\sh{B}$ sobre um mesmo aberto $U$, tem-se, portanto,
\[
\label{II.6.5.1.1}
  N_{\sh{B}/\sh{A}}(fg) = N_{\sh{B}/\sh{A}}(f)N_{\sh{B}/\sh{A}}(g)
\tag{6.5.1.1}
\]
para as seções correspondentes de $\sh{A}$ sobre $U$;
\[
\label{II.6.5.1.2}
  N_{\sh{B}/\sh{A}}(1_{\sh{B}}) = 1_{\sh{A}};
\tag{6.5.1.2}
\]
por fim, para toda seção $s$ de $\sh{A}$ sobre $U$,
\[
\label{II.6.5.1.3}
  N_{\sh{B}/\sh{A}}(s\cdot1_{\sh{B}}) = s^n
\tag{6.5.1.3}
\]
se o posto de $\sh{B}$ é constante e igual a $n$ (para $s=0_{\sh{A}}$, essa fórmula dá $N(0_{\sh{B}})=0_{\sh{A}}$ se $n\geq1$, e $N(0_{\sh{B}})=N(1_{\sh{B}})=1_{\sh{A}}$ se $n=0$).

Sob a hipótese~(I), para que $f\in\Gamma(U,\sh{B})$ seja invertível, é necessário e suficiente que $N(f)\in\Gamma(U,\sh{A})$ seja invertível.
Sob a hipótese~(II), essa condição é necessária;
ela é suficiente (por \sref{II.6.4.7}) quando se supõe que $\sh{B}\to\sh{B}\otimes_{\sh{A}}\sh{R}(X)$ é \emph{injetivo} e que se verifica a hipótese mais restritiva seguinte:
\begin{enumerate}
  \item[{\rm(II~\emph{bis})}] \emph{$(X,\sh{A})$ é um pré-esquema localmente noetheriano reduzido, $\sh{B}$ é um $\sh{A}$-módulo coerente tal que $\sh{B}\otimes_{\sh{A}}\sh{R}(X)$ é um $\sh{R}(X)$-módulo localmente livre e, para toda seção $f\in\Gamma(U,\sh{B})$ sobre um aberto $U$ tal que $\sh{B}\otimes_{\sh{A}}\sh{R}(X)|U$ seja de posto constante~$n$ sobre $\sh{R}(X)|U$, os $\sigma_j(f)$ ($1\leq j\leq n$) são seções de $\sh{A}$ sobre $U$.}
\end{enumerate}

(Observemos ainda que essa condição é satisfeita quando $X$ é \emph{normal}.)
\end{env}

\begin{env}[6.5.2]
\label{II.6.5.2}
Suponhamos satisfeita uma das hipóteses~(I) ou (II) de \sref{II.6.5.1}, e seja $\sh{L}'$ um $\sh{B}$-módulo \emph{invertível}.
Associaremos canonicamente a ele (a menos de isomorfismo único) um $\sh{A}$-módulo \emph{invertível} $\sh{L}$ da maneira seguinte.
Denotemos por $\sh{A}^*$ (resp. $\sh{B}^*$) o subfeixe de $\sh{A}$ (resp. de $\sh{B}$) tal que $\Gamma(U,\sh{A}^*)$ (resp. $\Gamma(U,\sh{B}^*)$) seja o conjunto dos
\oldpage[II]{127}
elementos invertíveis de $\Gamma(U,\sh{A})$ (resp. de $\Gamma(U,\sh{B})$) para todo aberto $U\subset X$;
esses são feixes de grupos multiplicativos, e $N_{\sh{B}/\sh{A}}$, restrito a $\sh{B}^*$, é um \emph{homomorfismo} $\sh{B}^*\to\sh{A}^*$ de feixes de grupos \sref{II.6.5.1}.
Seja $\mathfrak{L}$ o conjunto dos pares $(U_\lambda,\nu_\lambda)$ que têm a propriedade seguinte: $U_\lambda$ é um aberto de $X$, e $\nu_\lambda$ é um isomorfismo $\sh{L}'|U_\lambda\simto\sh{B}|U_\lambda$ de $(\sh{B}|U_\lambda)$-módulos.
Por hipótese, os $U_\lambda$ formam um recobrimento de $X$;
para dois índices quaisquer $\lambda$ e $\mu$, ponhamos $\omega_{\lambda\mu}=(\nu_\lambda|U_\lambda\cap U_\mu)\circ(\nu_\mu|U_\lambda\cap U_\mu)^{-1}$, automorfismo de $\sh{B}|U_\lambda\cap U_\mu$ que se identifica canonicamente com uma seção de $\sh{B}^*$ sobre $U_\lambda\cap U_\mu$, e $(\omega_{\lambda\mu})$ é um $1$-cociclo do recobrimento $\mathfrak{U}=(U_\lambda)$ com valores em $\sh{B}^*$ \sref[0]{0.5.4.7}.
O fato de que $N_{\sh{B}/\sh{A}}:\sh{B}^*\to\sh{A}^*$ é um homomorfismo implica então que $(N_{\sh{B}/\sh{A}}\omega_{\lambda\mu})$ é um $1$-cociclo de $\mathfrak{U}$ com valores em $\sh{A}^*$, ao qual corresponde, portanto (a menos de isomorfismo único), um $\sh{A}$-módulo invertível $\sh{L}$ \sref[0]{0.5.4.7}, que denotaremos por $N_{\sh{B}/\sh{A}}(\sh{L}')$ e chamaremos \emph{norma} do $\sh{B}$-módulo invertível $\sh{L}'$.

Seja $\mathfrak{M}$ um subconjunto de $\mathfrak{L}$ tal que os $U_\lambda$ correspondentes ainda formem um recobrimento de $X$, e seja $\mathfrak{B}$ esse recobrimento;
a restrição do cociclo $(\omega_{\lambda\mu})$ a $\mathfrak{B}$ define um $1$-cociclo $(N_{\sh{B}/\sh{A}}\omega_{\lambda\mu})$ de $\mathfrak{B}$ com valores em $\sh{A}^*$, restrição do $1$-cociclo $(N_{\sh{B}/\sh{A}}\omega_{\lambda\mu})$ de $\mathfrak{U}$;
é claro que há um isomorfismo canônico do $\sh{A}$-módulo invertível definido por esse $1$-cociclo de $\mathfrak{B}$ para $N_{\sh{B}/\sh{A}}(\sh{L}')$, o que permite definir esse $\sh{A}$-módulo invertível por um sub-recobrimento qualquer de $\mathfrak{U}$.
Essa possibilidade mostra imediatamente que, se $\sh{L}'_1$ e $\sh{L}'_2$ são dois $\sh{B}$-módulos invertíveis, então, em virtude de \sref{II.6.5.1.1} e \sref{II.6.5.1.2},
\[
\label{II.6.5.2.1}
  N(\sh{L}'_1\otimes_{\sh{B}}\sh{L}'_2) = N(\sh{L}'_1)\otimes_{\sh{A}}N(\sh{L}'_2)
\tag{6.5.2.1}
\]
e
\[
\label{II.6.5.2.2}
  N_{\sh{B}/\sh{A}}(\sh{B}) = \sh{A}
\tag{6.5.2.2}
\]
assim como
\[
\label{II.6.5.2.3}
  N({\sh{L}'}^{-1}) = (N(\sh{L}'))^{-1}
\tag{6.5.2.3}
\]
a menos de isomorfismos canônicos.
Além disso, resulta de \sref{II.6.5.1.3} que, se $\sh{L}$ é um $\sh{A}$-módulo invertível e se $\sh{B}$ é de posto constante~$n$ sobre $\sh{A}$ no caso~(I) (resp. $\sh{B}\otimes_{\sh{A}}\sh{R}(X)$ é de posto constante~$n$ sobre $\sh{R}(X)$ no caso~(II)), então se tem, a menos de isomorfismo canônico,
\[
\label{II.6.5.2.4}
  N_{\sh{B}/\sh{A}}(\sh{L}\otimes_{\sh{A}}\sh{B}) = \sh{L}^{\otimes n}.
\tag{6.5.2.4}
\]
\end{env}

\begin{env}[6.5.3]
\label{II.6.5.3}
Mostremos agora que $N_{\sh{B}/\sh{A}}$ é um \emph{funtor} covariante na categoria dos $\sh{B}$-módulos invertíveis.
Seja, com efeito, $h':\sh{L}'_1\to\sh{L}'_2$ um homomorfismo de $\sh{B}$-módulos invertíveis, e seja $\mathfrak{B}=(U_\lambda)$ um recobrimento aberto de $X$ tal que, para todo $\lambda$, se tenham $(\sh{B}|U_\lambda)$-isomorfismos $\eta_\lambda^{(1)}:\sh{L}'_1|U_\lambda\simto\sh{B}|U_\lambda$ e $\eta_\lambda^{(2)}:\sh{L}'_2|U_\lambda\simto\sh{B}|U_\lambda$;
há, portanto, para cada $\lambda$, um endomorfismo $h'_\lambda$ de $\sh{B}|U_\lambda$ tal que $h'_\lambda\circ\eta_\lambda^{(1)}=\eta_\lambda^{(2)}\circ(h'|U_\lambda)$, e pode-se evidentemente identificar $h'_\lambda$ com uma seção de $\sh{B}$ sobre $U_\lambda$ \sref[0]{0.5.1.1}.
Logo, para todo par $(\lambda,\mu)$ de índices, coincidem as restrições a $U_\lambda\cap U_\mu$ de $(\eta_\lambda^{(2)})^{-1}\circ h'_\lambda\circ\eta_\lambda^{(1)}$ e de $(\eta_\mu^{(2)})^{-1}\circ h'_\mu\circ\eta_\mu^{(1)}$.
Deduz-se, para os $1$-cociclos $(\omega_{\lambda\mu}^{(1)})$ e $(\omega_{\lambda\mu}^{(2)})$ com valores em $\sh{B}^*$ correspondentes a $\sh{L}'_1$ e $\sh{L}'_2$, a relação
\[
  \omega_{\lambda\mu}^{(2)}h'_\mu = h'_\lambda\omega_{\lambda\mu}^{(1)}.
\]
\oldpage[II]{128}
Se pusermos $h_\lambda=N(h'_\lambda)$, teremos, portanto, as relações análogas
\[
  N(\omega_{\lambda\mu}^{(2)})h_\mu = h_\lambda N(\omega_{\lambda\mu}^{(1)})
\]
e, por conseguinte, os $h_\lambda$ definem um homomorfismo $N(\sh{L}'_1)\to N(\sh{L}'_2)$ que denotaremos por $N_{\sh{B}/\sh{A}}(h')$ ou $N(h')$.
Sob a hipótese~(I), para que $h'$ seja um \emph{isomorfismo}, é necessário e suficiente (pois se trata de uma questão local) que $N_{\sh{B}/\sh{A}}(h')$ o seja.
Sob a hipótese~(II), essa condição é ainda necessária;
ela é suficiente se a hipótese~(II~\emph{bis}) é satisfeita e se $\sh{B}\to\sh{B}\otimes_{\sh{A}}\sh{R}(X)$ é injetivo.

Tomemos, em particular, $\sh{L}'_1=\sh{B}$;
os homomorfismos $\sh{B}\to\sh{L}'$ identificam-se então \sref[0]{0.5.1.1} com as seções de $\sh{L}'$ sobre $X$, donde uma aplicação canônica
\[
  N_{\sh{B}/\sh{A}} : \Gamma(X,\sh{L}')\to\Gamma(X,N_{\sh{B}/\sh{A}}(\sh{L}')).
\]

Resulta ainda de \sref{II.6.5.1.1} que, se $f'_1\in\Gamma(X,\sh{L}'_1)$ e $f'_2\in\Gamma(X,\sh{L}'_2)$, então
\[
\label{II.6.5.3.1}
  N(f'_1\otimes f'_2) = N(f'_1)\otimes N(f'_2).
\tag{6.5.3.1}
\]

Para todo $\sh{A}$-módulo invertível $\sh{L}$ e toda seção $f\in\Gamma(X,\sh{L})$, tem-se, levando em conta \sref{II.6.5.2.4},
\[
\label{II.6.5.3.2}
  N_{\sh{B}/\sh{A}}(f\otimes 1_{\sh{B}}) = f^{\otimes n}
\tag{6.5.3.2}
\]
quando $\sh{B}$ é de posto constante~$n$ sob a hipótese~(I) (resp. quando $\sh{B}\otimes_{\sh{A}}\sh{R}(X)$ é de posto constante~$n$ sob a hipótese~(II)).
Por fim, para que o homomorfismo $\sh{B}\to\sh{L}'$ correspondente a uma seção $f'$ de $\sh{L}'$ sobre $X$ seja um isomorfismo, é necessário e suficiente que $f'_x$ seja uma base de $\sh{L}'_x$ para todo $x\in X$;
sob a hipótese~(I), essa condição equivale, portanto, a dizer que $(N(f'))_x$ é uma base de $(N(\sh{L}'))_x$ para todo $x$;
sob a hipótese~(II), essa condição é ainda necessária, e é suficiente quando $\sh{B}$ satisfaz a hipótese~(II~\emph{bis}) e $\sh{B}\to\sh{B}\otimes_{\sh{A}}\sh{R}(X)$ é injetivo.
\end{env}

\begin{env}[6.5.4]
\label{II.6.5.4}
Sejam $(X,\sh{A})$ e $(X',\sh{A}')$ dois espaços anelados, $f:X'\to X$ um morfismo, $\sh{B}$ uma $\sh{A}$-álgebra e $\sh{B}'=f^*(\sh{B})$ a $\sh{A}'$-álgebra imagem inversa.
Suponhamos satisfeita uma das hipóteses seguintes:
\begin{enumerate}
  \item $\sh{B}$ satisfaz a hipótese~(I) de \sref{II.6.5.1}.
  \item $(X,\sh{A})$ e $\sh{B}$ satisfazem a hipótese~(II) de \sref{II.6.5.1}, $(X',\sh{A}')$ é um pré-esquema localmente noetheriano reduzido e, se denotarmos por $X_\alpha$ e $X'_\beta$ os subpré-esquemas fechados reduzidos respectivos de $X$ e $X'$ cujos espaços subjacentes são as componentes irredutíveis desses espaços, a restrição de $f$ a cada $X'_\beta$ é um morfismo \emph{dominante} de $X'_\beta$ em um $X_\alpha$.
\end{enumerate}

Nessas condições, $\sh{B}'$ satisfaz, respectivamente, a hipótese~(I) ou a hipótese~(II) de \sref{II.6.5.1};
a primeira afirmação é imediata.
Para estabelecer a segunda, basta ver que, para todo $x'\in X'$, os postos dos $\sh{B}'\otimes_{\sh{O}_{X'}}\sh{R}(X'_\beta)$ são \emph{os mesmos} para os $\beta$ tais que $x'\in X'_\beta$.
Ora, se a restrição de $f$ a $X'_\beta$ é um morfismo dominante em $X_\alpha$, o posto de $\sh{B}'\otimes_{\sh{O}_{X'}}\sh{R}(X'_\beta)$ é igual ao de $\sh{B}\otimes_{\sh{O}_X}\sh{R}(X_\alpha)$ (como se vê imediatamente reduzindo-nos ao
\oldpage[II]{129}
caso afim, como em \sref{II.6.4.9}), donde nossa afirmação, em virtude da hipótese~(II) e de \sref{II.6.5.1}.

Posto isso, resulta de \sref{II.6.4.8}, \sref{II.6.4.9} e \sref{II.6.4.10} que, se $s$ é uma seção de $\sh{B}$ sobre um aberto $U\subset X$, e $s'$ a seção correspondente de $\sh{B}'$ sobre $f^{-1}(U)$, então $N_{\sh{B}'/\sh{A}'}(s')$ é a seção de $\sh{A}'$ sobre $f^{-1}(U)$ que corresponde à seção $N_{\sh{B}/\sh{A}}(s)$ de $\sh{A}$ sobre $U$.

Se $\sh{M}$ é um $\sh{B}$-módulo invertível, deduz-se do que precede que, se $\sh{M}'=f^*(\sh{M})$ (que é um $\sh{B}'$-módulo invertível), então se tem, a menos de isomorfismo canônico, $N_{\sh{B}'/\sh{A}'}(\sh{M}')=f^*(N_{\sh{B}/\sh{A}}(\sh{M}))$.
\end{env}

\begin{env}[6.5.5]
\label{II.6.5.5}
Suponhamos, doravante, que $(X,\sh{A})$ seja um \emph{pré-esquema}.
O dado de uma $\sh{A}$-álgebra quase-coerente $\sh{B}$, que é um \emph{$\sh{A}$-módulo de tipo finito}, equivale então, como se sabe, ao de um morfismo \emph{finito} $g:X'\to X$ tal que $g_*(\sh{O}_{X'})=\sh{B}$, definido a menos de $X$-isomorfismo (\sref{II.6.1.2} e \sref{II.1.3.1}).
Além disso, dar um $\sh{O}_{X'}$-módulo quase-coerente $\sh{F}'$ equivale a dar um $\sh{B}$-módulo quase-coerente $\sh{F}$ tal que $g_*(\sh{F}')=\sh{F}$ \sref{II.1.4.3}, e, para que $\sh{F}'$ seja invertível, é necessário e suficiente que $\sh{F}$ o seja \sref{II.6.1.12}.
Para traduzir os resultados precedentes em termos do morfismo finito $g$, será preciso, portanto, supor ou que $g_*(\sh{O}_{X'})$ é um $\sh{O}_X$-módulo \emph{localmente livre} (de tipo finito), ou que $(X,\sh{O}_X)$ e $g_*(\sh{O}_{X'})$ satisfazem a hipótese~(II).
Para todo $\sh{O}_{X'}$-módulo invertível $\sh{L}'$, poremos então
\[
\label{II.6.5.5.1}
  N_{X'/X}(\sh{L}') = N_{g_*(\sh{O}_{X'})/\sh{O}_X}(g_*(\sh{L}'))
\tag{6.5.5.1}
\]
e diremos que essa é a \emph{norma} (relativa a $g$) de $\sh{L}'$.
Do mesmo modo, se $h':\sh{L}'_1\to\sh{L}'_2$ é um homomorfismo de $\sh{O}_{X'}$-módulos invertíveis, poremos
\[
\label{II.6.5.5.2}
  N_{X'/X}(h') = N_{g_*(\sh{O}_{X'})/\sh{O}_X}(g_*(h')) : N_{X'/X}(\sh{L}'_1)\to N_{X'/X}(\sh{L}'_2).
\tag{6.5.5.2}
\]

Em particular, para $\sh{L}'_1=\sh{O}_{X'}$, obtém-se assim uma aplicação canônica
\[
\label{II.6.5.5.3}
  N_{X'/X} : \Gamma(X',\sh{L}') \to \Gamma(X,N_{X'/X}(\sh{L}')).
\tag{6.5.5.3}
\]

Deixaremos ao leitor a maior parte dessas traduções e nos limitaremos a explicitar as seguintes:
\end{env}

\begin{proposition}[6.5.6]
\label{II.6.5.6}
Seja $g:X'\to X$ um morfismo finito, e suponhamos ou que $g_*(\sh{O}_{X'})$ é um $\sh{O}_X$-módulo localmente livre, ou que $(X,\sh{O}_X)$ e $g_*(\sh{O}_{X'})$ satisfazem (II~\emph{bis}) (o que ocorrerá, em particular, quando $X$ for localmente noetheriano e normal).
Para que um homomorfismo $h':\sh{L}'_1\to\sh{L}'_2$ de $\sh{O}_{X'}$-módulos invertíveis seja um isomorfismo, é necessário e suficiente, na primeira hipótese, que $N_{X'/X}(h')$ seja um isomorfismo;
na segunda hipótese, essa condição é ainda necessária, e é suficiente quando o homomorfismo $g_*(\sh{O}_{X'})\to g_*(\sh{O}_{X'})\otimes_{\sh{O}_X}\sh{R}(X)$ é injetivo.
\end{proposition}

\begin{proof}
Observemos que se utiliza aqui o fato de que, para que $g_*(h')$ seja um isomorfismo, é necessário e suficiente que $h'$ o seja \sref{II.1.4.2}.
\end{proof}

\begin{corollary}[6.5.7]
\label{II.6.5.7}
Seja $g:X'\to X$ um morfismo finito, e suponhamos ou que $g_*(\sh{O}_{X'})$ é um $\sh{O}_X$-módulo localmente livre, ou que $(X,\sh{O}_X)$ e $g_*(\sh{O}_{X'})$ satisfazem (II~\emph{bis}) e que $g_*(\sh{O}_{X'})\to g_*(\sh{O}_{X'})\otimes_{\sh{O}_X}\sh{R}(X)$ é injetivo.
Seja $\sh{L}'$ um $\sh{O}_{X'}$-módulo invertível, $f'$ uma seção de $\sh{L}'$ sobre $X'$, e $f=N_{X'/X}(f')$ a seção correspondente de $\sh{L}=N_{X'/X}(\sh{L}')$ sobre $X$ que
\oldpage[II]{130}
lhe corresponde \sref{II.6.5.5.3}.
Então, $g(X'\setminus X'_{f'})=X\setminus X_f$, e $X_f$ é o maior aberto $U$ de $X$ tal que $g^{-1}(U)\subset X'_{f'}$.
\end{corollary}

\begin{proof}
Com efeito, $g(X'\setminus X'_{f'})$ é fechado em $X$ \sref{II.6.1.10};
basta, portanto, demonstrar a última afirmação.
Ora, a relação $U\subset X_f$ equivale ao fato de que o homomorfismo $\sh{O}_X|U\to\sh{L}|U$ definido por $f|U$ é um isomorfismo.
Em virtude de \sref{II.6.5.6}, isso equivale a dizer que o homomorfismo $\sh{O}_{X'}|g^{-1}(U)\to\sh{L}'|g^{-1}(U)$ definido por $f'|g^{-1}(U)$ é um isomorfismo, isto é, à relação $g^{-1}(U)\subset X'_{f'}$.
\end{proof}

\begin{proposition}[6.5.8]
\label{II.6.5.8}
Sejam $g:X'\to X$ um morfismo finito e $f:Y\to X$ um morfismo;
sejam $Y'=X'_{(Y)}$, $g'=g_{(Y)}$ e $f'=f_{(X')}$, de modo que o diagrama
\[
  \xymatrix{
    X' \ar[d]_{g}
    & Y' \ar[l]_{f'} \ar[d]^{g'}
  \\X
    & Y \ar[l]^{f}
  }
\]
seja comutativo.

Suponhamos que $g_*(\sh{O}_{X'})$ seja um $\sh{O}_X$-módulo localmente livre, ou que $(X,\sh{O}_X)$ e $g_*(\sh{O}_{X'})$ satisfaçam (II).
Suponhamos, além disso, que $Y$ seja um pré-esquema localmente noetheriano reduzido e que a restrição de $f$ a toda componente irredutível de $Y$ seja um morfismo dominante em uma componente irredutível de $X$.
Então, para todo $\sh{O}_{X'}$-módulo invertível $\sh{L}'$, tem-se
\[
  N_{Y'/Y}({f'}^*(\sh{L}')) = f^*(N_{X'/X}(\sh{L}'))
\]
a menos de isomorfismo canônico.
\end{proposition}

\begin{proof}
Observemos que $f^*(g_*(\sh{L}'))=g'_*({f'}^*(\sh{L}'))$, em virtude de \sref{II.1.5.2}, e, em particular, $g'_*(\sh{O}_{Y'})=f^*(g_*(\sh{O}_{X'}))$;
se $g_*(\sh{O}_{X'})$ é localmente livre, o mesmo ocorre com $g'_*(\sh{O}_{Y'})$.
A conclusão resulta então das definições e de \sref{II.6.5.4}.
\end{proof}

\begin{remark}[6.5.9]
\label{II.6.5.9}
Generalizaremos posteriormente a noção de norma desenvolvida acima e a relacionaremos com a noção de imagem direta de um divisor.
\end{remark}


\subsection{Aplicação: critérios de amplitude}
\label{subsection:II.6.6}

\begin{proposition}[6.6.1]
\label{II.6.6.1}
Sejam $Y$ um pré-esquema, $f:X\to Y$ um morfismo quase-compacto e $g:X'\to X$ um morfismo finito e sobrejetivo.
Suponhamos que $g_*(\sh{O}_{X'})$ seja um $\sh{O}_X$-módulo localmente livre, ou que $(X,\sh{O}_X)$ e $g_*(\sh{O}_{X'})$ satisfaçam (II~\emph{bis}).
Então, para todo $\sh{O}_{X'}$-módulo invertível $\sh{L}'$ amplo relativamente a $f\circ g$, $\sh{L}=N_{X'/X}(\sh{L}')$ é amplo relativamente a $f$.
\end{proposition}

\begin{proof}
Podemos supor que $Y$ seja afim \sref{II.4.6.4}, e então, em virtude de \sref{II.4.6.6}, o enunciado é equivalente ao seguinte:
\end{proof}

\begin{corollary}[6.6.2]
\label{II.6.6.2}
Sejam $X$ um pré-esquema quase-compacto e $g:X'\to X$ um morfismo finito sobrejetivo, tal que $g_*(\sh{O}_{X'})$ seja um $\sh{O}_X$-módulo localmente livre, ou que $(X,\sh{O}_X)$ e $g_*(\sh{O}_{X'})$ satisfaçam (II~\emph{bis}).
Então, para todo $\sh{O}_{X'}$-módulo amplo $\sh{L}'$, $\sh{L}=N_{X'/X}(\sh{L}')$ é amplo.
\end{corollary}

\begin{proof}
Na segunda hipótese, podemos supor, além disso, que o homomorfismo canônico $g_*(\sh{O}_{X'})\to g_*(\sh{O}_{X'})\otimes_{\sh{O}_X}\sh{R}(X)$ seja \emph{injetivo}.
Com efeito, caso contrário, seja $\sh{T}$
\oldpage[II]{131}
o núcleo desse homomorfismo, que é um ideal coerente de $\sh{B}=g_*(\sh{O}_{X'})$ \sref[I]{I.6.1.1}, e ponhamos $X''=\Spec(\sh{B}/\sh{T})$;
temos, portanto, um diagrama comutativo
\[
  \xymatrix{
    X'' \ar[rr]^{h} \ar[dr]_{g'}
    && X' \ar[dl]^{g}
  \\&X
  }
\]
em que $h$ é uma imersão fechada \sref{II.1.4.10}.
Além disso, sabemos que o suporte de $\sh{T}$ é um conjunto fechado \sref[0]{0.5.2.2} cujo interior em $X$ é vazio \sref[I]{I.7.4.6}; daí concluímos que, para o ponto genérico $x$ de uma componente irredutível de $X$, existe uma vizinhança aberta afim $U$ de $x$ tal que $\sh{B}|U=(\sh{B}/\sh{T})|U$.
Como $g$ é, por hipótese, sobrejetivo, concluímos que $x\in g'(X'')$;
logo $g'$ é dominante e, sendo um morfismo finito, é \emph{sobrejetivo} \sref{II.6.1.10};
por definição,
\[
  g'_*(\sh{O}_{X''})\otimes_{\sh{O}_X}\sh{R}(X)
  = (\sh{B}/\sh{T})\otimes_{\sh{O}_X}\sh{R}(X)
  = g_*(\sh{O}_{X'})\otimes_{\sh{O}_X}\sh{R}(X)
\]
assim, $(X,\sh{O}_X)$ e $g'_*(\sh{O}_{X''})$ satisfazem (II~\emph{bis}) e, além disso, $g'_*(\sh{O}_{X''})\to g'_*(\sh{O}_{X''})\otimes_{\sh{O}_X}\sh{R}(X)$ é injetivo.
Finalmente, $h^*(\sh{L}')=\sh{L}''$ é um $\sh{O}_{X''}$-módulo amplo \sref{II.4.6.13}[(i~\emph{bis})], e temos $N_{X''/X}(\sh{L}'')=N_{X'/X}(\sh{L}')$.
Com efeito, para definir esses dois $\sh{O}_X$-módulos invertíveis, podemos utilizar o mesmo recobrimento aberto afim $(U_\lambda)$ de $X$, de modo que as restrições de $g_*(\sh{L}')$ e $g'_*(\sh{L}'')$ a $U_\lambda$ sejam isomorfas, respectivamente, a $\sh{B}|U_\lambda$ e $(\sh{B}/\sh{T})|U_\lambda$.
Vê-se imediatamente que a todo isomorfismo $\eta_\lambda:g_*(\sh{L}')|U_\lambda\to\sh{B}|U_\lambda$ corresponde canonicamente um isomorfismo
\[
  \eta'_\lambda : g'_*(\sh{L}'')|U_\lambda \to (\sh{B}/\sh{T})|U_\lambda
\]
de modo que, se $(\omega_{\lambda\mu})$ e $(\omega'_{\lambda\mu})$ são os $1$-cociclos correspondentes aos sistemas de isomorfismos $(\eta_\lambda)$ e $(\eta'_\lambda)$ \sref{II.6.5.2}, $\omega'_{\lambda\mu}$ é a imagem canônica, em $\Gamma(U_\lambda\cap U_\mu,\sh{B}/\sh{T})$, de $\omega_{\lambda\mu}\in\Gamma(U_\lambda\cap U_\mu,\sh{B})$.
Em virtude da definição de $\sh{T}$, concluímos, portanto, que
\[
  N_{\sh{B}/\sh{A}}(\omega_{\lambda\mu}) = N_{(\sh{B}/\sh{T})/\sh{A}}(\omega'_{\lambda\mu})
\]
(onde $\sh{A}=\sh{O}_X$), donde a igualdade anunciada.

Suponhamos, pois, que o homomorfismo $g_*(\sh{O}_{X'})\to g_*(\sh{O}_{X'})\otimes_{\sh{O}_X}\sh{R}(X)$ seja injetivo quando estamos na hipótese~(II~\emph{bis}).
Basta demonstrar que, quando $f$ percorre as seções de $\sh{L}^{\otimes n}$ (para $n>0$) sobre $X$, os $X_f$ formam uma base da topologia de $X$ \sref{II.4.5.2}.
Ora, sejam $x\in X$ e $U$ uma vizinhança qualquer de $x$;
como $g^{-1}(x)$ é finito \sref{II.6.1.7} e $\sh{L}'$ é amplo, existem um inteiro $n>0$ e uma seção $f'$ de ${\sh{L}'}^{\otimes n}$ sobre $X'$ tais que $X'_{f'}$ seja uma vizinhança de $g^{-1}(x)$ contida em $g^{-1}(U)$ \sref{II.4.5.4}.
Como
\[
  \sh{L}^{\otimes n} = N_{X'/X}({\sh{L}'}^{\otimes n}),
\]
basta tomar $f=N_{X'/X}(f')$;
com efeito, então $X\setminus X_f=g(X'\setminus X'_{f'})$ \sref{II.6.5.7}, e, portanto, $x\in X_f\subset U$.
\end{proof}

\begin{corollary}[6.6.3]
\label{II.6.6.3}
Sob as hipóteses de \sref{II.6.6.1}, para que um $\sh{O}_X$-módulo invertível $\sh{L}$ seja amplo relativamente a $f$, é necessário e suficiente que $\sh{L}'=g^*(\sh{L})$ seja amplo relativamente a $f\circ g$.
\end{corollary}

\begin{proof}
A condição é necessária, pois $g$ é afim \sref{II.5.1.12}.
Para demonstrar que a condição é suficiente, podemos supor que $Y$ seja afim \sref{II.4.6.4}, de modo que $X$ e $X'$ sejam quase-compactos e $\sh{L}'$ seja amplo \sref{II.4.6.6}; resta demonstrar que $\sh{L}$ é amplo. Ora, o conjunto dos pontos
\oldpage[II]{132}
$x\in X$ que admitem uma vizinhança na qual $g_*(\sh{O}_{X'})$ (resp. $g_*(\sh{O}_{X'})\otimes_{\sh{O}_X}\sh{R}(X)$) tem posto dado~$n$, na primeira (resp. segunda) hipótese, é simultaneamente aberto e fechado em $X$; logo $X$ é o pré-esquema soma de um número finito desses abertos, e podemos, portanto, supor que seja igual a um deles \sref{II.4.6.17}.
Mas então $N_{X'/X}(\sh{L}')=\sh{L}^{\otimes n}$; assim, $\sh{L}^{\otimes n}$ é amplo em virtude de \sref{II.6.6.2}, e, portanto, o mesmo ocorre com $\sh{L}$ \sref{II.4.5.6}.
\end{proof}

\begin{corollary}[6.6.4]
\label{II.6.6.4}
Suponhamos satisfeitas as hipóteses de \sref{II.6.6.1} e suponhamos, além disso, que $f:X\to Y$ seja de tipo finito.
Então, para que $f$ seja quase-projetivo, é necessário e suficiente que $f\circ g$ seja quase-projetivo.
Se supusermos, além disso, que $Y$ seja um esquema quase-compacto ou um pré-esquema cujo espaço subjacente seja noetheriano, então, para que $f$ seja projetivo, é necessário e suficiente que $f\circ g$ seja projetivo.
\end{corollary}

\begin{proof}
A hipótese implica que $f\circ g$ é de tipo finito.
Tendo em conta a definição dos morfismos quase-projetivos \sref{II.5.3.1}, a primeira afirmação resulta de \sref{II.6.6.1} e \sref{II.6.6.3}.
Tendo em conta esse resultado e \sref{II.5.5.3}[(ii)], resta demonstrar que, quando $f$ é suposto quase-projetivo, para que $f$ seja próprio, é necessário e suficiente que $f\circ g$ seja próprio.
Ora, $f$ é então separado \sref{II.5.3.1} e de tipo finito;
como $g$ é sobrejetivo, nossa afirmação resulta de \sref{II.5.4.2}[(ii)] e \sref{II.5.4.3}[(ii)].
\end{proof}

Em particular:

\begin{corollary}[6.6.5]
\label{II.6.6.5}
Sejam $X$ um pré-esquema de tipo finito sobre um corpo $K$ e $K'$ uma extensão de grau finito de $K$.
Para que $X$ seja projetivo (resp. quase-projetivo) sobre $K$, é necessário e suficiente que $X'=X\otimes_K K'$ seja projetivo (resp. quase-projetivo) sobre $K'$.
\end{corollary}

\begin{proof}
A condição é, com efeito, necessária (\sref{II.5.3.4}[(iii)] e \sref{II.5.5.5}[(iii)]).
Reciprocamente, suponhamos que ela seja satisfeita, e seja $g:X'\to X$ a projeção canônica.
É claro que $g$ é um morfismo finito \sref{II.6.1.5}[(iii)] e sobrejetivo \sref[I]{I.3.5.2}[(ii)].
Além disso, $g_*(\sh{O}_{X'})$ é um $\sh{O}_X$-módulo localmente livre, por ser isomorfo a $\sh{O}_X\otimes_K K'$ \sref{II.1.5.2}.
Resulta então da hipótese e de \sref{II.6.1.11} e \sref{II.5.5.5}[(ii)] que $X'$ é projetivo (resp. quase-projetivo) \emph{sobre $K$};
deduzimos então de \sref{II.6.6.4} que $X$ é projetivo (resp. quase-projetivo) sobre $K$.
\end{proof}

No Capítulo~V, demonstraremos que o enunciado de \sref{II.6.6.5} continua válido quando $K'$ é uma extensão \emph{arbitrária} de $K$.

O restante deste número é dedicado à demonstração do critério de \sref{II.6.6.11}, que é um refinamento bastante técnico de \sref{II.6.6.1};
pode ser omitido em uma primeira leitura.

\begin{lemma}[6.6.6]
\label{II.6.6.6}
Sejam $X$ um pré-esquema noetheriano reduzido e $\sh{E}$ um $\sh{O}_X$-módulo coerente tal que $\sh{E}\otimes_{\sh{O}_X}\sh{R}(X)$ seja um $\sh{R}(X)$-módulo localmente livre de posto constante~$n$.
Existe então um pré-esquema noetheriano reduzido $Z$ e um morfismo finito birracional $h:Z\to X$ que possui a seguinte propriedade: os morfismos de feixes de conjuntos $\sigma_i:\shHom_{\sh{O}_X}(\sh{E},\sh{E})\to\sh{R}(X)$ (para $1\leq i\leq n$) (cf. \sref{II.6.4.10}) levam $\shHom_{\sh{O}_X}(\sh{E},\sh{E})$ à $\sh{O}_X$-álgebra coerente $h_*(\sh{O}_Z)$.
\end{lemma}

\begin{proof}
Consideremos um aberto afim $U$ de $X$, de anel $A(U)=A$;
sejam $E=\Gamma(U,\sh{E})$ e $C_U$ a subálgebra de $R(U)$ gerada pelos $\sigma_i(u)$ quando $u$ percorre $\Hom_A(E,E)$;
vimos em \sref{II.6.4.7.1} que essa $A$-álgebra é um \emph{$A$-módulo de tipo finito}.
Além disso, é claro que a formação das álgebras $C_U$ comuta com as operações de restrição de um aberto afim $U$ a um aberto afim $U'\subset U$.
Definimos, assim, uma sub-$\sh{O}_X$-álgebra \emph{finita} $\sh{C}$ de $\sh{R}(X)$ tal que $\Gamma(U,\sh{C})=C_U$ para todo aberto afim $U$ de $X$. Nós
\oldpage[II]{133}
tomaremos $Z=\Spec(\sh{C})$ e, para $h$, o morfismo estrutural, que é, portanto, finito \sref{II.6.1.2};
como $\sh{C}$ é reduzida, $Z$ é um pré-esquema noetheriano reduzido \sref{II.1.3.8}.
Finalmente, o anel total de frações de $C_U$ é $R(U)$, por definição, e, como $C_U$ está contido no fecho integral de $A(U)$ em $R(U)$, há uma correspondência biunívoca entre os ideais primos minimais de $A(U)$ e os ideais primos minimais de $C_U$ \cite[t.~I, p.~259]{I-13}, o que demonstra que $h$ é birracional e termina a demonstração.
\end{proof}

\begin{corollary}[6.6.7]
\label{II.6.6.7}
Sob as hipóteses de \sref{II.6.6.6}, seja $W$ um aberto de $X$ tal que, para todo $x\in W$, ou $X$ seja normal no ponto $x$, ou $\sh{E}_x$ seja um $\sh{O}_x$-módulo livre.
Podemos então supor $h$ definido de modo que a restrição de $h$ a $h^{-1}(W)$ seja um isomorfismo de $h^{-1}(W)$ sobre $W$.
\end{corollary}

\begin{proof}
Uma ou outra hipótese implica que, se $U\subset W$ é um aberto afim, então, com as notações de \sref{II.6.6.6}, $(\sigma_i(u))_x\in A_x$ para todo $x\in U$ \sref{II.6.4.3}, logo $\sigma_i(u)\in A$, e a conclusão resulta da definição de $h$ dada em \sref{II.6.6.6}.
\end{proof}

\begin{env}[6.6.8]
\label{II.6.6.8}
Sejam $X$ um pré-esquema noetheriano reduzido e $g:X'\to X$ um morfismo finito sobrejetivo, de modo que $\sh{B}=g_*(\sh{O}_{X'})$ seja uma $\sh{O}_X$-álgebra \emph{coerente};
suponhamos, além disso, que $\sh{B}\otimes_{\sh{O}_X}\sh{R}(X)$ seja um $\sh{R}(X)$-módulo \emph{localmente livre de posto constante~$n$}.
Podemos então aplicar o Lema~\sref{II.6.6.6}, tomando $\sh{E}=\sh{B}$, donde, com as notações de \sref{II.6.6.6}, obtemos um homomorfismo de feixes de monoides multiplicativos $\sigma_n:\shHom_{\sh{O}_X}(\sh{B},\sh{B})\to h_*(\sh{O}_Z)$; compondo esse homomorfismo com o homomorfismo canônico $\sh{B}\to\shHom_{\sh{O}_X}(\sh{B},\sh{B})$ \sref{II.6.5.1}, obtemos, portanto, um homomorfismo de feixes de monoides multiplicativos:
\[
\label{II.6.6.8.1}
  N' : \sh{B} = g_*(\sh{O}_{X'}) \to h_*(\sh{O}_Z) = \sh{C}.
\tag{6.6.8.1}
\]

Posto isso, para todo $\sh{O}_{X'}$-módulo invertível $\sh{L}'$, $g_*(\sh{L}')$ é um $\sh{B}$-módulo invertível \sref{II.6.1.12}, e o método de \sref{II.6.5.2} permite associar funtorialmente a $\sh{L}'$ um $\sh{C}$-módulo invertível, que denotaremos por $N'(g_*(\sh{L}'))$.
\end{env}

\begin{lemma}[6.6.9]
\label{II.6.6.9}
Sejam $X$ um pré-esquema noetheriano reduzido e $g:X'\to X$ um morfismo finito sobrejetivo tal que $g_*(\sh{O}_{X'})\otimes_{\sh{O}_X}\sh{R}(X)$ seja um $\sh{R}(X)$-módulo localmente livre de posto constante.
Existe então um pré-esquema noetheriano reduzido $Z$ e um morfismo finito birracional $h:Z\to X$ que possui a seguinte propriedade: para todo $\sh{O}_{X'}$-módulo amplo $\sh{L}'$, o $\sh{O}_Z$-módulo invertível $\sh{M}$ tal que $h_*(\sh{M})=N'(g_*(\sh{L}'))$ (com as notações de \sref{II.6.6.8}) é amplo.
\end{lemma}

\begin{proof}
Suponhamos, primeiramente, que o homomorfismo $\sh{B}\to\sh{B}\otimes_{\sh{O}_X}\sh{R}(X)$ seja \emph{injetivo}.
Definamos $Z$ e $h$ como em \sref{II.6.6.6} (com $\sh{E}=g_*(\sh{O}_{X'})$).
Seja $z\in Z$;
devemos demonstrar que existem um inteiro $m>0$ e uma seção $t$ de $\sh{M}^{\otimes m}$ sobre $Z$ tais que $Z_t$ seja um aberto afim que contenha $z$ \sref{II.4.5.2}.
Sejam $x=h(z)$ e $U$ um aberto afim de $X$ que contenha $x$;
então $h^{-1}(U)$ é uma vizinhança aberta afim de $z$, e basta encontrar $t$ tal que $z\in Z_t\subset h^{-1}(U)$, pois $Z_t$ será então necessariamente afim \sref{II.5.5.8}.
Existem, por hipótese, um inteiro $n>0$ e uma seção $s'$ de ${\sh{L}'}^{\otimes n}$ sobre $X'$ tais que
\[
\label{II.6.6.9.1}
  g^{-1}(x) \subset X'_{s'} \subset g^{-1}(U)
\tag{6.6.9.1}
\]
em virtude de \sref{II.4.5.4}.
Por definição, $s'$ é também uma seção de $g_*({\sh{L}'}^{\otimes n})$ sobre $X$, e a ela corresponde, como em \sref{II.6.5.2}, uma seção
\[
  s=N'(s')
  \quad\text{de}\quad
  N'\bigl(g_*({\sh{L}'}^{\otimes n})\bigr)
  =N'(g_*(\sh{L}'))^{\otimes n}
\]
sobre
\oldpage[II]{134}
$X$.
Veremos que, se $t$ é a seção $s$ considerada como seção de $\sh{M}^{\otimes n}$ sobre $Z$, então $t$ satisfaz o que se pede.
Ponhamos
\[
\label{II.6.6.9.2}
  V = X\setminus g(X'\setminus X'_{s'})
\tag{6.6.9.2}
\]
que é um aberto de $X$ que contém $x$ e está contido em $U$, em virtude de \sref{II.6.6.9.1} e \sref{II.6.1.10}.
Vamos demonstrar que
\[
\label{II.6.6.9.3}
  h^{-1}(V) \subset Z_t \subset h^{-1}(U),
\tag{6.6.9.3}
\]
o que terminará a demonstração.
Isso equivale a dizer que o conjunto $T$ dos $y\in X$ tais que $s_y$ seja invertível contém $V$ e está contido em $U$.
Para isso, consideremos, primeiramente, um aberto afim $W$ contido em $V$;
então $g^{-1}(W)$ é um aberto afim em $X'$ e, em virtude de \sref{II.6.6.9.2}, $s'_{y'}$ é invertível para todo $y'\in g^{-1}(W)$;
pelas hipóteses sobre $X$ e $\sh{B}$, podemos aplicar os resultados de \sref{II.6.4.7}, e vemos que, se $y=g(y')$, então $s_y$ é invertível;
em outras palavras, $V\subset T$.
Por outro lado, resulta também de \sref{II.6.4.7} que, reciprocamente, se $s_y$ é invertível, o mesmo ocorre com $s'_{y'}$, o que implica $y'\in g^{-1}(U)$ por \sref{II.6.6.9.1} e, portanto, $y\in U$, donde $T\subset U$ neste caso.

Passamos daí ao caso geral pelo mesmo raciocínio que em \sref{II.6.6.2}, substituindo $X'$ por $X''$ de modo que $g'_*(\sh{O}_{X''})\to g'_*(\sh{O}_{X''})\otimes_{\sh{O}_X}\sh{R}(X)$ seja injetivo, e $\sh{L}'$ por um $\sh{O}_{X''}$-módulo amplo $\sh{L}''$ tal que $N'(g_*(\sh{L}'))=N'(g'_*(\sh{L}''))$.
O Lema~\sref{II.6.6.9} fica assim demonstrado em todos os casos (com uma escolha conveniente de $h$).
\end{proof}

\begin{corollary}[6.6.10]
\label{II.6.6.10}
Suponhamos satisfeitas as hipóteses de \sref{II.6.6.9};
para todo $\sh{O}_X$-módulo invertível $\sh{L}$ tal que $g^*(\sh{L})$ seja amplo, $h^*(\sh{L})$ é amplo.
\end{corollary}

\begin{proof}
Se pusermos $\sh{L}'=g^*(\sh{L})$, teremos $g_*(\sh{L}')=\sh{L}\otimes_{\sh{O}_X}\sh{B}$ \sref[0]{0.5.4.10}, logo
\[
  N'(g_*(\sh{L}')) = (\sh{L}\otimes_{\sh{O}_X}\sh{C})^{\otimes n}
\]
(pelo mesmo raciocínio que para \sref{II.6.5.2.4}).
Concluímos, portanto, que $\sh{M}=(h^*(\sh{L}))^{\otimes n}$ e, como $\sh{M}$ é amplo, o mesmo ocorre com $h^*(\sh{L})$ \sref{II.4.5.6}.
\end{proof}

\begin{proposition}[6.6.11]
\label{II.6.6.11}
Sejam $Y$ um esquema afim, $X$ um pré-esquema noetheriano reduzido, $f:X\to Y$ um morfismo quase-compacto e $g:X'\to X$ um morfismo finito sobrejetivo.
Seja $W$ um subconjunto aberto de $X$ tal que, para todo $x\in W$, ou $X$ seja normal no ponto $x$, ou exista uma vizinhança aberta $T\subset W$ de $x$ tal que $(g_*(\sh{O}_{X'}))|T$ seja um $(\sh{O}_X|T)$-módulo localmente livre.
Existe então um $Y$-pré-esquema reduzido $Z$ e um $Y$-morfismo \emph{finito e birracional} $h:Z\to X$ tal que a restrição de $h$ a $h^{-1}(W)$ seja um isomorfismo de $h^{-1}(W)$ sobre $W$ e que possua a seguinte propriedade: para todo $\sh{O}_X$-módulo invertível $\sh{L}$ tal que $g^*(\sh{L})$ seja amplo relativamente a $f\circ g$, $h^*(\sh{L})$ é amplo relativamente a $f\circ h$.
\end{proposition}

\begin{proof}
Como $Y$ é afim, $g^*(\sh{L})$ é amplo, e trata-se, portanto, de demonstrar que, para uma escolha conveniente de $h$, $h^*(\sh{L})$ é amplo \sref{II.4.6.6}.
Vamos demonstrar que podemos substituir $g$ por um morfismo finito sobrejetivo $g':X''\to X$ tal que ${g'}^*(\sh{L})$ seja amplo e $g'_*(\sh{O}_{X''})\otimes_{\sh{O}_X}\sh{R}(X)$ seja um $\sh{R}(X)$-módulo \emph{localmente livre de posto constante};
estaremos então reduzidos às condições de \sref{II.6.6.10}, e a proposição estará demonstrada.

Para isso, seja $\sh{B}=g_*(\sh{O}_{X'})$;
designemos por $X_i$ (para $1\leq i\leq n$) os subpré-esquemas fechados reduzidos de $X$ cujos espaços subjacentes são as componentes irredutíveis de $X$ \sref[I]{I.5.2.1};
eles são integrais por hipótese.
Seja $X'_i$ o subpré-esquema fechado
\oldpage[II]{135}
$g^{-1}(X_i)$ de $X'$, e seja $g_i:X'_i\to X_i$ a restrição de $g$ a $X'_i$, que é finita \sref{II.6.1.5}[(iii)] e sobrejetiva;
seja $k_i$ o \emph{posto} da $\sh{O}_{X_i}$-álgebra $\sh{B}_i=(g_i)_*(\sh{O}_{X'_i})$.
Como $\sh{B}_i\otimes_{\sh{O}_{X_i}}\sh{R}(X_i)$ é um pré-feixe constante \sref[I]{I.7.3.5}, o posto $k_i$ é também o posto da $(\sh{O}_X|U)$-álgebra $\sh{B}|U$ para todo aberto $U$ de $X$ que intersecte somente a componente irredutível $X_i$.
Se $T$ é um subconjunto aberto de $X$ tal que $\sh{B}|T$ seja isomorfo a $\sh{O}_X^m|T$, resulta da observação precedente que os números $k_i$ são iguais a $m$ para todos os índices $i$ tais que $T\cap X_i\neq\emp$.
Seja então $U$ o aberto de $X$ formado pelos pontos que admitem uma vizinhança na qual $\sh{B}$ é um $\sh{O}_X$-módulo localmente livre, e sejam $U_j$ ($1\leq j\leq s$) suas componentes conexas, que são abertas em $X$ e em número finito (pois $U$ é noetheriano);
designemos por $V_j$ o subpré-esquema fechado de $X'$ que é o \emph{fecho} do subpré-esquema induzido sobre o aberto $g^{-1}(U_j)$ \sref[I]{I.9.5.11}.
De acordo com o que precede, para todos os índices $i$ tais que $X_i\cap U_j\neq\emp$, os postos $k_i$ são todos iguais a um mesmo inteiro $m_j$;
observemos, além disso, que um mesmo $X_i$ não pode intersectar dois $U_j$ de índices distintos.
Sejam $i_\lambda$ os índices $i$ tais que $X_i\cap U=\emp$.
Consideremos o produto $k$ de todos os $k_i$, ponhamos $n_i=k/k_i$, e seja $X''$ o pré-esquema definido da seguinte maneira.
Para cada $j$ ($1\leq j\leq s$), consideram-se $k/m_j$ pré-esquemas isomorfos a $V_j$, e, para cada $\lambda$, $k/k_{i_\lambda}=n_{i_\lambda}$ pré-esquemas isomorfos a $X'_{i_\lambda}$;
$X''$ é a \emph{soma} de todos esses pré-esquemas.
Define-se um morfismo $g'':X''\to X'$ cuja restrição a cada um dos somandos de $X''$ é a injeção canônica;
é claro que $g''$ é um morfismo finito dominante, portanto sobrejetivo (pois um morfismo finito é fechado \sref{II.6.1.10});
ponhamos $g'=g\circ g''$, que é um morfismo finito sobrejetivo $X''\to X$.
Tem-se
\[
  {g'}^*(\sh{L})={g''}^*(g^*(\sh{L})),
\]
logo ${g'}^*(\sh{L})$ é um $\sh{O}_{X''}$-módulo amplo \sref{II.5.1.12}.
É então claro que, para esse novo pré-esquema $X''$, os postos definidos como os $k_i$ para $X'$ são todos \emph{iguais a $k$};
levando em conta \sref[I]{I.7.3.3}, conclui-se imediatamente que, para todo aberto afim $T$ de $X$,
\[
  (g'_*(\sh{O}_{X''})\otimes_{\sh{O}_X}\sh{R}(X))|T
\]
é um $(\sh{R}(X)|T)$-módulo isomorfo a $(\sh{R}(X)|T)^k$.
\end{proof}

\begin{corollary}[6.6.12]
\label{II.6.6.12}
Se, no enunciado de \sref{II.6.6.11}, tem-se $W=X$, então, para que um $\sh{O}_X$-módulo invertível $\sh{L}$ seja amplo relativamente a $f$, é necessário e suficiente que $g^*(\sh{L})$ seja amplo relativamente a $f\circ g$.
\end{corollary}

\begin{remark}[6.6.13]
\label{II.6.6.13}
No Capítulo~III, veremos que, se $Y$ é noetheriano, $f$ é de tipo finito e a restrição de $f$ ao subpré-esquema fechado reduzido de $X$ cujo espaço subjacente é $X\setminus W$ é \emph{própria}, então a conclusão de \sref{II.6.6.12} ainda é válida.
Mas daremos, no Capítulo~V, exemplos de esquemas algébricos $X$ sobre um corpo $K$ (não sendo próprio o morfismo estrutural $X\to\Spec(K)$) cuja normalização $X'$ é quase-afim, mas que não são quase-afins (de modo que $\sh{O}_X$ não é amplo, embora $\sh{O}_{X'}$ o seja \sref{II.5.1.2} e o morfismo $X'\to X$ seja finito sobrejetivo \sref{II.6.3.10}).
Veremos no número seguinte que essa circunstância não pode ocorrer quando se substitui «quase-afim» por «afim».
\end{remark}

\subsection{O teorema de Chevalley}
\label{subsection:II.6.7}

Vamos estabelecer (com o auxílio do critério de Serre \sref{II.5.2.1}) o teorema seguinte, demonstrado por C.~Chevalley por outros métodos no caso dos esquemas \emph{algébricos}.
\oldpage[II]{136}
\begin{theorem}[6.7.1]
\label{II.6.7.1}
Sejam $X$ um esquema afim, $Y$ um pré-esquema noetheriano e $f:X\to Y$ um morfismo finito sobrejetivo.
Então $Y$ é um esquema afim.
\end{theorem}

\begin{proof}
É claro que
\[
  f_\red:X_\red\to Y_\red
\]
é finito \sref{II.6.1.5}[(vi)];
como $X_\red$ é um esquema afim, e dizer que $Y$ é afim equivale a dizer que $Y_\red$ é afim, pois $Y$ é noetheriano \sref[I]{I.6.1.7}, vê-se que se pode supor $Y$ \emph{reduzido}.
Para toda parte fechada $Y'$ de $Y$, existe então um único subpré-esquema reduzido de $Y$ que tem $Y'$ como espaço subjacente \sref[I]{I.5.1.2};
sua imagem inversa $f^{-1}(Y')$, canonicamente isomorfa a $X\times_Y Y'$ \sref[I]{I.4.4.1}, é afim como subpré-esquema fechado de $X$, e a restrição de $f$ a $f^{-1}(Y')$, que se identifica com $f\times_Y1_{Y'}$, é um morfismo finito sobrejetivo \sref{II.6.1.5}[(iii)].
Em virtude do princípio de indução noetheriana \sref[0]{0.2.2.2}, pode-se, portanto (levando em conta \sref[I]{I.6.1.7}), reduzir à demonstração do teorema sob a hipótese de que, para \emph{toda} parte fechada $Y'\neq Y$, todo subpré-esquema fechado de $Y$ que tenha $Y'$ como espaço subjacente seja afim.
Conclui-se daí que, \emph{para todo $\sh{O}_Y$-módulo coerente $\sh{F}$ cujo suporte (fechado) $Z$ seja distinto de $Y$, tem-se $\HH^1(Y,\sh{F})=0$}.
Com efeito, existe um subpré-esquema fechado $Y'$ de $Y$ que tem $Z$ como espaço subjacente e tal que, se $j:Y'\to Y$ é a injeção canônica, tenha-se
\[
  \sh{F}=j_*(j^*(\sh{F}))
\]
\sref[I]{I.9.3.5};
por conseguinte,
\[
  \HH^1(Y,\sh{F})=\HH^1(Y',j^*(\sh{F}))=0
\]
por \sref{II.5.2.3} e \sref[I]{I.5.1.9.2}.

Suponhamos primeiro que $Y$ não seja irredutível, e sejam $Y'$ uma componente irredutível de $Y$ e $Y''=Y\setminus Y'$;
designaremos também por $Y'$ o subpré-esquema fechado reduzido de $Y$ que tem $Y'$ como espaço subjacente, e por $j$ a injeção canônica $Y'\to Y$.
Seja $\sh{F}$ um $\sh{O}_Y$-módulo coerente, e consideremos o homomorfismo canônico
\[
  \rho:\sh{F}\to\sh{F}'=j_*(j^*(\sh{F}))
\]
\sref[0]{0.4.4.3};
$\sh{F}'$ é um $\sh{O}_Y$-módulo coerente em virtude de \sref[0]{0.5.3.10} e \sref[0]{0.5.3.12}, pois
\[
  j_*(\sh{O}_{Y'})=\sh{O}_Y/\sh{J},
\]
onde $\sh{J}$ designa o feixe de ideais de $\sh{O}_Y$ que define o subpré-esquema $Y'$.
Por conseguinte, $\sh{G}=\Ker\rho$ e $\sh{K}=\Im\rho$ são também $\sh{O}_Y$-módulos coerentes \sref[0]{0.5.3.4};
ora, por definição, a fibra $\sh{F}'_y$ de $\sh{F}'$ no ponto genérico $y$ de $Y'$ é igual a $\sh{F}_y$, pois $y$ é interior a $Y'$ e, portanto, $\sh{J}_y=0$, já que $Y$ é reduzido.
Conclui-se que $y$ não pertence ao suporte (fechado) de $\sh{G}$;
por outro lado, o suporte de $\sh{F}'$ (e, \emph{a fortiori}, o de $\sh{K}$) está contido em $Y'$;
em outras palavras, os suportes de $\sh{G}$ e de $\sh{K}$ são \emph{distintos de $Y$}.
Deduz-se que
\[
  \HH^1(Y,\sh{G})=\HH^1(Y,\sh{K})=0,
\]
e a sequência exata de cohomologia aplicada à sequência exata
\[
  0\to\sh{G}\to\sh{F}\to\sh{K}\to0
\]
mostra que $\HH^1(Y,\sh{F})=0$.
Conclui-se então pelo critério de Serre \sref{II.5.2.1}.

Suponhamos, portanto, $Y$ irredutível e, por conseguinte, \emph{íntegro}.
Pode-se também supor que $X$ seja \emph{íntegro}:
com efeito, se designamos por $X_i$ os subpré-esquemas fechados reduzidos de $X$ que têm como espaços subjacentes as componentes irredutíveis de $X$ \sref[I]{I.5.2.1}, e por $f_i$ a restrição de $f$ a $X_i$, ao menos um dos $f_i$ é dominante e, como é um morfismo finito \sref{II.6.1.5}, é sobrejetivo \sref{II.6.1.10};
como, por outro lado, $X_i$ é um esquema afim, vê-se que se pode substituir $X$ por $X_i$ no enunciado.

\begin{lemma}[6.7.1.1]
\label{II.6.7.1.1}
Sejam $X$ e $Y$ dois pré-esquemas íntegros noetherianos, $x$ (resp. $y$) o ponto genérico de $X$ (resp. de $Y$), e $f:X\to Y$ um morfismo finito sobrejetivo.
Seja $\sh{L}$ um $\sh{O}_X$-módulo invertível tal que existam uma vizinhança aberta afim $U$ de $y$ e uma seção $g\in\Gamma(X,\sh{L})$ para as quais
\oldpage[II]{137}
\[
  x\in X_g\subset f^{-1}(U).
\]
Existem então inteiros $m,n>0$, um homomorfismo
\[
  u:\sh{O}_Y^m\to f_*(\sh{L}^{\otimes n}),
\]
e uma vizinhança aberta $V$ de $y$ tais que a restrição $u|V$ seja um isomorfismo de $\sh{O}_Y^m|V$ sobre $f_*(\sh{L}^{\otimes n})|V$.
\end{lemma}

\begin{proof}
Seja $C$ o anel (íntegro) de $U$, e seja $k=\sh{O}_y$ seu corpo de frações:
como $f$ é finito, $U'=f^{-1}(U)$ é afim \sref{II.1.3.2};
seja $D$ seu anel (íntegro), de corpo de frações $K=\sh{O}_x$;
por hipótese, $D$ é um $C$-módulo de tipo finito \sref{II.6.1.4}, logo $K$ é uma extensão de posto finito de $k$.
A fibra
\[
  f^{-1}(y)=X\times_Y\Spec(\kres(y))=X\times_Y\Spec(\sh{O}_y)
\]
identifica-se com $\Spec(K)$ \sref[I]{I.3.6.6};
sejam $s_i$ ($1\leq i\leq m$) elementos de $D$ que formam uma base de $K$ sobre $k$.
Existe $n>0$ tal que as seções
\[
  (s_i|X_g)g^{\otimes n}
\]
de $\sh{L}^{\otimes n}$ sobre $X_g$ prolonguem-se a seções $b_i$ ($1\leq i\leq m$) de $\sh{L}^{\otimes n}$ sobre $X$ \sref[I]{I.9.3.1}.
Os $b_i$ são também, por definição, seções de $f_*(\sh{L}^{\otimes n})$ sobre $Y$ e definem, portanto, um homomorfismo
\[
  u:\sh{O}_Y^m\to f_*(\sh{L}^{\otimes n})
\]
\sref[0]{0.5.1.1};
veremos que $u$ responde à questão.
Tem-se
\[
  \sh{L}^{\otimes n}|U'=\widetilde{M},
\]
onde $M$ é um $D$-módulo de tipo finito;
se $\vphi$ é a injeção $C\to D$ correspondente ao morfismo $f^{-1}(U)\to U$ dado pela restrição de $f$, então $M_{[\vphi]}$ é um $C$-módulo de tipo finito.
Por conseguinte,
\[
  f_*(\sh{L}^{\otimes n})|U=(M_{[\vphi]})\supertilde
\]
é coerente \sref[I]{I.1.6.3} e, como $U$ é um aberto afim arbitrário de $Y$, $f_*(\sh{L}^{\otimes n})$ é coerente.
Além disso, $u|U=\widetilde{\theta}$, onde $\theta$ é um $C$-homomorfismo
\[
  C^m\to M_{[\vphi]},
\]
e $u_y=\theta_y$ é o homomorfismo
\[
  \theta\otimes1:k^m=C^m\otimes_Ck\to M_{[\vphi]}\otimes_Ck;
\]
este último é, por definição, um \emph{isomorfismo}, pois os $(b_i)_x$ formam uma base de $M_{[\vphi]}\otimes K$ sobre $k$, sendo $g_x$, por hipótese, um gerador de $(\sh{L}^{\otimes n})_x$.
Conclui-se que os suportes dos $\sh{O}_Y$-módulos $\Ker u$ e $\Coker u$ não contêm $y$;
como esses $\sh{O}_Y$-módulos são coerentes \sref[0]{0.5.3.3}, seus suportes são fechados \sref[0]{0.5.2.2}, donde resulta o lema.
\end{proof}

Posto isso, as hipóteses do lema \sref{II.6.7.1.1} são satisfeitas no caso considerado tomando-se $\sh{L}=\sh{O}_X$, pois $X$ é afim \sref[I]{I.1.1.10};
ponhamos $\sh{A}=\sh{O}_Y$ e $\sh{B}=f_*(\sh{O}_X)$.
Em virtude do critério de Serre \sref{II.5.2.1.1}, basta provar que, para todo $\sh{O}_Y$-módulo coerente $\sh{F}$, tem-se $\HH^1(Y,\sh{F})=0$;
basta mesmo prová-lo quando $\sh{F}\subset\sh{O}_Y$, o que implica que $\sh{F}$ é sem torção, pois $Y$ é íntegro;
na verdade, mostraremos que $\HH^1(Y,\sh{F})=0$ para \emph{todo} $\sh{O}_Y$-módulo coerente \emph{sem torção} $\sh{F}$.
Ora, o homomorfismo $u:\sh{A}^m\to\sh{B}$ define um homomorfismo
\[
  v:\sh{G}=\shHom_{\sh{A}}(\sh{B},\sh{F})
  \to
  \shHom_{\sh{A}}(\sh{A}^m,\sh{F})
  =\sh{F}^m.
\]

Mostremos primeiro que $v$ é \emph{injetivo}:
por hipótese, $\sh{T}=\Coker u$ tem um suporte que não intersecta $V$, sendo, portanto, um $\sh{O}_Y$-módulo \emph{de torção} \sref[I]{I.7.4.6};
a sequência exata
\[
  \sh{A}^m\to\sh{B}\to\sh{T}\to0
\]
dá, pela exatidão à esquerda do funtor $\shHom_{\sh{A}}$, a sequência exata
\[
  0\to\shHom_{\sh{A}}(\sh{T},\sh{F})
  \to\sh{G}\xrightarrow{v}\sh{F}^m.
\]
Mas, como $\sh{F}$ é sem torção, tem-se
\[
  \shHom_{\sh{A}}(\sh{T},\sh{F})=0
\]
\sref[0]{0.5.2.6}, donde resulta nossa afirmação.
Tem-se, portanto, a sequência exata
\[
  0\to\sh{G}\to\sh{F}^m\to\Coker v\to0
\]
\oldpage[II]{138}
onde $\sh{G}$ e $\Coker v$ são $\sh{O}_Y$-módulos coerentes (\sref[0]{0.5.3.4} e \sref[0]{0.5.3.5});
em virtude da sequência exata de coomologia, bastará mostrar que
\[
  \HH^1(Y,\sh{G})=\HH^1(Y,\Coker v)=0
\]
para daí deduzir que
\[
  \HH^1(Y,\sh{F}^m)=(\HH^1(Y,\sh{F}))^m=0,
\]
e, por conseguinte, que $\HH^1(Y,\sh{F})=0$.
Ora, a restrição $v|V$ é um isomorfismo, portanto o suporte de $\Coker v$ é distinto de $Y$, donde
\[
  \HH^1(Y,\Coker v)=0
\]
pela hipótese inicial.
Por outro lado, $\sh{G}$ é um $\sh{B}$-módulo \emph{coerente} \sref[I]{I.9.6.4};
como $X$ é afim sobre $Y$, existe um $\sh{O}_X$-módulo quase-coerente $\sh{K}$ tal que $\sh{G}$ seja isomorfo a $f_*(\sh{K})$ \sref{II.1.4.3};
como $X$ é afim, tem-se $\HH^1(X,\sh{K})=0$ \sref[I]{I.5.1.9.2}, logo também $\HH^1(Y,\sh{G})=0$ por \sref{II.5.2.3}, o que conclui a demonstração do teorema \sref{II.6.7.1}.
\end{proof}

\begin{corollary}[6.7.2]
\label{II.6.7.2}
Sejam $X$ um pré-esquema noetheriano e $(X_i)_{1\leq i\leq n}$ um recobrimento finito do espaço $X$ formado por conjuntos fechados.
Para que $X$ seja afim, é necessário e suficiente que, para cada $i$, exista um subpré-esquema fechado de $X$ que tenha $X_i$ como espaço subjacente e seja afim.
\end{corollary}

\begin{proof}
Com efeito, se assim for, seja $X'$ o esquema \emph{soma} dos $X_i$;
é claro que $X'$ é afim, e define-se um morfismo sobrejetivo $f:X'\to X$ tomando como restrição de $f$ a $X_i$ a injeção canônica.
Em virtude de \sref{II.6.7.1}, tudo se reduz a verificar que $f$ é \emph{finito}, e isso foi visto em \sref{II.6.1.5}.
\end{proof}

\begin{corollary}[6.7.3]
\label{II.6.7.3}
Para que um pré-esquema noetheriano $X$ seja afim, é necessário e suficiente que sejam afins os subpré-esquemas fechados reduzidos cujos espaços subjacentes são as componentes irredutíveis de $X$.
\end{corollary}

\section{Critérios valorativos}
\label{section:II.7}

Neste parágrafo, damos critérios valorativos de separação e da propriedade de ser próprio para um morfismo, isto é, critérios que fazem intervir um esquema auxiliar variável da forma $\Spec(A)$, onde $A$ é um anel de valoração.
Mediante hipóteses «noetherianas» convenientes, pode-se, nesses critérios, restringir ao caso em que $A$ é um anel de valoração \emph{discreta}.
Este será, sem dúvida, o único caso que teremos de aplicar no que segue, e só introduzimos os anéis de valoração arbitrários, no caso geral, para estabelecer a ligação com os desenvolvimentos clássicos.

\subsection{Recordação sobre os anéis de valoração}
\label{subsection:II.7.1}

\begin{env}[7.1.1]
\label{II.7.1.1}
Entre as diversas propriedades equivalentes que caracterizam os anéis de valoração, utilizaremos a seguinte:
diz-se que um anel $A$ é um \emph{anel de valoração} se é um anel íntegro que não é um corpo e se, no conjunto dos anéis locais contidos no corpo de frações $K$ de $A$ e distintos de $K$, $A$ é \emph{maximal} para a relação de dominação \sref[I]{I.8.1.1}.
Recordemos que um anel de valoração é \emph{integralmente fechado}.
Se $A$ é um anel de valoração, $A_\mathfrak{p}$ é também um anel de valoração para todo ideal primo $\mathfrak{p}\neq0$ de $A$.
\end{env}

\begin{env}[7.1.2]
\label{II.7.1.2}
Sejam $K$ um corpo e $A$ um subanel local de $K$ que não seja um corpo;
\oldpage[II]{139}
existe então um anel de valoração que tem $K$ como corpo de frações e domina $A$ (\cite[p.~1-07, lema~2]{I-1}).

Por outro lado, sejam $B$ um anel de valoração, $k$ seu corpo residual, $K$ seu corpo de frações e $L$ uma extensão de $k$.
Existe então um anel de valoração \emph{completo} $C$ que domina $B$ e cujo corpo residual é $L$.
Com efeito, $L$ é uma extensão algébrica de uma extensão transcendente pura
\[
  L'=k(T_\mu)_{\mu\in M};
\]
sabe-se que se pode prolongar a valoração de $K$ correspondente a $B$ a uma valoração de
\[
  K'=K(T_\mu)_{\mu\in M}
\]
de modo que $L'$ seja o corpo residual dessa valoração (\cite[p.~98]{II-24});
substituindo $B$ pelo completamento do anel dessa valoração prolongada, vê-se que se pode limitar ao caso em que $B$ é completo e $L$ é um fecho algébrico de $k$.
Se $\overline{K}$ é um fecho algébrico de $K$, pode-se então prolongar a $\overline{K}$ a valoração que define $B$, e o corpo residual correspondente é um fecho algébrico de $k$, como se vê levantando a $\overline{K}$ os coeficientes de um polinômio mônico de $k[T]$.
Fica-se, portanto, finalmente reduzido ao caso em que $L=k$, e basta então tomar para $C$ o completamento de $B$ para responder à questão.
\end{env}

\begin{env}[7.1.3]
\label{II.7.1.3}
Sejam $K$ um corpo e $A$ um subanel de $K$;
o fecho integral $A'$ de $A$ em $K$ é a interseção dos anéis de valoração que contêm $A$ e têm $K$ como corpo de frações (\cite[p.~51, th.~2]{I-11}).
A proposição \sref{II.7.1.2} exprime-se sob a forma geométrica equivalente seguinte.
\end{env}

\begin{proposition}[7.1.4]
\label{II.7.1.4}
Sejam $Y$ um pré-esquema, $p:X\to Y$ um morfismo, $x$ um ponto de $X$, $y=p(x)$, e $y'\neq y$ uma especialização \sref[0]{0.2.1.2} de $y$.
Existe então um esquema local $Y'$, espectro de um anel de valoração, e um morfismo separado $f:Y'\to Y$ tais que, se $a$ designa o único ponto fechado de $Y'$ e $b$ o ponto genérico de $Y'$, tenha-se
\[
  f(a)=y'
  \qquad\text{e}\qquad
  f(b)=y.
\]
Pode-se, além disso, supor satisfeita uma das duas propriedades adicionais seguintes:
\begin{enumerate}
  \item[\rm{(i)}] $Y'$ é o espectro de um anel de valoração completo cujo corpo residual é algebricamente fechado, e existe um $\kres(y)$-homomorfismo $\kres(x)\to\kres(b)$.
  \item[\rm{(ii)}] Existe um $\kres(y)$-isomorfismo $\kres(x)\simto\kres(b)$.
\end{enumerate}
\end{proposition}

\begin{proof}
Seja $Y_1$ o subpré-esquema fechado reduzido de $Y$ que tem $\overline{\{y\}}$ como espaço subjacente \sref[I]{I.5.2.1}, e seja $X_1$ o subpré-esquema fechado imagem inversa $p^{-1}(Y_1)$;
como $y'\in\overline{\{y\}}$ por hipótese e $\kres(x)$ é o mesmo em $X$ e em $X_1$, pode-se supor que $Y$ seja \emph{íntegro}, de ponto genérico $y$.
$\sh{O}_{y'}$ é então um anel local íntegro que não é um corpo, cujo corpo de frações é $\sh{O}_y=\kres(y)$, e $\kres(x)$ é uma extensão de $\kres(y)$.

Para realizar as condições $f(a)=y'$ e $f(b)=y$, bem como a condição adicional~(i), escolha-se primeiro um anel de valoração $B$ que domine $\sh{O}_{y'}$ e tenha $\kres(x)$ como corpo de frações, e depois um anel de valoração completo $A'$ que domine $B$ e cujo corpo residual seja algebricamente fechado.
Então
\[
  \kres(x)=\Frac(B)\longrightarrow\Frac(A')=\kres(b)
\]
é o $\kres(y)$-homomorfismo requerido.
Para a condição~(ii), tome-se diretamente um anel de valoração $A'$ que domine $\sh{O}_{y'}$ e tenha $\kres(x)$ como corpo de frações.
Em ambos os casos, põe-se $Y'=\Spec(A')$, e a existência dos anéis indicados é garantida por \sref{II.7.1.2}.
\end{proof}

\begin{env}[7.1.5]
\label{II.7.1.5}
Recordemos que um anel local $A$ é dito \emph{de dimensão~1} se existe um ideal primo distinto do ideal maximal $\mathfrak{m}$ e se todo ideal primo de $A$ distinto de $\mathfrak{m}$ é um ideal primo \emph{minimal};
quando $A$ é \emph{íntegro}, isso equivale a dizer que $\mathfrak{m}$ e $(0)$ são os únicos ideais primos e que $\mathfrak{m}\neq(0)$;
em outros termos, $Y=\Spec(A)$ reduz-se a dois
\oldpage[II]{140}
pontos $a$ e $b$: $a$ é o único ponto \emph{fechado}, tem-se $\mathfrak{j}_a=\mathfrak{m}$, e $\kres(a)=k$ é o \emph{corpo residual} $k=A/\mathfrak{m}$;
$b$ é o \emph{ponto genérico} de $Y$, $\mathfrak{j}_b=(0)$, sendo o conjunto $\{b\}$ o único aberto de $Y$ distinto de $\emp$ e de $Y$ (aberto que é, portanto, \emph{denso em toda parte}), e $\kres(b)=K$ é o \emph{corpo de frações} de $A$.
\end{env}

\begin{env}[7.1.6]
\label{II.7.1.6}
Para um anel local noetheriano $A$ de dimensão~1, sabe-se (\cite[pp.~2-08 e 17-01]{I-1}) que as seguintes condições são equivalentes:
\begin{enumerate}
    \item[\rm{(a)}] $A$ é normal;
    \item[\rm{(b)}] $A$ é regular;
    \item[\rm{(c)}] $A$ é um anel de valoração;
\end{enumerate}
além disso, $A$ é então um \emph{anel de valoração discreta}.
As Proposições~\sref{II.7.1.2} e \sref{II.7.1.3} têm então os seguintes análogos para os anéis de valoração discreta:
\end{env}

\begin{proposition}[7.1.7]
\label{II.7.1.7}
Sejam $A$ um anel local integral noetheriano que não seja um corpo, $K$ seu corpo de frações e $L$ uma extensão de tipo finito de $K$;
existe então um anel de valoração discreta que domina $A$ e tem $L$ como corpo de frações.
\end{proposition}

\begin{proof}
Suponhamos, antes de tudo, que $L=K$.
Sejam $\mathfrak{m}$ o ideal maximal de $A$, $(x_1,\ldots,x_n)$ um sistema de geradores não nulos de $\mathfrak{m}$ e $B$ o subanel $A[x_2/x_1,\ldots,x_n/x_1]$ de $K$, que é noetheriano.
É imediato que o ideal $\mathfrak{m}B$ de $B$ é idêntico ao ideal principal $x_1B$;
se $\mathfrak{p}$ é um ideal primo minimal sobre $x_1B$, então $\mathfrak{p}$ tem altura~1 (\cite[t.~I, p.~277]{I-13});
em outras palavras, $B_\mathfrak{p}$ é um anel local noetheriano \emph{de dimensão~1};
é claro que $\mathfrak{p}B_\mathfrak{p}\cap A$ é um ideal de $A$ que contém $\mathfrak{m}$ e não contém $1$, sendo, portanto, igual a $\mathfrak{m}$, e assim $B_\mathfrak{p}$ \emph{domina} $A$ \sref[I]{I.8.1.1}.
Decorre do Teorema de Krull--Akizuki (\cite[p.~293]{II-25}) que o fecho integral $C$ de $B_\mathfrak{p}$ é um anel noetheriano (embora $C$ não seja necessariamente um $B_\mathfrak{p}$-módulo de tipo finito);
se $\mathfrak{n}$ é um ideal maximal de $C$, então $C_\mathfrak{n}$ é um anel local noetheriano normal de dimensão~1 (\cite[p.~295]{II-25}) e, portanto, um anel de valoração discreta que domina $B_\mathfrak{p}$ e, \emph{a fortiori}, $A$.

Agora, se $L$ é uma extensão de tipo finito de $K$, pode-se, pelo que precede, reduzir ao caso em que $A$ já é um anel de valoração discreta.
Seja $w$ uma valoração de $K$ associada a $A$;
existe uma valoração discreta $w'$ de $L$ que \emph{prolonga} $w$: com efeito, por indução sobre o número de geradores de $L$, reduz-se ao caso em que $L=K(\alpha)$, e então a proposição é clássica (\cite[p.~106]{II-24}).
\end{proof}

\begin{corollary}[7.1.8]
\label{II.7.1.8}
Sejam $A$ um anel integral noetheriano, $K$ seu corpo de frações e $L$ uma extensão de tipo finito de $K$.
Então o fecho integral de $A$ em $L$ é a interseção dos anéis de valoração discreta que têm $L$ como corpo de frações e contêm $A$.
\end{corollary}

\begin{proof}
Com efeito, tal anel de valoração discreta, sendo normal, contém \emph{a fortiori} todo elemento de $L$ integral sobre $A$.
Basta, portanto, provar que, se $x\in L$ não é integral sobre $A$, existe um anel de valoração discreta $C$ que tem $L$ como corpo de frações, contém $A$ e não contém $x$.
A hipótese sobre $x$ implica que $x\not\in B=A[1/x]$ ou, em outras palavras, que $1/x$ não é invertível no anel noetheriano $B$.
Existe, portanto, um ideal primo $\mathfrak{p}$ de $B$ que contém $1/x$.
O anel local integral $B_\mathfrak{p}$ é noetheriano e está contido em $L$, que é uma extensão de tipo finito do corpo de frações de $B_\mathfrak{p}$ (este último contendo $K$).
Por \sref{II.7.1.7}, existe, portanto, um anel de valoração discreta $C$ que domina $B_\mathfrak{p}$ e tem $L$ como corpo de frações;
como $1/x\in\mathfrak{p}B_\mathfrak{p}$ pertence ao ideal maximal de $C$, tem-se $x\not\in C$, o que conclui a demonstração.
\end{proof}

A forma geométrica de \sref{II.7.1.7} é a seguinte:
\oldpage[II]{141}
\begin{proposition}[7.1.9]
\label{II.7.1.9}
Sejam $Y$ um pré-esquema localmente noetheriano, $p:X\to Y$ um morfismo localmente de tipo finito, $x$ um ponto de $X$, $y=p(x)$ e $y'\neq y$ uma especialização de $y$.
Existe então um esquema local $Y'$, espectro de um anel de valoração discreta, um morfismo separado $f:Y'\to Y$ e uma $Y$-aplicação racional $g$ de $Y'$ em $X$ tais que, denotando-se por $a$ o ponto fechado de $Y'$ e por $b$ o ponto genérico de $Y'$, tem-se $f(a)=y'$, $f(b)=y$, $g(b)=x$, e tais que, no diagrama comutativo
\[
\label{II.7.1.9.1}
    \xymatrix{
        & \kres(x) \ar[dl]_{\gamma}
    \\  \kres(b)
        & \kres(y) \ar[u]_{\pi} \ar[l]^{\vphi}
    }
\tag{7.1.9.1}
\]
(onde $\pi$, $\vphi$ e $\gamma$ são os homomorfismos correspondentes, respectivamente, a $p$, $f$ e $g$), o homomorfismo $\gamma$ é bijetivo.
\end{proposition}

\begin{proof}
Como em \sref{II.7.1.4}, pode-se reduzir ao caso em que $Y$ é integral, com ponto genérico $y$ (levando em conta \sref[I]{I.6.3.4}[iv]), e, como a questão é local em $X$ e $Y$, pode-se supor que $p$ é de tipo finito;
estamos então na situação de \sref{II.7.1.4}, com a propriedade adicional de que $\kres(x)$ é uma extensão \emph{de tipo finito} de $\kres(y)$ \sref[I]{I.6.4.11} e de que $\sh{O}_{y'}$ é noetheriano;
isso permite aplicar \sref{II.7.1.7} e tomar $Y'=\Spec(A')$, onde $A'$ é um anel de valoração discreta que domina $\sh{O}_{y'}$ e cujo corpo de frações é $\kres(x)$.
Definimos, assim, um diagrama comutativo \sref{II.7.1.9.1} no qual $\gamma$ é uma bijeção, e $\pi$ e $\vphi$ correspondem aos morfismos $p$ e $f$.
Além disso, como $X$ e $Y$ são localmente noetherianos \sref[I]{I.6.6.2} e $Y'$ é integral, existe uma única $Y$-aplicação racional $g$ de $Y'$ em $X$ à qual corresponde o isomorfismo $\gamma$ \sref[I]{I.7.1.15}, o que conclui a demonstração.
\end{proof}

\subsection{Critério valuativo de separação}
\label{subsection:II.7.2}

\begin{proposition}[7.2.1]
\label{II.7.2.1}
Sejam $X$ e $Y$ dois pré-esquemas e $f:X\to Y$ um morfismo quase-compacto.
As duas condições seguintes são equivalentes:
\begin{enumerate}
    \item[\rm{(a)}] O morfismo $f$ é fechado.
    \item[\rm{(b)}] Para todo $x\in X$ e toda especialização $y'\neq y$ de $y=f(x)$, existe uma especialização $x'$ de $x$ tal que $f(x')=y'$.
\end{enumerate}
\end{proposition}

\begin{proof}
A condição~(b) exprime que $f(\overline{\{x\}})=\overline{\{y\}}$ e é, portanto, consequência de~(a).
Para mostrar que (b) implica (a), consideremos um subconjunto fechado $X'$ do espaço subjacente $X$;
seja $Y'=\overline{f(X')}$, e mostremos que $Y'=f(X')$.
Consideremos os subpré-esquemas fechados reduzidos de $X$ e $Y$ cujos espaços subjacentes são, respectivamente, $X'$ e $Y'$ \sref[I]{I.5.2.1};
existe então um morfismo $f':X'\to Y'$ tal que o diagrama
\[
    \xymatrix{
        X' \ar[r]^{f'} \ar[d]
        & Y' \ar[d]
    \\  X \ar[r]_{f}
        & Y
    }
\]
é comutativo \sref[I]{I.5.2.2} e, como $f$ é quase-compacto, $f'$ também o é.
Ficamos, portanto, reduzidos a demonstrar que, se $f$ é um morfismo quase-compacto e \emph{dominante}, então
\oldpage[II]{142}
a condição~(b) implica que $f(X)=Y$.
Ora, seja $y'$ um ponto de $Y$, e seja $y$ o ponto genérico de um componente irredutível de $Y$ que contém $y'$;
em virtude de~(b), basta mostrar que $f^{-1}(y)$ não é vazio.
Mas sabe-se que essa propriedade decorre do fato de $f$ ser quase-compacto e dominante \sref[I]{I.6.6.5}.
\end{proof}

\begin{corollary}[7.2.2]
\label{II.7.2.2}
Seja $f:X\to Y$ uma imersão quase-compacta.
Para que o espaço subjacente $X$ seja fechado em $Y$, é necessário e suficiente que ele contenha toda especialização, em $Y$, de cada um de seus pontos.
\end{corollary}

\begin{proposition}[7.2.3]
\label{II.7.2.3}
Sejam $Y$ um pré-esquema (resp. um pré-esquema localmente noetheriano) e $f:X\to Y$ um morfismo (resp. um morfismo localmente de tipo finito).
As seguintes condições são equivalentes:
\begin{enumerate}
    \item[\rm{(a)}] $f$ é separado.
    \item[\rm{(b)}] O morfismo diagonal $X\to X\times_Y X$ é quase-compacto e, para todo $Y$-pré-esquema da forma $Y'=\Spec(A)$, sendo $A$ um anel de valoração (resp. um anel de valoração discreta), quaisquer dois $Y$-morfismos de $Y'$ em $X$ que coincidam no ponto genérico de $Y'$ são iguais.
    \item[\rm{(c)}] O morfismo diagonal $X\to X\times_Y X$ é quase-compacto e, para todo $Y$-pré-esquema da forma $Y'=\Spec(A)$, sendo $A$ um anel de valoração (resp. um anel de valoração discreta), quaisquer duas $Y'$-seções de $X'=X_{(Y')}$ que coincidam no ponto genérico de $Y'$ são iguais.
\end{enumerate}
\end{proposition}

\begin{proof}
A equivalência entre (b) e (c) decorre da correspondência bijetiva entre os $Y$-morfismos de $Y'$ em $X$ e as $Y'$-seções de $X'$ \sref[I]{I.3.3.14}.
Se $X$ é separado sobre $Y$, a condição~(b) é satisfeita, em virtude de \sref[I]{I.7.2.2.1}, pois $Y'$ é integral.
Resta mostrar que (b) implica que o morfismo diagonal $\Delta:X\to X\times_Y X$ é fechado, o que equivale a mostrar que ele satisfaz o critério de \sref{II.7.2.2}.
Sejam, pois, $z$ um ponto da diagonal $\Delta(X)$ e $z'\neq z$ uma especialização de $z$ em $X\times_Y X$.
Existe então \sref{II.7.1.4} um anel de valoração $A$ e um morfismo $f$ de $Y'=\Spec(A)$ em $X\times_Y X$ tal que $f$ envia o ponto fechado $a$ de $Y'$ a $z'$ e o ponto genérico $b$ de $Y'$ a $z$;
esse morfismo faz de $Y'$ um $(X\times_Y X)$-pré-esquema e, \emph{a fortiori}, um $Y$-pré-esquema.
Compondo com $f$ as duas projeções de $X\times_Y X$, obtêm-se dois $Y$-morfismos $g_1$ e $g_2$ de $Y'$ em $X$ que, por hipótese, coincidem no ponto $b$;
eles são, portanto, iguais a um único morfismo $g$, o que implica \sref[I]{I.5.3.1} que $f$ se fatora como $f=\Delta\circ g$ e, por conseguinte, $z'\in\Delta(X)$.
Se supusermos que $Y$ é localmente noetheriano e $f$ é localmente de tipo finito, então $X\times_Y X$ é localmente noetheriano \sref[I]{I.6.6.7};
pode-se, portanto, retomar o raciocínio precedente supondo que $A$ é um anel de valoração discreta, em virtude de \sref{II.7.1.9}.
\end{proof}

\begin{remark}[7.2.4]
\label{II.7.2.4}
\begin{enumerate}
    \item[\rm{(i)}] A hipótese de que o morfismo $\Delta$ seja quase-compacto é sempre satisfeita quando $Y$ é localmente noetheriano e $f$ é localmente de tipo finito, pois $X\times_Y X$ é então localmente noetheriano \sref[I]{I.6.6.4}[i].
        No caso geral, ela significa também que, para toda cobertura $(U_\alpha)$ de $X$ por abertos afins, os conjuntos $U_\alpha\cap U_\beta$ são \emph{quase-compactos}.
    \item[\rm{(ii)}] Para que $f$ seja separado, \emph{basta} exigir a condição~(b) ou (c) apenas para anéis de valoração $A$ \emph{completos} cujos corpos residuais sejam \emph{algebricamente fechados};
        isso decorre da demonstração de \sref{II.7.2.3} e de \sref{II.7.1.4}.
\end{enumerate}
\end{remark}
\oldpage[II]{143}

\subsection{Critério valuativo de propriedade}
\label{subsection:II.7.3}

\begin{proposition}[7.3.1]
\label{II.7.3.1}
Sejam $A$ um anel de valoração, $Y=\Spec(A)$, $b$ o ponto genérico de $Y$, $X$ um \emph{esquema} integral e $f:X\to Y$ um morfismo \emph{fechado} tal que $f^{-1}(b)$ consista em um único ponto $x$ e que o homomorfismo correspondente $\kres(b)\to\kres(x)$ seja bijetivo.
Então $f$ é um isomorfismo.
\end{proposition}

\begin{proof}
Como $f$ é fechado e dominante, tem-se $f(X)=Y$;
basta \sref[I]{I.4.2.2} provar que, para todo $y'\neq b$ em $Y$, existe \emph{um único} ponto $x'$ tal que $f(x')=y'$ e que o homomorfismo correspondente $\sh{O}_{y'}\to\sh{O}_{x'}$ é bijetivo, pois então $f$ será um homeomorfismo.
Ora, se $f(x')=y'$, então $\sh{O}_{x'}$ é um anel local contido em $K=\kres(x)=\kres(b)$ e que domina $\sh{O}_{y'}$;
este último é o anel local $A_{y'}$ e, portanto, é um anel de valoração \sref{II.7.1.1} que tem $K$ como corpo de frações.
Além disso, $\sh{O}_{x'}\neq K$, pois $x'$ não é o ponto genérico de $X$ \sref[0]{0.2.1.3};
conclui-se, portanto, que $\sh{O}_{x'}=\sh{O}_{y'}$.
Como $X$ é um esquema integral, a igualdade $\sh{O}_{x'}=\sh{O}_{x''}$ implica $x'=x''$ \sref[I]{I.8.2.2}, o que conclui a demonstração.
\end{proof}

\begin{env}[7.3.2]
\label{II.7.3.2}
Sejam $A$ um anel de valoração, $K$ seu corpo de frações, $Y=\Spec(A)$ e $b$ o ponto genérico de $Y$, de modo que $\sh{O}_b=\kres(b)$ é igual a $K$;
seja $f:X\to Y$ um morfismo.
Sabe-se \sref[I]{I.7.1.4} que as \emph{$Y$-seções racionais} de $X$ estão em correspondência bijetiva com os \emph{germes} em $b$ de $Y$-seções definidas em vizinhanças de $b$, donde se obtém uma aplicação canônica
\[
\label{II.7.3.2.1}
    \Gamma_\mathrm{rat}(X/Y) \to \Gamma(f^{-1}(b)/\Spec(K))
\tag{7.3.2.1}
\]
sendo os elementos de $\Gamma(f^{-1}(b)/\Spec(K))$ identificados, por definição \sref[I]{I.3.4.5}, com os \emph{pontos de $f^{-1}(b)=X\otimes_A K$ racionais sobre $K$}.
Quando $f$ é \emph{separado}, decorre de \sref[I]{I.5.4.7} que a aplicação \sref{II.7.3.2.1} é \emph{injetiva}, pois $Y$ é um esquema integral.

Compondo \sref{II.7.3.2.1} com a aplicação canônica $\Gamma(X/Y)\to\Gamma_\mathrm{rat}(X/Y)$ \sref[I]{I.7.1.2}, obtém-se uma aplicação canônica
\[
\label{II.7.3.2.2}
    \Gamma(X/Y) \to \Gamma(f^{-1}(b)/\Spec(K)).
\tag{7.3.2.2}
\]
Quando $f$ é \emph{separado}, essa aplicação é novamente \emph{injetiva} \sref[I]{I.5.4.7}.
\end{env}

\begin{proposition}[7.3.3]
\label{II.7.3.3}
Sejam $A$ um anel de valoração com corpo de frações $K$, $Y=\Spec(A)$, $b$ o ponto genérico de $Y$ e $f:X\to Y$ um morfismo \emph{separado} e \emph{fechado}.
Então a aplicação canônica \sref{II.7.3.2.2} é bijetiva (o que equivale a dizer que ela é \emph{sobrejetiva} e implica que as $Y$-seções racionais de $X$ estão \emph{definidas em toda parte}).
\end{proposition}

\begin{proof}
Seja, pois, $x$ um ponto de $f^{-1}(b)$ que é \emph{racional} sobre $K$.
Como $f$ é separado, também o é o morfismo $f^{-1}(b)\to\Spec(K)$ correspondente a $f$ \sref[I]{I.5.5.1}[iv] e, como toda seção de $f^{-1}(b)$ é uma imersão fechada \sref[I]{I.5.4.6}, $\{x\}$ é fechado \emph{em $f^{-1}(b)$}.
Consideremos o subpré-esquema fechado reduzido $X'$ de $X$ que tem como espaço subjacente o fecho $\overline{\{x\}}$ de $\{x\}$ \emph{em $X$}.
É claro que a restrição de $f$ a $X'$ satisfaz as hipóteses de \sref{II.7.3.1} e é, portanto, um \emph{isomorfismo} de $X'$ sobre $Y$, cujo isomorfismo inverso é a $Y$-seção procurada de $X$.
\end{proof}

\begin{env}[7.3.4]
\label{II.7.3.4}
Para enunciar os dois resultados seguintes, utilizaremos uma terminologia que será justificada e discutida no Capítulo~IV: se $F$ é um subconjunto
\oldpage[II]{144}
de um pré-esquema $Y$, define-se a \emph{codimensão} de $F$ em $Y$, denotada por $\codim_Y F$, como o ínfimo dos inteiros $\dim(\sh{O}_z)$ quando $z$ percorre $F$.
\end{env}

\begin{corollary}[7.3.5]
\label{II.7.3.5}
Sejam $Y$ um pré-esquema localmente noetheriano reduzido e $N$ o conjunto dos pontos $y\in Y$ nos quais $Y$ não é regular \sref[0]{0.4.1.4};
suponha-se que $\codim_Y N\geq2$.
Seja $f:X\to Y$ um morfismo de tipo finito, \emph{separado} e \emph{fechado}, e seja $g$ uma $Y$-seção racional de $X$;
se $Y'$ é o conjunto dos pontos de $Y$ nos quais $g$ não está definida (conjunto que é \emph{fechado} \sref[I]{I.7.2.1}), então $\codim_Y Y'\geq2$.
\end{corollary}

\begin{proof}
Basta provar que $g$ está definida em todo ponto $z\in Y$ tal que $\dim\sh{O}_z\leq1$.
Se $\dim\sh{O}_z=0$, então $z$ é o ponto genérico de um componente irredutível de $Y$ \sref[I]{I.1.1.14} e, portanto, pertence a todo aberto denso em toda parte de $Y$ e, em particular, ao domínio de definição de $g$.
Suponhamos, pois, que $\dim\sh{O}_z=1$; por hipótese, $\sh{O}_z$ é então um anel local noetheriano regular e, portanto, por \sref{II.7.1.6}, um \emph{anel de valoração discreta}.
Seja $Z=\Spec(\sh{O}_z)$;
como $U=Y\setmin Y'$ é denso em toda parte, $U\cap Z$ não é vazio \sref[I]{I.2.4.2};
seja $g'$ a aplicação racional de $Z$ em $X$ induzida por $g$ \sref[I]{I.7.2.8};
basta mostrar que $g'$ é um \emph{morfismo} \sref[I]{I.7.2.9}.
Ora, $g'$ pode ser considerada uma $Z$-seção racional do $Z$-pré-esquema $f^{-1}(Z)=X\times_Y Z$;
é claro que o morfismo $f^{-1}(Z)\to Z$ correspondente a $f$ é fechado, e decorre de \sref[I]{I.5.5.1}[i] que ele é separado;
conclui-se, portanto, de \sref{II.7.3.3} que $g'$ está definida em toda parte;
como $Z$ é reduzido e $X$ é separado sobre $Y$, $g'$ é um morfismo \sref[I]{I.7.2.2}.
\end{proof}

\begin{corollary}[7.3.6]
\label{II.7.3.6}
Sejam $S$ um pré-esquema localmente noetheriano e $X$ e $Y$ dois $S$-pré-esquemas;
suponha-se que $Y$ seja reduzido e, além disso, que o conjunto $N$ dos pontos $y\in Y$ nos quais $Y$ não é regular satisfaça $\codim_Y N\geq2$;
suponha-se, por fim, que o morfismo estrutural $X\to S$ seja próprio.
Seja $f$ uma $S$-aplicação racional de $Y$ em $X$, e seja $Y'$ o conjunto dos pontos de $Y$ nos quais $f$ não está definida;
então $\codim_Y Y'\geq2$.
\end{corollary}

\begin{proof}
Sabe-se \sref[I]{I.7.1.2} que se podem identificar as $S$-aplicações racionais de $Y$ em $X$ com as $Y$-seções racionais de $X\times_S Y$;
como o morfismo estrutural $X\times_S Y\to Y$ é fechado \sref{II.5.4.1}, pode-se aplicar \sref{II.7.3.5}, donde resulta o corolário.
\end{proof}

\begin{remark}[7.3.7]
\label{II.7.3.7}
As hipóteses sobre $Y$ em \sref{II.7.3.5} e \sref{II.7.3.6} são satisfeitas, em particular, quando $Y$ é \emph{normal} \sref[0]{0.4.1.4}, em virtude de \sref{II.7.1.6}.
\end{remark}

Podemos caracterizar os morfismos universalmente fechados (resp. os morfismos próprios) por uma recíproca de \sref{II.7.3.3}:
\begin{theorem}[7.3.8]
\label{II.7.3.8}
Sejam $Y$ um pré-esquema (resp. um pré-esquema localmente noetheriano) e $f:X\to Y$ um morfismo quase-compacto separado (resp. um morfismo de tipo finito).
As seguintes condições são equivalentes:
\begin{enumerate}
    \item[\rm{(a)}] $f$ é universalmente fechado (resp. próprio).
    \item[\rm{(b)}] Para todo $Y$-esquema da forma $Y'=\Spec(A)$, sendo $A$ um anel de valoração (resp. um anel de valoração discreta) com corpo de frações $K$, a aplicação canônica
        \[
            \Hom_Y(Y',X) \to \Hom_Y(\Spec(K),X)
        \]
        correspondente à injeção canônica $A\to K$ é sobrejetiva (resp. bijetiva).
\oldpage[II]{145}
  \item[\rm{(c)}] Para todo $Y$-esquema da forma $Y'=\Spec(A)$, em que $A$ é um anel de valoração (resp. um anel de valoração discreta), a aplicação canônica \sref{II.7.3.2.2} relativa ao $Y'$-pré-esquema $X_{(Y')}$ é sobrejetiva (resp. bijetiva).
\end{enumerate}
\end{theorem}

\begin{proof}
A equivalência entre (b) e (c) decorre imediatamente de \sref[I]{I.3.3.14};
(a) implica (c), pois (a) implica, em qualquer das duas hipóteses, que $f_{(Y')}$ é separado \sref[I]{I.5.5.1}[iv] e fechado, e basta aplicar \sref{II.7.3.3}.
Resta provar que (b) implica (a).

Consideremos primeiro o caso em que $Y$ é arbitrário e $f$ é separado e quase-compacto.
Se $f$ satisfaz a condição~(b), então também a satisfaz
\[
  f_{(Y'')}:X_{(Y'')}\longrightarrow Y'',
\]
em que $Y''$ é um $Y$-pré-esquema arbitrário, em virtude da equivalência entre (b) e (c) e do fato de que
\[
  X_{(Y'')}\times_{Y''}Y'=X\times_Y Y'
\]
para todo morfismo $Y'\to Y''$ \sref[I]{I.3.3.9.1};
como, além disso, $f_{(Y'')}$ é separado e quase-compacto quando $f$ o é (\sref[I]{I.5.5.1}[iv] e \sref[I]{I.6.6.4}[iii]), reduzimo-nos a provar que (b) implica que $f$ é \emph{fechado}.
Para isso, basta verificar a condição~(b) de \sref{II.7.2.1}.
Sejam, pois, $x\in X$ e $y'$ uma especialização de $y=f(x)$, distinta de $y$;
em virtude de \sref{II.7.1.4}, existem um esquema $Y'$, espectro de um anel de valoração, e um morfismo separado $g:Y'\to Y$ tais que, se $a$ designa o ponto fechado e $b$ o ponto genérico de $Y'$, tem-se $g(a)=y'$, $g(b)=y$, e existe um $\kres(y)$-homomorfismo $\kres(x)\to\kres(b)$.
Este último corresponde canonicamente a um $Y$-morfismo $\Spec(\kres(b))\to X$ \sref[I]{I.2.4.6}, e decorre então de (b) que existe um $Y$-morfismo $h:Y'\to X$ ao qual corresponde o morfismo precedente.
Tem-se então $h(b)=x$;
se pusermos $h(a)=x'$, então $x'$ é uma especialização de $x$ e
\[
  f(x')=f(h(a))=g(a)=y'.
\]

Se agora $Y$ é localmente noetheriano e $f$ de tipo finito, a hipótese~(b) implica primeiro que $f$ é \emph{separado}, em virtude de \sref{II.7.2.3}, pois o morfismo diagonal $X\to X\times_Y X$ é quase-compacto \sref{II.7.2.4}.
Além disso, para provar que $f$ é próprio, basta mostrar que
\[
  f_{(Y'')}:X_{(Y'')}\longrightarrow Y''
\]
é \emph{fechado} para todo $Y$-pré-esquema $Y''$ \emph{de tipo finito}, levando em conta \sref{II.5.6.3}.
Como $Y''$ é então localmente noetheriano, pode-se retomar o raciocínio do primeiro caso tomando para $Y'$ o espectro de um anel de valoração discreta e aplicando \sref{II.7.1.9} em vez de \sref{II.7.1.4}.
\end{proof}

\begin{remarks}[7.3.9]
\label{II.7.3.9}
\begin{enumerate}
  \item[\rm{(i)}] Quando $Y$ é um pré-esquema arbitrário e $f$ um morfismo separado, para que $f$ seja universalmente fechado, \emph{basta} que a condição~(b) ou a condição~(c) seja verificada para os anéis de valoração $A$ \emph{completos} cujo corpo residual seja \emph{algebricamente fechado};
  isso decorre, com efeito, da demonstração precedente e de \sref{II.7.1.4}.

  \item[\rm{(ii)}] Do critério~(c) de \sref{II.7.3.8} deduz-se uma nova demonstração, mais próxima dos métodos clássicos, do fato de que um morfismo projetivo $X\to Y$ é fechado \sref{II.5.5.3}.
  Com efeito, pode-se supor que $Y$ é afim e, por conseguinte, identificar $X$ com um subpré-esquema fechado de um fibrado projetivo $\bb{P}_Y^n$ \sref{II.5.3.3};
  para provar que $X\to Y$ é fechado, basta verificar que o mesmo vale para o morfismo estrutural $\bb{P}_Y^n\to Y$, e o critério~(c) de \sref{II.7.3.8}, juntamente com \sref{II.4.1.3.1}, mostra que nos reduzimos a provar o seguinte fato:
  \emph{se $Y$ é o espectro de um anel de valoração $A$, de corpo de frações $K$, todo ponto de $\bb{P}_Y^n$ com valores em $K$ provém, por restrição ao ponto genérico de $Y$, de um ponto de $\bb{P}_Y^n$ com valores em $A$.}
  Ora, todo $\sh{O}_Y$-módulo invertível é trivial \sref[I]{I.2.4.8};
  portanto, resulta de \sref{II.4.2.6} que um ponto de $\bb{P}_Y^n$ com valores em $K$ se identifica com uma classe de sistemas $(\zeta c_0,\zeta c_1,\ldots,\zeta c_n)$ com entradas em $K$, em que $\zeta\ne0$ e os $c_i$ são elementos de $K$, não todos nulos.
  Multiplicando os $c_i$ por um elemento de $A$ de
\oldpage[II]{146}
  valoração conveniente, pode-se supor que todos os $c_i$ pertencem a $A$ e que ao menos um deles é invertível.
  Mas então, por \sref{II.4.2.6}, o sistema $(c_0,\ldots,c_n)$ também define um ponto de $\bb{P}_Y^n$ com valores em $A$, o que demonstra a afirmação.

  \item[\rm{(iii)}] Os critérios \sref{II.7.2.3} e \sref{II.7.3.8} são especialmente cômodos quando se considera o dado de um $Y$-pré-esquema $X$ como equivalente ao dado do funtor que, a um $Y$-pré-esquema $Y'$, associa
  \[
    X(Y')=\Hom_Y(Y',X);
  \]
  esses critérios nos permitirão, por exemplo, provar que, sob certas condições, os ``esquemas de Picard'' são próprios.
\end{enumerate}
\end{remarks}

\begin{corollary}[7.3.10]
\label{II.7.3.10}
Sejam $Y$ um esquema integral (resp. um esquema integral localmente noetheriano), $X$ um esquema integral e $f:X\to Y$ um morfismo dominante.
\begin{enumerate}
  \item[\rm{(i)}] Se $f$ é quase-compacto e universalmente fechado, todo anel de valoração cujo corpo de frações seja o corpo $R(X)$ das funções racionais sobre $X$ e que domine um anel local de $Y$ domina também um anel local de $X$.
  \item[\rm{(ii)}] Reciprocamente, suponha-se que $f$ seja de tipo finito e que a propriedade enunciada em~(i) seja verificada por todo anel de valoração (resp. todo anel de valoração discreta) que tenha $R(X)$ como corpo de frações.
  Então $f$ é próprio.
\end{enumerate}
\end{corollary}

\begin{proof}
Observemos antes de tudo que as hipóteses implicam, em todos os casos, que $f$ é separado \sref[I]{I.5.5.9}.
\begin{enumerate}
  \item[\rm{(i)}] Sejam $K=R(Y)$, $L=R(X)$, $y$ um ponto de $Y$ e $A$ um anel de valoração que domine $\sh{O}_y$ e tenha $L$ como corpo de frações;
  a injeção $\sh{O}_y\to A$ define, portanto, um morfismo $h$ de $Y'=\Spec(A)$ em $Y$ \sref[I]{I.2.4.4} tal que $h(a)=y$, denotando por $a$ o ponto fechado de $Y'$;
  além disso, se $\eta$ é o ponto genérico de $Y$, que é também o de $\Spec(\sh{O}_y)$, tem-se $h(b)=\eta$, denotando por $b$ o ponto genérico de $Y'$ (pois $K\subset L$ por hipótese).
  Se $\xi$ é o ponto genérico de $X$, tem-se por hipótese $\kres(\xi)=\kres(b)=L$, donde resulta um $Y$-morfismo $g:\Spec(L)\to X$ tal que $g(b)=\xi$;
  em virtude de \sref{II.7.3.8}, $g$ provém de um $Y$-morfismo $g':Y'\to X$.
  Se $x=g'(a)$, é claro que $A$ domina $\sh{O}_x$.

  \item[\rm{(ii)}] Como a questão é local sobre $Y$, pode-se supor que $Y$ seja afim (resp. afim e noetheriano).
  Como $f$ é de tipo finito, pode-se aplicar, em ambos os casos, o lema de Chow \sref{II.5.6.1}.
  Existem, portanto, um morfismo projetivo $p:P\to Y$, uma imersão $j:X'\to P$ e um morfismo projetivo, sobrejetivo e birracional $g:X'\to X$, sendo $X'$ integral, tais que o diagrama
  \[
    \xymatrix{
      P \ar[d]_p
      & X' \ar[l]_j \ar[d]^g
      \\ Y
      & X \ar[l]^f
    }
  \]
  é comutativo.
  Basta provar que $j$ é uma imersão \emph{fechada}, pois então $f\circ g=p\circ j$ será um morfismo projetivo, e portanto próprio, e, como $g$ é sobrejetivo, $f$ será também próprio \sref{II.5.4.3}.
  Seja $Z$ o subpré-esquema fechado reduzido de $P$ cujo espaço subjacente é $\overline{j(X')}$ \sref[I]{I.5.2.1};
  como $X'$ é integral, $j$ se fatora como $i\circ h$, onde $i:Z\to P$ é a injeção canônica, $h:X'\to Z$ uma imersão aberta dominante \sref[I]{I.5.2.3}, e $Z$ é integral;
\oldpage[II]{147}
  além disso, $Z$ é projetivo sobre $Y$, e vê-se que podemos limitar-nos ao caso em que $P$ é \emph{integral} e $j$ \emph{dominante} e \emph{birracional}, reduzindo-se tudo a mostrar que $j$ é \emph{sobrejetivo}.
  Com efeito, seja $z\in P$;
  $\sh{O}_z$ é um anel local íntegro (resp. íntegro e noetheriano) cujo corpo de frações é
  \[
    L=R(P)=R(X')=R(X).
  \]
  Podemos limitar-nos ao caso em que $z$ não é o ponto genérico de $P$.
  Existe então, por \sref{II.7.1.2} e \sref{II.7.1.7}, um anel de valoração (resp. um anel de valoração discreta) $A$ que domina $\sh{O}_z$ e tem $L$ como corpo de frações.
  \emph{A fortiori}, $A$ domina $\sh{O}_y$, onde $y=p(z)$, e, por hipótese, existe então um $x\in X$ tal que $A$ domina $\sh{O}_x$.
  Como $g$ é próprio, a primeira parte da demonstração mostra que $A$ domina também $\sh{O}_{x'}$ para algum $x'\in X'$;
  segue-se que $\sh{O}_z$ e $\sh{O}_{j(x')}=\sh{O}_{x'}$ são aparentados \sref[I]{I.8.1.4} e, como $P$ é um esquema, isso implica $z=j(x')$ \sref[I]{I.8.2.2} e conclui a demonstração.
\end{enumerate}
\end{proof}

\begin{corollary}[7.3.11]
\label{II.7.3.11}
Sejam $X$ e $Y$ dois esquemas integrais, e seja $f:X\to Y$ um morfismo dominante, quase-compacto e universalmente fechado.
Suponha-se, além disso, que $Y$ seja afim, com anel (integral) $B$.
Então $\Gamma(X,\sh{O}_X)$ é canonicamente isomorfo a um subanel do fecho integral de $B$ em $R(X)$.
\end{corollary}

\begin{proof}
Com efeito, por \sref[I]{I.8.2.1.1}, $\Gamma(X,\sh{O}_X)$ identifica-se com a interseção dos $\sh{O}_x$ para $x\in X$;
em virtude de \sref{II.7.3.10}, \sref{II.7.1.2} e \sref{II.7.1.3}, $\Gamma(X,\sh{O}_X)$ fica, portanto, contido na interseção dos anéis de valoração que contêm $B$ e têm $R(X)$ como corpo de frações;
a conclusão segue então de \sref{II.7.1.3}.
\end{proof}

\begin{remarks}[7.3.12]
\label{II.7.3.12}
Sob as hipóteses de \sref{II.7.3.11}, e supondo-se que $R(X)$ seja uma extensão de tipo finito de $R(Y)$, pode-se, em muitos casos, concluir que $\Gamma(X,\sh{O}_X)$ é um \emph{módulo finitamente gerado} sobre o anel $B=\Gamma(Y,\sh{O}_Y)$.
Isto ocorrerá, por exemplo, quando $B$ for uma \emph{álgebra de tipo finito sobre um corpo}, pois se sabe então que o fecho integral de $B$ numa extensão de tipo finito do seu corpo de frações é um $B$-módulo finitamente gerado (\cite[t.~I, p.~267, th.~9]{I-13});
a conclusão resulta então de \sref{II.7.3.11} e do fato de $B$ ser noetheriano.

Em particular, \emph{um esquema $X$ próprio e afim sobre um corpo $K$ é finito}.
Com efeito, em virtude de \sref{II.1.6.4}, \sref{II.5.4.6} e \sref[I]{I.6.4.4}[(c)], pode-se reduzir ao caso em que $X$ é \emph{reduzido}.
Além disso, basta provar que cada um dos subpré-esquemas fechados de $X$ cujo espaço subjacente é um componente irredutível de $X$ (que são em número finito) é finito sobre $K$; assim, levando em conta \sref{II.5.4.5}, reduz-se finalmente ao caso em que $X$ é \emph{integral}.
Mas então o resultado decorre das observações precedentes.

No capítulo~III, voltaremos a demonstrar a proposição precedente por outros métodos e como consequência de resultados mais gerais, mostrando que, se $f:X\to Y$ é próprio e $Y$ localmente noetheriano, $f_*(\sh{F})$ é coerente para todo $\sh{O}_X$-módulo coerente $\sh{F}$ \sref[III]{III.4.4.2}.

Observemos, por fim, que o critério~\sref{II.7.3.10} é tomado como \emph{definição} dos morfismos próprios na geometria algébrica clássica.
Nós o mencionamos aqui apenas a título de observação, pois o critério~\sref{II.7.3.8} parece mais manejável em todas as aplicações que conhecemos.
\end{remarks}

\oldpage[II]{148}

\subsection{Curvas algébricas e corpos de funções de dimensão 1}
\label{subsection:II.7.4}

O objetivo deste número é mostrar como se formula, na linguagem dos esquemas, a noção clássica de curva algébrica (tal como é exposta, por exemplo, no livro de C.~Chevalley \cite{II-23}).
Em todo este número, \emph{denotamos por $k$ um corpo; todos os esquemas considerados são $k$-esquemas de tipo finito, e todos os morfismos são $k$-morfismos}.

\begin{proposition}[7.4.1]
\label{II.7.4.1}
Seja $X$ um pré-esquema de tipo finito sobre $k$ (portanto, noetheriano);
sejam $x_i$ ($1\leq i\leq n$) os pontos genéricos das componentes irredutíveis $X_i$ de $X$, e seja $K_i=\kres(x_i)$ ($1\leq i\leq n$).
As condições seguintes são equivalentes:
\begin{enumerate}
    \item[\rm{(a)}] Cada um dos $K_i$ é uma extensão de $k$ cujo grau de transcendência é igual a $1$.
    \item[\rm{(b)}] Para todo ponto fechado $x$ de $X$, o anel local $\sh{O}_x$ tem dimensão $1$ \sref{II.7.1.5}.
    \item[\rm{(c)}] Os subconjuntos fechados irredutíveis de $X$ distintos dos $X_i$ são exatamente os pontos fechados de $X$.
\end{enumerate}
\end{proposition}

\begin{proof}
Como $X$ é quase-compacto, todo subconjunto fechado irredutível $F$ de $X$ contém um ponto fechado \sref[0]{0.2.1.3}.
Em virtude de \sref[I]{I.2.4.2}, há uma correspondência biunívoca entre os ideais primos de $\sh{O}_x$ e os subconjuntos fechados irredutíveis de $X$ que contêm $x$ \sref[I]{I.1.1.14};
daí decorre imediatamente a equivalência entre (b) e (c).
Por outro lado, se $\mathfrak{p}_\alpha$ ($1\leq\alpha\leq r$) são os ideais primos minimais do anel local noetheriano $\sh{O}_x$, então os anéis locais $\sh{O}_x/\mathfrak{p}_\alpha$ são integrais e têm por corpos de frações os $K_i$ tais que $x\in X_i$.
Além disso, sabe-se (\cite[p.~4-06, th.~2]{I-1}) que a dimensão de uma $k$-álgebra local integral de tipo finito é igual ao grau de transcendência sobre $k$ de seu corpo de frações.
Por fim, a dimensão de $\sh{O}_x$ é o supremo das dimensões dos anéis $\sh{O}_x/\mathfrak{p}_\alpha$;
ora, a condição~(a) implica que essas dimensões são iguais a $1$, logo (a) implica (b);
reciprocamente, se $\sh{O}_x$ tem dimensão $1$, nenhum dos $\mathfrak{p}_\alpha$ pode ser igual ao ideal maximal de $\sh{O}_x$, pois, caso contrário, $\sh{O}_x$ teria dimensão $0$;
portanto, cada um dos anéis $\sh{O}_x/\mathfrak{p}_\alpha$ tem dimensão $1$, o que mostra que (b) implica (a).
\end{proof}

Observemos que, sob as condições de \sref{II.7.4.1}, o conjunto $X$ é \emph{vazio} ou \emph{infinito}, como decorre imediatamente de \sref[I]{I.6.4.4}.

\begin{definition}[7.4.2]
\label{II.7.4.2}
Chama-se \emph{curva algébrica sobre $k$} um \emph{esquema} algébrico não vazio sobre $k$ que satisfaz as condições de \sref{II.7.4.1}.
\end{definition}

Na linguagem da dimensão, que será introduzida no capítulo~IV, isso se expressará dizendo que uma curva algébrica sobre $k$ é um $k$-esquema algébrico não vazio \emph{cujas componentes irredutíveis são todas de dimensão $1$}.

Observemos que, se $X$ é uma curva algébrica sobre $k$, então os subpré-esquemas fechados reduzidos $X_i$ ($1\leq i\leq n$) de $X$ cujos espaços subjacentes são as componentes irredutíveis de $X$ também são curvas algébricas sobre $k$.

\begin{corollary}[7.4.3]
\label{II.7.4.3}
Seja $X$ uma curva algébrica irredutível.
O único ponto não fechado de $X$ é seu ponto genérico.
Os subconjuntos fechados de $X$ distintos de $X$ são os conjuntos finitos de pontos fechados;
estes são também os únicos subconjuntos de $X$ que não são densos em toda parte.
\end{corollary}

\begin{proof}
Se um ponto $x\in X$ não é fechado, seu fecho em $X$ é um subconjunto fechado irredutível de $X$ e, portanto, é necessariamente todo $X$, em virtude de \sref{II.7.4.1}; assim, $x$ é o ponto genérico de $X$.
Um subconjunto fechado $F$ de $X$ distinto de $X$ não pode conter
\oldpage[II]{149}
o ponto genérico de $X$ e, portanto, todos os seus pontos são fechados (em $X$ e, \emph{a fortiori}, em $F$);
considerando o subpré-esquema fechado reduzido de $X$ cujo espaço subjacente é $F$ \sref[I]{I.5.2.1}, resulta então de \sref[I]{I.6.2.2} que $F$ é finito e discreto.
O fecho em $X$ de todo subconjunto infinito de $X$ é, portanto, necessariamente igual a $X$.
\end{proof}

Quando $X$ é uma curva algébrica qualquer, aplicando-se \sref{II.7.4.3} às componentes irredutíveis de $X$, vê-se que os únicos pontos não fechados de $X$ são os pontos genéricos dessas componentes.

\begin{corollary}[7.4.4]
\label{II.7.4.4}
Sejam $X$ e $Y$ duas curvas algébricas irredutíveis sobre $k$, e seja $f:X\to Y$ um $k$-morfismo.
Para que $f$ seja dominante, é necessário e suficiente que $f^{-1}(y)$ seja finito para todo $y\in Y$.
\end{corollary}

\begin{proof}
Com efeito, se $f$ não é dominante, então $f(X)$ é necessariamente um subconjunto \emph{finito} de $Y$, em virtude de \sref{II.7.4.3}; logo, não é possível que $f^{-1}(y)$ seja finito para todo ponto de $Y$, pois, caso contrário, $X$ seria finito, o que é absurdo \sref{II.7.4.1}.
Reciprocamente, se $f$ é dominante, então, para todo $y\in Y$ distinto do ponto genérico $\eta$ de $Y$, $f^{-1}(y)$ é fechado em $X$, pois $\{y\}$ é fechado em $Y$ \sref{II.7.4.3};
além disso, por hipótese, $f^{-1}(y)$ não contém o ponto genérico $\xi$ de $X$ e, portanto, é finito, em virtude de \sref{II.7.4.3}.
Finalmente, para ver que, quando $f$ é dominante, $f^{-1}(\eta)$ é finito, observamos que a fibra $f^{-1}(\eta)$ é um esquema irredutível de tipo finito sobre $\kres(\eta)$ e tem $\xi$ como ponto genérico (\sref[I]{I.6.3.9} e \sref[I]{I.6.4.11}).
Como $\kres(\xi)$ e $\kres(\eta)$ são extensões de tipo finito de $k$, ambas de grau de transcendência $1$, $\kres(\xi)$ é necessariamente uma extensão finita de $\kres(\eta)$; logo, $\xi$ é fechado em $f^{-1}(\eta)$ \sref[I]{I.6.4.2}, e $f^{-1}(\eta)$ consiste, portanto, no único ponto $\xi$.
\end{proof}

Veremos no Capítulo~III que um morfismo \emph{próprio} $f:X\to Y$ de pré-esquemas noetherianos, tal que $f^{-1}(y)$ seja finito para todo $y\in Y$, é necessariamente \emph{finito};
decorrerá, portanto, de \sref{II.7.4.4} que um morfismo dominante próprio de uma curva algébrica irredutível em uma curva algébrica é \emph{finito}.

\begin{corollary}[7.4.5]
\label{II.7.4.5}
Seja $X$ uma curva algébrica sobre $k$.
Para que $X$ seja regular, é necessário e suficiente que $X$ seja normal ou, equivalentemente, que os anéis locais de seus pontos fechados sejam anéis de valoração discreta.
\end{corollary}

\begin{proof}
Decorre imediatamente da condição~(b) de \sref{II.7.4.1} e de \sref{II.7.1.6}.
\end{proof}

\begin{corollary}[7.4.6]
\label{II.7.4.6}
Sejam $X$ uma curva algébrica reduzida e $\sh{A}$ uma $\sh{R}(X)$-álgebra coerente reduzida;
então o fecho integral $X'$ de $X$ relativamente a $\sh{A}$ \sref{II.6.3.4} é uma curva algébrica normal, e o morfismo canônico $X'\to X$ é finito.
\end{corollary}

\begin{proof}
O fato de $X'\to X$ ser finito decorre de \sref{II.6.3.10};
$X'$ é, portanto, um $k$-esquema algébrico;
além disso, se $x_i$ ($1\leq i\leq n$) são os pontos genéricos das componentes irredutíveis de $X$, e $x'_j$ ($1\leq j\leq m$) os pontos genéricos das componentes irredutíveis de $X'$, cada $\kres(x'_j)$ é uma extensão algébrica finita de um dos $\kres(x_i)$ \sref{II.6.3.6} e, portanto, tem grau de transcendência $1$ sobre $k$.
Assim, $X'$ é de fato uma curva algébrica sobre $k$ e, além disso, sabemos que $X'$ é soma de um número finito de esquemas integrais e normais (\sref{II.6.3.6} e \sref{II.6.3.7}).
\end{proof}

\begin{env}[7.4.7]
\label{II.7.4.7}
Diz-se que uma curva algébrica $X$ sobre $k$ é \emph{completa} se é \emph{própria} sobre $k$.
\end{env}

\begin{corollary}[7.4.8]
\label{II.7.4.8}
Para que uma curva algébrica reduzida $X$ sobre $k$ seja completa, é necessário e suficiente que sua normalização $X'$ seja completa.
\end{corollary}

\oldpage[II]{150}
\begin{proof}
Com efeito, o morfismo canônico $f:X'\to X$ é finito \sref{II.7.4.6} e, portanto, próprio \sref{II.6.1.11} e sobrejetivo \sref{II.6.3.8};
se $g:X\to\Spec(k)$ é o morfismo estrutural, então $g$ e $g\circ f$ são próprios simultaneamente, como resulta de \sref{II.5.4.2}[(ii)] e \sref{II.5.4.3}[(ii)], pois $g$ é separado por hipótese.
\end{proof}

\begin{proposition}[7.4.9]
\label{II.7.4.9}
Sejam $X$ uma curva algébrica \emph{normal} sobre $k$ e $Y$ um $k$-esquema algébrico \emph{próprio}.
Então toda $k$-aplicação racional de $X$ em $Y$ está definida em toda parte ou, em outras palavras, é um morfismo.
\end{proposition}

\begin{proof}
Com efeito, decorre de \sref{II.7.3.7} que, nos pontos $x\in X$ onde tal aplicação não está definida, a dimensão de $\sh{O}_x$ deveria ser $\geq2$; logo, o conjunto desses pontos é vazio;
a última afirmação decorre de \sref[I]{I.7.2.3}.
\end{proof}

\begin{corollary}[7.4.10]
\label{II.7.4.10}
Uma curva algébrica normal sobre $k$ é quase-projetiva sobre $k$.
\end{corollary}

\begin{proof}
Como $X$ é soma de um número finito de curvas algébricas integrais e normais \sref{II.6.3.8}, podemos nos restringir ao caso em que $X$ é integral \sref{II.5.3.6}.
Como $X$ é quase-compacto, ele é recoberto por um número finito de abertos afins $U_i$ ($1\leq i\leq n$) e, como cada um deles é de tipo finito sobre $k$, para cada $i$ existem um inteiro $n_i$ e uma $k$-imersão $f_i:U_i\to\bb{P}_k^{n_i}$ \sref{II.5.3.3} e \sref{II.5.3.4}[(i)].
Como $U_i$ é denso em $X$, decorre de \sref{II.7.4.9} que $f_i$ se prolonga a um $k$-morfismo $g_i:X\to\bb{P}_k^{n_i}$, de onde obtemos um $k$-morfismo $g=(g_1,\ldots,g_n)_k$ de $X$ no produto $P$ dos $\bb{P}_k^{n_i}$ sobre $k$.
Além disso, para cada $i$, como a restrição de $g_i$ a $U_i$ é uma imersão, o mesmo vale para a restrição de $g$ a $U_i$ \sref[I]{I.5.3.14}.
Como os $U_i$ recobrem $X$ e $g$ é separado \sref[I]{I.5.5.1}[(v)], $g$ é uma imersão de $X$ em $P$ \sref[I]{I.8.2.8}.
Como o morfismo de Segre \sref{II.4.3.3} fornece uma imersão de $P$ em $\bb{P}_k^N$, isso conclui a prova de que $X$ é quase-projetiva.
\end{proof}

\begin{corollary}[7.4.11]
\label{II.7.4.11}
Toda curva algébrica normal $X$ é isomorfa ao esquema induzido sobre um aberto denso em toda parte de uma curva algébrica normal e completa $\widehat{X}$, determinada a menos de isomorfismo único.
\end{corollary}

\begin{proof}
Se $X_1$ e $X_2$ são duas curvas normais e completas, decorre imediatamente de \sref{II.7.4.9} que todo isomorfismo de um aberto $U_1$ denso em $X_1$ sobre um aberto $U_2$ denso em $X_2$ se prolonga de maneira única a um isomorfismo de $X_1$ sobre $X_2$;
daí resulta a afirmação de unicidade.
Para provar a existência de $\widehat{X}$, basta observar que podemos considerar $X$ como um subesquema de $\bb{P}_k^n$ \sref{II.7.4.10}.
Seja $\overline{X}$ o fecho de $X$ em $\bb{P}_k^n$ \sref[I]{I.9.5.11};
como $X$ é induzido por $\overline{X}$ sobre um aberto denso em $\overline{X}$ \sref[I]{I.9.5.10}, os pontos genéricos $x_i$ das componentes irredutíveis de $X$ são também os das componentes irredutíveis de $\overline{X}$, e os $\kres(x_i)$ são os mesmos para ambos os esquemas; portanto, por \sref{II.7.4.1}, $\overline{X}$ é uma curva algébrica sobre $k$, reduzida \sref[I]{I.9.5.9} e projetiva sobre $k$ \sref{II.5.5.1}, e assim completa \sref{II.5.5.3}.
Tomemos então como $\widehat{X}$ a \emph{normalização} de $\overline{X}$, que continua completa \sref{II.7.4.8};
além disso, se $h:\widehat{X}\to\overline{X}$ é o morfismo canônico, a restrição de $h$ a $h^{-1}(X)$ é um isomorfismo sobre $X$, pois $X$ é normal \sref{II.6.3.4}, e, como $h^{-1}(X)$ contém os pontos genéricos das componentes irredutíveis de $\widehat{X}$ \sref{II.6.3.8}, é denso em $\widehat{X}$, o que conclui a demonstração.
\end{proof}

\oldpage[II]{151}
\begin{remark}[7.4.12]
\label{II.7.4.12}
Mostraremos no Capítulo~V que a conclusão de \sref{II.7.4.10} vale mesmo sem supor que a curva seja normal (nem sequer reduzida);
mostraremos também que, para que uma curva algébrica (reduzida ou não) seja afim, é necessário e suficiente que suas componentes irredutíveis (reduzidas) não sejam completas.
\end{remark}

\begin{corollary}[7.4.13]
\label{II.7.4.13}
Sejam $X$ uma curva normal irredutível sobre $k$, com corpo de funções $K=R(X)$, e $Y$ uma curva integral completa sobre $k$, com corpo de funções $L=R(Y)$.
Existe uma correspondência biunívoca canônica entre os $k$-morfismos \emph{dominantes} $X\to Y$ e os $k$-monomorfismos $L\to K$.
\end{corollary}

\begin{proof}
Em virtude de \sref{II.7.4.9}, as $k$-aplicações racionais de $X$ em $Y$ se identificam com os $k$-morfismos $u:X\to Y$.
Como os morfismos dominantes $u:X\to Y$ se caracterizam pela condição $u(x)=y$ (denotando por $x$ e $y$ os pontos genéricos respectivos de $X$ e $Y$), o corolário resulta dessas observações e de \sref[I]{I.7.1.13}.
\end{proof}

\begin{env}[7.4.14]
\label{II.7.4.14}
Pode-se precisar o resultado de \sref{II.7.4.13} quando se toma como $Y$ a \emph{reta projetiva} $\bb{P}_k^1=\Proj(k[T_0,T])$, onde $T_0$ e $T$ são duas indeterminadas.
Então $Y$ é um esquema integral \sref{II.2.4.4}, e o esquema induzido sobre o aberto $D_+(T_0)$ de $Y$ é isomorfo a $\Spec(k[T])$ \sref{II.2.3.6}; portanto, o ponto genérico de $Y$ é o ideal $(0)$ de $k[T]$, e seu corpo de funções racionais é $k(T)$, o que mostra que $Y$ é uma curva algébrica completa sobre $k$.
Além disso, o único ideal primo homogêneo de $S=k[T_0,T]$ que contém $T_0$ e é distinto de $S_+$ é o ideal principal $(T_0)$; por conseguinte, o complemento de $D_+(T_0)$ em $Y=\bb{P}_k^1$ reduz-se a \emph{um ponto fechado}, chamado «ponto no infinito», que denotaremos por $\infty$ (para um estudo geral das relações entre fibrados vetoriais e fibrados projetivos, veja \sref{subsection:II.8.4}).
Com estas notações:
\end{env}

\begin{corollary}[7.4.15]
\label{II.7.4.15}
Seja $X$ uma curva normal irredutível sobre $k$, com corpo de funções $K=R(X)$.
Existe uma correspondência biunívoca canônica entre o conjunto $K$ e o conjunto dos morfismos $u$ de $X$ em $\bb{P}_k^1$ distintos do morfismo constante de valor $\infty$.
Para que tal morfismo seja dominante, é necessário e suficiente que o elemento correspondente de $K$ seja transcendente sobre $k$.
\end{corollary}

\begin{proof}
O enunciado resulta imediatamente de \sref{II.7.4.9} e do lema seguinte:
\begin{lemma}[7.4.15.1]
\label{II.7.4.15.1}
Seja $X$ um pré-esquema integral sobre $k$, e seja $K=R(X)$ seu corpo de funções racionais.
Existe uma correspondência biunívoca canônica entre o conjunto $K$ e o conjunto das aplicações racionais $u$ de $X$ em $\bb{P}_k^1$ distintas do morfismo constante de valor $\infty$.
Para que tal aplicação racional seja dominante, é necessário e suficiente que o elemento correspondente de $K$ seja transcendente sobre $k$.
\end{lemma}
Antes de tudo, as aplicações racionais de $X$ em $\bb{P}_k^1$ correspondem biunivocamente aos pontos de $\bb{P}_k^1$ a valores na extensão $K$ de $k$ \sref[I]{I.7.1.12}.
Se tal ponto está situado no ponto genérico de $\bb{P}_k^1$ \sref[I]{I.3.4.5}, a aplicação racional correspondente é evidentemente dominante.
No caso contrário, como todo ponto de $\bb{P}_k^1$ distinto do ponto genérico é fechado \sref{II.7.4.3}, a imagem do domínio de definição $U$ de $u$ pelo único morfismo $U\to\bb{P}_k^1$ da classe $u$ \sref[I]{I.7.2.2} reduz-se a um ponto fechado $y$ de $\bb{P}_k^1$, e esse morfismo (que não está necessariamente definido em todo $X$) não é, portanto, dominante;
por abuso de linguagem, dizemos então que a aplicação racional $u$ é «constante, de valor $y$».
Resta pôr em correspondência biunívoca os pontos de $\bb{P}_k^1$ a valores em $K$ que não estão situados em $\infty$ \sref[I]{I.3.4.5} e os elementos de $K$, e verificar que tal ponto está situado no ponto genérico de $\bb{P}_k^1$ se, e somente se, corresponde a um
\oldpage[II]{152}
elemento transcendente sobre $k$.
Mas isto é imediato \sref{II.4.2.6}[exemplo~1].
\end{proof}

\begin{corollary}[7.4.16]
\label{II.7.4.16}
Sejam $X$ e $Y$ duas curvas algébricas sobre $k$, normais, completas e irredutíveis;
sejam $K=R(X)$ e $L=R(Y)$ seus corpos de funções.
Existe uma correspondência biunívoca canônica entre o conjunto dos $k$-isomorfismos $X\simto Y$ e o conjunto dos $k$-isomorfismos $L\simto K$.
\end{corollary}

\begin{proof}
É uma consequência evidente de \sref{II.7.4.13}.
\end{proof}

\begin{env}[7.4.17]
\label{II.7.4.17}
Este corolário \sref{II.7.4.16} mostra que uma curva algébrica sobre $k$, normal, completa e irredutível, é \emph{determinada por seu corpo $K$ de funções racionais a menos de isomorfismo único};
por definição, $K$ é uma extensão de tipo finito de $k$, de grau de transcendência $1$ (o que classicamente se denomina um \emph{corpo de funções algébricas de uma variável}).
Além disso:
\end{env}

\begin{proposition}[7.4.18]
\label{II.7.4.18}
Para toda extensão $K$ de $k$, de tipo finito e de grau de transcendência $1$, existe uma curva algébrica $X$ (determinada a menos de isomorfismo único), normal, completa e irredutível, tal que $R(X)=K$.
O conjunto dos anéis locais de $X$ identifica-se \sref[I]{I.8.2.1} com o conjunto formado pelo corpo $K$ e pelos anéis de valoração que contêm $k$ e têm $K$ como corpo de frações.
\end{proposition}

\begin{proof}
Com efeito, $K$ é uma extensão de grau finito de uma extensão puramente transcendente $k(T)$ de $k$, que se identifica, como se viu, com o corpo das funções racionais da reta projetiva $Y=\bb{P}_k^1$.
Seja $X$ o \emph{fecho integral} de $Y$ relativamente a $K$ \sref{II.6.3.4};
então $X$ é uma curva algébrica normal sobre $k$, com corpo de funções $K$ \sref{II.6.3.7}, e é completa, pois o morfismo $X\to Y$ é finito \sref{II.7.4.6}.
Os anéis locais $\sh{O}_x$ de $X$ são o corpo $K$, quando $x$ é o ponto genérico, ou anéis de valoração discreta que contêm $k$ e têm $K$ como corpo de frações, quando $x$ é distinto do ponto genérico \sref{II.7.4.5}.
Reciprocamente, seja $A$ um tal anel;
como o morfismo $X\to\Spec(k)$ é próprio e $A$ domina $k$, $A$ também domina um anel local $\sh{O}_x$ de $X$ \sref{II.7.3.10};
como este último é um anel de valoração que tem $K$ como corpo de frações, ele é necessariamente igual a $A$.
\end{proof}

\begin{remarks}[7.4.19]
\label{II.7.4.19}
De \sref{II.7.4.16} e \sref{II.7.4.18} resulta que \emph{o dado de uma curva algébrica sobre $k$, normal, completa e irredutível, é essencialmente equivalente ao dado de uma extensão $K$ de $k$, de tipo finito e de grau de transcendência $1$}.
Observe-se que, se $k'$ é uma extensão do corpo de base $k$, $X\otimes_k k'$ continuará sendo uma curva algébrica completa sobre $k'$ \sref{II.5.4.2}[(iii)], mas, em geral, não será nem reduzida nem irredutível.
Será, contudo, reduzida e irredutível se $K$ for uma extensão separável de $k$ e $k$ for algebricamente fechado em $K$ (o que, numa terminologia clássica que não seguiremos, se exprime dizendo que $K$ é uma «extensão regular» de $k$).
Mas, mesmo nesse caso, pode ocorrer que $X\otimes_k k'$ não seja normal.
O leitor encontrará detalhes sobre essas questões no Capítulo~IV.
\end{remarks}

\section{Esquemas de explosão; cones projetantes; fecho projetivo}
\label{section:II.8}

\subsection{Pré-esquemas de explosão}
\label{subsection:II.8.1}
\label{II.8.1}

\begin{env}[8.1.1]
\label{II.8.1.1}
Seja $Y$ um pré-esquema e, para cada inteiro $n\geq0$, seja $\sh{I}_n$ um feixe quase-coerente de ideais de $\sh{O}_Y$; suponhamos verificadas as condições seguintes:
\[
\label{II.8.1.1.1}
  \sh{I}_0=\sh{O}_Y,\qquad \sh{I}_n\subset\sh{I}_m\quad\text{para }m\leq n,
\tag{8.1.1.1}
\]
\[
\label{II.8.1.1.2}
  \sh{I}_m\sh{I}_n\subset\sh{I}_{m+n}\quad\text{quaisquer que sejam }m,n.
\tag{8.1.1.2}
\]

\oldpage[II]{153}
Observemos que essas hipóteses implicam
\[
\label{II.8.1.1.3}
  \sh{I}_1^n\subset\sh{I}_n.
\tag{8.1.1.3}
\]

Ponhamos
\[
\label{II.8.1.1.4}
  \sh{S}=\bigoplus_{n\geq0}\sh{I}_n.
\tag{8.1.1.4}
\]

Resulta de \sref{II.8.1.1.1} e \sref{II.8.1.1.2} que $\sh{S}$ é uma $\sh{O}_Y$-álgebra graduada quase-coerente e define, portanto, um $Y$-esquema $X=\Proj(\sh{S})$.
Se $\sh{J}$ é um feixe invertível de ideais de $\sh{O}_Y$, então $\sh{I}_n\otimes_{\sh{O}_Y}\sh{J}^{\otimes n}$ identifica-se canonicamente com $\sh{I}_n\sh{J}^n$.
Se substituímos então os $\sh{I}_n$ pelos $\sh{I}_n\sh{J}^n$ e, com isso, $\sh{S}$ pela $\sh{O}_Y$-álgebra quase-coerente $\sh{S}_{(\sh{J})}$, então $X_{(\sh{J})}=\Proj(\sh{S}_{(\sh{J})})$ é canonicamente isomorfo a $X$ \sref{II.3.1.8}.
\end{env}

\begin{env}[8.1.2]
\label{II.8.1.2}
Suponhamos que $Y$ seja \emph{localmente integral}, de modo que o feixe $\sh{R}(Y)$ das funções racionais é uma $\sh{O}_Y$-álgebra quase-coerente \sref[I]{I.7.3.7}.
Dizemos que um $\sh{O}_Y$-submódulo $\sh{I}$ de $\sh{R}(Y)$ é um \emph{ideal fracionário} de $\sh{R}(Y)$ se for \emph{finitamente gerado como $\sh{O}_Y$-módulo} \sref[0]{0.5.2.1}.
Suponhamos dado, para todo $n\geq0$, um ideal fracionário quase-coerente $\sh{I}_n$ de $\sh{R}(Y)$, tal que $\sh{I}_0=\sh{O}_Y$ e que se verifique a condição \sref{II.8.1.1.2} (mas não necessariamente a segunda condição de \sref{II.8.1.1.1});
podemos então definir novamente, pela fórmula \sref{II.8.1.1.4}, uma $\sh{O}_Y$-álgebra graduada quase-coerente e o $Y$-esquema correspondente $X=\Proj(\sh{S})$;
teremos novamente um isomorfismo canônico $X\to X_{(\sh{J})}$ para todo ideal fracionário \emph{invertível} $\sh{J}$ de $\sh{R}(Y)$.
\end{env}

\begin{definition}[8.1.3]
\label{II.8.1.3}
Seja $Y$ um pré-esquema (resp. um pré-esquema localmente integral), e seja $\sh{I}$ um ideal quase-coerente de $\sh{O}_Y$ (resp. um ideal fracionário quase-coerente de $\sh{R}(Y)$).
Dizemos que o $Y$-esquema
\[
  X=\Proj\left(\bigoplus_{n\geq0}\sh{I}^n\right)
\]
é obtido pela explosão do ideal $\sh{I}$, ou que é o pré-esquema de explosão de $Y$ relativamente a $\sh{I}$.
Quando $\sh{I}$ é um ideal quase-coerente de $\sh{O}_Y$ e $Y'$ é o subpré-esquema fechado de $Y$ definido por $\sh{I}$, dizemos também que $X$ é o $Y$-esquema obtido pela explosão de $Y'$.
\end{definition}

Por definição, $\sh{S}=\bigoplus_{n\geq0}\sh{I}^n$ é então gerada por $\sh{S}_1=\sh{I}$;
se $\sh{I}$ é um $\sh{O}_Y$-módulo finitamente gerado, então $X$ é, portanto, \emph{projetivo} sobre $Y$ \sref{II.5.5.2}.
Sem hipótese alguma sobre $\sh{I}$, o $\sh{O}_X$-módulo $\sh{O}_X(1)$ é \emph{invertível} \sref{II.3.2.5} e \emph{muito amplo}, em virtude de \sref{II.4.4.3} aplicado ao morfismo estrutural $X\to Y$.

Observemos que, se $f:X\to Y$ é o morfismo estrutural, a restrição de $f$ a $f^{-1}(Y\setminus Y')$ é um \emph{isomorfismo} sobre $Y\setminus Y'$ quando $\sh{I}$ é um \emph{ideal de $\sh{O}_Y$} e $Y'$ é o subpré-esquema fechado que ele define:
com efeito, como a questão é local em $Y$, basta supor que $\sh{I}=\sh{O}_Y$, e nossa afirmação resulta de \sref{II.3.1.7}.

Quando substituímos $\sh{I}$ por $\sh{I}^d$ ($d>0$), o $Y$-esquema de explosão $X$ é substituído por um $Y$-esquema $X'$ canonicamente isomorfo \sref{II.8.1.1};
do mesmo modo, para todo ideal (resp. ideal fracionário) \emph{invertível} $\sh{J}$, o pré-esquema de explosão $X_{(\sh{J})}$ relativamente ao ideal $\sh{I}\sh{J}$ é canonicamente isomorfo a $X$ \sref{II.8.1.1}.

Em particular, quando $\sh{I}$ é um ideal (resp. um ideal fracionário) \emph{invertível}, o $Y$-esquema obtido pela explosão de $\sh{I}$ é \emph{isomorfo a $Y$} \sref{II.3.1.7}.

\begin{proposition}[8.1.4]
\label{II.8.1.4}
Seja $Y$ um pré-esquema integral.
\begin{enumerate}
    \item[\rm{(i)}] Para toda sequência $(\sh{I}_n)$ de ideais fracionários quase-coerentes de $\sh{R}(Y)$ que satisfaça \sref{II.8.1.1.2}
\oldpage[II]{154}
        e tal que $\sh{I}_0=\sh{O}_Y$, o $Y$-esquema $X=\Proj(\bigoplus_{n\geq0}\sh{I}_n)$ é integral, e o morfismo estrutural $f:X\to Y$ é dominante.
    \item[\rm{(ii)}] Seja $\sh{I}$ um ideal fracionário quase-coerente de $\sh{R}(Y)$, e seja $X$ o $Y$-esquema obtido pela explosão de $Y$ relativamente a $\sh{I}$.
        Se $\sh{I}\neq0$, então o morfismo estrutural $f:X\to Y$ é birracional e sobrejetivo.
\end{enumerate}
\end{proposition}

\begin{proof}
\medskip\noindent
\begin{enumerate}
    \item[\rm{(i)}] Isto resulta do fato de que $\sh{S}=\bigoplus_{n\geq0}\sh{I}_n$ é uma $\sh{O}_Y$-álgebra \emph{integral} (\sref{II.3.1.12} e \sref{II.3.1.14}), pois, para todo $y\in Y$, $\sh{O}_y$ é um anel integral \sref[I]{I.5.1.4}.
    \item[\rm{(ii)}] Por (i), $X$ é integral; se, além disso, $x$ e $y$ são os pontos genéricos de $X$ e $Y$, respectivamente, tem-se $f(x)=y$, e resta demonstrar que $\kres(x)$ tem dimensão 1 sobre $\kres(y)$.
        Mas $x$ também é o ponto genérico da fibra $f^{-1}(y)$; se $\psi$ é o morfismo canônico $Z\to Y$, onde $Z=\Spec(\kres(y))$, então o pré-esquema $f^{-1}(y)$ pode ser identificado com $\Proj(\sh{S}')$, onde $\sh{S}'=\psi^*(\sh{S})$ \sref{II.3.5.3}.
        Mas é claro que $\sh{S}'=\bigoplus_{n\geq0}(\sh{I}_y)^n$, e, como $\sh{I}$ é um ideal fracionário quase-coerente de $\sh{R}(Y)$ que não é nulo, $\sh{I}_y\neq0$ \sref[I]{I.7.3.6}, de onde $\sh{I}_y=\kres(y)$;
        então $\Proj(\sh{S}')$ pode ser identificado com $\Spec(\kres(y))$ \sref{II.3.1.7}, de onde segue a conclusão.
\end{enumerate}
\end{proof}

Demonstraremos uma recíproca de \sref{II.8.1.4} em \sref[III]{III.2.3.8}.

\begin{env}[8.1.5]
\label{II.8.1.5}
Voltemos à situação e às notações de \sref{II.8.1.1}.
Por definição, os homomorfismos de inclusão $\sh{I}_{n+1}\to\sh{I}_n$ \sref{II.8.1.1.1} definem, para cada $k\in\bb{Z}$, um homomorfismo injetivo de grau zero de $\sh{S}$-módulos graduados
\[
\label{II.8.1.5.1}
  u_k:\sh{S}_+(k+1)\to\sh{S}(k);
\tag{8.1.5.1}
\]
como $\sh{S}_+(k+1)$ e $\sh{S}(k+1)$ são canonicamente \textbf{(TN)}-isomorfos, a $u_k$ corresponde canonicamente um homomorfismo injetivo de $\sh{O}_X$-módulos \sref{II.3.4.2}:
\[
\label{II.8.1.5.2}
  \widetilde{u}_k:\sh{O}_X(k+1)\to\sh{O}_X(k).
\tag{8.1.5.2}
\]

Recordemos, por outro lado, \sref{II.3.2.6} que foram definidos homomorfismos canônicos
\[
\label{II.8.1.5.3}
  \lambda:\sh{O}_X(h)\otimes_{\sh{O}_X}\sh{O}_X(k)\to\sh{O}_X(h+k)
\tag{8.1.5.3}
\]
e, como o diagrama
\[
  \xymatrix{
    \sh{S}(h)\otimes_{\sh{S}}\sh{S}(k)\otimes_{\sh{S}}\sh{S}(l)
      \ar[r]
      \ar[d]
  & \sh{S}(h+k)\otimes_{\sh{S}}\sh{S}(l)
      \ar[d]
  \\\sh{S}(h)\otimes_{\sh{S}}\sh{S}(k+l)
      \ar[r]
  & \sh{S}(h+k+l)
  }
\]
é comutativo, resulta da funtorialidade dos $\lambda$ \sref{II.3.2.6} que os homomorfismos \sref{II.8.1.5.3} definem sobre
\[
\label{II.8.1.5.4}
  \sh{S}_X=\bigoplus_{n\in\bb{Z}}\sh{O}_X(n)
\tag{8.1.5.4}
\]
\oldpage[II]{155}
uma estrutura de $\sh{O}_X$-álgebra graduada quase-coerente.

Além disso, o diagrama
\[
  \xymatrix{
    \sh{S}(h)\otimes_{\sh{S}}\sh{S}_+(k+1)
      \ar[r]
      \ar[d]_{1\otimes u_k}
  & \sh{S}_+(h+k+1)
      \ar[d]^{u_{k+h}}
  \\\sh{S}(h)\otimes_{\sh{S}}\sh{S}(k)
      \ar[r]
  & \sh{S}(h+k)
  }
\]
é comutativo; a funtorialidade dos $\lambda$ mostra então que temos um diagrama comutativo
\[
\label{II.8.1.5.5}
  \xymatrix{
    \sh{O}_X(h)\otimes_{\sh{O}_X}\sh{O}_X(k+1)
      \ar[r]^{\lambda}
      \ar[d]_{1\otimes\widetilde{u}_k}
  & \sh{O}_X(h+k+1)
      \ar[d]^{\widetilde{u}_{k+h}}
  \\\sh{O}_X(h)\otimes_{\sh{O}_X}\sh{O}_X(k)
      \ar[r]^{\lambda}
  & \sh{O}_X(h+k)
  }
\tag{8.1.5.5}
\]
sendo as setas horizontais os homomorfismos canônicos.
Podemos, portanto, dizer que os $\widetilde{u}_k$ definem um \emph{homomorfismo injetivo} (de grau zero) \emph{de $\sh{S}_X$-módulos graduados}
\[
\label{II.8.1.5.6}
  \widetilde{u}:\sh{S}_X(1)\to\sh{S}_X.
\tag{8.1.5.6}
\]
\end{env}

\begin{env}[8.1.6]
\label{II.8.1.6}
Mantendo as notações de \sref{II.8.1.5}, observemos agora que, para $n\geq0$, o homomorfismo composto $\widetilde{v}_n=\widetilde{u}_{n-1}\circ\widetilde{u}_{n-2}\circ\ldots\circ\widetilde{u}_0$ é um homomorfismo \emph{injetivo} $\sh{O}_X(n)\to\sh{O}_X$;
denotaremos por $\sh{I}_{n,X}$ sua imagem, que é, portanto, um ideal quase-coerente de $\sh{O}_X$, \emph{isomorfo} a $\sh{O}_X(n)$.
Além disso, o diagrama
\[
  \xymatrix{
    \sh{O}_X(m)\otimes_{\sh{O}_X}\sh{O}_X(n)
      \ar[r]^{\lambda}
      \ar[d]_{\widetilde{v}_m\otimes\widetilde{v}_n}
  & \sh{O}_X(m+n)
      \ar[d]^{\widetilde{v}_{m+n}}
  \\\sh{O}_X
      \ar[r]_{\id}
  & \sh{O}_X
  }
\]
é comutativo para $m\geq0$, $n\geq0$.
Deduzimos assim as inclusões seguintes:
\[
\label{II.8.1.6.1}
  \sh{I}_{0,X}=\sh{O}_X,\quad \sh{I}_{n,X}\subset\sh{I}_{m,X}
  \qquad\mbox{para $0\leq m\leq n$;}
\tag{8.1.6.1}
\]
\[
\label{II.8.1.6.2}
  \sh{I}_{m,X}\sh{I}_{n,X}\subset\sh{I}_{m+n,X}
  \qquad\qquad\mbox{para $m\geq0$, $n\geq0$.}
\tag{8.1.6.2}
\]
\end{env}

\oldpage[II]{156}

\begin{proposition}[8.1.7]
\label{II.8.1.7}
Seja $Y$ um pré-esquema, $\sh{I}$ um ideal quase-coerente de $\sh{O}_Y$, e $X=\Proj\left(\bigoplus_{n\geq0}\sh{I}^n\right)$ o $Y$-esquema obtido pela explosão de $\sh{I}$.
Temos então, para todo $n>0$, um isomorfismo canônico
\[
\label{II.8.1.7.1}
  \sh{O}_X(n)\simto\sh{I}^n\sh{O}_X=\sh{I}_{n,X}
\tag{8.1.7.1}
\]
(cf. \sref[0]{0.4.3.5}) e, portanto, $\sh{I}^n\sh{O}_X$ é um $\sh{O}_X$-módulo invertível muito amplo se $n>0$.
\end{proposition}

\begin{proof}
A última afirmação é imediata, pois $\sh{O}_X(1)$ é invertível \sref{II.3.2.5} e muito amplo para $Y$ por definição (\sref{II.4.4.3} e \sref{II.4.4.9}).
Por outro lado, por definição, a imagem de $v_n$ é precisamente $\sh{I}^n\sh{S}$, e \sref{II.8.1.7.1} resulta então da exatidão do funtor $\sh{M}\mapsto\widetilde{\sh{M}}$ \sref{II.3.2.4} e da fórmula~\sref{II.3.2.4.1}.
\end{proof}

\begin{corollary}[8.1.8]
\label{II.8.1.8}
Sob as hipóteses de \sref{II.8.1.7}, se $f:X\to Y$ é o morfismo estrutural e $Y'$ o subpré-esquema fechado de $Y$ definido por $\sh{I}$, então o subpré-esquema fechado $X'=f^{-1}(Y')$ de $X$ é definido por $\sh{I}\sh{O}_X$ (canonicamente isomorfo a $\sh{O}_X(1)$), de onde se obtém uma sequência exata curta canônica
\[
\label{II.8.1.8.1}
  0\to\sh{O}_X(1)\to\sh{O}_X\to\sh{O}_{X'}\to0.
\tag{8.1.8.1}
\]
\end{corollary}

\begin{proof}
Isto resulta de \sref{II.8.1.7.1} e de \sref[I]{I.4.4.5}.
\end{proof}

\begin{env}[8.1.9]
\label{II.8.1.9}
Sob as hipóteses de \sref{II.8.1.7}, podemos precisar a estrutura dos $\sh{I}_{n,X}$.
Observemos, com efeito, que o homomorfismo
\[
  \widetilde{u}_{-1}:\sh{O}_X\to\sh{O}_X(-1)
\]
corresponde canonicamente a uma seção $s$ de $\sh{O}_X(-1)$ sobre $X$, que chamaremos de \emph{seção canônica} (relativa a $\sh{I}$) \sref[0]{0.5.1.1}.
Por outro lado, no diagrama~\sref{II.8.1.5.5}, as setas horizontais são isomorfismos \sref{II.3.2.7}; substituindo nesse diagrama $h$ por $k$ e $k$ por $-1$, obtemos $\widetilde{u}_{k-1}=1_k\otimes\widetilde{u}_{-1}$ (onde $1_k$ denota a identidade de $\sh{O}_X(k)$), isto é, o homomorfismo $\widetilde{u}_k$ nada mais é que a \emph{multiplicação tensorial pela seção canônica $s$} (para todo $k\in\bb{Z}$).
O homomorfismo $\widetilde{u}$ \sref{II.8.1.5.6} também se interpreta da mesma maneira.

Portanto, para todo $n\geq0$, o homomorfismo $\widetilde{v}_n:\sh{O}_X(n)\to\sh{O}_X$ nada mais é que a multiplicação tensorial por $s^{\otimes n}$; deduz-se então:
\end{env}

\begin{corollary}[8.1.10]
\label{II.8.1.10}
Com as notações de \sref{II.8.1.8}, o espaço subjacente a $X'$ é o conjunto dos pontos $x\in X$ tais que $s(x)=0$, onde $s$ denota a seção canônica de $\sh{O}_X(-1)$.
\end{corollary}

\begin{proof}
Com efeito, se $c_x$ é um gerador da fibra $(\sh{O}_X(1))_x$ em um ponto $x$, então $s_x\otimes c_x$ identifica-se canonicamente com um gerador da fibra de $\sh{I}_{1,X}$ no ponto $x$ e, portanto, é invertível se, e somente se, $s_x\notin\mathfrak{m}_x(\sh{O}_X(-1))_x$, ou, equivalentemente, se, e somente se, $s(x)\neq0$.
\end{proof}

\begin{proposition}[8.1.11]
\label{II.8.1.11}
Seja $Y$ um pré-esquema integral, $\sh{I}$ um ideal fracionário quase-coerente de $\sh{R}(Y)$, e $X$ o $Y$-esquema obtido pela explosão de $\sh{I}$.
Então $\sh{I}\sh{O}_X$ é um $\sh{O}_X$-módulo invertível muito amplo para $Y$.
\end{proposition}

\begin{proof}
Como a questão é local em $Y$ \sref{II.4.4.5}, podemos reduzir ao caso em que $Y=\Spec(A)$, onde $A$ é um anel integral com corpo de frações $K$, e $\sh{I}=\widetilde{\mathfrak{I}}$, sendo $\mathfrak{I}$ um ideal fracionário de $K$;
existe então um elemento $a\neq0$ de $A$ tal que $a\mathfrak{I}\subset A$.
Seja $S=\bigoplus_{n\geq0}\mathfrak{I}^n$;
a aplicação $x\mapsto ax$ é um $A$-isomorfismo de $\mathfrak{I}^{n+1}=(S(1))_n$ sobre $a\mathfrak{I}^{n+1}=a\mathfrak{I}S_n\subset\mathfrak{I}^n=S_n$,
\oldpage[II]{157}
e, portanto, define um \textbf{(TN)}-isomorfismo de grau zero de $S$-módulos graduados $S_+(1)\to a\mathfrak{I}S$.
Por outro lado, $x\mapsto a^{-1}x$ é um isomorfismo de grau zero de $S$-módulos graduados $a\mathfrak{I}S\simto\mathfrak{I}S$.
Obtemos assim, por composição \sref{II.3.2.4}, um isomorfismo de $\sh{O}_X$-módulos $\sh{O}_X(1)\simto\sh{I}\sh{O}_X$ e, como $S$ é gerada por $S_1=\mathfrak{I}$, $\sh{O}_X(1)$ é invertível \sref{II.3.2.5} e muito amplo (\sref{II.4.4.3} e \sref{II.4.4.9}), o que prova a afirmação.
\end{proof}

\subsection{Resultados preliminares sobre a localização de anéis graduados}
\label{subsection:II.8.2}

\begin{env}[8.2.1]
\label{II.8.2.1}
Seja $S$ um anel graduado, que, por enquanto, não supomos concentrado em graus não negativos.
Poremos
\[
\label{II.8.2.1.1}
  S^\geq=\bigoplus_{n\geq0}S_n,
  \qquad
  S^\leq=\bigoplus_{n\leq0}S_n
\tag{8.2.1.1}
\]
que são subanéis graduados de $S$, concentrados, respectivamente, em graus não negativos e não positivos.
Se $f$ é um elemento homogêneo de grau $d$ (positivo ou negativo) de $S$, o anel de frações $S_f=S'$ fica novamente munido de uma estrutura de anel graduado, tomando-se como $S'_n$ ($n\in\bb{Z}$) o conjunto dos $x/f^k$, onde $x\in S_{n+kd}$ ($k\geq0$);
poremos, além disso, $S_{(f)}=S'_0$, e escreveremos $S_f^\geq$ e $S_f^\leq$ em lugar de $S^{\prime\geq}$ e $S^{\prime\leq}$, respectivamente.
Se $d>0$, tem-se
\[
\label{II.8.2.1.2}
  (S^\geq)_f=S_f
\tag{8.2.1.2}
\]
pois, se $x\in S_{n+kd}$ e $n+kd<0$, pode-se escrever $x/f^k=xf^h/f^{h+k}$, e tem-se $n+(h+k)d>0$ para $h$ suficientemente grande e positivo.
Conclui-se, por definição, que
\[
\label{II.8.2.1.3}
  (S^\geq)_{(f)}=(S_f^\geq)_0=S_{(f)}.
\tag{8.2.1.3}
\]

Se $M$ é um $S$-módulo graduado, poremos analogamente
\[
\label{II.8.2.1.4}
  M^\geq=\bigoplus_{n\geq0}M_n,
  \qquad
  M^\leq=\bigoplus_{n\leq0}M_n
\tag{8.2.1.4}
\]
que são, respectivamente, um $S^\geq$-módulo graduado e um $S^\leq$-módulo graduado, e têm por interseção o $S_0$-módulo $M_0$.
Se $f\in S_d$, definimos igualmente $M_f$ como o $S_f$-módulo graduado cujos elementos de grau $n$ são os $z/f^k$, onde $z\in M_{n+kd}$ ($k\geq0$);
denotaremos por $M_{(f)}$ o conjunto dos elementos de grau zero de $M_f$, que é um $S_{(f)}$-módulo, e escreveremos $M_f^\geq$ e $M_f^\leq$ em lugar de $(M_f)^\geq$ e $(M_f)^\leq$, respectivamente.
Se $d>0$, vê-se, como anteriormente, que
\[
\label{II.8.2.1.5}
  (M^\geq)_f=M_f
\tag{8.2.1.5}
\]
e
\[
\label{II.8.2.1.6}
  (M^\geq)_{(f)}=(M_f^\geq)_0=M_{(f)}.
\tag{8.2.1.6}
\]
\end{env}

\begin{env}[8.2.2]
\label{II.8.2.2}
Seja $\bb{z}$ uma indeterminada, que chamaremos de \emph{variável de homogeneização}.
Se $S$ é um anel graduado (com graus inteiros arbitrários), então a álgebra de polinômios\footnote{Não se deve confundir aqui com o uso da notação $\widehat{S}$ para designar o completamento separado de um anel.}
\[
\label{II.8.2.2.1}
  \widehat{S}=S[\bb{z}]
\tag{8.2.2.1}
\]
\oldpage[II]{158}
é uma $S$-álgebra graduada, quando se toma como grau de $f\bb{z}^n$ ($n\geq0$), com $f$ homogêneo,
\[
\label{II.8.2.2.2}
  \deg(f\bb{z}^n)=n+\deg f.
\tag{8.2.2.2}
\]
\end{env}

\begin{lemma}[8.2.3]
\label{II.8.2.3}
\begin{enumerate}
  \item[\rm{(i)}] Existem isomorfismos canônicos de anéis (não graduados)
    \[
    \label{II.8.2.3.1}
      \widehat{S}_{(\bb{z})}
      \simto
      \widehat{S}/(\bb{z}-1)\widehat{S}
      \simto
      S.
    \tag{8.2.3.1}
    \]
  \item[\rm{(ii)}] Existe um isomorfismo canônico de anéis (não graduados)
    \[
    \label{II.8.2.3.2}
      \widehat{S}_{(f)} \simto S_f^\leq
    \tag{8.2.3.2}
    \]
    para todo $f\in S_d$ com $d>0$.
\end{enumerate}
\end{lemma}

\begin{proof}
O primeiro dos isomorfismos de \sref{II.8.2.3.1} foi definido em \sref{II.2.2.5}, e o segundo é trivial;
o isomorfismo $\widehat{S}_{(\bb{z})}\simto S$ assim definido faz, portanto, corresponder à fração $x\bb{z}^n/\bb{z}^{n+k}$, onde $\deg(x)=k$ ($k\geq-n$), o elemento $x$.
O homomorfismo \sref{II.8.2.3.2} faz corresponder à fração $x\bb{z}^n/f^k$, onde $\deg(x)=kd-n$, o elemento $x/f^k$, de grau $-n$ em $S_f^\leq$, e é novamente evidente que se obtém de fato um isomorfismo.
\end{proof}

\begin{env}[8.2.4]
\label{II.8.2.4}
Seja $M$ um $S$-módulo graduado.
É claro que o $S$-módulo
\[
\label{II.8.2.4.1}
  \widehat{M}=M\otimes_S\widehat{S}=M\otimes_SS[\bb{z}]
\tag{8.2.4.1}
\]
é a soma direta dos $S$-módulos $M\otimes_SS\bb{z}^n$, e portanto dos grupos abelianos $M_k\otimes S\bb{z}^n$ ($k\in\bb{Z}$, $n\geq0$);
munimos $\widehat{M}$ de uma estrutura de $\widehat{S}$-módulo graduado tomando
\[
\label{II.8.2.4.2}
  \deg(x\otimes\bb{z}^n)=n+\deg x
\tag{8.2.4.2}
\]
para todo $x$ homogêneo em $M$.
Deixamos ao leitor a tarefa de demonstrar o análogo de \sref{II.8.2.3}:
\end{env}

\begin{lemma}[8.2.5]
\label{II.8.2.5}
\begin{enumerate}
  \item[\rm{(i)}] Existe um di-isomorfismo canônico de módulos (não graduados)
    \[
    \label{II.8.2.5.1}
      \widehat{M}_{(\bb{z})} \simto M.
    \tag{8.2.5.1}
    \]
  \item[\rm{(ii)}] Para todo $f\in S_d$ ($d>0$), existe um di-isomorfismo de módulos (não graduados)
    \[
    \label{II.8.2.5.2}
      \widehat{M}_{(f)} \simto M_f^\leq.
    \tag{8.2.5.2}
    \]
\end{enumerate}
\end{lemma}

\begin{env}[8.2.6]
\label{II.8.2.6}
Seja $S$ um anel graduado em graus \emph{positivos}, e consideremos em $S$ a sequência decrescente de ideais graduados
\[
\label{II.8.2.6.1}
  S_{[n]} = \bigoplus_{m\geq n} S_m
  \qquad\mbox{($n\geq0$)}
\tag{8.2.6.1}
\]
(em particular, tem-se $S_{[0]}=S$ e $S_{[1]}=S_+$).
Como é claro que $S_{[m]}S_{[n]}\subset S_{[m+n]}$, pode-se definir um \emph{anel graduado} $S^\natural$ tomando
\[
\label{II.8.2.6.2}
  S^\natural = \bigoplus_{n\geq0} S_n^\natural
  \quad
  \text{com}
  \quad
  S_n^\natural = S_{[n]}.
\tag{8.2.6.2}
\]
$S_0^\natural$ é então o anel $S$ considerado como anel \emph{não graduado}, e $S^\natural$ é, portanto, uma $S_0^\natural$-álgebra.
Para todo elemento homogêneo $f\in S_d$ ($d>0$), denotaremos por $f^\natural$ o elemento $f$ considerado como pertencente a $S_{[d]}=S_d^\natural$.
Com estas notações:
\end{env}

\oldpage[II]{159}
\begin{lemma}[8.2.7]
\label{II.8.2.7}
Seja $S$ um anel graduado em graus positivos, e seja $f$ um elemento homogêneo de $S_d$ ($d>0$).
Há isomorfismos canônicos de anéis
\[
\label{II.8.2.7.1}
  S_f \simto \bigoplus_{n\in\bb{Z}} S(n)_{(f)}
\tag{8.2.7.1}
\]
\[
\label{II.8.2.7.2}
  (S_f^\geq)_{f/1} \simto S_f
\tag{8.2.7.2}
\]
\[
\label{II.8.2.7.3}
  S_{(f^\natural)}^\natural \simto S_f^\geq
\tag{8.2.7.3}
\]
dos quais os dois primeiros são isomorfismos de anéis graduados.
\end{lemma}

\begin{proof}
É imediato, por definição, que $(S_f)_n=(S(n)_f)_0$, de onde decorre o isomorfismo de \sref{II.8.2.7.1}, que não é senão a identidade.
Por outro lado, como $f/1$ é invertível em $S_f$, existe um isomorfismo canônico $S_f\simto(S_f^\geq)_{f/1}=(S_f)_{f/1}$, em virtude de \sref{II.8.2.1.2} aplicado a $S_f$;
o isomorfismo inverso é, por definição, o isomorfismo de \sref{II.8.2.7.2}.
Finalmente, se $x=\sum_{m\geq n}y_m$ é um elemento de $S_{[n]}$, com $n=kd$, ao elemento $x/(f^\natural)^k$ corresponde o elemento $\sum_m y_m/f^k$ de $S_f^\geq$, e verifica-se imediatamente que isso define um isomorfismo \sref{II.8.2.7.3}.
\end{proof}

\begin{env}[8.2.8]
\label{II.8.2.8}
Se $M$ é um $S$-módulo graduado, poremos analogamente, para todo $n\in\bb{Z}$,
\[
\label{II.8.2.8.1}
  M_{[n]} = \bigoplus_{m\geq n} M_m
\tag{8.2.8.1}
\]
e, como $S_{[m]}M_{[n]} \subset M_{[m+n]}$ ($m\geq0$), podemos definir um $S^\natural$-módulo graduado $M^\natural$ tomando
\[
\label{II.8.2.8.2}
  M^\natural = \bigoplus_{n\in\bb{Z}} M_n^\natural
  \quad
  \text{com}
  \quad
  M_n^\natural = M_{[n]}.
\tag{8.2.8.2}
\]
\end{env}

Deixamos novamente ao leitor a tarefa de demonstrar o seguinte lema:
\begin{lemma}[8.2.9]
\label{II.8.2.9}
Com as notações de \sref{II.8.2.7} e \sref{II.8.2.8}, há di-isomorfismos canônicos de módulos
\[
\label{II.8.2.9.1}
  M_f \simto \bigoplus_{n\in\bb{Z}} M(n)_{(f)}
\tag{8.2.9.1}
\]
\[
\label{II.8.2.9.2}
  (M_f^\geq)_{f/1} \simto M_f
\tag{8.2.9.2}
\]
\[
\label{II.8.2.9.3}
  M_{(f^\natural)}^\natural \simto M_f^\geq
\tag{8.2.9.3}
\]
dos quais os dois primeiros são di-isomorfismos de módulos graduados.
\end{lemma}

\begin{lemma}[8.2.10]
\label{II.8.2.10}
Seja $S$ um anel graduado em graus positivos.
\begin{enumerate}
  \item[\rm{(i)}] Para que $S^\natural$ seja uma $S_0^\natural$-álgebra de tipo finito (resp. um anel noetheriano), é necessário e suficiente que $S$ seja uma $S_0$-álgebra de tipo finito (resp. um anel noetheriano).
  \item[\rm{(ii)}] Para que $S_{n+1}^\natural = S_1^\natural S_n^\natural$ ($n\geq n_0$), é necessário e suficiente que $S_{n+1}=S_1S_n$ ($n\geq n_0$).
  \item[\rm{(iii)}] Para que $S_n^\natural = (S_1^\natural)^n$ ($n\geq n_0$), é necessário e suficiente que $S_n=S_1^n$ ($n\geq n_0$).
  \item[\rm{(iv)}] Se $(f_\alpha)$ é um conjunto de elementos homogêneos de $S_+$ tal que $S_+$ está contido no radical em $S$ do ideal de $S$ gerado pelos $f_\alpha$, então $S_+^\natural$ está contido no radical em $S^\natural$ do ideal de $S^\natural$ gerado pelos $f_\alpha^\natural$.
\end{enumerate}
\end{lemma}

\begin{proof}
\medskip\noindent
\begin{enumerate}
  \item[\rm{(i)}] Se $S^\natural$ é uma $S_0^\natural$-álgebra de tipo finito, então $S_+=S_1^\natural$ é um módulo finitamente gerado sobre $S=S_0^\natural$, por \sref{II.2.1.6}[i], e, portanto, $S$ é uma $S_0$-álgebra de tipo finito \sref{II.2.1.4};
    se $S^\natural$ é um anel noetheriano, então $S_0^\natural=S$ também o é \sref{II.2.1.5}.
    Reciprocamente, se $S$ é uma $S_0$-álgebra
\oldpage[II]{160}
    de tipo finito, sabemos por \sref{II.2.1.6}[ii] que existem $h>0$ e $m_0>0$ tais que $S_{n+h}=S_hS_n$ para $n\geq m_0$;
    podemos evidentemente supor que $m_0\geq h$.
    Além disso, os $S_m$ são $S_0$-módulos finitamente gerados \sref{II.2.1.6}[i].
    Logo, se $n\geq m_0+h$, então $S_n^\natural = S_hS_{n-h}^\natural = S_h^\natural S_{n-h}^\natural$;
    e, se $m<m_0+h$, pondo $E = S_{m_0}+\ldots+S_{m_0+h-1}$, tem-se
    \[
      S_m^\natural = S_m + \ldots + S_{m_0+h-1} + S_hE + S_h^2E + \ldots.
    \]
    Para $1\leq m\leq m_0$, seja $G_m$ a união dos sistemas finitos de geradores dos $S_0$-módulos $S_i$ para $m\leq i\leq m_0+h-1$, considerada como um subconjunto de $S_{[m]}$.
    Para $m_0+1\leq m\leq m_0+h-1$, seja $G_m$ a união dos sistemas finitos de geradores dos $S_0$-módulos $S_i$ para $m\leq i\leq m_0+h-1$ e de $S_hE$, considerada como um subconjunto de $S_{[m]}$.
    É claro que $S_m^\natural=S_0^\natural G_m$ para $1\leq m\leq m_0+h-1$ e, portanto, a união $G$ dos $G_m$ para $1\leq m\leq m_0+h-1$ é um sistema de geradores da $S_0^\natural$-álgebra $S^\natural$.
    Concluímos, assim, que, se $S=S_0^\natural$ é um anel noetheriano, então $S^\natural$ também o é.
  \item[\rm{(ii)}] É claro que, se $S_{n+1}=S_1S_n$ para $n\geq n_0$, então $S_{n+1}^\natural=S_1S_n^\natural$, e, \emph{a fortiori}, $S_{n+1}^\natural=S_1^\natural S_n^\natural$ para $n\geq n_0$.
    Reciprocamente, esta última igualdade pode ser escrita como
    \[
      S_{n+1} + S_{n+2} + \ldots
      =
      (S_1 + S_2 + \ldots)(S_n + S_{n+1} + \ldots)
    \]
    e a comparação dos termos de grau $n+1$ (em $S$) nos dois membros dá $S_{n+1}=S_1S_n$.
  \item[\rm{(iii)}] Se $S_n=S_1^n$ para $n\geq n_0$, então $S_n^\natural=S_1^n+S_1^{n+1}+\ldots$;
    como $S_1^\natural$ contém $S_1+S_1^2+\ldots$, tem-se $S_n^\natural\subset S_1^{\natural n}$ e, portanto, $S_n^\natural=S_1^{\natural n}$ para $n\geq n_0$.
    Reciprocamente, os únicos termos de $S_1^{\natural n}=(S_1+S_2+\ldots)^n$ que são de grau $n$ em $S$ são os de $S_1^n$;
    a igualdade $S_n^\natural=S_1^{\natural n}$ implica, portanto, que $S_n=S_1^n$.
  \item[\rm{(iv)}] Basta mostrar que, se um elemento $g\in S_{k+h}$ é considerado como um elemento de $S_k^\natural$ ($k>0$, $h\geq0$), então existe um inteiro $n>0$ tal que $g^n$ é uma combinação linear (em $S_{kn}^\natural$) dos $f_\alpha^\natural$ com coeficientes em $S^\natural$.
    Por hipótese, existe um inteiro $m_0$ tal que, para $m\geq m_0$, tem-se, \emph{em $S$}, $g^m = \sum_\alpha c_{\alpha m}f_\alpha$, onde os índices $\alpha$ que aparecem aqui são \emph{independentes de $m$};
    além disso, podemos evidentemente supor que os $c_{\alpha m}$ sejam homogêneos, com
    \[
      \deg(c_{\alpha m}) = m(k+h)-\deg f_\alpha
    \]
    em $S$.
    Tomemos então $m_0$ suficientemente grande para assegurar que $km_0>\deg f_\alpha$ para todos os $f_\alpha$ que aparecem em $g^{m_0}$;
    para todo $\alpha$, seja $c'_{\alpha m}$ o elemento $c_{\alpha m}$ considerado como tendo grau $km-\deg f_\alpha$ \emph{em $S^\natural$};
    tem-se então, em $S^\natural$, $g^m = \sum_\alpha c'_{\alpha m}f_\alpha^\natural$, o que conclui a demonstração.
\end{enumerate}
\end{proof}

\begin{env}[8.2.11]
\label{II.8.2.11}
Consideremos a $S_0$-álgebra graduada
\[
\label{II.8.2.11.1}
  S^\natural \otimes_S S_0
  =
  S^\natural/S_+S^\natural
  =
  \bigoplus_{n\geq0} S_{[n]}/S_+S_{[n]}.
\tag{8.2.11.1}
\]

Como $S_n$ é um $S_0$-módulo quociente de $S_{[n]}/S_+S_{[n]}$, há um homomorfismo canônico de $S_0$-álgebras graduadas
\[
\label{II.8.2.11.2}
  S^\natural \otimes_S S_0 \to S
\tag{8.2.11.2}
\]
que é evidentemente \emph{sobrejetivo} e, portanto, corresponde \sref{II.2.9.2} a uma \emph{imersão fechada} canônica
\[
\label{II.8.2.11.3}
  \Proj(S) \to \Proj(S^\natural \otimes_S S_0).
\tag{8.2.11.3}
\]
\end{env}

\oldpage[II]{161}
\begin{proposition}[8.2.12]
\label{II.8.2.12}
O morfismo canônico \sref{II.8.2.11.3} é bijetivo.
Para que o homomorfismo \sref{II.8.2.11.2} seja (TN)-bijetivo, é necessário e suficiente que exista um $n_0$ tal que $S_{n+1}=S_1S_n$ para $n\geq n_0$.
Se esta última condição for satisfeita, então \sref{II.8.2.11.3} é um isomorfismo;
a recíproca é verdadeira quando $S$ é noetheriano.
\end{proposition}

\begin{proof}
Para provar a primeira afirmação, basta \sref{II.2.8.3} mostrar que o núcleo $\mathfrak{I}$ do homomorfismo \sref{II.8.2.11.2} é constituído por elementos \emph{nilpotentes}.
Mas, se $f\in S_{[n]}$ é um elemento cuja classe módulo $S_+S_{[n]}$ pertence a esse núcleo, isso implica que $f\in S_{[n+1]}$;
então $f^{n+1}$, considerado como um elemento de $S_{[n(n+1)]}$, pertence também a $S_+S_{[n(n+1)]}$, pois pode ser escrito como $f\cdot f^n$;
portanto, a classe de $f^{n+1}$ módulo $S_+S_{[n(n+1)]}$ é nula, o que prova nossa afirmação.
Como a hipótese de que $S_{n+1}=S_1S_n$ para $n\geq n_0$ é equivalente a $S_{n+1}^\natural=S_1^\natural S_n^\natural$ para $n\geq n_0$ \sref{II.8.2.10}[ii], essa hipótese equivale, por definição, ao fato de \sref{II.8.2.11.2} ser (TN)-injetivo e, portanto, (TN)-bijetivo, de modo que \sref{II.8.2.11.3} é um isomorfismo, por \sref{II.2.9.1}.
Reciprocamente, se \sref{II.8.2.11.3} for um isomorfismo, então o feixe $\widetilde{\mathfrak{I}}$ sobre $\Proj(S^\natural\otimes_S S_0)$ é nulo \sref{II.2.9.2}[i];
como $S^\natural\otimes_S S_0$ é noetheriano, por ser um quociente de $S^\natural$ \sref{II.8.2.10}[i], concluímos de \sref{II.2.7.3} que $\mathfrak{I}$ satisfaz a condição (TN) e, portanto, $S_{n+1}^\natural=S_1^\natural S_n^\natural$ para $n\geq n_0$, e isso conclui a demonstração, por \sref{II.8.2.10}[ii].
\end{proof}

\begin{env}[8.2.13]
\label{II.8.2.13}
Consideremos agora as injeções canônicas $(S_+)^n\to S_{[n]}$, que definem um homomorfismo injetivo de grau zero de anéis graduados
\[
\label{II.8.2.13.1}
  \bigoplus_{n\geq0} (S_+)^n \to S^\natural.
\tag{8.2.13.1}
\]
\end{env}

\begin{proposition}[8.2.14]
\label{II.8.2.14}
Para que o homomorfismo \sref{II.8.2.13.1} seja um (TN)-isomorfismo, é necessário e suficiente que exista um $n_0$ tal que $S_n=S_1^n$ para todo $n\geq n_0$.
Quando isso ocorre, o morfismo correspondente a \sref{II.8.2.13.1} está definido em toda parte e é um isomorfismo
\[
  \Proj(S^\natural) \simto \Proj(\bigoplus_{n\geq0}(S_+)^n);
\]
a recíproca é verdadeira quando $S$ é noetheriano.
\end{proposition}

\begin{proof}
As duas primeiras afirmações são evidentes, em virtude de \sref{II.8.2.10}[iii] e \sref{II.2.9.1}.
A terceira decorrerá de \sref{II.8.2.10}[i e iii] e do lema seguinte:
\begin{lemma}[8.2.14.1]
\label{II.8.2.14.1}
Seja $T$ um anel graduado em graus positivos que seja também uma $T_0$-álgebra de tipo finito.
Se o morfismo correspondente ao homomorfismo injetivo $\bigoplus_{n\geq0}T_1^n\to T$ estiver definido em toda parte e for também um isomorfismo $\Proj(T)\to\Proj(\bigoplus_{n\geq0}T_1^n)$, então existe um $n_0$ tal que $T_n=T_1^n$ para $n\geq n_0$.
\end{lemma}

Sejam $g_i$ ($1\leq i\leq r$) geradores do $T_0$-módulo $T_1$.
A hipótese implica, antes de tudo, que os $D_+(g_i)$ recobrem $\Proj(T)$ \sref{II.2.8.1}.
Seja $(h_j)_{1\leq j\leq s}$ um sistema de elementos homogêneos de $T_+$, com $\deg(h_j)=n_j$, que forme, juntamente com os $g_i$, um sistema de geradores do ideal $T_+$, ou, equivalentemente \sref{II.2.1.3}, um sistema de geradores de $T$ como $T_0$-álgebra;
se pusermos $T'=\bigoplus_{n\geq0}T_1^n$, então o elemento $h_j/g_i^{n_j}$ do anel $T_{(g_i)}$ deve, por hipótese, pertencer ao subanel $T'_{(g_i)}$, e, portanto, existe um inteiro $k$ tal que $T_1^k h_j\subset T_1^{k+n_j}$ para todo $j$.
Concluímos, assim, por indução sobre o expoente $r$, que $T_1^k h_j^r \subset T'$ para todo $r\geq1$, e, pela definição dos $h_j$, temos, portanto, $T_1^k T\subset T'$.
Por outro lado, para todo $j$, existe um inteiro $m_j$ tal que $h_j^{m_j}$ pertence ao ideal de $T$ gerado pelos $g_i$ \sref{II.2.3.14}, de modo que $h_j^{m_j}\in T_1 T$, e
\oldpage[II]{162}
$h_j^{m_j k}\in T_1^k T\subset T'$.
Há, portanto, um inteiro $m_0\geq k$ tal que $h_j^m\in T_1^{m n_j}$ para $m\geq m_0$.
Assim, se $q$ é o maior dos inteiros $n_j$, então $n_0=q s m_0+k$ é o número procurado.
Com efeito, um elemento de $T_n$, para $n\geq n_0$, é uma soma de monômios pertencentes a $T_1^\alpha u$, onde $u$ é um produto de potências dos $h_j$;
se $\alpha\geq k$, decorre do que precede que $T_1^\alpha u\subset T_1^n$;
no caso contrário, um dos expoentes dos $h_j$ é $\geq m_0$, de modo que $u\in T_1^\beta v$, onde $\beta\geq k$ e $v$ é novamente um produto de potências dos $h_j$;
podemos então reduzir ao caso anterior e, assim, concluímos que $T_1^\alpha u\subset T_1^n$ em todos os casos.
\end{proof}

\begin{remark}[8.2.15]
\label{II.8.2.15}
A condição $S_n=S_1^n$ para $n\geq n_0$ implica evidentemente que $S_{n+1}=S_1S_n$ para $n\geq n_0$, mas a recíproca não é necessariamente verdadeira, mesmo supondo que o anel $S$ seja noetheriano.
Por exemplo, seja $K$ um corpo, $A=K[\bb{x}]$ e $B=K[\bb{y}]/\bb{y}^2K[\bb{y}]$, onde $\bb{x}$ e $\bb{y}$ são indeterminadas, tomando-se $\bb{x}$ de grau 1 e $\bb{y}$ de grau 2, e seja $S=A\otimes_K B$, de modo que $S$ é uma álgebra graduada sobre $K$ que tem uma base constituída pelos elementos $1$, $\bb{x}^n$ ($n\geq1$) e $\bb{x}^n\bb{y}$ ($n\geq0$).
É imediato que $S_{n+1}=S_1S_n$ para $n\geq2$, mas $S_1^n=K\bb{x}^n$, ao passo que $S_n=K\bb{x}^n+K\bb{x}^{n-2}\bb{y}$ para $n\geq2$.
\end{remark}

\subsection{Cones projetantes}
\label{subsection:II.8.3}

\begin{env}[8.3.1]
\label{II.8.3.1}
Seja $Y$ um pré-esquema;
em todo este número, consideraremos apenas \emph{$Y$-pré-esquemas} e \emph{$Y$-morfismos}.
Seja $\sh{S}$ uma $\sh{O}_Y$-álgebra graduada quase-coerente em graus \emph{positivos};
\emph{suporemos, além disso, que $\sh{S}_0=\sh{O}_Y$}.
De acordo com as notações introduzidas em \sref{II.8.2.2}, poremos
\[
\label{II.8.3.1.1}
  \widehat{\sh{S}} = \sh{S}[\bb{z}] = \sh{S} \otimes_{\sh{O}_Y} \sh{O}_Y[\bb{z}]
\tag{8.3.1.1}
\]
que consideramos como uma $\sh{O}_Y$-álgebra graduada em graus positivos, definindo os graus pela fórmula \sref{II.8.2.2.2}, de modo que, para todo aberto afim $U$ de $Y$, tem-se
\[
  \Gamma(U,\widehat{\sh{S}}) = (\Gamma(U,\sh{S}))[\bb{z}].
\]
No que segue, poremos
\[
\label{II.8.3.1.2}
  X = \Proj(\sh{S}),
  \quad
  C = \Spec(\sh{S}),
  \quad
  \widehat{C} = \Proj(\widehat{\sh{S}})
\tag{8.3.1.2}
\]
(onde, na definição de $C$, $\sh{S}$ é considerada como uma $\sh{O}_Y$-álgebra não graduada), e diremos que $C$ (resp. $\widehat{C}$) é o \emph{cone afim} (resp. o \emph{cone projetivo}) definido por $\sh{S}$;
diremos também ``cone'' em vez de ``cone afim''.
Por abuso de linguagem, diremos ainda que $C$ (resp. $\widehat{C}$) é o \emph{cone projetante afim de $X$} (resp. o \emph{cone projetante projetivo de $X$}), ficando subentendido que o pré-esquema $X$ é dado sob a forma $\Proj(\sh{S})$;
por fim, diremos que $\widehat{C}$ é o \emph{fecho projetivo} de $C$ (entendendo-se que o dado de $\sh{S}$ figura na estrutura de $C$).
\end{env}
\begin{proposition}[8.3.2]
\label{II.8.3.2}
Existem $Y$-morfismos canônicos
\[
\label{II.8.3.2.1}
  Y \xrightarrow{\varepsilon} C \xrightarrow{i} \widehat{C}
\tag{8.3.2.1}
\]
\[
\label{II.8.3.2.2}
  X \xrightarrow{j} \widehat{C}
\tag{8.3.2.2}
\]

\oldpage[II]{163}
tais que $\varepsilon$ e $j$ são imersões fechadas, e $i$ é um morfismo afim que é uma imersão aberta dominante, para a qual
\[
\label{II.8.3.2.3}
  i(C) = \widehat{C}\setmin j(X);
\tag{8.3.2.3}
\]
além disso, $\widehat{C}$ é o menor subpré-esquema fechado de $\widehat{C}$ que contém $i(C)$.
\end{proposition}

\begin{proof}
Para definir $i$, consideremos o aberto de $\widehat{C}$ dado por
\[
\label{II.8.3.2.4}
  \widehat{C}_{\bb{z}} = \Spec(\widehat{\sh{S}}/(\bb{z}-1)\widehat{\sh{S}})
\tag{8.3.2.4}
\]
\sref{II.3.1.4}, identificando-se canonicamente $\bb{z}$ com uma seção de $\widehat{\sh{S}}$ sobre $Y$.
O isomorfismo $i: C\simto\widehat{C}_{\bb{z}}$ corresponde então ao isomorfismo canônico \sref{II.8.2.3.1}
\[
  \widehat{\sh{S}}/(\bb{z}-1)\widehat{\sh{S}} \simto \sh{S}.
\]

O morfismo $\varepsilon$ corresponde ao homomorfismo de aumentação $\sh{S}\to\sh{S}_0=\sh{O}_Y$, cujo núcleo é $\sh{S}_+$ \sref{II.1.2.7}; como esse homomorfismo é sobrejetivo, $\varepsilon$ é uma imersão fechada \sref{II.1.4.10}.
Por fim, $j$ corresponde \sref{II.3.5.1} ao homomorfismo sobrejetivo de grau zero $\widehat{\sh{S}}\to\sh{S}$, cuja restrição a $\sh{S}$ é a identidade e que é nulo sobre $\bb{z}\widehat{\sh{S}}$, que é seu núcleo;
$j$ está definido em toda parte e é uma imersão fechada, por \sref{II.3.6.2}.

Para demonstrar as outras afirmações de \sref{II.8.3.2}, podemos evidentemente nos restringir ao caso em que $Y=\Spec(A)$ é afim e $\sh{S}=\widetilde{S}$, sendo $S$ uma $A$-álgebra graduada, donde $\widehat{\sh{S}}=(\widehat{S})\supertilde$;
os elementos homogêneos $f$ de $S_+$ podem então ser identificados com seções de $\widehat{\sh{S}}$ sobre $Y$, e o aberto de $\widehat{C}$, denotado por $D_+(f)$ em \sref{II.2.3.3}, pode também ser escrito $\widehat{C}_f$ \sref{II.3.1.4};
do mesmo modo, o aberto de $C$ denotado por $D(f)$ em \sref[I]{I.1.1.1} pode ser escrito $C_f$ \sref[0]{0.5.5.2}.
Posto isso, resulta de \sref{II.2.3.14} e da definição de $\widehat{S}$ que, neste caso, os abertos $\widehat{C}_{\bb{z}}=i(C)$ e $\widehat{C}_f$ (com $f$ homogêneo em $S_+$) formam um \emph{recobrimento} de $\widehat{C}$.
Além disso, com essas notações,
\[
\label{II.8.3.2.5}
  i^{-1}(\widehat{C}_f) = C_f;
\tag{8.3.2.5}
\]
com efeito, $\widehat{C}_f\cap i(C) = \widehat{C}_f\cap\widehat{C}_{\bb{z}} = \widehat{C}_{f\bb{z}} = \Spec(\widehat{S}_{(f\bb{z})})$.
Mas, se $d=\deg(f)$, então $\widehat{S}_{(f\bb{z})}$ é canonicamente isomorfo a $(\widehat{S}_{(\bb{z})})_{f/\bb{z}^d}$ \sref{II.2.2.2}, e resulta da definição do isomorfismo em \sref{II.8.2.3.1} que a imagem de $(\widehat{S}_{(\bb{z})})_{f/\bb{z}^d}$ pelo isomorfismo correspondente dos anéis de frações é exatamente $S_f$.
Como $C_f=\Spec(S_f)$, isso demonstra \sref{II.8.3.2.5} e mostra, ao mesmo tempo, que o morfismo $i$ é afim;
além disso, a restrição de $i$ a $C_f$, considerada como um morfismo para $\widehat{C}_f$, corresponde \sref[I]{I.1.7.3} ao homomorfismo canônico $\widehat{S}_{(f)}\to\widehat{S}_{(f\bb{z})}$, e, pelo que precede e por \sref{II.8.2.3.2}, podemos enunciar o resultado seguinte:
\begin{env}[8.3.2.6]
\label{II.8.3.2.6}
Se $Y=\Spec(A)$ é afim e $\sh{S}=\widetilde{S}$, então, para todo $f$ homogêneo em $S_+$, $\widehat{C}_f$ identifica-se canonicamente com $\Spec(S_f^\leq)$, e o morfismo $C_f\to\widehat{C}_f$ obtido por restrição de $i$ corresponde então à injeção canônica $S_f^\leq\to S_f$.
\end{env}

Observemos agora que (para $Y$ afim) o complementar de $\widehat{C}_{\bb{z}}$ em $\widehat{C}=\Proj(\widehat{S})$
\oldpage[II]{164}
é, por definição, o conjunto dos ideais primos homogêneos de $\widehat{S}$ que contêm $\bb{z}$, isto é, exatamente $j(X)$, pela definição de $j$, o que demonstra \sref{II.8.3.2.3}.

Por fim, para demonstrar a última afirmação de \sref{II.8.3.2}, podemos supor que $Y$ seja afim.
Com as notações precedentes, observemos que, no anel $\widehat{S}$, $\bb{z}$ não é divisor de zero;
como o fecho de $i(C)$ é $\widehat{C}$, basta demonstrar o lema seguinte:
\begin{lemma}[8.3.2.7]
\label{II.8.3.2.7}
Sejam $T$ um anel graduado em graus positivos, $Z=\Proj(T)$ e $g$ um elemento homogêneo de $T$ de grau $d>0$.
Se $g$ não é divisor de zero em $T$, então $Z$ é o menor subpré-esquema fechado de $Z$ que contém $Z_g=D_+(g)$.
\end{lemma}

Por \sref[I]{I.4.1.9}, a questão é local em $Z$;
para todo elemento homogêneo $h\in T_e$ ($e>0$), basta, portanto, demonstrar que $Z_h$ é o menor subpré-esquema fechado de $Z_h$ que contém $Z_{gh}$;
resulta das definições e de \sref[I]{I.4.3.2} que essa condição equivale a exigir que o homomorfismo canônico $T_{(h)}\to T_{(gh)}$ seja \emph{injetivo}.
Mas esse homomorfismo pode ser identificado com o homomorfismo canônico $T_{(h)}\to(T_{(h)})_{g^e/h^d}$ \sref{II.2.2.3}.
Ora, como $g^e$ não é divisor de zero em $T$, $g^e/h^d$ não é divisor de zero em $T_h$ (nem, \emph{a fortiori}, em $T_{(h)}$), pois a relação $(g^e/h^d)(t/h^m)=0$ (para $t\in T$ e $m>0$) implica a existência de um $n>0$ tal que $h^ng^et=0$, donde $h^nt=0$ e, portanto, $t/h^m=0$ em $T_h$.
Isso conclui a demonstração \sref[0]{0.1.2.2}.
\end{proof}
\begin{env}[8.3.3]
\label{II.8.3.3}
Identificaremos frequentemente o cone afim $C$ com o subpré-esquema induzido pelo cone projetivo $\widehat{C}$ sobre o aberto $i(C)$ por meio da imersão aberta $i$.
O subpré-esquema fechado de $C$ associado à imersão fechada $\varepsilon$ é chamado de \emph{pré-esquema dos vértices} de $C$;
dizemos também que $\varepsilon$, que é uma $Y$-seção de $C$, é a \emph{seção vértice}, ou a \emph{seção nula}, de $C$;
podemos identificar $Y$ com o pré-esquema dos vértices de $C$ por meio de $\varepsilon$.
Além disso, $i\circ\varepsilon$ é uma $Y$-seção de $\widehat{C}$, e portanto também uma imersão fechada \sref[I]{I.5.4.6}, correspondente ao homomorfismo canônico sobrejetivo de grau zero $\widehat{\sh{S}} = \sh{S}[\bb{z}] \to \sh{O}_Y[\bb{z}]$ \sref{II.3.1.7}, cujo núcleo é $\sh{S}_+[\bb{z}] = \sh{S}_+\widehat{\sh{S}}$;
o subpré-esquema de $\widehat{C}$ associado a essa imersão fechada também é chamado de \emph{pré-esquema dos vértices} de $\widehat{C}$, e $i\circ\varepsilon$ de \emph{seção vértice} de $\widehat{C}$;
ele pode ser identificado com $Y$ por meio de $i\circ\varepsilon$.
Finalmente, o subpré-esquema fechado de $\widehat{C}$ associado a $j$ é chamado de \emph{lugar no infinito} de $\widehat{C}$, e pode ser identificado com $X$ por meio de $j$.
\end{env}

\begin{env}[8.3.4]
\label{II.8.3.4}
Os subpré-esquemas de $C$ (resp. $\widehat{C}$) induzidos sobre os \emph{abertos}
\[
\label{II.8.3.4.1}
  E = C \setmin \varepsilon(Y),
  \qquad
  \widehat{E} = \widehat{C} \setmin i(\varepsilon(Y))
\tag{8.3.4.1}
\]
são chamados, respectivamente (por abuso de linguagem), de \emph{cone afim sem vértice} e \emph{cone projetivo sem vértice} definidos por $\sh{S}$;
observemos que, apesar dessa terminologia, \emph{$E$ não é necessariamente afim} sobre $Y$, nem $\widehat{E}$ projetivo sobre $Y$ \sref{II.8.4.3}.
Quando identificamos $C$ com $i(C)$, temos, portanto, para os espaços subjacentes
\[
\label{II.8.3.4.2}
  C \cup \widehat{E} = \widehat{C},
  \qquad
  C \cap \widehat{E} = E
\tag{8.3.4.2}
\]
de modo que $\widehat{C}$ pode ser considerado como obtido por \emph{colagem} dos subpré-esquemas abertos $C$ e $\widehat{E}$;
além disso, por \sref{II.8.3.2.3},
\[
\label{II.8.3.4.3}
  E = \widehat{E} \setmin j(X).
\tag{8.3.4.3}
\]

\oldpage[II]{165}
Se $Y=\Spec(A)$ é afim, então, com as notações de \sref{II.8.3.2},
\[
\label{II.8.3.4.4}
  E = \bigcup C_f,
  \qquad
  \widehat{E} = \bigcup \widehat{C}_f,
  \qquad
  C_f = C \cap \widehat{C}_f
\tag{8.3.4.4}
\]
onde $f$ percorre o conjunto dos elementos homogêneos de $S_+$ (ou apenas um subconjunto $M$ desse conjunto, tal que $M$ gere um ideal de $S_+$ cujo radical em $S_+$ seja o próprio $S_+$, ou, equivalentemente, tal que os $X_f$, para $f\in M$, recubram $X$ \sref{II.2.3.14}).
A colagem de $C$ e $\widehat{C}_f$ ao longo de $C_f$ fica então determinada pelos morfismos de injeção $C_f\to C$ e $C_f\to\widehat{C}_f$, que, como vimos \sref{II.8.3.2.6}, correspondem respectivamente aos homomorfismos canônicos $S\to S_f$ e $S_f^\leq\to S_f$.
\end{env}
\begin{proposition}[8.3.5]
\label{II.8.3.5}
Com as notações de \sref{II.8.3.1} e \sref{II.8.3.4}, o morfismo associado \sref{II.3.5.1} à injeção canônica $\vphi:\sh{S} \to \widehat{\sh{S}} = \sh{S}[\bb{z}]$ é um morfismo afim sobrejetivo (chamado de retração canônica)
\[
\label{II.8.3.5.1}
  p:\widehat{E} \to X
\tag{8.3.5.1}
\]
tal que
\[
\label{II.8.3.5.2}
  p \circ j = 1_X.
\tag{8.3.5.2}
\]
\end{proposition}

\begin{proof}
Para demonstrar a proposição, podemos nos restringir ao caso em que $Y$ é afim.
Levando em conta a expressão \sref{II.8.3.4.4} de $\widehat{E}$, o fato de que o domínio de definição $G(\vphi)$ de $p$ é igual a $\widehat{E}$ decorrerá da primeira das seguintes afirmações:
\begin{env}[8.3.5.3]
\label{II.8.3.5.3}
Se $Y=\Spec(A)$ é afim e $\sh{S}=\widetilde{S}$, então, para todo $f\in S_+$ homogêneo,
\[
\label{II.8.3.5.4}
  p^{-1}(X_f) = \widehat{C}_f
\tag{8.3.5.4}
\]
e a restrição de $p$ a $\widehat{C}_f=\Spec(S_f^\leq)$, considerada como morfismo de $\widehat{C}_f$ para $X_f$, corresponde à injeção canônica $S_{(f)}\to S_f^\leq$.
Se, além disso, $f\in S_1$, então $\widehat{C}_f$ é isomorfo a $X_f\otimes_{\bb{Z}}\bb{Z}[T]$ (onde $T$ é uma indeterminada).
\end{env}

De fato, a fórmula~\sref{II.8.3.5.4} é precisamente um caso particular de \sref{II.2.8.1.1}, e a segunda afirmação é exatamente a definição de $\Proj(\vphi)$ quando $Y$ é afim \sref{II.2.8.1}.
A fórmula~\sref{II.8.3.5.2} e a sobrejetividade de $p$ decorrem então do fato de que a composição $\sh{S}\to\widehat{\sh{S}}\to\sh{S}$ dos homomorfismos canônicos é a identidade em $\sh{S}$.
Finalmente, a última afirmação de \sref{II.8.3.5.3} decorre do fato de que $S_f^\leq$ é isomorfo a $S_{(f)}[T]$ quando $f\in S_1$ \sref{II.2.2.1}.
\end{proof}
\begin{corollary}[8.3.6]
\label{II.8.3.6}
A restrição
\[
\label{II.8.3.6.1}
  \pi: E \to X
\tag{8.3.6.1}
\]
de $p$ a $E$ é um morfismo afim sobrejetivo.
Se $Y$ é afim e $f$ homogêneo em $S_+$, então
\[
\label{II.8.3.6.2}
  \pi^{-1}(X_f) = C_f
\tag{8.3.6.2}
\]
e a restrição de $\pi$ a $C_f$ corresponde à injeção canônica $S_{(f)}\to S_f$.
Se, além disso, $f\in S_1$, então $C_f$ é isomorfo a $X_f\otimes_{\bb{Z}}\bb{Z}[T,T^{-1}]$ (onde $T$ é uma indeterminada).
\end{corollary}

\begin{proof}
A fórmula~\sref{II.8.3.6.2} decorre imediatamente de \sref{II.8.3.5.3} e \sref{II.8.3.2.5}, e mostra a sobrejetividade de $\pi$;
já vimos que a imersão $i$, restrita a $C_f$, corresponde
\oldpage[II]{166}
à injeção $S_f^\leq\to S_f$ \sref{II.8.3.2}.
Finalmente, a última afirmação é consequência do fato de que, para $f\in S_1$, $S_f$ é isomorfo a $S_{(f)}[T,T^{-1}]$ \sref{II.2.2.1}.
\end{proof}
\begin{remark}[8.3.7]
\label{II.8.3.7}
Quando $Y$ é afim, os elementos do espaço subjacente a $E$ são os ideais primos $\mathfrak{p}$ de $S$ (não necessariamente homogêneos) que não contêm $S_+$, pela definição da imersão $\varepsilon$ \sref{II.8.3.2}.
Para tal ideal $\mathfrak{p}$, os $\mathfrak{p}\cap S_n$ satisfazem evidentemente as condições de \sref{II.2.1.9}, de modo que existe um único ideal primo \emph{homogêneo} $\mathfrak{q}$ de $S$ tal que $\mathfrak{q}\cap S_n=\mathfrak{p}\cap S_n$ para todo $n$;
a aplicação $\pi:E\to X$ dos espaços subjacentes é então interpretada pela relação
\[
\label{II.8.3.7.1}
  \pi(\mathfrak{p}) = \mathfrak{q}.
\tag{8.3.7.1}
\]

De fato, para demonstrar essa relação, basta considerar um elemento $f$ homogêneo em $S_+$ tal que $\mathfrak{p}\in D(f)$, e observar que $\mathfrak{q}_{(f)}$ é a imagem inversa de $\mathfrak{p}_f$ pela injeção $S_{(f)}\to S_f$.
\end{remark}

\begin{corollary}[8.3.8]
\label{II.8.3.8}
Se $\sh{S}$ é gerada por $\sh{S}_1$, os morfismos $p$ e $\pi$ são de tipo finito;
para todo $x\in X$, a fibra $p^{-1}(x)$ é isomorfa a $\Spec(\kres(x)[T])$, e a fibra $\pi^{-1}(x)$ é isomorfa a $\Spec(\kres(x)[T,T^{-1}])$.
\end{corollary}

\begin{proof}
Isso decorre imediatamente de \sref{II.8.3.5} e \sref{II.8.3.6}, observando que, quando $Y$ é afim e $S$ é gerado por $S_1$, os $X_f$, para $f\in S_1$, formam um recobrimento de $X$ \sref{II.2.3.14}.
\end{proof}
\begin{remark}[8.3.9]
\label{II.8.3.9}
O cone afim sem vértice correspondente à $\sh{O}_Y$-álgebra graduada $\sh{O}_Y[T]$ (onde $T$ é uma indeterminada) identifica-se com $G_m=\Spec(\sh{O}_Y[T,T^{-1}])$, pois é precisamente $C_T$, como vimos em \sref{II.8.3.2} (veja \sref{II.8.4.4} para um resultado mais geral).
Esse pré-esquema é canonicamente munido de uma estrutura de \textquotedblleft\emph{$Y$-esquema em grupos comutativos}\textquotedblright.
Essa noção será desenvolvida em detalhe mais adiante, mas pode ser resumida brevemente aqui como segue.
Um $Y$-esquema em grupos é um $Y$-esquema $G$ munido de dois $Y$-morfismos, $p:G\times_Y G\to G$ e $s:G\to G$, que satisfazem condições formalmente análogas aos axiomas da lei de composição e da aplicação de inversão de um grupo: o diagrama
\[
  \xymatrix{
    G \times G \times G
      \ar[r]^{p \times 1}
      \ar[d]_{1 \times p}
  & G \times G
      \ar[d]^{p}
  \\G \times G
      \ar[r]_{p}
  & G
  }
\]
deve ser comutativo (\textquotedblleft associatividade\textquotedblright), e deve valer uma condição correspondente ao fato de que, nos grupos, as aplicações
\[
  (x,y)
  \mapsto
  (x,x^{-1},y)
  \mapsto
  (x,x^{-1}y)
  \mapsto
  x(x^{-1}y)
\]
e
\[
  (x,y)
  \mapsto
  (x,x^{-1},y)
  \mapsto
  (x,yx^{-1})
  \mapsto
  (yx^{-1})x
\]
devem ambas se reduzir a $(x,y)\mapsto y$;
a sequência de morfismos correspondente, por exemplo, à primeira aplicação composta é
\[
  G \times G
  \xrightarrow{(1,s) \times 1}
  G \times G \times G
  \xrightarrow{1 \times p}
  G\times G
  \xrightarrow{p}
  G
\]
e o leitor escreverá de modo análogo a segunda sequência.

\oldpage[II]{167}
É imediato \sref[I]{I.3.4.3} que dar uma estrutura de $Y$-esquema em grupos sobre um $Y$-esquema $G$ equivale a dar, para \emph{todo} $Y$-pré-esquema $Z$, uma estrutura de \emph{grupo} no conjunto $\Hom_Y(Z,G)$, de modo que, para todo $Y$-morfismo $Z\to Z'$, a aplicação correspondente $\Hom_Y(Z',G)\to\Hom_Y(Z,G)$ seja um homomorfismo de grupos.
No caso particular de $G_m$ considerado aqui, $\Hom_Y(Z,G_m)$ identifica-se com o conjunto das $Z$-seções de $Z\times_Y G_m$ \sref[I]{I.3.3.14}, portanto com o conjunto das $Z$-seções de $\Spec(\sh{O}_Z[T,T^{-1}])$;
por fim, o mesmo raciocínio de \sref[I]{I.3.3.15} mostra que esse conjunto se identifica canonicamente com o conjunto dos elementos \emph{invertíveis} do anel $\Gamma(Z,\sh{O}_Z)$, e a estrutura de grupo nesse conjunto provém da multiplicação no anel $\Gamma(Z,\sh{O}_Z)$.
O leitor verificará que os morfismos $p$ e $s$ acima são obtidos da seguinte maneira: correspondem, por \sref{II.1.2.7} e \sref{II.1.4.6}, aos homomorfismos de $\sh{O}_Y$-álgebras
\begin{align*}
  \pi: &\sh{O}_Y[T,T^{-1}] \to \sh{O}_Y[T,T^{-1},T',T'^{-1}]
\\\sigma: &\sh{O}_Y[T,T^{-1}] \to \sh{O}_Y[T,T^{-1}]
\end{align*}
e são inteiramente definidos pelos dados $\pi(T)=TT'$ e $\sigma(T)=T^{-1}$.

Dito isso, $G_m$ pode ser considerado como um \textquotedblleft\emph{domínio universal de operadores}\textquotedblright\ para todo \emph{cone afim} $C=\Spec(\sh{S})$, onde $\sh{S}$ é uma $\sh{O}_Y$-álgebra quase-coerente graduada em graus positivos.
Isso significa que se pode definir canonicamente um $Y$-morfismo $G_m\times_Y C\to C$ com as propriedades formais de uma lei externa de um conjunto munido de um grupo de operadores;
ou, como acima para esquemas em grupos, que se pode dar, para todo $Y$-pré-esquema $Z$, uma lei externa em $\Hom_Y(Z,C)$ cujo conjunto de operadores seja o grupo $\Hom_Y(Z,G_m)$, com os axiomas usuais dos conjuntos munidos de um grupo de operadores e uma condição de compatibilidade em relação aos $Y$-morfismos $Z\to Z'$.
No caso presente, o morfismo $G_m\times_Y C\to C$ é definido por um homomorfismo de $\sh{O}_Y$-álgebras $\sh{S} \to \sh{S}\otimes_{\sh{O}_Y}\sh{O}_Y[T,T^{-1}] = \sh{S}[T,T^{-1}]$, que associa a cada seção $s_n\in\Gamma(U,\sh{S}_n)$ (onde $U$ é um aberto de $Y$) a seção $s_n T^n\in\Gamma(U,\sh{S}\otimes_{\sh{O}_Y}\sh{O}_Y[T,T^{-1}])$.

Reciprocamente, suponhamos dada uma $\sh{O}_Y$-álgebra quase-coerente $\sh{S}$, \emph{a priori não graduada}, e, sobre $C=\Spec(\sh{S})$, uma estrutura de \textquotedblleft\emph{$Y$-esquema em conjuntos munidos de um grupo de operadores}\textquotedblright\ cujo domínio de operadores seja o $Y$-esquema em grupos $G_m$;
obtém-se então canonicamente uma \emph{graduação} de $\sh{O}_Y$-álgebra em $\sh{S}$.
De fato, dar um $Y$-morfismo $G_m\times_Y C\to C$ equivale a dar um homomorfismo de $\sh{O}_Y$-álgebras $\psi:\sh{S}\to\sh{S}[T,T^{-1}]$, que se escreve $\psi=\sum_{n\in\bb{Z}}\psi_n T^n$, onde os $\psi_n:\sh{S}\to\sh{S}$ são homomorfismos de $\sh{O}_Y$-módulos (com $\psi_n(s)=0$ salvo para um número finito de índices $n$ em uma vizinhança de cada ponto, para toda seção $s\in\Gamma(U,\sh{S})$ sobre qualquer aberto $U$ de $Y$).
Verifica-se então que os axiomas dos conjuntos munidos de um grupo de operadores implicam que os $\psi_n(\sh{S})=\sh{S}_n$ definem uma graduação (em graus inteiros arbitrários) de $\sh{O}_Y$-álgebra em $\sh{S}$, sendo os $\psi_n$ os projetores correspondentes.
Obtém-se assim uma noção de estrutura de \textquotedblleft\emph{cone afim}\textquotedblright\ em todo $Y$-esquema afim, definida de maneira \textquotedblleft geométrica\textquotedblright\ sem recorrer a uma graduação prévia.
\oldpage[II]{168}
Não desenvolveremos mais esse ponto de vista aqui, e deixamos ao leitor a tarefa de formular precisamente as definições e os resultados correspondentes às indicações anteriores.
\end{remark}
\subsection{Fecho projetivo de um fibrado vetorial}
\label{subsection:II.8.4}

\begin{env}[8.4.1]
\label{II.8.4.1}
Sejam $Y$ um pré-esquema e $\sh{E}$ um $\sh{O}_Y$-módulo quase-coerente.
Tomando como $\sh{S}$ a $\sh{O}_Y$-álgebra graduada $\bb{S}_{\sh{O}_Y}(\sh{E})$, a definição~\sref{II.8.3.1.1} mostra que $\widehat{\sh{S}}$ se identifica com $\bb{S}_{\sh{O}_Y}(\sh{E}\oplus\sh{O}_Y)$.
Como o cone afim $\Spec(\sh{S})$ definido por $\sh{S}$ é, por definição, $\bb{V}(\sh{E})$, e $\Proj(\sh{S})$ é, por definição, $\bb{P}(\sh{E})$, obtemos:
\end{env}
\begin{proposition}[8.4.2]
\label{II.8.4.2}
O fecho projetivo de um fibrado vetorial $\bb{V}(\sh{E})$ sobre $Y$ é canonicamente isomorfo a $\bb{P}(\sh{E}\oplus\sh{O}_Y)$, e o lugar no infinito deste último é canonicamente isomorfo a $\bb{P}(\sh{E})$.
\end{proposition}
\begin{remark}[8.4.3]
\label{II.8.4.3}
Tomemos, por exemplo, $\sh{E}=\sh{O}_Y^r$ com $r\geq2$;
os cones sem vértice $E$ e $\widehat{E}$ definidos por $\sh{S}$ não são então nem afins nem projetivos sobre $Y$ se $Y\neq\emp$.
A segunda afirmação é imediata, pois $\widehat{C}=\bb{P}(\sh{O}_Y^{r+1})$ é projetivo sobre $Y$, e os espaços subjacentes a $E$ e $\widehat{E}$ são abertos não fechados de $\widehat{C}$; portanto, as imersões canônicas $E\to\widehat{C}$ e $\widehat{E}\to\widehat{C}$ não são projetivas \sref{II.5.5.3}, e concluímos por \sref{II.5.5.5}[v].
Por outro lado, supondo $Y=\Spec(A)$ afim e, por exemplo, $r=2$, temos $C=\Spec(A[T_1,T_2])$, e $E$ é o pré-esquema induzido por $C$ no aberto $D(T_1)\cup D(T_2)$;
mas já vimos que este último não é afim \sref[I]{I.5.5.11};
\emph{a fortiori}, $\widehat{E}$ não pode ser afim, pois $E$ é o aberto onde a seção $\bb{z}$ sobre $\widehat{E}$ não se anula \sref{II.8.3.2}.

No entanto:
\end{remark}
\begin{proposition}[8.4.4]
\label{II.8.4.4}
Se $\sh{L}$ é um $\sh{O}_Y$-módulo invertível, existem isomorfismos canônicos para os cones sem vértice $E$ e $\widehat{E}$ correspondentes a $C=\bb{V}(\sh{L})$:
\[
\label{eq:2.8.4.4.1}
  \Spec\left(
    \bigoplus_{n\in\bb{Z}}\sh{L}^{\otimes n}
  \right)
  \simto
  E
\tag{8.4.4.1}
\]
\[
\label{eq:2.8.4.4.2}
  \bb{V}(\sh{L}^{-1})
  \simto
  \widehat{E}.
\tag{8.4.4.2}
\]

Além disso, existe um isomorfismo canônico do fecho projetivo de $\bb{V}(\sh{L})$ sobre o de $\bb{V}(\sh{L}^{-1})$ que transforma a seção nula (resp. o lugar no infinito) do primeiro no lugar no infinito (resp. na seção nula) do segundo.
\end{proposition}

\begin{proof}
Aqui temos $\sh{S}=\bigoplus_{n\geq0}\sh{L}^{\otimes n}$;
a injeção canônica
\[
  \sh{S} \to \bigoplus_{n\in\bb{Z}} \sh{L}^{\otimes n}
\]
define um morfismo dominante canônico
\[
\label{eq:2.8.4.4.3}
  \Spec\left(
    \bigoplus_{n\in\bb{Z}} \sh{L}^{\otimes n}
  \right)
  \to
  \bb{V}(\sh{L})
  =
  \Spec\left(
    \bigoplus_{n\geq0} \sh{L}^{\otimes n}
  \right)
\tag{8.4.4.3}
\]
e basta provar que esse morfismo é um isomorfismo do esquema $\Spec(\bigoplus_{n\in\bb{Z}} \sh{L}^{\otimes n})$ sobre $E$.
Como a questão é local em $Y$, podemos supor que $Y=\Spec(A)$ é afim
\oldpage[II]{169}
e que $\sh{L}=\sh{O}_Y$, de modo que $\sh{S}=(A[T])\supertilde$ e $\bigoplus_{n\in\bb{Z}}\sh{L}^{\otimes n} = (A[T,T^{-1}])\supertilde$.
Mas $A[T,T^{-1}]$ é o anel de frações $A[T]_T$ de $A[T]$, e portanto \eref{eq:2.8.4.4.3} identifica $\Spec(\bigoplus_{n\in\bb{Z}}\sh{L}^{\otimes n})$ com o pré-esquema induzido por $C=\bb{V}(\sh{L})$ no aberto $D(T)$;
o complementar $V(T)$ desse aberto em $C$ é o espaço subjacente ao subpré-esquema fechado de $C$ definido pelo ideal $TA[T]$, que é precisamente a seção nula de $C$, de onde $E=D(T)$.

Por outro lado, o isomorfismo \eref{eq:2.8.4.4.2} será consequência da última afirmação, pois $\bb{V}(\sh{L}^{-1})$ é o complementar do lugar no infinito de seu fecho projetivo, e $\widehat{E}$ é o complementar da seção nula do fecho projetivo de $C=\bb{V}(\sh{L})$.
Esses fechos projetivos são, respectivamente, $\bb{P}(\sh{L}^{-1}\oplus\sh{O}_Y)$ e $\bb{P}(\sh{L}\oplus\sh{O}_Y)$;
mas podemos escrever $\sh{L}\oplus\sh{O}_Y = \sh{L}\otimes(\sh{L}^{-1}\oplus\sh{O}_Y)$.
A existência do isomorfismo canônico procurado decorre então de \sref{II.4.1.4}, e tudo se reduz a verificar que esse isomorfismo troca as seções nulas e os lugares no infinito.
Para isso, podemos nos reduzir ao caso em que $Y=\Spec(A)$ é afim, $L=Ac$ e $L^{-1}=Ac'$, com o isomorfismo canônico $L\otimes L^{-1}\to A$ levando $c\otimes c'$ ao elemento $1$ de $A$.
Então $\bb{S}(L\oplus A)$ é o produto tensorial de $A[\bb{z}]$ com $\bigoplus_{n\geq0}Ac^{\otimes n}$, e $\bb{S}(L^{-1}\oplus A)$ é o produto tensorial de $A[\bb{z}]$ com $\bigoplus_{n\geq0}Ac^{'\otimes n}$, e o isomorfismo definido em \sref{II.4.1.4} leva $\bb{z}^h\otimes c^{'\otimes(n-h)}$ ao elemento $\bb{z}^{n-h}\otimes c^{\otimes h}$.
Ora, em $\bb{P}(\sh{L}^{-1}\oplus\sh{O}_Y)$, o lugar no infinito é o conjunto dos pontos onde a seção $\bb{z}$ se anula, e a seção nula é o conjunto dos pontos onde a seção $c'$ se anula;
como temos definições análogas para $\bb{P}(\sh{L}\oplus\sh{O}_Y)$, a conclusão decorre imediatamente da descrição anterior.
\end{proof}
\subsection{Comportamento funtorial}
\label{subsection:II.8.5}

\begin{env}[8.5.1]
\label{II.8.5.1}
Sejam $Y$ e $Y'$ pré-esquemas, $q:Y'\to Y$ um morfismo, e $\sh{S}$ (resp. $\sh{S}'$) uma $\sh{O}_Y$-álgebra quase-coerente graduada em graus \emph{positivos} (resp. uma $\sh{O}_{Y'}$-álgebra quase-coerente graduada em graus \emph{positivos}).
Consideremos um $q$-morfismo de álgebras graduadas
\[
\label{eq:2.8.5.1.1}
  \vphi: \sh{S} \to \sh{S}'.
\tag{8.5.1.1}
\]

Sabemos \sref{II.1.5.6} que a ele corresponde canonicamente um morfismo
\[
  \Phi = \Spec(\vphi): \Spec(\sh{S}') \to \Spec(\sh{S})
\]
tal que o diagrama
\[
\label{eq:2.8.5.1.2}
  \xymatrix{
    C'
      \ar[r]^{\Phi}
      \ar[d]
  & C
      \ar[d]
  \\Y'
      \ar[r]_{q}
  & Y
  }
\tag{8.5.1.2}
\]
é comutativo, onde escrevemos $C=\Spec(\sh{S})$ e $C'=\Spec(\sh{S}')$.
\emph{Suponhamos, além disso, que $\sh{S}_0=\sh{O}_Y$ e $\sh{S}'_0=\sh{O}_{Y'}$};
sejam $\varepsilon:Y\to C$ e $\varepsilon':Y'\to C'$ as imersões canônicas \sref{II.8.3.2};
temos então o diagrama comutativo
\[
\label{eq:2.8.5.1.3}
  \xymatrix{
    Y'
      \ar[r]^{q}
      \ar[d]_{\varepsilon'}
  & Y
      \ar[d]^{\varepsilon}
  \\C'
      \ar[r]_{\Phi}
  & C
  }
\tag{8.5.1.3}
\]
\oldpage[II]{170}
que corresponde ao diagrama
\[
  \xymatrix{
    \sh{S}
      \ar[r]^{\vphi}
      \ar[d]
  & \sh{S}'
      \ar[d]
  \\\sh{O}_Y
      \ar[r]
  & \sh{O}_{Y'}
  }
\]
onde as setas verticais são os homomorfismos de aumentação; a comutatividade decorre, portanto, da hipótese de que $\vphi$ é um homomorfismo de álgebras \emph{graduadas}.
\end{env}
\begin{proposition}[8.5.2]
\label{II.8.5.2}
Se $E$ (resp. $E'$) é o cone afim sem vértice definido por $\sh{S}$ (resp. $\sh{S}'$), então $\Phi^{-1}(E)\subset E'$;
se, além disso, $\Proj(\vphi):G(\vphi)\to\Proj(\sh{S})$ está definido em toda parte (ou, equivalentemente, se $G(\vphi)=\Proj(\sh{S}')$), então $\Phi^{-1}(E)=E'$, e reciprocamente.
\end{proposition}

\begin{proof}
A primeira afirmação decorre da comutatividade de \eref{eq:2.8.5.1.3}.
Para provar a segunda, podemos nos reduzir ao caso em que $Y=\Spec(A)$ e $Y'=\Spec(A')$ são afins, $\sh{S}=\widetilde{S}$ e $\sh{S}'=\widetilde{S'}$.
Para todo $f$ homogêneo de $S_+$, escrevendo $f'=\vphi(f)$, temos $\Phi^{-1}(C_f)=C'_{f'}$ \sref[I]{I.2.2.4.1};
dizer que $G(\vphi)=\Proj(S')$ significa que o radical (em $S'_+$) do ideal gerado pelos $f'=\vphi(f)$ é o próprio $S'_+$ (\sref{II.2.8.1} e \sref{II.2.3.14}), o que equivale a dizer que os $C'_{f'}$ recobrem $E'$ \sref{II.8.3.4.4}.
\end{proof}
\begin{env}[8.5.3]
\label{II.8.5.3}
O $q$-morfismo $\vphi$ prolonga-se canonicamente a um $q$-morfismo de álgebras graduadas
\[
\label{II.8.5.3.1}
  \widehat{\vphi}: \widehat{\sh{S}} \to \widehat{\sh{S}'}
\tag{8.5.3.1}
\]
pondo $\widehat{\vphi}(\bb{z})=\bb{z}$.
Daí se deduz um morfismo
\[
  \widehat{\Phi} = \Proj(\widehat{\vphi}) : G(\widehat{\vphi}) \to \widehat{C} = \Proj(\widehat{\sh{S}})
\]
tal que o diagrama
\[
  \xymatrix{
    G(\widehat{\vphi})
      \ar[r]^{\widehat{\Phi}}
      \ar[d]
  & \widehat{C}
      \ar[d]
  \\Y'
      \ar[r]_{q}
  & Y
  }
\]
é comutativo \sref{II.3.5.6}.
Resulta imediatamente das definições que, se denotarmos por $i:C\to\widehat{C}$ e $i':C'\to\widehat{C'}$ as imersões abertas canônicas \sref{II.8.3.2}, então $i'(C')\subset G(\widehat{\vphi})$, e o diagrama
\[
\label{eq:2.8.5.3.2}
  \xymatrix{
    C'
      \ar[r]^{\Phi}
      \ar[d]_{i'}
  & C
      \ar[d]^{i}
  \\G(\widehat{\vphi})
      \ar[r]_{\widehat{\Phi}}
  & \widehat{C}
  }
\tag{8.5.3.2}
\]
é comutativo.
Finalmente, se pusermos $X=\Proj(\sh{S})$ e $X'=\Proj(\sh{S}')$, e se $j:X\to\widehat{C}$ e $j':X'\to\widehat{C'}$ forem as imersões fechadas canônicas \sref{II.8.3.2}, resulta ainda da definição dessas imersões que $j'(G(\vphi))\subset G(\widehat{\vphi})$, e que o diagrama
\oldpage[II]{171}
\[
\label{II.8.5.3.3}
  \xymatrix{
    G(\vphi)
      \ar[r]^{\Proj(\vphi)}
      \ar[d]_{j'}
  & X
      \ar[d]^{j}
  \\G(\widehat{\vphi})
      \ar[r]_{\widehat{\Phi}}
  & \widehat{C}
  }
\tag{8.5.3.3}
\]
é comutativo.
\end{env}
\begin{proposition}[8.5.4]
\label{II.8.5.4}
Se $\widehat{E}$ (resp. $\widehat{E'}$) é o cone projetivo sem vértice definido por $\sh{S}$ (resp. por $\sh{S}'$), então $\widehat{\Phi}^{-1}(\widehat{E}) \subset \widehat{E'}$;
além disso, se $p:\widehat{E}\to X$ e $p':\widehat{E'}\to X'$ são as retrações canônicas, então $p'(\widehat{\Phi}^{-1}(\widehat{E})) \subset G(\vphi)$, e o diagrama
\[
\label{II.8.5.4.1}
  \xymatrix{
    \widehat{\Phi}^{-1}(\widehat{E})
      \ar[r]^{\widehat{\Phi}}
      \ar[d]_{p'}
  & \widehat{E}
      \ar[d]^{p}
  \\G(\vphi)
      \ar[r]_{\Proj(\vphi)}
  & X
  }
\tag{8.5.4.1}
\]
é comutativo.
Se $\Proj(\vphi)$ está definido em toda parte, o mesmo vale para $\widehat{\Phi}$, e temos $\widehat{\Phi}^{-1}(\widehat{E}) = \widehat{E'}$.
\end{proposition}

\begin{proof}
A primeira afirmação decorre da comutatividade dos diagramas \eref{eq:2.8.5.1.3} e \eref{eq:2.8.5.3.2}, e as duas afirmações seguintes, da definição das retrações canônicas \sref{II.8.3.5} e da definição de $\widehat{\vphi}$.
Para ver que $\widehat{\Phi}$ está definido em toda parte sempre que $\Proj(\vphi)$ o está, podemos nos limitar ao caso em que $Y=\Spec(A)$ e $Y'=\Spec(A')$ são afins, e em que $\sh{S}=\widetilde{S}$ e $\sh{S}'=\widetilde{S'}$;
a hipótese é que, quando $f$ percorre o conjunto dos elementos homogêneos de $S_+$, o radical em $S'_+$ do ideal gerado em $S'_+$ pelos $\vphi(f)$ é o próprio $S'_+$;
concluímos imediatamente que o radical em $(S'[\bb{z}])_+$ do ideal gerado por $\bb{z}$ e pelos $\vphi(f)$ é o próprio $(S'[\bb{z}])_+$, donde resulta nossa afirmação;
isso mostra ao mesmo tempo que $\widehat{E'}$ é a união dos $\widehat{C'}_{\vphi(f)}$, e, portanto, é igual a $\widehat{\Phi}^{-1}(\widehat{E})$.
\end{proof}
\begin{corollary}[8.5.5]
\label{II.8.5.5}
Sempre que $\Proj(\vphi)$ está definido em toda parte, a imagem inversa por $\widehat{\Phi}$ do espaço subjacente ao lugar no infinito (resp. ao pré-esquema dos vértices) de $\widehat{C}$ é o espaço subjacente ao lugar no infinito (resp. ao pré-esquema dos vértices) de $\widehat{C'}$.
\end{corollary}

\begin{proof}
Isto decorre imediatamente de \sref{II.8.5.4} e \sref{II.8.5.2}, levando em conta as relações \sref{II.8.3.4.1} e \sref{II.8.3.4.2}.
\end{proof}
\subsection{Um isomorfismo canônico para os cones sem vértice}
\label{subsection:II.8.6}

\begin{env}[8.6.1]
\label{II.8.6.1}
Sejam $Y$ um pré-esquema e $\sh{S}$ uma $\sh{O}_Y$-álgebra quase-coerente graduada em graus positivos \emph{tal que $\sh{S}_0=\sh{O}_Y$}, e seja $X$ o $Y$-esquema $\Proj(\sh{S})$.
Vamos aplicar os resultados de \sref{subsection:II.8.5} ao caso em que $Y'=X$ e $q:X\to Y$ é o morfismo estrutural;
seja
\[
\label{II.8.6.1.1}
  \sh{S}_X=\bigoplus_{n\in\bb{Z}}\sh{O}_X(n)
  \tag{8.6.1.1}
\]
\oldpage[II]{172}
que é uma $\sh{O}_X$-álgebra graduada quase-coerente, cuja multiplicação é definida por meio dos homomorfismos canônicos \sref{II.3.2.6.1}
\[
  \sh{O}_X(m)\otimes_{\sh{O}_X}\sh{O}_X(n)\to\sh{O}_X(m+n)
\]
cuja associatividade é assegurada pelo diagrama comutativo de \sref{II.2.5.11.4}.
Seja $\sh{S}'$ a $\sh{O}_X$-subálgebra quase-coerente graduada em graus positivos $\sh{S}_X^\geq=\bigoplus_{n\geq0}\sh{O}_X(n)$ de $\sh{S}_X$.

Finalmente, consideremos o $q$-morfismo canônico
\[
\label{II.8.6.1.2}
  \alpha: \sh{S} \to \sh{S}_X^\geq
\tag{8.6.1.2}
\]
definido em \sref{II.3.3.2.3} como um homomorfismo $\sh{S}\to q_*(\sh{S}_X)$, mas que evidentemente envia $\sh{S}$ a $q_*(\sh{S}_X^\geq)$.
Ponhamos
\[
\label{II.8.6.1.3}
  C_X = \Spec(\sh{S}_X^\geq),
  \quad
  \widehat{C}_X = \Proj(\sh{S}_X^\geq[\bb{z}]),
  \quad
  X' = \Proj(\sh{S}_X^\geq)
\tag{8.6.1.3}
\]
e denotemos por $E_X$ e $\widehat{E}_X$ o cone afim sem vértice e o cone projetivo sem vértice correspondentes (respectivamente);
denotemos por $\varepsilon_X:X\to C_X$, $i_X:C_X\to\widehat{C}_X$, $j_X:X'\to\widehat{C}_X$, $p_X:\widehat{E}_X\to X'$ e $\pi_X:E_X\to X'$ os morfismos canônicos definidos em \sref{subsection:II.8.3}.
\end{env}
\begin{proposition}[8.6.2]
\label{II.8.6.2}
O morfismo estrutural $u:X'\to X$ é um \emph{isomorfismo}, e o morfismo $\Proj(\alpha)$ está definido em toda parte e é idêntico a $u$.
O morfismo $\Proj(\widehat{\alpha}):\widehat{C}_X\to\widehat{C}$ está definido em toda parte, e suas restrições a $\widehat{E}_X$ e $E_X$ são \emph{isomorfismos} sobre $\widehat{E}$ e $E$, respectivamente.
Finalmente, se identificarmos $X'$ com $X$ por meio de $u$, os morfismos $p_X$ e $\pi_X$ identificam-se com os morfismos estruturais dos $X$-pré-esquemas $\widehat{E}_X$ e $E_X$.
\end{proposition}

\begin{proof}
Podemos evidentemente nos limitar ao caso em que $Y=\Spec(A)$ é afim e $\sh{S}=\widetilde{S}$;
então $X$ é a união dos abertos afins $X_f$, onde $f$ percorre o conjunto dos elementos homogêneos de $S_+$, e o anel de cada $X_f$ é $S_{(f)}$.
Resulta de \sref{II.8.2.7.1} que
\[
\label{II.8.6.2.1}
  \Gamma(X_f, \sh{S}_X^\geq) = S_f^\geq.
\tag{8.6.2.1}
\]

Portanto, $u^{-1}(X_f)=\Proj(S_f^\geq)$.
Mas, se $f\in S_d$ ($d>0$), então $\Proj(S_f^\geq)$ é canonicamente isomorfo a $\Proj((S_f^\geq)^{(d)})$ \sref{II.2.4.7}, e sabemos também que $(S_f^\geq)^{(d)}=(S^{(d)})_f^\geq$ pode ser identificado com $S_{(f)}[T]$ \sref{II.2.2.1} pela aplicação $T\mapsto f/1$;
concluímos então \sref{II.3.1.7} que o morfismo estrutural $u^{-1}(X_f)\to X_f$ é um isomorfismo, donde resulta a primeira afirmação.
Para demonstrar a segunda, observemos que a restrição $u^{-1}(X_f)\cap G(\alpha)\to X=\Proj(S)$ de $\Proj(\alpha)$ corresponde à aplicação canônica $x\mapsto x/1$ de $S$ em $S_f^\geq$ \sref{II.2.6.2};
deduzimos, antes de tudo, que $G(\alpha)=X'$, e depois, levando em conta que $u^{-1}(X_f)=(u^{-1}(X_f))_{f/1}$, resulta de \sref{II.2.8.1.1} que a imagem de $u^{-1}(X_f)$ por $\Proj(\alpha)$ está contida em $X_f$, e que a restrição de $\Proj(\alpha)$ a $u^{-1}(X_f)$, considerada como morfismo para $X_f=\Spec(S_{(f)})$, é de fato idêntica à de $u$.
Finalmente, aplicando \sref{II.8.3.5.4} a $p_X$ em vez de a $p$, vemos que $p_X^{-1}(u^{-1}(X_f)) = \Spec((S_f^\geq)_{f/1}^\leq)$, e esse aberto é, por \sref{II.8.5.4.1}, a imagem inversa por $\Proj(\widehat{\alpha})$ de $p^{-1}(X_f)=\Spec(S_f^\leq)$ \sref{II.8.3.5.3}.
Levando em conta \sref{II.8.2.3.2}, a restrição de $\Proj(\widehat{\alpha})$ a $p_X^{-1}(u^{-1}(X_f))$ corresponde ao isomorfismo inverso de \sref{II.8.2.7.2}, restrito a $S_f^\leq$, donde resulta a terceira afirmação;
a última afirmação é evidente por definição.

\oldpage[II]{173}
Observemos também que resulta do diagrama comutativo \eref{eq:2.8.5.3.2} que \emph{a restrição a $C_X$ de $\Proj(\widehat{\alpha})$ não é outra senão o morfismo $\Spec(\alpha)$}.
\end{proof}
\begin{corollary}[8.6.3]
\label{II.8.6.3}
Considerados como $X$-esquemas, $\widehat{E}_X$ é canonicamente isomorfo a $\Spec(\sh{S}_X^\leq)$, e $E_X$ a $\Spec(\sh{S}_X)$.
\end{corollary}

\begin{proof}
Como sabemos que os morfismos $p_X$ e $\pi_X$ são afins (\sref{II.8.3.5} e \sref{II.8.3.6}), basta (em virtude de \sref{II.1.3.1}) demonstrar o corolário no caso em que $Y=\Spec(A)$ é afim e $\sh{S}=\widetilde{S}$.
A primeira afirmação resulta da existência dos isomorfismos canônicos \sref{II.8.2.7.2} $(S_f^\geq)_{f/1}^\leq \simto S_f^\leq$ e do fato de que esses isomorfismos são compatíveis com a passagem de $f$ a $fg$ (onde $f$ e $g$ são homogêneos em $S_+$).
Do mesmo modo, aplicando \sref{II.8.3.6.2} a $\pi_X$ em vez de a $\pi$, vemos que $\pi_X^{-1}(u^{-1}(X_f)) = \Spec((S_f^\geq)_{f/1})$ para $f$ homogêneo em $S_+$, e a segunda afirmação resulta então da existência dos isomorfismos canônicos \sref{II.8.2.7.2} $(S_f^\geq)_{f/1} \simto S_f$.

Pode-se então dizer que $\widehat{C}_X$, considerado como um $X$-esquema, é obtido por \emph{colagem} dos $X$-esquemas afins $C_X=\Spec(\sh{S}_X^\geq)$ e $\widehat{E}_X=\Spec(\sh{S}_X^\leq)$ sobre $X$, onde a interseção dos dois $X$-esquemas afins é o aberto $E_X=\Spec(\sh{S}_X)$.
\end{proof}
\begin{corollary}[8.6.4]
\label{II.8.6.4}
Suponhamos que $\sh{O}_X(1)$ seja um $\sh{O}_X$-módulo invertível e que $\sh{S}_X$ seja isomorfa a $\bigoplus_{n\in\bb{Z}}(\sh{O}_X(1))^{\otimes n}$ (o que ocorrerá, em particular, sempre que $\sh{S}$ for gerada por $\sh{S}_1$ (\sref{II.3.2.5} e \sref{II.3.2.7})).
Então o cone projetivo sem vértice $\widehat{E}$ identifica-se com o fibrado vetorial de posto 1 $\bb{V}(\sh{O}_X(-1))$ sobre $X$, e o cone afim sem vértice $E$ com o subpré-esquema desse fibrado vetorial induzido sobre o complementar da seção nula.
Com essa identificação, a retração canônica $\widehat{E}\to X$ identifica-se com o morfismo estrutural do $X$-esquema $\bb{V}(\sh{O}_X(-1))$.
Finalmente, existe um $Y$-morfismo canônico $\bb{V}(\sh{O}_X(1))\to C$, cuja restrição ao complementar da seção nula de $\bb{V}(\sh{O}_X(1))$ é um isomorfismo desse complementar sobre o cone afim sem vértice $E$.
\end{corollary}

\begin{proof}
Se pusermos $\sh{L}=\sh{O}_X(1)$, então $\sh{S}_X^\geq$ é idêntica a $\bb{S}_{\sh{O}_X}(\sh{L})$, de modo que $\widehat{E}_X$ se identifica canonicamente com $\bb{V}(\sh{L}^{-1})$, em virtude de \sref{II.8.6.3}, e $C_X$ com $\bb{V}(\sh{L})$.
O morfismo $\bb{V}(\sh{L})\to C$ é a restrição de $\Proj(\widehat{\alpha})$, e as afirmações do corolário são então casos particulares de \sref{II.8.6.2}.
\end{proof}

Observamos que a imagem inversa pelo morfismo $\bb{V}(\sh{O}_X(1))\to C$ do espaço subjacente ao pré-esquema dos vértices de $C$ é o espaço subjacente à seção nula de $\bb{V}(\sh{O}_X(1))$ \sref{II.8.5.5};
mas, em geral, os subpré-esquemas correspondentes de $C$ e de $\bb{V}(\sh{O}_X(1))$ não são isomorfos.
Essa questão será estudada abaixo.
\subsection{Explosão de cones projetantes}
\label{subsection:II.8.7}

\begin{env}[8.7.1]
\label{II.8.7.1}
Nas condições de \sref{II.8.6.1}, temos, pondo $r=\Proj(\widehat{\alpha})$, um diagrama comutativo
\[
\label{II.8.7.1.1}
  \xymatrix{
    X
      \ar[r]^{i_X\circ\varepsilon_X}
      \ar[d]_{q}
  & \widehat{C}_X
      \ar[d]^{r}
  \\Y
      \ar[r]_{i\circ\varepsilon}
  & \widehat{C}
  }
\tag{8.7.1.1}
\]
\oldpage[II]{174}
em virtude de \eref{eq:2.8.5.1.3} e \eref{eq:2.8.5.3.2};
além disso, a restrição de $r$ ao complementar $\widehat{C}_X\setmin i_X(\varepsilon_X(X))$ da seção nula é um \emph{isomorfismo} sobre o complementar $\widehat{C}\setmin i(\varepsilon(Y))$ da seção nula, em virtude de \sref{II.8.6.2}.
Se supusermos, para simplificar, que $Y$ é afim e que $\sh{S}$ é de tipo finito e gerada por $\sh{S}_1$, então $X$ é projetivo sobre $Y$ e $\widehat{C}_X$ é projetivo sobre $X$ \sref{II.5.5.1}; logo, $\widehat{C}_X$ é projetivo sobre $Y$ \sref{II.5.5.5}[ii], e \emph{a fortiori} sobre $\widehat{C}$ \sref{II.5.5.5}[v].
Temos assim um $Y$-morfismo projetivo $r:\widehat{C}_X\to\widehat{C}$ (cuja restrição a $C_X$ é um $Y$-morfismo projetivo $C_X\to C$) que \emph{contrai $X$ sobre $Y$} e induz um \emph{isomorfismo} quando restringido aos \emph{complementares de $X$ e de $Y$}.
Temos, portanto, uma relação entre $C_X$ e $C$, análoga à que existe entre um pré-esquema obtido por explosão e seu pré-esquema de partida \sref{II.8.1.3}.
Mostraremos, com efeito, que $C_X$ pode ser identificado com o espectro homogêneo de uma $\sh{O}_C$-álgebra graduada.
\end{env}
\begin{env}[8.7.2]
\label{II.8.7.2}
Com as notações de \sref{II.8.6.1}, consideremos, para todo $n\geq0$, o ideal quase-coerente
\[
\label{II.8.7.2.1}
  \sh{S}_{[n]} = \bigoplus_{m\geq n}\sh{S}_m
\tag{8.7.2.1}
\]
da $\sh{O}_Y$-álgebra graduada $\sh{S}$.
É claro que
\[
\label{II.8.7.2.2}
  \sh{S}_{[0]} = \sh{S},
  \qquad
  \sh{S}_{[n]} \subset \sh{S}_{[m]}
  \qquad\qquad
  \mbox{para $m\leq n$}
\tag{8.7.2.2}
\]
\[
\label{II.8.7.2.3}
  \sh{S}_{[m]} \sh{S}_{[n]} \subset \sh{S}_{[m+n]}.
\tag{8.7.2.3}
\]

Consideremos o $\sh{O}_C$-módulo associado a $\sh{S}_{[n]}$, que é um ideal quase-coerente de $\sh{O}_C = \widetilde{\sh{S}}$ \sref{II.1.4.4}
\[
\label{II.8.7.2.4}
  \sh{I}_n = (\sh{S}_{[n]})\supertilde.
\tag{8.7.2.4}
\]

Deduzimos, portanto, de \sref{II.8.7.2.2} e \sref{II.8.7.2.3}, utilizando \sref{II.1.4.4} e \sref{II.1.4.8.1}, as fórmulas análogas
\[
\label{II.8.7.2.5}
  \sh{I}_{0} = \sh{O}_C,
  \qquad
  \sh{I}_{n} \subset \sh{I}_{m}
  \qquad
  \text{para }m\leq n
\tag{8.7.2.5}
\]
\[
\label{II.8.7.2.6}
  \sh{I}_{m} \sh{I}_{n} \subset \sh{I}_{m+n}.
\tag{8.7.2.6}
\]

Estamos, assim, nas condições de \sref{II.8.1.1}, o que nos leva a introduzir a $\sh{O}_C$-álgebra graduada quase-coerente
\[
\label{II.8.7.2.7}
  \sh{S}^\natural
  =
  \bigoplus_{n\geq0}\sh{I}_n
  =
  \left(
    \bigoplus_{n\geq0} \sh{S}_{[n]}
  \right)\supertilde.
\tag{8.7.2.7}
\]
\end{env}
\begin{proposition}[8.7.3]
\label{II.8.7.3}
Existe um $C$-isomorfismo canônico
\[
\label{II.8.7.3.1}
  h: C_X \simto \Proj(\sh{S}^\natural).
\tag{8.7.3.1}
\]
\end{proposition}

\begin{proof}
Suponhamos, antes de tudo, que $Y=\Spec(A)$ seja afim, de modo que $\sh{S}=\widetilde{S}$, com $S$ uma $A$-álgebra graduada em graus positivos, e $C=\Spec(S)$.
A definição~\sref{II.8.2.6.2} mostra então, com as notações de \sref{II.8.2.6}, que $\sh{S}^\natural=(S^\natural)\supertilde$.
Para definir \sref{II.8.7.3.1}, consideremos um elemento homogêneo $f\in S_d$ ($d>0$) e o elemento correspondente $f^\natural\in S^\natural$ \sref{II.8.2.6};
o $S$-isomorfismo de \sref{II.8.2.7.3} define então um $C$-isomorfismo
\[
\label{II.8.7.3.2}
  \Spec(S_f^\geq) \simto \Spec(S_{(f^\natural)}^\natural).
\tag{8.7.3.2}
\]

\oldpage[II]{175}
Mas, com as notações de \sref{II.8.6.2}, se $v:C_X\to X$ é o morfismo estrutural, resulta de \sref{II.8.6.2.1} que $v^{-1}(X_f)=\Spec(S_f^\geq)$.
Temos também $\Spec(S_{(f^\natural)}^\natural)=D_+(f^\natural)$, o que significa que \sref{II.8.7.3.2} define um isomorfismo $v^{-1}(X_f)\to D_+(f^\natural)$.
Além disso, se $g\in S_e$ ($e>0$), então o diagrama
\[
  \xymatrix{
    v^{-1}(X_{fg})
      \ar[r]^{\sim}
      \ar[d]
  & D_+(f^\natural g^\natural)
      \ar[d]
  \\v^{-1}(X_f)
      \ar[r]^{\sim}
  & D_+(f^\natural)
  }
\]
é comutativo, pela definição do isomorfismo de \sref{II.8.2.7.3}.
Finalmente, por definição, $S_+$ é gerado pelos elementos homogêneos $f$, e decorre então de \sref{II.8.2.10}[iv] e de \sref{II.2.3.14} que os $D_+(f^\natural)$ formam um recobrimento de $\Proj(S^\natural)$, e que os $v^{-1}(X_f)$ formam um recobrimento de $C_X$, pois os $X_f$ formam um recobrimento de $X$;
neste caso, definimos assim o isomorfismo \sref{II.8.7.3.1}.

Para demonstrar \sref{II.8.7.3} no caso geral, basta mostrar que, se $U$ e $U'$ são abertos afins de $Y$, correspondentes aos anéis $A$ e $A'$ (respectivamente), e tais que $U'\subset U$, então, pondo $\sh{S}|U=\widetilde{S}$ e $\sh{S}|U'=\widetilde{S'}$, o diagrama
\[
\label{II.8.7.3.3}
  \xymatrix{
    C_{U'}
      \ar[r]
      \ar[d]
  & \Proj(S^{'\natural})
      \ar[d]
  \\C_U
      \ar[r]
  & \Proj(S^\natural)
  }
\tag{8.7.3.3}
\]
é comutativo.
Mas $S'$ identifica-se canonicamente com $S\otimes_A A'$, e, portanto, $S^{'\natural}$ identifica-se canonicamente com
\[
  S^\natural \otimes_S S' = S^\natural \otimes_A A';
\]
assim, $\Proj(S^{'\natural}) = \Proj(S^\natural)\times_U U'$ \sref{II.2.8.10};
do mesmo modo, se $X=\Proj(S)$ e $X'=\Proj(S')$, então $X'=X\times_U U'$ e $\sh{S}_{X'}=\sh{S}_X\otimes_{\sh{O}_U}\sh{O}_{U'}$ \sref{II.3.5.4}, ou, o que equivale, $\sh{S}_{X'}=j^*(\sh{S}_X)$, onde $j$ é a projeção $X'\to X$.
Temos então, por \sref{II.1.5.2}, que $C_{U'} = C_U\times_X X' = C_U\times_U U'$, e a comutatividade de \sref{II.8.7.3.3} é imediata.
\end{proof}
\begin{remark}[8.7.4]
\label{II.8.7.4}
\begin{enumerate}
  \item[\rm{(i)}] O final da demonstração de \sref{II.8.7.3} generaliza-se imediatamente da maneira seguinte.
    Seja $g:Y'\to Y$ um morfismo, $\sh{S}'=g^*(\sh{S})$, e $X'=\Proj(\sh{S}')$;
    temos então um diagrama comutativo
    \[
    \label{II.8.7.4.1}
      \xymatrix{
        C_{X'}
          \ar[r]
          \ar[d]
      & \Proj(\sh{S}^{'\natural})
          \ar[d]
      \\C_X
          \ar[r]
      & \Proj(\sh{S}^\natural)
      }
    \tag{8.7.4.1}
    \]

    Seja agora $\vphi:\sh{S}''\to\sh{S}$ um homomorfismo de $\sh{O}_Y$-álgebras graduadas tal que, pondo $X''=\Proj(\sh{S}'')$, o morfismo $u=\Proj(\vphi):X\to X''$ esteja definido em toda parte;
    temos também
\oldpage[II]{176}
    um $Y$-morfismo $v:C\to C''$ (com $C''=\Spec(\sh{S}'')$) tal que $\sh{A}(v)=\vphi$, e, como $\vphi$ é um homomorfismo de álgebras graduadas, $\vphi$ induz um $v$-morfismo de álgebras graduadas $\psi:\sh{S}^{\prime\prime\natural}\to\sh{S}^\natural$ \sref{II.1.4.1}.
    Além disso, resulta de \sref{II.8.2.10}[iv] e da hipótese sobre $\vphi$ que $\Proj(\psi)$ está definido em toda parte.
    Finalmente, levando em conta \sref{II.3.5.6.1}, existe um $u$-morfismo canônico $\sh{S}_{X''}\to\sh{S}_X$, donde se obtém, por \sref{II.1.5.6}, um morfismo $w:C_X\to C_{X''}$.
    Nessas condições, o diagrama
    \[
    \label{II.8.7.4.2}
      \xymatrix{
        C_X
          \ar[r]^-{\sim}
          \ar[d]_{w}
      & \Proj(\sh{S}^\natural)
          \ar[d]^{\Proj(\psi)}
      \\C_{X''}
          \ar[r]^-{\sim}
      & \Proj(\sh{S}^{\prime\prime\natural})
      }
    \tag{8.7.4.2}
    \]
    é comutativo, como se verifica imediatamente reduzindo ao caso em que $Y$ é afim.
  \item[\rm{(ii)}] Observemos que, por \sref{II.8.7.2.5} e \sref{II.8.7.2.6}, temos $\sh{I}_1^m\subset\sh{I}_m\subset\sh{I}_1$ para todo $m>0$.
    Mas, por definição, $\sh{I}_1=(\sh{S}_+)\supertilde$, e, portanto, $\sh{I}_1$ define o subpré-esquema fechado $\varepsilon(Y)$ em $C$ (\sref{II.1.4.10} e \sref{II.8.3.2});
    concluímos assim que, para todo $m>0$, \emph{o suporte de $\sh{O}_C/\sh{I}_m$ está contido no espaço subjacente ao pré-esquema dos vértices $\varepsilon(Y)$};
    sobre a imagem inversa do cone afim sem vértice $E$, o morfismo estrutural $\Proj(\sh{S}^\natural)\to C$ restringe-se, portanto, a um \emph{isomorfismo} (por \sref{II.8.7.3} e \sref{II.8.7.1}).
    Além disso, identificando canonicamente $C$ com um aberto de $\widehat{C}$ \sref{II.8.3.3}, podemos evidentemente prolongar os ideais $\sh{I}_m$ de $\sh{O}_C$ a ideais $\sh{J}_m$ de $\sh{O}_{\widehat{C}}$, exigindo que coincidam com $\sh{O}_{\widehat{C}}$ no aberto $\widehat{E}$ de $\widehat{C}$.
    Se definirmos $\sh{T}=\bigoplus_{n\geq 0}\sh{J}_n$, que é uma $\sh{O}_{\widehat{C}}$-álgebra graduada quase-coerente, podemos prolongar o isomorfismo \sref{II.8.7.3.1} a um $\widehat{C}$-isomorfismo
    \[
    \label{II.8.7.4.3}
      \widehat{C}_X \simto \Proj(\sh{T}).
    \tag{8.7.4.3}
    \]

    Com efeito, sobre $\widehat{E}$, resulta do que precede que $\Proj(\sh{T})$ se identifica canonicamente com $\widehat{E}$, e definimos então o isomorfismo \sref{II.8.7.4.3} sobre $\widehat{E}$ exigindo que coincida com o isomorfismo canônico $\widehat{E}_X\to\widehat{E}$ \sref{II.8.6.2};
    é claro que esse isomorfismo e \sref{II.8.7.3.1} coincidem então sobre $E$.
\end{enumerate}
\end{remark}
\begin{corollary}[8.7.5]
\label{II.8.7.5}
Suponhamos que exista $n_0>0$ tal que
\[
\label{II.8.7.5.1}
  \sh{S}_{n+1} = \sh{S}_1\sh{S}_n
  \qquad
  \text{para }n\geq n_0.
\tag{8.7.5.1}
\]

Então, o pré-esquema dos vértices de $C_X$ (isomorfo a $X$) é a imagem inversa, pelo morfismo canônico $r:C_X\to C$, do pré-esquema dos vértices de $C$ (isomorfo a $Y$).
Reciprocamente, se essa propriedade se verifica e se, além disso, supomos que $Y$ é noetheriano e que $\sh{S}$ é de tipo finito, então existe $n_0>0$ tal que \sref{II.8.7.5.1} se verifica.
\end{corollary}

\begin{proof}
Como a primeira afirmação é local sobre $Y$, podemos supor que $Y=\Spec(A)$ seja afim, de modo que $\sh{S}=\widetilde{S}$, com $S$ uma $A$-álgebra graduada em graus positivos.
A afirmação decorre então de \sref{II.8.2.12}, pois $\Proj(S^\natural\otimes_S S_0) = C_X\times_C\varepsilon(Y)$ (pela identificação de \sref{II.8.7.3.1}), ou, em outros termos, porque esse pré-esquema é a imagem inversa de $\varepsilon(Y)$ em $C_X$ \sref[I]{I.4.4.1}.
A recíproca também decorre de \sref{II.8.2.12} quando $Y$ é noetheriano e afim e $S$ é de tipo finito.
\oldpage[II]{177}
Se $Y$ é noetheriano (mas não necessariamente afim) e $\sh{S}$ é de tipo finito, então existe uma cobertura finita de $Y$ por abertos afins noetherianos $U_i$, e deduzimos então do que precede que, para todo $i$, existe um inteiro $n_i$ tal que $\sh{S}_{n+1}|U_i = (\sh{S}_1|U_i)(\sh{S}_n|U_i)$ para $n\geq n_i$;
o maior dos $n_i$ garante então que \sref{II.8.7.5.1} se verifica.
\end{proof}
\begin{env}[8.7.6]
\label{II.8.7.6}
Consideremos agora o $C$-pré-esquema $Z$ obtido pela \emph{explosão} do \emph{subpré-esquema vértice $\varepsilon(Y)$} no cone afim $C$;
pela Definição~\sref{II.8.1.3}, ele é precisamente o pré-esquema $\Proj(\bigoplus_{n\geq0}\sh{S}_+^n)$;
a injeção canônica
\[
\label{II.8.7.6.1}
  \iota: \bigoplus_{n\geq0} \sh{S}_+^n \to \sh{S}^\natural
\tag{8.7.6.1}
\]
define (pela identificação de \sref{II.8.7.3}) um $C$-morfismo dominante canônico
\[
\label{II.8.7.6.2}
  G(\iota)\to Z
  \tag{8.7.6.2}
\]
onde $G(\iota)$ é um aberto de $C_X$ \sref{II.3.5.1};
notemos que se pode ter $G(\iota)\neq C_X$, como mostra o exemplo em que $Y=\Spec(K)$, sendo $K$ um corpo, e $\sh{S}=\widetilde{S}$, com $S=K[\bb{y}]$, onde $\bb{y}$ é uma indeterminada \emph{de grau 2};
se $R_n$ denota o conjunto $(S_+)^n$, considerado como um subconjunto de $S_{[n]}=S_n^\natural$, então $S_+^\natural$ não é o radical em $S^\natural$ do ideal gerado pela união dos $R_n$ (cf. \sref{II.2.3.14}).
\end{env}
\begin{corollary}[8.7.7]
\label{II.8.7.7}
Suponhamos que exista $n_0>0$ tal que
\[
\label{II.8.7.7.1}
  \sh{S}_n=\sh{S}_1^n
  \quad
  \text{para }n\geq n_0.
  \tag{8.7.7.1}
\]

Então, o morfismo canônico \sref{II.8.7.6.2} está definido em toda parte e é um isomorfismo $C_X\simto Z$.
Reciprocamente, se essa propriedade se verifica e se, além disso, supomos que $Y$ é noetheriano e que $\sh{S}$ é de tipo finito, então existe $n_0$ tal que \sref{II.8.7.7.1} se verifica.
\end{corollary}

\begin{proof}
A primeira afirmação é local sobre $Y$ e decorre, portanto, de \sref{II.8.2.14};
o mesmo vale para a recíproca, raciocinando como em \sref{II.8.7.5}.
\end{proof}
\begin{remark}[8.7.8]
\label{II.8.7.8}
Como a condição \sref{II.8.7.7.1} implica \sref{II.8.7.5.1}, vemos que, quando ela se verifica, não só $C_X$ se identifica com o pré-esquema obtido pela explosão do vértice (identificado com $Y$) do cone afim $C$, mas também o vértice (identificado com $X$) de $C_X$ se identifica com o subpré-esquema fechado que é a imagem inversa do vértice $Y$ de $C$.
Além disso, a hipótese~\sref{II.8.7.7.1} implica que, sobre $X=\Proj(\sh{S})$, os $\sh{O}_X$-módulos $\sh{O}_X(n)$ são invertíveis (\sref{II.3.2.5} e \sref{II.3.2.9}), e que $\sh{O}_X(n)=\sh{L}^{\otimes n}$, com $\sh{L}=\sh{O}_X(1)$ (\sref{II.3.2.7} e \sref{II.3.2.9});
pela Definição~\sref{II.8.6.1.1}, $C_X$ é, portanto, o \emph{fibrado vetorial} $\bb{V}(\sh{L})$ sobre $X$, e seu vértice é a \emph{seção nula} desse fibrado vetorial.
\end{remark}
\subsection{Feixes amplos e contrações}
\label{subsection:II.8.8}

\begin{env}[8.8.1]
\label{II.8.8.1}
Sejam $Y$ um pré-esquema, $f:X\to Y$ um morfismo \emph{separado} e \emph{quase-compacto}, e $\sh{L}$ um $\sh{O}_X$-módulo invertível \emph{amplo relativamente a $f$}.
Consideremos a $\sh{O}_Y$-álgebra graduada em graus positivos
\[
\label{II.8.8.1.1}
  \sh{S} = \sh{O}_Y \oplus \bigoplus_{n\geq1}f_*(\sh{L}^{\otimes n})
\tag{8.8.1.1}
\]
\oldpage[II]{178}
que é quase-coerente \sref[I]{I.9.2.2}[a].
Há um homomorfismo canônico de $\sh{O}_X$-álgebras graduadas
\[
\label{II.8.8.1.2}
  \tau: f^*(\sh{S}) \to \bigoplus_{n\geq0}\sh{L}^{\otimes n}
\tag{8.8.1.2}
\]
que, nos graus $\geq1$, coincide com o homomorfismo canônico $\sigma:f^*(f_*(\sh{L}^{\otimes n})) \to \sh{L}^{\otimes n}$ \sref[0]{0.4.4.3}, e é a identidade no grau 0.
A hipótese de que $\sh{L}$ é $f$-amplo implica então (\sref{II.4.6.3} e \sref{II.3.6.1}) que o $Y$-morfismo correspondente
\[
\label{II.8.8.1.3}
  r = r_{\sh{L},\tau} : X \to P = \Proj(\sh{S})
\tag{8.8.1.3}
\]
está definido em toda parte e é uma \emph{imersão aberta dominante}, e que
\[
\label{II.8.8.1.4}
  r^*(\sh{O}_P(n)) = \sh{L}^{\otimes n}
  \qquad\qquad
  \mbox{para todo $n\in\bb{Z}$.}
\tag{8.8.1.4}
\]
\end{env}
\begin{proposition}[8.8.2]
\label{II.8.8.2}
Seja $C=\Spec(\sh{S})$ o cone afim definido por $\sh{S}$;
se $\sh{L}$ é $f$-amplo, existe um $Y$-morfismo canônico
\[
\label{II.8.8.2.1}
  g : V = \bb{V}(\sh{L}) \to C
\tag{8.8.2.1}
\]
tal que o diagrama
\[
\label{II.8.8.2.2}
  \xymatrix{
    X
      \ar[r]^{j}
      \ar[d]_{f}
  & \bb{V}(\sh{L})
      \ar[r]^{\pi}
      \ar[d]^{g}
  & X
      \ar[d]^{f}
  \\Y
      \ar[r]_{\varepsilon}
  & C
      \ar[r]_{\psi}
  & Y
  }
  \tag{8.8.2.2}
\]
é comutativo, onde $\psi$ e $\pi$ são os morfismos estruturais, e $j$ e $\varepsilon$ as imersões canônicas que enviam, respectivamente, $X$ e $Y$ à seção nula de $\bb{V}(\sh{L})$ e ao pré-esquema dos vértices de $C$.
Além disso, a restrição de $g$ a $\bb{V}(\sh{L})\setmin j(X)$ é uma imersão aberta
\[
\label{II.8.8.2.3}
  \bb{V}(\sh{L})\setmin j(X)\to E=C\setmin\varepsilon(Y)
  \tag{8.8.2.3}
\]
no cone afim sem vértice $E$ correspondente a $\sh{S}$.
\end{proposition}

\begin{proof}
Com as notações de \sref{II.8.8.1}, sejam $\sh{S}_P^\geq = \bigoplus_{n\geq0}\sh{O}_P(n)$ e $C_P=\Spec(\sh{S}_P^\geq)$.
Sabemos por \sref{II.8.6.2} que existe um morfismo canônico $h=\Spec(\alpha):C_P\to C$ tal que o diagrama
\[
\label{II.8.8.2.4}
  \xymatrix{
    C_P
      \ar[r]
      \ar[d]_{h}
  & P
      \ar[d]^{p}
  \\C
      \ar[r]_{\psi}
  & Y
  }
\tag{8.8.2.4}
\]
é comutativo; além disso, se $\varepsilon_P:P\to C_P$ é a imersão canônica, então o diagrama
\[
\label{II.8.8.2.5}
  \xymatrix{
    P
      \ar[r]^{\varepsilon_P}
      \ar[d]_{p}
  & C_P
      \ar[d]^{h}
  \\Y
      \ar[r]_{\varepsilon}
  & C
  }
\tag{8.8.2.5}
\]
é comutativo \sref{II.8.7.1.1} e, finalmente, a restrição de $h$ ao cone afim sem vértice $E_P$ é um \emph{isomorfismo} $E_P\simto E$ \sref{II.8.6.2}.
Resulta de \sref{II.8.8.1.4} que
\[
  r^*(\sh{S}_P^\geq)=\bb{S}_{\sh{O}_X}(\sh{L})
\]
\oldpage[II]{179}
e, portanto, existe um $P$-morfismo canônico $q:\bb{V}(\sh{L})\to C_P$, com o diagrama comutativo
\[
\label{II.8.8.2.6}
  \xymatrix{
    \bb{V}(\sh{L})
      \ar[r]^{\pi}
      \ar[d]_{q}
  & X
      \ar[d]^{r}
  \\C_P
      \ar[r]
  & P
  }
\tag{8.8.2.6}
\]
que identifica $\bb{V}(\sh{L})$ com o produto $C_P\times_P X$ \sref{II.1.5.2};
como $r$ é uma imersão aberta, $q$ também o é \sref[I]{I.4.3.2}.
Além disso, a restrição de $q$ a $\bb{V}(\sh{L})\setmin j(X)$ envia este pré-esquema a $E_P$, por \sref{II.8.5.2}, e o diagrama
\[
\label{II.8.8.2.7}
  \xymatrix{
    X
      \ar[r]^{j}
      \ar[d]_{r}
  & \bb{V}(\sh{L})
      \ar[d]^{q}
  \\P
      \ar[r]_{\varepsilon_P}
  & C_P
  }
\tag{8.8.2.7}
\]
é comutativo (pois é um caso particular de \eref{eq:2.8.5.1.3}).
As afirmações de \sref{II.8.8.2} decorrem imediatamente desses fatos, tomando para $g$ o morfismo composto $h\circ q$.
\end{proof}
\begin{remark}[8.8.3]
\label{II.8.8.3}
Suponhamos ainda que $Y$ seja um pré-esquema \emph{noetheriano} e que $f$ seja um morfismo \emph{próprio}.
Como $r$ é então \emph{próprio} \sref{II.5.4.4}, portanto fechado, e é também uma imersão aberta dominante, $r$ é necessariamente um \emph{isomorfismo} $X\simto P$.
Veremos ainda, no Capítulo~III \sref[III]{III.2.3.5.1}, que $\sh{S}$ é então necessariamente uma $\sh{O}_Y$-álgebra \emph{de tipo finito}.
Segue-se que $\sh{S}^\natural$ é uma $\sh{S}_0^\natural$-álgebra \emph{de tipo finito} (\sref{II.8.2.10}[i] e \sref{II.8.7.2.7});
como $C_P$ é $C$-isomorfo a $\Proj(\sh{S}^\natural)$ \sref{II.8.7.3}, vemos que o morfismo $h:C_P\to C$ é \emph{projetivo};
como o morfismo $r$ é um isomorfismo, $q:\bb{V}(\sh{L})\to C_P$ também o é, e concluímos que o morfismo $g:\bb{V}(\sh{L})\to C$ é \emph{projetivo}.
Além disso, como a restrição de $h$ a $E_P$ é um isomorfismo sobre $E$ e $q$ é um isomorfismo, a restrição \sref{II.8.8.2.3} de $g$ é um \emph{isomorfismo} $\bb{V}(\sh{L})\setmin j(X)\simto E$.

Se supusermos ainda que $\sh{L}$ é \emph{muito amplo} para $f$, veremos também no Capítulo~III \sref[III]{III.2.3.5.1} que existe um inteiro $n_0>0$ tal que $\sh{S}_n=\sh{S}_1^n$ para $n\geq n_0$.
Concluiremos então, por \sref{II.8.7.7}, que $\bb{V}(\sh{L})$ se identifica com o pré-esquema $Z$ obtido pela explosão do subpré-esquema dos vértices (identificado com $Y$) do cone afim $C$,
e que a \emph{seção nula} de $\bb{V}(\sh{L})$ (identificada com $X$) é a \emph{imagem inversa} do subpré-esquema dos vértices $Y$ de $C$.

Uma parte dos resultados precedentes pode, de fato, ser estabelecida sem a hipótese noetheriana:
\end{remark}
\begin{corollary}[8.8.4]
\label{II.8.8.4}
Sejam $Y$ um pré-esquema (resp. um esquema quase-compacto), $f:X\to Y$ um morfismo próprio e $\sh{L}$ um $\sh{O}_X$-módulo invertível amplo relativamente a $f$.
Então o morfismo de \sref{II.8.8.2.1} é próprio (resp. projetivo), e sua restrição \sref{II.8.8.2.3} é um isomorfismo.
\end{corollary}

\begin{proof}
Para provar que $g$ é próprio, podemos reduzir-nos ao caso em que $Y$ é afim, e basta então considerar o caso em que $Y$ é um esquema quase-compacto.
O mesmo raciocínio que em \sref{II.8.8.3} mostra primeiro que $r$ é um \emph{isomorfismo} $X\simto P$;
portanto, $q$ também é um isomorfismo e, como a restrição de $h$ a $E_P$ é um isomorfismo $E_P\simto E$, já vemos que \sref{II.8.8.2.3} é um isomorfismo.
Resta apenas provar que $g$ é \emph{projetivo}.

Como $f$ é de tipo finito por hipótese, podemos aplicar \sref{II.3.8.5} ao homomorfismo
\oldpage[II]{180}
$\tau$ de \sref{II.8.8.1.2}:
existem um inteiro $d>0$ e um $\sh{O}_Y$-submódulo quase-coerente de tipo finito $\sh{E}$ de $\sh{S}_d$ tais que, se $\sh{S}'$ é a $\sh{O}_Y$-subálgebra de $\sh{S}$ gerada por $\sh{E}$ e $\tau'=\tau\circ f^*(\vphi)$ (onde $\vphi$ é a injeção canônica $\sh{S}'\to\sh{S}$), então $r'=r_{\sh{L},\tau'}$ é uma imersão
\[
  X\to P'=\Proj(\sh{S}').
\]
Além disso, como $\vphi$ é injetiva, $r'$ é também uma \emph{imersão dominante} \sref{II.3.7.6};
o mesmo raciocínio que para $r$ mostra, portanto, que $r'$ é uma \emph{imersão fechada sobrejetiva};
como $r'$ se fatora como $X\xrightarrow{r}\Proj(\sh{S})\xrightarrow{\Phi}\Proj(\sh{S}')$, onde $\Phi=\Proj(\vphi)$, concluímos que $\Phi$ também é uma \emph{imersão fechada sobrejetiva}.
Mas isso implica que $\Phi$ é um \emph{isomorfismo};
com efeito, podemos reduzir-nos ao caso em que $Y=\Spec(A)$ é afim e $\sh{S}=\widetilde{S}$, $\sh{S}'=\widetilde{S'}$, onde $S$ é uma $A$-álgebra graduada e $S'$ uma subálgebra graduada de $S$.
Para todo elemento homogêneo $t\in S'$, $S'_{(t)}$ é então um subanel de $S_{(t)}$;
voltando à definição de $\Proj(\vphi)$ \sref{II.2.8.1}, vemos que basta provar que, se $B'$ é um subanel de um anel $B$ e o morfismo $\Spec(B)\to\Spec(B')$ correspondente à injeção canônica $B'\to B$ é uma imersão fechada, esse morfismo é necessariamente um \emph{isomorfismo};
mas isso decorre de \sref[I]{I.4.2.3}.
Além disso, $\Phi^*(\sh{O}_{P'}(n))=\sh{O}_P(n)$ (\sref{II.3.5.2}[ii] e \sref{II.3.5.4}), e portanto $r^{'*}(\sh{O}_{P'}(n))$ é isomorfo a $\sh{L}^{\otimes n}$ \sref{II.4.6.3}.
Seja $\sh{S}''=\sh{S}^{'(d)}$, de modo que \sref{II.3.1.8}[i] $X$ se identifica canonicamente com $P''=\Proj(\sh{S}'')$, e $\sh{L}''=\sh{L}^{\otimes d}$ com $\sh{O}_{P''}(1)$ \sref{II.3.2.9}[ii].

Agora, se $C''=\Spec(\sh{S}'')$, então $\sh{S}_{P''}^\geq=\bigoplus_{n\geq 0}\sh{O}_{P''}(n)$ se identifica com $\bigoplus_{n\geq 0}\sh{L}^{\prime\prime\otimes n}$, e portanto $C_{P''}=\Spec(\sh{S}_{P''}^\geq)$ com $\bb{V}(\sh{L}'')$;
sabemos ainda por \sref{II.8.7.3} que $C_{P''}$ é $C''$-isomorfo a $\Proj(\sh{S}^{\prime\prime\natural})$;
por definição de $\sh{S}''$, $\sh{S}^{\prime\prime\natural}$ é gerada por $\sh{S}_1^{\prime\prime\natural}$, e $\sh{S}_1^{\prime\prime\natural}$ é de tipo finito sobre $\sh{S}_0^{\prime\prime\natural}=\sh{S}''$ (\sref{II.8.2.10}[i e iii]), de modo que $\Proj(\sh{S}^{\prime\prime\natural})$ é \emph{projetivo} sobre $C''$ \sref{II.5.5.1}.
Consideremos o diagrama
\[
\label{II.8.8.4.1}
  \xymatrix{
    \bb{V}(\sh{L})
      \ar[r]^-{g}
      \ar[d]_{u}
  & \Spec(\sh{S}) = C
      \ar[d]^{v}
  \\\bb{V}(\sh{L}'')
      \ar[r]_{g''}
  & \Spec(\sh{S}'') = C''
  }
  \tag{8.8.4.1}
\]
onde $g$ e $g''$ correspondem, por \sref{II.1.5.6}, aos $f$-morfismos canônicos
\[
  \sh{S}\to\bigoplus_{n\geq0}\sh{L}^{\otimes n}
  \quad\mbox{e}\quad
  \sh{S}''\to\bigoplus_{n\geq0}\sh{L}^{\prime\prime\otimes n}
\]
de \sref{II.3.3.2.3} (veja-se \sref{II.8.8.5} abaixo), e $v$ e $u$ aos morfismos de inclusão $\sh{S}''\to\sh{S}$ e $\bigoplus_{n\geq0}\sh{L}^{\otimes nd}\to\bigoplus_{n\geq0}\sh{L}^{\otimes n}$, respectivamente;
é imediato por \sref{II.3.3.2} que esse diagrama é comutativo.
Acabamos de ver que $g''$ é um morfismo projetivo;
por outro lado, $u$ é um morfismo \emph{finito}.
Com efeito, como a questão é local sobre $X$, podemos supor que $X$ é afim de anel $A$ e que $\sh{L}=\sh{O}_X$;
tudo se reduz então a observar que o anel $A[T]$ é um módulo de tipo finito sobre seu subanel $A[T^d]$ (onde $T$ é uma indeterminada).
Como $Y$ é um esquema quase-compacto e $C''$ é afim sobre $Y$, sabemos que $C''$ também é um esquema quase-compacto,
\oldpage[II]{181}
e portanto $g''\circ u$ é um morfismo projetivo \sref{II.5.5.5}[ii];
pela comutatividade de \sref{II.8.8.4.1}, $v\circ g$ também é projetivo e, como $v$ é afim, portanto separado, concluímos finalmente que $g$ é projetivo \sref{II.5.5.5}[v].
\end{proof}
\begin{env}[8.8.5]
\label{II.8.8.5}
Retomemos a situação de \sref{II.8.8.1}.
Veremos que o morfismo $g:\bb{V}(\sh{L})\to C$ também pode ser definido de uma maneira válida para todo $\sh{O}_X$-módulo invertível $\sh{L}$ (não necessariamente amplo).
Para isso, consideremos o $f$-morfismo
\[
\label{II.8.8.5.1}
  \tau^\flat:\sh{S}\to\bigoplus_{n\geq0}\sh{L}^{\otimes n}
  \tag{8.8.5.1}
\]
correspondente ao morfismo $\tau$ de \sref{II.8.8.1.2}.
Isso induz, por \sref{II.1.5.6}, um morfismo $g':V\to C$ tal que, se $\pi:V\to X$ e $\psi:C\to Y$ são os morfismos estruturais, os diagramas
\[
\label{II.8.8.5.2}
  \xymatrix{
    X
      \ar[d]_{f}
  & V
      \ar[l]^{\pi}
      \ar[d]^{g'}
  \\Y
  & C
      \ar[l]_{\psi}
  }
  \qquad
  \xymatrix{
    X
      \ar[r]^{j}
      \ar[d]_{f}
  & V
      \ar[d]^{g'}
  \\Y
      \ar[r]_{\varepsilon}
  & C
  }
\tag{8.8.5.2}
\]
são comutativos (\eref{eq:2.8.5.1.2} e \eref{eq:2.8.5.1.3}).
Mostraremos que (se supusermos que $\sh{L}$ é $f$-amplo) \emph{os morfismos $g$ e $g'$ são idênticos}.

Como a questão é local sobre $Y$, podemos supor que $Y=\Spec(A)$ é afim e, por \sref{II.8.8.1.3}, identificar $X$ com um aberto de $P=\Proj(S)$, onde $S=A\oplus\bigoplus_{n\geq1}\Gamma(X,\sh{L}^{\otimes n})$;
deduzimos então, por \sref{II.8.8.1.4}, que $\Gamma(X,\sh{O}_P(n))=\Gamma(X,\sh{L}^{\otimes n})$ para todo $n\in\bb{Z}$.
Levando em conta a definição de $h=\Spec(\alpha)$, onde $\alpha$ é o $p$-morfismo canônico $\widetilde{S}\to\sh{S}_P^\geq$ \sref{II.8.6.1.2}, é preciso mostrar que a restrição a $X$ de $\alpha^\sharp:p^*(\widetilde{S})\to\sh{S}_P^\geq$ é idêntica a $\tau$.
Levando em conta \sref[0]{0.4.4.3}, basta mostrar que, se compusermos o homomorfismo canônico $\alpha_n:S_n\to\Gamma(P,\sh{O}_P(n))$ com o homomorfismo de restrição $\Gamma(P,\sh{O}_P(n))\to\Gamma(X,\sh{O}_P(n))=\Gamma(X,\sh{L}^{\otimes n})$, obtemos a identidade para todo $n>0$;
mas isso resulta imediatamente da definição da álgebra $S$ e da de $\alpha_n$ \sref{II.2.6.2}.
\end{env}
\begin{proposition}[8.8.6]
\label{II.8.8.6}
Suponhamos (com as notações de \sref{II.8.8.5}) que, se escrevermos $f=(f_0,\lambda)$, o homomorfismo $\lambda:\sh{O}_Y\to f_*(\sh{O}_X)$ seja bijetivo;
então:
\begin{enumerate}
  \item[\rm{(i)}] se escrevermos $g=(g_0,\mu)$, então $\mu:\sh{O}_C\to g_*(\sh{O}_V)$ é um isomorfismo; e
  \item[\rm{(ii)}] se $X$ é integral (resp. localmente integral e normal), então $C$ é integral (resp. normal).
\end{enumerate}
\end{proposition}

\begin{proof}
Com efeito, o $f$-morfismo $\tau^\flat$ é então um \emph{isomorfismo}
\[
  \tau^\flat:\sh{S}=\psi_*(\sh{O}_C)\to f_*(\pi_*(\sh{O}_V))=\psi_*(g_*(\sh{O}_V))
\]
e o $Y$-morfismo $g$ pode ser considerado como aquele para o qual o homomorfismo $\sh{A}(g)$ \sref{II.1.1.2} é igual a $\tau^\flat$.
Para ver que $\mu$ é um isomorfismo de $\sh{O}_C$-módulos, basta, por \sref{II.1.4.2}, ver que $\sh{A}(\mu):\psi_*(\sh{O}_C)\to\psi_*(g_*(\sh{O}_V))$ é um isomorfismo.
Mas, pela Definição~\sref{II.1.1.2}, temos $\sh{A}(\mu)=\sh{A}(g)$, de onde segue a conclusão de (i).

Para provar (ii), podemos reduzir-nos ao caso em que $Y$ é afim, e então $\sh{S}=\widetilde{S}$, com
\oldpage[II]{182}
$S=\bigoplus_{n\geq0}\Gamma(X,\sh{L}^{\otimes n})$;
a hipótese de que $X$ é integral implica que o anel $S$ é íntegro \sref[I]{I.7.4.4}, e portanto $C$ também é integral \sref[I]{I.5.1.4}.
Para mostrar que $C$ é normal, utilizaremos o lema seguinte:

\begin{lemma}[8.8.6.1]
\label{II.8.8.6.1}
Seja $Z$ um pré-esquema integral normal.
Então o anel $\Gamma(Z,\sh{O}_Z)$ é íntegro e integralmente fechado.
\end{lemma}

\begin{proof}
Segue de \sref[I]{I.8.2.1.1} que $\Gamma(Z,\sh{O}_Z)$ é a interseção, no corpo de funções racionais $R(Z)$, dos anéis integralmente fechados $\sh{O}_z$ para todos os $z\in Z$.
\end{proof}

Posto isso, mostremos primeiro que $V$ é \emph{localmente integral} e \emph{normal};
para isso, podemos reduzir-nos ao caso em que $X=\Spec(A)$ é afim, de anel $A$ íntegro e integralmente fechado \sref{II.6.3.8}, e em que $\sh{L}=\sh{O}_X$.
Como então $V=\Spec(A[T])$, e $A[T]$ é íntegro e integralmente fechado \cite[p.~99]{II-24}, isso demonstra nossa afirmação.
Para todo aberto afim $U$ de $C$, $g^{-1}(U)$ é quase-compacto, pois o morfismo $g$ é quase-compacto;
como $V$ é localmente integral, as componentes conexas de $g^{-1}(U)$ são pré-esquemas integrais abertos em $g^{-1}(U)$ e, portanto, são em número finito; além disso, como $V$ é normal, esses pré-esquemas também são normais \sref{II.6.3.8}.
Então $\Gamma(U,\sh{O}_C)$, que é igual a $\Gamma(g^{-1}(U),\sh{O}_V)$ por (i), é o produto direto finito de anéis íntegros e integralmente fechados \sref{II.8.8.6.1}, o que mostra que $C$ é normal \sref{II.6.3.4}.
\end{proof}
\subsection{O critério de amplitude de Grauert: enunciado}
\label{subsection:II.8.9}

Propomo-nos mostrar que as propriedades demonstradas em \sref{II.8.8.2} \emph{caracterizam} os módulos amplos relativamente a $f$ entre os $\sh{O}_X$-módulos e, mais precisamente, demonstrar o seguinte critério:
\begin{theorem}[8.9.1]
\label{II.8.9.1}
\emph{(Critério de Grauert).}
Sejam $Y$ um pré-esquema, $p:X\to Y$ um morfismo separado e quase-compacto, e $\sh{L}$ um $\sh{O}_X$-módulo invertível.
Para que $\sh{L}$ seja amplo relativamente a $p$, é necessário e suficiente que existam um $Y$-pré-esquema $C$, uma $Y$-seção $\varepsilon:Y\to C$ de $C$ e um $Y$-morfismo $q:\bb{V}(\sh{L})\to C$ que satisfaçam as seguintes propriedades:
\begin{enumerate}
  \item[\rm{(i)}] o diagrama
    \[
      \label{II.8.9.1.1}
        \xymatrix{
          X
            \ar[r]^{j}
            \ar[d]_{p}
        & \bb{V}(\sh{L})
            \ar[d]^{q}
        \\Y
            \ar[r]_{\varepsilon}
        & C
        }
      \tag{8.9.1.1}
    \]
    é comutativo, onde $j$ é a seção nula do fibrado vetorial $\bb{V}(\sh{L})$; e
  \item[\rm{(ii)}] a restrição de $q$ a $\bb{V}(\sh{L})\setmin j(X)$ é uma imersão aberta quase-compacta
    \[
      \bb{V}(\sh{L})\setmin j(X)\to C
    \]
    cuja imagem não intersecta $\varepsilon(Y)$.
\end{enumerate}
\end{theorem}

Observemos que, se $C$ é \emph{separado} sobre $Y$, podemos, na condição~(ii), suprimir a hipótese de que a imersão aberta seja quase-compacta;
para ver que essa propriedade (a quase-compacidade) é, de fato, consequência das outras condições, podemos reduzir-nos ao caso em que $Y$ é afim, e a afirmação decorre então de \sref[I]{I.5.5.1}[i] e \sref[I]{I.5.5.10}.
Podemos também suprimir
\oldpage[II]{183}
a mesma hipótese se supusermos que $X$ é noetheriano, pois então $V$ também é noetheriano, e a afirmação decorre de \sref[I]{I.6.3.5}.

\begin{corollary}[8.9.2]
\label{II.8.9.2}
Se o morfismo $p:X\to Y$ é próprio, podemos, no enunciado do Teorema~\sref{II.8.9.1}, supor que $q$ é próprio e substituir “imersão aberta” por “isomorfismo”.
\end{corollary}

De maneira mais sugestiva, podemos dizer (sempre que $p:X\to Y$ seja próprio) que \emph{$\sh{L}$ é amplo relativamente a $p$ se, e somente se, podemos “\emph{contrair}” a seção nula do fibrado vetorial $\bb{V}(\sh{L})$ ao pré-esquema de base $Y$.}
Um caso particular importante é aquele em que $Y$ é o espectro de um corpo e em que a operação de “contração” consiste em contrair a seção nula de $\bb{V}(\sh{L})$ \emph{a um único ponto}.

\begin{env}[8.9.3]
\label{II.8.9.3}
Tanto a necessidade das condições do Teorema \sref{II.8.9.1} como o Corolário \sref{II.8.9.2} decorrem imediatamente de \sref{II.8.8.2} e \sref{II.8.8.4}.

Para mostrar que as condições de \sref{II.8.9.1} são suficientes, consideremos uma situação ligeiramente mais geral.
Para isso, ponhamos (com as notações de \sref{II.8.8.2})
\[
  \sh{S}'=\bigoplus_{n\geq 0}\sh{L}^{\otimes n}
\]
e
\[
  V=\bb{V}(\sh{L})=\Spec(\sh{S}').
\]

O subpré-esquema fechado $j(X)$, seção nula de $\bb{V}(\sh{L})$, é definido pelo feixe quase-coerente de ideais $\sh{J}=(\sh{S}'_+)\supertilde$ de $\sh{O}_V$ \sref{II.1.4.10}.
Esse $\sh{O}_V$-módulo é \emph{inversível}, pois essa propriedade é local sobre $X$, e isso se reduz a observar que o ideal $TA[T]$ de um anel de polinômios $A[T]$ é um $A[T]$-módulo livre cíclico.
Além disso, é imediato (novamente porque a questão é local sobre $X$) que
\[
  \sh{L}=j^*(\sh{J})
\]
e
\[
  j_*(\sh{L})=\sh{J}/\sh{J}^2.
\]

Por outro lado, se
\[
  \pi:\bb{V}(\sh{L})\to X
\]
é o morfismo estrutural, então $\pi_*(\sh{J})=\sh{S}'_+$ e $\pi_*(\sh{J}/\sh{J}^2)=\sh{L}$;
há, portanto, homomorfismos canônicos $\sh{L}\to\pi_*(\sh{J})\to\sh{L}$, sendo o primeiro a injeção canônica $\sh{L}\to\sh{S}'_+$ e o segundo a projeção canônica de $\sh{S}'_+$ sobre $\sh{S}'_1=\sh{L}$, e sua composição é a identidade.
Podemos também incluir canonicamente $\pi_*(\sh{J}) = \sh{S}'_+ = \bigoplus_{n\geq1}\sh{L}^{\otimes n}$ no \emph{produto} $\prod_{n\geq1}\sh{L}^{\otimes n} = \varprojlim_n \pi_*(\sh{J}/\sh{J}^{n+1})$ (pois $\pi_*(\sh{J}/\sh{J}^{n+1}) = \sh{L}\oplus\sh{L}^{\otimes 2}\oplus\ldots\oplus\sh{L}^{\otimes n}$), e temos assim homomorfismos canônicos
\[
\label{II.8.9.3.1}
  \sh{L}\to\varprojlim\pi_*(\sh{J}/\sh{J}^{n+1})\to\sh{L}
  \tag{8.9.3.1}
\]
cuja composição é a identidade.

Com isso, a generalização de \sref{II.8.9.1} que vamos demonstrar é a seguinte:
\end{env}
\begin{proposition}[8.9.4]
\label{II.8.9.4}
Sejam $Y$ um pré-esquema, $V$ um $Y$-pré-esquema e $X$ um subpré-esquema fechado de $V$ definido por um ideal $\sh{J}$ de $\sh{O}_V$, que é um $\sh{O}_V$-módulo \emph{invertível};
se $j:X\to V$ é
\oldpage[II]{184}
a injeção canônica, ponhamos então $\sh{L} = j^*(\sh{J}) = \sh{J}\otimes_{\sh{O}_V}\sh{O}_X$, de modo que $j_*(\sh{L})=\sh{J}/\sh{J}^2$.
Suponhamos que o morfismo estrutural $p:X\to Y$ seja separado e quase-compacto, e que as condições seguintes sejam satisfeitas:
\begin{enumerate}
  \item[\rm{(i)}] existe um $Y$-morfismo $\pi:V\to X$ de tipo finito tal que $\pi\circ j=1_X$, e, portanto, $\pi_*(\sh{J}/\sh{J}^2)=\sh{L}$;
  \item[\rm{(ii)}] existe um homomorfismo de $\sh{O}_X$-módulos $\vphi:\sh{L}\to\varprojlim\pi_*(\sh{J}/\sh{J}^{n+1})$ tal que a composição
    \[
      \sh{L}
      \xrightarrow{\vphi}
      \varprojlim \pi_*(\sh{J}/\sh{J}^{n+1})
      \xrightarrow{\alpha}
      \pi_*(\sh{J}/\sh{J}^2) = \sh{L}
    \]
    (onde $\alpha$ é o homomorfismo canônico) é a identidade;
  \item[\rm{(iii)}] existem um $Y$-pré-esquema $C$, uma $Y$-seção $\varepsilon$ de $C$ e um $Y$-morfismo $q:V\to C$ tais que o diagrama
    \[
    \label{II.8.9.4.1}
      \xymatrix{
        X
          \ar[r]^{j}
          \ar[d]_{p}
      & V
          \ar[d]^{q}
      \\Y
          \ar[r]_{\varepsilon}
      & C
      }
    \tag{8.9.4.1}
    \]
    é comutativo; e
  \item[\rm{(iv)}] a restrição de $q$ a $W=V\setmin j(X)$ é uma imersão aberta quase-compacta em $C$, cuja imagem não intersecta $\varepsilon(Y)$.
\end{enumerate}
Então $\sh{L}$ é amplo relativamente a $p$.
\end{proposition}


\subsection{O critério de amplitude de Grauert: demonstração}
\label{subsection:II.8.10}

\begin{lemma}[8.10.1]
\label{II.8.10.1}
Sejam $\pi:V\to X$ um morfismo, $j:X\to V$ uma $X$-seção de $V$ que é também uma imersão fechada, e $\sh{J}$ um ideal quase-coerente de $\sh{O}_V$ que define o subpré-esquema fechado de $V$ associado a $j$.
Então valem todas as afirmações seguintes.
\begin{enumerate}
  \item[\rm{(i)}] Para todo $n\geq0$, $\pi_*(\sh{O}_V/\sh{J}^{n+1})$ e $\pi_*(\sh{J}/\sh{J}^{n+1})$ são $\sh{O}_X$-módulos quase-coerentes, e $\pi_*(\sh{O}_V/\sh{J})=\sh{O}_X$ e $\pi_*(\sh{J}/\sh{J}^2)=j^*(\sh{J})$.
  \item[\rm{(ii)}] Se $X=\{\xi\}=\Spec(k)$, onde $k$ é um corpo, então $\varprojlim\pi_*(\sh{O}_V/\sh{J}^{n+1})$ é isomorfo ao completamento separado do anel local $\sh{O}_{j(\xi)}$ para a topologia $\mathfrak{m}_{j(\xi)}$-ádica.
  \item[\rm{(iii)}] Suponhamos que $\sh{J}$ seja um $\sh{O}_V$-módulo invertível (o que implica que
    \[
      \sh{L} = j^*(\sh{J}) = \pi_*(\sh{J}/\sh{J}^2)
    \]
    é um $\sh{O}_X$-módulo invertível), e que exista um homomorfismo $\vphi:\sh{L}\to\varprojlim\pi_*(\sh{J}/\sh{J}^{n+1})$ tal que a composição $\sh{L} \xrightarrow{\vphi} \varprojlim\pi_*(\sh{J}/\sh{J}^{n+1}) \xrightarrow{\alpha} \pi_*(\sh{J}/\sh{J}^2)$ (onde $\alpha$ é o homomorfismo canônico) seja a identidade.
    Se pusermos $\sh{S}=\bigoplus_{n\geq0}\sh{L}^{\otimes n}$, então $\vphi$ induz canonicamente um isomorfismo de $\sh{O}_X$-álgebras entre o completamento $\widehat{\sh{S}}$ de $\sh{S}$ relativamente à sua filtração canônica (sendo esse completamento isomorfo ao produto $\prod_{n\geq0}\sh{L}^{\otimes n}$) e $\varprojlim\pi_*(\sh{O}_V/\sh{J}^{n+1})$.
\end{enumerate}
\end{lemma}

\begin{proof}
Observemos, antes de tudo, que o suporte do $\sh{O}_V$-módulo $\sh{O}_V/\sh{J}^{n+1}$ é $j(X)$, e que o suporte de $\sh{J}/\sh{J}^{n+1}$ está contido em $j(X)$.
No caso de (ii), $j(X)$ é um ponto fechado $j(\xi)$ de $V$,
\oldpage[II]{185}
e, por definição, $\pi_*(\sh{O}_V/\sh{J}^{n+1})$ é a fibra de $\sh{O}_V/\sh{J}^{n+1}$ no ponto $j(\xi)$, isto é, pondo $C=\sh{O}_{j(\xi)}$ e denotando por $\mathfrak{m}$ o ideal maximal de $C$, o $C$-módulo $C/\mathfrak{m}^{n+1}$;
a afirmação (ii) é então evidente.

Para demonstrar (i), observemos que a questão é local em $X$;
podemos, portanto, restringir-nos ao caso em que $X$ é afim.
Seja $U$ um aberto afim de $V$;
então $j(X)\cap U$ é um aberto afim de $j(X)$, de modo que $U_0=\pi(j(X)\cap U)$ é um aberto afim de $X$ isomorfo a $j(X)\cap U$;
para todo aberto afim $W_0\subset U_0$ de $X$, $W=\pi^{-1}(W_0)\cap U$ é um aberto afim de $V$, pois $X$ é um esquema \sref[I]{I.5.5.10};
em particular, $U'=U\cap\pi^{-1}(U_0)$ é um aberto afim de $V$, e é claro que $\pi(U')=U_0$ e $j(U_0)=j(X)\cap U$.
Ora, por definição,
\[
  \Gamma(W_0,\pi_*(\sh{O}_V/\sh{J}^{n+1}))
  =\Gamma(\pi^{-1}(W_0),\sh{O}_V/\sh{J}^{n+1});
\]
mas, como todo ponto de $\pi^{-1}(W_0)$ que não pertence a $j(W_0)$ possui uma vizinhança aberta em $\pi^{-1}(W_0)$ que não intersecta $j(X)$ e na qual $\sh{O}_V/\sh{J}^{n+1}$ é, portanto, nulo, é claro que as seções de $\sh{O}_V/\sh{J}^{n+1}$ sobre $\pi^{-1}(W_0)$ e sobre $W$ se correspondem biunivocamente.
Em outros termos, se $\pi'$ é a restrição de $\pi$ a $U'$, os $(\sh{O}_X|U_0)$-módulos
\[
  \pi_*(\sh{O}_V/\sh{J}^{n+1})|U_0
  \quad\hbox{e}\quad
  \pi'_*\bigl((\sh{O}_V/\sh{J}^{n+1})|U'\bigr)
\]
são idênticos.
Como $U'$ e $U_0$ são afins e os $U_0$ recobrem $X$, concluímos então \sref[I]{I.1.6.3} que $\pi_*(\sh{O}_V/\sh{J}^{n+1})$ é quase-coerente, e a demonstração é idêntica para $\pi_*(\sh{J}/\sh{J}^{n+1})$.

Para demonstrar finalmente (iii), observemos que $\sh{S}$ é precisamente $\bb{S}_{\sh{O}_X}(\sh{L})$;
por conseguinte, $\vphi$ induz canonicamente um homomorfismo de $\sh{O}_X$-álgebras
\[
  \psi:\sh{S}\longrightarrow
  \varprojlim_n\pi_*(\sh{O}_V/\sh{J}^{n+1})
\]
\sref{II.1.7.4};
além disso, esse homomorfismo envia $\sh{L}^{\otimes n}$ a $\varprojlim_m\pi_*(\sh{J}^n/\sh{J}^{m+1})$, de modo que é contínuo para as topologias consideradas e, por conseguinte, se prolonga em um homomorfismo
\[
  \widehat{\psi}:\widehat{\sh{S}}\longrightarrow
  \varprojlim_n\pi_*(\sh{O}_V/\sh{J}^{n+1}).
\]
Para ver que se trata efetivamente de um isomorfismo, podemos, como na demonstração de (i), reduzir-nos ao caso em que $X=\Spec(A)$ e $V=\Spec(B)$ são afins, com $\sh{J}=\widetilde{\mathfrak{J}}$, onde $\mathfrak{J}$ é um ideal de $B$;
a $\pi$ corresponde uma injeção $A\to B$ que identifica $A$ com um subanel de $B$ complementar a $\mathfrak{J}$, e $\sh{L}$ (resp. $\pi_*(\sh{O}_V/\sh{J}^{n+1})$) é o $\sh{O}_X$-módulo quase-coerente associado ao $A$-módulo $L=\mathfrak{J}/\mathfrak{J}^2$ (resp. $B/\mathfrak{J}^{n+1}$).
Como $\sh{J}$ é um $\sh{O}_V$-módulo \emph{invertível}, podemos ainda supor que $\mathfrak{J}=Bt$, onde $t$ não é divisor de zero em $B$.
Da relação $B=A\oplus Bt$ deduzimos então que, para todo $n>0$,
\[
  B=A\oplus At\oplus At^2\oplus\cdots\oplus At^n\oplus Bt^{n+1}
\]
e, portanto, existe um $A$-isomorfismo canônico entre o anel de séries formais $A[[T]]$ e $C=\varprojlim B/\mathfrak{J}^{n+1}$ que envia $T$ a $t$.
Temos também $L=A\overline{t}$, onde $\overline{t}$ é a classe de $t$ módulo $Bt^2$, e, por hipótese, o homomorfismo $\vphi$ envia $\overline{t}$ a um elemento $t'\in C$ congruente a $t$ módulo $Ct^2$.
Deduzimos então, por indução em $n$, que
\[
  A\oplus At'\oplus\cdots\oplus A(t')^n\oplus Ct^{n+1}
  =A\oplus At\oplus\cdots\oplus At^n\oplus Ct^{n+1},
\]
o que prova que o homomorfismo $\widehat{\psi}$ corresponde efetivamente a um isomorfismo entre $\prod_{n\geq0}L^{\otimes n}$ e $C$.
\end{proof}

\begin{lemma}[8.10.2]
\label{II.8.10.2}
Sob as hipóteses do Lema~\sref{II.8.10.1}, seja $g:X'\to X$ um morfismo,
\oldpage[II]{186}
seja $V'=V\times_X X'$, e sejam $\pi':V'\to X'$ e $g':V'\to V$ as projeções canônicas, de modo que temos o diagrama comutativo
\[
  \xymatrix{
    V
      \ar[d]_{\pi}
  & V'
      \ar[l]^{g'}
      \ar[d]^{\pi'}
  \\X
  & X'
      \ar[l]_{g}
  }
\]

Então $j'=j\times1_{X'}$ é uma $X'$-seção de $V'$ que é, além disso, uma imersão fechada, e $\sh{J}'=g^{'*}(\sh{J})\sh{O}_{V'}$ é o ideal quase-coerente de $\sh{O}_{V'}$ que define o subpré-esquema fechado de $V'$ associado a $j'$.
Além disso,
\[
  \pi'_*(\sh{O}_{V'}/\sh{J}'^{\,n+1})
  =g^*(\pi_*(\sh{O}_V/\sh{J}^{n+1})).
\]
Por fim, $\sh{J}'$ é um $\sh{O}_{V'}$-módulo canonicamente isomorfo a $g^{'*}(\sh{J})$ e, em particular, é invertível se $\sh{J}$ é um $\sh{O}_V$-módulo invertível.
\end{lemma}

\begin{proof}
O fato de $j'$ ser uma imersão fechada decorre de \sref[I]{I.4.3.1}, e $j'$ é uma $X'$-seção de $V'$ pela funtorialidade da mudança de base de pré-esquemas.
Além disso, se $Z$ (resp. $Z'$) é o subpré-esquema fechado de $V$ (resp. $V'$) associado a $j$ (resp. $j'$), então $Z'=g^{'-1}(Z)$ \sref[I]{I.4.3.1}, e a segunda afirmação decorre de \sref[I]{I.4.4.5}.
Para demonstrar as demais afirmações, vemos, como em \sref{II.8.10.1}, que podemos nos reduzir ao caso em que $X$, $V$ e $X'$ (e, portanto, também $V'$) sejam afins;
conservamos as notações da demonstração de \sref{II.8.10.1} e pomos $X'=\Spec(A')$.
Então $V'=\Spec(B')$, onde $B'=B\otimes_A A'$, e $\sh{J}'=\widetilde{\mathfrak{J}'}$, onde $\mathfrak{J}'=\Im(\mathfrak{J}\otimes_A A')$.
Portanto,
\[
  B'/\mathfrak{J}'^{\,n+1}
  =(B/\mathfrak{J}^{n+1})\otimes_A A';
\]
além disso, como $\mathfrak{J}$ é um fator direto de $B$ como $A$-módulo, $\mathfrak{J}\otimes_A A'$ é um fator direto de $B'$ como $A'$-módulo e, por conseguinte, se identifica canonicamente com $\mathfrak{J}'$.
\end{proof}

\begin{corollary}[8.10.3]
\label{II.8.10.3}
Suponhamos satisfeitas as hipóteses do Lema~\sref{II.8.10.1}, e suponhamos, além disso, que $\pi$ seja de tipo finito e que $\sh{J}$ seja um $\sh{O}_V$-módulo invertível.
Então, para todo $x\in X$, o anel local da fibra $\pi^{-1}(x)$ no ponto $j(x)$ é um anel regular (portanto, íntegro) de dimensão 1, cujo completamento é isomorfo ao anel de séries formais $\kres(x)[[T]]$ (onde $T$ é uma indeterminada);
além disso, existe exatamente uma componente irredutível de $\pi^{-1}(x)$ que contém $j(x)$.
\end{corollary}

\begin{proof}
Como $\pi^{-1}(x)=V\times_X\Spec(\kres(x))$, \sref{II.8.10.2} nos reduz ao caso em que $X$ é o espectro de um corpo $K$.
Como $\pi$ é de tipo finito \sref[I]{I.6.3.4}[iv], $\sh{O}_{j(x)}$ é então um anel local noetheriano e, portanto, separado para a topologia $\mathfrak{m}_{j(x)}$-ádica \sref[0]{0.7.3.5};
de \sref{II.8.10.1}[ii e iii] segue-se que o completamento desse anel é isomorfo a $K[[T]]$ e, por conseguinte, $\sh{O}_{j(x)}$ é regular e de dimensão 1 (\cite[p.~17-01, th.~1]{I-1});
por fim, como $\sh{O}_{j(x)}$ é íntegro, $j(x)$ pertence exatamente a uma das componentes irredutíveis (em número finito) de $\pi^{-1}(x)$ \sref[I]{I.5.1.4}.
\end{proof}

\begin{corollary}[8.10.4]
\label{II.8.10.4}
Suponhamos satisfeitas as hipóteses do Lema~\sref{II.8.10.1}, e suponhamos, além disso, que $\sh{J}$ seja um $\sh{O}_V$-módulo invertível.
Seja $W=V\setmin j(X)$;
para todo ideal quase-coerente $\sh{K}$ de $\sh{O}_X$, ponhamos $\sh{K}_V=\pi^*(\sh{K})\sh{O}_V$ e $\sh{K}_W=\sh{K}_V|W$.
Então $\sh{K}_V$ é o maior ideal quase-coerente de $\sh{O}_V$ cuja restrição a $W$ é $\sh{K}_W$.
\end{corollary}

\begin{proof}
Com efeito, vemos, como em \sref{II.8.10.1}, que a questão é local sobre $X$ e sobre $V$;
podemos, portanto, conservar as notações da demonstração de \sref{II.8.10.1}, com $\mathfrak{J}=Bt$, onde $t$ não é divisor de zero em $B$.
Além disso, $W=\Spec(B_t)$ e $\sh{K}=\widetilde{\mathfrak{K}}$, onde $\mathfrak{K}$ é um ideal de $A$;
daí $\pi^*(\sh{K})\sh{O}_V=(\mathfrak{K}.B)\supertilde$ \sref[I]{I.1.6.9}, $\sh{K}_W=(\mathfrak{K}.B_t)\supertilde$, e o maior ideal
\oldpage[II]{187}
de $B$ cuja imagem canônica em $B_t$ seja $\mathfrak{K}.B_t$ é a imagem inversa de $\mathfrak{K}.B_t$, isto é, o conjunto dos $s\in B$ tais que, para algum inteiro $n>0$, se tenha $t^ns\in\mathfrak{K}.B$.
É preciso demonstrar que esta última relação implica que $s\in\mathfrak{K}.B$, ou ainda que a imagem canônica de $t$ não é divisor de zero em $B/\mathfrak{K}B=(A/\mathfrak{K})\otimes_A B$, o que decorre de \sref{II.8.10.2} aplicado a $X'=\Spec(A/\mathfrak{K})$.
\end{proof}

\begin{corollary}[8.10.5]
\label{II.8.10.5}
Suponhamos satisfeitas as hipóteses de \sref{II.8.10.3};
sejam $W=V\setmin j(X)$, $x$ um ponto de $X$, $\sh{K}$ um ideal quase-coerente de $\sh{O}_X$, e $z$ o ponto genérico da componente irredutível de $\pi^{-1}(x)$ que contém $j(x)$ \sref{II.8.10.3}.
\begin{enumerate}
    \item[\rm{(i)}] Seja $g$ uma seção de $\sh{O}_V$ sobre $V$ tal que $g|W$ seja uma seção de $\sh{K}_W$ sobre $W$ (com as notações de \sref{II.8.10.4}).
        Então $g$ é uma seção de $\sh{K}_V$;
        se, além disso, $g(z)\neq0$, e se, para todo inteiro $m>0$, denotarmos por $g_m^x$ o germe no ponto $x$ da imagem canônica $g_m$ de $g$ em $\Gamma(X,\pi_*(\sh{O}_V/\sh{J}^{m+1}))$, então existe um inteiro $m>0$ tal que a imagem de $g_m^x$ em
        \[
            (\pi_*(\sh{O}_V/\sh{J}^{m+1}))_x \otimes_{\sh{O}_x} \kres(x)
        \]
        seja diferente de zero.
    \item[\rm{(ii)}] Suponhamos, além disso, satisfeitas as condições de \sref{II.8.10.1}[iii].
        Então, se existe uma seção $g$ de $\sh{K}_V$ sobre $V$ tal que $g(z)\neq0$, existe um inteiro $n\geq0$ e uma seção $f$ de $\sh{K}.\sh{L}^{\otimes n}=\sh{K}\otimes\sh{L}^{\otimes n}\subset\sh{L}^{\otimes n}$ tal que $f(x)\neq0$.
        Se $g$ é uma seção de $\sh{J}$, pode-se tomar $n>0$.
\end{enumerate}
\end{corollary}

\begin{proof}
\begin{enumerate}
    \item[\rm{(i)}] Como o ideal de $\sh{O}_W$ gerado por $g|W$ está contido em $\sh{K}_W$ por hipótese, o ideal de $\sh{O}_V$ gerado por $g$ está contido em $\sh{K}_V$ por \sref{II.8.10.4}; em outras palavras, $g$ é uma seção de $\sh{K}_V$.
        Para demonstrar a segunda afirmação de (i), podemos supor novamente que $X$ e $V$ sejam afins e conservar as notações de \sref{II.8.10.1};
        a fibra $\pi^{-1}(x)$ é então afim, de anel $B'=B\otimes_A\kres(x)$, e existe em $B'$ um elemento $t'$ que não é divisor de zero e tal que $B'=\kres(x)\oplus B't'$.
        Como $j(x)$ é uma especialização de $z$ e $g(z)\neq0$, tem-se necessariamente $g_{j(x)}\neq0$.
        Mas $\sh{O}_{j(x)}$ é um anel local separado \sref{II.8.10.3} e, portanto, se injeta em seu completamento; por conseguinte, a imagem de $g$ nesse completamento não é nula.
        Ora, esse completamento é isomorfo a $\varprojlim_n(B'/B't'^{\,n+1})$ \sref{II.8.10.3};
        se $g'=g\otimes1\in B'$, existe então um inteiro $m$ tal que $g'\notin B't'^{\,m+1}$, ou ainda, a imagem $g'_m$ de $g'$ em $B'/B't'^{\,m+1}$ é diferente de zero.
        Mas, como $g'_m$ é precisamente a imagem de $g_m^x$, nossa afirmação está demonstrada.
    \item[\rm{(ii)}] Em virtude de \sref{II.8.10.1}[iii], $\pi_*(\sh{O}_V/\sh{J}^{m+1})$ é isomorfo à soma direta dos $\sh{L}^{\otimes k}$ para $0\leq k\leq m$;
        denotemos por $f_k$ a seção de $\sh{L}^{\otimes k}$ sobre $X$ que é a componente do elemento de $\bigoplus_{k=0}^m\Gamma(X,\sh{L}^{\otimes k})$ correspondente a $g_m$ por esse isomorfismo.
        Escolhendo $m$ como em (i), há, portanto, um índice $k$ tal que $f_k(x)\neq0$, por (i).
        Para ver que $f_k$ é uma seção de $\sh{K}\sh{L}^{\otimes k}$, basta considerar, como acima, o caso em que $X$ e $V$ sejam afins, e isso decorre imediatamente do fato de que $g\in\mathfrak{K}.B$ (com as notações de \sref{II.8.10.4}).
        A última afirmação decorre do fato de que a hipótese $g\in\Gamma(V,\sh{J})$ implica que $f_0=0$.
\end{enumerate}
\end{proof}

\begin{env}[8.10.6]
\label{II.8.10.6}
\emph{Demonstração de \sref{II.8.9.4}.}
A questão é local sobre $Y$ \sref{II.4.6.4};
como $\varepsilon$ é uma $Y$-seção, podemos, portanto, substituir $C$ por uma vizinhança aberta afim $U$ de um ponto de $\varepsilon(Y)$ tal que $\varepsilon(Y)\cap U$ seja fechado em $U$.
Em outras palavras, podemos supor que $C$ seja afim e que $Y$ seja um subpré-esquema fechado de $C$ (e, portanto, também afim) definido por um feixe quase-coerente
\oldpage[II]{188}
$\sh{I}$ de ideais de $\sh{O}_C$.
Como $p$ é separado e quase-compacto, $X$ é então um esquema quase-compacto, e ficamos reduzidos a demonstrar que $\sh{L}$ é \emph{amplo} \sref{II.4.6.4}.
Em virtude do critério~\sref{II.4.5.2}[a], devemos, portanto, demonstrar o seguinte:
para todo ideal quase-coerente $\sh{K}$ de $\sh{O}_X$ e todo ponto $x\in X$ que não pertença ao suporte de $\sh{O}_X/\sh{K}$, existe um inteiro $n>0$ e uma seção $f$ de $\sh{K}\otimes\sh{L}^{\otimes n}$ sobre $X$ tal que $f(x)\neq0$.

Para isso, ponhamos
\begin{align*}
  \sh{K}_V &= \pi^*(\sh{K})\sh{O}_V,\\
  \sh{K}_W &= \sh{K}_V|W,
\end{align*}
onde $W=V\setmin j(X)$;
como a restrição de $q$ a $W$ é uma imersão aberta quase-compacta em $C$, segue de \sref[I]{I.9.4.2} que $\sh{K}_W$ é a restrição a $W$ de um ideal quase-coerente $\sh{K}'_V$ de $\sh{O}_V$ da forma
\[
    \sh{K}'_V = q^*(\sh{K}_C)\sh{O}_V,
\]
onde $\sh{K}_C$ é um ideal quase-coerente de $\sh{O}_C$.
Além disso, como, por hipótese, $q^{-1}(Y)\subset j(X)$, e como $Y$ é definido pelo ideal $\sh{I}$, a restrição a $W$ de $q^*(\sh{I})\sh{O}_V$ é igual à de $\sh{O}_V$; portanto, $\sh{K}_W$ é também a restrição a $W$ de $q^*(\sh{I}\sh{K}_C)\sh{O}_V$, e podemos supor que $\sh{K}_C\subset\sh{I}$, de onde
\[
\label{II.8.10.6.1}
    \sh{K}'_V \subset q^*(\sh{I})\sh{O}_V \subset \sh{J}
\tag{8.10.6.1}
\]
levando em conta \sref[I]{I.4.4.6} e a comutatividade de \sref{II.8.9.4.1}.
Além disso, deduz-se de \sref{II.8.10.4} que
\[
\label{II.8.10.6.2}
    \sh{K}'_V \subset \sh{K}_V.
\tag{8.10.6.2}
\]
Assim sendo, segue de \sref{II.8.10.3} que $j(x)$ pertence a uma única componente irredutível de $\pi^{-1}(x)$;
seja $z$ o ponto genérico dessa componente, e ponhamos $z'=q(z)$.
Em virtude de \sref{II.8.10.5}, a demonstração estará concluída (levando em conta \sref{II.8.10.6.1} e \sref{II.8.10.6.2}) se demonstrarmos a existência de uma seção $g$ de $\sh{K}'_V$ sobre $V$ tal que $g(z)\neq0$.
Mas, por hipótese, $\sh{K}$ tem uma restrição igual à de $\sh{O}_X$ em uma vizinhança aberta de $x$;
por outro lado, segue de \sref{II.8.10.3} que $z\neq j(x)$, e, portanto, $z\in W$, de modo que $(\sh{K}_W)_z=\sh{O}_{V,z}$ e, por definição, $(\sh{K}_C)_{z'}=\sh{O}_{C,z'}$.
Como $C$ é afim, existe, pois, uma seção $g'$ de $\sh{K}_C$ sobre $C$ tal que $g'(z')\neq0$; tomando para $g$ a seção de $\sh{K}'_V$ que corresponde canonicamente a $g'$, temos efetivamente $g(z)\neq0$, o que conclui a demonstração.
\end{env}

\begin{remark}[8.10.7]
\label{II.8.10.7}
Não sabemos se, no enunciado de \sref{II.8.9.4}, a condição~(ii) é ou não supérflua.
Em todo caso, a conclusão deixa de ser válida se não se supuser a existência de um $Y$-morfismo $\pi:V\to X$ tal que $\pi\circ j=1_X$;
indiquemos brevemente como se pode, de fato, construir um contraexemplo, cujos detalhes só serão desenvolvidos mais adiante.
Tomamos $Y=\Spec(k)$, onde $k$ é um corpo, e $C=\Spec(A)$, onde $A=k[T_1,T_2]$, com a $Y$-seção $\varepsilon$ correspondente ao homomorfismo de aumentação $A\to k$.
Denotamos por $C'$ o esquema obtido de $C$ pela explosão do ponto fechado $a=\varepsilon(Y)$ de $C$;
se $D$ é a imagem inversa de $a$ em $C'$, consideramos em $D$ um ponto fechado $b$,
\oldpage[II]{189}
e denotamos por $V$ o esquema obtido de $C'$ pela explosão de $b$;
$X$ é o subpré-esquema fechado de $V$ dado pela imagem inversa de $a$ pelo morfismo estrutural $q:V\to C$.
Mostremos agora que $X$ é a união de duas componentes irredutíveis, $X_1$ e $X_2$, onde $X_1$ é a imagem inversa de $b$ em $V$.
É imediato que o ideal $\sh{J}$ de $\sh{O}_V$ que define $X$ continua sendo invertível, mas pode-se provar que $j^*(\sh{J})=\sh{L}$ (onde $j$ é a injeção canônica $X\to V$) não é amplo, considerando o «grau» da imagem inversa de $\sh{L}$ sobre $X_1$, que seria $>0$ se $\sh{L}$ fosse amplo, mas que se pode demonstrar (por um cálculo elementar de interseções) ser igual a $0$.
\end{remark}

\subsection{Unicidade das contrações}
\label{subsection:II.8.11}

\begin{lemma}[8.11.1]
\label{II.8.11.1}
Sejam $U$ e $V$ dois pré-esquemas, e seja $h=(h_0,\lambda):U\to V$ um morfismo sobrejetivo.
Suponhamos que:
\begin{enumerate}
  \item[$1^\circ$] $\lambda:\sh{O}_V\to h_*(\sh{O}_U)=(h_0)_*(\sh{O}_U)$ é um isomorfismo.
  \item[$2^\circ$] O espaço subjacente a $V$ se identifica com o espaço quociente do espaço subjacente a $U$ pela relação
  \[
      R:h_0(x)=h_0(y)
  \]
  (condição que é sempre satisfeita quando o morfismo $h$ é \emph{aberto} ou \emph{fechado}, ou, \emph{a fortiori}, quando $h$ é \emph{próprio}).
\end{enumerate}
Então, para todo pré-esquema $W$, a aplicação
\[
\label{II.8.11.1.1}
    \Hom(V,W)\longrightarrow\Hom(U,W)
\tag{8.11.1.1}
\]
que a cada morfismo $v=(v_0,\nu)$ de $V$ em $W$ associa o morfismo $u=v\circ h=(u_0,\mu)$, é uma bijeção de $\Hom(V,W)$ sobre o conjunto dos $u$ tais que $u_0$ seja constante em cada fibra $h_0^{-1}(x)$.
\end{lemma}

\begin{proof}
É claro que, se $u=v\circ h$, de modo que $u_0=v_0\circ h_0$, então $u_0$ é constante em cada conjunto $h_0^{-1}(x)$.
Reciprocamente, se $u$ tem essa propriedade, mostraremos que existe um único $v\in\Hom(V,W)$ tal que $u=v\circ h$.
A existência e a unicidade da aplicação contínua $v_0:V\to W$ tal que $u_0=v_0\circ h_0$ decorrem das hipóteses, pois $h_0$ se identifica com a aplicação canônica de $U$ sobre $U/R$.
Podemos também, substituindo $V$ por um pré-esquema isomorfo se necessário, supor que $\lambda$ é a identidade;
por hipótese, $\mu$ é então um homomorfismo
\[
    \mu:\sh{O}_W\longrightarrow(u_0)_*(\sh{O}_U)
    =(v_0)_*((h_0)_*(\sh{O}_U))
\]
tal que o homomorfismo correspondente $\mu^\sharp:u_0^*(\sh{O}_W)\to\sh{O}_U$ é \emph{local} em cada fibra.
Como $(v_0)_*((h_0)_*(\sh{O}_U))=(v_0)_*(\sh{O}_V)$, temos necessariamente $\nu=\mu$, e tudo se reduz a demonstrar que o homomorfismo correspondente $\nu^\sharp:v_0^*(\sh{O}_W)\to\sh{O}_V$ é local em cada fibra.
Ora, todo $y\in V$ é da forma $h_0(x)$ para algum $x\in U$;
ponhamos $z=v_0(y)=u_0(x)$.
Então, por \sref[0]{0.3.5.5}, o homomorfismo $\mu_x^\sharp$ se fatora como
\[
    \mu_x^\sharp:\sh{O}_z
    \xrightarrow{\nu_y^\sharp}\sh{O}_y
    \xrightarrow{\lambda_x^\sharp}\sh{O}_x.
\]
Por hipótese, $\lambda_x^\sharp$ e $\mu_x^\sharp$ são homomorfismos locais;
portanto, $\lambda_x^\sharp$ transforma todo elemento invertível de $\sh{O}_y$ em um elemento invertível de $\sh{O}_x$;
se $\nu_y^\sharp$ transformasse um elemento não invertível de $\sh{O}_z$ em um elemento invertível de $\sh{O}_y$, então $\mu_x^\sharp$ transformaria esse elemento de $\sh{O}_z$ em um elemento invertível de $\sh{O}_x$, contradizendo a hipótese, de onde resulta o lema.
\end{proof}

\begin{corollary}[8.11.2]
\label{II.8.11.2}
Sejam $U$ um pré-esquema integral e $V$ um pré-esquema normal;
então todo morfismo $h:U\to V$ que seja universalmente fechado, birracional e radicial é um isomorfismo.
\end{corollary}

\begin{proof}
Se $h=(h_0,\lambda)$, segue das hipóteses que $h_0$ é injetiva e fechada, e que $h_0(U)$
\oldpage[II]{190}
é densa em $V$, de modo que $h_0$ é um \emph{homeomorfismo} de $U$ sobre $V$.
Para demonstrar o corolário, bastará provar que $\lambda:\sh{O}_V\to(h_0)_*(\sh{O}_U)$ é um isomorfismo: poderemos então aplicar \sref{II.8.11.1}, que provará que a aplicação \sref{II.8.11.1.1} é bijetiva (pois cada uma das fibras $h_0^{-1}(x)$ se reduz a um único ponto);
portanto, $h$ será um isomorfismo.
Como a questão é evidentemente local sobre $V$, podemos supor que $V=\Spec(A)$ é afim, de anel $A$ íntegro e integralmente fechado \sref{II.8.8.6.1};
$h$ corresponde então \sref[I]{I.2.2.4} a um homomorfismo $\vphi:A\to\Gamma(U,\sh{O}_U)$, e tudo se reduz a provar que $\vphi$ é um isomorfismo.
Ora, se $K$ é o corpo de frações de $A$, então $\Gamma(U,\sh{O}_U)$ tem, por hipótese, $K$ como corpo de frações, e $A$ é um subanel de $\Gamma(U,\sh{O}_U)$, sendo $\vphi$ a injeção canônica \sref[I]{I.8.2.7}.
Como o morfismo $h$ satisfaz as hipóteses de \sref{II.7.3.11}, $\Gamma(U,\sh{O}_U)$ é um subanel do fecho integral de $A$ em $K$ e, portanto, é idêntico a $A$ por hipótese.
\end{proof}

\begin{remark}[8.11.3]
\label{II.8.11.3}
Veremos no Capítulo~III \sref[III]{III.4.4.11} que, quando $V$ é um pré-esquema \emph{localmente noetheriano}, todo morfismo $h:U\to V$ que seja próprio e quase-finito (em particular, todo morfismo que satisfaça as hipóteses de \sref{II.8.11.2}) é necessariamente \emph{finito}.
A conclusão de \sref{II.8.11.2} decorre então, nesse caso, de \sref{II.6.1.15}.
\end{remark}

\begin{env}[8.11.4]
\label{II.8.11.4}
Veremos agora que, no critério de Grauert \sref{II.8.9.1}, pode-se frequentemente afirmar que o pré-esquema $C$ e a «contração» $q$ estão determinados de maneira \emph{essencialmente única}.
\end{env}

\begin{lemma}[8.11.5]
\label{II.8.11.5}
Sejam $Y$ um pré-esquema, $p:X\to Y$ um morfismo próprio, $\sh{L}$ um $\sh{O}_X$-módulo invertível amplo relativamente a $p$, $C$ um $Y$-pré-esquema, $\varepsilon:Y\to C$ uma $Y$-seção e $q:V=\bb{V}(\sh{L})\to C$ um $Y$-morfismo, tais que o diagrama de \sref{II.8.9.1.1} seja comutativo.
Suponhamos, além disso, que, se $p=(p_0,\theta)$, então $\theta:\sh{O}_Y\to p_*(\sh{O}_X)$ seja um isomorfismo.
Sejam $\sh{S}'=\bigoplus_{n\geq0}p_*(\sh{L}^{\otimes n})$ e $C'=\Spec(\sh{S}')$, e seja $q':\bb{V}(\sh{L})\to C'$ o $Y$-morfismo canônico \sref{II.8.8.5}.
Então existe um único $Y$-morfismo $u:C'\to C$ tal que $q=u\circ q'$.
\end{lemma}

\begin{proof}
A hipótese sobre $\theta$ implica, em particular, que $p$ é sobrejetivo;
como, em virtude de \sref{II.8.8.4}, a restrição de $q'$ a $\bb{V}(\sh{L})\setmin j(X)$ é um \emph{isomorfismo} sobre $C'\setmin\varepsilon'(Y)$ (onde $\varepsilon'$ é a seção vértice de $C'$), segue de \sref{II.8.8.4} que $q'$ é \emph{próprio e sobrejetivo};
além disso, em virtude de \sref{II.8.8.6}, se pusermos $q'=(q'_0,\tau)$, então $\tau:\sh{O}_{C'}\to q'_*(\sh{O}_V)$ é um isomorfismo.
Estamos, portanto, nas condições de aplicação de \sref{II.8.11.1}, e teremos demonstrado o lema se provarmos que $q$ é constante em toda fibra $q'{}^{-1}(z')$, onde $z'\in C'$.
Ora, essa condição é trivialmente satisfeita para $z'\notin\varepsilon'(Y)$.
Se $z'\in\varepsilon'(Y)$, existe um único $y\in Y$ tal que $z'=\varepsilon'(y)$ e, em virtude da comutatividade de \sref{II.8.8.5.2} e do fato de que $q'$ envia $\bb{V}(\sh{L})\setmin j(X)$ em $C'\setmin\varepsilon'(Y)$, tem-se $q'{}^{-1}(z')=j(p^{-1}(y))$;
a comutatividade do diagrama de \sref{II.8.9.1.1} demonstra então nossa afirmação.
\end{proof}

\begin{corollary}[8.11.6]
\label{II.8.11.6}
Sob as hipóteses de \sref{II.8.11.5}, suponhamos, além disso, que $q$ seja próprio e que a restrição de $q$ a $\bb{V}(\sh{L})\setmin j(X)$ seja um isomorfismo sobre $C\setmin\varepsilon(Y)$.
Então o morfismo $u$ é universalmente fechado, sobrejetivo e radicial, e sua restrição a $C'\setmin\varepsilon'(Y)$ é um isomorfismo sobre $C\setmin\varepsilon(Y)$.
\end{corollary}

\begin{proof}
Como $q'$ é um isomorfismo de $\bb{V}(\sh{L})\setmin j(X)$ sobre $C'\setmin\varepsilon'(Y)$ \sref{II.8.8.4}, a última afirmação decorre imediatamente de $q=u\circ q'$.
Além disso, a comutatividade dos diagramas
\oldpage[II]{191}
\sref{II.8.8.5.2} e \sref{II.8.9.1.1} mostra que a restrição de $u$ ao subpré-esquema fechado $\varepsilon'(Y)$ de $C'$ é um isomorfismo sobre o subpré-esquema fechado $\varepsilon(Y)$ de $C$, de onde resulta imediatamente que, para todo $z'\in\varepsilon'(Y)$, se $z=u(z')$, então $u$ define um isomorfismo de $\kres(z)$ sobre $\kres(z')$.
Estas observações provam que $u$ é bijetivo e radicial;
além disso, se $\psi:C\to Y$ e $\psi':C'\to Y$ são os morfismos estruturais, então $\psi'=\psi\circ u$ e, como $\psi'$ é separado \sref{II.1.2.4}, o mesmo vale para $u$ \sref[I]{I.5.5.1}[v].
Vimos, por outro lado, na demonstração do lema \sref{II.8.11.5}, que $q'$ é sobrejetivo;
como $q=u\circ q'$ é próprio, concluímos finalmente de \sref{II.5.4.3} e \sref{II.5.4.9} que $u$ é universalmente fechado.
\end{proof}

\begin{proposition}[8.11.7]
\label{II.8.11.7}
Sejam $Y$ um pré-esquema, $X$ um pré-esquema íntegro, $p:X\to Y$ um morfismo próprio, $\sh{L}$ um $\sh{O}_X$-módulo invertível amplo relativamente a $p$, $C$ um $Y$-pré-esquema normal, $\varepsilon:Y\to C$ uma $Y$-seção e $q:V=\bb{V}(\sh{L})\to C$ um $Y$-morfismo, tais que o diagrama de \sref{II.8.9.1.1} seja comutativo.
Suponhamos também que $q$ seja próprio e que sua restrição a $V\setmin j(X)$ seja um isomorfismo sobre $C\setmin\varepsilon(Y)$.
Suponhamos, além disso, que, se $p=(p_0,\theta)$, então $\theta:\sh{O}_Y\to p_*(\sh{O}_X)$ seja um isomorfismo.
Sejam $\sh{S}'=\bigoplus_{n\geq0}p_*(\sh{L}^{\otimes n})$ e $C'=\Spec(\sh{S}')$, e seja $q':\bb{V}(\sh{L})\to C'$ o $Y$-morfismo canônico \sref{II.8.8.5}.
Então o único $Y$-morfismo $u:C'\to C$ tal que $q=u\circ q'$ é um isomorfismo.
\end{proposition}

\begin{proof}
Decorre de \sref{II.8.8.6} que $C'$ é íntegro;
como $u$ é um homeomorfismo dos espaços subjacentes $C'\to C$ (pois $u$ é bijetivo e fechado por \sref{II.8.11.6}), $C$ é irredutível, portanto íntegro, e, visto que a restrição de $u$ a um aberto não vazio de $C'$ é um isomorfismo sobre um aberto de $C$, $u$ é birracional.
Como $C$ é, por hipótese, normal, basta aplicar \sref{II.8.11.2} para obter a conclusão.
\end{proof}

\begin{remarks}[8.11.8]
\label{II.8.11.8}
\begin{enumerate}
  \item[\rm{(i)}] Observemos que a hipótese de que $C$ é normal implica que o mesmo vale para $X$.
  Com efeito, $C'=\Spec(\sh{S}')$ é então normal, por ser isomorfo a $C$, e íntegro em virtude de \sref{II.8.8.6};
  concluímos daí que $\Proj(\sh{S}')$ é \emph{normal}.
  Com efeito, a questão é local sobre $Y$;
  se $Y$ é afim, $\sh{S}'=\widetilde{S'}$, o anel $S'=\Gamma(C',\sh{O}_{C'})$ é íntegro e integralmente fechado \sref{II.8.8.6.1}; portanto, para todo elemento homogêneo $f\in S'_+$, o anel graduado $S'_f$ é íntegro e integralmente fechado \cite[t.~I, p.~257 e 261]{I-13}, e, por conseguinte, o mesmo vale para o anel $S'_{(f)}$ de seus termos de grau zero, pois a interseção de $S'_f$ com o corpo de frações de $S'_{(f)}$ é igual a $S'_{(f)}$;
  isto prova nossa afirmação \sref{II.6.3.4}.
  Por fim, como $X$ é isomorfo a um subpré-esquema aberto de $\Proj(\sh{S}')$ \sref{II.8.8.1}, $X$ é de fato normal.
  Podemos, portanto, exprimir a proposição \sref{II.8.11.7} da seguinte forma:
  \emph{Se $X$ é íntegro e normal, $p=(p_0,\theta):X\to Y$ é um morfismo próprio tal que $\theta:\sh{O}_Y\to p_*(\sh{O}_X)$ seja um isomorfismo, então, para todo $\sh{O}_X$-módulo $\sh{L}$ amplo relativamente a $p$, existe uma única maneira de contrair a seção nula de $V=\bb{V}(\sh{L})$ de modo a obter um $Y$-esquema normal $C$ e um $Y$-morfismo próprio $q:V\to C$.}
  \item[\rm{(ii)}] Quando $p$ é próprio, a hipótese $p_*(\sh{O}_X)=\sh{O}_Y$ pode ser considerada uma hipótese auxiliar, que não constitui uma verdadeira restrição à generalidade do resultado.
  Com efeito, se ela não for verificada, basta substituir $Y$ pelo $Y$-esquema $Y'=\Spec(p_*(\sh{O}_X))$ e considerar $X$ como um $Y'$-esquema.
  Voltaremos a este método geral no Capítulo~III, \S~4.
\end{enumerate}
\end{remarks}

\subsection{Feixes quase-coerentes sobre os cones projetantes}
\label{subsection:II.8.12}

\begin{env}[8.12.1]
\label{II.8.12.1}
Retomemos as hipóteses e notações de \sref{II.8.3.1}.
Seja $\sh{M}$ um $\sh{S}$-módulo graduado quase-coerente; para evitar qualquer confusão, designaremos por $\widetilde{\sh{M}}$ o $\sh{O}_C$-módulo
\oldpage[II]{192}
quase-coerente associado a $\sh{M}$ \sref{II.1.4.3}, quando $\sh{M}$ é considerado como um $\sh{S}$-módulo \emph{não graduado}, e por $\shProj_0(\sh{M})$ o $\sh{O}_X$-módulo quase-coerente associado a $\sh{M}$, sendo $\sh{M}$ desta vez considerado como um $\sh{S}$-módulo graduado (em outras palavras, o $\sh{O}_X$-módulo denotado por $\widetilde{\sh{M}}$ em \sref{II.3.2.2}).
Poremos, além disso,
\[
\label{II.8.12.1.1}
  \sh{M}_X=\shProj(\sh{M})=\bigoplus_{n\in\bb{Z}}\shProj_0(\sh{M}(n));
\tag{8.12.1.1}
\]
estando a $\sh{O}_X$-álgebra graduada quase-coerente $\sh{S}_X$ definida por \sref{II.8.6.1.1}, $\shProj(\sh{M})$ é munido de uma estrutura de $\sh{S}_X$-módulo graduado (quase-coerente), por meio dos homomorfismos canônicos \sref{II.3.2.6.1}
\[
\label{II.8.12.1.2}
  \sh{O}_X(m)\otimes_{\sh{O}_X}\shProj_0(\sh{M}(n))\to\shProj_0(\sh{S}(m)\otimes_\sh{S}\sh{M}(n))\to\shProj_0(\sh{M}(m+n)),
\tag{8.12.1.2}
\]
sendo a verificação dos axiomas de módulo feita com o auxílio do diagrama comutativo de \sref{II.2.5.11.4}.

Se $Y=\Spec(A)$ é afim, $\sh{S}=\widetilde{S}$ e $\sh{M}=\widetilde{M}$, onde $S$ é uma $A$-álgebra graduada e $M$ um $S$-módulo graduado, então, para todo elemento homogêneo $f\in S_+$, tem-se
\[
\label{II.8.12.1.3}
  \Gamma(X_f,\shProj(\widetilde{M}))=M_f
\tag{8.12.1.3}
\]
em virtude das definições e de \sref{II.8.2.9.1}.

Consideremos agora o $\widehat{\sh{S}}$-módulo graduado quase-coerente
\[
\label{II.8.12.1.4}
  \widehat{\sh{M}}=\sh{M}\otimes_\sh{S}\widehat{\sh{S}}
\tag{8.12.1.4}
\]
($\widehat{\sh{S}}$ sendo definido por \sref{II.8.3.1.1}); dele se deduz um $\sh{O}_{\widehat{C}}$-módulo graduado quase-coerente $\shProj_0(\widehat{\sh{M}})$, que denotaremos também por
\[
\label{II.8.12.1.5}
  \sh{M}^\square=\shProj_0(\widehat{\sh{M}}).
\tag{8.12.1.5}
\]
É claro \sref{II.3.2.4} que $\sh{M}^\square$ é um funtor aditivo exato em $\sh{M}$, que comuta com as somas diretas e com os limites indutivos.
\end{env}

\begin{proposition}[8.12.2]
\label{II.8.12.2}
Com as notações de \sref{II.8.3.2}, têm-se isomorfismos canônicos funtoriais
\[
\label{II.8.12.2.1}
  i^*(\sh{M}^\square)\simto\widetilde{\sh{M}},
  \qquad
  j^*(\sh{M}^\square)\simto\shProj_0(\sh{M}).
\tag{8.12.2.1}
\]
\end{proposition}

\begin{proof}
Com efeito, $i^*(\sh{M}^\square)$ se identifica canonicamente com $(\widehat{\sh{M}}/(\bb{z}-1)\widehat{\sh{M}})\supertilde$ sobre $\Spec(\widehat{\sh{S}}/(\bb{z}-1)\widehat{\sh{S}})$ em virtude de \sref{II.3.2.3};
o primeiro dos isomorfismos canônicos de \sref{II.8.12.2.1} deduz-se então imediatamente \sref{II.1.4.1} do isomorfismo canônico $\widehat{\sh{M}}/(\bb{z}-1)\widehat{\sh{M}}\simto\sh{M}$.
Por outro lado, a imersão canônica $j:X\to\widehat{C}$ corresponde ao homomorfismo canônico $\widehat{\sh{S}}\to\sh{S}$ de núcleo $\bb{z}\widehat{\sh{S}}$ \sref{II.8.3.2};
o segundo homomorfismo de \sref{II.8.12.2.1} é o caso particular do homomorfismo canônico de \sref{II.3.5.2}[ii], pelo fato de que se tem aqui $\widehat{\sh{M}}\otimes_{\widehat{\sh{S}}}\sh{S}=\sh{M}$;
para verificar que se trata de um isomorfismo, podemos nos limitar ao caso em que $Y=\Spec(A)$ é afim, $\sh{S}=\widetilde{S}$ e $\sh{M}=\widetilde{M}$;
reportando-nos a \sref{II.2.8.8}, a verificação de que, para todo $f$ homogêneo em $S_+$, o homomorfismo precedente, restrito a $X_f$, se reduz a um isomorfismo é então imediata.
\end{proof}
\oldpage[II]{193}

Por abuso de linguagem, diremos ainda, em razão da existência do primeiro isomorfismo de \sref{II.8.12.2.1}, que $\sh{M}^\square$ é o \emph{fecho projetivo} do $\sh{O}_C$-módulo $\widetilde{\sh{M}}$ (ficando subentendido que a graduação do $\sh{S}$-módulo $\sh{M}$ faz parte dos dados do $\sh{O}_C$-módulo $\widetilde{\sh{M}}$).

\begin{env}[8.12.3]
\label{II.8.12.3}
Com as notações de \sref{II.8.3.5}, tem-se um homomorfismo canônico funtorial
\[
\label{II.8.12.3.1}
  p^*(\shProj_0(\sh{M}))\to\sh{M}^\square|\widehat{E}.
\tag{8.12.3.1}
\]
Com efeito, trata-se de um caso particular do homomorfismo $\nu^\sharp$ definido de maneira geral em \sref{II.3.5.6}.
Se $Y=\Spec(A)$ é afim, $\sh{S}=\widetilde{S}$ e $\sh{M}=\widetilde{M}$, vê-se, reportando-se a \sref{II.2.8.8}, que a restrição de \sref{II.8.12.3.1} a $p^{-1}(X_f)=\widehat{C}_f$ (para um $f$ homogêneo em $S_+$) corresponde ao homomorfismo canônico
\[
\label{II.8.12.3.2}
  M_{(f)}\otimes_{S_{(f)}}S_f^\leq\to M_f^\leq
\tag{8.12.3.2}
\]
levando em conta \sref{II.8.2.3.2} e \sref{II.8.2.5.2}.
\end{env}

\begin{env}[8.12.4]
\label{II.8.12.4}
Coloquemo-nos nas condições de \sref{II.8.5.1}, cujas hipóteses e notações conservamos.
Decorre de \sref{II.1.5.6} que, para todo $\sh{S}$-módulo graduado quase-coerente $\sh{M}$, tem-se, por um lado, um isomorfismo canônico
\[
\label{II.8.12.4.1}
  \Phi^*(\widetilde{\sh{M}})\simto(q^*(\sh{M})\otimes_{q^*(\sh{S})}\sh{S}')\supertilde
\tag{8.12.4.1}
\]
de $\sh{O}_{C'}$-módulos;
e, por outro lado, \sref{II.3.5.6} acarreta a existência de um $\Proj(\vphi)$-morfismo canônico
\[
\label{II.8.12.4.2}
  \shProj_0(\sh{M})\to\shProj_0(q^*(\sh{M})\otimes_{q^*(\sh{S})}\sh{S}')|G(\vphi)
\tag{8.12.4.2}
\]
e também de um $\widehat{\Phi}$-morfismo canônico
\[
\label{II.8.12.4.3}
  \shProj_0(\widehat{\sh{M}})\to\shProj_0(q^*(\widehat{\sh{M}})\otimes_{q^*(\widehat{\sh{S}})}\widehat{\sh{S}'})|G(\widehat{\vphi}).
\tag{8.12.4.3}
\]
\end{env}

\begin{env}[8.12.5]
\label{II.8.12.5}
Consideremos agora a situação de \sref{II.8.6.1}, com as mesmas notações;
tomamos, portanto, no que precede, $Y'=X$, sendo o morfismo $q:X\to Y$ o morfismo estrutural, e $\vphi$ o $q$-morfismo canônico \sref{II.8.6.1.2}.
Tem-se então um isomorfismo canônico
\[
\label{II.8.12.5.1}
  q^*(\sh{M})\otimes_{q^*(\sh{S})}\sh{S}_X^\geq\simto\sh{M}_X^\geq
\tag{8.12.5.1}
\]
pondo $\sh{M}_X^\geq=\bigoplus_{n\geq0}\shProj_0(\sh{M}(n))$.
Podemos, com efeito, nos limitar ao caso em que $Y=\Spec(A)$ é afim, $\sh{S}=\widetilde{S}$ e $\sh{M}=\widetilde{M}$, e definir o isomorfismo de \sref{II.8.12.5.1} em cada um dos abertos afins $X_f$ ($f$ homogêneo em $S_+$), verificando a compatibilidade com a passagem a um múltiplo homogêneo de $f$.
Ora, a restrição a $X_f$ do primeiro membro de \sref{II.8.12.5.1} é $\widetilde{M}'=((M\otimes_A S_{(f)})\otimes_{S\otimes_A S_{(f)}}S_f^\geq)\supertilde$ por \sref{II.8.6.2.1};
como se tem um isomorfismo canônico de $M\otimes_A S_{(f)}$ sobre $M\otimes_S(S\otimes_A S_{(f)})$, deduz-se um isomorfismo de $\widetilde{M}'$ sobre $(M\otimes_S S_f^\geq)\supertilde$, e este último é canonicamente isomorfo, por \sref{II.8.2.9.1}, à restrição a $X_f$ do segundo membro de \sref{II.8.12.5.1}, e satisfaz às condições de compatibilidade requeridas.
\oldpage[II]{194}

Substituindo $\sh{M}$ por $\widehat{\sh{M}}$, $\sh{S}$ por $\widehat{\sh{S}}$ e $\sh{S}_X$ por $(\sh{S}_X^\geq)^\wedge$ no raciocínio precedente, tem-se igualmente um isomorfismo canônico
\[
\label{II.8.12.5.2}
  q^*(\widehat{\sh{M}})\otimes_{q^*(\widehat{\sh{S}})}(\sh{S}_X^\geq)^\wedge\simto(\sh{M}_X^\geq)^\wedge.
\tag{8.12.5.2}
\]
Se lembrarmos \sref{II.8.6.2} que o morfismo estrutural $u:\Proj(\sh{S}_X^\geq)\to X$ é um isomorfismo, deduzimos primeiramente do que precede que se tem um $u$-isomorfismo canônico
\[
\label{II.8.12.5.3}
  \shProj_0\sh{M}\simto\shProj_0(\sh{M}_X^\geq)
\tag{8.12.5.3}
\]
como caso particular de \sref{II.8.12.4.2}.
Constata-se, com efeito, com as notações da demonstração de \sref{II.8.6.2}, que isso equivale a ver que o homomorfismo canônico $M_{(f)}\otimes_{S_{(f)}}(S_f^\geq)^{(d)}\to(M_f^\geq)^{(d)}$ é um isomorfismo quando $f\in S_d$, o que é imediato.

Em segundo lugar, o isomorfismo de \sref{II.8.12.5.2} permite, aplicando desta vez \sref{II.8.12.4.3} ao morfismo canônico $r=\Proj(\widehat{\alpha}):\widehat{C}_X\to\widehat{C}$, obter um $r$-morfismo canônico
\[
\label{II.8.12.5.4}
  \sh{M}^\square\to(\sh{M}_X^\geq)^\square.
\tag{8.12.5.4}
\]
Lembremos agora \sref{II.8.6.2} que as restrições de $r$ aos cones sem vértice $\widehat{E}_X$ e $E_X$ são \emph{isomorfismos} sobre $\widehat{E}$ e $E$, respectivamente.
Além disso:
\end{env}

\begin{proposition}[8.12.6]
\label{II.8.12.6}
As restrições a $\widehat{E}_X$ e a $E_X$ do $r$-morfismo canônico de \sref{II.8.12.5.4} são isomorfismos
\[
\label{II.8.12.6.1}
  \sh{M}^\square|\widehat{E}\simto(\sh{M}_X^\geq)^\square|\widehat{E}_X
\tag{8.12.6.1}
\]
\[
\label{II.8.12.6.2}
  \widetilde{\sh{M}}|E\simto(\sh{M}_X^\geq)\supertilde|E_X.
\tag{8.12.6.2}
\]
\end{proposition}

\begin{proof}
Reduzimo-nos ao caso em que $Y$ é afim, como na demonstração de \sref{II.8.6.2}, cujas notações adotamos;
reportando-nos às definições \sref{II.2.8.8}, devemos mostrar que o homomorfismo canônico
\[
  \widehat{M}_{(f)}\otimes_{\widehat{S}_{(f)}}(S_f^\geq)_{(f/1)}^\wedge\to(M\otimes_S S_f^\geq)_{(f/1)}^\wedge
\]
é um isomorfismo;
mas, em virtude de \sref{II.8.2.3.2} e \sref{II.8.2.5.2}, o primeiro membro se identifica canonicamente com $M_f^\leq\otimes_{S_f^\leq}(S_f^\geq)_{f/1}^\leq$, portanto com $M_f^\leq$ em virtude de \sref{II.8.2.7.2}, e o segundo com $(M_f^\geq)_{f/1}^\leq$, portanto também com $M_f^\leq$ em virtude de \sref{II.8.2.9.2}, de onde a conclusão no que diz respeito a \sref{II.8.12.6.1};
a fórmula \sref{II.8.12.6.2} decorre então de \sref{II.8.12.6.1} e de \sref{II.8.12.2.1}.
\end{proof}

\begin{corollary}[8.12.7]
\label{II.8.12.7}
Com as identificações de \sref{II.8.6.3}, a restrição de $(\sh{M}_X^\geq)^\square$ a $\widehat{E}_X$ se identifica com $(\sh{M}_X^\leq)\supertilde$, e a restrição de $(\sh{M}_X^\geq)^\square$ a $E_X$ se identifica com $\widetilde{\sh{M}}_X$.
\end{corollary}

\begin{proof}
Com efeito, reduzimo-nos ao caso afim, e isso decorre da identificação de $(M_f^\geq)_{f/1}^\leq$ com $M_f^\leq$ e de $(M_f^\geq)_{f/1}$ com $M_f$ \sref{II.8.2.9.2}.
\end{proof}

\begin{proposition}[8.12.8]
\label{II.8.12.8}
Sob as hipóteses de \sref{II.8.6.4}, o homomorfismo canônico de \sref{II.8.12.3.1} é um isomorfismo.
\end{proposition}

\begin{proof}
Levando em conta o fato de que $\Proj(\sh{S}_X^\geq)\to X$ é um isomorfismo \sref{II.8.6.2}, e os
\oldpage[II]{195}
isomorfismos de \sref{II.8.12.5.3} e \sref{II.8.12.6.1}, ficamos reduzidos a demonstrar a proposição correspondente para o homomorfismo canônico $p_X^*(\shProj_0(\sh{M}_X^\geq))\to(\sh{M}_X^\geq)^\square|\widehat{E}_X$; em outras palavras, ficamos reduzidos ao caso em que $\sh{S}_1$ é um $\sh{O}_Y$-módulo invertível e em que $\sh{S}$ é gerada por $\sh{S}_1$.
Com as notações de \sref{II.8.12.3}, tem-se então, para um $f\in S_1$, $S_f^\leq=S_{(f)}[1/f]$, e o homomorfismo canônico $M_{(f)}\otimes_{S_{(f)}}S_f^\leq\to M_f^\leq$ é um isomorfismo pela definição de $M_f^\leq$.
\end{proof}

\begin{env}[8.12.9]
\label{II.8.12.9}
Consideremos agora os $\sh{S}$-módulos quase-coerentes
\[
  \sh{M}_{[n]}=\bigoplus_{m\geq n}\sh{M}_m
\]
e (com as notações de \sref{II.8.7.2}) o $\sh{S}^\natural$-módulo graduado quase-coerente
\[
\label{II.8.12.9.1}
  \sh{M}^\natural=\left(\bigoplus_{n\geq0}\sh{M}_{[n]}\right)\supertilde.
\tag{8.12.9.1}
\]
Vimos em \sref{II.8.7.3} que existe um $C$-isomorfismo canônico $h:C_X\simto\Proj(\sh{S}^\natural)$.
Além disso:
\end{env}

\begin{proposition}[8.12.10]
\label{II.8.12.10}
Existe um $h$-isomorfismo canônico
\[
\label{II.8.12.10.1}
  \shProj_0(\sh{M}^\natural)\simto\widetilde{\sh{M}}_X.
\tag{8.12.10.1}
\]
\end{proposition}

\begin{proof}
Raciocina-se como em \sref{II.8.7.3}, utilizando desta vez a existência do di-isomorfismo de \sref{II.8.2.9.3} em vez do de \sref{II.8.2.7.3}.
Deixamos os detalhes ao leitor.
\end{proof}

\subsection{Fecho projetivo de subfeixes e subesquemas fechados}
\label{subsection:II.8.13}

\begin{env}[8.13.1]
\label{II.8.13.1}
Sendo as hipóteses e notações as de \sref{II.8.12.1}, consideremos um sub-$\sh{S}$-módulo quase-coerente $\sh{N}$ de $\sh{M}$, \emph{não necessariamente graduado}.
Podemos, portanto, considerar o $\sh{O}_C$-módulo quase-coerente $\widetilde{\sh{N}}$ associado a $\sh{N}$, que é um sub-$\sh{O}_C$-módulo de $\widetilde{\sh{M}}$.
Vimos, por outro lado, em \sref{II.8.12.2.1}, que $\widetilde{\sh{M}}$ se identifica com a restrição de $\sh{M}^\square$ a $C$.
Como a injeção canônica $i:C\to\widehat{C}$ é um morfismo afim \sref{II.8.3.2} e, \emph{a fortiori}, quase-compacto, o \emph{prolongamento canônico} $(\widetilde{\sh{N}})^-$, maior sub-$\sh{O}_{\widehat{C}}$-módulo contido em $\sh{M}^\square$ e que induz $\widetilde{\sh{N}}$ sobre $C$, é um $\sh{O}_{\widehat{C}}$-módulo quase-coerente \sref[I]{I.9.4.2}.
Daremos uma descrição explícita dele com o auxílio de um $\widehat{\sh{S}}$-módulo graduado.
\end{env}

\begin{env}[8.13.2]
\label{II.8.13.2}
Para isso, consideremos, para todo inteiro $n\geq0$, o homomorfismo $\bigoplus_{i\leq n}\sh{M}_i\to\sh{M}$ que, para todo aberto $U$ de $Y$, associa à família
\[
  (s_i)\in\bigoplus_{i\leq n}\Gamma(U,\sh{M}_i)
\]
a seção $\sum_i s_i\in\Gamma(U,\sh{M})$.
Designemos por $\sh{N}'_n$ a imagem inversa de $\sh{N}$ por esse homomorfismo, que é um sub-$\sh{O}_Y$-módulo quase-coerente de $\bigoplus_{i\leq n}\sh{M}_i$.
Consideremos agora o homomorfismo $\bigoplus_{i\leq n}\sh{M}_i\to\widehat{\sh{M}}=\sh{M}[\bb{z}]$ que associa a $(s_i)$ a seção $\sum_{i\leq n}s_i\bb{z}^{n-i}\in\Gamma(U,\widehat{\sh{M}}_n)$, e seja $\sh{N}_n$ a imagem de $\sh{N}'_n$ por esse homomorfismo;
verifica-se imediatamente que $\overline{\sh{N}}=\bigoplus_{n\geq0}\sh{N}_n$ é um sub-$\widehat{\sh{S}}$-módulo graduado (quase-coerente) de $\widehat{\sh{M}}$;
diz-se que $\overline{\sh{N}}$ é deduzido de $\sh{N}$ por \emph{homogeneização}, mediante a «variável homogeneizante» $\bb{z}$.
Observaremos
\oldpage[II]{196}
que, se $\sh{N}$ já é um sub-$\sh{S}$-módulo \emph{graduado} de $\sh{M}$, então $\overline{\sh{N}}$ se identifica com a soma direta das componentes $\widehat{\sh{N}}_n$ de grau $n\geq0$ em $\widehat{\sh{N}}=\sh{N}[\bb{z}]$.
\end{env}

\begin{proposition}[8.13.3]
\label{II.8.13.3}
O $\sh{O}_{\widehat{C}}$-módulo $\shProj_0(\overline{\sh{N}})$ é o prolongamento canônico $(\widetilde{\sh{N}})^-$ de $\widetilde{\sh{N}}$ a $\widehat{C}$.
\end{proposition}

\begin{proof}
A questão é local em $Y$ e $\widehat{C}$, em virtude da definição do prolongamento canônico \sref[I]{I.9.4.1}.
Podemos, portanto, supor desde já que $Y=\Spec(A)$ seja afim, com $\sh{S}=\widetilde{S}$, $\sh{M}=\widetilde{M}$ e $\sh{N}=\widetilde{N}$, onde $N$ é um sub-$S$-módulo de $M$ não necessariamente graduado.
Além disso, por \sref{II.8.3.2.6}, $\widehat{C}$ é a união dos abertos afins $\widehat{C}_z=C$ e $\widehat{C}_f=\Spec(S_f^\leq)$ (com $f$ homogêneo em $S_+$).
Basta, pois, mostrar que: (1) a restrição de $\shProj_0(\overline{\sh{N}})$ a $C$ é $\widetilde{\sh{N}}$; (2) a restrição de $\shProj_0(\overline{\sh{N}})$ a cada $\widehat{C}_f$ é o prolongamento canônico da restrição de $\widetilde{\sh{N}}$ a $C\cap\widehat{C}_f=\Spec(S_f)$ \sref{II.8.3.2.6}.
Para o primeiro ponto, observemos que $\shProj_0(\overline{\sh{N}})|C$ se identifica com $(\overline{N}_{(\bb{z})})\supertilde$ \sref{II.8.3.2.4};
mas $\overline{N}_{(\bb{z})}$ se identifica canonicamente \sref{II.2.2.5} com a imagem de $\overline{N}$ em $\widehat{M}/(\bb{z}-1)\widehat{M}$ e, pelo isomorfismo canônico deste último com $M$ \sref{II.8.2.5}, essa imagem se identifica com $N$, pela definição de $\overline{N}$ dada em \sref{II.8.13.2}.

Para provar o segundo ponto, observemos que a injeção $i:C\cap\widehat{C}_f\to\widehat{C}_f$ corresponde à injeção canônica $S_f^\leq\to S_f$ \sref{II.8.3.2.6};
temos também $\Gamma(\widehat{C}_f,\sh{M}^\square)=M_f^\leq$, $\Gamma(\widehat{C}_f,i_*(\widetilde{\sh{N}}))=N_f$ e, por \sref{II.8.12.2.1}, $\Gamma(\widehat{C}_f,i_*(i^*(\sh{M}^\square)))=M_f$.
Tendo em conta \sref[I]{I.9.4.2}, ficamos, portanto, reduzidos a mostrar que $\overline{N}_{(f)}\subset\widehat{M}_{(f)}=M_f^\leq$ se identifica canonicamente com a imagem inversa de $N_f$ pela injeção canônica $M_f^\leq\to M_f$.
Com efeito, seja $d=\deg(f)>0$, e suponhamos que um elemento $(\sum_{k\leq md}x_k)/f^m$ de $M_f$ (com $x_k\in M_k$) seja da forma $y/f^m$, com $y\in N$.
Multiplicando $y$ e os $x_k$ por uma mesma potência adequada $f^h$, podemos desde já supor que $\sum_{k\leq md}x_k=y$.
Mas, na identificação de \sref{II.8.2.5.2}, $(\sum_{k\leq md}x_k)/f^m$ corresponde a $\sum_{k\leq md}x_k\bb{z}^{md-k}/f^m$, e este é de fato um elemento de $\overline{N}_{(f)}$, pois $\sum_{k\leq md}x_k\in N$;
a recíproca é evidente.
\end{proof}

\begin{remark}[8.13.4]
\label{II.8.13.4}
\begin{enumerate}
    \item[\rm{(i)}] O caso de aplicação mais importante de \sref{II.8.13.3} é aquele em que $\sh{M}=\sh{S}$, sendo então $\widetilde{\sh{N}}$ um feixe \emph{arbitrário} quase-coerente de ideais $\sh{J}$ de $\sh{O}_C$ \sref{II.1.4.3}, correspondente bijetivamente a um \emph{subpré-esquema fechado} $Z$ de $C$.
        Então o prolongamento canônico $(\sh{J})^-$ de $\sh{J}$ é o feixe quase-coerente de ideais de $\sh{O}_{\widehat{C}}$ que define o \emph{fecho} $\overline{Z}$ de $Z$ em $\widehat{C}$ \sref[I]{I.9.5.10};
        a Proposição~\sref{II.8.13.3} fornece uma maneira canônica de definir $\overline{Z}$ mediante um ideal graduado em $\widehat{\sh{S}}=\sh{S}[\bb{z}]$.
    \item[\rm{(ii)}] Suponhamos, para simplificar, que $Y$ seja afim e adotemos a notação da demonstração de \sref{II.8.13.3}.
        Para cada $x\in N$ não nulo, seja $d(x)$ o maior grau das componentes homogêneas $x_i$ de $x$ em $M$;
        por definição, $\overline{N}$ é o submódulo de $\widehat{M}$ formado por $0$ e pelos elementos da forma $h(x,k)=\bb{z}^k\sum_{i\leq d(x)}x_i\bb{z}^{d(x)-i}$ (para $k\geq0$ inteiro);
        portanto, é gerado, como módulo sobre $\widehat{S}=S[\bb{z}]$, pelos elementos da forma
        \[
            h(x,0) = \sum_{i\leq d(x)}x_i\bb{z}^{d(x)-i}.
        \]
\oldpage[II]{197}
        Diz-se que $h(x,0)$ é deduzido de $x$ por \emph{homogeneização}, mediante a «variável homogeneizante» $\bb{z}$.
        Mas, como $h(x,0)$ não depende aditivamente de $x$ (nem, \emph{a fortiori}, $S$-linearmente), \emph{evitaremos supor} (mesmo quando $M=S$) que os $h(x,0)$ formam um \emph{sistema de geradores} do $\widehat{S}$-módulo graduado $\overline{N}$ quando $x$ percorre um \emph{sistema de geradores} do $S$-módulo $N$.
        Isso ocorre, contudo, quando $N$ é um $S$-módulo \emph{cíclico livre} (o único caso considerado na geometria algébrica elementar), pois, se $t$ é uma base de $N$, então $h(t,0)$ gera o $\widehat{S}$-módulo $\overline{N}$.
\end{enumerate}
\end{remark}


\subsection{Complementos sobre os feixes associados a $\mathcal{S}$-módulos graduados}
\label{subsection:II.8.14}

\begin{env}[8.14.1]
\label{II.8.14.1}
Sejam $Y$ um pré-esquema, $\sh{S}$ uma $\sh{O}_Y$-álgebra quase-coerente \emph{graduada em graus positivos}, $X=\Proj(\sh{S})$ e $q:X\to Y$ o morfismo estrutural (que é separado, por \sref{II.3.1.3}).
Com a notação de \sref{II.8.12.1}, definimos o funtor $\sh{M}\mapsto\sh{M}_X=\shProj(\sh{M})$ da categoria dos $\sh{S}$-módulos graduados quase-coerentes para a categoria dos $\sh{S}_X$-módulos graduados quase-coerentes;
é também claro \sref{II.3.2.4} que esse funtor é \emph{aditivo} e \emph{exato}, e que comuta com limites indutivos.

Observemos, além disso, que decorre imediatamente da definição \sref{II.8.12.1.1} que
\[
\label{II.8.14.1.1}
    \shProj(\sh{M}(n)) = (\shProj(\sh{M}))(n)
    \quad\mbox{para todo $n\in\bb{Z}$.}
\tag{8.14.1.1}
\]
\end{env}

\begin{env}[8.14.2]
\label{II.8.14.2}
Vamos primeiro estender os homomorfismos canônicos $\lambda$ e $\mu$, definidos em \sref{II.3.2.6}, aos $\sh{S}_X$-módulos da forma $\shProj(\sh{M})$.
Para isso, observemos que, para quaisquer $m\in\bb{Z}$ e $n\in\bb{Z}$, temos, por \sref{II.2.1.2.1}, um homomorfismo canônico de $\sh{O}_X$-módulos
\[
\label{II.8.14.2.1}
    \lambda_{mn}:
    \shProj_0(\sh{M}(m))\otimes_{\sh{O}_X}\shProj_0(\sh{N}(n))
    \to
    \shProj_0((\sh{M}\otimes_{\sh{S}}\sh{N})(m+n)),
\tag{8.14.2.1}
\]
e, igualmente, por \sref{II.2.1.2.2}, um homomorfismo canônico de $\sh{O}_X$-módulos
\[
\label{II.8.14.2.2}
    \mu_{mn}:
    \shProj_0((\shHom_{\sh{S}}(\sh{M},\sh{N}))(n-m))
    \to
    \shHom_{\sh{O}_X}(\shProj_0(\sh{M}(m)),\shProj_0(\sh{N}(n)))
\tag{8.14.2.2}
\]
para quaisquer $\sh{S}$-módulos graduados quase-coerentes $\sh{M}$ e $\sh{N}$.
Isso induz um homomorfismo
\[
    \mu_k:
    \shProj_0((\shHom_{\sh{S}}(\sh{M},\sh{N}))(k))
    \to
    (\shHom_{\sh{S}_X}(\shProj(\sh{M}),\shProj(\sh{N})))_k
\]
que envia cada $u\in\Gamma(U,\shProj_0((\shHom_{\sh{S}}(\sh{M},\sh{N}))(k)))$ ao homomorfismo $\mu_k(u)$, de grau $k$, de $\bb{Z}$-módulos graduados $\Gamma(U,\shProj(\sh{M}))\to\Gamma(U,\shProj(\sh{N}))$ (onde $U$ é aberto em $X$), cuja restrição a cada $\Gamma(U,\shProj_0(\sh{M}(m)))$ coincide com $\mu_{m,m+k}(u)$;
além disso, voltando à definição dos $\mu_{mn}$ \sref{II.2.5.12.1}, vê-se imediatamente que $\mu_k(u)$ é de fato um homomorfismo de grau $k$ de $\Gamma(U,\sh{S}_X)$-módulos graduados e que, ademais, os $\mu_k$ definem um homomorfismo de \emph{$\sh{S}_X$-módulos graduados}
\[
\label{II.8.14.2.3}
    \shProj(\shHom_{\sh{S}}(\sh{M},\sh{N}))
    \to
    \shHom_{\sh{S}_X}(\shProj(\sh{M}),\shProj(\sh{N})).
\tag{8.14.2.3}
\]

Do mesmo modo, levando em conta o diagrama de associatividade \sref{II.2.5.11.4}, os homomorfismos \sref{II.8.14.2.1} fornecem um homomorfismo de \emph{$\sh{S}_X$-módulos graduados}
\[
\label{II.8.14.2.4}
    \lambda:
    \shProj(\sh{M})\otimes_{\sh{S}_X}\shProj(\sh{N})
    \to
    \shProj(\sh{M}\otimes_{\sh{S}}\sh{N}).
\tag{8.14.2.4}
\]
\end{env}
\oldpage[II]{198}

\begin{proposition}[8.14.3]
\label{II.8.14.3}
O homomorfismo \sref{II.8.14.2.4} é bijetivo;
o mesmo vale para \sref{II.8.14.2.3} sempre que o $\sh{S}$-módulo graduado $\sh{M}$ admita uma apresentação finita \sref{II.3.1.1}.
\end{proposition}

\begin{proof}
A questão é claramente local em $X$ e $Y$;
podemos, portanto, supor que $Y=\Spec(A)$ seja afim, com $\sh{S}=\widetilde{S}$, $\sh{M}=\widetilde{M}$ e $\sh{N}=\widetilde{N}$, onde $S$ é uma $A$-álgebra graduada em graus positivos e $M$ e $N$ são $S$-módulos graduados.
Se $f$ é um elemento homogêneo de $S_+$, os homomorfismos \sref{II.8.14.2.1} e \sref{II.8.14.2.2}, restritos ao aberto afim $D_+(f)$, correspondem aos homomorfismos canônicos \sref{II.2.5.11.1} e \sref{II.2.5.12.1}:
\begin{align*}
    M(m)_{(f)}\otimes_{S_{(f)}}N(n)_{(f)}
    &\to
    (M\otimes_S N)(m+n)_{(f)}
\\  (\Hom_S(M,N))(n-m)_{(f)}
    &\to
    \Hom_{S_{(f)}}(M(m)_{(f)},N(n)_{(f)}).
\end{align*}

Se nos reportarmos às definições desses homomorfismos, veremos então (levando em conta \sref{II.8.2.9.1}) que a restrição de \sref{II.8.14.2.4} a $D_+(f)$ corresponde ao homomorfismo canônico
\[
    M_f\otimes_{S_f}N_f \to (M\otimes_S N)_f
\]
definido em \sref[0]{0.1.3.4}, e sabemos que este último homomorfismo é um isomorfismo.
Do mesmo modo, a restrição de \sref{II.8.14.2.3} a $D_+(f)$ corresponde ao homomorfismo canônico \sref[0]{0.1.3.5}
\[
    (\Hom_S(M,N))_f \to \Hom_{S_f}(M_f,N_f)
\]
levando em conta que, como $M$ é de tipo finito, o módulo $\Hom_S(M,N)$, soma direta dos subgrupos formados pelos homomorfismos \emph{homogêneos} de $S$-módulos \sref{II.2.1.2}, coincide com o conjunto de \emph{todos} os homomorfismos $M\to N$ de $S$-módulos.
A hipótese de que $M$ admite uma apresentação finita implica então \sref[0]{0.1.3.5} que o homomorfismo canônico considerado é, de fato, um isomorfismo.
\end{proof}

\begin{proposition}[8.14.4]
\label{II.8.14.4}
Se $U$ é um aberto quase-compacto de $X$, existe um inteiro $d$ tal que, para todo inteiro $n$ múltiplo de $d$, $\sh{O}_X(n)|U$ é invertível, tendo por inverso $\sh{O}_X(-n)|U$.
\end{proposition}

\begin{proof}
Como $q(U)$ é quase-compacto, ele é recoberto por um número finito de abertos afins $V_i$, e assim todo $x\in U$ está contido em algum aberto afim da forma $D_+(f)$, onde $f$ é um elemento homogêneo de grau $>0$ de um dos anéis $\Gamma(V_i,\sh{S})$.
Como $U$ é quase-compacto, podemos recobri-lo por um número finito de tais abertos $D_+(f_j)$;
seja $d$ um múltiplo comum dos graus dos $f_j$.
Esse $d$ tem a propriedade requerida, por \sref{II.2.5.17}.
\end{proof}

\begin{env}[8.14.5]
\label{II.8.14.5}
Sob as hipóteses e com a notação de \sref{II.8.14.1}, definimos, em \sref{II.3.3.2}, homomorfismos canônicos de $\sh{O}_Y$-módulos
\[
\label{II.8.14.5.1}
    \alpha_n: \sh{M}_n \to q_*(\shProj_0(\sh{M}(n)))
    \qquad (n\in\bb{Z}).
\tag{8.14.5.1}
\]

Generalizando a notação de \sref{II.3.3.1}, pomos, para todo \emph{$\sh{S}_X$-módulo graduado $\sh{F}$},
\[
\label{II.8.14.5.2}
    \boldsymbol{\Gamma}_*(\sh{F}) = \bigoplus_{n\in\bb{Z}}q_*(\sh{F}_n).
\tag{8.14.5.2}
\]

Em particular, $\boldsymbol{\Gamma}_*(\sh{S}_X)=\bigoplus_{n\in\bb{Z}}q_*(\sh{O}_X(n))$ é a $\sh{O}_Y$-álgebra graduada denotada por $\boldsymbol{\Gamma}_*(\sh{O}_X)$ em \sref{II.3.3.1.2};
é claro que $\boldsymbol{\Gamma}_*(\sh{F})$ é um \emph{$\boldsymbol{\Gamma}_*(\sh{S}_X)$-módulo graduado} \sref[0]{0.4.2.2}.
Quando
\oldpage[II]{199}
tomamos $\sh{M}=\sh{S}$ nos homomorfismos \sref{II.8.14.5.1}, obtemos o homomorfismo de $\sh{O}_Y$-álgebras graduadas
\[
\label{II.8.14.5.3}
    \alpha:\sh{S}\to\boldsymbol{\Gamma}_*(\sh{S}_X)
\tag{8.14.5.3}
\]
definido anteriormente em \sref{II.3.3.2}, e que faz de $\boldsymbol{\Gamma}_*(\sh{F})$ um \emph{$\sh{S}$-módulo graduado};
os homomorfismos \sref{II.8.14.5.1} definem então um homomorfismo (de grau $0$) de \emph{$\sh{S}$-módulos graduados}
\[
\label{II.8.14.5.4}
    \alpha:\sh{M}\to\boldsymbol{\Gamma}_*(\shProj(\sh{M})).
\tag{8.14.5.4}
\]
\end{env}

\begin{env}[8.14.6]
\label{II.8.14.6}
Em geral, para um $\sh{S}_X$-módulo graduado quase-coerente $\sh{F}$, o $\sh{S}$-módulo graduado $\boldsymbol{\Gamma}_*(\sh{F})$ não é necessariamente quase-coerente.
Consideremos um aberto $X'$ de $X$ tal que a restrição $q':X'\to Y$ de $q$ a $X'$ seja um morfismo \emph{quase-compacto}.
Como $q'$ é, além disso, separado, $q'_*(\sh{F}')$ é então um $\sh{O}_Y$-módulo quase-coerente para todo $\sh{O}_{X'}$-módulo quase-coerente $\sh{F}'$ \sref[I]{I.9.2.2}[b].
Pomos
\[
\label{II.8.14.6.1}
    \sh{S}_{X'} = \sh{S}_X|X' = \bigoplus_{n\in\bb{Z}}\sh{O}_X(n)|X'
\tag{8.14.6.1}
\]
e, para todo $\sh{S}_{X'}$-módulo graduado $\sh{F}'$,
\[
\label{II.8.14.6.2}
    \boldsymbol{\Gamma}'_*(\sh{F}') = \bigoplus_{n\in\bb{Z}}q'_*(\sh{F}'_n).
\tag{8.14.6.2}
\]

A observação anterior mostra então que, se $\sh{F}'$ é um $\sh{S}_{X'}$-módulo quase-coerente, $\boldsymbol{\Gamma}'_*(\sh{F}')$ é um $\sh{S}$-módulo graduado \emph{quase-coerente} \sref[I]{I.9.6.1}.

Observemos também que a injeção canônica $j:X'\to X$ é \emph{quase-compacta}, porque $q'=q\circ j$ é quase-compacto e $q$ é separado \sref[I]{I.6.6.4}[v].
Então $\sh{F}=j_*(\sh{F}')$ é um $\sh{S}_X$-módulo graduado quase-coerente para todo $\sh{S}_{X'}$-módulo graduado quase-coerente $\sh{F}'$, e das definições anteriores decorre que
\[
\label{II.8.14.6.3}
    \boldsymbol{\Gamma}'_*(\sh{F}') = \boldsymbol{\Gamma}_*(\sh{F}).
\tag{8.14.6.3}
\]

Sob as mesmas hipóteses sobre $X'$, para todo $\sh{S}$-módulo graduado quase-coerente $\sh{M}$, pomos
\[
\label{II.8.14.6.4}
    \shProj'(\sh{M}) = \shProj(\sh{M})|X'
\tag{8.14.6.4}
\]
que é um $\sh{S}_{X'}$-módulo graduado quase-coerente.
O homomorfismo canônico
\[
    \shProj(\sh{M}) \to j_*(\shProj'(\sh{M}))
\]
\sref[0]{0.4.4.3} fornece, portanto, um homomorfismo canônico $\boldsymbol{\Gamma}_*(\shProj(\sh{M}))\to\boldsymbol{\Gamma}'_*(\shProj'(\sh{M}))$ de $\sh{S}$-módulos graduados e, por composição com \sref{II.8.14.5.4}, obtemos um homomorfismo canônico funtorial (de grau $0$) de $\sh{S}$-módulos graduados quase-coerentes
\[
\label{II.8.14.6.5}
    \alpha': \sh{M} \to \boldsymbol{\Gamma}'_*(\shProj'(\sh{M})).
\tag{8.14.6.5}
\]
\end{env}

\begin{env}[8.14.7]
\label{II.8.14.7}
Mantendo as hipóteses sobre $X'$ de \sref{II.8.14.6}, seja $\sh{F}'$ um \emph{$\sh{S}_{X'}$-módulo graduado quase-coerente}, de modo que $\shProj'(\boldsymbol{\Gamma}'_*(\sh{F}'))$ também seja um $\sh{S}_{X'}$-módulo graduado \emph{quase-coerente}.
\oldpage[II]{200}
Vamos definir um homomorfismo canônico funtorial (de grau $0$) de $\sh{S}_{X'}$-módulos graduados
\[
\label{II.8.14.7.1}
    \beta': \shProj'(\boldsymbol{\Gamma}'_*(\sh{F}')) \to \sh{F}'.
\tag{8.14.7.1}
\]

Suponhamos, antes de tudo, que $Y=\Spec(A)$ seja afim e que $\sh{S}=\widetilde{S}$, onde $S$ é uma $A$-álgebra graduada em graus positivos;
então $\boldsymbol{\Gamma}'_*(\sh{F}')=\widetilde{M}$, onde $M=\bigoplus_{n\in\bb{Z}}\Gamma(X',\sh{F}'_n)$ é um $S$-módulo graduado.
Seja $f\in S_d$ tal que $D_+(f)\subset X'$;
por definição \sref{II.2.6.2}, $\alpha_d(f)$ restrita a $D_+(f)$ é a seção de $\sh{O}_X(d)$ sobre $D_+(f)$ correspondente ao elemento $f/1$ de $(S(d))_{(f)}$ e, portanto, é invertível;
por conseguinte, também o é $\alpha_{nd}(f^n)$ para todo $n>0$.
Assim, define-se imediatamente um $S_f$-homomorfismo (de grau $0$) de módulos graduados $\beta_f:M_f\to\Gamma(D_+(f),\sh{F}')$ enviando cada elemento $z/f^n\in M_f$ (onde $z\in M$) à seção $(z|D_+(f))(\alpha_{nd}(f^n)|D_+(f))^{-1}$ de $\sh{F}'$ sobre $D_+(f)$.
Além disso, temos um diagrama comutativo correspondente a \sref{II.2.6.4.1}, donde resulta a definição de $\beta'$ nesse caso.
Para passar ao caso geral, é preciso considerar uma $A$-álgebra $A'$, a $A'$-álgebra graduada $S'=S\otimes_A A'$, e utilizar o diagrama comutativo análogo a \sref{II.2.8.13.2};
deixamos os pormenores ao leitor.
\end{env}

\begin{proposition}[8.14.8]
\label{II.8.14.8}
Se $X'$ é um aberto de $X=\Proj(\sh{S})$ tal que $q':X'\to Y$ seja quase-compacto, o homomorfismo $\beta'$ definido em \sref{II.8.14.7} é bijetivo.
\end{proposition}

\begin{proof}
Podemos claramente nos restringir ao caso em que $Y$ seja afim, e tudo se reduz então a provar (com a notação de \sref{II.8.14.7}) que o homomorfismo $\beta_f:M_f\to\Gamma(D_+(f),\sh{F}')$ é um isomorfismo.
Mas substituir $f$ por uma de suas potências não altera nem $D_+(f)$ nem $\beta_f$;
como $X'$ é \emph{quase-compacto} por hipótese, podemos sempre supor, por \sref{II.8.14.4}, que o feixe $\sh{O}_X(d)$ seja \emph{invertível}.
Como $X'$ é um esquema (pois $q'$ é separado), a proposição é então exatamente \sref[I]{I.9.3.1}.
\end{proof}

\begin{corollary}[8.14.9]
\label{II.8.14.9}
Sob as hipóteses de \sref{II.8.14.8}, todo $\sh{S}_{X'}$-módulo graduado quase-coerente é isomorfo a um $\sh{S}_{X'}$-módulo graduado da forma $\shProj'(\sh{M})$, onde $\sh{M}$ é um $\sh{S}$-módulo graduado quase-coerente.
Além disso, se $\sh{F}'$ é de tipo finito e supomos que $Y$ seja um esquema quase-compacto ou um pré-esquema cujo espaço subjacente seja noetheriano, podemos supor que $\sh{M}$ seja de tipo finito.
\end{corollary}

\begin{proof}
A demonstração, a partir de \sref{II.8.14.8}, segue exatamente o mesmo caminho que a demonstração de \sref{II.3.4.5} a partir de \sref{II.3.4.4}, e deixamos os pormenores ao leitor.
\end{proof}

\begin{proposition}[8.14.10]
\label{II.8.14.10}
Sob as hipóteses de \sref{II.8.14.7}, sejam $\sh{M}$ um $\sh{S}$-módulo graduado quase-coerente e $\sh{F}'$ um $\sh{S}_{X'}$-módulo graduado quase-coerente;
os homomorfismos compostos
\[
\label{II.8.14.10.1}
    \shProj'(\sh{M})
    \xrightarrow{\shProj'(\alpha')}
    \shProj'(\boldsymbol{\Gamma}'_*(\shProj'(\sh{M})))
    \xrightarrow{\beta'}
    \shProj'(\sh{M})
\tag{8.14.10.1}
\]
\[
\label{II.8.14.10.2}
    \boldsymbol{\Gamma}'_*(\sh{F}')
    \xrightarrow{\alpha'}
    \boldsymbol{\Gamma}'_*(\shProj'(\boldsymbol{\Gamma}'_*(\sh{F}')))
    \xrightarrow{\boldsymbol{\Gamma}'_*(\beta')}
    \boldsymbol{\Gamma}'_*(\sh{F}')
\tag{8.14.10.2}
\]
são os isomorfismos identidade.
\end{proposition}

\begin{proof}
A questão é local em $Y$, e a demonstração se faz como em \sref{II.2.6.5};
deixamos os pormenores ao leitor.
\end{proof}

\begin{remark}[8.14.11]
\label{II.8.14.11}
No capítulo~III \sref[III]{III.2.3.1}, veremos que, quando $Y$ é \emph{localmente noetheriano} e $\sh{S}$ é uma $\sh{O}_Y$-álgebra graduada quase-coerente \emph{de tipo finito} (caso em que
\oldpage[II]{201}
podemos tomar $X'=X$), então o homomorfismo $\alpha$ \sref{II.8.14.5.4} é \textbf{(TN)}-\emph{bijetivo} para todo $\sh{S}$-módulo graduado quase-coerente $\sh{M}$ que satisfaça a condição~\textbf{(TF)}.
\end{remark}

\begin{remark}[8.14.12]
\label{II.8.14.12}
A situação descrita em \sref{II.8.14.4} é um caso particular da seguinte.
Seja $X$ um espaço anelado e seja $\sh{S}$ uma $\sh{O}_X$-álgebra graduada (em graus positivos e negativos);
suponhamos que exista um inteiro $d>0$ tal que $\sh{S}_d$ e $\sh{S}_{-d}$ sejam \emph{invertíveis}, e que o homomorfismo canônico
\[
\label{II.8.14.12.1}
    \sh{S}_d\otimes_{\sh{O}_X}\sh{S}_{-d} \to \sh{O}_X
\tag{8.14.12.1}
\]
seja um \emph{isomorfismo} (de modo que $\sh{S}_{-d}$ se identifique com $\sh{S}_d^{-1}$).
Dizemos então que a $\sh{O}_X$-álgebra graduada $\sh{S}$ é \emph{periódica}, \emph{de período $d$}.
Essa terminologia provém da seguinte propriedade:
\emph{sob as hipóteses precedentes, para todo $\sh{S}$-módulo graduado $\sh{F}$, o homomorfismo canônico}
\[
\label{II.8.14.12.2}
    \sh{S}_d\otimes_{\sh{O}_X}\sh{F}_n \to \sh{F}_{n+d}
\tag{8.14.12.2}
\]
\emph{é um isomorfismo para todo $n\in\bb{Z}$.}
Com efeito, a questão é local sobre $X$, e podemos supor que $\sh{S}_d$ possua uma seção \emph{invertível} $s$ sobre $X$, cuja inversa $s'$ seja uma seção de $\sh{S}_{-d}$.
O homomorfismo $\sh{F}_{n+d}\to\sh{S}_d\otimes\sh{F}_n$ que envia cada seção $z\in\Gamma(U,\sh{F}_{n+d})$ à seção $(s|U)\otimes(s'|U)z$ de $\sh{S}_d\otimes\sh{F}_n$ sobre $U$ é então o inverso de \sref{II.8.14.12.2}, donde a afirmação.
Isso induz, para todo $k\in\bb{Z}$, um isomorfismo canônico
\[
    (\sh{S}_d)^{\otimes k}\otimes\sh{F}_n \simto \sh{F}_{n+kd}.
\]
Por conseguinte, \emph{dar um $\sh{S}$-módulo graduado $\sh{F}$ equivale a dar os $\sh{S}_0$-módulos $\sh{F}_i$ ($0\leq i\leq d-1$) e os homomorfismos canônicos}
\[
    \sh{S}_i\otimes\sh{F}_j \to \sh{F}_{i+j}
    \qquad
    \mbox{para $0\leq i,j\leq d-1$}
\]
(pondo $\sh{F}_{i+j}=\sh{S}_d\otimes_{\sh{S}_0}\sh{F}_{i+j-d}$ quando $i+j\geq d$).
Naturalmente, para que esses homomorfismos definam uma estrutura bem determinada de $\sh{S}$-módulo sobre a soma direta dos $(\sh{S}_d)^{\otimes k}\otimes\sh{F}_i$ ($k\in\bb{Z}$, $0\leq i\leq d-1$), eles devem satisfazer certas condições de associatividade que não explicitaremos.

No caso $d=1$ (que é o considerado em \sref{subsection:II.3.3}), podemos dizer, portanto, que a categoria dos $\sh{S}$-módulos graduados (resp. dos $\sh{S}$-módulos graduados quase-coerentes, se $X$ é um pré-esquema e $\sh{S}$ é quase-coerente) é \emph{equivalente} à categoria dos $\sh{S}_0$-módulos arbitrários (resp. quase-coerentes);
é nesse sentido que podemos considerar os resultados deste número como generalizações dos do §\hyperref[section:II.3]{3}.
Além disso, vemos que, sob condições de finitude adequadas, os resultados deste número (juntamente com \sref{II.8.14.11}) reduzem, em certa medida, o estudo das álgebras graduadas quase-coerentes sobre um pré-esquema, e dos módulos graduados «módulo \textbf{(TN)}» sobre tais álgebras, ao estudo do caso particular em que as álgebras consideradas são \emph{periódicas} (e em que a condição~\textbf{(TN)} para $\sh{M}$ \sref{II.3.4.2} implica, portanto, que $\sh{M}=0$).
\end{remark}

\begin{remark}[8.14.13]
\label{II.8.14.13}
Sob as hipóteses de \sref{II.8.14.1}, seja $d$ um inteiro $>0$;
definimos um $Y$-isomorfismo canônico $h$ de $X$ sobre $X^{(d)}=\Proj(\sh{S}^{(d)})$ \sref{II.3.1.8}.
Para todo
\oldpage[II]{202}
$\sh{S}$-módulo graduado quase-coerente $\sh{M}$ e todo inteiro $k$ tal que $0\leq k\leq d-1$, temos também (com a notação de \sref{II.3.1.1}) um $h$-isomorfismo canônico
\[
\label{II.8.14.13.1}
    (\shProj(\sh{M}))^{(d,k)} \xleftarrow{\sim} \shProj(\sh{M}^{(d,k)}).
\tag{8.14.13.1}
\]

Suponhamos, inicialmente, que $Y=\Spec(A)$ seja afim, $\sh{S}=\widetilde{S}$ e $\sh{M}=\widetilde{M}$, onde $S$ é uma $A$-álgebra graduada em graus positivos e $M$ um $S$-módulo graduado.
Sabemos que, para todo $f\in S_e$ ($e>0$), $h$ envia $D_+(f)$ sobre $D_+(f^d)$ e corresponde ao isomorfismo canônico $S_{(f^d)}\to S_{(f)}$ \sref{II.2.2.2}.
A restrição de \sref{II.8.14.13.1} a $D_+(f^d)$ corresponde então ao di-isomorfismo canônico $M_{f^d}\to M_f$, restrito aos elementos de $M_{f^d}$ cujo grau é congruente a $k$ módulo $d$.
Deixamos ao leitor a tarefa de verificar que esses isomorfismos são compatíveis com a passagem de $f$ a um de seus múltiplos homogêneos $fg$, e depois que há uma compatibilidade análoga com a passagem de $S$ a uma $A'$-álgebra graduada $S'=S\otimes_A A'$, onde $A'$ é uma $A$-álgebra.
Em particular, obtemos assim um $h$-isomorfismo
\[
\label{II.8.14.13.2}
    (\sh{S}^{(d)})_{X^{(d)}} \simto (\sh{S}_X)^{(d)}
\tag{8.14.13.2}
\]
que respeita as estruturas multiplicativas da fonte e do alvo e que, graças a \sref{II.8.14.13.1}, se torna um $h$-di-isomorfismo de um $(\sh{S}^{(d)})_{X^{(d)}}$-módulo graduado sobre um $(\sh{S}_X)^{(d)}$-módulo graduado.
Analogamente, temos um $h$-isomorfismo
\[
\label{II.8.14.13.3}
    \shProj_0(\sh{S}^{(d,k)}(n)) \simto \sh{O}_X(nd+k),
\tag{8.14.13.3}
\]
que completa o resultado de \sref{II.3.2.9}[ii].

O isomorfismo de \sref{II.8.14.13.1} induz imediatamente um isomorfismo de $\sh{S}^{(d)}$-módulos graduados
\[
\label{II.8.14.13.4}
    \boldsymbol{\Gamma}_*^{(d)}(\shProj(\sh{M}^{(d,k)})) \simto \boldsymbol{\Gamma}_*((\shProj(\sh{M}))^{(d,k)})
\tag{8.14.13.4}
\]
onde $\boldsymbol{\Gamma}_*^{(d)}$ corresponde ao morfismo estrutural $q^{(d)}:X^{(d)}\to Y$;
verifica-se imediatamente que o homomorfismo canônico $\alpha$ \sref{II.8.14.5.4} e o homomorfismo análogo $\alpha^{(d)}$ para $X^{(d)}$ tornam comutativo o diagrama seguinte:
\[
\label{II.8.14.13.5}
    \xymatrix{
        &\sh{M}^{(d,k)} \ar[dl]_{\alpha^{(d)}} \ar[dr]^{\alpha}&
    \\  \boldsymbol{\Gamma}_*^{(d)}(\shProj(\sh{M}^{(d,k)})) \ar[rr]^{\sim}
        && \boldsymbol{\Gamma}_*((\shProj(\sh{M}))^{(d,k)})
    }
\tag{8.14.13.5}
\]
A verificação é feita supondo que $Y$ seja afim e calculando as restrições das imagens por $\alpha^{(d)}$ e $\alpha$ de um mesmo elemento de $M^{(d,k)}$ aos abertos $D_+(f^d)$ e $D_+(f)$ (com a mesma notação acima).
\end{remark}

\begin{proposition}[8.14.14]
\label{II.8.14.14}
Seja $Y$ um pré-esquema quase-compacto, seja $\sh{S}$ uma $\sh{O}_Y$-álgebra graduada quase-coerente de tipo finito e seja $\sh{M}$ um $\sh{S}$-módulo graduado quase-coerente que satisfaça a condição~\textbf{(TF)};
seja $X=\Proj(\sh{S})$.
Então $\sh{S}_X$ é uma $\sh{O}_X$-álgebra graduada periódica \sref{II.8.14.12}, e existe
\oldpage[II]{203}
algum período $d$ de $\sh{S}_X$ tal que os $(\shProj(\sh{M}))^{(d,k)}$ ($0\leq k\leq d-1$) sejam $(\sh{S}_X)^{(d)}$-módulos de tipo finito.
\end{proposition}

\begin{proof}
Com efeito, \sref{II.3.1.10} prova que existe um $d$ tal que $\sh{S}^{(d)}$ seja gerada por $\sh{S}_d=(\sh{S}^{(d)})_1$, sendo esta última um $\sh{S}_0$-módulo de tipo finito.
Para provar a primeira afirmação, podemos, portanto, em virtude de \sref{II.8.14.13.2}, reduzir-nos ao caso $d=1$, e então a proposição decorre de \sref{II.3.2.7}.
Além disso, levando em conta \sref{II.8.14.13.1}, a segunda afirmação é consequência de \sref{II.2.1.6}[iii] e \sref{II.3.4.3}.
\end{proof}
\clearpage
\phantomsection
\addcontentsline{toc}{section}{Bibliografia}
\renewcommand{\refname}{Bibliografia}
\noindent\textit{Nota editorial. Para uso autônomo, reproduzem-se aqui as entradas [1], [9],
[11] e [13], citadas neste volume, a partir da bibliografia da série precedente.
A edição original de EGA~II imprime somente a continuação [23]--[26] na página impressa~205.}
\par\medskip

\begin{thebibliography}{26}

\bibitem[1]{I-1}
H.~Cartan e C.~Chevalley,
\emph{Séminaire de l'École Normale Supérieure, 8th year (1955--56):
Géométrie algébrique}.

\bibitem[9]{I-9}
M.~Nagata,
\emph{A general theory of algebraic geometry over Dedekind domains},
Amer. Math. J. \textbf{78} (1956), 78--116, e \textbf{80} (1958),
382--420.

\bibitem[11]{I-11}
P.~Samuel,
\emph{Commutative algebra} (notas de D.~Herzig), Cornell University,
1953.

\bibitem[13]{I-13}
P.~Samuel e O.~Zariski,
\emph{Commutative algebra}, 2 vols., Van Nostrand, Nova York, 1958--60.

\oldpage[II]{205}
\bibitem[23]{II-23}
C.~Chevalley,
\emph{Introduction to the theory of algebraic functions of one variable}.
Math. Surveys, n.º~6, Amer. Math. Soc., Nova York, 1951.

\bibitem[24]{II-24}
P.~Jaffard,
\emph{Les systèmes d'idéaux}, Paris (Dunod), 1960.

\bibitem[25]{II-25}
M.~Nagata,
\emph{On the derived normal rings of Noetherian integral domains},
Mem. Coll. Sci. Kyoto, sér.~A, t.~XXIX (1955), pp.~293--303.

\bibitem[26]{II-26}
M.~Nagata,
\emph{Existence theorems for non projective complete algebraic varieties},
Ill. J. Math., t.~II (1958), pp.~490--498.

\end{thebibliography}
\clearpage
\oldpage[II]{206}
\clearpage
\oldpage[II]{207}
\phantomsection
\addcontentsline{toc}{section}{Índice de notações}
\section*{Índice de notações}

\begingroup
\small
\setlength{\parindent}{0pt}
\setlength{\parskip}{2pt}

$\sh{A}(X),\sh{A}(\sh{F}),\sh{A}(\sh{B})$
($X$ um $S$-pré-esquema, $\sh{F}$ um $\sh{O}_X$-módulo,
$\sh{B}$ uma $\sh{O}_X$-álgebra): \hyperref[II.1.1.1]{1.1.1}.

$\sh{A}(h)$ ($h$ um $S$-morfismo): \hyperref[II.1.1.2]{1.1.2}.

$\Spec(\sh{B})$ ($\sh{B}$ uma $\sh{O}_S$-álgebra): \hyperref[II.1.3.1]{1.3.1}.

$\widetilde{\sh{M}}$ ($\sh{M}$ um $\sh{B}$-módulo,
$\sh{B}$ uma $\sh{O}_S$-álgebra): \hyperref[II.1.4.3]{1.4.3}.

$\bb{T}(E),\bb{S}(E),\bb{S}_A(E)$ ($E$ um $A$-módulo): \hyperref[II.1.7.1]{1.7.1}.

$\bb{S}(\sh{E}),\bb{S}_{\sh{A}}(\sh{E})$
($\sh{E}$ um $\sh{A}$-módulo): \hyperref[II.1.7.4]{1.7.4}.

$\bb{V}(\sh{E})$ ($\sh{E}$ um $\sh{O}_S$-módulo quase-coerente): \hyperref[II.1.7.8]{1.7.8}.

$S_n,S_+,M_n,S^{(d)},M^{(d,k)},M^{(d)}$
($S$ um anel graduado, $M$ um $S$-módulo graduado): \hyperref[II.2.1.1]{2.1.1}.

$M(n),S(n)$ ($S$ um anel graduado, $M$ um $S$-módulo graduado): \hyperref[II.2.1.1]{2.1.1}.

$\Hom_S(M,N)$ ($S$ um anel graduado, $M,N$ $S$-módulos graduados): \hyperref[II.2.1.2]{2.1.2}.

$\mathfrak N_+,\mathfrak r_+(\mathfrak J)$
($\mathfrak J$ um ideal de $S_+$): \hyperref[II.2.1.10]{2.1.10}.

$S_{(f)},M_{(f)}$ ($f$ um elemento homogêneo de $S_+$): \hyperref[II.2.2.1]{2.2.1}.

$S_{(T)},M_{(T)},S_{(\mathfrak p)},M_{(\mathfrak p)}$
($T$ uma parte multiplicativa de $S_+$ formada por elementos homogêneos,
$\mathfrak p$ um ideal primo graduado de $S_+$): \hyperref[II.2.2.7]{2.2.7}.

$\Proj(S)$ ($S$ um anel graduado): \hyperref[II.2.3.1]{2.3.1}.

$V_+(E)$ ($E$ uma parte de $S_+$): \hyperref[II.2.3.2]{2.3.2}.

$D_+(f)$ ($f$ um elemento de $S$): \hyperref[II.2.3.3]{2.3.3}.

$\mathfrak j_+(Y)$ ($Y$ uma parte de $\Proj(S)$): \hyperref[II.2.3.10]{2.3.10}.

$\widetilde M$ ($M$ um $S$-módulo graduado): \hyperref[II.2.5.2]{2.5.2}.

$\widetilde u$ ($u$ um homomorfismo de grau zero): \hyperref[II.2.5.4]{2.5.4}.

$\sh{O}_X(n),\sh{F}(n)$
($X=\Proj(S)$, $\sh{F}$ um $\sh{O}_X$-módulo): \hyperref[II.2.5.10]{2.5.10}.

$\lambda$ (homomorfismo canônico): \hyperref[II.2.5.11]{2.5.11}.

$\mu$ (homomorfismo canônico): \hyperref[II.2.5.12]{2.5.12}.

$X^{(d)}$ ($X=\Proj(S)$, $d$ um inteiro $>0$): \hyperref[II.2.5.16]{2.5.16}.

$\Gamma_*(\sh{F})$ ($\sh{F}$ um $\sh{O}_X$-módulo,
$X=\Proj(S)$): \hyperref[II.2.6.1]{2.6.1}.

$\alpha_n,\alpha,\alpha_{\sh{M}}$ (homomorfismos canônicos): \hyperref[II.2.6.2]{2.6.2}.

$\beta,\beta_{\sh{F}}$ (homomorfismos canônicos): \hyperref[II.2.6.4]{2.6.4}.

$G(\vphi),{}^a\vphi$ (por abuso de linguagem;
$\vphi$ um homomorfismo de anéis graduados): \hyperref[II.2.8.1]{2.8.1}.

$\Proj(\vphi)$ ($\vphi$ um homomorfismo de anéis graduados): \hyperref[II.2.8.2]{2.8.2}.

$\sh{S}^{(d)},\sh{M}^{(d,k)},\sh{M}^{(d)},\sh{M}(n)$
($\sh{S}$ uma $\sh{O}_Y$-álgebra graduada,
$\sh{M}$ um $\sh{S}$-módulo graduado): \hyperref[II.3.1.1]{3.1.1}.

$\Proj(\sh{S})$ ($\sh{S}$ uma $\sh{O}_Y$-álgebra graduada quase-coerente): \hyperref[II.3.1.3]{3.1.3}.

$X_f$ ($X=\Proj(\sh{S})$, $f$ uma seção de $\sh{S}_d$ sobre $Y$): \hyperref[II.3.1.4]{3.1.4}.

$\widetilde{\sh{M}}$ ($\sh{M}$ um $\sh{S}$-módulo graduado): \hyperref[II.3.2.2]{3.2.2}.

$\sh{O}_X(n),\sh{F}(n)$
($X=\Proj(\sh{S})$, $\sh{F}$ um $\sh{O}_X$-módulo): \hyperref[II.3.2.5]{3.2.5}.

$\lambda,\mu$ (homomorfismos canônicos): \hyperref[II.3.2.6]{3.2.6}.

$\Gamma_*(\sh{F})$ ($\sh{F}$ um $\sh{O}_X$-módulo,
$X=\Proj(\sh{S})$): \hyperref[II.3.3.1]{3.3.1}.

$\alpha_n,\alpha,\alpha_{\sh{M}}$ (homomorfismos canônicos): \hyperref[II.3.3.2]{3.3.2}.
\oldpage[II]{208}
$\beta,\beta_{\sh{F}}$ (homomorfismos canônicos): \hyperref[II.3.3.4]{3.3.4}.

$G(\vphi),\Proj(\vphi)$
($\vphi$ um homomorfismo de $\sh{O}_Y$-álgebras graduadas): \hyperref[II.3.5.1]{3.5.1}.

$\Proj(u)$ ($u$ um $g$-morfismo de uma $\sh{O}_Y$-álgebra graduada
em uma $\sh{O}_{Y'}$-álgebra graduada): \hyperref[II.3.5.6]{3.5.6}.

$r_{\sh{L},\psi}$ ($\sh{L}$ um $\sh{O}_X$-módulo invertível): \hyperref[II.3.7.1]{3.7.1}.

$\bb{P}(\sh{E}),\bb{P}(E),\bb{P}_Y^n,\bb{P}_A^n$: \hyperref[II.4.1.1]{4.1.1}.

$\bb{P}(u)$ ($u$ um homomorfismo sobrejetivo de $\sh{O}_Y$-módulos): \hyperref[II.4.1.2]{4.1.2}.

$\Hyp_Y(X,\sh{E})$ ($\sh{E}$ um $\sh{O}_Y$-módulo quase-coerente): \hyperref[II.4.2.5]{4.2.5}.

$\varsigma$ (morfismo de Segre): \hyperref[II.4.3.1]{4.3.1}.

$\sh{F}(n)$ ($\sh{F}$ um $\sh{O}_X$-módulo, $X$ munido de um
$\sh{O}_X$-módulo invertível): \hyperref[II.4.5.1]{4.5.1}.

$\varepsilon$ (automorfismo identidade): \hyperref[II.4.5.1]{4.5.1}.

$P(u,T),\sigma_i(u)$ ($u$ um endomorfismo de módulo livre de base finita): \hyperref[II.6.4.1]{6.4.1}.

$P(u,T),\sigma_i(u),\det u$
($u$ um endomorfismo de módulo de tipo finito sobre um anel íntegro ou um
anel noetheriano reduzido): \hyperref[II.6.4.2]{6.4.2} e \hyperref[II.6.4.7]{6.4.7}.

$\sigma_i(u),\det u$ ($u$ um endomorfismo de um $\sh{A}$-módulo localmente livre): \hyperref[II.6.4.8]{6.4.8}.

$\sigma_i,\det$ (homomorfismos de feixes): \hyperref[II.6.4.8]{6.4.8}, \hyperref[II.6.4.9]{6.4.9} e \hyperref[II.6.4.10]{6.4.10}.

$\sigma_i(u),\det u$
($u$ um endomorfismo de um $\sh{O}_X$-módulo de tipo finito sobre um pré-esquema
localmente íntegro ou um pré-esquema localmente noetheriano reduzido):
\hyperref[II.6.4.9]{6.4.9} e \hyperref[II.6.4.10]{6.4.10}.

$N_{\sh{B}/\sh{A}}(f)$
($f$ uma seção da $\sh{A}$-álgebra $\sh{B}$): \hyperref[II.6.5.1]{6.5.1}.

$N_{\sh{B}/\sh{A}}(\sh{L}')$
($\sh{L}'$ um $\sh{B}$-módulo invertível): \hyperref[II.6.5.2]{6.5.2}.

$N_{\sh{B}/\sh{A}}(h')$
($h'$ um homomorfismo de $\sh{B}$-módulos invertíveis): \hyperref[II.6.5.3]{6.5.3}.

$N_{X'/X}(\sh{L}')$
($X'$ finito sobre $X$, $\sh{L}'$ um $\sh{O}_{X'}$-módulo invertível): \hyperref[II.6.5.5]{6.5.5}.

$N_{X'/X}(h')$
($h'$ um homomorfismo de $\sh{O}_{X'}$-módulos invertíveis): \hyperref[II.6.5.5]{6.5.5}.

$\sh{S}_X$ ($X=\Proj(\sh{S})$, $\sh{S}$ uma soma de ideais de
$\sh{O}_Y$): \hyperref[II.8.1.5]{8.1.5}.

$S^\geq,S^\leq,S_f^\geq,S_f^\leq,
M^\geq,M^\leq,M_f^\geq,M_f^\leq$
($S$ um anel graduado, $M$ um $S$-módulo graduado,
$f$ um elemento homogêneo de $S_+$): \hyperref[II.8.2.1]{8.2.1}.

$\widehat S$ ($S$ um anel graduado): \hyperref[II.8.2.2]{8.2.2}.

$\widehat M$ ($M$ um $S$-módulo graduado): \hyperref[II.8.2.4]{8.2.4}.

$S_{[n]},S^\natural,f^\natural$
($S$ um anel graduado, $f$ homogêneo em $S$): \hyperref[II.8.2.6]{8.2.6}.

$M_{[n]},M^\natural$ ($M$ um $S$-módulo graduado): \hyperref[II.8.2.8]{8.2.8}.

$\widehat{\sh{S}}$ ($\sh{S}$ uma $\sh{O}_Y$-álgebra graduada): \hyperref[II.8.3.1]{8.3.1}.

$\bb{G}_m$ (esquema de grupos): \hyperref[II.8.3.9]{8.3.9}.

$\sh{S}_X$ ($X=\Proj(\sh{S})$, $\sh{S}$ uma $\sh{O}_Y$-álgebra graduada): \hyperref[II.8.6.1]{8.6.1}.

$\sh{S}_X,C_X,\widehat C_X,E_X,\widehat E_X$
($X=\Proj(\sh{S})$, $\sh{S}$ uma $\sh{O}_Y$-álgebra graduada
com $\sh{S}_0=\sh{O}_Y$): \hyperref[II.8.6.1]{8.6.1}.

$\sh{S}_{[n]},\sh{S}^{\natural}$
($\sh{S}$ uma $\sh{O}_Y$-álgebra graduada): \hyperref[II.8.7.2]{8.7.2}.

$\shProj_0(\sh{M}),\shProj(\sh{M}),\sh{M}_X,
\widehat{\sh{M}},\sh{M}^{\square}$
($\sh{M}$ um $\sh{S}$-módulo graduado): \hyperref[II.8.12.1]{8.12.1}.

$\sh{M}_X^\geq$
($X=\Proj(\sh{S})$, $\sh{M}$ um $\sh{S}$-módulo graduado): \hyperref[II.8.12.5]{8.12.5}.

$\sh{M}_{[n]},\sh{M}^{\natural}$
($\sh{M}$ um $\sh{S}$-módulo graduado): \hyperref[II.8.12.9]{8.12.9}.

$(\widetilde{\sh{N}})^-$
($\sh{N}$ um sub-$\sh{S}$-módulo de um $\sh{S}$-módulo graduado): \hyperref[II.8.13.1]{8.13.1}.

$\overline{\sh{N}}$
($\sh{N}$ um sub-$\sh{S}$-módulo de um $\sh{S}$-módulo graduado): \hyperref[II.8.13.2]{8.13.2}.

$\Gamma_*(\sh{F})$
($X=\Proj(\sh{S})$, $\sh{F}$ um $\sh{S}_X$-módulo graduado): \hyperref[II.8.14.5]{8.14.5}.

$\sh{S}_{X'},\Gamma'_*(\sh{F}'),\shProj'(\sh{M}),
\alpha',\beta'$
($X'$ aberto em $X=\Proj(\sh{S})$, $\sh{F}'$ um $\sh{S}_{X'}$-módulo graduado,
$\sh{M}$ um $\sh{S}$-módulo graduado): \hyperref[II.8.14.6]{8.14.6}.

\endgroup

\clearpage
\oldpage[II]{209}
\phantomsection
\addcontentsline{toc}{section}{Índice terminológico}
\section*{Índice terminológico}

\begingroup
\small
\setlength{\parindent}{0pt}
\setlength{\parskip}{2pt}
\newcommand{\egaTerm}[2]{\textbf{#1}: #2.\par}

\egaTerm{Afim (morfismo)}{\hyperref[II.1.6.1]{1.6.1}}
\egaTerm{Afim (pré-esquema) sobre um pré-esquema}{\hyperref[II.1.2.1]{1.2.1}}
\egaTerm{Álgebra graduada sobre um anel graduado}{\hyperref[II.2.1.2]{2.1.2}}
\egaTerm{Álgebra simétrica de um $A$-módulo}{\hyperref[II.1.7.1]{1.7.1}}
\egaTerm{$\sh{A}$-álgebra simétrica de um $\sh{A}$-módulo}{\hyperref[II.1.7.4]{1.7.4}}
\egaTerm{$\sh{O}_Y$-álgebra graduada essencialmente reduzida}{\hyperref[II.3.1.12]{3.1.12}}
\egaTerm{$\sh{O}_Y$-álgebra graduada íntegra}{\hyperref[II.3.1.12]{3.1.12}}
\egaTerm{Amplo ($\sh{O}_X$-módulo invertível)}{\hyperref[II.4.5.3]{4.5.3}}
\egaTerm{Amplo relativamente a $f$; $f$-amplo; $Y$-amplo ($\sh{O}_X$-módulo invertível)}{\hyperref[II.4.6.1]{4.6.1}}
\egaTerm{Anel graduado essencialmente íntegro}{\hyperref[II.2.1.11]{2.1.11}}
\egaTerm{Anel graduado essencialmente reduzido}{\hyperref[II.2.1.10]{2.1.10}}
\egaTerm{Associado a um $\sh{B}$-módulo (feixe)}{\hyperref[II.1.4.3]{1.4.3}}
\egaTerm{Associado a um $S$-módulo graduado (feixe)}{\hyperref[II.2.5.3]{2.5.3}}
\egaTerm{Associado a um $\sh{S}$-módulo graduado (feixe)}{\hyperref[II.3.2.2]{3.2.2}}
\egaTerm{Associado a uma $\sh{O}_S$-álgebra ($S$-esquema)}{\hyperref[II.1.3.1]{1.3.1}}

\egaTerm{Condição (TF); condição (TN)}{\hyperref[II.2.7.2]{2.7.2}}
\egaTerm{Condição (TF); condição (TN)}{\hyperref[II.3.4.2]{3.4.2}}
\egaTerm{Cone afim, cone projetivo definido por uma $\sh{O}_Y$-álgebra graduada}{\hyperref[II.8.3.1]{8.3.1}}
\egaTerm{Cone afim sem vértice, cone projetivo sem vértice definido por uma $\sh{O}_Y$-álgebra graduada}{\hyperref[II.8.3.4]{8.3.4}}
\egaTerm{Cone projetante afim, cone projetante projetivo de um pré-esquema $\Proj(\sh{S})$}{\hyperref[II.8.3.1]{8.3.1}}
\egaTerm{Curva algébrica sobre um corpo}{\hyperref[II.7.4.2]{7.4.2}}
\egaTerm{Curva algébrica completa}{\hyperref[II.7.4.7]{7.4.7}}
\egaTerm{Critério de Serre}{\hyperref[II.5.2.1]{5.2.1}}

\egaTerm{Determinante de um endomorfismo de um módulo de tipo finito sobre um anel íntegro}{\hyperref[II.6.4.2]{6.4.2}}
\egaTerm{Determinante de um endomorfismo de um módulo de tipo finito sobre um anel noetheriano reduzido}{\hyperref[II.6.4.7]{6.4.7}}
\egaTerm{Determinante de um endomorfismo de um $\sh{A}$-módulo localmente livre}{\hyperref[II.6.4.8]{6.4.8}}
\egaTerm{Determinante de um endomorfismo de um $\sh{O}_X$-módulo de tipo finito sobre um pré-esquema $X$ localmente íntegro}{\hyperref[II.6.4.9]{6.4.9}}
\egaTerm{Determinante de um endomorfismo de um $\sh{O}_X$-módulo de tipo finito sobre um pré-esquema $X$ localmente noetheriano e reduzido}{\hyperref[II.6.4.10]{6.4.10}}

\egaTerm{Explosão (pré-esquema)}{\hyperref[II.8.1.3]{8.1.3}}
\egaTerm{Inteiro (morfismo)}{\hyperref[II.6.1.1]{6.1.1}}
\egaTerm{Inteiro (pré-esquema) sobre um pré-esquema}{\hyperref[II.6.1.1]{6.1.1}}
\egaTerm{Inteira ($\sh{O}_S$-álgebra quase-coerente)}{\hyperref[II.6.1.2]{6.1.2}}
\egaTerm{Inteira (seção de uma $\sh{A}$-álgebra) sobre $\sh{A}$}{\hyperref[II.6.3.1]{6.3.1}}

\egaTerm{Feixe fundamental de um fibrado projetivo}{\hyperref[II.4.1.1]{4.1.1}}
\egaTerm{Fecho integral de um feixe de anéis $\sh{A}$ em uma $\sh{A}$-álgebra}{\hyperref[II.6.3.2]{6.3.2}}
\egaTerm{Fecho integral de um pré-esquema $X$ relativamente a uma $\sh{O}_X$-álgebra}{\hyperref[II.6.3.4]{6.3.4}}
\egaTerm{Fecho projetivo de um cone afim}{\hyperref[II.8.3.1]{8.3.1}}
\oldpage[II]{210}
\egaTerm{Fibra geométrica racional de um fibrado projetivo sobre uma extensão $K$ de $\kres(y)$}{\hyperref[II.4.2.6]{4.2.6}}
\egaTerm{Fibrado projetivo definido por um $\sh{O}_Y$-módulo}{\hyperref[II.4.1.1]{4.1.1}}
\egaTerm{Fibrado vetorial definido por um $\sh{O}_S$-módulo}{\hyperref[II.1.7.8]{1.7.8}}
\egaTerm{Finito (morfismo)}{\hyperref[II.6.1.1]{6.1.1}}
\egaTerm{Finito (pré-esquema) sobre um pré-esquema}{\hyperref[II.6.1.1]{6.1.1}}
\egaTerm{Finita ($\sh{O}_S$-álgebra quase-coerente)}{\hyperref[II.6.1.2]{6.1.2}}
\egaTerm{Funções simétricas elementares de um endomorfismo de um módulo localmente livre}{\hyperref[II.6.4.1]{6.4.1}}
\egaTerm{Funções simétricas elementares de um endomorfismo de um módulo de tipo finito sobre um anel íntegro}{\hyperref[II.6.4.2]{6.4.2}}
\egaTerm{Funções simétricas elementares de um endomorfismo de um módulo de tipo finito sobre um anel noetheriano reduzido}{\hyperref[II.6.4.7]{6.4.7}}
\egaTerm{Funções simétricas elementares de um endomorfismo de um $\sh{A}$-módulo localmente livre de posto constante}{\hyperref[II.6.4.8]{6.4.8}}
\egaTerm{Funções simétricas elementares de um endomorfismo de um $\sh{O}_X$-módulo de tipo finito sobre um pré-esquema $X$ localmente íntegro}{\hyperref[II.6.4.9]{6.4.9}}
\egaTerm{Funções simétricas elementares de um endomorfismo de um $\sh{O}_X$-módulo de tipo finito sobre um pré-esquema $X$ localmente noetheriano reduzido}{\hyperref[II.6.4.10]{6.4.10}}
\egaTerm{Funções simétricas elementares de uma seção de uma $\sh{A}$-álgebra}{\hyperref[II.6.5.1]{6.5.1}}

\egaTerm{Homogeneização de um sub-$\sh{S}$-módulo de um $\sh{S}$-módulo graduado}{\hyperref[II.8.13.2]{8.13.2}}
\egaTerm{Homomorfismo de anéis graduados}{\hyperref[II.2.1.2]{2.1.2}}
\egaTerm{Homomorfismo de grau $k$ de módulos graduados}{\hyperref[II.2.1.2]{2.1.2}}

\egaTerm{Ideal de $S_+$; ideal primo graduado de $S_+$ ($S$ um anel graduado)}{\hyperref[II.2.1.10]{2.1.10}}

\egaTerm{Lema de Chow}{\hyperref[II.5.6.1]{5.6.1}}
\egaTerm{Lugar no infinito de um cone projetivo}{\hyperref[II.8.3.3]{8.3.3}}

\egaTerm{Morfismo canônico $G(\varepsilon)\to\Proj(S)$, $S=\bigoplus_{n\geq0}\Gamma(X,\sh{L}^{\otimes n})$}{\hyperref[II.4.5.1]{4.5.1}}
\egaTerm{Morfismo canônico $X\to\Spec(\Gamma(X,\sh{O}_X))$}{\hyperref[II.5.1.1]{5.1.1}}
\egaTerm{Morfismo de Segre}{\hyperref[II.4.3.1]{4.3.1}}

\egaTerm{Nilradical de $S_+$ ($S$ um anel graduado)}{\hyperref[II.2.1.10]{2.1.10}}
\egaTerm{Nilradical de $\sh{S}_+$ ($\sh{S}$ uma $\sh{O}_Y$-álgebra graduada quase-coerente)}{\hyperref[II.3.1.12]{3.1.12}}
\egaTerm{Norma de uma seção de uma $\sh{A}$-álgebra}{\hyperref[II.6.5.1]{6.5.1}}
\egaTerm{Norma de um $\sh{B}$-módulo invertível}{\hyperref[II.6.5.2]{6.5.2}}
\egaTerm{Norma de uma seção de um $\sh{B}$-módulo invertível}{\hyperref[II.6.5.3]{6.5.3}}
\egaTerm{Norma de um $\sh{O}_{X'}$-módulo invertível}{\hyperref[II.6.5.5]{6.5.5}}

\egaTerm{Periódica (álgebra graduada)}{\hyperref[II.8.14.12]{8.14.12}}
\egaTerm{Polinômio característico de um endomorfismo de um módulo de tipo finito sobre um anel íntegro}{\hyperref[II.6.4.2]{6.4.2}}
\egaTerm{Polinômio característico de um endomorfismo de um módulo de tipo finito sobre um anel noetheriano reduzido}{\hyperref[II.6.4.7]{6.4.7}}
\egaTerm{Pré-esquema dos vértices de um cone afim (projetivo)}{\hyperref[II.8.3.3]{8.3.3}}
\egaTerm{Apresentação finita (módulo graduado que possui uma)}{\hyperref[II.2.1.1]{2.1.1}}
\egaTerm{Apresentação finita ($\sh{S}$-módulo graduado que possui uma)}{\hyperref[II.3.1.1]{3.1.1}}
\egaTerm{Produto tensorial graduado de dois módulos graduados}{\hyperref[II.2.1.2]{2.1.2}}
\egaTerm{Projetivo (morfismo)}{\hyperref[II.5.5.2]{5.5.2}}
\egaTerm{Projetivo (pré-esquema) sobre um pré-esquema}{\hyperref[II.5.5.2]{5.5.2}}
\egaTerm{Próprio (morfismo)}{\hyperref[II.5.4.1]{5.4.1}}
\egaTerm{Própria (parte)}{\hyperref[II.5.4.10]{5.4.10}}
\egaTerm{Próprio (pré-esquema) sobre um pré-esquema}{\hyperref[II.5.4.1]{5.4.1}}

\egaTerm{Quase-afim (morfismo)}{\hyperref[II.5.1.1]{5.1.1}}
\egaTerm{Quase-afim (pré-esquema) sobre um pré-esquema}{\hyperref[II.5.1.1]{5.1.1}}
\egaTerm{Quase-afim (esquema)}{\hyperref[II.5.1.1]{5.1.1}}
\oldpage[II]{211}
\markboth{CAPÍTULO II}{CAPÍTULO II}
\egaTerm{Quase-finito (morfismo)}{\hyperref[II.6.2.3]{6.2.3}}
\egaTerm{Quase-finito (pré-esquema) sobre um pré-esquema}{\hyperref[II.6.2.3]{6.2.3}}
\egaTerm{Quase-projetivo (morfismo)}{\hyperref[II.5.3.1]{5.3.1}}
\egaTerm{Quase-projetivo (pré-esquema) sobre um pré-esquema}{\hyperref[II.5.3.1]{5.3.1}}

\egaTerm{Radical de um ideal de $S_+$ ($S$ anel graduado)}{\hyperref[II.2.1.10]{2.1.10}}
\egaTerm{Retração canônica de um cone projetivo sem vértice sobre seu lugar no infinito}{\hyperref[II.8.3.5]{8.3.5}}

\egaTerm{Seção canônica de $\sh{O}_X(-1)$ ($X$ pré-esquema obtido por explosão)}{\hyperref[II.8.1.9]{8.1.9}}
\egaTerm{Seção nula de um cone afim}{\hyperref[II.8.3.3]{8.3.3}}
\egaTerm{Seção nula de um fibrado vetorial}{\hyperref[II.1.7.9]{1.7.9}}
\egaTerm{Seção vértice de um cone afim (projetivo)}{\hyperref[II.8.3.3]{8.3.3}}
\egaTerm{Espectro de uma $\sh{O}_S$-álgebra}{\hyperref[II.1.3.1]{1.3.1}}
\egaTerm{Espectro primo homogêneo de um anel graduado}{\hyperref[II.2.3.1]{2.3.1}}
\egaTerm{Espectro homogêneo de uma $\sh{O}_S$-álgebra graduada quase-coerente}{\hyperref[II.3.1.3]{3.1.3}}

\egaTerm{(TN)-injetivo, (TN)-sobrejetivo, (TN)-bijetivo (homomorfismo de módulos graduados)}{\hyperref[II.2.7.2]{2.7.2}}
\egaTerm{(TN)-isomorfismo de módulos graduados}{\hyperref[II.2.7.2]{2.7.2}}
\egaTerm{(TN)-injetivo, (TN)-sobrejetivo, (TN)-bijetivo (homomorfismo de anéis graduados)}{\hyperref[II.2.9.1]{2.9.1}}
\egaTerm{(TN)-isomorfismo de anéis graduados}{\hyperref[II.2.9.1]{2.9.1}}
\egaTerm{(TN)-injetivo, (TN)-sobrejetivo, (TN)-bijetivo (homomorfismo de $\sh{S}$-módulos graduados)}{\hyperref[II.3.4.2]{3.4.2}}
\egaTerm{(TN)-isomorfismo de $\sh{S}$-módulos graduados}{\hyperref[II.3.4.2]{3.4.2}}
\egaTerm{(TN)-injetivo, (TN)-sobrejetivo, (TN)-bijetivo (homomorfismo de $\sh{O}_Y$-álgebras graduadas)}{\hyperref[II.3.6.1]{3.6.1}}
\egaTerm{(TN)-isomorfismo de $\sh{O}_Y$-álgebras graduadas}{\hyperref[II.3.6.1]{3.6.1}}
\egaTerm{Topologia espectral sobre $\Proj(S)$}{\hyperref[II.2.3.3]{2.3.3}}
\egaTerm{Muito amplo para $q$, muito amplo para $Y$ ($\sh{O}_X$-módulo invertível)}{\hyperref[II.4.4.2]{4.4.2}}
\egaTerm{De tipo finito ($\sh{S}$-módulo graduado)}{\hyperref[II.3.1.1]{3.1.1}}

\egaTerm{Universalmente fechado (morfismo)}{\hyperref[II.5.4.9]{5.4.9}}

\endgroup

\clearpage
\oldpage[II]{212}
\thispagestyle{empty}
\mbox{}

\clearpage
\oldpage[II]{213}
\markboth{SUMÁRIO}{SUMÁRIO}
\phantomsection
\addcontentsline{toc}{section}{Sumário original}
\section*{Sumário}
\thispagestyle{plain}

\begingroup
\small
\renewcommand{\arraystretch}{1.04}
\begin{longtable}{@{}p{0.09\textwidth}p{0.69\textwidth}r@{}}
\textbf{Unidade} & \textbf{Título} & \textbf{Página original}\\
\hline
\textbf{II} & \textbf{Estudo global elementar de algumas classes de morfismos} & \textbf{5}\\
\hyperref[section:II.1]{\textsection1} & Morfismos afins & 5\\
\hyperref[subsection:II.1.1]{1.1} & $S$-pré-esquemas e $\sh{O}_S$-álgebras & 5\\
\hyperref[subsection:II.1.2]{1.2} & Pré-esquemas afins sobre um pré-esquema & 6\\
\hyperref[subsection:II.1.3]{1.3} & Pré-esquema afim sobre $S$ associado a uma $\sh{O}_S$-álgebra & 8\\
\hyperref[subsection:II.1.4]{1.4} & Feixes quase-coerentes sobre um pré-esquema afim sobre $S$ & 9\\
\hyperref[subsection:II.1.5]{1.5} & Mudança do pré-esquema de base & 12\\
\hyperref[subsection:II.1.6]{1.6} & Morfismos afins & 14\\
\hyperref[subsection:II.1.7]{1.7} & Fibrado vetorial associado a um feixe de módulos & 14\\
\hyperref[section:II.2]{\textsection2} & Espectros primos homogêneos & 19\\
\hyperref[subsection:II.2.1]{2.1} & Generalidades sobre anéis e módulos graduados & 19\\
\hyperref[subsection:II.2.2]{2.2} & Anéis de frações de um anel graduado & 23\\
\hyperref[subsection:II.2.3]{2.3} & Espectro primo homogêneo de um anel graduado & 25\\
\hyperref[subsection:II.2.4]{2.4} & A estrutura de esquema sobre $\Proj(S)$ & 28\\
\hyperref[subsection:II.2.5]{2.5} & Feixe associado a um módulo graduado & 30\\
\hyperref[subsection:II.2.6]{2.6} & $S$-módulo graduado associado a um feixe sobre $\Proj(S)$ & 36\\
\hyperref[subsection:II.2.7]{2.7} & Condições de finitude & 38\\
\hyperref[subsection:II.2.8]{2.8} & Comportamentos funtoriais & 41\\
\hyperref[subsection:II.2.9]{2.9} & Subpré-esquemas fechados de um esquema $\Proj(S)$ & 48\\
\hyperref[section:II.3]{\textsection3} & Espectro homogêneo de um feixe de álgebras graduadas & 49\\
\hyperref[subsection:II.3.1]{3.1} & Espectro homogêneo de uma $\sh{O}_Y$-álgebra graduada quase-coerente & 49\\
\hyperref[subsection:II.3.2]{3.2} & Feixe sobre $\Proj(\sh{S})$ associado a um $\sh{S}$-módulo graduado & 54\\
\hyperref[subsection:II.3.3]{3.3} & $\sh{S}$-módulo graduado associado a um feixe sobre $\Proj(\sh{S})$ & 56\\
\hyperref[subsection:II.3.4]{3.4} & Condições de finitude & 59\\
\hyperref[subsection:II.3.5]{3.5} & Comportamentos funtoriais & 61\\
\hyperref[subsection:II.3.6]{3.6} & Subpré-esquemas fechados de um pré-esquema $\Proj(\sh{S})$ & 64\\
\hyperref[subsection:II.3.7]{3.7} & Morfismos de um pré-esquema em um espectro homogêneo & 65\\
\hyperref[subsection:II.3.8]{3.8} & Critérios de imersão em um espectro homogêneo & 69\\
\end{longtable}
\oldpage[II]{214}
\markboth{A. GROTHENDIECK --- Cap. II}{A. GROTHENDIECK --- Cap. II}
\begin{longtable}{@{}p{0.09\textwidth}p{0.69\textwidth}r@{}}
\hyperref[section:II.4]{\textsection4} & Fibrados projetivos; feixes amplos & 71\\
\hyperref[subsection:II.4.1]{4.1} & Definição dos fibrados projetivos & 71\\
\hyperref[subsection:II.4.2]{4.2} & Morfismos de um pré-esquema em um fibrado projetivo & 72\\
\hyperref[subsection:II.4.3]{4.3} & O morfismo de Segre & 76\\
\hyperref[subsection:II.4.4]{4.4} & Imersões em fibrados projetivos; feixes muito amplos & 78\\
\hyperref[subsection:II.4.5]{4.5} & Feixes amplos & 83\\
\hyperref[subsection:II.4.6]{4.6} & Feixes relativamente amplos & 89\\
\hyperref[section:II.5]{\textsection5} & Morfismos quase-afins; morfismos quase-projetivos; morfismos próprios; morfismos projetivos & 94\\
\hyperref[subsection:II.5.1]{5.1} & Morfismos quase-afins & 94\\
\hyperref[subsection:II.5.2]{5.2} & O critério de Serre & 97\\
\hyperref[subsection:II.5.3]{5.3} & Morfismos quase-projetivos & 99\\
\hyperref[subsection:II.5.4]{5.4} & Morfismos próprios e morfismos universalmente fechados & 100\\
\hyperref[subsection:II.5.5]{5.5} & Morfismos projetivos & 103\\
\hyperref[subsection:II.5.6]{5.6} & O lema de Chow & 106\\
\hyperref[section:II.6]{\textsection6} & Morfismos inteiros e morfismos finitos & 110\\
\hyperref[subsection:II.6.1]{6.1} & Pré-esquemas inteiros sobre outro & 110\\
\hyperref[subsection:II.6.2]{6.2} & Morfismos quase-finitos & 114\\
\hyperref[subsection:II.6.3]{6.3} & Fecho integral de um pré-esquema & 116\\
\hyperref[subsection:II.6.4]{6.4} & Determinante de um endomorfismo de um $\sh{O}_X$-módulo & 120\\
\hyperref[subsection:II.6.5]{6.5} & Norma de um feixe invertível & 125\\
\hyperref[subsection:II.6.6]{6.6} & Aplicação: critérios de amplitude & 130\\
\hyperref[subsection:II.6.7]{6.7} & O teorema de Chevalley & 135\\
\hyperref[section:II.7]{\textsection7} & Critérios valorativos & 138\\
\hyperref[subsection:II.7.1]{7.1} & Recordação sobre os anéis de valoração & 138\\
\hyperref[subsection:II.7.2]{7.2} & Critério valorativo de separação & 141\\
\hyperref[subsection:II.7.3]{7.3} & Critério valorativo da propriedade de ser próprio & 143\\
\hyperref[subsection:II.7.4]{7.4} & Curvas algébricas e corpos de funções de dimensão um & 148\\
\hyperref[section:II.8]{\textsection8} & Esquemas de explosão; cones projetantes; fecho projetivo & 152\\
\hyperref[subsection:II.8.1]{8.1} & Pré-esquemas de explosão & 152\\
\hyperref[subsection:II.8.2]{8.2} & Resultados preliminares sobre a localização de anéis graduados & 157\\
\hyperref[subsection:II.8.3]{8.3} & Cones projetantes & 162\\
\hyperref[subsection:II.8.4]{8.4} & Fecho projetivo de um fibrado vetorial & 168\\
\hyperref[subsection:II.8.5]{8.5} & Comportamentos funtoriais & 169\\
\hyperref[subsection:II.8.6]{8.6} & Um isomorfismo canônico para os cones sem vértice & 171\\
\end{longtable}
\oldpage[II]{215}
\markboth{SUMÁRIO}{SUMÁRIO}
\begin{longtable}{@{}p{0.09\textwidth}p{0.69\textwidth}r@{}}
\hyperref[subsection:II.8.7]{8.7} & Explosão de cones projetantes & 173\\
\hyperref[subsection:II.8.8]{8.8} & Feixes amplos e contrações & 177\\
\hyperref[subsection:II.8.9]{8.9} & O critério de amplitude de Grauert: enunciado & 182\\
\hyperref[subsection:II.8.10]{8.10} & O critério de amplitude de Grauert: demonstração & 184\\
\hyperref[subsection:II.8.11]{8.11} & Unicidade das contrações & 189\\
\hyperref[subsection:II.8.12]{8.12} & Feixes quase-coerentes sobre os cones projetantes & 191\\
\hyperref[subsection:II.8.13]{8.13} & Fecho projetivo de subfeixes e subesquemas fechados & 195\\
\hyperref[subsection:II.8.14]{8.14} & Complementos sobre os feixes associados a $\sh{S}$-módulos graduados & 197\\
\multicolumn{2}{@{}l}{Bibliografia (continuação)} & 205\\
\multicolumn{2}{@{}l}{Índice de notações} & 207\\
\multicolumn{2}{@{}l}{Índice terminológico} & 209\\
\multicolumn{2}{@{}l}{Errata e adendos (Lista 1)} & 217\\
\end{longtable}

\begin{flushright}
\emph{Recebido em 22 de julho de 1960.}
\end{flushright}
\endgroup
\oldpage[II]{216}
\oldpage[II]{217}
\phantomsection
\addcontentsline{toc}{section}{Errata e adendos (Lista 1)}
\section*{Errata e adendos}
\begin{center}
\large (Lista 1)
\end{center}

\subsection*{Capítulo 0}

\paragraph{(0, 2.1.3).}
Na 15ª linha a contar de baixo da p.~21, substitua «inter-seção» por
«interseção».

\paragraph{(0, 2.1.6).}
Na linha 13 da p.~22, substitua
\[
  Z_i\cap C\left(\bigcup_{j\ne i}Z_j\right)
  \quad\hbox{por}\quad
  \bigcup_{j\ne i}Z_j.
\]

\paragraph{(0, 3.1.6).}
Na 17ª linha a contar de baixo da p.~25, substitua
$\Gamma(U,\sh{F})$ por $\Gamma(V,\sh{F})$.

\paragraph{(0, 4.1.1).}
Na linha 12 da p.~36, substitua $\sh{G}\to\sh{F}$ por
$\sh{B}\to\sh{A}$.

\paragraph{(0, 4.2.1).}
Na 11ª linha a contar de baixo da p.~39, substitua «feixes» por
«feixe».

\paragraph{(0, 5.1.2).}
Após a linha 14 da p.~45, acrescente:

Se $\sh{G}$ é um $\sh{O}_Y$-módulo gerado por suas seções sobre $Y$,
então $f^*(\sh{G})$ é gerado por suas seções sobre $X$, pois $f^*$ é
um funtor exato à direita.

\paragraph{(0, 5.2.4).}
Na 18ª linha a contar de baixo da p.~46, substitua $f_U$ por $f_U^*$.

\paragraph{(0, 5.3.4).}
Antes da 16ª linha a contar de baixo da p.~47, acrescente:

Se
\[
  \sh{A}\to\sh{B}\to\sh{C}\to\sh{D}\to\sh{E}
\]
é uma sequência exata de $\sh{O}_X$-módulos, e se
$\sh{A},\sh{B},\sh{D},\sh{E}$ são coerentes, então $\sh{C}$ é
coerente.

\paragraph{(0, 5.4.1).}
Após a linha 2 da p.~49, acrescente:

Suponha que $\sh{O}_X$ seja coerente, e seja $\sh{F}$ um
$\sh{O}_X$-módulo coerente. Se, em um ponto $x\in X$, $\sh{F}_x$ é um
$\sh{O}_x$-módulo livre de posto $n$, existe uma vizinhança $U$ de $x$
tal que $\sh{F}|U$ é localmente livre de posto $n$; com efeito,
$\sh{F}_x$ é então isomorfo a $\sh{O}_x^n$, e a afirmação decorre de
(0, 5.2.7).

\paragraph{(0, 6.6.2).}
Substitua as linhas 10 e 11 da p.~59 por:

Com efeito, isto decorre de (0, 6.4.1, d).

\paragraph{(0, 6.7.1).}
Antes da 10ª linha a contar de baixo da p.~59, acrescente:

Se $f$ é um morfismo plano, também dizemos que $X$ é
\emph{plano sobre $Y$}, ou \emph{$Y$-plano}.

\paragraph{(0, 7.2.9).}
Na 11ª linha a contar de baixo da p.~65, substitua $M_n$ por
$M_{n-1}$.

\subsection*{Capítulo I}

\paragraph{(I, 1.2.1).}
Na linha 18 da p.~83, substitua $K$ por $k$ (duas vezes).

\paragraph{(I, 1.2.7).}
Na 16ª linha a contar de baixo da p.~84, substitua $A$ por $A'$ (duas
vezes) e $\rad$ por $\rad'$; na 17ª linha a contar de baixo, substitua
$X$ por $X'$.

\paragraph{(I, 1.7.3) e (I, 2.2.4).}
Os enunciados dados sob esses números podem ser generalizados, de
acordo com uma observação de J.~Tate, como segue.

\subsection*{1.8. Morfismos de espaços localmente anelados para esquemas afins}
\label{ega2:addendum:subsection:I.1.8}

\begin{proposition}[1.8.1]
\label{ega2:addendum:I.1.8.1}
Sejam $(S,\sh{O}_S)$ um esquema afim e $(X,\sh{O}_X)$ um espaço
localmente anelado. Há uma bijeção canônica do conjunto dos
homomorfismos de anéis
\oldpage[II]{218}
\[
  \Gamma(S,\sh{O}_S)\to\Gamma(X,\sh{O}_X)
\]
no conjunto dos morfismos de espaços anelados
$(\psi,\theta):(X,\sh{O}_X)\to(S,\sh{O}_S)$ tais que, para todo
$x\in X$, $\theta_x^\sharp:\sh{O}_{\psi(x)}\to\sh{O}_x$ seja um
homomorfismo local.
\end{proposition}

\begin{proof}
Observemos primeiro que, se $(X,\sh{O}_X)$ e $(S,\sh{O}_S)$ são dois
espaços anelados quaisquer, um morfismo $(\psi,\theta)$ de
$(X,\sh{O}_X)$ em $(S,\sh{O}_S)$ define canonicamente um homomorfismo
de anéis
\[
  \Gamma(\theta):\Gamma(S,\sh{O}_S)\to\Gamma(X,\sh{O}_X),
\]
e, portanto, uma primeira aplicação
\begin{equation}
\label{ega2:addendum:I.1.8.1.1}
  \rho:\Hom((X,\sh{O}_X),(S,\sh{O}_S))
  \to\Hom(\Gamma(S,\sh{O}_S),\Gamma(X,\sh{O}_X)).
\tag{1.8.1.1}
\end{equation}
Reciprocamente, sob as hipóteses da proposição, ponhamos
$A=\Gamma(S,\sh{O}_S)$ e consideremos um homomorfismo de anéis
$\vphi:A\to\Gamma(X,\sh{O}_X)$. Para todo $x\in X$, o conjunto dos
$f\in A$ tais que $\vphi(f)(x)=0$ é um ideal primo de $A$, pois
$\sh{O}_x/\mathfrak{m}_x=\kres(x)$ é um corpo. Logo, é um elemento de
$S=\Spec(A)$, que também denotaremos por ${}^a\vphi(x)$. Além disso,
para todo $f\in A$, por definição (0, 5.5.2),
\[
  {}^a\vphi^{-1}(D(f))=X_f,
\]
o que prova que ${}^a\vphi$ é uma aplicação contínua $X\to S$.

Definamos agora um homomorfismo de $\sh{O}_S$-módulos
\[
  \widetilde{\vphi}:\sh{O}_S\to{}^a\vphi_*(\sh{O}_X).
\]
Para todo $f\in A$, tem-se
$\Gamma(D(f),\sh{O}_S)=A_f$ (I, 1.3.6); a $s/f\in A_f$ associamos
\[
  (\vphi(s)|X_f)(\vphi(f)|X_f)^{-1}
  \in\Gamma(X_f,\sh{O}_X)
  =\Gamma(D(f),{}^a\vphi_*(\sh{O}_X)).
\]
A passagem de $D(f)$ a $D(fg)$ mostra imediatamente que isto define
um homomorfismo de feixes e, portanto, um morfismo de espaços anelados
$({}^a\vphi,\widetilde{\vphi})$. Com a mesma notação, escrevendo
$y={}^a\vphi(x)$, tem-se (0, 3.7.1)
\[
  \widetilde{\vphi}_x^\sharp(s_y/f_y)
  =(\vphi(s)_x)(\vphi(f)_x)^{-1}.
\]
Como $s_y\in\mathfrak{m}_y$ equivale, por definição, a
$\vphi(s)_x\in\mathfrak{m}_x$, o homomorfismo
$\widetilde{\vphi}_x^\sharp:\sh{O}_y\to\sh{O}_x$ é local. Definimos,
assim, uma segunda aplicação
\[
  \sigma:\Hom(\Gamma(S,\sh{O}_S),\Gamma(X,\sh{O}_X))
  \to\mathfrak{L},
\]
onde $\mathfrak{L}$ denota o conjunto dos morfismos
$(\psi,\theta):(X,\sh{O}_X)\to(S,\sh{O}_S)$ para os quais
$\theta_x^\sharp$ é local para todo $x\in X$.

Resta provar que $\sigma$ e a restrição de $\rho$ a $\mathfrak{L}$ são
inversas. A definição de $\widetilde{\vphi}$ dá imediatamente
$\Gamma(\widetilde{\vphi})=\vphi$, e, portanto, $\rho\circ\sigma$ é a
identidade. Para ver que $\sigma\circ\rho$ é a identidade, partamos de
$(\psi,\theta)\in\mathfrak{L}$ e ponhamos $\vphi=\Gamma(\theta)$. A
hipótese sobre $\theta_x^\sharp$ fornece, passando aos corpos
residuais, um monomorfismo
\[
  \theta^x:\kres(\psi(x))\to\kres(x)
\]
tal que, para todo $f\in A=\Gamma(S,\sh{O}_S)$,
\[
  \theta^x(f(\psi(x)))=\vphi(f)(x).
\]
Assim, $f(\psi(x))=0$ equivale a $\vphi(f)(x)=0$, o que prova que
${}^a\vphi=\psi$. Por outro lado, as definições tornam comutativo o
diagrama
\[
  \xymatrix{
    A\ar[r]^\vphi\ar[d] &
    \Gamma(X,\sh{O}_X)\ar[d]\\
    A_{\psi(x)}\ar[r]^{\theta_x^\sharp} &
    \sh{O}_x
  }
\]
e também é comutativo o diagrama análogo em que $\theta_x^\sharp$ é
substituído por $\widetilde{\vphi}_x^\sharp$. De (0, 1.2.4), segue que
$\widetilde{\vphi}_x^\sharp=\theta_x^\sharp$ e, consequentemente,
$\widetilde{\vphi}=\theta$.
\end{proof}

\begin{env}[1.8.2]
\label{ega2:addendum:I.1.8.2}
Quando $(X,\sh{O}_X)$ e $(Y,\sh{O}_Y)$ são espaços localmente
anelados, consideraremos os morfismos
$(\psi,\theta):(X,\sh{O}_X)\to(Y,\sh{O}_Y)$ tais que, para todo
$x\in X$, $\theta_x^\sharp:\sh{O}_{\psi(x)}\to\sh{O}_x$ seja um
homomorfismo local.

Sempre que, doravante, falarmos de um
\oldpage[II]{219}
\emph{morfismo de espaços localmente anelados}, entenderemos sempre
um morfismo desse tipo. Com esta definição dos morfismos, os espaços
localmente anelados formam uma categoria. Para dois objetos $X,Y$
dessa categoria, $\Hom(X,Y)$ denotará, portanto, o conjunto dos
morfismos de espaços localmente anelados de $X$ em $Y$ (o conjunto
denotado por $\mathfrak{L}$ em \sref{I.1.8.1}). Quando quisermos
considerar o conjunto dos morfismos de espaços anelados de $X$ em
$Y$, escreveremos $\Hom_{\mathrm{rs}}(X,Y)$ para evitar qualquer
ambiguidade. A aplicação \sref{I.1.8.1.1} escreve-se, assim,
\begin{equation}
\label{ega2:addendum:I.1.8.2.1}
  \rho:\Hom_{\mathrm{rs}}(X,Y)
  \to\Hom(\Gamma(Y,\sh{O}_Y),\Gamma(X,\sh{O}_X)),
\tag{1.8.2.1}
\end{equation}
e sua restrição
\begin{equation}
\label{ega2:addendum:I.1.8.2.2}
  \rho':\Hom(X,Y)
  \to\Hom(\Gamma(Y,\sh{O}_Y),\Gamma(X,\sh{O}_X))
\tag{1.8.2.2}
\end{equation}
é funtorial em $X$ e $Y$ na categoria dos espaços localmente
anelados.
\end{env}

\begin{corollary}[1.8.3]
\label{ega2:addendum:I.1.8.3}
Seja $(Y,\sh{O}_Y)$ um espaço localmente anelado. Para que $Y$ seja um
esquema afim, é necessário e suficiente que, para todo espaço
localmente anelado $(X,\sh{O}_X)$, a aplicação \sref{I.1.8.2.2} seja
bijetiva.
\end{corollary}

\begin{proof}
A Proposição~\sref{I.1.8.1} mostra que a condição é necessária.
Reciprocamente, suponha que ela valha e ponha $A=\Gamma(Y,\sh{O}_Y)$.
A hipótese e \sref{I.1.8.1} mostram que os funtores
\[
  X\longmapsto\Hom(X,Y)
  \quad\hbox{e}\quad
  X\longmapsto\Hom(X,\Spec(A))
\]
dos espaços localmente anelados para os conjuntos são isomorfos. Isso
implica a existência de um isomorfismo canônico
$Y\to\Spec(A)$ (cf.~0, 8).
\end{proof}

\begin{env}[1.8.4]
\label{ega2:addendum:I.1.8.4}
Seja $S=\Spec(A)$ um esquema afim. Denote por $(S',A')$ o espaço
anelado cujo espaço subjacente consiste em um ponto e cujo feixe
estrutural $A'$ é o feixe (necessariamente simples) sobre $S'$ definido
pelo anel $A$. Seja $\pi:S\to S'$ a única aplicação. Para todo aberto
$U$ de $S$, a aplicação canônica
\[
  \Gamma(S',A')=\Gamma(S,\sh{O}_S)\to\Gamma(U,\sh{O}_S)
\]
define um $\pi$-morfismo $\iota:A'\to\sh{O}_S$ de feixes de anéis.
Obtém-se, assim, um morfismo canônico de espaços anelados
\[
  i=(\pi,\iota):(S,\sh{O}_S)\to(S',A').
\]
Para todo $A$-módulo $M$, denote por $M'$ o feixe simples sobre $S'$
definido por $M$, que é evidentemente um $A'$-módulo. É claro que
$i_*(\widetilde M)=M'$ (I, 1.3.7).
\end{env}

\begin{lemma}[1.8.5]
\label{ega2:addendum:I.1.8.5}
Com a notação de \sref{I.1.8.4}, para todo $A$-módulo $M$, o
$\sh{O}_S$-homomorfismo canônico e funtorial (0, 4.4.3.3)
\begin{equation}
\label{ega2:addendum:I.1.8.5.1}
  i^*(i_*(\widetilde M))\to\widetilde M
\tag{1.8.5.1}
\end{equation}
é um isomorfismo.
\end{lemma}

\begin{proof}
Ambos os membros de \sref{I.1.8.5.1} são exatos à direita em $M$ (o
funtor $M\mapsto i_*(\widetilde M)$ sendo evidentemente exato) e
comutam com somas diretas. Considerando $M$ como o conúcleo de um
homomorfismo $A^{(I)}\to A^{(J)}$, reduzimo-nos ao caso $M=A$, no qual
a afirmação é imediata.
\end{proof}

\begin{corollary}[1.8.6]
\label{ega2:addendum:I.1.8.6}
Seja $(X,\sh{O}_X)$ um espaço anelado e seja $u:X\to S$ um morfismo de
\oldpage[II]{220}
espaços anelados. Para todo $A$-módulo $M$, há, com a notação de
\sref{I.1.8.4}, um isomorfismo canônico e funtorial de
$\sh{O}_X$-módulos
\begin{equation}
\label{ega2:addendum:I.1.8.6.1}
  u^*(\widetilde M)\isoto u^*(i^*(M')).
\tag{1.8.6.1}
\end{equation}
\end{corollary}

\begin{corollary}[1.8.7]
\label{ega2:addendum:I.1.8.7}
Sob as hipóteses de \sref{I.1.8.6}, para todo $A$-módulo $M$ e todo
$\sh{O}_X$-módulo $\sh{F}$, há um isomorfismo canônico, funtorial em
$M$ e $\sh{F}$,
\begin{equation}
\label{ega2:addendum:I.1.8.7.1}
  \Hom_{\sh{O}_S}(\widetilde M,u_*(\sh{F}))
  \isoto\Hom_A(M,\Gamma(X,\sh{F})).
\tag{1.8.7.1}
\end{equation}
\end{corollary}

\begin{proof}
Por (0, 4.4.3) e pelo Lema~\sref{I.1.8.5}, há um isomorfismo canônico
de bifuntores
\[
  \Hom_{\sh{O}_S}(\widetilde M,u_*(\sh{F}))
  \isoto\Hom_{A'}(M',i_*(u_*(\sh{F}))).
\]
O segundo membro é claramente $\Hom_A(M,\Gamma(X,\sh{F}))$. Sob o
homomorfismo canônico \sref{I.1.8.7.1}, a um
$\sh{O}_S$-homomorfismo
$h:\widetilde M\to u_*(\sh{F})$ (equivalentemente, a um $u$-morfismo
$\widetilde M\to\sh{F}$) corresponde
\[
  \Gamma(h):M\to\Gamma(S,u_*(\sh{F}))=\Gamma(X,\sh{F}).
\]
\end{proof}

\begin{env}[1.8.8]
\label{ega2:addendum:I.1.8.8}
Com a notação de \sref{I.1.8.4}, decorre claramente de (0, 4.1.1) que
todo morfismo de espaços anelados $(\psi,\theta):X\to S'$ equivale ao
dado de um homomorfismo de anéis $A\to\Gamma(X,\sh{O}_X)$. Assim, a
Proposição~\sref{I.1.8.1} pode ser interpretada como uma bijeção
canônica
\[
  \Hom(X,S)\isoto\Hom(X,S'),
\]
onde o segundo membro significa morfismos de espaços anelados, pois
$A$ não é, em geral, um anel local.

Mais geralmente, sejam $X,Y$ espaços localmente anelados, e seja
$(Y',A')$ o espaço anelado cujo espaço subjacente consiste em um ponto
e cujo feixe de anéis é o feixe simples definido por
$\Gamma(Y,\sh{O}_Y)$. Então \sref{I.1.8.2.1} pode ser interpretada
como uma aplicação
\begin{equation}
\label{ega2:addendum:I.1.8.8.1}
  \rho:\Hom_{\mathrm{rs}}(X,Y)\to\Hom(X,Y').
\tag{1.8.8.1}
\end{equation}
O Corolário~\sref{I.1.8.3} diz que os esquemas afins se caracterizam,
entre os espaços localmente anelados, como aqueles para os quais a
restrição de $\rho$ a $\Hom(X,Y)$,
\begin{equation}
\label{ega2:addendum:I.1.8.8.2}
  \rho':\Hom(X,Y)\to\Hom(X,Y'),
\tag{1.8.8.2}
\end{equation}
é bijetiva para todo espaço localmente anelado $X$. Em um capítulo
posterior, generalizaremos esta definição, associando assim a um
espaço anelado arbitrário $Z$ (e não apenas a um cujo espaço
subjacente consiste em um ponto) um espaço localmente anelado que
voltaremos a chamar $\Spec(Z)$. Esse será o ponto de partida de uma
teoria «relativa» dos pré-esquemas sobre espaços anelados arbitrários,
estendendo os resultados do Capítulo~I.
\end{env}

\begin{env}[1.8.9]
\label{ega2:addendum:I.1.8.9}
Os pares $(X,\sh{F})$, em que $X$ é um espaço localmente anelado e
$\sh{F}$ um $\sh{O}_X$-módulo, formam uma categoria: um morfismo
$(X,\sh{F})\to(Y,\sh{G})$ é um par $(u,h)$ formado por um morfismo de
\oldpage[II]{221}
espaços localmente anelados $u:X\to Y$ e de um $u$-morfismo
$h:\sh{G}\to\sh{F}$ de módulos; esses
morfismos (para $(X,\sh{F})$ e $(Y,\sh{G})$ fixados) formam um conjunto
que denotaremos por $\Hom((X,\sh{F}),(Y,\sh{G}))$; a aplicação
$(u,h)\mapsto(\rho'(u),\Gamma(h))$ é uma aplicação canônica
\begin{equation}
\label{ega2:addendum:I.1.8.9.1}
\begin{split}
  \Hom((X,\sh{F}),(Y,\sh{G}))\to
  \Hom\bigl(
    &(\Gamma(Y,\sh{O}_Y),\Gamma(Y,\sh{G})),\\
    &(\Gamma(X,\sh{O}_X),\Gamma(X,\sh{F}))
  \bigr),
\end{split}
\tag{1.8.9.1}
\end{equation}
funtorial em $(X,\sh{F})$ e $(Y,\sh{G})$, sendo o segundo membro
o conjunto dos di-homomorfismos correspondentes aos anéis e aos
módulos considerados (0, 1.0.2).
\end{env}

\begin{corollary}[1.8.10]
\label{ega2:addendum:I.1.8.10}
Sejam $Y$ um espaço localmente anelado e $\sh{G}$ um
$\sh{O}_Y$-módulo. Para que $Y$ seja um esquema afim e $\sh{G}$ um
$\sh{O}_Y$-módulo quase-coerente, é necessário e suficiente que, para
todo par $(X,\sh{F})$ formado por um espaço localmente anelado $X$ e
um $\sh{O}_X$-módulo $\sh{F}$, a aplicação canônica
\sref{I.1.8.9.1} seja bijetiva.
\end{corollary}

Deixamos ao leitor a tarefa de desenvolver o raciocínio, modelado no
de \sref{I.1.8.3}, no qual se utilizam \sref{I.1.8.1} e
\sref{I.1.8.7}.

\begin{remark}[1.8.11]
\label{ega2:addendum:I.1.8.11}
Os enunciados (I, 1.7.3), (I, 1.7.4) e (I, 2.2.4), bem como a
definição (I, 1.6.1), são casos particulares de \sref{I.1.8.1}; do
mesmo modo, (I, 2.2.5) decorre de \sref{I.1.8.7}. O
Corolário~\sref{I.1.8.7} também fornece (I, 1.6.3) (e, portanto,
(I, 1.6.4)) como caso particular. Com efeito, se $X$ é afim e
$\Gamma(X,\sh{F})=N$, os funtores
\[
  M\longmapsto
  \Hom_{\sh{O}_S}(\widetilde M,u_*(\widetilde N))
  \quad\hbox{e}\quad
  M\longmapsto
  \Hom_{\sh{O}_S}(\widetilde M,(N_{[\vphi]})\supertilde)
\]
(onde $\vphi:A\to\Gamma(X,\sh{O}_X)$ corresponde a $u$) são
isomorfos por \sref{I.1.8.7} e (I, 1.3.8). Finalmente, (I, 1.6.5) (e,
portanto, (I, 1.6.6)) decorre de \sref{I.1.8.6} e do fato de que, para
todo $f\in A$, os $A_f$-módulos
\[
  N'\otimes_{A'}A_f
  \quad\hbox{e}\quad
  (N'\otimes_{A'}A)_f
\]
(com as notações de (I, 1.6.5)) são canonicamente isomorfos.
\end{remark}

\paragraph{Inserir depois de (I, 3.2.8):}

\begin{env}[3.2.9]
\label{ega2:addendum:I.3.2.9}
Decorre de \sref{I.1.8.1} que (I, 3.2.2) pode ser precisado da seguinte
forma. O pré-esquema
\[
  Z=\Spec(B\otimes_A C)
\]
não é apenas um produto de $X=\Spec(B)$ e $Y=\Spec(C)$ na categoria
dos $S$-pré-esquemas, mas também na categoria dos espaços localmente
anelados sobre $S$ (com a definição dos $S$-morfismos modelada em
(I, 2.5.2)). A demonstração de (I, 3.2.6) mostra ainda que, para
$S$-pré-esquemas quaisquer $X,Y$, o pré-esquema $X\times_S Y$ não é
apenas o produto deles na categoria dos $S$-pré-esquemas, mas também
na categoria dos espaços localmente anelados sobre o pré-esquema $S$.
\end{env}

\paragraph{(I, 4.4.5).}
Na linha 14 a contar de baixo da p.~126, substituir $B$ por $A$; na
linha 13 a contar de baixo, substituir «$A$-álgebra» por
«$B$-álgebra».

\paragraph{(I, 5.3.5).}
Na linha 2 a contar de baixo da p.~132, substituir $Y\to S$ por
$g:Y\to S$.

\paragraph{(I, 6.3.2.1).}
Na linha 6 a contar de baixo da p.~144, ler:
\[
  D(g_i)\subset W.
\]

\paragraph{(I, 7.3.8).}
Substituir o texto atual, que está errado, pelo seguinte:

Sejam $X,Y$ dois pré-esquemas integrais; então $\sh{R}(X)$ e
$\sh{R}(Y)$ são, respectivamente, um $\sh{O}_X$-módulo quase-coerente
e um $\sh{O}_Y$-módulo quase-coerente (I, 7.3.3). Se $f:X\to Y$ é um
morfismo dominante, existe um homomorfismo canônico de
$\sh{O}_X$-módulos
\oldpage[II]{222}
\begin{equation}
\label{ega2:addendum:I.7.3.8.1}
  \tau:f^*(\sh{R}(Y))\to\sh{R}(X).
\tag{7.3.8.1}
\end{equation}

Suponhamos primeiro que $X=\Spec(A)$ e $Y=\Spec(B)$ sejam afins, com
$A,B$ domínios de integridade. Então $f$ corresponde a um
homomorfismo injetivo $B\to A$, que se estende a um monomorfismo
$L\to K$ do corpo de frações $L$ de $B$ no corpo de frações $K$ de
$A$. O homomorfismo \sref{I.7.3.8.1} corresponde então ao
homomorfismo canônico
\[
  L\otimes_B A\to K
\]
de (I, 1.6.5).

No caso geral, para todo par de abertos afins não vazios
$U\subset X$, $V\subset Y$ tais que $f(U)\subset V$, a construção
anterior define um homomorfismo $\tau_{U,V}$. Se $U'\subset U$,
$V'\subset V$ e $f(U')\subset V'$, então $\tau_{U,V}$ prolonga
$\tau_{U',V'}$; a afirmação segue por colagem. Se $x,y$ são os pontos
genéricos de $X,Y$, respectivamente, então $f(x)=y$ e
\[
  \bigl(f^*(\sh{R}(Y))\bigr)_x
  =\sh{O}_y\otimes_{\sh{O}_y}\sh{O}_x
  =\sh{O}_x
\]
por (0, 4.3.1), de modo que $\tau_x$ é um isomorfismo.

\paragraph{(I, 9.5.2).}
Nas três últimas linhas da p.~176 e nas cinco primeiras da p.~177,
substituir em toda parte $X'$ por $Y'$ e $X''$ por $Y''$.

\paragraph{(I, 10.5.4).}
Na linha 11 a contar de baixo da p.~187, substituir (10.3.4) por
(10.3.5).

\paragraph{(I, 10.6.3).}
Na linha 3 a contar de baixo da p.~189, substituir $X$ por
$\mathfrak X$.

\paragraph{(I, 10.14.2).}
Na linha 5 da p.~210, substituir «$\sh{O}_{\mathfrak X}$-módulo
coerente» por «ideal coerente de $\sh{O}_{\mathfrak X}$»; na linha 7,
substituir «$\sh{O}_{\mathfrak X}$-módulos coerentes» por «ideais
coerentes de $\sh{O}_{\mathfrak X}$».
\end{document}
